\documentclass{article}
\usepackage[paperwidth=16cm,paperheight=14cm,margin=0cm]{geometry}
\usepackage[utf8]{inputenc}\usepackage[T1]{fontenc}
\usepackage{mathpazo}\usepackage{amsmath,amssymb}
\usepackage[most]{tcolorbox}
\pagestyle{empty}\setlength{\parindent}{0pt}

\definecolor{crel}{HTML}{1F6FA8}
\definecolor{cpui}{HTML}{6A4C93}
\definecolor{clit}{HTML}{2E7D6B}
\definecolor{cgeo}{HTML}{C0503A}
\definecolor{cdon}{HTML}{C77F0A}
\definecolor{cgra}{HTML}{8A5A2B}
\begin{document}
\begin{tcolorbox}[enhanced,width=16cm,height=14cm,colback=white,colframe=crel,boxrule=2pt,arc=0mm,
left=6mm,right=6mm,top=0mm,bottom=4mm,sharp corners,
overlay={\fill[crel] (frame.north west) rectangle ([yshift=-1.25cm]frame.north east);
\node[anchor=west,text=white,font=\sffamily\bfseries\LARGE] at ([xshift=5mm,yshift=-0.62cm]frame.north west) {Question 1};
\node[anchor=east,text=white,font=\sffamily\itshape\large] at ([xshift=-5mm,yshift=-0.62cm]frame.north east) {Relatifs et fractions};}]
\vspace*{1.75cm}
\begin{minipage}[t][6.9cm][c]{\linewidth}\LARGE\centering
\begin{minipage}{0.95\linewidth}\raggedright\linespread{1.25}\selectfont Que vaut $(-2)\times(+3)\times(-4)$~?\end{minipage}
\end{minipage}

\vfill
\begin{tcolorbox}[enhanced,nobeforeafter,width=7.1cm,height=2.1cm,valign=center,colback=crel!6,colframe=crel,arc=2.5mm,boxrule=0.9pt,left=3mm]{\Large\bfseries\color{crel}a)}\quad{\Large $24$}\end{tcolorbox}\hfill\begin{tcolorbox}[enhanced,nobeforeafter,width=7.1cm,height=2.1cm,valign=center,colback=crel!6,colframe=crel,arc=2.5mm,boxrule=0.9pt,left=3mm]{\Large\bfseries\color{crel}b)}\quad{\Large $-3$}\end{tcolorbox}\par\vspace{2.5mm}
\begin{tcolorbox}[enhanced,nobeforeafter,width=7.1cm,height=2.1cm,valign=center,colback=crel!6,colframe=crel,arc=2.5mm,boxrule=0.9pt,left=3mm]{\Large\bfseries\color{crel}c)}\quad{\Large $-2$}\end{tcolorbox}\hfill\begin{tcolorbox}[enhanced,nobeforeafter,width=7.1cm,height=2.1cm,valign=center,colback=crel!6,colframe=crel,arc=2.5mm,boxrule=0.9pt,left=3mm]{\Large\bfseries\color{crel}d)}\quad{\Large $-24$}\end{tcolorbox}
\end{tcolorbox}
\newpage
\begin{tcolorbox}[enhanced,width=16cm,height=14cm,colback=white,colframe=crel,boxrule=2pt,arc=0mm,
left=6mm,right=6mm,top=0mm,bottom=4mm,sharp corners,
overlay={\fill[crel] (frame.north west) rectangle ([yshift=-1.25cm]frame.north east);
\node[anchor=west,text=white,font=\sffamily\bfseries\LARGE] at ([xshift=5mm,yshift=-0.62cm]frame.north west) {Question 2};
\node[anchor=east,text=white,font=\sffamily\itshape\large] at ([xshift=-5mm,yshift=-0.62cm]frame.north east) {Relatifs et fractions};}]
\vspace*{1.75cm}
\begin{minipage}[t][6.9cm][c]{\linewidth}\LARGE\centering
\begin{minipage}{0.95\linewidth}\raggedright\linespread{1.25}\selectfont Que vaut $(+5)\times(-2)\times(+3)$~?\end{minipage}
\end{minipage}

\vfill
\begin{tcolorbox}[enhanced,nobeforeafter,width=7.1cm,height=2.1cm,valign=center,colback=crel!6,colframe=crel,arc=2.5mm,boxrule=0.9pt,left=3mm]{\Large\bfseries\color{crel}a)}\quad{\Large $6$}\end{tcolorbox}\hfill\begin{tcolorbox}[enhanced,nobeforeafter,width=7.1cm,height=2.1cm,valign=center,colback=crel!6,colframe=crel,arc=2.5mm,boxrule=0.9pt,left=3mm]{\Large\bfseries\color{crel}b)}\quad{\Large $-30$}\end{tcolorbox}\par\vspace{2.5mm}
\begin{tcolorbox}[enhanced,nobeforeafter,width=7.1cm,height=2.1cm,valign=center,colback=crel!6,colframe=crel,arc=2.5mm,boxrule=0.9pt,left=3mm]{\Large\bfseries\color{crel}c)}\quad{\Large $30$}\end{tcolorbox}\hfill\begin{tcolorbox}[enhanced,nobeforeafter,width=7.1cm,height=2.1cm,valign=center,colback=crel!6,colframe=crel,arc=2.5mm,boxrule=0.9pt,left=3mm]{\Large\bfseries\color{crel}d)}\quad{\Large $-13$}\end{tcolorbox}
\end{tcolorbox}
\newpage
\begin{tcolorbox}[enhanced,width=16cm,height=14cm,colback=white,colframe=crel,boxrule=2pt,arc=0mm,
left=6mm,right=6mm,top=0mm,bottom=4mm,sharp corners,
overlay={\fill[crel] (frame.north west) rectangle ([yshift=-1.25cm]frame.north east);
\node[anchor=west,text=white,font=\sffamily\bfseries\LARGE] at ([xshift=5mm,yshift=-0.62cm]frame.north west) {Question 3};
\node[anchor=east,text=white,font=\sffamily\itshape\large] at ([xshift=-5mm,yshift=-0.62cm]frame.north east) {Relatifs et fractions};}]
\vspace*{1.75cm}
\begin{minipage}[t][6.9cm][c]{\linewidth}\LARGE\centering
\begin{minipage}{0.95\linewidth}\raggedright\linespread{1.25}\selectfont Que vaut $(-36)\div(+4)$~?\end{minipage}
\end{minipage}

\vfill
\begin{tcolorbox}[enhanced,nobeforeafter,width=7.1cm,height=2.1cm,valign=center,colback=crel!6,colframe=crel,arc=2.5mm,boxrule=0.9pt,left=3mm]{\Large\bfseries\color{crel}a)}\quad{\Large $9$}\end{tcolorbox}\hfill\begin{tcolorbox}[enhanced,nobeforeafter,width=7.1cm,height=2.1cm,valign=center,colback=crel!6,colframe=crel,arc=2.5mm,boxrule=0.9pt,left=3mm]{\Large\bfseries\color{crel}b)}\quad{\Large $-8$}\end{tcolorbox}\par\vspace{2.5mm}
\begin{tcolorbox}[enhanced,nobeforeafter,width=7.1cm,height=2.1cm,valign=center,colback=crel!6,colframe=crel,arc=2.5mm,boxrule=0.9pt,left=3mm]{\Large\bfseries\color{crel}c)}\quad{\Large $-18$}\end{tcolorbox}\hfill\begin{tcolorbox}[enhanced,nobeforeafter,width=7.1cm,height=2.1cm,valign=center,colback=crel!6,colframe=crel,arc=2.5mm,boxrule=0.9pt,left=3mm]{\Large\bfseries\color{crel}d)}\quad{\Large $-9$}\end{tcolorbox}
\end{tcolorbox}
\newpage
\begin{tcolorbox}[enhanced,width=16cm,height=14cm,colback=white,colframe=crel,boxrule=2pt,arc=0mm,
left=6mm,right=6mm,top=0mm,bottom=4mm,sharp corners,
overlay={\fill[crel] (frame.north west) rectangle ([yshift=-1.25cm]frame.north east);
\node[anchor=west,text=white,font=\sffamily\bfseries\LARGE] at ([xshift=5mm,yshift=-0.62cm]frame.north west) {Question 4};
\node[anchor=east,text=white,font=\sffamily\itshape\large] at ([xshift=-5mm,yshift=-0.62cm]frame.north east) {Relatifs et fractions};}]
\vspace*{1.75cm}
\begin{minipage}[t][6.9cm][c]{\linewidth}\LARGE\centering
\begin{minipage}{0.95\linewidth}\raggedright\linespread{1.25}\selectfont Que vaut $(+42)\div(-6)$~?\end{minipage}
\end{minipage}

\vfill
\begin{tcolorbox}[enhanced,nobeforeafter,width=7.1cm,height=2.1cm,valign=center,colback=crel!6,colframe=crel,arc=2.5mm,boxrule=0.9pt,left=3mm]{\Large\bfseries\color{crel}a)}\quad{\Large $7$}\end{tcolorbox}\hfill\begin{tcolorbox}[enhanced,nobeforeafter,width=7.1cm,height=2.1cm,valign=center,colback=crel!6,colframe=crel,arc=2.5mm,boxrule=0.9pt,left=3mm]{\Large\bfseries\color{crel}b)}\quad{\Large $-14$}\end{tcolorbox}\par\vspace{2.5mm}
\begin{tcolorbox}[enhanced,nobeforeafter,width=7.1cm,height=2.1cm,valign=center,colback=crel!6,colframe=crel,arc=2.5mm,boxrule=0.9pt,left=3mm]{\Large\bfseries\color{crel}c)}\quad{\Large $-6$}\end{tcolorbox}\hfill\begin{tcolorbox}[enhanced,nobeforeafter,width=7.1cm,height=2.1cm,valign=center,colback=crel!6,colframe=crel,arc=2.5mm,boxrule=0.9pt,left=3mm]{\Large\bfseries\color{crel}d)}\quad{\Large $-7$}\end{tcolorbox}
\end{tcolorbox}
\newpage
\begin{tcolorbox}[enhanced,width=16cm,height=14cm,colback=white,colframe=crel,boxrule=2pt,arc=0mm,
left=6mm,right=6mm,top=0mm,bottom=4mm,sharp corners,
overlay={\fill[crel] (frame.north west) rectangle ([yshift=-1.25cm]frame.north east);
\node[anchor=west,text=white,font=\sffamily\bfseries\LARGE] at ([xshift=5mm,yshift=-0.62cm]frame.north west) {Question 5};
\node[anchor=east,text=white,font=\sffamily\itshape\large] at ([xshift=-5mm,yshift=-0.62cm]frame.north east) {Relatifs et fractions};}]
\vspace*{1.75cm}
\begin{minipage}[t][6.9cm][c]{\linewidth}\LARGE\centering
\begin{minipage}{0.95\linewidth}\raggedright\linespread{1.25}\selectfont Que vaut $5-3\times(-2)$~?\end{minipage}
\end{minipage}

\vfill
\begin{tcolorbox}[enhanced,nobeforeafter,width=7.1cm,height=2.1cm,valign=center,colback=crel!6,colframe=crel,arc=2.5mm,boxrule=0.9pt,left=3mm]{\Large\bfseries\color{crel}a)}\quad{\Large $-11$}\end{tcolorbox}\hfill\begin{tcolorbox}[enhanced,nobeforeafter,width=7.1cm,height=2.1cm,valign=center,colback=crel!6,colframe=crel,arc=2.5mm,boxrule=0.9pt,left=3mm]{\Large\bfseries\color{crel}b)}\quad{\Large $-4$}\end{tcolorbox}\par\vspace{2.5mm}
\begin{tcolorbox}[enhanced,nobeforeafter,width=7.1cm,height=2.1cm,valign=center,colback=crel!6,colframe=crel,arc=2.5mm,boxrule=0.9pt,left=3mm]{\Large\bfseries\color{crel}c)}\quad{\Large $-1$}\end{tcolorbox}\hfill\begin{tcolorbox}[enhanced,nobeforeafter,width=7.1cm,height=2.1cm,valign=center,colback=crel!6,colframe=crel,arc=2.5mm,boxrule=0.9pt,left=3mm]{\Large\bfseries\color{crel}d)}\quad{\Large $11$}\end{tcolorbox}
\end{tcolorbox}
\newpage
\begin{tcolorbox}[enhanced,width=16cm,height=14cm,colback=white,colframe=crel,boxrule=2pt,arc=0mm,
left=6mm,right=6mm,top=0mm,bottom=4mm,sharp corners,
overlay={\fill[crel] (frame.north west) rectangle ([yshift=-1.25cm]frame.north east);
\node[anchor=west,text=white,font=\sffamily\bfseries\LARGE] at ([xshift=5mm,yshift=-0.62cm]frame.north west) {Question 6};
\node[anchor=east,text=white,font=\sffamily\itshape\large] at ([xshift=-5mm,yshift=-0.62cm]frame.north east) {Relatifs et fractions};}]
\vspace*{1.75cm}
\begin{minipage}[t][6.9cm][c]{\linewidth}\LARGE\centering
\begin{minipage}{0.95\linewidth}\raggedright\linespread{1.25}\selectfont Que vaut $7-4\times(-3)$~?\end{minipage}
\end{minipage}

\vfill
\begin{tcolorbox}[enhanced,nobeforeafter,width=7.1cm,height=2.1cm,valign=center,colback=crel!6,colframe=crel,arc=2.5mm,boxrule=0.9pt,left=3mm]{\Large\bfseries\color{crel}a)}\quad{\Large $-5$}\end{tcolorbox}\hfill\begin{tcolorbox}[enhanced,nobeforeafter,width=7.1cm,height=2.1cm,valign=center,colback=crel!6,colframe=crel,arc=2.5mm,boxrule=0.9pt,left=3mm]{\Large\bfseries\color{crel}b)}\quad{\Large $-9$}\end{tcolorbox}\par\vspace{2.5mm}
\begin{tcolorbox}[enhanced,nobeforeafter,width=7.1cm,height=2.1cm,valign=center,colback=crel!6,colframe=crel,arc=2.5mm,boxrule=0.9pt,left=3mm]{\Large\bfseries\color{crel}c)}\quad{\Large $19$}\end{tcolorbox}\hfill\begin{tcolorbox}[enhanced,nobeforeafter,width=7.1cm,height=2.1cm,valign=center,colback=crel!6,colframe=crel,arc=2.5mm,boxrule=0.9pt,left=3mm]{\Large\bfseries\color{crel}d)}\quad{\Large $-19$}\end{tcolorbox}
\end{tcolorbox}
\newpage
\begin{tcolorbox}[enhanced,width=16cm,height=14cm,colback=white,colframe=crel,boxrule=2pt,arc=0mm,
left=6mm,right=6mm,top=0mm,bottom=4mm,sharp corners,
overlay={\fill[crel] (frame.north west) rectangle ([yshift=-1.25cm]frame.north east);
\node[anchor=west,text=white,font=\sffamily\bfseries\LARGE] at ([xshift=5mm,yshift=-0.62cm]frame.north west) {Question 7};
\node[anchor=east,text=white,font=\sffamily\itshape\large] at ([xshift=-5mm,yshift=-0.62cm]frame.north east) {Relatifs et fractions};}]
\vspace*{1.75cm}
\begin{minipage}[t][6.9cm][c]{\linewidth}\LARGE\centering
\begin{minipage}{0.95\linewidth}\raggedright\linespread{1.25}\selectfont Que vaut $\dfrac{1}{3}+\dfrac{1}{4}$~?\end{minipage}
\end{minipage}

\vfill
\begin{tcolorbox}[enhanced,nobeforeafter,width=7.1cm,height=2.1cm,valign=center,colback=crel!6,colframe=crel,arc=2.5mm,boxrule=0.9pt,left=3mm]{\Large\bfseries\color{crel}a)}\quad{\Large $\dfrac{1}{12}$}\end{tcolorbox}\hfill\begin{tcolorbox}[enhanced,nobeforeafter,width=7.1cm,height=2.1cm,valign=center,colback=crel!6,colframe=crel,arc=2.5mm,boxrule=0.9pt,left=3mm]{\Large\bfseries\color{crel}b)}\quad{\Large $\dfrac{1}{6}$}\end{tcolorbox}\par\vspace{2.5mm}
\begin{tcolorbox}[enhanced,nobeforeafter,width=7.1cm,height=2.1cm,valign=center,colback=crel!6,colframe=crel,arc=2.5mm,boxrule=0.9pt,left=3mm]{\Large\bfseries\color{crel}c)}\quad{\Large $\dfrac{2}{7}$}\end{tcolorbox}\hfill\begin{tcolorbox}[enhanced,nobeforeafter,width=7.1cm,height=2.1cm,valign=center,colback=crel!6,colframe=crel,arc=2.5mm,boxrule=0.9pt,left=3mm]{\Large\bfseries\color{crel}d)}\quad{\Large $\dfrac{7}{12}$}\end{tcolorbox}
\end{tcolorbox}
\newpage
\begin{tcolorbox}[enhanced,width=16cm,height=14cm,colback=white,colframe=crel,boxrule=2pt,arc=0mm,
left=6mm,right=6mm,top=0mm,bottom=4mm,sharp corners,
overlay={\fill[crel] (frame.north west) rectangle ([yshift=-1.25cm]frame.north east);
\node[anchor=west,text=white,font=\sffamily\bfseries\LARGE] at ([xshift=5mm,yshift=-0.62cm]frame.north west) {Question 8};
\node[anchor=east,text=white,font=\sffamily\itshape\large] at ([xshift=-5mm,yshift=-0.62cm]frame.north east) {Relatifs et fractions};}]
\vspace*{1.75cm}
\begin{minipage}[t][6.9cm][c]{\linewidth}\LARGE\centering
\begin{minipage}{0.95\linewidth}\raggedright\linespread{1.25}\selectfont Que vaut $\dfrac{2}{5}+\dfrac{1}{2}$~?\end{minipage}
\end{minipage}

\vfill
\begin{tcolorbox}[enhanced,nobeforeafter,width=7.1cm,height=2.1cm,valign=center,colback=crel!6,colframe=crel,arc=2.5mm,boxrule=0.9pt,left=3mm]{\Large\bfseries\color{crel}a)}\quad{\Large $\dfrac{1}{5}$}\end{tcolorbox}\hfill\begin{tcolorbox}[enhanced,nobeforeafter,width=7.1cm,height=2.1cm,valign=center,colback=crel!6,colframe=crel,arc=2.5mm,boxrule=0.9pt,left=3mm]{\Large\bfseries\color{crel}b)}\quad{\Large $\dfrac{3}{7}$}\end{tcolorbox}\par\vspace{2.5mm}
\begin{tcolorbox}[enhanced,nobeforeafter,width=7.1cm,height=2.1cm,valign=center,colback=crel!6,colframe=crel,arc=2.5mm,boxrule=0.9pt,left=3mm]{\Large\bfseries\color{crel}c)}\quad{\Large $\dfrac{3}{10}$}\end{tcolorbox}\hfill\begin{tcolorbox}[enhanced,nobeforeafter,width=7.1cm,height=2.1cm,valign=center,colback=crel!6,colframe=crel,arc=2.5mm,boxrule=0.9pt,left=3mm]{\Large\bfseries\color{crel}d)}\quad{\Large $\dfrac{9}{10}$}\end{tcolorbox}
\end{tcolorbox}
\newpage
\begin{tcolorbox}[enhanced,width=16cm,height=14cm,colback=white,colframe=crel,boxrule=2pt,arc=0mm,
left=6mm,right=6mm,top=0mm,bottom=4mm,sharp corners,
overlay={\fill[crel] (frame.north west) rectangle ([yshift=-1.25cm]frame.north east);
\node[anchor=west,text=white,font=\sffamily\bfseries\LARGE] at ([xshift=5mm,yshift=-0.62cm]frame.north west) {Question 9};
\node[anchor=east,text=white,font=\sffamily\itshape\large] at ([xshift=-5mm,yshift=-0.62cm]frame.north east) {Relatifs et fractions};}]
\vspace*{1.75cm}
\begin{minipage}[t][6.9cm][c]{\linewidth}\LARGE\centering
\begin{minipage}{0.95\linewidth}\raggedright\linespread{1.25}\selectfont Que vaut $\dfrac{2}{3}\times\dfrac{9}{4}$~?\end{minipage}
\end{minipage}

\vfill
\begin{tcolorbox}[enhanced,nobeforeafter,width=7.1cm,height=2.1cm,valign=center,colback=crel!6,colframe=crel,arc=2.5mm,boxrule=0.9pt,left=3mm]{\Large\bfseries\color{crel}a)}\quad{\Large $\dfrac{8}{27}$}\end{tcolorbox}\hfill\begin{tcolorbox}[enhanced,nobeforeafter,width=7.1cm,height=2.1cm,valign=center,colback=crel!6,colframe=crel,arc=2.5mm,boxrule=0.9pt,left=3mm]{\Large\bfseries\color{crel}b)}\quad{\Large $-\dfrac{3}{2}$}\end{tcolorbox}\par\vspace{2.5mm}
\begin{tcolorbox}[enhanced,nobeforeafter,width=7.1cm,height=2.1cm,valign=center,colback=crel!6,colframe=crel,arc=2.5mm,boxrule=0.9pt,left=3mm]{\Large\bfseries\color{crel}c)}\quad{\Large $\dfrac{3}{2}$}\end{tcolorbox}\hfill\begin{tcolorbox}[enhanced,nobeforeafter,width=7.1cm,height=2.1cm,valign=center,colback=crel!6,colframe=crel,arc=2.5mm,boxrule=0.9pt,left=3mm]{\Large\bfseries\color{crel}d)}\quad{\Large $\dfrac{11}{7}$}\end{tcolorbox}
\end{tcolorbox}
\newpage
\begin{tcolorbox}[enhanced,width=16cm,height=14cm,colback=white,colframe=crel,boxrule=2pt,arc=0mm,
left=6mm,right=6mm,top=0mm,bottom=4mm,sharp corners,
overlay={\fill[crel] (frame.north west) rectangle ([yshift=-1.25cm]frame.north east);
\node[anchor=west,text=white,font=\sffamily\bfseries\LARGE] at ([xshift=5mm,yshift=-0.62cm]frame.north west) {Question 10};
\node[anchor=east,text=white,font=\sffamily\itshape\large] at ([xshift=-5mm,yshift=-0.62cm]frame.north east) {Relatifs et fractions};}]
\vspace*{1.75cm}
\begin{minipage}[t][6.9cm][c]{\linewidth}\LARGE\centering
\begin{minipage}{0.95\linewidth}\raggedright\linespread{1.25}\selectfont Que vaut $\left(-\dfrac{3}{5}\right)\times\dfrac{10}{9}$~?\end{minipage}
\end{minipage}

\vfill
\begin{tcolorbox}[enhanced,nobeforeafter,width=7.1cm,height=2.1cm,valign=center,colback=crel!6,colframe=crel,arc=2.5mm,boxrule=0.9pt,left=3mm]{\Large\bfseries\color{crel}a)}\quad{\Large $\dfrac{2}{3}$}\end{tcolorbox}\hfill\begin{tcolorbox}[enhanced,nobeforeafter,width=7.1cm,height=2.1cm,valign=center,colback=crel!6,colframe=crel,arc=2.5mm,boxrule=0.9pt,left=3mm]{\Large\bfseries\color{crel}b)}\quad{\Large $\dfrac{1}{2}$}\end{tcolorbox}\par\vspace{2.5mm}
\begin{tcolorbox}[enhanced,nobeforeafter,width=7.1cm,height=2.1cm,valign=center,colback=crel!6,colframe=crel,arc=2.5mm,boxrule=0.9pt,left=3mm]{\Large\bfseries\color{crel}c)}\quad{\Large $-\dfrac{27}{50}$}\end{tcolorbox}\hfill\begin{tcolorbox}[enhanced,nobeforeafter,width=7.1cm,height=2.1cm,valign=center,colback=crel!6,colframe=crel,arc=2.5mm,boxrule=0.9pt,left=3mm]{\Large\bfseries\color{crel}d)}\quad{\Large $-\dfrac{2}{3}$}\end{tcolorbox}
\end{tcolorbox}
\newpage
\begin{tcolorbox}[enhanced,width=16cm,height=14cm,colback=white,colframe=crel,boxrule=2pt,arc=0mm,
left=6mm,right=6mm,top=0mm,bottom=4mm,sharp corners,
overlay={\fill[crel] (frame.north west) rectangle ([yshift=-1.25cm]frame.north east);
\node[anchor=west,text=white,font=\sffamily\bfseries\LARGE] at ([xshift=5mm,yshift=-0.62cm]frame.north west) {Question 11};
\node[anchor=east,text=white,font=\sffamily\itshape\large] at ([xshift=-5mm,yshift=-0.62cm]frame.north east) {Relatifs et fractions};}]
\vspace*{1.75cm}
\begin{minipage}[t][6.9cm][c]{\linewidth}\LARGE\centering
\begin{minipage}{0.95\linewidth}\raggedright\linespread{1.25}\selectfont Que vaut $\dfrac{2}{3}\div\dfrac{4}{5}$~?\end{minipage}
\end{minipage}

\vfill
\begin{tcolorbox}[enhanced,nobeforeafter,width=7.1cm,height=2.1cm,valign=center,colback=crel!6,colframe=crel,arc=2.5mm,boxrule=0.9pt,left=3mm]{\Large\bfseries\color{crel}a)}\quad{\Large $\dfrac{5}{6}$}\end{tcolorbox}\hfill\begin{tcolorbox}[enhanced,nobeforeafter,width=7.1cm,height=2.1cm,valign=center,colback=crel!6,colframe=crel,arc=2.5mm,boxrule=0.9pt,left=3mm]{\Large\bfseries\color{crel}b)}\quad{\Large $\dfrac{6}{5}$}\end{tcolorbox}\par\vspace{2.5mm}
\begin{tcolorbox}[enhanced,nobeforeafter,width=7.1cm,height=2.1cm,valign=center,colback=crel!6,colframe=crel,arc=2.5mm,boxrule=0.9pt,left=3mm]{\Large\bfseries\color{crel}c)}\quad{\Large $-\dfrac{5}{6}$}\end{tcolorbox}\hfill\begin{tcolorbox}[enhanced,nobeforeafter,width=7.1cm,height=2.1cm,valign=center,colback=crel!6,colframe=crel,arc=2.5mm,boxrule=0.9pt,left=3mm]{\Large\bfseries\color{crel}d)}\quad{\Large $\dfrac{8}{15}$}\end{tcolorbox}
\end{tcolorbox}
\newpage
\begin{tcolorbox}[enhanced,width=16cm,height=14cm,colback=white,colframe=crel,boxrule=2pt,arc=0mm,
left=6mm,right=6mm,top=0mm,bottom=4mm,sharp corners,
overlay={\fill[crel] (frame.north west) rectangle ([yshift=-1.25cm]frame.north east);
\node[anchor=west,text=white,font=\sffamily\bfseries\LARGE] at ([xshift=5mm,yshift=-0.62cm]frame.north west) {Question 12};
\node[anchor=east,text=white,font=\sffamily\itshape\large] at ([xshift=-5mm,yshift=-0.62cm]frame.north east) {Relatifs et fractions};}]
\vspace*{1.75cm}
\begin{minipage}[t][6.9cm][c]{\linewidth}\LARGE\centering
\begin{minipage}{0.95\linewidth}\raggedright\linespread{1.25}\selectfont Que vaut $\dfrac{3}{4}\div\dfrac{9}{8}$~?\end{minipage}
\end{minipage}

\vfill
\begin{tcolorbox}[enhanced,nobeforeafter,width=7.1cm,height=2.1cm,valign=center,colback=crel!6,colframe=crel,arc=2.5mm,boxrule=0.9pt,left=3mm]{\Large\bfseries\color{crel}a)}\quad{\Large $\dfrac{27}{32}$}\end{tcolorbox}\hfill\begin{tcolorbox}[enhanced,nobeforeafter,width=7.1cm,height=2.1cm,valign=center,colback=crel!6,colframe=crel,arc=2.5mm,boxrule=0.9pt,left=3mm]{\Large\bfseries\color{crel}b)}\quad{\Large $\dfrac{2}{3}$}\end{tcolorbox}\par\vspace{2.5mm}
\begin{tcolorbox}[enhanced,nobeforeafter,width=7.1cm,height=2.1cm,valign=center,colback=crel!6,colframe=crel,arc=2.5mm,boxrule=0.9pt,left=3mm]{\Large\bfseries\color{crel}c)}\quad{\Large $-\dfrac{2}{3}$}\end{tcolorbox}\hfill\begin{tcolorbox}[enhanced,nobeforeafter,width=7.1cm,height=2.1cm,valign=center,colback=crel!6,colframe=crel,arc=2.5mm,boxrule=0.9pt,left=3mm]{\Large\bfseries\color{crel}d)}\quad{\Large $\dfrac{3}{2}$}\end{tcolorbox}
\end{tcolorbox}
\newpage
\begin{tcolorbox}[enhanced,width=16cm,height=14cm,colback=white,colframe=crel,boxrule=2pt,arc=0mm,
left=6mm,right=6mm,top=0mm,bottom=4mm,sharp corners,
overlay={\fill[crel] (frame.north west) rectangle ([yshift=-1.25cm]frame.north east);
\node[anchor=west,text=white,font=\sffamily\bfseries\LARGE] at ([xshift=5mm,yshift=-0.62cm]frame.north west) {Question 13};
\node[anchor=east,text=white,font=\sffamily\itshape\large] at ([xshift=-5mm,yshift=-0.62cm]frame.north east) {Relatifs et fractions};}]
\vspace*{1.75cm}
\begin{minipage}[t][6.9cm][c]{\linewidth}\LARGE\centering
\begin{minipage}{0.95\linewidth}\raggedright\linespread{1.25}\selectfont Quel est l'inverse de $\dfrac{3}{4}$~?\end{minipage}
\end{minipage}

\vfill
\begin{tcolorbox}[enhanced,nobeforeafter,width=7.1cm,height=2.1cm,valign=center,colback=crel!6,colframe=crel,arc=2.5mm,boxrule=0.9pt,left=3mm]{\Large\bfseries\color{crel}a)}\quad{\Large $\dfrac{3}{4}$}\end{tcolorbox}\hfill\begin{tcolorbox}[enhanced,nobeforeafter,width=7.1cm,height=2.1cm,valign=center,colback=crel!6,colframe=crel,arc=2.5mm,boxrule=0.9pt,left=3mm]{\Large\bfseries\color{crel}b)}\quad{\Large $-\dfrac{4}{3}$}\end{tcolorbox}\par\vspace{2.5mm}
\begin{tcolorbox}[enhanced,nobeforeafter,width=7.1cm,height=2.1cm,valign=center,colback=crel!6,colframe=crel,arc=2.5mm,boxrule=0.9pt,left=3mm]{\Large\bfseries\color{crel}c)}\quad{\Large $\dfrac{4}{3}$}\end{tcolorbox}\hfill\begin{tcolorbox}[enhanced,nobeforeafter,width=7.1cm,height=2.1cm,valign=center,colback=crel!6,colframe=crel,arc=2.5mm,boxrule=0.9pt,left=3mm]{\Large\bfseries\color{crel}d)}\quad{\Large $-\dfrac{3}{4}$}\end{tcolorbox}
\end{tcolorbox}
\newpage
\begin{tcolorbox}[enhanced,width=16cm,height=14cm,colback=white,colframe=crel,boxrule=2pt,arc=0mm,
left=6mm,right=6mm,top=0mm,bottom=4mm,sharp corners,
overlay={\fill[crel] (frame.north west) rectangle ([yshift=-1.25cm]frame.north east);
\node[anchor=west,text=white,font=\sffamily\bfseries\LARGE] at ([xshift=5mm,yshift=-0.62cm]frame.north west) {Question 14};
\node[anchor=east,text=white,font=\sffamily\itshape\large] at ([xshift=-5mm,yshift=-0.62cm]frame.north east) {Relatifs et fractions};}]
\vspace*{1.75cm}
\begin{minipage}[t][6.9cm][c]{\linewidth}\LARGE\centering
\begin{minipage}{0.95\linewidth}\raggedright\linespread{1.25}\selectfont Quel est l'inverse de $-\dfrac{2}{5}$~?\end{minipage}
\end{minipage}

\vfill
\begin{tcolorbox}[enhanced,nobeforeafter,width=7.1cm,height=2.1cm,valign=center,colback=crel!6,colframe=crel,arc=2.5mm,boxrule=0.9pt,left=3mm]{\Large\bfseries\color{crel}a)}\quad{\Large $-\dfrac{2}{5}$}\end{tcolorbox}\hfill\begin{tcolorbox}[enhanced,nobeforeafter,width=7.1cm,height=2.1cm,valign=center,colback=crel!6,colframe=crel,arc=2.5mm,boxrule=0.9pt,left=3mm]{\Large\bfseries\color{crel}b)}\quad{\Large $\dfrac{2}{5}$}\end{tcolorbox}\par\vspace{2.5mm}
\begin{tcolorbox}[enhanced,nobeforeafter,width=7.1cm,height=2.1cm,valign=center,colback=crel!6,colframe=crel,arc=2.5mm,boxrule=0.9pt,left=3mm]{\Large\bfseries\color{crel}c)}\quad{\Large $-\dfrac{5}{2}$}\end{tcolorbox}\hfill\begin{tcolorbox}[enhanced,nobeforeafter,width=7.1cm,height=2.1cm,valign=center,colback=crel!6,colframe=crel,arc=2.5mm,boxrule=0.9pt,left=3mm]{\Large\bfseries\color{crel}d)}\quad{\Large $\dfrac{5}{2}$}\end{tcolorbox}
\end{tcolorbox}
\newpage
\begin{tcolorbox}[enhanced,width=16cm,height=14cm,colback=white,colframe=crel,boxrule=2pt,arc=0mm,
left=6mm,right=6mm,top=0mm,bottom=4mm,sharp corners,
overlay={\fill[crel] (frame.north west) rectangle ([yshift=-1.25cm]frame.north east);
\node[anchor=west,text=white,font=\sffamily\bfseries\LARGE] at ([xshift=5mm,yshift=-0.62cm]frame.north west) {Question 15};
\node[anchor=east,text=white,font=\sffamily\itshape\large] at ([xshift=-5mm,yshift=-0.62cm]frame.north east) {Relatifs et fractions};}]
\vspace*{1.75cm}
\begin{minipage}[t][6.9cm][c]{\linewidth}\LARGE\centering
\begin{minipage}{0.95\linewidth}\raggedright\linespread{1.25}\selectfont Que vaut $\dfrac{1}{4}-\dfrac{2}{3}$~?\end{minipage}
\end{minipage}

\vfill
\begin{tcolorbox}[enhanced,nobeforeafter,width=7.1cm,height=2.1cm,valign=center,colback=crel!6,colframe=crel,arc=2.5mm,boxrule=0.9pt,left=3mm]{\Large\bfseries\color{crel}a)}\quad{\Large $\dfrac{5}{12}$}\end{tcolorbox}\hfill\begin{tcolorbox}[enhanced,nobeforeafter,width=7.1cm,height=2.1cm,valign=center,colback=crel!6,colframe=crel,arc=2.5mm,boxrule=0.9pt,left=3mm]{\Large\bfseries\color{crel}b)}\quad{\Large $-1$}\end{tcolorbox}\par\vspace{2.5mm}
\begin{tcolorbox}[enhanced,nobeforeafter,width=7.1cm,height=2.1cm,valign=center,colback=crel!6,colframe=crel,arc=2.5mm,boxrule=0.9pt,left=3mm]{\Large\bfseries\color{crel}c)}\quad{\Large $-\dfrac{1}{12}$}\end{tcolorbox}\hfill\begin{tcolorbox}[enhanced,nobeforeafter,width=7.1cm,height=2.1cm,valign=center,colback=crel!6,colframe=crel,arc=2.5mm,boxrule=0.9pt,left=3mm]{\Large\bfseries\color{crel}d)}\quad{\Large $-\dfrac{5}{12}$}\end{tcolorbox}
\end{tcolorbox}
\newpage
\begin{tcolorbox}[enhanced,width=16cm,height=14cm,colback=white,colframe=crel,boxrule=2pt,arc=0mm,
left=6mm,right=6mm,top=0mm,bottom=4mm,sharp corners,
overlay={\fill[crel] (frame.north west) rectangle ([yshift=-1.25cm]frame.north east);
\node[anchor=west,text=white,font=\sffamily\bfseries\LARGE] at ([xshift=5mm,yshift=-0.62cm]frame.north west) {Question 16};
\node[anchor=east,text=white,font=\sffamily\itshape\large] at ([xshift=-5mm,yshift=-0.62cm]frame.north east) {Relatifs et fractions};}]
\vspace*{1.75cm}
\begin{minipage}[t][6.9cm][c]{\linewidth}\LARGE\centering
\begin{minipage}{0.95\linewidth}\raggedright\linespread{1.25}\selectfont Que vaut $\dfrac{3}{5}-\dfrac{7}{10}$~?\end{minipage}
\end{minipage}

\vfill
\begin{tcolorbox}[enhanced,nobeforeafter,width=7.1cm,height=2.1cm,valign=center,colback=crel!6,colframe=crel,arc=2.5mm,boxrule=0.9pt,left=3mm]{\Large\bfseries\color{crel}a)}\quad{\Large $\dfrac{1}{10}$}\end{tcolorbox}\hfill\begin{tcolorbox}[enhanced,nobeforeafter,width=7.1cm,height=2.1cm,valign=center,colback=crel!6,colframe=crel,arc=2.5mm,boxrule=0.9pt,left=3mm]{\Large\bfseries\color{crel}b)}\quad{\Large $-\dfrac{4}{5}$}\end{tcolorbox}\par\vspace{2.5mm}
\begin{tcolorbox}[enhanced,nobeforeafter,width=7.1cm,height=2.1cm,valign=center,colback=crel!6,colframe=crel,arc=2.5mm,boxrule=0.9pt,left=3mm]{\Large\bfseries\color{crel}c)}\quad{\Large $-\dfrac{1}{10}$}\end{tcolorbox}\hfill\begin{tcolorbox}[enhanced,nobeforeafter,width=7.1cm,height=2.1cm,valign=center,colback=crel!6,colframe=crel,arc=2.5mm,boxrule=0.9pt,left=3mm]{\Large\bfseries\color{crel}d)}\quad{\Large $-\dfrac{2}{25}$}\end{tcolorbox}
\end{tcolorbox}
\newpage
\begin{tcolorbox}[enhanced,width=16cm,height=14cm,colback=white,colframe=cpui,boxrule=2pt,arc=0mm,
left=6mm,right=6mm,top=0mm,bottom=4mm,sharp corners,
overlay={\fill[cpui] (frame.north west) rectangle ([yshift=-1.25cm]frame.north east);
\node[anchor=west,text=white,font=\sffamily\bfseries\LARGE] at ([xshift=5mm,yshift=-0.62cm]frame.north west) {Question 17};
\node[anchor=east,text=white,font=\sffamily\itshape\large] at ([xshift=-5mm,yshift=-0.62cm]frame.north east) {Puissances et racines carrées};}]
\vspace*{1.75cm}
\begin{minipage}[t][6.9cm][c]{\linewidth}\LARGE\centering
\begin{minipage}{0.95\linewidth}\raggedright\linespread{1.25}\selectfont Que vaut $2^{5}$~?\end{minipage}
\end{minipage}

\vfill
\begin{tcolorbox}[enhanced,nobeforeafter,width=7.1cm,height=2.1cm,valign=center,colback=cpui!6,colframe=cpui,arc=2.5mm,boxrule=0.9pt,left=3mm]{\Large\bfseries\color{cpui}a)}\quad{\Large $16$}\end{tcolorbox}\hfill\begin{tcolorbox}[enhanced,nobeforeafter,width=7.1cm,height=2.1cm,valign=center,colback=cpui!6,colframe=cpui,arc=2.5mm,boxrule=0.9pt,left=3mm]{\Large\bfseries\color{cpui}b)}\quad{\Large $32$}\end{tcolorbox}\par\vspace{2.5mm}
\begin{tcolorbox}[enhanced,nobeforeafter,width=7.1cm,height=2.1cm,valign=center,colback=cpui!6,colframe=cpui,arc=2.5mm,boxrule=0.9pt,left=3mm]{\Large\bfseries\color{cpui}c)}\quad{\Large $10$}\end{tcolorbox}\hfill\begin{tcolorbox}[enhanced,nobeforeafter,width=7.1cm,height=2.1cm,valign=center,colback=cpui!6,colframe=cpui,arc=2.5mm,boxrule=0.9pt,left=3mm]{\Large\bfseries\color{cpui}d)}\quad{\Large $-32$}\end{tcolorbox}
\end{tcolorbox}
\newpage
\begin{tcolorbox}[enhanced,width=16cm,height=14cm,colback=white,colframe=cpui,boxrule=2pt,arc=0mm,
left=6mm,right=6mm,top=0mm,bottom=4mm,sharp corners,
overlay={\fill[cpui] (frame.north west) rectangle ([yshift=-1.25cm]frame.north east);
\node[anchor=west,text=white,font=\sffamily\bfseries\LARGE] at ([xshift=5mm,yshift=-0.62cm]frame.north west) {Question 18};
\node[anchor=east,text=white,font=\sffamily\itshape\large] at ([xshift=-5mm,yshift=-0.62cm]frame.north east) {Puissances et racines carrées};}]
\vspace*{1.75cm}
\begin{minipage}[t][6.9cm][c]{\linewidth}\LARGE\centering
\begin{minipage}{0.95\linewidth}\raggedright\linespread{1.25}\selectfont Que vaut $3^{4}$~?\end{minipage}
\end{minipage}

\vfill
\begin{tcolorbox}[enhanced,nobeforeafter,width=7.1cm,height=2.1cm,valign=center,colback=cpui!6,colframe=cpui,arc=2.5mm,boxrule=0.9pt,left=3mm]{\Large\bfseries\color{cpui}a)}\quad{\Large $27$}\end{tcolorbox}\hfill\begin{tcolorbox}[enhanced,nobeforeafter,width=7.1cm,height=2.1cm,valign=center,colback=cpui!6,colframe=cpui,arc=2.5mm,boxrule=0.9pt,left=3mm]{\Large\bfseries\color{cpui}b)}\quad{\Large $12$}\end{tcolorbox}\par\vspace{2.5mm}
\begin{tcolorbox}[enhanced,nobeforeafter,width=7.1cm,height=2.1cm,valign=center,colback=cpui!6,colframe=cpui,arc=2.5mm,boxrule=0.9pt,left=3mm]{\Large\bfseries\color{cpui}c)}\quad{\Large $-81$}\end{tcolorbox}\hfill\begin{tcolorbox}[enhanced,nobeforeafter,width=7.1cm,height=2.1cm,valign=center,colback=cpui!6,colframe=cpui,arc=2.5mm,boxrule=0.9pt,left=3mm]{\Large\bfseries\color{cpui}d)}\quad{\Large $81$}\end{tcolorbox}
\end{tcolorbox}
\newpage
\begin{tcolorbox}[enhanced,width=16cm,height=14cm,colback=white,colframe=cpui,boxrule=2pt,arc=0mm,
left=6mm,right=6mm,top=0mm,bottom=4mm,sharp corners,
overlay={\fill[cpui] (frame.north west) rectangle ([yshift=-1.25cm]frame.north east);
\node[anchor=west,text=white,font=\sffamily\bfseries\LARGE] at ([xshift=5mm,yshift=-0.62cm]frame.north west) {Question 19};
\node[anchor=east,text=white,font=\sffamily\itshape\large] at ([xshift=-5mm,yshift=-0.62cm]frame.north east) {Puissances et racines carrées};}]
\vspace*{1.75cm}
\begin{minipage}[t][6.9cm][c]{\linewidth}\LARGE\centering
\begin{minipage}{0.95\linewidth}\raggedright\linespread{1.25}\selectfont Écrire $10^{3}\times10^{4}$ sous la forme d'une seule puissance de $10$.\end{minipage}
\end{minipage}

\vfill
\begin{tcolorbox}[enhanced,nobeforeafter,width=7.1cm,height=2.1cm,valign=center,colback=cpui!6,colframe=cpui,arc=2.5mm,boxrule=0.9pt,left=3mm]{\Large\bfseries\color{cpui}a)}\quad{\Large $100^{7}$}\end{tcolorbox}\hfill\begin{tcolorbox}[enhanced,nobeforeafter,width=7.1cm,height=2.1cm,valign=center,colback=cpui!6,colframe=cpui,arc=2.5mm,boxrule=0.9pt,left=3mm]{\Large\bfseries\color{cpui}b)}\quad{\Large $10^{12}$}\end{tcolorbox}\par\vspace{2.5mm}
\begin{tcolorbox}[enhanced,nobeforeafter,width=7.1cm,height=2.1cm,valign=center,colback=cpui!6,colframe=cpui,arc=2.5mm,boxrule=0.9pt,left=3mm]{\Large\bfseries\color{cpui}c)}\quad{\Large $10^{-1}$}\end{tcolorbox}\hfill\begin{tcolorbox}[enhanced,nobeforeafter,width=7.1cm,height=2.1cm,valign=center,colback=cpui!6,colframe=cpui,arc=2.5mm,boxrule=0.9pt,left=3mm]{\Large\bfseries\color{cpui}d)}\quad{\Large $10^{7}$}\end{tcolorbox}
\end{tcolorbox}
\newpage
\begin{tcolorbox}[enhanced,width=16cm,height=14cm,colback=white,colframe=cpui,boxrule=2pt,arc=0mm,
left=6mm,right=6mm,top=0mm,bottom=4mm,sharp corners,
overlay={\fill[cpui] (frame.north west) rectangle ([yshift=-1.25cm]frame.north east);
\node[anchor=west,text=white,font=\sffamily\bfseries\LARGE] at ([xshift=5mm,yshift=-0.62cm]frame.north west) {Question 20};
\node[anchor=east,text=white,font=\sffamily\itshape\large] at ([xshift=-5mm,yshift=-0.62cm]frame.north east) {Puissances et racines carrées};}]
\vspace*{1.75cm}
\begin{minipage}[t][6.9cm][c]{\linewidth}\LARGE\centering
\begin{minipage}{0.95\linewidth}\raggedright\linespread{1.25}\selectfont Écrire $10^{5}\times10^{-2}$ sous la forme d'une seule puissance de $10$.\end{minipage}
\end{minipage}

\vfill
\begin{tcolorbox}[enhanced,nobeforeafter,width=7.1cm,height=2.1cm,valign=center,colback=cpui!6,colframe=cpui,arc=2.5mm,boxrule=0.9pt,left=3mm]{\Large\bfseries\color{cpui}a)}\quad{\Large $10^{7}$}\end{tcolorbox}\hfill\begin{tcolorbox}[enhanced,nobeforeafter,width=7.1cm,height=2.1cm,valign=center,colback=cpui!6,colframe=cpui,arc=2.5mm,boxrule=0.9pt,left=3mm]{\Large\bfseries\color{cpui}b)}\quad{\Large $10^{3}$}\end{tcolorbox}\par\vspace{2.5mm}
\begin{tcolorbox}[enhanced,nobeforeafter,width=7.1cm,height=2.1cm,valign=center,colback=cpui!6,colframe=cpui,arc=2.5mm,boxrule=0.9pt,left=3mm]{\Large\bfseries\color{cpui}c)}\quad{\Large $10^{-10}$}\end{tcolorbox}\hfill\begin{tcolorbox}[enhanced,nobeforeafter,width=7.1cm,height=2.1cm,valign=center,colback=cpui!6,colframe=cpui,arc=2.5mm,boxrule=0.9pt,left=3mm]{\Large\bfseries\color{cpui}d)}\quad{\Large $100^{3}$}\end{tcolorbox}
\end{tcolorbox}
\newpage
\begin{tcolorbox}[enhanced,width=16cm,height=14cm,colback=white,colframe=cpui,boxrule=2pt,arc=0mm,
left=6mm,right=6mm,top=0mm,bottom=4mm,sharp corners,
overlay={\fill[cpui] (frame.north west) rectangle ([yshift=-1.25cm]frame.north east);
\node[anchor=west,text=white,font=\sffamily\bfseries\LARGE] at ([xshift=5mm,yshift=-0.62cm]frame.north west) {Question 21};
\node[anchor=east,text=white,font=\sffamily\itshape\large] at ([xshift=-5mm,yshift=-0.62cm]frame.north east) {Puissances et racines carrées};}]
\vspace*{1.75cm}
\begin{minipage}[t][6.9cm][c]{\linewidth}\LARGE\centering
\begin{minipage}{0.95\linewidth}\raggedright\linespread{1.25}\selectfont Quelle est l'écriture décimale de $10^{-2}$~?\end{minipage}
\end{minipage}

\vfill
\begin{tcolorbox}[enhanced,nobeforeafter,width=7.1cm,height=2.1cm,valign=center,colback=cpui!6,colframe=cpui,arc=2.5mm,boxrule=0.9pt,left=3mm]{\Large\bfseries\color{cpui}a)}\quad{\Large $0{,}01$}\end{tcolorbox}\hfill\begin{tcolorbox}[enhanced,nobeforeafter,width=7.1cm,height=2.1cm,valign=center,colback=cpui!6,colframe=cpui,arc=2.5mm,boxrule=0.9pt,left=3mm]{\Large\bfseries\color{cpui}b)}\quad{\Large $0{,}1$}\end{tcolorbox}\par\vspace{2.5mm}
\begin{tcolorbox}[enhanced,nobeforeafter,width=7.1cm,height=2.1cm,valign=center,colback=cpui!6,colframe=cpui,arc=2.5mm,boxrule=0.9pt,left=3mm]{\Large\bfseries\color{cpui}c)}\quad{\Large $0{,}001$}\end{tcolorbox}\hfill\begin{tcolorbox}[enhanced,nobeforeafter,width=7.1cm,height=2.1cm,valign=center,colback=cpui!6,colframe=cpui,arc=2.5mm,boxrule=0.9pt,left=3mm]{\Large\bfseries\color{cpui}d)}\quad{\Large $-100$}\end{tcolorbox}
\end{tcolorbox}
\newpage
\begin{tcolorbox}[enhanced,width=16cm,height=14cm,colback=white,colframe=cpui,boxrule=2pt,arc=0mm,
left=6mm,right=6mm,top=0mm,bottom=4mm,sharp corners,
overlay={\fill[cpui] (frame.north west) rectangle ([yshift=-1.25cm]frame.north east);
\node[anchor=west,text=white,font=\sffamily\bfseries\LARGE] at ([xshift=5mm,yshift=-0.62cm]frame.north west) {Question 22};
\node[anchor=east,text=white,font=\sffamily\itshape\large] at ([xshift=-5mm,yshift=-0.62cm]frame.north east) {Puissances et racines carrées};}]
\vspace*{1.75cm}
\begin{minipage}[t][6.9cm][c]{\linewidth}\LARGE\centering
\begin{minipage}{0.95\linewidth}\raggedright\linespread{1.25}\selectfont Quelle est l'écriture décimale de $10^{-3}$~?\end{minipage}
\end{minipage}

\vfill
\begin{tcolorbox}[enhanced,nobeforeafter,width=7.1cm,height=2.1cm,valign=center,colback=cpui!6,colframe=cpui,arc=2.5mm,boxrule=0.9pt,left=3mm]{\Large\bfseries\color{cpui}a)}\quad{\Large $0{,}01$}\end{tcolorbox}\hfill\begin{tcolorbox}[enhanced,nobeforeafter,width=7.1cm,height=2.1cm,valign=center,colback=cpui!6,colframe=cpui,arc=2.5mm,boxrule=0.9pt,left=3mm]{\Large\bfseries\color{cpui}b)}\quad{\Large $0{,}0001$}\end{tcolorbox}\par\vspace{2.5mm}
\begin{tcolorbox}[enhanced,nobeforeafter,width=7.1cm,height=2.1cm,valign=center,colback=cpui!6,colframe=cpui,arc=2.5mm,boxrule=0.9pt,left=3mm]{\Large\bfseries\color{cpui}c)}\quad{\Large $-1\,000$}\end{tcolorbox}\hfill\begin{tcolorbox}[enhanced,nobeforeafter,width=7.1cm,height=2.1cm,valign=center,colback=cpui!6,colframe=cpui,arc=2.5mm,boxrule=0.9pt,left=3mm]{\Large\bfseries\color{cpui}d)}\quad{\Large $0{,}001$}\end{tcolorbox}
\end{tcolorbox}
\newpage
\begin{tcolorbox}[enhanced,width=16cm,height=14cm,colback=white,colframe=cpui,boxrule=2pt,arc=0mm,
left=6mm,right=6mm,top=0mm,bottom=4mm,sharp corners,
overlay={\fill[cpui] (frame.north west) rectangle ([yshift=-1.25cm]frame.north east);
\node[anchor=west,text=white,font=\sffamily\bfseries\LARGE] at ([xshift=5mm,yshift=-0.62cm]frame.north west) {Question 23};
\node[anchor=east,text=white,font=\sffamily\itshape\large] at ([xshift=-5mm,yshift=-0.62cm]frame.north east) {Puissances et racines carrées};}]
\vspace*{1.75cm}
\begin{minipage}[t][6.9cm][c]{\linewidth}\LARGE\centering
\begin{minipage}{0.95\linewidth}\raggedright\linespread{1.25}\selectfont Quelle est la notation scientifique de $3\,400\,000$~?\end{minipage}
\end{minipage}

\vfill
\begin{tcolorbox}[enhanced,nobeforeafter,width=7.1cm,height=2.1cm,valign=center,colback=cpui!6,colframe=cpui,arc=2.5mm,boxrule=0.9pt,left=3mm]{\Large\bfseries\color{cpui}a)}\quad{\Large $3{,}4\times10^{-6}$}\end{tcolorbox}\hfill\begin{tcolorbox}[enhanced,nobeforeafter,width=7.1cm,height=2.1cm,valign=center,colback=cpui!6,colframe=cpui,arc=2.5mm,boxrule=0.9pt,left=3mm]{\Large\bfseries\color{cpui}b)}\quad{\Large $3{,}4\times10^{6}$}\end{tcolorbox}\par\vspace{2.5mm}
\begin{tcolorbox}[enhanced,nobeforeafter,width=7.1cm,height=2.1cm,valign=center,colback=cpui!6,colframe=cpui,arc=2.5mm,boxrule=0.9pt,left=3mm]{\Large\bfseries\color{cpui}c)}\quad{\Large $3{,}4\times10^{7}$}\end{tcolorbox}\hfill\begin{tcolorbox}[enhanced,nobeforeafter,width=7.1cm,height=2.1cm,valign=center,colback=cpui!6,colframe=cpui,arc=2.5mm,boxrule=0.9pt,left=3mm]{\Large\bfseries\color{cpui}d)}\quad{\Large $34\times10^{6}$}\end{tcolorbox}
\end{tcolorbox}
\newpage
\begin{tcolorbox}[enhanced,width=16cm,height=14cm,colback=white,colframe=cpui,boxrule=2pt,arc=0mm,
left=6mm,right=6mm,top=0mm,bottom=4mm,sharp corners,
overlay={\fill[cpui] (frame.north west) rectangle ([yshift=-1.25cm]frame.north east);
\node[anchor=west,text=white,font=\sffamily\bfseries\LARGE] at ([xshift=5mm,yshift=-0.62cm]frame.north west) {Question 24};
\node[anchor=east,text=white,font=\sffamily\itshape\large] at ([xshift=-5mm,yshift=-0.62cm]frame.north east) {Puissances et racines carrées};}]
\vspace*{1.75cm}
\begin{minipage}[t][6.9cm][c]{\linewidth}\LARGE\centering
\begin{minipage}{0.95\linewidth}\raggedright\linespread{1.25}\selectfont Quelle est la notation scientifique de $0{,}00052$~?\end{minipage}
\end{minipage}

\vfill
\begin{tcolorbox}[enhanced,nobeforeafter,width=7.1cm,height=2.1cm,valign=center,colback=cpui!6,colframe=cpui,arc=2.5mm,boxrule=0.9pt,left=3mm]{\Large\bfseries\color{cpui}a)}\quad{\Large $5{,}2\times10^{-3}$}\end{tcolorbox}\hfill\begin{tcolorbox}[enhanced,nobeforeafter,width=7.1cm,height=2.1cm,valign=center,colback=cpui!6,colframe=cpui,arc=2.5mm,boxrule=0.9pt,left=3mm]{\Large\bfseries\color{cpui}b)}\quad{\Large $5{,}2\times10^{-4}$}\end{tcolorbox}\par\vspace{2.5mm}
\begin{tcolorbox}[enhanced,nobeforeafter,width=7.1cm,height=2.1cm,valign=center,colback=cpui!6,colframe=cpui,arc=2.5mm,boxrule=0.9pt,left=3mm]{\Large\bfseries\color{cpui}c)}\quad{\Large $5{,}2\times10^{4}$}\end{tcolorbox}\hfill\begin{tcolorbox}[enhanced,nobeforeafter,width=7.1cm,height=2.1cm,valign=center,colback=cpui!6,colframe=cpui,arc=2.5mm,boxrule=0.9pt,left=3mm]{\Large\bfseries\color{cpui}d)}\quad{\Large $52\times10^{-4}$}\end{tcolorbox}
\end{tcolorbox}
\newpage
\begin{tcolorbox}[enhanced,width=16cm,height=14cm,colback=white,colframe=cpui,boxrule=2pt,arc=0mm,
left=6mm,right=6mm,top=0mm,bottom=4mm,sharp corners,
overlay={\fill[cpui] (frame.north west) rectangle ([yshift=-1.25cm]frame.north east);
\node[anchor=west,text=white,font=\sffamily\bfseries\LARGE] at ([xshift=5mm,yshift=-0.62cm]frame.north west) {Question 25};
\node[anchor=east,text=white,font=\sffamily\itshape\large] at ([xshift=-5mm,yshift=-0.62cm]frame.north east) {Puissances et racines carrées};}]
\vspace*{1.75cm}
\begin{minipage}[t][6.9cm][c]{\linewidth}\LARGE\centering
\begin{minipage}{0.95\linewidth}\raggedright\linespread{1.25}\selectfont Que vaut $-3^{2}$~?\end{minipage}
\end{minipage}

\vfill
\begin{tcolorbox}[enhanced,nobeforeafter,width=7.1cm,height=2.1cm,valign=center,colback=cpui!6,colframe=cpui,arc=2.5mm,boxrule=0.9pt,left=3mm]{\Large\bfseries\color{cpui}a)}\quad{\Large $-6$}\end{tcolorbox}\hfill\begin{tcolorbox}[enhanced,nobeforeafter,width=7.1cm,height=2.1cm,valign=center,colback=cpui!6,colframe=cpui,arc=2.5mm,boxrule=0.9pt,left=3mm]{\Large\bfseries\color{cpui}b)}\quad{\Large $9$}\end{tcolorbox}\par\vspace{2.5mm}
\begin{tcolorbox}[enhanced,nobeforeafter,width=7.1cm,height=2.1cm,valign=center,colback=cpui!6,colframe=cpui,arc=2.5mm,boxrule=0.9pt,left=3mm]{\Large\bfseries\color{cpui}c)}\quad{\Large $6$}\end{tcolorbox}\hfill\begin{tcolorbox}[enhanced,nobeforeafter,width=7.1cm,height=2.1cm,valign=center,colback=cpui!6,colframe=cpui,arc=2.5mm,boxrule=0.9pt,left=3mm]{\Large\bfseries\color{cpui}d)}\quad{\Large $-9$}\end{tcolorbox}
\end{tcolorbox}
\newpage
\begin{tcolorbox}[enhanced,width=16cm,height=14cm,colback=white,colframe=cpui,boxrule=2pt,arc=0mm,
left=6mm,right=6mm,top=0mm,bottom=4mm,sharp corners,
overlay={\fill[cpui] (frame.north west) rectangle ([yshift=-1.25cm]frame.north east);
\node[anchor=west,text=white,font=\sffamily\bfseries\LARGE] at ([xshift=5mm,yshift=-0.62cm]frame.north west) {Question 26};
\node[anchor=east,text=white,font=\sffamily\itshape\large] at ([xshift=-5mm,yshift=-0.62cm]frame.north east) {Puissances et racines carrées};}]
\vspace*{1.75cm}
\begin{minipage}[t][6.9cm][c]{\linewidth}\LARGE\centering
\begin{minipage}{0.95\linewidth}\raggedright\linespread{1.25}\selectfont Que vaut $-2^{4}$~?\end{minipage}
\end{minipage}

\vfill
\begin{tcolorbox}[enhanced,nobeforeafter,width=7.1cm,height=2.1cm,valign=center,colback=cpui!6,colframe=cpui,arc=2.5mm,boxrule=0.9pt,left=3mm]{\Large\bfseries\color{cpui}a)}\quad{\Large $16$}\end{tcolorbox}\hfill\begin{tcolorbox}[enhanced,nobeforeafter,width=7.1cm,height=2.1cm,valign=center,colback=cpui!6,colframe=cpui,arc=2.5mm,boxrule=0.9pt,left=3mm]{\Large\bfseries\color{cpui}b)}\quad{\Large $-16$}\end{tcolorbox}\par\vspace{2.5mm}
\begin{tcolorbox}[enhanced,nobeforeafter,width=7.1cm,height=2.1cm,valign=center,colback=cpui!6,colframe=cpui,arc=2.5mm,boxrule=0.9pt,left=3mm]{\Large\bfseries\color{cpui}c)}\quad{\Large $-8$}\end{tcolorbox}\hfill\begin{tcolorbox}[enhanced,nobeforeafter,width=7.1cm,height=2.1cm,valign=center,colback=cpui!6,colframe=cpui,arc=2.5mm,boxrule=0.9pt,left=3mm]{\Large\bfseries\color{cpui}d)}\quad{\Large $8$}\end{tcolorbox}
\end{tcolorbox}
\newpage
\begin{tcolorbox}[enhanced,width=16cm,height=14cm,colback=white,colframe=cpui,boxrule=2pt,arc=0mm,
left=6mm,right=6mm,top=0mm,bottom=4mm,sharp corners,
overlay={\fill[cpui] (frame.north west) rectangle ([yshift=-1.25cm]frame.north east);
\node[anchor=west,text=white,font=\sffamily\bfseries\LARGE] at ([xshift=5mm,yshift=-0.62cm]frame.north west) {Question 27};
\node[anchor=east,text=white,font=\sffamily\itshape\large] at ([xshift=-5mm,yshift=-0.62cm]frame.north east) {Puissances et racines carrées};}]
\vspace*{1.75cm}
\begin{minipage}[t][6.9cm][c]{\linewidth}\LARGE\centering
\begin{minipage}{0.95\linewidth}\raggedright\linespread{1.25}\selectfont Écrire $\dfrac{10^{7}}{10^{3}}$ sous la forme d'une seule puissance de $10$.\end{minipage}
\end{minipage}

\vfill
\begin{tcolorbox}[enhanced,nobeforeafter,width=7.1cm,height=2.1cm,valign=center,colback=cpui!6,colframe=cpui,arc=2.5mm,boxrule=0.9pt,left=3mm]{\Large\bfseries\color{cpui}a)}\quad{\Large $10^{-4}$}\end{tcolorbox}\hfill\begin{tcolorbox}[enhanced,nobeforeafter,width=7.1cm,height=2.1cm,valign=center,colback=cpui!6,colframe=cpui,arc=2.5mm,boxrule=0.9pt,left=3mm]{\Large\bfseries\color{cpui}b)}\quad{\Large $10^{4}$}\end{tcolorbox}\par\vspace{2.5mm}
\begin{tcolorbox}[enhanced,nobeforeafter,width=7.1cm,height=2.1cm,valign=center,colback=cpui!6,colframe=cpui,arc=2.5mm,boxrule=0.9pt,left=3mm]{\Large\bfseries\color{cpui}c)}\quad{\Large $10^{21}$}\end{tcolorbox}\hfill\begin{tcolorbox}[enhanced,nobeforeafter,width=7.1cm,height=2.1cm,valign=center,colback=cpui!6,colframe=cpui,arc=2.5mm,boxrule=0.9pt,left=3mm]{\Large\bfseries\color{cpui}d)}\quad{\Large $10^{10}$}\end{tcolorbox}
\end{tcolorbox}
\newpage
\begin{tcolorbox}[enhanced,width=16cm,height=14cm,colback=white,colframe=cpui,boxrule=2pt,arc=0mm,
left=6mm,right=6mm,top=0mm,bottom=4mm,sharp corners,
overlay={\fill[cpui] (frame.north west) rectangle ([yshift=-1.25cm]frame.north east);
\node[anchor=west,text=white,font=\sffamily\bfseries\LARGE] at ([xshift=5mm,yshift=-0.62cm]frame.north west) {Question 28};
\node[anchor=east,text=white,font=\sffamily\itshape\large] at ([xshift=-5mm,yshift=-0.62cm]frame.north east) {Puissances et racines carrées};}]
\vspace*{1.75cm}
\begin{minipage}[t][6.9cm][c]{\linewidth}\LARGE\centering
\begin{minipage}{0.95\linewidth}\raggedright\linespread{1.25}\selectfont Écrire $\dfrac{10^{5}}{10^{8}}$ sous la forme d'une seule puissance de $10$.\end{minipage}
\end{minipage}

\vfill
\begin{tcolorbox}[enhanced,nobeforeafter,width=7.1cm,height=2.1cm,valign=center,colback=cpui!6,colframe=cpui,arc=2.5mm,boxrule=0.9pt,left=3mm]{\Large\bfseries\color{cpui}a)}\quad{\Large $10^{-3}$}\end{tcolorbox}\hfill\begin{tcolorbox}[enhanced,nobeforeafter,width=7.1cm,height=2.1cm,valign=center,colback=cpui!6,colframe=cpui,arc=2.5mm,boxrule=0.9pt,left=3mm]{\Large\bfseries\color{cpui}b)}\quad{\Large $10^{13}$}\end{tcolorbox}\par\vspace{2.5mm}
\begin{tcolorbox}[enhanced,nobeforeafter,width=7.1cm,height=2.1cm,valign=center,colback=cpui!6,colframe=cpui,arc=2.5mm,boxrule=0.9pt,left=3mm]{\Large\bfseries\color{cpui}c)}\quad{\Large $10^{3}$}\end{tcolorbox}\hfill\begin{tcolorbox}[enhanced,nobeforeafter,width=7.1cm,height=2.1cm,valign=center,colback=cpui!6,colframe=cpui,arc=2.5mm,boxrule=0.9pt,left=3mm]{\Large\bfseries\color{cpui}d)}\quad{\Large $10^{40}$}\end{tcolorbox}
\end{tcolorbox}
\newpage
\begin{tcolorbox}[enhanced,width=16cm,height=14cm,colback=white,colframe=cpui,boxrule=2pt,arc=0mm,
left=6mm,right=6mm,top=0mm,bottom=4mm,sharp corners,
overlay={\fill[cpui] (frame.north west) rectangle ([yshift=-1.25cm]frame.north east);
\node[anchor=west,text=white,font=\sffamily\bfseries\LARGE] at ([xshift=5mm,yshift=-0.62cm]frame.north west) {Question 29};
\node[anchor=east,text=white,font=\sffamily\itshape\large] at ([xshift=-5mm,yshift=-0.62cm]frame.north east) {Puissances et racines carrées};}]
\vspace*{1.75cm}
\begin{minipage}[t][6.9cm][c]{\linewidth}\LARGE\centering
\begin{minipage}{0.95\linewidth}\raggedright\linespread{1.25}\selectfont Que vaut $\sqrt{49}$~?\end{minipage}
\end{minipage}

\vfill
\begin{tcolorbox}[enhanced,nobeforeafter,width=7.1cm,height=2.1cm,valign=center,colback=cpui!6,colframe=cpui,arc=2.5mm,boxrule=0.9pt,left=3mm]{\Large\bfseries\color{cpui}a)}\quad{\Large $24$}\end{tcolorbox}\hfill\begin{tcolorbox}[enhanced,nobeforeafter,width=7.1cm,height=2.1cm,valign=center,colback=cpui!6,colframe=cpui,arc=2.5mm,boxrule=0.9pt,left=3mm]{\Large\bfseries\color{cpui}b)}\quad{\Large $7$}\end{tcolorbox}\par\vspace{2.5mm}
\begin{tcolorbox}[enhanced,nobeforeafter,width=7.1cm,height=2.1cm,valign=center,colback=cpui!6,colframe=cpui,arc=2.5mm,boxrule=0.9pt,left=3mm]{\Large\bfseries\color{cpui}c)}\quad{\Large $8$}\end{tcolorbox}\hfill\begin{tcolorbox}[enhanced,nobeforeafter,width=7.1cm,height=2.1cm,valign=center,colback=cpui!6,colframe=cpui,arc=2.5mm,boxrule=0.9pt,left=3mm]{\Large\bfseries\color{cpui}d)}\quad{\Large $14$}\end{tcolorbox}
\end{tcolorbox}
\newpage
\begin{tcolorbox}[enhanced,width=16cm,height=14cm,colback=white,colframe=cpui,boxrule=2pt,arc=0mm,
left=6mm,right=6mm,top=0mm,bottom=4mm,sharp corners,
overlay={\fill[cpui] (frame.north west) rectangle ([yshift=-1.25cm]frame.north east);
\node[anchor=west,text=white,font=\sffamily\bfseries\LARGE] at ([xshift=5mm,yshift=-0.62cm]frame.north west) {Question 30};
\node[anchor=east,text=white,font=\sffamily\itshape\large] at ([xshift=-5mm,yshift=-0.62cm]frame.north east) {Puissances et racines carrées};}]
\vspace*{1.75cm}
\begin{minipage}[t][6.9cm][c]{\linewidth}\LARGE\centering
\begin{minipage}{0.95\linewidth}\raggedright\linespread{1.25}\selectfont Que vaut $\sqrt{144}$~?\end{minipage}
\end{minipage}

\vfill
\begin{tcolorbox}[enhanced,nobeforeafter,width=7.1cm,height=2.1cm,valign=center,colback=cpui!6,colframe=cpui,arc=2.5mm,boxrule=0.9pt,left=3mm]{\Large\bfseries\color{cpui}a)}\quad{\Large $72$}\end{tcolorbox}\hfill\begin{tcolorbox}[enhanced,nobeforeafter,width=7.1cm,height=2.1cm,valign=center,colback=cpui!6,colframe=cpui,arc=2.5mm,boxrule=0.9pt,left=3mm]{\Large\bfseries\color{cpui}b)}\quad{\Large $13$}\end{tcolorbox}\par\vspace{2.5mm}
\begin{tcolorbox}[enhanced,nobeforeafter,width=7.1cm,height=2.1cm,valign=center,colback=cpui!6,colframe=cpui,arc=2.5mm,boxrule=0.9pt,left=3mm]{\Large\bfseries\color{cpui}c)}\quad{\Large $12$}\end{tcolorbox}\hfill\begin{tcolorbox}[enhanced,nobeforeafter,width=7.1cm,height=2.1cm,valign=center,colback=cpui!6,colframe=cpui,arc=2.5mm,boxrule=0.9pt,left=3mm]{\Large\bfseries\color{cpui}d)}\quad{\Large $24$}\end{tcolorbox}
\end{tcolorbox}
\newpage
\begin{tcolorbox}[enhanced,width=16cm,height=14cm,colback=white,colframe=cpui,boxrule=2pt,arc=0mm,
left=6mm,right=6mm,top=0mm,bottom=4mm,sharp corners,
overlay={\fill[cpui] (frame.north west) rectangle ([yshift=-1.25cm]frame.north east);
\node[anchor=west,text=white,font=\sffamily\bfseries\LARGE] at ([xshift=5mm,yshift=-0.62cm]frame.north west) {Question 31};
\node[anchor=east,text=white,font=\sffamily\itshape\large] at ([xshift=-5mm,yshift=-0.62cm]frame.north east) {Puissances et racines carrées};}]
\vspace*{1.75cm}
\begin{minipage}[t][6.9cm][c]{\linewidth}\LARGE\centering
\begin{minipage}{0.95\linewidth}\raggedright\linespread{1.25}\selectfont Entre quels nombres entiers consécutifs se trouve $\sqrt{20}$~?\end{minipage}
\end{minipage}

\vfill
\begin{tcolorbox}[enhanced,nobeforeafter,width=7.1cm,height=2.1cm,valign=center,colback=cpui!6,colframe=cpui,arc=2.5mm,boxrule=0.9pt,left=3mm]{\Large\bfseries\color{cpui}a)}\quad{\Large $4\text{ et }5$}\end{tcolorbox}\hfill\begin{tcolorbox}[enhanced,nobeforeafter,width=7.1cm,height=2.1cm,valign=center,colback=cpui!6,colframe=cpui,arc=2.5mm,boxrule=0.9pt,left=3mm]{\Large\bfseries\color{cpui}b)}\quad{\Large $10\text{ et }11$}\end{tcolorbox}\par\vspace{2.5mm}
\begin{tcolorbox}[enhanced,nobeforeafter,width=7.1cm,height=2.1cm,valign=center,colback=cpui!6,colframe=cpui,arc=2.5mm,boxrule=0.9pt,left=3mm]{\Large\bfseries\color{cpui}c)}\quad{\Large $5\text{ et }6$}\end{tcolorbox}\hfill\begin{tcolorbox}[enhanced,nobeforeafter,width=7.1cm,height=2.1cm,valign=center,colback=cpui!6,colframe=cpui,arc=2.5mm,boxrule=0.9pt,left=3mm]{\Large\bfseries\color{cpui}d)}\quad{\Large $3\text{ et }4$}\end{tcolorbox}
\end{tcolorbox}
\newpage
\begin{tcolorbox}[enhanced,width=16cm,height=14cm,colback=white,colframe=cpui,boxrule=2pt,arc=0mm,
left=6mm,right=6mm,top=0mm,bottom=4mm,sharp corners,
overlay={\fill[cpui] (frame.north west) rectangle ([yshift=-1.25cm]frame.north east);
\node[anchor=west,text=white,font=\sffamily\bfseries\LARGE] at ([xshift=5mm,yshift=-0.62cm]frame.north west) {Question 32};
\node[anchor=east,text=white,font=\sffamily\itshape\large] at ([xshift=-5mm,yshift=-0.62cm]frame.north east) {Puissances et racines carrées};}]
\vspace*{1.75cm}
\begin{minipage}[t][6.9cm][c]{\linewidth}\LARGE\centering
\begin{minipage}{0.95\linewidth}\raggedright\linespread{1.25}\selectfont Entre quels nombres entiers consécutifs se trouve $\sqrt{50}$~?\end{minipage}
\end{minipage}

\vfill
\begin{tcolorbox}[enhanced,nobeforeafter,width=7.1cm,height=2.1cm,valign=center,colback=cpui!6,colframe=cpui,arc=2.5mm,boxrule=0.9pt,left=3mm]{\Large\bfseries\color{cpui}a)}\quad{\Large $6\text{ et }7$}\end{tcolorbox}\hfill\begin{tcolorbox}[enhanced,nobeforeafter,width=7.1cm,height=2.1cm,valign=center,colback=cpui!6,colframe=cpui,arc=2.5mm,boxrule=0.9pt,left=3mm]{\Large\bfseries\color{cpui}b)}\quad{\Large $8\text{ et }9$}\end{tcolorbox}\par\vspace{2.5mm}
\begin{tcolorbox}[enhanced,nobeforeafter,width=7.1cm,height=2.1cm,valign=center,colback=cpui!6,colframe=cpui,arc=2.5mm,boxrule=0.9pt,left=3mm]{\Large\bfseries\color{cpui}c)}\quad{\Large $25\text{ et }26$}\end{tcolorbox}\hfill\begin{tcolorbox}[enhanced,nobeforeafter,width=7.1cm,height=2.1cm,valign=center,colback=cpui!6,colframe=cpui,arc=2.5mm,boxrule=0.9pt,left=3mm]{\Large\bfseries\color{cpui}d)}\quad{\Large $7\text{ et }8$}\end{tcolorbox}
\end{tcolorbox}
\newpage
\begin{tcolorbox}[enhanced,width=16cm,height=14cm,colback=white,colframe=clit,boxrule=2pt,arc=0mm,
left=6mm,right=6mm,top=0mm,bottom=4mm,sharp corners,
overlay={\fill[clit] (frame.north west) rectangle ([yshift=-1.25cm]frame.north east);
\node[anchor=west,text=white,font=\sffamily\bfseries\LARGE] at ([xshift=5mm,yshift=-0.62cm]frame.north west) {Question 33};
\node[anchor=east,text=white,font=\sffamily\itshape\large] at ([xshift=-5mm,yshift=-0.62cm]frame.north east) {Calcul littéral et équations};}]
\vspace*{1.75cm}
\begin{minipage}[t][6.9cm][c]{\linewidth}\LARGE\centering
\begin{minipage}{0.95\linewidth}\raggedright\linespread{1.25}\selectfont Développer $3(2x+5)$.\end{minipage}
\end{minipage}

\vfill
\begin{tcolorbox}[enhanced,nobeforeafter,width=7.1cm,height=2.1cm,valign=center,colback=clit!6,colframe=clit,arc=2.5mm,boxrule=0.9pt,left=3mm]{\Large\bfseries\color{clit}a)}\quad{\Large $6x+15$}\end{tcolorbox}\hfill\begin{tcolorbox}[enhanced,nobeforeafter,width=7.1cm,height=2.1cm,valign=center,colback=clit!6,colframe=clit,arc=2.5mm,boxrule=0.9pt,left=3mm]{\Large\bfseries\color{clit}b)}\quad{\Large $2x+15$}\end{tcolorbox}\par\vspace{2.5mm}
\begin{tcolorbox}[enhanced,nobeforeafter,width=7.1cm,height=2.1cm,valign=center,colback=clit!6,colframe=clit,arc=2.5mm,boxrule=0.9pt,left=3mm]{\Large\bfseries\color{clit}c)}\quad{\Large $6x+5$}\end{tcolorbox}\hfill\begin{tcolorbox}[enhanced,nobeforeafter,width=7.1cm,height=2.1cm,valign=center,colback=clit!6,colframe=clit,arc=2.5mm,boxrule=0.9pt,left=3mm]{\Large\bfseries\color{clit}d)}\quad{\Large $6x-15$}\end{tcolorbox}
\end{tcolorbox}
\newpage
\begin{tcolorbox}[enhanced,width=16cm,height=14cm,colback=white,colframe=clit,boxrule=2pt,arc=0mm,
left=6mm,right=6mm,top=0mm,bottom=4mm,sharp corners,
overlay={\fill[clit] (frame.north west) rectangle ([yshift=-1.25cm]frame.north east);
\node[anchor=west,text=white,font=\sffamily\bfseries\LARGE] at ([xshift=5mm,yshift=-0.62cm]frame.north west) {Question 34};
\node[anchor=east,text=white,font=\sffamily\itshape\large] at ([xshift=-5mm,yshift=-0.62cm]frame.north east) {Calcul littéral et équations};}]
\vspace*{1.75cm}
\begin{minipage}[t][6.9cm][c]{\linewidth}\LARGE\centering
\begin{minipage}{0.95\linewidth}\raggedright\linespread{1.25}\selectfont Développer $-2(4x-3)$.\end{minipage}
\end{minipage}

\vfill
\begin{tcolorbox}[enhanced,nobeforeafter,width=7.1cm,height=2.1cm,valign=center,colback=clit!6,colframe=clit,arc=2.5mm,boxrule=0.9pt,left=3mm]{\Large\bfseries\color{clit}a)}\quad{\Large $-8x-6$}\end{tcolorbox}\hfill\begin{tcolorbox}[enhanced,nobeforeafter,width=7.1cm,height=2.1cm,valign=center,colback=clit!6,colframe=clit,arc=2.5mm,boxrule=0.9pt,left=3mm]{\Large\bfseries\color{clit}b)}\quad{\Large $-8x+6$}\end{tcolorbox}\par\vspace{2.5mm}
\begin{tcolorbox}[enhanced,nobeforeafter,width=7.1cm,height=2.1cm,valign=center,colback=clit!6,colframe=clit,arc=2.5mm,boxrule=0.9pt,left=3mm]{\Large\bfseries\color{clit}c)}\quad{\Large $4x+6$}\end{tcolorbox}\hfill\begin{tcolorbox}[enhanced,nobeforeafter,width=7.1cm,height=2.1cm,valign=center,colback=clit!6,colframe=clit,arc=2.5mm,boxrule=0.9pt,left=3mm]{\Large\bfseries\color{clit}d)}\quad{\Large $-8x-3$}\end{tcolorbox}
\end{tcolorbox}
\newpage
\begin{tcolorbox}[enhanced,width=16cm,height=14cm,colback=white,colframe=clit,boxrule=2pt,arc=0mm,
left=6mm,right=6mm,top=0mm,bottom=4mm,sharp corners,
overlay={\fill[clit] (frame.north west) rectangle ([yshift=-1.25cm]frame.north east);
\node[anchor=west,text=white,font=\sffamily\bfseries\LARGE] at ([xshift=5mm,yshift=-0.62cm]frame.north west) {Question 35};
\node[anchor=east,text=white,font=\sffamily\itshape\large] at ([xshift=-5mm,yshift=-0.62cm]frame.north east) {Calcul littéral et équations};}]
\vspace*{1.75cm}
\begin{minipage}[t][6.9cm][c]{\linewidth}\LARGE\centering
\begin{minipage}{0.95\linewidth}\raggedright\linespread{1.25}\selectfont Développer et réduire $(x+2)(x+3)$.\end{minipage}
\end{minipage}

\vfill
\begin{tcolorbox}[enhanced,nobeforeafter,width=7.1cm,height=2.1cm,valign=center,colback=clit!6,colframe=clit,arc=2.5mm,boxrule=0.9pt,left=3mm]{\Large\bfseries\color{clit}a)}\quad{\Large $x^2+6x+5$}\end{tcolorbox}\hfill\begin{tcolorbox}[enhanced,nobeforeafter,width=7.1cm,height=2.1cm,valign=center,colback=clit!6,colframe=clit,arc=2.5mm,boxrule=0.9pt,left=3mm]{\Large\bfseries\color{clit}b)}\quad{\Large $x^2+5x+5$}\end{tcolorbox}\par\vspace{2.5mm}
\begin{tcolorbox}[enhanced,nobeforeafter,width=7.1cm,height=2.1cm,valign=center,colback=clit!6,colframe=clit,arc=2.5mm,boxrule=0.9pt,left=3mm]{\Large\bfseries\color{clit}c)}\quad{\Large $x^2+6$}\end{tcolorbox}\hfill\begin{tcolorbox}[enhanced,nobeforeafter,width=7.1cm,height=2.1cm,valign=center,colback=clit!6,colframe=clit,arc=2.5mm,boxrule=0.9pt,left=3mm]{\Large\bfseries\color{clit}d)}\quad{\Large $x^2+5x+6$}\end{tcolorbox}
\end{tcolorbox}
\newpage
\begin{tcolorbox}[enhanced,width=16cm,height=14cm,colback=white,colframe=clit,boxrule=2pt,arc=0mm,
left=6mm,right=6mm,top=0mm,bottom=4mm,sharp corners,
overlay={\fill[clit] (frame.north west) rectangle ([yshift=-1.25cm]frame.north east);
\node[anchor=west,text=white,font=\sffamily\bfseries\LARGE] at ([xshift=5mm,yshift=-0.62cm]frame.north west) {Question 36};
\node[anchor=east,text=white,font=\sffamily\itshape\large] at ([xshift=-5mm,yshift=-0.62cm]frame.north east) {Calcul littéral et équations};}]
\vspace*{1.75cm}
\begin{minipage}[t][6.9cm][c]{\linewidth}\LARGE\centering
\begin{minipage}{0.95\linewidth}\raggedright\linespread{1.25}\selectfont Développer et réduire $(x+4)(x-1)$.\end{minipage}
\end{minipage}

\vfill
\begin{tcolorbox}[enhanced,nobeforeafter,width=7.1cm,height=2.1cm,valign=center,colback=clit!6,colframe=clit,arc=2.5mm,boxrule=0.9pt,left=3mm]{\Large\bfseries\color{clit}a)}\quad{\Large $x^2+3x+3$}\end{tcolorbox}\hfill\begin{tcolorbox}[enhanced,nobeforeafter,width=7.1cm,height=2.1cm,valign=center,colback=clit!6,colframe=clit,arc=2.5mm,boxrule=0.9pt,left=3mm]{\Large\bfseries\color{clit}b)}\quad{\Large $x^2-4$}\end{tcolorbox}\par\vspace{2.5mm}
\begin{tcolorbox}[enhanced,nobeforeafter,width=7.1cm,height=2.1cm,valign=center,colback=clit!6,colframe=clit,arc=2.5mm,boxrule=0.9pt,left=3mm]{\Large\bfseries\color{clit}c)}\quad{\Large $x^2-4x+3$}\end{tcolorbox}\hfill\begin{tcolorbox}[enhanced,nobeforeafter,width=7.1cm,height=2.1cm,valign=center,colback=clit!6,colframe=clit,arc=2.5mm,boxrule=0.9pt,left=3mm]{\Large\bfseries\color{clit}d)}\quad{\Large $x^2+3x-4$}\end{tcolorbox}
\end{tcolorbox}
\newpage
\begin{tcolorbox}[enhanced,width=16cm,height=14cm,colback=white,colframe=clit,boxrule=2pt,arc=0mm,
left=6mm,right=6mm,top=0mm,bottom=4mm,sharp corners,
overlay={\fill[clit] (frame.north west) rectangle ([yshift=-1.25cm]frame.north east);
\node[anchor=west,text=white,font=\sffamily\bfseries\LARGE] at ([xshift=5mm,yshift=-0.62cm]frame.north west) {Question 37};
\node[anchor=east,text=white,font=\sffamily\itshape\large] at ([xshift=-5mm,yshift=-0.62cm]frame.north east) {Calcul littéral et équations};}]
\vspace*{1.75cm}
\begin{minipage}[t][6.9cm][c]{\linewidth}\LARGE\centering
\begin{minipage}{0.95\linewidth}\raggedright\linespread{1.25}\selectfont Réduire l'expression $5x+3-2x+7$.\end{minipage}
\end{minipage}

\vfill
\begin{tcolorbox}[enhanced,nobeforeafter,width=7.1cm,height=2.1cm,valign=center,colback=clit!6,colframe=clit,arc=2.5mm,boxrule=0.9pt,left=3mm]{\Large\bfseries\color{clit}a)}\quad{\Large $7x+10$}\end{tcolorbox}\hfill\begin{tcolorbox}[enhanced,nobeforeafter,width=7.1cm,height=2.1cm,valign=center,colback=clit!6,colframe=clit,arc=2.5mm,boxrule=0.9pt,left=3mm]{\Large\bfseries\color{clit}b)}\quad{\Large $3x-4$}\end{tcolorbox}\par\vspace{2.5mm}
\begin{tcolorbox}[enhanced,nobeforeafter,width=7.1cm,height=2.1cm,valign=center,colback=clit!6,colframe=clit,arc=2.5mm,boxrule=0.9pt,left=3mm]{\Large\bfseries\color{clit}c)}\quad{\Large $13x$}\end{tcolorbox}\hfill\begin{tcolorbox}[enhanced,nobeforeafter,width=7.1cm,height=2.1cm,valign=center,colback=clit!6,colframe=clit,arc=2.5mm,boxrule=0.9pt,left=3mm]{\Large\bfseries\color{clit}d)}\quad{\Large $3x+10$}\end{tcolorbox}
\end{tcolorbox}
\newpage
\begin{tcolorbox}[enhanced,width=16cm,height=14cm,colback=white,colframe=clit,boxrule=2pt,arc=0mm,
left=6mm,right=6mm,top=0mm,bottom=4mm,sharp corners,
overlay={\fill[clit] (frame.north west) rectangle ([yshift=-1.25cm]frame.north east);
\node[anchor=west,text=white,font=\sffamily\bfseries\LARGE] at ([xshift=5mm,yshift=-0.62cm]frame.north west) {Question 38};
\node[anchor=east,text=white,font=\sffamily\itshape\large] at ([xshift=-5mm,yshift=-0.62cm]frame.north east) {Calcul littéral et équations};}]
\vspace*{1.75cm}
\begin{minipage}[t][6.9cm][c]{\linewidth}\LARGE\centering
\begin{minipage}{0.95\linewidth}\raggedright\linespread{1.25}\selectfont Réduire l'expression $4x-6+3x+2$.\end{minipage}
\end{minipage}

\vfill
\begin{tcolorbox}[enhanced,nobeforeafter,width=7.1cm,height=2.1cm,valign=center,colback=clit!6,colframe=clit,arc=2.5mm,boxrule=0.9pt,left=3mm]{\Large\bfseries\color{clit}a)}\quad{\Large $3x$}\end{tcolorbox}\hfill\begin{tcolorbox}[enhanced,nobeforeafter,width=7.1cm,height=2.1cm,valign=center,colback=clit!6,colframe=clit,arc=2.5mm,boxrule=0.9pt,left=3mm]{\Large\bfseries\color{clit}b)}\quad{\Large $7x-4$}\end{tcolorbox}\par\vspace{2.5mm}
\begin{tcolorbox}[enhanced,nobeforeafter,width=7.1cm,height=2.1cm,valign=center,colback=clit!6,colframe=clit,arc=2.5mm,boxrule=0.9pt,left=3mm]{\Large\bfseries\color{clit}c)}\quad{\Large $7x-8$}\end{tcolorbox}\hfill\begin{tcolorbox}[enhanced,nobeforeafter,width=7.1cm,height=2.1cm,valign=center,colback=clit!6,colframe=clit,arc=2.5mm,boxrule=0.9pt,left=3mm]{\Large\bfseries\color{clit}d)}\quad{\Large $x-4$}\end{tcolorbox}
\end{tcolorbox}
\newpage
\begin{tcolorbox}[enhanced,width=16cm,height=14cm,colback=white,colframe=clit,boxrule=2pt,arc=0mm,
left=6mm,right=6mm,top=0mm,bottom=4mm,sharp corners,
overlay={\fill[clit] (frame.north west) rectangle ([yshift=-1.25cm]frame.north east);
\node[anchor=west,text=white,font=\sffamily\bfseries\LARGE] at ([xshift=5mm,yshift=-0.62cm]frame.north west) {Question 39};
\node[anchor=east,text=white,font=\sffamily\itshape\large] at ([xshift=-5mm,yshift=-0.62cm]frame.north east) {Calcul littéral et équations};}]
\vspace*{1.75cm}
\begin{minipage}[t][6.9cm][c]{\linewidth}\LARGE\centering
\begin{minipage}{0.95\linewidth}\raggedright\linespread{1.25}\selectfont Factoriser $6x+15$.\end{minipage}
\end{minipage}

\vfill
\begin{tcolorbox}[enhanced,nobeforeafter,width=7.1cm,height=2.1cm,valign=center,colback=clit!6,colframe=clit,arc=2.5mm,boxrule=0.9pt,left=3mm]{\Large\bfseries\color{clit}a)}\quad{\Large $6(x+5)$}\end{tcolorbox}\hfill\begin{tcolorbox}[enhanced,nobeforeafter,width=7.1cm,height=2.1cm,valign=center,colback=clit!6,colframe=clit,arc=2.5mm,boxrule=0.9pt,left=3mm]{\Large\bfseries\color{clit}b)}\quad{\Large $3(2x+15)$}\end{tcolorbox}\par\vspace{2.5mm}
\begin{tcolorbox}[enhanced,nobeforeafter,width=7.1cm,height=2.1cm,valign=center,colback=clit!6,colframe=clit,arc=2.5mm,boxrule=0.9pt,left=3mm]{\Large\bfseries\color{clit}c)}\quad{\Large $3(6x+5)$}\end{tcolorbox}\hfill\begin{tcolorbox}[enhanced,nobeforeafter,width=7.1cm,height=2.1cm,valign=center,colback=clit!6,colframe=clit,arc=2.5mm,boxrule=0.9pt,left=3mm]{\Large\bfseries\color{clit}d)}\quad{\Large $3(2x+5)$}\end{tcolorbox}
\end{tcolorbox}
\newpage
\begin{tcolorbox}[enhanced,width=16cm,height=14cm,colback=white,colframe=clit,boxrule=2pt,arc=0mm,
left=6mm,right=6mm,top=0mm,bottom=4mm,sharp corners,
overlay={\fill[clit] (frame.north west) rectangle ([yshift=-1.25cm]frame.north east);
\node[anchor=west,text=white,font=\sffamily\bfseries\LARGE] at ([xshift=5mm,yshift=-0.62cm]frame.north west) {Question 40};
\node[anchor=east,text=white,font=\sffamily\itshape\large] at ([xshift=-5mm,yshift=-0.62cm]frame.north east) {Calcul littéral et équations};}]
\vspace*{1.75cm}
\begin{minipage}[t][6.9cm][c]{\linewidth}\LARGE\centering
\begin{minipage}{0.95\linewidth}\raggedright\linespread{1.25}\selectfont Factoriser $12x-4$.\end{minipage}
\end{minipage}

\vfill
\begin{tcolorbox}[enhanced,nobeforeafter,width=7.1cm,height=2.1cm,valign=center,colback=clit!6,colframe=clit,arc=2.5mm,boxrule=0.9pt,left=3mm]{\Large\bfseries\color{clit}a)}\quad{\Large $4(3x-1)$}\end{tcolorbox}\hfill\begin{tcolorbox}[enhanced,nobeforeafter,width=7.1cm,height=2.1cm,valign=center,colback=clit!6,colframe=clit,arc=2.5mm,boxrule=0.9pt,left=3mm]{\Large\bfseries\color{clit}b)}\quad{\Large $4(3x-4)$}\end{tcolorbox}\par\vspace{2.5mm}
\begin{tcolorbox}[enhanced,nobeforeafter,width=7.1cm,height=2.1cm,valign=center,colback=clit!6,colframe=clit,arc=2.5mm,boxrule=0.9pt,left=3mm]{\Large\bfseries\color{clit}c)}\quad{\Large $4(12x-1)$}\end{tcolorbox}\hfill\begin{tcolorbox}[enhanced,nobeforeafter,width=7.1cm,height=2.1cm,valign=center,colback=clit!6,colframe=clit,arc=2.5mm,boxrule=0.9pt,left=3mm]{\Large\bfseries\color{clit}d)}\quad{\Large $12(x-1)$}\end{tcolorbox}
\end{tcolorbox}
\newpage
\begin{tcolorbox}[enhanced,width=16cm,height=14cm,colback=white,colframe=clit,boxrule=2pt,arc=0mm,
left=6mm,right=6mm,top=0mm,bottom=4mm,sharp corners,
overlay={\fill[clit] (frame.north west) rectangle ([yshift=-1.25cm]frame.north east);
\node[anchor=west,text=white,font=\sffamily\bfseries\LARGE] at ([xshift=5mm,yshift=-0.62cm]frame.north west) {Question 41};
\node[anchor=east,text=white,font=\sffamily\itshape\large] at ([xshift=-5mm,yshift=-0.62cm]frame.north east) {Calcul littéral et équations};}]
\vspace*{1.75cm}
\begin{minipage}[t][6.9cm][c]{\linewidth}\LARGE\centering
\begin{minipage}{0.95\linewidth}\raggedright\linespread{1.25}\selectfont Que vaut $3x^2-2x+1$ pour $x=-2$~?\end{minipage}
\end{minipage}

\vfill
\begin{tcolorbox}[enhanced,nobeforeafter,width=7.1cm,height=2.1cm,valign=center,colback=clit!6,colframe=clit,arc=2.5mm,boxrule=0.9pt,left=3mm]{\Large\bfseries\color{clit}a)}\quad{\Large $18$}\end{tcolorbox}\hfill\begin{tcolorbox}[enhanced,nobeforeafter,width=7.1cm,height=2.1cm,valign=center,colback=clit!6,colframe=clit,arc=2.5mm,boxrule=0.9pt,left=3mm]{\Large\bfseries\color{clit}b)}\quad{\Large $17$}\end{tcolorbox}\par\vspace{2.5mm}
\begin{tcolorbox}[enhanced,nobeforeafter,width=7.1cm,height=2.1cm,valign=center,colback=clit!6,colframe=clit,arc=2.5mm,boxrule=0.9pt,left=3mm]{\Large\bfseries\color{clit}c)}\quad{\Large $-7$}\end{tcolorbox}\hfill\begin{tcolorbox}[enhanced,nobeforeafter,width=7.1cm,height=2.1cm,valign=center,colback=clit!6,colframe=clit,arc=2.5mm,boxrule=0.9pt,left=3mm]{\Large\bfseries\color{clit}d)}\quad{\Large $-17$}\end{tcolorbox}
\end{tcolorbox}
\newpage
\begin{tcolorbox}[enhanced,width=16cm,height=14cm,colback=white,colframe=clit,boxrule=2pt,arc=0mm,
left=6mm,right=6mm,top=0mm,bottom=4mm,sharp corners,
overlay={\fill[clit] (frame.north west) rectangle ([yshift=-1.25cm]frame.north east);
\node[anchor=west,text=white,font=\sffamily\bfseries\LARGE] at ([xshift=5mm,yshift=-0.62cm]frame.north west) {Question 42};
\node[anchor=east,text=white,font=\sffamily\itshape\large] at ([xshift=-5mm,yshift=-0.62cm]frame.north east) {Calcul littéral et équations};}]
\vspace*{1.75cm}
\begin{minipage}[t][6.9cm][c]{\linewidth}\LARGE\centering
\begin{minipage}{0.95\linewidth}\raggedright\linespread{1.25}\selectfont Que vaut $2x^2+5x-3$ pour $x=-1$~?\end{minipage}
\end{minipage}

\vfill
\begin{tcolorbox}[enhanced,nobeforeafter,width=7.1cm,height=2.1cm,valign=center,colback=clit!6,colframe=clit,arc=2.5mm,boxrule=0.9pt,left=3mm]{\Large\bfseries\color{clit}a)}\quad{\Large $-10$}\end{tcolorbox}\hfill\begin{tcolorbox}[enhanced,nobeforeafter,width=7.1cm,height=2.1cm,valign=center,colback=clit!6,colframe=clit,arc=2.5mm,boxrule=0.9pt,left=3mm]{\Large\bfseries\color{clit}b)}\quad{\Large $-6$}\end{tcolorbox}\par\vspace{2.5mm}
\begin{tcolorbox}[enhanced,nobeforeafter,width=7.1cm,height=2.1cm,valign=center,colback=clit!6,colframe=clit,arc=2.5mm,boxrule=0.9pt,left=3mm]{\Large\bfseries\color{clit}c)}\quad{\Large $6$}\end{tcolorbox}\hfill\begin{tcolorbox}[enhanced,nobeforeafter,width=7.1cm,height=2.1cm,valign=center,colback=clit!6,colframe=clit,arc=2.5mm,boxrule=0.9pt,left=3mm]{\Large\bfseries\color{clit}d)}\quad{\Large $-12$}\end{tcolorbox}
\end{tcolorbox}
\newpage
\begin{tcolorbox}[enhanced,width=16cm,height=14cm,colback=white,colframe=clit,boxrule=2pt,arc=0mm,
left=6mm,right=6mm,top=0mm,bottom=4mm,sharp corners,
overlay={\fill[clit] (frame.north west) rectangle ([yshift=-1.25cm]frame.north east);
\node[anchor=west,text=white,font=\sffamily\bfseries\LARGE] at ([xshift=5mm,yshift=-0.62cm]frame.north west) {Question 43};
\node[anchor=east,text=white,font=\sffamily\itshape\large] at ([xshift=-5mm,yshift=-0.62cm]frame.north east) {Calcul littéral et équations};}]
\vspace*{1.75cm}
\begin{minipage}[t][6.9cm][c]{\linewidth}\LARGE\centering
\begin{minipage}{0.95\linewidth}\raggedright\linespread{1.25}\selectfont Quelle est la solution de l'équation $3x+5=20$~?\end{minipage}
\end{minipage}

\vfill
\begin{tcolorbox}[enhanced,nobeforeafter,width=7.1cm,height=2.1cm,valign=center,colback=clit!6,colframe=clit,arc=2.5mm,boxrule=0.9pt,left=3mm]{\Large\bfseries\color{clit}a)}\quad{\Large $x=\dfrac{25}{3}$}\end{tcolorbox}\hfill\begin{tcolorbox}[enhanced,nobeforeafter,width=7.1cm,height=2.1cm,valign=center,colback=clit!6,colframe=clit,arc=2.5mm,boxrule=0.9pt,left=3mm]{\Large\bfseries\color{clit}b)}\quad{\Large $x=12$}\end{tcolorbox}\par\vspace{2.5mm}
\begin{tcolorbox}[enhanced,nobeforeafter,width=7.1cm,height=2.1cm,valign=center,colback=clit!6,colframe=clit,arc=2.5mm,boxrule=0.9pt,left=3mm]{\Large\bfseries\color{clit}c)}\quad{\Large $x=5$}\end{tcolorbox}\hfill\begin{tcolorbox}[enhanced,nobeforeafter,width=7.1cm,height=2.1cm,valign=center,colback=clit!6,colframe=clit,arc=2.5mm,boxrule=0.9pt,left=3mm]{\Large\bfseries\color{clit}d)}\quad{\Large $x=\dfrac{1}{5}$}\end{tcolorbox}
\end{tcolorbox}
\newpage
\begin{tcolorbox}[enhanced,width=16cm,height=14cm,colback=white,colframe=clit,boxrule=2pt,arc=0mm,
left=6mm,right=6mm,top=0mm,bottom=4mm,sharp corners,
overlay={\fill[clit] (frame.north west) rectangle ([yshift=-1.25cm]frame.north east);
\node[anchor=west,text=white,font=\sffamily\bfseries\LARGE] at ([xshift=5mm,yshift=-0.62cm]frame.north west) {Question 44};
\node[anchor=east,text=white,font=\sffamily\itshape\large] at ([xshift=-5mm,yshift=-0.62cm]frame.north east) {Calcul littéral et équations};}]
\vspace*{1.75cm}
\begin{minipage}[t][6.9cm][c]{\linewidth}\LARGE\centering
\begin{minipage}{0.95\linewidth}\raggedright\linespread{1.25}\selectfont Quelle est la solution de l'équation $4x-7=21$~?\end{minipage}
\end{minipage}

\vfill
\begin{tcolorbox}[enhanced,nobeforeafter,width=7.1cm,height=2.1cm,valign=center,colback=clit!6,colframe=clit,arc=2.5mm,boxrule=0.9pt,left=3mm]{\Large\bfseries\color{clit}a)}\quad{\Large $x=\dfrac{1}{7}$}\end{tcolorbox}\hfill\begin{tcolorbox}[enhanced,nobeforeafter,width=7.1cm,height=2.1cm,valign=center,colback=clit!6,colframe=clit,arc=2.5mm,boxrule=0.9pt,left=3mm]{\Large\bfseries\color{clit}b)}\quad{\Large $x=24$}\end{tcolorbox}\par\vspace{2.5mm}
\begin{tcolorbox}[enhanced,nobeforeafter,width=7.1cm,height=2.1cm,valign=center,colback=clit!6,colframe=clit,arc=2.5mm,boxrule=0.9pt,left=3mm]{\Large\bfseries\color{clit}c)}\quad{\Large $x=\dfrac{7}{2}$}\end{tcolorbox}\hfill\begin{tcolorbox}[enhanced,nobeforeafter,width=7.1cm,height=2.1cm,valign=center,colback=clit!6,colframe=clit,arc=2.5mm,boxrule=0.9pt,left=3mm]{\Large\bfseries\color{clit}d)}\quad{\Large $x=7$}\end{tcolorbox}
\end{tcolorbox}
\newpage
\begin{tcolorbox}[enhanced,width=16cm,height=14cm,colback=white,colframe=clit,boxrule=2pt,arc=0mm,
left=6mm,right=6mm,top=0mm,bottom=4mm,sharp corners,
overlay={\fill[clit] (frame.north west) rectangle ([yshift=-1.25cm]frame.north east);
\node[anchor=west,text=white,font=\sffamily\bfseries\LARGE] at ([xshift=5mm,yshift=-0.62cm]frame.north west) {Question 45};
\node[anchor=east,text=white,font=\sffamily\itshape\large] at ([xshift=-5mm,yshift=-0.62cm]frame.north east) {Calcul littéral et équations};}]
\vspace*{1.75cm}
\begin{minipage}[t][6.9cm][c]{\linewidth}\LARGE\centering
\begin{minipage}{0.95\linewidth}\raggedright\linespread{1.25}\selectfont Quelle est la solution de l'équation $5x+3=2x+12$~?\end{minipage}
\end{minipage}

\vfill
\begin{tcolorbox}[enhanced,nobeforeafter,width=7.1cm,height=2.1cm,valign=center,colback=clit!6,colframe=clit,arc=2.5mm,boxrule=0.9pt,left=3mm]{\Large\bfseries\color{clit}a)}\quad{\Large $x=5$}\end{tcolorbox}\hfill\begin{tcolorbox}[enhanced,nobeforeafter,width=7.1cm,height=2.1cm,valign=center,colback=clit!6,colframe=clit,arc=2.5mm,boxrule=0.9pt,left=3mm]{\Large\bfseries\color{clit}b)}\quad{\Large $x=\dfrac{9}{7}$}\end{tcolorbox}\par\vspace{2.5mm}
\begin{tcolorbox}[enhanced,nobeforeafter,width=7.1cm,height=2.1cm,valign=center,colback=clit!6,colframe=clit,arc=2.5mm,boxrule=0.9pt,left=3mm]{\Large\bfseries\color{clit}c)}\quad{\Large $x=-3$}\end{tcolorbox}\hfill\begin{tcolorbox}[enhanced,nobeforeafter,width=7.1cm,height=2.1cm,valign=center,colback=clit!6,colframe=clit,arc=2.5mm,boxrule=0.9pt,left=3mm]{\Large\bfseries\color{clit}d)}\quad{\Large $x=3$}\end{tcolorbox}
\end{tcolorbox}
\newpage
\begin{tcolorbox}[enhanced,width=16cm,height=14cm,colback=white,colframe=clit,boxrule=2pt,arc=0mm,
left=6mm,right=6mm,top=0mm,bottom=4mm,sharp corners,
overlay={\fill[clit] (frame.north west) rectangle ([yshift=-1.25cm]frame.north east);
\node[anchor=west,text=white,font=\sffamily\bfseries\LARGE] at ([xshift=5mm,yshift=-0.62cm]frame.north west) {Question 46};
\node[anchor=east,text=white,font=\sffamily\itshape\large] at ([xshift=-5mm,yshift=-0.62cm]frame.north east) {Calcul littéral et équations};}]
\vspace*{1.75cm}
\begin{minipage}[t][6.9cm][c]{\linewidth}\LARGE\centering
\begin{minipage}{0.95\linewidth}\raggedright\linespread{1.25}\selectfont Quelle est la solution de l'équation $7x-4=3x+8$~?\end{minipage}
\end{minipage}

\vfill
\begin{tcolorbox}[enhanced,nobeforeafter,width=7.1cm,height=2.1cm,valign=center,colback=clit!6,colframe=clit,arc=2.5mm,boxrule=0.9pt,left=3mm]{\Large\bfseries\color{clit}a)}\quad{\Large $x=-3$}\end{tcolorbox}\hfill\begin{tcolorbox}[enhanced,nobeforeafter,width=7.1cm,height=2.1cm,valign=center,colback=clit!6,colframe=clit,arc=2.5mm,boxrule=0.9pt,left=3mm]{\Large\bfseries\color{clit}b)}\quad{\Large $x=1$}\end{tcolorbox}\par\vspace{2.5mm}
\begin{tcolorbox}[enhanced,nobeforeafter,width=7.1cm,height=2.1cm,valign=center,colback=clit!6,colframe=clit,arc=2.5mm,boxrule=0.9pt,left=3mm]{\Large\bfseries\color{clit}c)}\quad{\Large $x=3$}\end{tcolorbox}\hfill\begin{tcolorbox}[enhanced,nobeforeafter,width=7.1cm,height=2.1cm,valign=center,colback=clit!6,colframe=clit,arc=2.5mm,boxrule=0.9pt,left=3mm]{\Large\bfseries\color{clit}d)}\quad{\Large $x=\dfrac{6}{5}$}\end{tcolorbox}
\end{tcolorbox}
\newpage
\begin{tcolorbox}[enhanced,width=16cm,height=14cm,colback=white,colframe=clit,boxrule=2pt,arc=0mm,
left=6mm,right=6mm,top=0mm,bottom=4mm,sharp corners,
overlay={\fill[clit] (frame.north west) rectangle ([yshift=-1.25cm]frame.north east);
\node[anchor=west,text=white,font=\sffamily\bfseries\LARGE] at ([xshift=5mm,yshift=-0.62cm]frame.north west) {Question 47};
\node[anchor=east,text=white,font=\sffamily\itshape\large] at ([xshift=-5mm,yshift=-0.62cm]frame.north east) {Calcul littéral et équations};}]
\vspace*{1.75cm}
\begin{minipage}[t][6.9cm][c]{\linewidth}\LARGE\centering
\begin{minipage}{0.95\linewidth}\raggedright\linespread{1.25}\selectfont Le triple d'un nombre, augmenté de 5, est égal à $26$. Quel est ce nombre~?\end{minipage}
\end{minipage}

\vfill
\begin{tcolorbox}[enhanced,nobeforeafter,width=7.1cm,height=2.1cm,valign=center,colback=clit!6,colframe=clit,arc=2.5mm,boxrule=0.9pt,left=3mm]{\Large\bfseries\color{clit}a)}\quad{\Large $\dfrac{31}{3}$}\end{tcolorbox}\hfill\begin{tcolorbox}[enhanced,nobeforeafter,width=7.1cm,height=2.1cm,valign=center,colback=clit!6,colframe=clit,arc=2.5mm,boxrule=0.9pt,left=3mm]{\Large\bfseries\color{clit}b)}\quad{\Large $21$}\end{tcolorbox}\par\vspace{2.5mm}
\begin{tcolorbox}[enhanced,nobeforeafter,width=7.1cm,height=2.1cm,valign=center,colback=clit!6,colframe=clit,arc=2.5mm,boxrule=0.9pt,left=3mm]{\Large\bfseries\color{clit}c)}\quad{\Large $\dfrac{26}{3}$}\end{tcolorbox}\hfill\begin{tcolorbox}[enhanced,nobeforeafter,width=7.1cm,height=2.1cm,valign=center,colback=clit!6,colframe=clit,arc=2.5mm,boxrule=0.9pt,left=3mm]{\Large\bfseries\color{clit}d)}\quad{\Large $7$}\end{tcolorbox}
\end{tcolorbox}
\newpage
\begin{tcolorbox}[enhanced,width=16cm,height=14cm,colback=white,colframe=clit,boxrule=2pt,arc=0mm,
left=6mm,right=6mm,top=0mm,bottom=4mm,sharp corners,
overlay={\fill[clit] (frame.north west) rectangle ([yshift=-1.25cm]frame.north east);
\node[anchor=west,text=white,font=\sffamily\bfseries\LARGE] at ([xshift=5mm,yshift=-0.62cm]frame.north west) {Question 48};
\node[anchor=east,text=white,font=\sffamily\itshape\large] at ([xshift=-5mm,yshift=-0.62cm]frame.north east) {Calcul littéral et équations};}]
\vspace*{1.75cm}
\begin{minipage}[t][6.9cm][c]{\linewidth}\LARGE\centering
\begin{minipage}{0.95\linewidth}\raggedright\linespread{1.25}\selectfont Le double d'un nombre, diminué de 7, est égal à $15$. Quel est ce nombre~?\end{minipage}
\end{minipage}

\vfill
\begin{tcolorbox}[enhanced,nobeforeafter,width=7.1cm,height=2.1cm,valign=center,colback=clit!6,colframe=clit,arc=2.5mm,boxrule=0.9pt,left=3mm]{\Large\bfseries\color{clit}a)}\quad{\Large $11$}\end{tcolorbox}\hfill\begin{tcolorbox}[enhanced,nobeforeafter,width=7.1cm,height=2.1cm,valign=center,colback=clit!6,colframe=clit,arc=2.5mm,boxrule=0.9pt,left=3mm]{\Large\bfseries\color{clit}b)}\quad{\Large $4$}\end{tcolorbox}\par\vspace{2.5mm}
\begin{tcolorbox}[enhanced,nobeforeafter,width=7.1cm,height=2.1cm,valign=center,colback=clit!6,colframe=clit,arc=2.5mm,boxrule=0.9pt,left=3mm]{\Large\bfseries\color{clit}c)}\quad{\Large $\dfrac{15}{2}$}\end{tcolorbox}\hfill\begin{tcolorbox}[enhanced,nobeforeafter,width=7.1cm,height=2.1cm,valign=center,colback=clit!6,colframe=clit,arc=2.5mm,boxrule=0.9pt,left=3mm]{\Large\bfseries\color{clit}d)}\quad{\Large $22$}\end{tcolorbox}
\end{tcolorbox}
\newpage
\begin{tcolorbox}[enhanced,width=16cm,height=14cm,colback=white,colframe=cgeo,boxrule=2pt,arc=0mm,
left=6mm,right=6mm,top=0mm,bottom=4mm,sharp corners,
overlay={\fill[cgeo] (frame.north west) rectangle ([yshift=-1.25cm]frame.north east);
\node[anchor=west,text=white,font=\sffamily\bfseries\LARGE] at ([xshift=5mm,yshift=-0.62cm]frame.north west) {Question 49};
\node[anchor=east,text=white,font=\sffamily\itshape\large] at ([xshift=-5mm,yshift=-0.62cm]frame.north east) {Géométrie};}]
\vspace*{1.75cm}
\begin{minipage}[t][6.9cm][c]{\linewidth}\LARGE\centering
\begin{minipage}{0.95\linewidth}\raggedright\linespread{1.25}\selectfont Le triangle $ABC$ est rectangle en $A$, avec $AB=3$ cm et $AC=4$ cm. Combien mesure $BC$~?\end{minipage}
\end{minipage}

\vfill
\begin{tcolorbox}[enhanced,nobeforeafter,width=7.1cm,height=2.1cm,valign=center,colback=cgeo!6,colframe=cgeo,arc=2.5mm,boxrule=0.9pt,left=3mm]{\Large\bfseries\color{cgeo}a)}\quad{\Large $25~\text{cm}$}\end{tcolorbox}\hfill\begin{tcolorbox}[enhanced,nobeforeafter,width=7.1cm,height=2.1cm,valign=center,colback=cgeo!6,colframe=cgeo,arc=2.5mm,boxrule=0.9pt,left=3mm]{\Large\bfseries\color{cgeo}b)}\quad{\Large $5~\text{cm}$}\end{tcolorbox}\par\vspace{2.5mm}
\begin{tcolorbox}[enhanced,nobeforeafter,width=7.1cm,height=2.1cm,valign=center,colback=cgeo!6,colframe=cgeo,arc=2.5mm,boxrule=0.9pt,left=3mm]{\Large\bfseries\color{cgeo}c)}\quad{\Large $7~\text{cm}$}\end{tcolorbox}\hfill\begin{tcolorbox}[enhanced,nobeforeafter,width=7.1cm,height=2.1cm,valign=center,colback=cgeo!6,colframe=cgeo,arc=2.5mm,boxrule=0.9pt,left=3mm]{\Large\bfseries\color{cgeo}d)}\quad{\Large $\sqrt{7}~\text{cm}$}\end{tcolorbox}
\end{tcolorbox}
\newpage
\begin{tcolorbox}[enhanced,width=16cm,height=14cm,colback=white,colframe=cgeo,boxrule=2pt,arc=0mm,
left=6mm,right=6mm,top=0mm,bottom=4mm,sharp corners,
overlay={\fill[cgeo] (frame.north west) rectangle ([yshift=-1.25cm]frame.north east);
\node[anchor=west,text=white,font=\sffamily\bfseries\LARGE] at ([xshift=5mm,yshift=-0.62cm]frame.north west) {Question 50};
\node[anchor=east,text=white,font=\sffamily\itshape\large] at ([xshift=-5mm,yshift=-0.62cm]frame.north east) {Géométrie};}]
\vspace*{1.75cm}
\begin{minipage}[t][6.9cm][c]{\linewidth}\LARGE\centering
\begin{minipage}{0.95\linewidth}\raggedright\linespread{1.25}\selectfont Le triangle $ABC$ est rectangle en $A$, avec $AB=6$ cm et $AC=8$ cm. Combien mesure $BC$~?\end{minipage}
\end{minipage}

\vfill
\begin{tcolorbox}[enhanced,nobeforeafter,width=7.1cm,height=2.1cm,valign=center,colback=cgeo!6,colframe=cgeo,arc=2.5mm,boxrule=0.9pt,left=3mm]{\Large\bfseries\color{cgeo}a)}\quad{\Large $10~\text{cm}$}\end{tcolorbox}\hfill\begin{tcolorbox}[enhanced,nobeforeafter,width=7.1cm,height=2.1cm,valign=center,colback=cgeo!6,colframe=cgeo,arc=2.5mm,boxrule=0.9pt,left=3mm]{\Large\bfseries\color{cgeo}b)}\quad{\Large $100~\text{cm}$}\end{tcolorbox}\par\vspace{2.5mm}
\begin{tcolorbox}[enhanced,nobeforeafter,width=7.1cm,height=2.1cm,valign=center,colback=cgeo!6,colframe=cgeo,arc=2.5mm,boxrule=0.9pt,left=3mm]{\Large\bfseries\color{cgeo}c)}\quad{\Large $14~\text{cm}$}\end{tcolorbox}\hfill\begin{tcolorbox}[enhanced,nobeforeafter,width=7.1cm,height=2.1cm,valign=center,colback=cgeo!6,colframe=cgeo,arc=2.5mm,boxrule=0.9pt,left=3mm]{\Large\bfseries\color{cgeo}d)}\quad{\Large $\sqrt{28}~\text{cm}$}\end{tcolorbox}
\end{tcolorbox}
\newpage
\begin{tcolorbox}[enhanced,width=16cm,height=14cm,colback=white,colframe=cgeo,boxrule=2pt,arc=0mm,
left=6mm,right=6mm,top=0mm,bottom=4mm,sharp corners,
overlay={\fill[cgeo] (frame.north west) rectangle ([yshift=-1.25cm]frame.north east);
\node[anchor=west,text=white,font=\sffamily\bfseries\LARGE] at ([xshift=5mm,yshift=-0.62cm]frame.north west) {Question 51};
\node[anchor=east,text=white,font=\sffamily\itshape\large] at ([xshift=-5mm,yshift=-0.62cm]frame.north east) {Géométrie};}]
\vspace*{1.75cm}
\begin{minipage}[t][6.9cm][c]{\linewidth}\LARGE\centering
\begin{minipage}{0.95\linewidth}\raggedright\linespread{1.25}\selectfont Le triangle $ABC$ est rectangle en $A$, avec $BC=10$ cm et $AB=6$ cm. Combien mesure $AC$~?\end{minipage}
\end{minipage}

\vfill
\begin{tcolorbox}[enhanced,nobeforeafter,width=7.1cm,height=2.1cm,valign=center,colback=cgeo!6,colframe=cgeo,arc=2.5mm,boxrule=0.9pt,left=3mm]{\Large\bfseries\color{cgeo}a)}\quad{\Large $4~\text{cm}$}\end{tcolorbox}\hfill\begin{tcolorbox}[enhanced,nobeforeafter,width=7.1cm,height=2.1cm,valign=center,colback=cgeo!6,colframe=cgeo,arc=2.5mm,boxrule=0.9pt,left=3mm]{\Large\bfseries\color{cgeo}b)}\quad{\Large $8~\text{cm}$}\end{tcolorbox}\par\vspace{2.5mm}
\begin{tcolorbox}[enhanced,nobeforeafter,width=7.1cm,height=2.1cm,valign=center,colback=cgeo!6,colframe=cgeo,arc=2.5mm,boxrule=0.9pt,left=3mm]{\Large\bfseries\color{cgeo}c)}\quad{\Large $64~\text{cm}$}\end{tcolorbox}\hfill\begin{tcolorbox}[enhanced,nobeforeafter,width=7.1cm,height=2.1cm,valign=center,colback=cgeo!6,colframe=cgeo,arc=2.5mm,boxrule=0.9pt,left=3mm]{\Large\bfseries\color{cgeo}d)}\quad{\Large $\sqrt{136}~\text{cm}$}\end{tcolorbox}
\end{tcolorbox}
\newpage
\begin{tcolorbox}[enhanced,width=16cm,height=14cm,colback=white,colframe=cgeo,boxrule=2pt,arc=0mm,
left=6mm,right=6mm,top=0mm,bottom=4mm,sharp corners,
overlay={\fill[cgeo] (frame.north west) rectangle ([yshift=-1.25cm]frame.north east);
\node[anchor=west,text=white,font=\sffamily\bfseries\LARGE] at ([xshift=5mm,yshift=-0.62cm]frame.north west) {Question 52};
\node[anchor=east,text=white,font=\sffamily\itshape\large] at ([xshift=-5mm,yshift=-0.62cm]frame.north east) {Géométrie};}]
\vspace*{1.75cm}
\begin{minipage}[t][6.9cm][c]{\linewidth}\LARGE\centering
\begin{minipage}{0.95\linewidth}\raggedright\linespread{1.25}\selectfont Le triangle $ABC$ est rectangle en $A$, avec $BC=13$ cm et $AB=5$ cm. Combien mesure $AC$~?\end{minipage}
\end{minipage}

\vfill
\begin{tcolorbox}[enhanced,nobeforeafter,width=7.1cm,height=2.1cm,valign=center,colback=cgeo!6,colframe=cgeo,arc=2.5mm,boxrule=0.9pt,left=3mm]{\Large\bfseries\color{cgeo}a)}\quad{\Large $144~\text{cm}$}\end{tcolorbox}\hfill\begin{tcolorbox}[enhanced,nobeforeafter,width=7.1cm,height=2.1cm,valign=center,colback=cgeo!6,colframe=cgeo,arc=2.5mm,boxrule=0.9pt,left=3mm]{\Large\bfseries\color{cgeo}b)}\quad{\Large $\sqrt{194}~\text{cm}$}\end{tcolorbox}\par\vspace{2.5mm}
\begin{tcolorbox}[enhanced,nobeforeafter,width=7.1cm,height=2.1cm,valign=center,colback=cgeo!6,colframe=cgeo,arc=2.5mm,boxrule=0.9pt,left=3mm]{\Large\bfseries\color{cgeo}c)}\quad{\Large $8~\text{cm}$}\end{tcolorbox}\hfill\begin{tcolorbox}[enhanced,nobeforeafter,width=7.1cm,height=2.1cm,valign=center,colback=cgeo!6,colframe=cgeo,arc=2.5mm,boxrule=0.9pt,left=3mm]{\Large\bfseries\color{cgeo}d)}\quad{\Large $12~\text{cm}$}\end{tcolorbox}
\end{tcolorbox}
\newpage
\begin{tcolorbox}[enhanced,width=16cm,height=14cm,colback=white,colframe=cgeo,boxrule=2pt,arc=0mm,
left=6mm,right=6mm,top=0mm,bottom=4mm,sharp corners,
overlay={\fill[cgeo] (frame.north west) rectangle ([yshift=-1.25cm]frame.north east);
\node[anchor=west,text=white,font=\sffamily\bfseries\LARGE] at ([xshift=5mm,yshift=-0.62cm]frame.north west) {Question 53};
\node[anchor=east,text=white,font=\sffamily\itshape\large] at ([xshift=-5mm,yshift=-0.62cm]frame.north east) {Géométrie};}]
\vspace*{1.75cm}
\begin{minipage}[t][6.9cm][c]{\linewidth}\LARGE\centering
\begin{minipage}{0.95\linewidth}\raggedright\linespread{1.25}\selectfont Un triangle a des côtés de $6$ cm, $8$ cm et $10$ cm. Est-il rectangle~?\end{minipage}
\end{minipage}

\vfill
\begin{tcolorbox}[enhanced,nobeforeafter,width=7.1cm,height=2.1cm,valign=center,colback=cgeo!6,colframe=cgeo,arc=2.5mm,boxrule=0.9pt,left=3mm]{\Large\bfseries\color{cgeo}a)}\quad{\Large oui}\end{tcolorbox}\hfill\begin{tcolorbox}[enhanced,nobeforeafter,width=7.1cm,height=2.1cm,valign=center,colback=cgeo!6,colframe=cgeo,arc=2.5mm,boxrule=0.9pt,left=3mm]{\Large\bfseries\color{cgeo}b)}\quad{\Large s'il est isocèle}\end{tcolorbox}\par\vspace{2.5mm}
\begin{tcolorbox}[enhanced,nobeforeafter,width=7.1cm,height=2.1cm,valign=center,colback=cgeo!6,colframe=cgeo,arc=2.5mm,boxrule=0.9pt,left=3mm]{\Large\bfseries\color{cgeo}c)}\quad{\Large non}\end{tcolorbox}\hfill\begin{tcolorbox}[enhanced,nobeforeafter,width=7.1cm,height=2.1cm,valign=center,colback=cgeo!6,colframe=cgeo,arc=2.5mm,boxrule=0.9pt,left=3mm]{\Large\bfseries\color{cgeo}d)}\quad{\Large on ne peut pas savoir}\end{tcolorbox}
\end{tcolorbox}
\newpage
\begin{tcolorbox}[enhanced,width=16cm,height=14cm,colback=white,colframe=cgeo,boxrule=2pt,arc=0mm,
left=6mm,right=6mm,top=0mm,bottom=4mm,sharp corners,
overlay={\fill[cgeo] (frame.north west) rectangle ([yshift=-1.25cm]frame.north east);
\node[anchor=west,text=white,font=\sffamily\bfseries\LARGE] at ([xshift=5mm,yshift=-0.62cm]frame.north west) {Question 54};
\node[anchor=east,text=white,font=\sffamily\itshape\large] at ([xshift=-5mm,yshift=-0.62cm]frame.north east) {Géométrie};}]
\vspace*{1.75cm}
\begin{minipage}[t][6.9cm][c]{\linewidth}\LARGE\centering
\begin{minipage}{0.95\linewidth}\raggedright\linespread{1.25}\selectfont Un triangle a des côtés de $5$ cm, $7$ cm et $9$ cm. Est-il rectangle~?\end{minipage}
\end{minipage}

\vfill
\begin{tcolorbox}[enhanced,nobeforeafter,width=7.1cm,height=2.1cm,valign=center,colback=cgeo!6,colframe=cgeo,arc=2.5mm,boxrule=0.9pt,left=3mm]{\Large\bfseries\color{cgeo}a)}\quad{\Large s'il est isocèle}\end{tcolorbox}\hfill\begin{tcolorbox}[enhanced,nobeforeafter,width=7.1cm,height=2.1cm,valign=center,colback=cgeo!6,colframe=cgeo,arc=2.5mm,boxrule=0.9pt,left=3mm]{\Large\bfseries\color{cgeo}b)}\quad{\Large oui}\end{tcolorbox}\par\vspace{2.5mm}
\begin{tcolorbox}[enhanced,nobeforeafter,width=7.1cm,height=2.1cm,valign=center,colback=cgeo!6,colframe=cgeo,arc=2.5mm,boxrule=0.9pt,left=3mm]{\Large\bfseries\color{cgeo}c)}\quad{\Large on ne peut pas savoir}\end{tcolorbox}\hfill\begin{tcolorbox}[enhanced,nobeforeafter,width=7.1cm,height=2.1cm,valign=center,colback=cgeo!6,colframe=cgeo,arc=2.5mm,boxrule=0.9pt,left=3mm]{\Large\bfseries\color{cgeo}d)}\quad{\Large non}\end{tcolorbox}
\end{tcolorbox}
\newpage
\begin{tcolorbox}[enhanced,width=16cm,height=14cm,colback=white,colframe=cgeo,boxrule=2pt,arc=0mm,
left=6mm,right=6mm,top=0mm,bottom=4mm,sharp corners,
overlay={\fill[cgeo] (frame.north west) rectangle ([yshift=-1.25cm]frame.north east);
\node[anchor=west,text=white,font=\sffamily\bfseries\LARGE] at ([xshift=5mm,yshift=-0.62cm]frame.north west) {Question 55};
\node[anchor=east,text=white,font=\sffamily\itshape\large] at ([xshift=-5mm,yshift=-0.62cm]frame.north east) {Géométrie};}]
\vspace*{1.75cm}
\begin{minipage}[t][6.9cm][c]{\linewidth}\Large\centering
\begin{minipage}{0.95\linewidth}\raggedright\linespread{1.25}\selectfont Dans le triangle $ABC$, $M$ est sur $[AB]$, $N$ est sur $[AC]$ et les droites $(MN)$ et $(BC)$ sont parallèles. $AM=3$ cm, $AB=6$ cm et $AN=4$ cm. Combien mesure $AC$~?\end{minipage}
\end{minipage}

\vfill
\begin{tcolorbox}[enhanced,nobeforeafter,width=7.1cm,height=2.1cm,valign=center,colback=cgeo!6,colframe=cgeo,arc=2.5mm,boxrule=0.9pt,left=3mm]{\Large\bfseries\color{cgeo}a)}\quad{\Large $2~\text{cm}$}\end{tcolorbox}\hfill\begin{tcolorbox}[enhanced,nobeforeafter,width=7.1cm,height=2.1cm,valign=center,colback=cgeo!6,colframe=cgeo,arc=2.5mm,boxrule=0.9pt,left=3mm]{\Large\bfseries\color{cgeo}b)}\quad{\Large $7~\text{cm}$}\end{tcolorbox}\par\vspace{2.5mm}
\begin{tcolorbox}[enhanced,nobeforeafter,width=7.1cm,height=2.1cm,valign=center,colback=cgeo!6,colframe=cgeo,arc=2.5mm,boxrule=0.9pt,left=3mm]{\Large\bfseries\color{cgeo}c)}\quad{\Large $4{,}5~\text{cm}$}\end{tcolorbox}\hfill\begin{tcolorbox}[enhanced,nobeforeafter,width=7.1cm,height=2.1cm,valign=center,colback=cgeo!6,colframe=cgeo,arc=2.5mm,boxrule=0.9pt,left=3mm]{\Large\bfseries\color{cgeo}d)}\quad{\Large $8~\text{cm}$}\end{tcolorbox}
\end{tcolorbox}
\newpage
\begin{tcolorbox}[enhanced,width=16cm,height=14cm,colback=white,colframe=cgeo,boxrule=2pt,arc=0mm,
left=6mm,right=6mm,top=0mm,bottom=4mm,sharp corners,
overlay={\fill[cgeo] (frame.north west) rectangle ([yshift=-1.25cm]frame.north east);
\node[anchor=west,text=white,font=\sffamily\bfseries\LARGE] at ([xshift=5mm,yshift=-0.62cm]frame.north west) {Question 56};
\node[anchor=east,text=white,font=\sffamily\itshape\large] at ([xshift=-5mm,yshift=-0.62cm]frame.north east) {Géométrie};}]
\vspace*{1.75cm}
\begin{minipage}[t][6.9cm][c]{\linewidth}\Large\centering
\begin{minipage}{0.95\linewidth}\raggedright\linespread{1.25}\selectfont Dans le triangle $ABC$, $M$ est sur $[AB]$, $N$ est sur $[AC]$ et les droites $(MN)$ et $(BC)$ sont parallèles. $AM=2$ cm, $AB=5$ cm et $AN=3$ cm. Combien mesure $AC$~?\end{minipage}
\end{minipage}

\vfill
\begin{tcolorbox}[enhanced,nobeforeafter,width=7.1cm,height=2.1cm,valign=center,colback=cgeo!6,colframe=cgeo,arc=2.5mm,boxrule=0.9pt,left=3mm]{\Large\bfseries\color{cgeo}a)}\quad{\Large $6~\text{cm}$}\end{tcolorbox}\hfill\begin{tcolorbox}[enhanced,nobeforeafter,width=7.1cm,height=2.1cm,valign=center,colback=cgeo!6,colframe=cgeo,arc=2.5mm,boxrule=0.9pt,left=3mm]{\Large\bfseries\color{cgeo}b)}\quad{\Large $1{,}2~\text{cm}$}\end{tcolorbox}\par\vspace{2.5mm}
\begin{tcolorbox}[enhanced,nobeforeafter,width=7.1cm,height=2.1cm,valign=center,colback=cgeo!6,colframe=cgeo,arc=2.5mm,boxrule=0.9pt,left=3mm]{\Large\bfseries\color{cgeo}c)}\quad{\Large $7{,}5~\text{cm}$}\end{tcolorbox}\hfill\begin{tcolorbox}[enhanced,nobeforeafter,width=7.1cm,height=2.1cm,valign=center,colback=cgeo!6,colframe=cgeo,arc=2.5mm,boxrule=0.9pt,left=3mm]{\Large\bfseries\color{cgeo}d)}\quad{\Large $\dfrac{10}{3}~\text{cm}$}\end{tcolorbox}
\end{tcolorbox}
\newpage
\begin{tcolorbox}[enhanced,width=16cm,height=14cm,colback=white,colframe=cgeo,boxrule=2pt,arc=0mm,
left=6mm,right=6mm,top=0mm,bottom=4mm,sharp corners,
overlay={\fill[cgeo] (frame.north west) rectangle ([yshift=-1.25cm]frame.north east);
\node[anchor=west,text=white,font=\sffamily\bfseries\LARGE] at ([xshift=5mm,yshift=-0.62cm]frame.north west) {Question 57};
\node[anchor=east,text=white,font=\sffamily\itshape\large] at ([xshift=-5mm,yshift=-0.62cm]frame.north east) {Géométrie};}]
\vspace*{1.75cm}
\begin{minipage}[t][6.9cm][c]{\linewidth}\LARGE\centering
\begin{minipage}{0.95\linewidth}\raggedright\linespread{1.25}\selectfont Le triangle $ABC$ est rectangle en $A$, avec $BC=5$ cm et $AB=3$ cm. Que vaut $\cos\widehat{ABC}$~?\end{minipage}
\end{minipage}

\vfill
\begin{tcolorbox}[enhanced,nobeforeafter,width=7.1cm,height=2.1cm,valign=center,colback=cgeo!6,colframe=cgeo,arc=2.5mm,boxrule=0.9pt,left=3mm]{\Large\bfseries\color{cgeo}a)}\quad{\Large $\dfrac{3}{2}$}\end{tcolorbox}\hfill\begin{tcolorbox}[enhanced,nobeforeafter,width=7.1cm,height=2.1cm,valign=center,colback=cgeo!6,colframe=cgeo,arc=2.5mm,boxrule=0.9pt,left=3mm]{\Large\bfseries\color{cgeo}b)}\quad{\Large $\dfrac{3}{5}$}\end{tcolorbox}\par\vspace{2.5mm}
\begin{tcolorbox}[enhanced,nobeforeafter,width=7.1cm,height=2.1cm,valign=center,colback=cgeo!6,colframe=cgeo,arc=2.5mm,boxrule=0.9pt,left=3mm]{\Large\bfseries\color{cgeo}c)}\quad{\Large $\dfrac{2}{5}$}\end{tcolorbox}\hfill\begin{tcolorbox}[enhanced,nobeforeafter,width=7.1cm,height=2.1cm,valign=center,colback=cgeo!6,colframe=cgeo,arc=2.5mm,boxrule=0.9pt,left=3mm]{\Large\bfseries\color{cgeo}d)}\quad{\Large $\dfrac{5}{3}$}\end{tcolorbox}
\end{tcolorbox}
\newpage
\begin{tcolorbox}[enhanced,width=16cm,height=14cm,colback=white,colframe=cgeo,boxrule=2pt,arc=0mm,
left=6mm,right=6mm,top=0mm,bottom=4mm,sharp corners,
overlay={\fill[cgeo] (frame.north west) rectangle ([yshift=-1.25cm]frame.north east);
\node[anchor=west,text=white,font=\sffamily\bfseries\LARGE] at ([xshift=5mm,yshift=-0.62cm]frame.north west) {Question 58};
\node[anchor=east,text=white,font=\sffamily\itshape\large] at ([xshift=-5mm,yshift=-0.62cm]frame.north east) {Géométrie};}]
\vspace*{1.75cm}
\begin{minipage}[t][6.9cm][c]{\linewidth}\LARGE\centering
\begin{minipage}{0.95\linewidth}\raggedright\linespread{1.25}\selectfont Le triangle $ABC$ est rectangle en $A$, avec $BC=10$ cm et $AB=8$ cm. Que vaut $\cos\widehat{ABC}$~?\end{minipage}
\end{minipage}

\vfill
\begin{tcolorbox}[enhanced,nobeforeafter,width=7.1cm,height=2.1cm,valign=center,colback=cgeo!6,colframe=cgeo,arc=2.5mm,boxrule=0.9pt,left=3mm]{\Large\bfseries\color{cgeo}a)}\quad{\Large $4$}\end{tcolorbox}\hfill\begin{tcolorbox}[enhanced,nobeforeafter,width=7.1cm,height=2.1cm,valign=center,colback=cgeo!6,colframe=cgeo,arc=2.5mm,boxrule=0.9pt,left=3mm]{\Large\bfseries\color{cgeo}b)}\quad{\Large $\dfrac{5}{4}$}\end{tcolorbox}\par\vspace{2.5mm}
\begin{tcolorbox}[enhanced,nobeforeafter,width=7.1cm,height=2.1cm,valign=center,colback=cgeo!6,colframe=cgeo,arc=2.5mm,boxrule=0.9pt,left=3mm]{\Large\bfseries\color{cgeo}c)}\quad{\Large $\dfrac{1}{5}$}\end{tcolorbox}\hfill\begin{tcolorbox}[enhanced,nobeforeafter,width=7.1cm,height=2.1cm,valign=center,colback=cgeo!6,colframe=cgeo,arc=2.5mm,boxrule=0.9pt,left=3mm]{\Large\bfseries\color{cgeo}d)}\quad{\Large $\dfrac{4}{5}$}\end{tcolorbox}
\end{tcolorbox}
\newpage
\begin{tcolorbox}[enhanced,width=16cm,height=14cm,colback=white,colframe=cgeo,boxrule=2pt,arc=0mm,
left=6mm,right=6mm,top=0mm,bottom=4mm,sharp corners,
overlay={\fill[cgeo] (frame.north west) rectangle ([yshift=-1.25cm]frame.north east);
\node[anchor=west,text=white,font=\sffamily\bfseries\LARGE] at ([xshift=5mm,yshift=-0.62cm]frame.north west) {Question 59};
\node[anchor=east,text=white,font=\sffamily\itshape\large] at ([xshift=-5mm,yshift=-0.62cm]frame.north east) {Géométrie};}]
\vspace*{1.75cm}
\begin{minipage}[t][6.9cm][c]{\linewidth}\LARGE\centering
\begin{minipage}{0.95\linewidth}\raggedright\linespread{1.25}\selectfont Le triangle $ABC$ est rectangle en $A$, avec $BC=10$ cm et $\cos\widehat{ABC}=0{,}6$. Combien mesure $AB$~?\end{minipage}
\end{minipage}

\vfill
\begin{tcolorbox}[enhanced,nobeforeafter,width=7.1cm,height=2.1cm,valign=center,colback=cgeo!6,colframe=cgeo,arc=2.5mm,boxrule=0.9pt,left=3mm]{\Large\bfseries\color{cgeo}a)}\quad{\Large $4~\text{cm}$}\end{tcolorbox}\hfill\begin{tcolorbox}[enhanced,nobeforeafter,width=7.1cm,height=2.1cm,valign=center,colback=cgeo!6,colframe=cgeo,arc=2.5mm,boxrule=0.9pt,left=3mm]{\Large\bfseries\color{cgeo}b)}\quad{\Large $7~\text{cm}$}\end{tcolorbox}\par\vspace{2.5mm}
\begin{tcolorbox}[enhanced,nobeforeafter,width=7.1cm,height=2.1cm,valign=center,colback=cgeo!6,colframe=cgeo,arc=2.5mm,boxrule=0.9pt,left=3mm]{\Large\bfseries\color{cgeo}c)}\quad{\Large $3~\text{cm}$}\end{tcolorbox}\hfill\begin{tcolorbox}[enhanced,nobeforeafter,width=7.1cm,height=2.1cm,valign=center,colback=cgeo!6,colframe=cgeo,arc=2.5mm,boxrule=0.9pt,left=3mm]{\Large\bfseries\color{cgeo}d)}\quad{\Large $6~\text{cm}$}\end{tcolorbox}
\end{tcolorbox}
\newpage
\begin{tcolorbox}[enhanced,width=16cm,height=14cm,colback=white,colframe=cgeo,boxrule=2pt,arc=0mm,
left=6mm,right=6mm,top=0mm,bottom=4mm,sharp corners,
overlay={\fill[cgeo] (frame.north west) rectangle ([yshift=-1.25cm]frame.north east);
\node[anchor=west,text=white,font=\sffamily\bfseries\LARGE] at ([xshift=5mm,yshift=-0.62cm]frame.north west) {Question 60};
\node[anchor=east,text=white,font=\sffamily\itshape\large] at ([xshift=-5mm,yshift=-0.62cm]frame.north east) {Géométrie};}]
\vspace*{1.75cm}
\begin{minipage}[t][6.9cm][c]{\linewidth}\LARGE\centering
\begin{minipage}{0.95\linewidth}\raggedright\linespread{1.25}\selectfont Le triangle $ABC$ est rectangle en $A$, avec $BC=8$ cm et $\cos\widehat{ABC}=0{,}5$. Combien mesure $AB$~?\end{minipage}
\end{minipage}

\vfill
\begin{tcolorbox}[enhanced,nobeforeafter,width=7.1cm,height=2.1cm,valign=center,colback=cgeo!6,colframe=cgeo,arc=2.5mm,boxrule=0.9pt,left=3mm]{\Large\bfseries\color{cgeo}a)}\quad{\Large $8~\text{cm}$}\end{tcolorbox}\hfill\begin{tcolorbox}[enhanced,nobeforeafter,width=7.1cm,height=2.1cm,valign=center,colback=cgeo!6,colframe=cgeo,arc=2.5mm,boxrule=0.9pt,left=3mm]{\Large\bfseries\color{cgeo}b)}\quad{\Large $5~\text{cm}$}\end{tcolorbox}\par\vspace{2.5mm}
\begin{tcolorbox}[enhanced,nobeforeafter,width=7.1cm,height=2.1cm,valign=center,colback=cgeo!6,colframe=cgeo,arc=2.5mm,boxrule=0.9pt,left=3mm]{\Large\bfseries\color{cgeo}c)}\quad{\Large $2~\text{cm}$}\end{tcolorbox}\hfill\begin{tcolorbox}[enhanced,nobeforeafter,width=7.1cm,height=2.1cm,valign=center,colback=cgeo!6,colframe=cgeo,arc=2.5mm,boxrule=0.9pt,left=3mm]{\Large\bfseries\color{cgeo}d)}\quad{\Large $4~\text{cm}$}\end{tcolorbox}
\end{tcolorbox}
\newpage
\begin{tcolorbox}[enhanced,width=16cm,height=14cm,colback=white,colframe=cgeo,boxrule=2pt,arc=0mm,
left=6mm,right=6mm,top=0mm,bottom=4mm,sharp corners,
overlay={\fill[cgeo] (frame.north west) rectangle ([yshift=-1.25cm]frame.north east);
\node[anchor=west,text=white,font=\sffamily\bfseries\LARGE] at ([xshift=5mm,yshift=-0.62cm]frame.north west) {Question 61};
\node[anchor=east,text=white,font=\sffamily\itshape\large] at ([xshift=-5mm,yshift=-0.62cm]frame.north east) {Géométrie};}]
\vspace*{1.75cm}
\begin{minipage}[t][6.9cm][c]{\linewidth}\Large\centering
\begin{minipage}{0.95\linewidth}\raggedright\linespread{1.25}\selectfont Dans un repère, $A(1\,;\,2)$, $B(4\,;\,3)$ et $C(0\,;\,-1)$. Quelles sont les coordonnées de l'image de $C$ par la translation qui transforme $A$ en $B$~?\end{minipage}
\end{minipage}

\vfill
\begin{tcolorbox}[enhanced,nobeforeafter,width=7.1cm,height=2.1cm,valign=center,colback=cgeo!6,colframe=cgeo,arc=2.5mm,boxrule=0.9pt,left=3mm]{\Large\bfseries\color{cgeo}a)}\quad{\Large $(1\,;\,2)$}\end{tcolorbox}\hfill\begin{tcolorbox}[enhanced,nobeforeafter,width=7.1cm,height=2.1cm,valign=center,colback=cgeo!6,colframe=cgeo,arc=2.5mm,boxrule=0.9pt,left=3mm]{\Large\bfseries\color{cgeo}b)}\quad{\Large $(-3\,;\,-2)$}\end{tcolorbox}\par\vspace{2.5mm}
\begin{tcolorbox}[enhanced,nobeforeafter,width=7.1cm,height=2.1cm,valign=center,colback=cgeo!6,colframe=cgeo,arc=2.5mm,boxrule=0.9pt,left=3mm]{\Large\bfseries\color{cgeo}c)}\quad{\Large $(3\,;\,0)$}\end{tcolorbox}\hfill\begin{tcolorbox}[enhanced,nobeforeafter,width=7.1cm,height=2.1cm,valign=center,colback=cgeo!6,colframe=cgeo,arc=2.5mm,boxrule=0.9pt,left=3mm]{\Large\bfseries\color{cgeo}d)}\quad{\Large $(4\,;\,2)$}\end{tcolorbox}
\end{tcolorbox}
\newpage
\begin{tcolorbox}[enhanced,width=16cm,height=14cm,colback=white,colframe=cgeo,boxrule=2pt,arc=0mm,
left=6mm,right=6mm,top=0mm,bottom=4mm,sharp corners,
overlay={\fill[cgeo] (frame.north west) rectangle ([yshift=-1.25cm]frame.north east);
\node[anchor=west,text=white,font=\sffamily\bfseries\LARGE] at ([xshift=5mm,yshift=-0.62cm]frame.north west) {Question 62};
\node[anchor=east,text=white,font=\sffamily\itshape\large] at ([xshift=-5mm,yshift=-0.62cm]frame.north east) {Géométrie};}]
\vspace*{1.75cm}
\begin{minipage}[t][6.9cm][c]{\linewidth}\Large\centering
\begin{minipage}{0.95\linewidth}\raggedright\linespread{1.25}\selectfont Dans un repère, $A(-2\,;\,1)$, $B(1\,;\,5)$ et $C(3\,;\,-2)$. Quelles sont les coordonnées de l'image de $C$ par la translation qui transforme $A$ en $B$~?\end{minipage}
\end{minipage}

\vfill
\begin{tcolorbox}[enhanced,nobeforeafter,width=7.1cm,height=2.1cm,valign=center,colback=cgeo!6,colframe=cgeo,arc=2.5mm,boxrule=0.9pt,left=3mm]{\Large\bfseries\color{cgeo}a)}\quad{\Large $(0\,;\,-6)$}\end{tcolorbox}\hfill\begin{tcolorbox}[enhanced,nobeforeafter,width=7.1cm,height=2.1cm,valign=center,colback=cgeo!6,colframe=cgeo,arc=2.5mm,boxrule=0.9pt,left=3mm]{\Large\bfseries\color{cgeo}b)}\quad{\Large $(6\,;\,2)$}\end{tcolorbox}\par\vspace{2.5mm}
\begin{tcolorbox}[enhanced,nobeforeafter,width=7.1cm,height=2.1cm,valign=center,colback=cgeo!6,colframe=cgeo,arc=2.5mm,boxrule=0.9pt,left=3mm]{\Large\bfseries\color{cgeo}c)}\quad{\Large $(4\,;\,3)$}\end{tcolorbox}\hfill\begin{tcolorbox}[enhanced,nobeforeafter,width=7.1cm,height=2.1cm,valign=center,colback=cgeo!6,colframe=cgeo,arc=2.5mm,boxrule=0.9pt,left=3mm]{\Large\bfseries\color{cgeo}d)}\quad{\Large $(7\,;\,1)$}\end{tcolorbox}
\end{tcolorbox}
\newpage
\begin{tcolorbox}[enhanced,width=16cm,height=14cm,colback=white,colframe=cgeo,boxrule=2pt,arc=0mm,
left=6mm,right=6mm,top=0mm,bottom=4mm,sharp corners,
overlay={\fill[cgeo] (frame.north west) rectangle ([yshift=-1.25cm]frame.north east);
\node[anchor=west,text=white,font=\sffamily\bfseries\LARGE] at ([xshift=5mm,yshift=-0.62cm]frame.north west) {Question 63};
\node[anchor=east,text=white,font=\sffamily\itshape\large] at ([xshift=-5mm,yshift=-0.62cm]frame.north east) {Géométrie};}]
\vspace*{1.75cm}
\begin{minipage}[t][6.9cm][c]{\linewidth}\LARGE\centering
\begin{minipage}{0.95\linewidth}\raggedright\linespread{1.25}\selectfont Une figure est agrandie à l'échelle $2$ (les longueurs sont multipliées par $2$). Par combien son aire est-elle multipliée~?\end{minipage}
\end{minipage}

\vfill
\begin{tcolorbox}[enhanced,nobeforeafter,width=7.1cm,height=2.1cm,valign=center,colback=cgeo!6,colframe=cgeo,arc=2.5mm,boxrule=0.9pt,left=3mm]{\Large\bfseries\color{cgeo}a)}\quad{\Large $2$}\end{tcolorbox}\hfill\begin{tcolorbox}[enhanced,nobeforeafter,width=7.1cm,height=2.1cm,valign=center,colback=cgeo!6,colframe=cgeo,arc=2.5mm,boxrule=0.9pt,left=3mm]{\Large\bfseries\color{cgeo}b)}\quad{\Large $8$}\end{tcolorbox}\par\vspace{2.5mm}
\begin{tcolorbox}[enhanced,nobeforeafter,width=7.1cm,height=2.1cm,valign=center,colback=cgeo!6,colframe=cgeo,arc=2.5mm,boxrule=0.9pt,left=3mm]{\Large\bfseries\color{cgeo}c)}\quad{\Large $4$}\end{tcolorbox}\hfill\begin{tcolorbox}[enhanced,nobeforeafter,width=7.1cm,height=2.1cm,valign=center,colback=cgeo!6,colframe=cgeo,arc=2.5mm,boxrule=0.9pt,left=3mm]{\Large\bfseries\color{cgeo}d)}\quad{\Large $5$}\end{tcolorbox}
\end{tcolorbox}
\newpage
\begin{tcolorbox}[enhanced,width=16cm,height=14cm,colback=white,colframe=cgeo,boxrule=2pt,arc=0mm,
left=6mm,right=6mm,top=0mm,bottom=4mm,sharp corners,
overlay={\fill[cgeo] (frame.north west) rectangle ([yshift=-1.25cm]frame.north east);
\node[anchor=west,text=white,font=\sffamily\bfseries\LARGE] at ([xshift=5mm,yshift=-0.62cm]frame.north west) {Question 64};
\node[anchor=east,text=white,font=\sffamily\itshape\large] at ([xshift=-5mm,yshift=-0.62cm]frame.north east) {Géométrie};}]
\vspace*{1.75cm}
\begin{minipage}[t][6.9cm][c]{\linewidth}\LARGE\centering
\begin{minipage}{0.95\linewidth}\raggedright\linespread{1.25}\selectfont Une figure est agrandie à l'échelle $3$ (les longueurs sont multipliées par $3$). Par combien son aire est-elle multipliée~?\end{minipage}
\end{minipage}

\vfill
\begin{tcolorbox}[enhanced,nobeforeafter,width=7.1cm,height=2.1cm,valign=center,colback=cgeo!6,colframe=cgeo,arc=2.5mm,boxrule=0.9pt,left=3mm]{\Large\bfseries\color{cgeo}a)}\quad{\Large $6$}\end{tcolorbox}\hfill\begin{tcolorbox}[enhanced,nobeforeafter,width=7.1cm,height=2.1cm,valign=center,colback=cgeo!6,colframe=cgeo,arc=2.5mm,boxrule=0.9pt,left=3mm]{\Large\bfseries\color{cgeo}b)}\quad{\Large $9$}\end{tcolorbox}\par\vspace{2.5mm}
\begin{tcolorbox}[enhanced,nobeforeafter,width=7.1cm,height=2.1cm,valign=center,colback=cgeo!6,colframe=cgeo,arc=2.5mm,boxrule=0.9pt,left=3mm]{\Large\bfseries\color{cgeo}c)}\quad{\Large $3$}\end{tcolorbox}\hfill\begin{tcolorbox}[enhanced,nobeforeafter,width=7.1cm,height=2.1cm,valign=center,colback=cgeo!6,colframe=cgeo,arc=2.5mm,boxrule=0.9pt,left=3mm]{\Large\bfseries\color{cgeo}d)}\quad{\Large $27$}\end{tcolorbox}
\end{tcolorbox}
\newpage
\begin{tcolorbox}[enhanced,width=16cm,height=14cm,colback=white,colframe=cdon,boxrule=2pt,arc=0mm,
left=6mm,right=6mm,top=0mm,bottom=4mm,sharp corners,
overlay={\fill[cdon] (frame.north west) rectangle ([yshift=-1.25cm]frame.north east);
\node[anchor=west,text=white,font=\sffamily\bfseries\LARGE] at ([xshift=5mm,yshift=-0.62cm]frame.north west) {Question 65};
\node[anchor=east,text=white,font=\sffamily\itshape\large] at ([xshift=-5mm,yshift=-0.62cm]frame.north east) {Proportionnalité, statistiques, probabilités};}]
\vspace*{1.75cm}
\begin{minipage}[t][6.9cm][c]{\linewidth}\LARGE\centering
\begin{minipage}{0.95\linewidth}\raggedright\linespread{1.25}\selectfont Ce tableau est un tableau de proportionnalité~: $\begin{array}{|c|c|}\hline 4 & 6\\\hline 10 &~?\\\hline\end{array}$. Quel nombre remplace «~?~»~?\end{minipage}
\end{minipage}

\vfill
\begin{tcolorbox}[enhanced,nobeforeafter,width=7.1cm,height=2.1cm,valign=center,colback=cdon!6,colframe=cdon,arc=2.5mm,boxrule=0.9pt,left=3mm]{\Large\bfseries\color{cdon}a)}\quad{\Large $\dfrac{12}{5}$}\end{tcolorbox}\hfill\begin{tcolorbox}[enhanced,nobeforeafter,width=7.1cm,height=2.1cm,valign=center,colback=cdon!6,colframe=cdon,arc=2.5mm,boxrule=0.9pt,left=3mm]{\Large\bfseries\color{cdon}b)}\quad{\Large $12$}\end{tcolorbox}\par\vspace{2.5mm}
\begin{tcolorbox}[enhanced,nobeforeafter,width=7.1cm,height=2.1cm,valign=center,colback=cdon!6,colframe=cdon,arc=2.5mm,boxrule=0.9pt,left=3mm]{\Large\bfseries\color{cdon}c)}\quad{\Large $15$}\end{tcolorbox}\hfill\begin{tcolorbox}[enhanced,nobeforeafter,width=7.1cm,height=2.1cm,valign=center,colback=cdon!6,colframe=cdon,arc=2.5mm,boxrule=0.9pt,left=3mm]{\Large\bfseries\color{cdon}d)}\quad{\Large $60$}\end{tcolorbox}
\end{tcolorbox}
\newpage
\begin{tcolorbox}[enhanced,width=16cm,height=14cm,colback=white,colframe=cdon,boxrule=2pt,arc=0mm,
left=6mm,right=6mm,top=0mm,bottom=4mm,sharp corners,
overlay={\fill[cdon] (frame.north west) rectangle ([yshift=-1.25cm]frame.north east);
\node[anchor=west,text=white,font=\sffamily\bfseries\LARGE] at ([xshift=5mm,yshift=-0.62cm]frame.north west) {Question 66};
\node[anchor=east,text=white,font=\sffamily\itshape\large] at ([xshift=-5mm,yshift=-0.62cm]frame.north east) {Proportionnalité, statistiques, probabilités};}]
\vspace*{1.75cm}
\begin{minipage}[t][6.9cm][c]{\linewidth}\LARGE\centering
\begin{minipage}{0.95\linewidth}\raggedright\linespread{1.25}\selectfont Ce tableau est un tableau de proportionnalité~: $\begin{array}{|c|c|}\hline 3 & 12\\\hline 7 &~?\\\hline\end{array}$. Quel nombre remplace «~?~»~?\end{minipage}
\end{minipage}

\vfill
\begin{tcolorbox}[enhanced,nobeforeafter,width=7.1cm,height=2.1cm,valign=center,colback=cdon!6,colframe=cdon,arc=2.5mm,boxrule=0.9pt,left=3mm]{\Large\bfseries\color{cdon}a)}\quad{\Large $16$}\end{tcolorbox}\hfill\begin{tcolorbox}[enhanced,nobeforeafter,width=7.1cm,height=2.1cm,valign=center,colback=cdon!6,colframe=cdon,arc=2.5mm,boxrule=0.9pt,left=3mm]{\Large\bfseries\color{cdon}b)}\quad{\Large $84$}\end{tcolorbox}\par\vspace{2.5mm}
\begin{tcolorbox}[enhanced,nobeforeafter,width=7.1cm,height=2.1cm,valign=center,colback=cdon!6,colframe=cdon,arc=2.5mm,boxrule=0.9pt,left=3mm]{\Large\bfseries\color{cdon}c)}\quad{\Large $28$}\end{tcolorbox}\hfill\begin{tcolorbox}[enhanced,nobeforeafter,width=7.1cm,height=2.1cm,valign=center,colback=cdon!6,colframe=cdon,arc=2.5mm,boxrule=0.9pt,left=3mm]{\Large\bfseries\color{cdon}d)}\quad{\Large $\dfrac{36}{7}$}\end{tcolorbox}
\end{tcolorbox}
\newpage
\begin{tcolorbox}[enhanced,width=16cm,height=14cm,colback=white,colframe=cdon,boxrule=2pt,arc=0mm,
left=6mm,right=6mm,top=0mm,bottom=4mm,sharp corners,
overlay={\fill[cdon] (frame.north west) rectangle ([yshift=-1.25cm]frame.north east);
\node[anchor=west,text=white,font=\sffamily\bfseries\LARGE] at ([xshift=5mm,yshift=-0.62cm]frame.north west) {Question 67};
\node[anchor=east,text=white,font=\sffamily\itshape\large] at ([xshift=-5mm,yshift=-0.62cm]frame.north east) {Proportionnalité, statistiques, probabilités};}]
\vspace*{1.75cm}
\begin{minipage}[t][6.9cm][c]{\linewidth}\LARGE\centering
\begin{minipage}{0.95\linewidth}\raggedright\linespread{1.25}\selectfont Que vaut $15\,\%$ de $80$~?\end{minipage}
\end{minipage}

\vfill
\begin{tcolorbox}[enhanced,nobeforeafter,width=7.1cm,height=2.1cm,valign=center,colback=cdon!6,colframe=cdon,arc=2.5mm,boxrule=0.9pt,left=3mm]{\Large\bfseries\color{cdon}a)}\quad{\Large $12$}\end{tcolorbox}\hfill\begin{tcolorbox}[enhanced,nobeforeafter,width=7.1cm,height=2.1cm,valign=center,colback=cdon!6,colframe=cdon,arc=2.5mm,boxrule=0.9pt,left=3mm]{\Large\bfseries\color{cdon}b)}\quad{\Large $120$}\end{tcolorbox}\par\vspace{2.5mm}
\begin{tcolorbox}[enhanced,nobeforeafter,width=7.1cm,height=2.1cm,valign=center,colback=cdon!6,colframe=cdon,arc=2.5mm,boxrule=0.9pt,left=3mm]{\Large\bfseries\color{cdon}c)}\quad{\Large $27$}\end{tcolorbox}\hfill\begin{tcolorbox}[enhanced,nobeforeafter,width=7.1cm,height=2.1cm,valign=center,colback=cdon!6,colframe=cdon,arc=2.5mm,boxrule=0.9pt,left=3mm]{\Large\bfseries\color{cdon}d)}\quad{\Large $68$}\end{tcolorbox}
\end{tcolorbox}
\newpage
\begin{tcolorbox}[enhanced,width=16cm,height=14cm,colback=white,colframe=cdon,boxrule=2pt,arc=0mm,
left=6mm,right=6mm,top=0mm,bottom=4mm,sharp corners,
overlay={\fill[cdon] (frame.north west) rectangle ([yshift=-1.25cm]frame.north east);
\node[anchor=west,text=white,font=\sffamily\bfseries\LARGE] at ([xshift=5mm,yshift=-0.62cm]frame.north west) {Question 68};
\node[anchor=east,text=white,font=\sffamily\itshape\large] at ([xshift=-5mm,yshift=-0.62cm]frame.north east) {Proportionnalité, statistiques, probabilités};}]
\vspace*{1.75cm}
\begin{minipage}[t][6.9cm][c]{\linewidth}\LARGE\centering
\begin{minipage}{0.95\linewidth}\raggedright\linespread{1.25}\selectfont Que vaut $30\,\%$ de $250$~?\end{minipage}
\end{minipage}

\vfill
\begin{tcolorbox}[enhanced,nobeforeafter,width=7.1cm,height=2.1cm,valign=center,colback=cdon!6,colframe=cdon,arc=2.5mm,boxrule=0.9pt,left=3mm]{\Large\bfseries\color{cdon}a)}\quad{\Large $750$}\end{tcolorbox}\hfill\begin{tcolorbox}[enhanced,nobeforeafter,width=7.1cm,height=2.1cm,valign=center,colback=cdon!6,colframe=cdon,arc=2.5mm,boxrule=0.9pt,left=3mm]{\Large\bfseries\color{cdon}b)}\quad{\Large $175$}\end{tcolorbox}\par\vspace{2.5mm}
\begin{tcolorbox}[enhanced,nobeforeafter,width=7.1cm,height=2.1cm,valign=center,colback=cdon!6,colframe=cdon,arc=2.5mm,boxrule=0.9pt,left=3mm]{\Large\bfseries\color{cdon}c)}\quad{\Large $75$}\end{tcolorbox}\hfill\begin{tcolorbox}[enhanced,nobeforeafter,width=7.1cm,height=2.1cm,valign=center,colback=cdon!6,colframe=cdon,arc=2.5mm,boxrule=0.9pt,left=3mm]{\Large\bfseries\color{cdon}d)}\quad{\Large $105$}\end{tcolorbox}
\end{tcolorbox}
\newpage
\begin{tcolorbox}[enhanced,width=16cm,height=14cm,colback=white,colframe=cdon,boxrule=2pt,arc=0mm,
left=6mm,right=6mm,top=0mm,bottom=4mm,sharp corners,
overlay={\fill[cdon] (frame.north west) rectangle ([yshift=-1.25cm]frame.north east);
\node[anchor=west,text=white,font=\sffamily\bfseries\LARGE] at ([xshift=5mm,yshift=-0.62cm]frame.north west) {Question 69};
\node[anchor=east,text=white,font=\sffamily\itshape\large] at ([xshift=-5mm,yshift=-0.62cm]frame.north east) {Proportionnalité, statistiques, probabilités};}]
\vspace*{1.75cm}
\begin{minipage}[t][6.9cm][c]{\linewidth}\LARGE\centering
\begin{minipage}{0.95\linewidth}\raggedright\linespread{1.25}\selectfont Dans une classe de $25$ élèves, $10$ sont demi-pensionnaires. Quel pourcentage des élèves cela représente-t-il~?\end{minipage}
\end{minipage}

\vfill
\begin{tcolorbox}[enhanced,nobeforeafter,width=7.1cm,height=2.1cm,valign=center,colback=cdon!6,colframe=cdon,arc=2.5mm,boxrule=0.9pt,left=3mm]{\Large\bfseries\color{cdon}a)}\quad{\Large $40\,\%$}\end{tcolorbox}\hfill\begin{tcolorbox}[enhanced,nobeforeafter,width=7.1cm,height=2.1cm,valign=center,colback=cdon!6,colframe=cdon,arc=2.5mm,boxrule=0.9pt,left=3mm]{\Large\bfseries\color{cdon}b)}\quad{\Large $60\,\%$}\end{tcolorbox}\par\vspace{2.5mm}
\begin{tcolorbox}[enhanced,nobeforeafter,width=7.1cm,height=2.1cm,valign=center,colback=cdon!6,colframe=cdon,arc=2.5mm,boxrule=0.9pt,left=3mm]{\Large\bfseries\color{cdon}c)}\quad{\Large $10\,\%$}\end{tcolorbox}\hfill\begin{tcolorbox}[enhanced,nobeforeafter,width=7.1cm,height=2.1cm,valign=center,colback=cdon!6,colframe=cdon,arc=2.5mm,boxrule=0.9pt,left=3mm]{\Large\bfseries\color{cdon}d)}\quad{\Large $4\,\%$}\end{tcolorbox}
\end{tcolorbox}
\newpage
\begin{tcolorbox}[enhanced,width=16cm,height=14cm,colback=white,colframe=cdon,boxrule=2pt,arc=0mm,
left=6mm,right=6mm,top=0mm,bottom=4mm,sharp corners,
overlay={\fill[cdon] (frame.north west) rectangle ([yshift=-1.25cm]frame.north east);
\node[anchor=west,text=white,font=\sffamily\bfseries\LARGE] at ([xshift=5mm,yshift=-0.62cm]frame.north west) {Question 70};
\node[anchor=east,text=white,font=\sffamily\itshape\large] at ([xshift=-5mm,yshift=-0.62cm]frame.north east) {Proportionnalité, statistiques, probabilités};}]
\vspace*{1.75cm}
\begin{minipage}[t][6.9cm][c]{\linewidth}\LARGE\centering
\begin{minipage}{0.95\linewidth}\raggedright\linespread{1.25}\selectfont Dans une classe de $20$ élèves, $7$ sont demi-pensionnaires. Quel pourcentage des élèves cela représente-t-il~?\end{minipage}
\end{minipage}

\vfill
\begin{tcolorbox}[enhanced,nobeforeafter,width=7.1cm,height=2.1cm,valign=center,colback=cdon!6,colframe=cdon,arc=2.5mm,boxrule=0.9pt,left=3mm]{\Large\bfseries\color{cdon}a)}\quad{\Large $7\,\%$}\end{tcolorbox}\hfill\begin{tcolorbox}[enhanced,nobeforeafter,width=7.1cm,height=2.1cm,valign=center,colback=cdon!6,colframe=cdon,arc=2.5mm,boxrule=0.9pt,left=3mm]{\Large\bfseries\color{cdon}b)}\quad{\Large $35\,\%$}\end{tcolorbox}\par\vspace{2.5mm}
\begin{tcolorbox}[enhanced,nobeforeafter,width=7.1cm,height=2.1cm,valign=center,colback=cdon!6,colframe=cdon,arc=2.5mm,boxrule=0.9pt,left=3mm]{\Large\bfseries\color{cdon}c)}\quad{\Large $3{,}5\,\%$}\end{tcolorbox}\hfill\begin{tcolorbox}[enhanced,nobeforeafter,width=7.1cm,height=2.1cm,valign=center,colback=cdon!6,colframe=cdon,arc=2.5mm,boxrule=0.9pt,left=3mm]{\Large\bfseries\color{cdon}d)}\quad{\Large $65\,\%$}\end{tcolorbox}
\end{tcolorbox}
\newpage
\begin{tcolorbox}[enhanced,width=16cm,height=14cm,colback=white,colframe=cdon,boxrule=2pt,arc=0mm,
left=6mm,right=6mm,top=0mm,bottom=4mm,sharp corners,
overlay={\fill[cdon] (frame.north west) rectangle ([yshift=-1.25cm]frame.north east);
\node[anchor=west,text=white,font=\sffamily\bfseries\LARGE] at ([xshift=5mm,yshift=-0.62cm]frame.north west) {Question 71};
\node[anchor=east,text=white,font=\sffamily\itshape\large] at ([xshift=-5mm,yshift=-0.62cm]frame.north east) {Proportionnalité, statistiques, probabilités};}]
\vspace*{1.75cm}
\begin{minipage}[t][6.9cm][c]{\linewidth}\LARGE\centering
\begin{minipage}{0.95\linewidth}\raggedright\linespread{1.25}\selectfont Quelle est la moyenne de la série $12\ ;\ 15\ ;\ 9\ ;\ 14$~?\end{minipage}
\end{minipage}

\vfill
\begin{tcolorbox}[enhanced,nobeforeafter,width=7.1cm,height=2.1cm,valign=center,colback=cdon!6,colframe=cdon,arc=2.5mm,boxrule=0.9pt,left=3mm]{\Large\bfseries\color{cdon}a)}\quad{\Large $12{,}5$}\end{tcolorbox}\hfill\begin{tcolorbox}[enhanced,nobeforeafter,width=7.1cm,height=2.1cm,valign=center,colback=cdon!6,colframe=cdon,arc=2.5mm,boxrule=0.9pt,left=3mm]{\Large\bfseries\color{cdon}b)}\quad{\Large $50$}\end{tcolorbox}\par\vspace{2.5mm}
\begin{tcolorbox}[enhanced,nobeforeafter,width=7.1cm,height=2.1cm,valign=center,colback=cdon!6,colframe=cdon,arc=2.5mm,boxrule=0.9pt,left=3mm]{\Large\bfseries\color{cdon}c)}\quad{\Large $12$}\end{tcolorbox}\hfill\begin{tcolorbox}[enhanced,nobeforeafter,width=7.1cm,height=2.1cm,valign=center,colback=cdon!6,colframe=cdon,arc=2.5mm,boxrule=0.9pt,left=3mm]{\Large\bfseries\color{cdon}d)}\quad{\Large $13{,}5$}\end{tcolorbox}
\end{tcolorbox}
\newpage
\begin{tcolorbox}[enhanced,width=16cm,height=14cm,colback=white,colframe=cdon,boxrule=2pt,arc=0mm,
left=6mm,right=6mm,top=0mm,bottom=4mm,sharp corners,
overlay={\fill[cdon] (frame.north west) rectangle ([yshift=-1.25cm]frame.north east);
\node[anchor=west,text=white,font=\sffamily\bfseries\LARGE] at ([xshift=5mm,yshift=-0.62cm]frame.north west) {Question 72};
\node[anchor=east,text=white,font=\sffamily\itshape\large] at ([xshift=-5mm,yshift=-0.62cm]frame.north east) {Proportionnalité, statistiques, probabilités};}]
\vspace*{1.75cm}
\begin{minipage}[t][6.9cm][c]{\linewidth}\LARGE\centering
\begin{minipage}{0.95\linewidth}\raggedright\linespread{1.25}\selectfont Quelle est la moyenne de la série $8\ ;\ 11\ ;\ 17$~?\end{minipage}
\end{minipage}

\vfill
\begin{tcolorbox}[enhanced,nobeforeafter,width=7.1cm,height=2.1cm,valign=center,colback=cdon!6,colframe=cdon,arc=2.5mm,boxrule=0.9pt,left=3mm]{\Large\bfseries\color{cdon}a)}\quad{\Large $13$}\end{tcolorbox}\hfill\begin{tcolorbox}[enhanced,nobeforeafter,width=7.1cm,height=2.1cm,valign=center,colback=cdon!6,colframe=cdon,arc=2.5mm,boxrule=0.9pt,left=3mm]{\Large\bfseries\color{cdon}b)}\quad{\Large $12{,}5$}\end{tcolorbox}\par\vspace{2.5mm}
\begin{tcolorbox}[enhanced,nobeforeafter,width=7.1cm,height=2.1cm,valign=center,colback=cdon!6,colframe=cdon,arc=2.5mm,boxrule=0.9pt,left=3mm]{\Large\bfseries\color{cdon}c)}\quad{\Large $36$}\end{tcolorbox}\hfill\begin{tcolorbox}[enhanced,nobeforeafter,width=7.1cm,height=2.1cm,valign=center,colback=cdon!6,colframe=cdon,arc=2.5mm,boxrule=0.9pt,left=3mm]{\Large\bfseries\color{cdon}d)}\quad{\Large $12$}\end{tcolorbox}
\end{tcolorbox}
\newpage
\begin{tcolorbox}[enhanced,width=16cm,height=14cm,colback=white,colframe=cdon,boxrule=2pt,arc=0mm,
left=6mm,right=6mm,top=0mm,bottom=4mm,sharp corners,
overlay={\fill[cdon] (frame.north west) rectangle ([yshift=-1.25cm]frame.north east);
\node[anchor=west,text=white,font=\sffamily\bfseries\LARGE] at ([xshift=5mm,yshift=-0.62cm]frame.north west) {Question 73};
\node[anchor=east,text=white,font=\sffamily\itshape\large] at ([xshift=-5mm,yshift=-0.62cm]frame.north east) {Proportionnalité, statistiques, probabilités};}]
\vspace*{1.75cm}
\begin{minipage}[t][6.9cm][c]{\linewidth}\LARGE\centering
\begin{minipage}{0.95\linewidth}\raggedright\linespread{1.25}\selectfont Quelle est la médiane de la série $4\ ;\ 9\ ;\ 2\ ;\ 7\ ;\ 5$~?\end{minipage}
\end{minipage}

\vfill
\begin{tcolorbox}[enhanced,nobeforeafter,width=7.1cm,height=2.1cm,valign=center,colback=cdon!6,colframe=cdon,arc=2.5mm,boxrule=0.9pt,left=3mm]{\Large\bfseries\color{cdon}a)}\quad{\Large $3$}\end{tcolorbox}\hfill\begin{tcolorbox}[enhanced,nobeforeafter,width=7.1cm,height=2.1cm,valign=center,colback=cdon!6,colframe=cdon,arc=2.5mm,boxrule=0.9pt,left=3mm]{\Large\bfseries\color{cdon}b)}\quad{\Large $5{,}4$}\end{tcolorbox}\par\vspace{2.5mm}
\begin{tcolorbox}[enhanced,nobeforeafter,width=7.1cm,height=2.1cm,valign=center,colback=cdon!6,colframe=cdon,arc=2.5mm,boxrule=0.9pt,left=3mm]{\Large\bfseries\color{cdon}c)}\quad{\Large $5$}\end{tcolorbox}\hfill\begin{tcolorbox}[enhanced,nobeforeafter,width=7.1cm,height=2.1cm,valign=center,colback=cdon!6,colframe=cdon,arc=2.5mm,boxrule=0.9pt,left=3mm]{\Large\bfseries\color{cdon}d)}\quad{\Large $2$}\end{tcolorbox}
\end{tcolorbox}
\newpage
\begin{tcolorbox}[enhanced,width=16cm,height=14cm,colback=white,colframe=cdon,boxrule=2pt,arc=0mm,
left=6mm,right=6mm,top=0mm,bottom=4mm,sharp corners,
overlay={\fill[cdon] (frame.north west) rectangle ([yshift=-1.25cm]frame.north east);
\node[anchor=west,text=white,font=\sffamily\bfseries\LARGE] at ([xshift=5mm,yshift=-0.62cm]frame.north west) {Question 74};
\node[anchor=east,text=white,font=\sffamily\itshape\large] at ([xshift=-5mm,yshift=-0.62cm]frame.north east) {Proportionnalité, statistiques, probabilités};}]
\vspace*{1.75cm}
\begin{minipage}[t][6.9cm][c]{\linewidth}\LARGE\centering
\begin{minipage}{0.95\linewidth}\raggedright\linespread{1.25}\selectfont Quelle est la médiane de la série $12\ ;\ 3\ ;\ 8\ ;\ 15\ ;\ 10\ ;\ 6$~?\end{minipage}
\end{minipage}

\vfill
\begin{tcolorbox}[enhanced,nobeforeafter,width=7.1cm,height=2.1cm,valign=center,colback=cdon!6,colframe=cdon,arc=2.5mm,boxrule=0.9pt,left=3mm]{\Large\bfseries\color{cdon}a)}\quad{\Large $3{,}5$}\end{tcolorbox}\hfill\begin{tcolorbox}[enhanced,nobeforeafter,width=7.1cm,height=2.1cm,valign=center,colback=cdon!6,colframe=cdon,arc=2.5mm,boxrule=0.9pt,left=3mm]{\Large\bfseries\color{cdon}b)}\quad{\Large $9$}\end{tcolorbox}\par\vspace{2.5mm}
\begin{tcolorbox}[enhanced,nobeforeafter,width=7.1cm,height=2.1cm,valign=center,colback=cdon!6,colframe=cdon,arc=2.5mm,boxrule=0.9pt,left=3mm]{\Large\bfseries\color{cdon}c)}\quad{\Large $10$}\end{tcolorbox}\hfill\begin{tcolorbox}[enhanced,nobeforeafter,width=7.1cm,height=2.1cm,valign=center,colback=cdon!6,colframe=cdon,arc=2.5mm,boxrule=0.9pt,left=3mm]{\Large\bfseries\color{cdon}d)}\quad{\Large $11{,}5$}\end{tcolorbox}
\end{tcolorbox}
\newpage
\begin{tcolorbox}[enhanced,width=16cm,height=14cm,colback=white,colframe=cdon,boxrule=2pt,arc=0mm,
left=6mm,right=6mm,top=0mm,bottom=4mm,sharp corners,
overlay={\fill[cdon] (frame.north west) rectangle ([yshift=-1.25cm]frame.north east);
\node[anchor=west,text=white,font=\sffamily\bfseries\LARGE] at ([xshift=5mm,yshift=-0.62cm]frame.north west) {Question 75};
\node[anchor=east,text=white,font=\sffamily\itshape\large] at ([xshift=-5mm,yshift=-0.62cm]frame.north east) {Proportionnalité, statistiques, probabilités};}]
\vspace*{1.75cm}
\begin{minipage}[t][6.9cm][c]{\linewidth}\LARGE\centering
\begin{minipage}{0.95\linewidth}\raggedright\linespread{1.25}\selectfont Quelle est l'étendue de la série $12\ ;\ 5\ ;\ 18\ ;\ 9$~?\end{minipage}
\end{minipage}

\vfill
\begin{tcolorbox}[enhanced,nobeforeafter,width=7.1cm,height=2.1cm,valign=center,colback=cdon!6,colframe=cdon,arc=2.5mm,boxrule=0.9pt,left=3mm]{\Large\bfseries\color{cdon}a)}\quad{\Large $13$}\end{tcolorbox}\hfill\begin{tcolorbox}[enhanced,nobeforeafter,width=7.1cm,height=2.1cm,valign=center,colback=cdon!6,colframe=cdon,arc=2.5mm,boxrule=0.9pt,left=3mm]{\Large\bfseries\color{cdon}b)}\quad{\Large $18$}\end{tcolorbox}\par\vspace{2.5mm}
\begin{tcolorbox}[enhanced,nobeforeafter,width=7.1cm,height=2.1cm,valign=center,colback=cdon!6,colframe=cdon,arc=2.5mm,boxrule=0.9pt,left=3mm]{\Large\bfseries\color{cdon}c)}\quad{\Large $23$}\end{tcolorbox}\hfill\begin{tcolorbox}[enhanced,nobeforeafter,width=7.1cm,height=2.1cm,valign=center,colback=cdon!6,colframe=cdon,arc=2.5mm,boxrule=0.9pt,left=3mm]{\Large\bfseries\color{cdon}d)}\quad{\Large $11$}\end{tcolorbox}
\end{tcolorbox}
\newpage
\begin{tcolorbox}[enhanced,width=16cm,height=14cm,colback=white,colframe=cdon,boxrule=2pt,arc=0mm,
left=6mm,right=6mm,top=0mm,bottom=4mm,sharp corners,
overlay={\fill[cdon] (frame.north west) rectangle ([yshift=-1.25cm]frame.north east);
\node[anchor=west,text=white,font=\sffamily\bfseries\LARGE] at ([xshift=5mm,yshift=-0.62cm]frame.north west) {Question 76};
\node[anchor=east,text=white,font=\sffamily\itshape\large] at ([xshift=-5mm,yshift=-0.62cm]frame.north east) {Proportionnalité, statistiques, probabilités};}]
\vspace*{1.75cm}
\begin{minipage}[t][6.9cm][c]{\linewidth}\LARGE\centering
\begin{minipage}{0.95\linewidth}\raggedright\linespread{1.25}\selectfont Quelle est l'étendue de la série $7\ ;\ 15\ ;\ 3\ ;\ 11\ ;\ 6$~?\end{minipage}
\end{minipage}

\vfill
\begin{tcolorbox}[enhanced,nobeforeafter,width=7.1cm,height=2.1cm,valign=center,colback=cdon!6,colframe=cdon,arc=2.5mm,boxrule=0.9pt,left=3mm]{\Large\bfseries\color{cdon}a)}\quad{\Large $18$}\end{tcolorbox}\hfill\begin{tcolorbox}[enhanced,nobeforeafter,width=7.1cm,height=2.1cm,valign=center,colback=cdon!6,colframe=cdon,arc=2.5mm,boxrule=0.9pt,left=3mm]{\Large\bfseries\color{cdon}b)}\quad{\Large $12$}\end{tcolorbox}\par\vspace{2.5mm}
\begin{tcolorbox}[enhanced,nobeforeafter,width=7.1cm,height=2.1cm,valign=center,colback=cdon!6,colframe=cdon,arc=2.5mm,boxrule=0.9pt,left=3mm]{\Large\bfseries\color{cdon}c)}\quad{\Large $15$}\end{tcolorbox}\hfill\begin{tcolorbox}[enhanced,nobeforeafter,width=7.1cm,height=2.1cm,valign=center,colback=cdon!6,colframe=cdon,arc=2.5mm,boxrule=0.9pt,left=3mm]{\Large\bfseries\color{cdon}d)}\quad{\Large $8{,}4$}\end{tcolorbox}
\end{tcolorbox}
\newpage
\begin{tcolorbox}[enhanced,width=16cm,height=14cm,colback=white,colframe=cdon,boxrule=2pt,arc=0mm,
left=6mm,right=6mm,top=0mm,bottom=4mm,sharp corners,
overlay={\fill[cdon] (frame.north west) rectangle ([yshift=-1.25cm]frame.north east);
\node[anchor=west,text=white,font=\sffamily\bfseries\LARGE] at ([xshift=5mm,yshift=-0.62cm]frame.north west) {Question 77};
\node[anchor=east,text=white,font=\sffamily\itshape\large] at ([xshift=-5mm,yshift=-0.62cm]frame.north east) {Proportionnalité, statistiques, probabilités};}]
\vspace*{1.75cm}
\begin{minipage}[t][6.9cm][c]{\linewidth}\Large\centering
\begin{minipage}{0.95\linewidth}\raggedright\linespread{1.25}\selectfont Un sac contient $3$ boules rouges, $5$ boules bleues et $2$ boules vertes. On tire une boule au hasard. Quelle est la probabilité de tirer une boule bleue~?\end{minipage}
\end{minipage}

\vfill
\begin{tcolorbox}[enhanced,nobeforeafter,width=7.1cm,height=2.1cm,valign=center,colback=cdon!6,colframe=cdon,arc=2.5mm,boxrule=0.9pt,left=3mm]{\Large\bfseries\color{cdon}a)}\quad{\Large $\dfrac{1}{3}$}\end{tcolorbox}\hfill\begin{tcolorbox}[enhanced,nobeforeafter,width=7.1cm,height=2.1cm,valign=center,colback=cdon!6,colframe=cdon,arc=2.5mm,boxrule=0.9pt,left=3mm]{\Large\bfseries\color{cdon}b)}\quad{\Large $1$}\end{tcolorbox}\par\vspace{2.5mm}
\begin{tcolorbox}[enhanced,nobeforeafter,width=7.1cm,height=2.1cm,valign=center,colback=cdon!6,colframe=cdon,arc=2.5mm,boxrule=0.9pt,left=3mm]{\Large\bfseries\color{cdon}c)}\quad{\Large $\dfrac{1}{2}$}\end{tcolorbox}\hfill\begin{tcolorbox}[enhanced,nobeforeafter,width=7.1cm,height=2.1cm,valign=center,colback=cdon!6,colframe=cdon,arc=2.5mm,boxrule=0.9pt,left=3mm]{\Large\bfseries\color{cdon}d)}\quad{\Large $\dfrac{1}{5}$}\end{tcolorbox}
\end{tcolorbox}
\newpage
\begin{tcolorbox}[enhanced,width=16cm,height=14cm,colback=white,colframe=cdon,boxrule=2pt,arc=0mm,
left=6mm,right=6mm,top=0mm,bottom=4mm,sharp corners,
overlay={\fill[cdon] (frame.north west) rectangle ([yshift=-1.25cm]frame.north east);
\node[anchor=west,text=white,font=\sffamily\bfseries\LARGE] at ([xshift=5mm,yshift=-0.62cm]frame.north west) {Question 78};
\node[anchor=east,text=white,font=\sffamily\itshape\large] at ([xshift=-5mm,yshift=-0.62cm]frame.north east) {Proportionnalité, statistiques, probabilités};}]
\vspace*{1.75cm}
\begin{minipage}[t][6.9cm][c]{\linewidth}\Large\centering
\begin{minipage}{0.95\linewidth}\raggedright\linespread{1.25}\selectfont Un sac contient $4$ boules rouges, $1$ boule bleue et $5$ boules vertes. On tire une boule au hasard. Quelle est la probabilité de tirer une boule rouge~?\end{minipage}
\end{minipage}

\vfill
\begin{tcolorbox}[enhanced,nobeforeafter,width=7.1cm,height=2.1cm,valign=center,colback=cdon!6,colframe=cdon,arc=2.5mm,boxrule=0.9pt,left=3mm]{\Large\bfseries\color{cdon}a)}\quad{\Large $\dfrac{2}{3}$}\end{tcolorbox}\hfill\begin{tcolorbox}[enhanced,nobeforeafter,width=7.1cm,height=2.1cm,valign=center,colback=cdon!6,colframe=cdon,arc=2.5mm,boxrule=0.9pt,left=3mm]{\Large\bfseries\color{cdon}b)}\quad{\Large $\dfrac{1}{3}$}\end{tcolorbox}\par\vspace{2.5mm}
\begin{tcolorbox}[enhanced,nobeforeafter,width=7.1cm,height=2.1cm,valign=center,colback=cdon!6,colframe=cdon,arc=2.5mm,boxrule=0.9pt,left=3mm]{\Large\bfseries\color{cdon}c)}\quad{\Large $\dfrac{2}{5}$}\end{tcolorbox}\hfill\begin{tcolorbox}[enhanced,nobeforeafter,width=7.1cm,height=2.1cm,valign=center,colback=cdon!6,colframe=cdon,arc=2.5mm,boxrule=0.9pt,left=3mm]{\Large\bfseries\color{cdon}d)}\quad{\Large $\dfrac{3}{5}$}\end{tcolorbox}
\end{tcolorbox}
\newpage
\begin{tcolorbox}[enhanced,width=16cm,height=14cm,colback=white,colframe=cdon,boxrule=2pt,arc=0mm,
left=6mm,right=6mm,top=0mm,bottom=4mm,sharp corners,
overlay={\fill[cdon] (frame.north west) rectangle ([yshift=-1.25cm]frame.north east);
\node[anchor=west,text=white,font=\sffamily\bfseries\LARGE] at ([xshift=5mm,yshift=-0.62cm]frame.north west) {Question 79};
\node[anchor=east,text=white,font=\sffamily\itshape\large] at ([xshift=-5mm,yshift=-0.62cm]frame.north east) {Proportionnalité, statistiques, probabilités};}]
\vspace*{1.75cm}
\begin{minipage}[t][6.9cm][c]{\linewidth}\LARGE\centering
\begin{minipage}{0.95\linewidth}\raggedright\linespread{1.25}\selectfont La probabilité qu'un événement $A$ se réalise est $0{,}3$. Quelle est la probabilité de l'événement contraire de $A$~?\end{minipage}
\end{minipage}

\vfill
\begin{tcolorbox}[enhanced,nobeforeafter,width=7.1cm,height=2.1cm,valign=center,colback=cdon!6,colframe=cdon,arc=2.5mm,boxrule=0.9pt,left=3mm]{\Large\bfseries\color{cdon}a)}\quad{\Large $0{,}3$}\end{tcolorbox}\hfill\begin{tcolorbox}[enhanced,nobeforeafter,width=7.1cm,height=2.1cm,valign=center,colback=cdon!6,colframe=cdon,arc=2.5mm,boxrule=0.9pt,left=3mm]{\Large\bfseries\color{cdon}b)}\quad{\Large $0{,}7$}\end{tcolorbox}\par\vspace{2.5mm}
\begin{tcolorbox}[enhanced,nobeforeafter,width=7.1cm,height=2.1cm,valign=center,colback=cdon!6,colframe=cdon,arc=2.5mm,boxrule=0.9pt,left=3mm]{\Large\bfseries\color{cdon}c)}\quad{\Large $0{,}2$}\end{tcolorbox}\hfill\begin{tcolorbox}[enhanced,nobeforeafter,width=7.1cm,height=2.1cm,valign=center,colback=cdon!6,colframe=cdon,arc=2.5mm,boxrule=0.9pt,left=3mm]{\Large\bfseries\color{cdon}d)}\quad{\Large $0{,}85$}\end{tcolorbox}
\end{tcolorbox}
\newpage
\begin{tcolorbox}[enhanced,width=16cm,height=14cm,colback=white,colframe=cdon,boxrule=2pt,arc=0mm,
left=6mm,right=6mm,top=0mm,bottom=4mm,sharp corners,
overlay={\fill[cdon] (frame.north west) rectangle ([yshift=-1.25cm]frame.north east);
\node[anchor=west,text=white,font=\sffamily\bfseries\LARGE] at ([xshift=5mm,yshift=-0.62cm]frame.north west) {Question 80};
\node[anchor=east,text=white,font=\sffamily\itshape\large] at ([xshift=-5mm,yshift=-0.62cm]frame.north east) {Proportionnalité, statistiques, probabilités};}]
\vspace*{1.75cm}
\begin{minipage}[t][6.9cm][c]{\linewidth}\LARGE\centering
\begin{minipage}{0.95\linewidth}\raggedright\linespread{1.25}\selectfont La probabilité qu'un événement $A$ se réalise est $0{,}65$. Quelle est la probabilité de l'événement contraire de $A$~?\end{minipage}
\end{minipage}

\vfill
\begin{tcolorbox}[enhanced,nobeforeafter,width=7.1cm,height=2.1cm,valign=center,colback=cdon!6,colframe=cdon,arc=2.5mm,boxrule=0.9pt,left=3mm]{\Large\bfseries\color{cdon}a)}\quad{\Large $0{,}15$}\end{tcolorbox}\hfill\begin{tcolorbox}[enhanced,nobeforeafter,width=7.1cm,height=2.1cm,valign=center,colback=cdon!6,colframe=cdon,arc=2.5mm,boxrule=0.9pt,left=3mm]{\Large\bfseries\color{cdon}b)}\quad{\Large $0{,}675$}\end{tcolorbox}\par\vspace{2.5mm}
\begin{tcolorbox}[enhanced,nobeforeafter,width=7.1cm,height=2.1cm,valign=center,colback=cdon!6,colframe=cdon,arc=2.5mm,boxrule=0.9pt,left=3mm]{\Large\bfseries\color{cdon}c)}\quad{\Large $0{,}35$}\end{tcolorbox}\hfill\begin{tcolorbox}[enhanced,nobeforeafter,width=7.1cm,height=2.1cm,valign=center,colback=cdon!6,colframe=cdon,arc=2.5mm,boxrule=0.9pt,left=3mm]{\Large\bfseries\color{cdon}d)}\quad{\Large $0{,}65$}\end{tcolorbox}
\end{tcolorbox}
\newpage
\begin{tcolorbox}[enhanced,width=16cm,height=14cm,colback=white,colframe=cgra,boxrule=2pt,arc=0mm,
left=6mm,right=6mm,top=0mm,bottom=4mm,sharp corners,
overlay={\fill[cgra] (frame.north west) rectangle ([yshift=-1.25cm]frame.north east);
\node[anchor=west,text=white,font=\sffamily\bfseries\LARGE] at ([xshift=5mm,yshift=-0.62cm]frame.north west) {Question 81};
\node[anchor=east,text=white,font=\sffamily\itshape\large] at ([xshift=-5mm,yshift=-0.62cm]frame.north east) {Grandeurs et mesures};}]
\vspace*{1.75cm}
\begin{minipage}[t][6.9cm][c]{\linewidth}\LARGE\centering
\begin{minipage}{0.95\linewidth}\raggedright\linespread{1.25}\selectfont Un véhicule parcourt $150$ km en 2 h. Quelle est sa vitesse moyenne~?\end{minipage}
\end{minipage}

\vfill
\begin{tcolorbox}[enhanced,nobeforeafter,width=7.1cm,height=2.1cm,valign=center,colback=cgra!6,colframe=cgra,arc=2.5mm,boxrule=0.9pt,left=3mm]{\Large\bfseries\color{cgra}a)}\quad{\Large $85~\text{km/h}$}\end{tcolorbox}\hfill\begin{tcolorbox}[enhanced,nobeforeafter,width=7.1cm,height=2.1cm,valign=center,colback=cgra!6,colframe=cgra,arc=2.5mm,boxrule=0.9pt,left=3mm]{\Large\bfseries\color{cgra}b)}\quad{\Large $300~\text{km/h}$}\end{tcolorbox}\par\vspace{2.5mm}
\begin{tcolorbox}[enhanced,nobeforeafter,width=7.1cm,height=2.1cm,valign=center,colback=cgra!6,colframe=cgra,arc=2.5mm,boxrule=0.9pt,left=3mm]{\Large\bfseries\color{cgra}c)}\quad{\Large $1{,}25~\text{km/h}$}\end{tcolorbox}\hfill\begin{tcolorbox}[enhanced,nobeforeafter,width=7.1cm,height=2.1cm,valign=center,colback=cgra!6,colframe=cgra,arc=2.5mm,boxrule=0.9pt,left=3mm]{\Large\bfseries\color{cgra}d)}\quad{\Large $75~\text{km/h}$}\end{tcolorbox}
\end{tcolorbox}
\newpage
\begin{tcolorbox}[enhanced,width=16cm,height=14cm,colback=white,colframe=cgra,boxrule=2pt,arc=0mm,
left=6mm,right=6mm,top=0mm,bottom=4mm,sharp corners,
overlay={\fill[cgra] (frame.north west) rectangle ([yshift=-1.25cm]frame.north east);
\node[anchor=west,text=white,font=\sffamily\bfseries\LARGE] at ([xshift=5mm,yshift=-0.62cm]frame.north west) {Question 82};
\node[anchor=east,text=white,font=\sffamily\itshape\large] at ([xshift=-5mm,yshift=-0.62cm]frame.north east) {Grandeurs et mesures};}]
\vspace*{1.75cm}
\begin{minipage}[t][6.9cm][c]{\linewidth}\LARGE\centering
\begin{minipage}{0.95\linewidth}\raggedright\linespread{1.25}\selectfont Un véhicule parcourt $240$ km en 3 h. Quelle est sa vitesse moyenne~?\end{minipage}
\end{minipage}

\vfill
\begin{tcolorbox}[enhanced,nobeforeafter,width=7.1cm,height=2.1cm,valign=center,colback=cgra!6,colframe=cgra,arc=2.5mm,boxrule=0.9pt,left=3mm]{\Large\bfseries\color{cgra}a)}\quad{\Large $80~\text{km/h}$}\end{tcolorbox}\hfill\begin{tcolorbox}[enhanced,nobeforeafter,width=7.1cm,height=2.1cm,valign=center,colback=cgra!6,colframe=cgra,arc=2.5mm,boxrule=0.9pt,left=3mm]{\Large\bfseries\color{cgra}b)}\quad{\Large $90~\text{km/h}$}\end{tcolorbox}\par\vspace{2.5mm}
\begin{tcolorbox}[enhanced,nobeforeafter,width=7.1cm,height=2.1cm,valign=center,colback=cgra!6,colframe=cgra,arc=2.5mm,boxrule=0.9pt,left=3mm]{\Large\bfseries\color{cgra}c)}\quad{\Large $720~\text{km/h}$}\end{tcolorbox}\hfill\begin{tcolorbox}[enhanced,nobeforeafter,width=7.1cm,height=2.1cm,valign=center,colback=cgra!6,colframe=cgra,arc=2.5mm,boxrule=0.9pt,left=3mm]{\Large\bfseries\color{cgra}d)}\quad{\Large $160~\text{km/h}$}\end{tcolorbox}
\end{tcolorbox}
\newpage
\begin{tcolorbox}[enhanced,width=16cm,height=14cm,colback=white,colframe=cgra,boxrule=2pt,arc=0mm,
left=6mm,right=6mm,top=0mm,bottom=4mm,sharp corners,
overlay={\fill[cgra] (frame.north west) rectangle ([yshift=-1.25cm]frame.north east);
\node[anchor=west,text=white,font=\sffamily\bfseries\LARGE] at ([xshift=5mm,yshift=-0.62cm]frame.north west) {Question 83};
\node[anchor=east,text=white,font=\sffamily\itshape\large] at ([xshift=-5mm,yshift=-0.62cm]frame.north east) {Grandeurs et mesures};}]
\vspace*{1.75cm}
\begin{minipage}[t][6.9cm][c]{\linewidth}\LARGE\centering
\begin{minipage}{0.95\linewidth}\raggedright\linespread{1.25}\selectfont Un cycliste roule à la vitesse constante de $20$ km/h pendant $30$ min. Quelle distance parcourt-il~?\end{minipage}
\end{minipage}

\vfill
\begin{tcolorbox}[enhanced,nobeforeafter,width=7.1cm,height=2.1cm,valign=center,colback=cgra!6,colframe=cgra,arc=2.5mm,boxrule=0.9pt,left=3mm]{\Large\bfseries\color{cgra}a)}\quad{\Large $10~\text{km}$}\end{tcolorbox}\hfill\begin{tcolorbox}[enhanced,nobeforeafter,width=7.1cm,height=2.1cm,valign=center,colback=cgra!6,colframe=cgra,arc=2.5mm,boxrule=0.9pt,left=3mm]{\Large\bfseries\color{cgra}b)}\quad{\Large $600~\text{km}$}\end{tcolorbox}\par\vspace{2.5mm}
\begin{tcolorbox}[enhanced,nobeforeafter,width=7.1cm,height=2.1cm,valign=center,colback=cgra!6,colframe=cgra,arc=2.5mm,boxrule=0.9pt,left=3mm]{\Large\bfseries\color{cgra}c)}\quad{\Large $6~\text{km}$}\end{tcolorbox}\hfill\begin{tcolorbox}[enhanced,nobeforeafter,width=7.1cm,height=2.1cm,valign=center,colback=cgra!6,colframe=cgra,arc=2.5mm,boxrule=0.9pt,left=3mm]{\Large\bfseries\color{cgra}d)}\quad{\Large $20~\text{km}$}\end{tcolorbox}
\end{tcolorbox}
\newpage
\begin{tcolorbox}[enhanced,width=16cm,height=14cm,colback=white,colframe=cgra,boxrule=2pt,arc=0mm,
left=6mm,right=6mm,top=0mm,bottom=4mm,sharp corners,
overlay={\fill[cgra] (frame.north west) rectangle ([yshift=-1.25cm]frame.north east);
\node[anchor=west,text=white,font=\sffamily\bfseries\LARGE] at ([xshift=5mm,yshift=-0.62cm]frame.north west) {Question 84};
\node[anchor=east,text=white,font=\sffamily\itshape\large] at ([xshift=-5mm,yshift=-0.62cm]frame.north east) {Grandeurs et mesures};}]
\vspace*{1.75cm}
\begin{minipage}[t][6.9cm][c]{\linewidth}\LARGE\centering
\begin{minipage}{0.95\linewidth}\raggedright\linespread{1.25}\selectfont Un cycliste roule à la vitesse constante de $15$ km/h pendant $48$ min. Quelle distance parcourt-il~?\end{minipage}
\end{minipage}

\vfill
\begin{tcolorbox}[enhanced,nobeforeafter,width=7.1cm,height=2.1cm,valign=center,colback=cgra!6,colframe=cgra,arc=2.5mm,boxrule=0.9pt,left=3mm]{\Large\bfseries\color{cgra}a)}\quad{\Large $7{,}2~\text{km}$}\end{tcolorbox}\hfill\begin{tcolorbox}[enhanced,nobeforeafter,width=7.1cm,height=2.1cm,valign=center,colback=cgra!6,colframe=cgra,arc=2.5mm,boxrule=0.9pt,left=3mm]{\Large\bfseries\color{cgra}b)}\quad{\Large $720~\text{km}$}\end{tcolorbox}\par\vspace{2.5mm}
\begin{tcolorbox}[enhanced,nobeforeafter,width=7.1cm,height=2.1cm,valign=center,colback=cgra!6,colframe=cgra,arc=2.5mm,boxrule=0.9pt,left=3mm]{\Large\bfseries\color{cgra}c)}\quad{\Large $12~\text{km}$}\end{tcolorbox}\hfill\begin{tcolorbox}[enhanced,nobeforeafter,width=7.1cm,height=2.1cm,valign=center,colback=cgra!6,colframe=cgra,arc=2.5mm,boxrule=0.9pt,left=3mm]{\Large\bfseries\color{cgra}d)}\quad{\Large $24~\text{km}$}\end{tcolorbox}
\end{tcolorbox}
\newpage
\begin{tcolorbox}[enhanced,width=16cm,height=14cm,colback=white,colframe=cgra,boxrule=2pt,arc=0mm,
left=6mm,right=6mm,top=0mm,bottom=4mm,sharp corners,
overlay={\fill[cgra] (frame.north west) rectangle ([yshift=-1.25cm]frame.north east);
\node[anchor=west,text=white,font=\sffamily\bfseries\LARGE] at ([xshift=5mm,yshift=-0.62cm]frame.north west) {Question 85};
\node[anchor=east,text=white,font=\sffamily\itshape\large] at ([xshift=-5mm,yshift=-0.62cm]frame.north east) {Grandeurs et mesures};}]
\vspace*{1.75cm}
\begin{minipage}[t][6.9cm][c]{\linewidth}\LARGE\centering
\begin{minipage}{0.95\linewidth}\raggedright\linespread{1.25}\selectfont Un train parcourt $240$ km à la vitesse moyenne de $80$ km/h. Combien de temps dure le trajet~?\end{minipage}
\end{minipage}

\vfill
\begin{tcolorbox}[enhanced,nobeforeafter,width=7.1cm,height=2.1cm,valign=center,colback=cgra!6,colframe=cgra,arc=2.5mm,boxrule=0.9pt,left=3mm]{\Large\bfseries\color{cgra}a)}\quad{\Large 3 h 10 min}\end{tcolorbox}\hfill\begin{tcolorbox}[enhanced,nobeforeafter,width=7.1cm,height=2.1cm,valign=center,colback=cgra!6,colframe=cgra,arc=2.5mm,boxrule=0.9pt,left=3mm]{\Large\bfseries\color{cgra}b)}\quad{\Large 2 h 45 min}\end{tcolorbox}\par\vspace{2.5mm}
\begin{tcolorbox}[enhanced,nobeforeafter,width=7.1cm,height=2.1cm,valign=center,colback=cgra!6,colframe=cgra,arc=2.5mm,boxrule=0.9pt,left=3mm]{\Large\bfseries\color{cgra}c)}\quad{\Large 3 h 30 min}\end{tcolorbox}\hfill\begin{tcolorbox}[enhanced,nobeforeafter,width=7.1cm,height=2.1cm,valign=center,colback=cgra!6,colframe=cgra,arc=2.5mm,boxrule=0.9pt,left=3mm]{\Large\bfseries\color{cgra}d)}\quad{\Large 3 h}\end{tcolorbox}
\end{tcolorbox}
\newpage
\begin{tcolorbox}[enhanced,width=16cm,height=14cm,colback=white,colframe=cgra,boxrule=2pt,arc=0mm,
left=6mm,right=6mm,top=0mm,bottom=4mm,sharp corners,
overlay={\fill[cgra] (frame.north west) rectangle ([yshift=-1.25cm]frame.north east);
\node[anchor=west,text=white,font=\sffamily\bfseries\LARGE] at ([xshift=5mm,yshift=-0.62cm]frame.north west) {Question 86};
\node[anchor=east,text=white,font=\sffamily\itshape\large] at ([xshift=-5mm,yshift=-0.62cm]frame.north east) {Grandeurs et mesures};}]
\vspace*{1.75cm}
\begin{minipage}[t][6.9cm][c]{\linewidth}\LARGE\centering
\begin{minipage}{0.95\linewidth}\raggedright\linespread{1.25}\selectfont Un train parcourt $150$ km à la vitesse moyenne de $60$ km/h. Combien de temps dure le trajet~?\end{minipage}
\end{minipage}

\vfill
\begin{tcolorbox}[enhanced,nobeforeafter,width=7.1cm,height=2.1cm,valign=center,colback=cgra!6,colframe=cgra,arc=2.5mm,boxrule=0.9pt,left=3mm]{\Large\bfseries\color{cgra}a)}\quad{\Large 2 h 15 min}\end{tcolorbox}\hfill\begin{tcolorbox}[enhanced,nobeforeafter,width=7.1cm,height=2.1cm,valign=center,colback=cgra!6,colframe=cgra,arc=2.5mm,boxrule=0.9pt,left=3mm]{\Large\bfseries\color{cgra}b)}\quad{\Large 2 h 30 min}\end{tcolorbox}\par\vspace{2.5mm}
\begin{tcolorbox}[enhanced,nobeforeafter,width=7.1cm,height=2.1cm,valign=center,colback=cgra!6,colframe=cgra,arc=2.5mm,boxrule=0.9pt,left=3mm]{\Large\bfseries\color{cgra}c)}\quad{\Large 2 h 50 min}\end{tcolorbox}\hfill\begin{tcolorbox}[enhanced,nobeforeafter,width=7.1cm,height=2.1cm,valign=center,colback=cgra!6,colframe=cgra,arc=2.5mm,boxrule=0.9pt,left=3mm]{\Large\bfseries\color{cgra}d)}\quad{\Large 3 h}\end{tcolorbox}
\end{tcolorbox}
\newpage
\begin{tcolorbox}[enhanced,width=16cm,height=14cm,colback=white,colframe=cgra,boxrule=2pt,arc=0mm,
left=6mm,right=6mm,top=0mm,bottom=4mm,sharp corners,
overlay={\fill[cgra] (frame.north west) rectangle ([yshift=-1.25cm]frame.north east);
\node[anchor=west,text=white,font=\sffamily\bfseries\LARGE] at ([xshift=5mm,yshift=-0.62cm]frame.north west) {Question 87};
\node[anchor=east,text=white,font=\sffamily\itshape\large] at ([xshift=-5mm,yshift=-0.62cm]frame.north east) {Grandeurs et mesures};}]
\vspace*{1.75cm}
\begin{minipage}[t][6.9cm][c]{\linewidth}\LARGE\centering
\begin{minipage}{0.95\linewidth}\raggedright\linespread{1.25}\selectfont Un véhicule roule à $36$ km/h. Quelle est sa vitesse en mètres par seconde~?\end{minipage}
\end{minipage}

\vfill
\begin{tcolorbox}[enhanced,nobeforeafter,width=7.1cm,height=2.1cm,valign=center,colback=cgra!6,colframe=cgra,arc=2.5mm,boxrule=0.9pt,left=3mm]{\Large\bfseries\color{cgra}a)}\quad{\Large $10~\text{m/s}$}\end{tcolorbox}\hfill\begin{tcolorbox}[enhanced,nobeforeafter,width=7.1cm,height=2.1cm,valign=center,colback=cgra!6,colframe=cgra,arc=2.5mm,boxrule=0.9pt,left=3mm]{\Large\bfseries\color{cgra}b)}\quad{\Large $36\,000~\text{m/s}$}\end{tcolorbox}\par\vspace{2.5mm}
\begin{tcolorbox}[enhanced,nobeforeafter,width=7.1cm,height=2.1cm,valign=center,colback=cgra!6,colframe=cgra,arc=2.5mm,boxrule=0.9pt,left=3mm]{\Large\bfseries\color{cgra}c)}\quad{\Large $0{,}6~\text{m/s}$}\end{tcolorbox}\hfill\begin{tcolorbox}[enhanced,nobeforeafter,width=7.1cm,height=2.1cm,valign=center,colback=cgra!6,colframe=cgra,arc=2.5mm,boxrule=0.9pt,left=3mm]{\Large\bfseries\color{cgra}d)}\quad{\Large $129{,}6~\text{m/s}$}\end{tcolorbox}
\end{tcolorbox}
\newpage
\begin{tcolorbox}[enhanced,width=16cm,height=14cm,colback=white,colframe=cgra,boxrule=2pt,arc=0mm,
left=6mm,right=6mm,top=0mm,bottom=4mm,sharp corners,
overlay={\fill[cgra] (frame.north west) rectangle ([yshift=-1.25cm]frame.north east);
\node[anchor=west,text=white,font=\sffamily\bfseries\LARGE] at ([xshift=5mm,yshift=-0.62cm]frame.north west) {Question 88};
\node[anchor=east,text=white,font=\sffamily\itshape\large] at ([xshift=-5mm,yshift=-0.62cm]frame.north east) {Grandeurs et mesures};}]
\vspace*{1.75cm}
\begin{minipage}[t][6.9cm][c]{\linewidth}\LARGE\centering
\begin{minipage}{0.95\linewidth}\raggedright\linespread{1.25}\selectfont Un véhicule roule à $72$ km/h. Quelle est sa vitesse en mètres par seconde~?\end{minipage}
\end{minipage}

\vfill
\begin{tcolorbox}[enhanced,nobeforeafter,width=7.1cm,height=2.1cm,valign=center,colback=cgra!6,colframe=cgra,arc=2.5mm,boxrule=0.9pt,left=3mm]{\Large\bfseries\color{cgra}a)}\quad{\Large $72\,000~\text{m/s}$}\end{tcolorbox}\hfill\begin{tcolorbox}[enhanced,nobeforeafter,width=7.1cm,height=2.1cm,valign=center,colback=cgra!6,colframe=cgra,arc=2.5mm,boxrule=0.9pt,left=3mm]{\Large\bfseries\color{cgra}b)}\quad{\Large $259{,}2~\text{m/s}$}\end{tcolorbox}\par\vspace{2.5mm}
\begin{tcolorbox}[enhanced,nobeforeafter,width=7.1cm,height=2.1cm,valign=center,colback=cgra!6,colframe=cgra,arc=2.5mm,boxrule=0.9pt,left=3mm]{\Large\bfseries\color{cgra}c)}\quad{\Large $20~\text{m/s}$}\end{tcolorbox}\hfill\begin{tcolorbox}[enhanced,nobeforeafter,width=7.1cm,height=2.1cm,valign=center,colback=cgra!6,colframe=cgra,arc=2.5mm,boxrule=0.9pt,left=3mm]{\Large\bfseries\color{cgra}d)}\quad{\Large $1{,}2~\text{m/s}$}\end{tcolorbox}
\end{tcolorbox}
\newpage
\begin{tcolorbox}[enhanced,width=16cm,height=14cm,colback=white,colframe=cgra,boxrule=2pt,arc=0mm,
left=6mm,right=6mm,top=0mm,bottom=4mm,sharp corners,
overlay={\fill[cgra] (frame.north west) rectangle ([yshift=-1.25cm]frame.north east);
\node[anchor=west,text=white,font=\sffamily\bfseries\LARGE] at ([xshift=5mm,yshift=-0.62cm]frame.north west) {Question 89};
\node[anchor=east,text=white,font=\sffamily\itshape\large] at ([xshift=-5mm,yshift=-0.62cm]frame.north east) {Grandeurs et mesures};}]
\vspace*{1.75cm}
\begin{minipage}[t][6.9cm][c]{\linewidth}\LARGE\centering
\begin{minipage}{0.95\linewidth}\raggedright\linespread{1.25}\selectfont Une pyramide a une base d'aire $12$ cm$^2$ et une hauteur de $5$ cm. Quel est son volume~?\end{minipage}
\end{minipage}

\vfill
\begin{tcolorbox}[enhanced,nobeforeafter,width=7.1cm,height=2.1cm,valign=center,colback=cgra!6,colframe=cgra,arc=2.5mm,boxrule=0.9pt,left=3mm]{\Large\bfseries\color{cgra}a)}\quad{\Large $30~\text{cm}^3$}\end{tcolorbox}\hfill\begin{tcolorbox}[enhanced,nobeforeafter,width=7.1cm,height=2.1cm,valign=center,colback=cgra!6,colframe=cgra,arc=2.5mm,boxrule=0.9pt,left=3mm]{\Large\bfseries\color{cgra}b)}\quad{\Large $20~\text{cm}^3$}\end{tcolorbox}\par\vspace{2.5mm}
\begin{tcolorbox}[enhanced,nobeforeafter,width=7.1cm,height=2.1cm,valign=center,colback=cgra!6,colframe=cgra,arc=2.5mm,boxrule=0.9pt,left=3mm]{\Large\bfseries\color{cgra}c)}\quad{\Large $60~\text{cm}^3$}\end{tcolorbox}\hfill\begin{tcolorbox}[enhanced,nobeforeafter,width=7.1cm,height=2.1cm,valign=center,colback=cgra!6,colframe=cgra,arc=2.5mm,boxrule=0.9pt,left=3mm]{\Large\bfseries\color{cgra}d)}\quad{\Large $17~\text{cm}^3$}\end{tcolorbox}
\end{tcolorbox}
\newpage
\begin{tcolorbox}[enhanced,width=16cm,height=14cm,colback=white,colframe=cgra,boxrule=2pt,arc=0mm,
left=6mm,right=6mm,top=0mm,bottom=4mm,sharp corners,
overlay={\fill[cgra] (frame.north west) rectangle ([yshift=-1.25cm]frame.north east);
\node[anchor=west,text=white,font=\sffamily\bfseries\LARGE] at ([xshift=5mm,yshift=-0.62cm]frame.north west) {Question 90};
\node[anchor=east,text=white,font=\sffamily\itshape\large] at ([xshift=-5mm,yshift=-0.62cm]frame.north east) {Grandeurs et mesures};}]
\vspace*{1.75cm}
\begin{minipage}[t][6.9cm][c]{\linewidth}\LARGE\centering
\begin{minipage}{0.95\linewidth}\raggedright\linespread{1.25}\selectfont Une pyramide a une base d'aire $30$ cm$^2$ et une hauteur de $4$ cm. Quel est son volume~?\end{minipage}
\end{minipage}

\vfill
\begin{tcolorbox}[enhanced,nobeforeafter,width=7.1cm,height=2.1cm,valign=center,colback=cgra!6,colframe=cgra,arc=2.5mm,boxrule=0.9pt,left=3mm]{\Large\bfseries\color{cgra}a)}\quad{\Large $120~\text{cm}^3$}\end{tcolorbox}\hfill\begin{tcolorbox}[enhanced,nobeforeafter,width=7.1cm,height=2.1cm,valign=center,colback=cgra!6,colframe=cgra,arc=2.5mm,boxrule=0.9pt,left=3mm]{\Large\bfseries\color{cgra}b)}\quad{\Large $40~\text{cm}^3$}\end{tcolorbox}\par\vspace{2.5mm}
\begin{tcolorbox}[enhanced,nobeforeafter,width=7.1cm,height=2.1cm,valign=center,colback=cgra!6,colframe=cgra,arc=2.5mm,boxrule=0.9pt,left=3mm]{\Large\bfseries\color{cgra}c)}\quad{\Large $34~\text{cm}^3$}\end{tcolorbox}\hfill\begin{tcolorbox}[enhanced,nobeforeafter,width=7.1cm,height=2.1cm,valign=center,colback=cgra!6,colframe=cgra,arc=2.5mm,boxrule=0.9pt,left=3mm]{\Large\bfseries\color{cgra}d)}\quad{\Large $60~\text{cm}^3$}\end{tcolorbox}
\end{tcolorbox}
\newpage
\begin{tcolorbox}[enhanced,width=16cm,height=14cm,colback=white,colframe=cgra,boxrule=2pt,arc=0mm,
left=6mm,right=6mm,top=0mm,bottom=4mm,sharp corners,
overlay={\fill[cgra] (frame.north west) rectangle ([yshift=-1.25cm]frame.north east);
\node[anchor=west,text=white,font=\sffamily\bfseries\LARGE] at ([xshift=5mm,yshift=-0.62cm]frame.north west) {Question 91};
\node[anchor=east,text=white,font=\sffamily\itshape\large] at ([xshift=-5mm,yshift=-0.62cm]frame.north east) {Grandeurs et mesures};}]
\vspace*{1.75cm}
\begin{minipage}[t][6.9cm][c]{\linewidth}\LARGE\centering
\begin{minipage}{0.95\linewidth}\raggedright\linespread{1.25}\selectfont Un cône de révolution a un rayon de base de $3$ cm et une hauteur de $4$ cm. Quel est son volume exact~?\end{minipage}
\end{minipage}

\vfill
\begin{tcolorbox}[enhanced,nobeforeafter,width=7.1cm,height=2.1cm,valign=center,colback=cgra!6,colframe=cgra,arc=2.5mm,boxrule=0.9pt,left=3mm]{\Large\bfseries\color{cgra}a)}\quad{\Large $12\pi~\text{cm}^3$}\end{tcolorbox}\hfill\begin{tcolorbox}[enhanced,nobeforeafter,width=7.1cm,height=2.1cm,valign=center,colback=cgra!6,colframe=cgra,arc=2.5mm,boxrule=0.9pt,left=3mm]{\Large\bfseries\color{cgra}b)}\quad{\Large $48\pi~\text{cm}^3$}\end{tcolorbox}\par\vspace{2.5mm}
\begin{tcolorbox}[enhanced,nobeforeafter,width=7.1cm,height=2.1cm,valign=center,colback=cgra!6,colframe=cgra,arc=2.5mm,boxrule=0.9pt,left=3mm]{\Large\bfseries\color{cgra}c)}\quad{\Large $4\pi~\text{cm}^3$}\end{tcolorbox}\hfill\begin{tcolorbox}[enhanced,nobeforeafter,width=7.1cm,height=2.1cm,valign=center,colback=cgra!6,colframe=cgra,arc=2.5mm,boxrule=0.9pt,left=3mm]{\Large\bfseries\color{cgra}d)}\quad{\Large $36\pi~\text{cm}^3$}\end{tcolorbox}
\end{tcolorbox}
\newpage
\begin{tcolorbox}[enhanced,width=16cm,height=14cm,colback=white,colframe=cgra,boxrule=2pt,arc=0mm,
left=6mm,right=6mm,top=0mm,bottom=4mm,sharp corners,
overlay={\fill[cgra] (frame.north west) rectangle ([yshift=-1.25cm]frame.north east);
\node[anchor=west,text=white,font=\sffamily\bfseries\LARGE] at ([xshift=5mm,yshift=-0.62cm]frame.north west) {Question 92};
\node[anchor=east,text=white,font=\sffamily\itshape\large] at ([xshift=-5mm,yshift=-0.62cm]frame.north east) {Grandeurs et mesures};}]
\vspace*{1.75cm}
\begin{minipage}[t][6.9cm][c]{\linewidth}\LARGE\centering
\begin{minipage}{0.95\linewidth}\raggedright\linespread{1.25}\selectfont Un cône de révolution a un rayon de base de $2$ cm et une hauteur de $6$ cm. Quel est son volume exact~?\end{minipage}
\end{minipage}

\vfill
\begin{tcolorbox}[enhanced,nobeforeafter,width=7.1cm,height=2.1cm,valign=center,colback=cgra!6,colframe=cgra,arc=2.5mm,boxrule=0.9pt,left=3mm]{\Large\bfseries\color{cgra}a)}\quad{\Large $4\pi~\text{cm}^3$}\end{tcolorbox}\hfill\begin{tcolorbox}[enhanced,nobeforeafter,width=7.1cm,height=2.1cm,valign=center,colback=cgra!6,colframe=cgra,arc=2.5mm,boxrule=0.9pt,left=3mm]{\Large\bfseries\color{cgra}b)}\quad{\Large $8\pi~\text{cm}^3$}\end{tcolorbox}\par\vspace{2.5mm}
\begin{tcolorbox}[enhanced,nobeforeafter,width=7.1cm,height=2.1cm,valign=center,colback=cgra!6,colframe=cgra,arc=2.5mm,boxrule=0.9pt,left=3mm]{\Large\bfseries\color{cgra}c)}\quad{\Large $24\pi~\text{cm}^3$}\end{tcolorbox}\hfill\begin{tcolorbox}[enhanced,nobeforeafter,width=7.1cm,height=2.1cm,valign=center,colback=cgra!6,colframe=cgra,arc=2.5mm,boxrule=0.9pt,left=3mm]{\Large\bfseries\color{cgra}d)}\quad{\Large $32\pi~\text{cm}^3$}\end{tcolorbox}
\end{tcolorbox}
\newpage
\begin{tcolorbox}[enhanced,width=16cm,height=14cm,colback=white,colframe=cgra,boxrule=2pt,arc=0mm,
left=6mm,right=6mm,top=0mm,bottom=4mm,sharp corners,
overlay={\fill[cgra] (frame.north west) rectangle ([yshift=-1.25cm]frame.north east);
\node[anchor=west,text=white,font=\sffamily\bfseries\LARGE] at ([xshift=5mm,yshift=-0.62cm]frame.north west) {Question 93};
\node[anchor=east,text=white,font=\sffamily\itshape\large] at ([xshift=-5mm,yshift=-0.62cm]frame.north east) {Grandeurs et mesures};}]
\vspace*{1.75cm}
\begin{minipage}[t][6.9cm][c]{\linewidth}\LARGE\centering
\begin{minipage}{0.95\linewidth}\raggedright\linespread{1.25}\selectfont Compléter~: $1{,}5~\text{L}=\;?\;\text{cm}^3$\end{minipage}
\end{minipage}

\vfill
\begin{tcolorbox}[enhanced,nobeforeafter,width=7.1cm,height=2.1cm,valign=center,colback=cgra!6,colframe=cgra,arc=2.5mm,boxrule=0.9pt,left=3mm]{\Large\bfseries\color{cgra}a)}\quad{\Large $15\,000~\text{cm}^3$}\end{tcolorbox}\hfill\begin{tcolorbox}[enhanced,nobeforeafter,width=7.1cm,height=2.1cm,valign=center,colback=cgra!6,colframe=cgra,arc=2.5mm,boxrule=0.9pt,left=3mm]{\Large\bfseries\color{cgra}b)}\quad{\Large $1\,500~\text{cm}^3$}\end{tcolorbox}\par\vspace{2.5mm}
\begin{tcolorbox}[enhanced,nobeforeafter,width=7.1cm,height=2.1cm,valign=center,colback=cgra!6,colframe=cgra,arc=2.5mm,boxrule=0.9pt,left=3mm]{\Large\bfseries\color{cgra}c)}\quad{\Large $150\,000~\text{cm}^3$}\end{tcolorbox}\hfill\begin{tcolorbox}[enhanced,nobeforeafter,width=7.1cm,height=2.1cm,valign=center,colback=cgra!6,colframe=cgra,arc=2.5mm,boxrule=0.9pt,left=3mm]{\Large\bfseries\color{cgra}d)}\quad{\Large $150~\text{cm}^3$}\end{tcolorbox}
\end{tcolorbox}
\newpage
\begin{tcolorbox}[enhanced,width=16cm,height=14cm,colback=white,colframe=cgra,boxrule=2pt,arc=0mm,
left=6mm,right=6mm,top=0mm,bottom=4mm,sharp corners,
overlay={\fill[cgra] (frame.north west) rectangle ([yshift=-1.25cm]frame.north east);
\node[anchor=west,text=white,font=\sffamily\bfseries\LARGE] at ([xshift=5mm,yshift=-0.62cm]frame.north west) {Question 94};
\node[anchor=east,text=white,font=\sffamily\itshape\large] at ([xshift=-5mm,yshift=-0.62cm]frame.north east) {Grandeurs et mesures};}]
\vspace*{1.75cm}
\begin{minipage}[t][6.9cm][c]{\linewidth}\LARGE\centering
\begin{minipage}{0.95\linewidth}\raggedright\linespread{1.25}\selectfont Compléter~: $750~\text{cm}^3=\;?\;\text{L}$\end{minipage}
\end{minipage}

\vfill
\begin{tcolorbox}[enhanced,nobeforeafter,width=7.1cm,height=2.1cm,valign=center,colback=cgra!6,colframe=cgra,arc=2.5mm,boxrule=0.9pt,left=3mm]{\Large\bfseries\color{cgra}a)}\quad{\Large $0{,}75~\text{L}$}\end{tcolorbox}\hfill\begin{tcolorbox}[enhanced,nobeforeafter,width=7.1cm,height=2.1cm,valign=center,colback=cgra!6,colframe=cgra,arc=2.5mm,boxrule=0.9pt,left=3mm]{\Large\bfseries\color{cgra}b)}\quad{\Large $75~\text{L}$}\end{tcolorbox}\par\vspace{2.5mm}
\begin{tcolorbox}[enhanced,nobeforeafter,width=7.1cm,height=2.1cm,valign=center,colback=cgra!6,colframe=cgra,arc=2.5mm,boxrule=0.9pt,left=3mm]{\Large\bfseries\color{cgra}c)}\quad{\Large $7{,}5~\text{L}$}\end{tcolorbox}\hfill\begin{tcolorbox}[enhanced,nobeforeafter,width=7.1cm,height=2.1cm,valign=center,colback=cgra!6,colframe=cgra,arc=2.5mm,boxrule=0.9pt,left=3mm]{\Large\bfseries\color{cgra}d)}\quad{\Large $0{,}075~\text{L}$}\end{tcolorbox}
\end{tcolorbox}
\newpage
\begin{tcolorbox}[enhanced,width=16cm,height=14cm,colback=white,colframe=cgra,boxrule=2pt,arc=0mm,
left=6mm,right=6mm,top=0mm,bottom=4mm,sharp corners,
overlay={\fill[cgra] (frame.north west) rectangle ([yshift=-1.25cm]frame.north east);
\node[anchor=west,text=white,font=\sffamily\bfseries\LARGE] at ([xshift=5mm,yshift=-0.62cm]frame.north west) {Question 95};
\node[anchor=east,text=white,font=\sffamily\itshape\large] at ([xshift=-5mm,yshift=-0.62cm]frame.north east) {Grandeurs et mesures};}]
\vspace*{1.75cm}
\begin{minipage}[t][6.9cm][c]{\linewidth}\LARGE\centering
\begin{minipage}{0.95\linewidth}\raggedright\linespread{1.25}\selectfont Compléter~: $3~\text{m}^2=\;?\;\text{cm}^2$\end{minipage}
\end{minipage}

\vfill
\begin{tcolorbox}[enhanced,nobeforeafter,width=7.1cm,height=2.1cm,valign=center,colback=cgra!6,colframe=cgra,arc=2.5mm,boxrule=0.9pt,left=3mm]{\Large\bfseries\color{cgra}a)}\quad{\Large $3\,000\,000~\text{cm}^2$}\end{tcolorbox}\hfill\begin{tcolorbox}[enhanced,nobeforeafter,width=7.1cm,height=2.1cm,valign=center,colback=cgra!6,colframe=cgra,arc=2.5mm,boxrule=0.9pt,left=3mm]{\Large\bfseries\color{cgra}b)}\quad{\Large $300~\text{cm}^2$}\end{tcolorbox}\par\vspace{2.5mm}
\begin{tcolorbox}[enhanced,nobeforeafter,width=7.1cm,height=2.1cm,valign=center,colback=cgra!6,colframe=cgra,arc=2.5mm,boxrule=0.9pt,left=3mm]{\Large\bfseries\color{cgra}c)}\quad{\Large $300\,000~\text{cm}^2$}\end{tcolorbox}\hfill\begin{tcolorbox}[enhanced,nobeforeafter,width=7.1cm,height=2.1cm,valign=center,colback=cgra!6,colframe=cgra,arc=2.5mm,boxrule=0.9pt,left=3mm]{\Large\bfseries\color{cgra}d)}\quad{\Large $30\,000~\text{cm}^2$}\end{tcolorbox}
\end{tcolorbox}
\newpage
\begin{tcolorbox}[enhanced,width=16cm,height=14cm,colback=white,colframe=cgra,boxrule=2pt,arc=0mm,
left=6mm,right=6mm,top=0mm,bottom=4mm,sharp corners,
overlay={\fill[cgra] (frame.north west) rectangle ([yshift=-1.25cm]frame.north east);
\node[anchor=west,text=white,font=\sffamily\bfseries\LARGE] at ([xshift=5mm,yshift=-0.62cm]frame.north west) {Question 96};
\node[anchor=east,text=white,font=\sffamily\itshape\large] at ([xshift=-5mm,yshift=-0.62cm]frame.north east) {Grandeurs et mesures};}]
\vspace*{1.75cm}
\begin{minipage}[t][6.9cm][c]{\linewidth}\LARGE\centering
\begin{minipage}{0.95\linewidth}\raggedright\linespread{1.25}\selectfont Compléter~: $2{,}5~\text{dm}^2=\;?\;\text{cm}^2$\end{minipage}
\end{minipage}

\vfill
\begin{tcolorbox}[enhanced,nobeforeafter,width=7.1cm,height=2.1cm,valign=center,colback=cgra!6,colframe=cgra,arc=2.5mm,boxrule=0.9pt,left=3mm]{\Large\bfseries\color{cgra}a)}\quad{\Large $2\,500~\text{cm}^2$}\end{tcolorbox}\hfill\begin{tcolorbox}[enhanced,nobeforeafter,width=7.1cm,height=2.1cm,valign=center,colback=cgra!6,colframe=cgra,arc=2.5mm,boxrule=0.9pt,left=3mm]{\Large\bfseries\color{cgra}b)}\quad{\Large $250~\text{cm}^2$}\end{tcolorbox}\par\vspace{2.5mm}
\begin{tcolorbox}[enhanced,nobeforeafter,width=7.1cm,height=2.1cm,valign=center,colback=cgra!6,colframe=cgra,arc=2.5mm,boxrule=0.9pt,left=3mm]{\Large\bfseries\color{cgra}c)}\quad{\Large $25\,000~\text{cm}^2$}\end{tcolorbox}\hfill\begin{tcolorbox}[enhanced,nobeforeafter,width=7.1cm,height=2.1cm,valign=center,colback=cgra!6,colframe=cgra,arc=2.5mm,boxrule=0.9pt,left=3mm]{\Large\bfseries\color{cgra}d)}\quad{\Large $25~\text{cm}^2$}\end{tcolorbox}
\end{tcolorbox}
\newpage
\begin{tcolorbox}[enhanced,width=16cm,height=14cm,colback=white,colframe=crel,boxrule=2pt,arc=0mm,
left=6mm,right=6mm,top=0mm,bottom=4mm,sharp corners,
overlay={\fill[crel] (frame.north west) rectangle ([yshift=-1.25cm]frame.north east);
\node[anchor=west,text=white,font=\sffamily\bfseries\LARGE] at ([xshift=5mm,yshift=-0.62cm]frame.north west) {Question 1 B};
\node[anchor=east,text=white,font=\sffamily\itshape\large] at ([xshift=-5mm,yshift=-0.62cm]frame.north east) {Relatifs et fractions};}]
\vspace*{1.75cm}
\begin{minipage}[t][6.9cm][c]{\linewidth}\LARGE\centering
\begin{minipage}{0.95\linewidth}\raggedright\linespread{1.25}\selectfont Que vaut $(-1)\times(-6)\times(-2)$~?\end{minipage}
\end{minipage}

\vfill
\begin{tcolorbox}[enhanced,nobeforeafter,width=7.1cm,height=2.1cm,valign=center,colback=crel!6,colframe=crel,arc=2.5mm,boxrule=0.9pt,left=3mm]{\Large\bfseries\color{crel}a)}\quad{\Large $12$}\end{tcolorbox}\hfill\begin{tcolorbox}[enhanced,nobeforeafter,width=7.1cm,height=2.1cm,valign=center,colback=crel!6,colframe=crel,arc=2.5mm,boxrule=0.9pt,left=3mm]{\Large\bfseries\color{crel}b)}\quad{\Large $-9$}\end{tcolorbox}\par\vspace{2.5mm}
\begin{tcolorbox}[enhanced,nobeforeafter,width=7.1cm,height=2.1cm,valign=center,colback=crel!6,colframe=crel,arc=2.5mm,boxrule=0.9pt,left=3mm]{\Large\bfseries\color{crel}c)}\quad{\Large $8$}\end{tcolorbox}\hfill\begin{tcolorbox}[enhanced,nobeforeafter,width=7.1cm,height=2.1cm,valign=center,colback=crel!6,colframe=crel,arc=2.5mm,boxrule=0.9pt,left=3mm]{\Large\bfseries\color{crel}d)}\quad{\Large $-12$}\end{tcolorbox}
\end{tcolorbox}
\newpage
\begin{tcolorbox}[enhanced,width=16cm,height=14cm,colback=white,colframe=crel,boxrule=2pt,arc=0mm,
left=6mm,right=6mm,top=0mm,bottom=4mm,sharp corners,
overlay={\fill[crel] (frame.north west) rectangle ([yshift=-1.25cm]frame.north east);
\node[anchor=west,text=white,font=\sffamily\bfseries\LARGE] at ([xshift=5mm,yshift=-0.62cm]frame.north west) {Question 2 B};
\node[anchor=east,text=white,font=\sffamily\itshape\large] at ([xshift=-5mm,yshift=-0.62cm]frame.north east) {Relatifs et fractions};}]
\vspace*{1.75cm}
\begin{minipage}[t][6.9cm][c]{\linewidth}\LARGE\centering
\begin{minipage}{0.95\linewidth}\raggedright\linespread{1.25}\selectfont Que vaut $(+4)\times(-3)\times(-2)$~?\end{minipage}
\end{minipage}

\vfill
\begin{tcolorbox}[enhanced,nobeforeafter,width=7.1cm,height=2.1cm,valign=center,colback=crel!6,colframe=crel,arc=2.5mm,boxrule=0.9pt,left=3mm]{\Large\bfseries\color{crel}a)}\quad{\Large $24$}\end{tcolorbox}\hfill\begin{tcolorbox}[enhanced,nobeforeafter,width=7.1cm,height=2.1cm,valign=center,colback=crel!6,colframe=crel,arc=2.5mm,boxrule=0.9pt,left=3mm]{\Large\bfseries\color{crel}b)}\quad{\Large $-1$}\end{tcolorbox}\par\vspace{2.5mm}
\begin{tcolorbox}[enhanced,nobeforeafter,width=7.1cm,height=2.1cm,valign=center,colback=crel!6,colframe=crel,arc=2.5mm,boxrule=0.9pt,left=3mm]{\Large\bfseries\color{crel}c)}\quad{\Large $-24$}\end{tcolorbox}\hfill\begin{tcolorbox}[enhanced,nobeforeafter,width=7.1cm,height=2.1cm,valign=center,colback=crel!6,colframe=crel,arc=2.5mm,boxrule=0.9pt,left=3mm]{\Large\bfseries\color{crel}d)}\quad{\Large $-10$}\end{tcolorbox}
\end{tcolorbox}
\newpage
\begin{tcolorbox}[enhanced,width=16cm,height=14cm,colback=white,colframe=crel,boxrule=2pt,arc=0mm,
left=6mm,right=6mm,top=0mm,bottom=4mm,sharp corners,
overlay={\fill[crel] (frame.north west) rectangle ([yshift=-1.25cm]frame.north east);
\node[anchor=west,text=white,font=\sffamily\bfseries\LARGE] at ([xshift=5mm,yshift=-0.62cm]frame.north west) {Question 3 B};
\node[anchor=east,text=white,font=\sffamily\itshape\large] at ([xshift=-5mm,yshift=-0.62cm]frame.north east) {Relatifs et fractions};}]
\vspace*{1.75cm}
\begin{minipage}[t][6.9cm][c]{\linewidth}\LARGE\centering
\begin{minipage}{0.95\linewidth}\raggedright\linespread{1.25}\selectfont Que vaut $(-56)\div(-8)$~?\end{minipage}
\end{minipage}

\vfill
\begin{tcolorbox}[enhanced,nobeforeafter,width=7.1cm,height=2.1cm,valign=center,colback=crel!6,colframe=crel,arc=2.5mm,boxrule=0.9pt,left=3mm]{\Large\bfseries\color{crel}a)}\quad{\Large $-7$}\end{tcolorbox}\hfill\begin{tcolorbox}[enhanced,nobeforeafter,width=7.1cm,height=2.1cm,valign=center,colback=crel!6,colframe=crel,arc=2.5mm,boxrule=0.9pt,left=3mm]{\Large\bfseries\color{crel}b)}\quad{\Large $8$}\end{tcolorbox}\par\vspace{2.5mm}
\begin{tcolorbox}[enhanced,nobeforeafter,width=7.1cm,height=2.1cm,valign=center,colback=crel!6,colframe=crel,arc=2.5mm,boxrule=0.9pt,left=3mm]{\Large\bfseries\color{crel}c)}\quad{\Large $7$}\end{tcolorbox}\hfill\begin{tcolorbox}[enhanced,nobeforeafter,width=7.1cm,height=2.1cm,valign=center,colback=crel!6,colframe=crel,arc=2.5mm,boxrule=0.9pt,left=3mm]{\Large\bfseries\color{crel}d)}\quad{\Large $14$}\end{tcolorbox}
\end{tcolorbox}
\newpage
\begin{tcolorbox}[enhanced,width=16cm,height=14cm,colback=white,colframe=crel,boxrule=2pt,arc=0mm,
left=6mm,right=6mm,top=0mm,bottom=4mm,sharp corners,
overlay={\fill[crel] (frame.north west) rectangle ([yshift=-1.25cm]frame.north east);
\node[anchor=west,text=white,font=\sffamily\bfseries\LARGE] at ([xshift=5mm,yshift=-0.62cm]frame.north west) {Question 4 B};
\node[anchor=east,text=white,font=\sffamily\itshape\large] at ([xshift=-5mm,yshift=-0.62cm]frame.north east) {Relatifs et fractions};}]
\vspace*{1.75cm}
\begin{minipage}[t][6.9cm][c]{\linewidth}\LARGE\centering
\begin{minipage}{0.95\linewidth}\raggedright\linespread{1.25}\selectfont Que vaut $(+45)\div(-9)$~?\end{minipage}
\end{minipage}

\vfill
\begin{tcolorbox}[enhanced,nobeforeafter,width=7.1cm,height=2.1cm,valign=center,colback=crel!6,colframe=crel,arc=2.5mm,boxrule=0.9pt,left=3mm]{\Large\bfseries\color{crel}a)}\quad{\Large $-10$}\end{tcolorbox}\hfill\begin{tcolorbox}[enhanced,nobeforeafter,width=7.1cm,height=2.1cm,valign=center,colback=crel!6,colframe=crel,arc=2.5mm,boxrule=0.9pt,left=3mm]{\Large\bfseries\color{crel}b)}\quad{\Large $-5$}\end{tcolorbox}\par\vspace{2.5mm}
\begin{tcolorbox}[enhanced,nobeforeafter,width=7.1cm,height=2.1cm,valign=center,colback=crel!6,colframe=crel,arc=2.5mm,boxrule=0.9pt,left=3mm]{\Large\bfseries\color{crel}c)}\quad{\Large $-4$}\end{tcolorbox}\hfill\begin{tcolorbox}[enhanced,nobeforeafter,width=7.1cm,height=2.1cm,valign=center,colback=crel!6,colframe=crel,arc=2.5mm,boxrule=0.9pt,left=3mm]{\Large\bfseries\color{crel}d)}\quad{\Large $5$}\end{tcolorbox}
\end{tcolorbox}
\newpage
\begin{tcolorbox}[enhanced,width=16cm,height=14cm,colback=white,colframe=crel,boxrule=2pt,arc=0mm,
left=6mm,right=6mm,top=0mm,bottom=4mm,sharp corners,
overlay={\fill[crel] (frame.north west) rectangle ([yshift=-1.25cm]frame.north east);
\node[anchor=west,text=white,font=\sffamily\bfseries\LARGE] at ([xshift=5mm,yshift=-0.62cm]frame.north west) {Question 5 B};
\node[anchor=east,text=white,font=\sffamily\itshape\large] at ([xshift=-5mm,yshift=-0.62cm]frame.north east) {Relatifs et fractions};}]
\vspace*{1.75cm}
\begin{minipage}[t][6.9cm][c]{\linewidth}\LARGE\centering
\begin{minipage}{0.95\linewidth}\raggedright\linespread{1.25}\selectfont Que vaut $-2-5\times(-4)$~?\end{minipage}
\end{minipage}

\vfill
\begin{tcolorbox}[enhanced,nobeforeafter,width=7.1cm,height=2.1cm,valign=center,colback=crel!6,colframe=crel,arc=2.5mm,boxrule=0.9pt,left=3mm]{\Large\bfseries\color{crel}a)}\quad{\Large $18$}\end{tcolorbox}\hfill\begin{tcolorbox}[enhanced,nobeforeafter,width=7.1cm,height=2.1cm,valign=center,colback=crel!6,colframe=crel,arc=2.5mm,boxrule=0.9pt,left=3mm]{\Large\bfseries\color{crel}b)}\quad{\Large $-22$}\end{tcolorbox}\par\vspace{2.5mm}
\begin{tcolorbox}[enhanced,nobeforeafter,width=7.1cm,height=2.1cm,valign=center,colback=crel!6,colframe=crel,arc=2.5mm,boxrule=0.9pt,left=3mm]{\Large\bfseries\color{crel}c)}\quad{\Large $-18$}\end{tcolorbox}\hfill\begin{tcolorbox}[enhanced,nobeforeafter,width=7.1cm,height=2.1cm,valign=center,colback=crel!6,colframe=crel,arc=2.5mm,boxrule=0.9pt,left=3mm]{\Large\bfseries\color{crel}d)}\quad{\Large $28$}\end{tcolorbox}
\end{tcolorbox}
\newpage
\begin{tcolorbox}[enhanced,width=16cm,height=14cm,colback=white,colframe=crel,boxrule=2pt,arc=0mm,
left=6mm,right=6mm,top=0mm,bottom=4mm,sharp corners,
overlay={\fill[crel] (frame.north west) rectangle ([yshift=-1.25cm]frame.north east);
\node[anchor=west,text=white,font=\sffamily\bfseries\LARGE] at ([xshift=5mm,yshift=-0.62cm]frame.north west) {Question 6 B};
\node[anchor=east,text=white,font=\sffamily\itshape\large] at ([xshift=-5mm,yshift=-0.62cm]frame.north east) {Relatifs et fractions};}]
\vspace*{1.75cm}
\begin{minipage}[t][6.9cm][c]{\linewidth}\LARGE\centering
\begin{minipage}{0.95\linewidth}\raggedright\linespread{1.25}\selectfont Que vaut $10-2\times(-6)$~?\end{minipage}
\end{minipage}

\vfill
\begin{tcolorbox}[enhanced,nobeforeafter,width=7.1cm,height=2.1cm,valign=center,colback=crel!6,colframe=crel,arc=2.5mm,boxrule=0.9pt,left=3mm]{\Large\bfseries\color{crel}a)}\quad{\Large $-2$}\end{tcolorbox}\hfill\begin{tcolorbox}[enhanced,nobeforeafter,width=7.1cm,height=2.1cm,valign=center,colback=crel!6,colframe=crel,arc=2.5mm,boxrule=0.9pt,left=3mm]{\Large\bfseries\color{crel}b)}\quad{\Large $-22$}\end{tcolorbox}\par\vspace{2.5mm}
\begin{tcolorbox}[enhanced,nobeforeafter,width=7.1cm,height=2.1cm,valign=center,colback=crel!6,colframe=crel,arc=2.5mm,boxrule=0.9pt,left=3mm]{\Large\bfseries\color{crel}c)}\quad{\Large $-48$}\end{tcolorbox}\hfill\begin{tcolorbox}[enhanced,nobeforeafter,width=7.1cm,height=2.1cm,valign=center,colback=crel!6,colframe=crel,arc=2.5mm,boxrule=0.9pt,left=3mm]{\Large\bfseries\color{crel}d)}\quad{\Large $22$}\end{tcolorbox}
\end{tcolorbox}
\newpage
\begin{tcolorbox}[enhanced,width=16cm,height=14cm,colback=white,colframe=crel,boxrule=2pt,arc=0mm,
left=6mm,right=6mm,top=0mm,bottom=4mm,sharp corners,
overlay={\fill[crel] (frame.north west) rectangle ([yshift=-1.25cm]frame.north east);
\node[anchor=west,text=white,font=\sffamily\bfseries\LARGE] at ([xshift=5mm,yshift=-0.62cm]frame.north west) {Question 7 B};
\node[anchor=east,text=white,font=\sffamily\itshape\large] at ([xshift=-5mm,yshift=-0.62cm]frame.north east) {Relatifs et fractions};}]
\vspace*{1.75cm}
\begin{minipage}[t][6.9cm][c]{\linewidth}\LARGE\centering
\begin{minipage}{0.95\linewidth}\raggedright\linespread{1.25}\selectfont Que vaut $\dfrac{3}{4}+\dfrac{1}{6}$~?\end{minipage}
\end{minipage}

\vfill
\begin{tcolorbox}[enhanced,nobeforeafter,width=7.1cm,height=2.1cm,valign=center,colback=crel!6,colframe=crel,arc=2.5mm,boxrule=0.9pt,left=3mm]{\Large\bfseries\color{crel}a)}\quad{\Large $\dfrac{11}{12}$}\end{tcolorbox}\hfill\begin{tcolorbox}[enhanced,nobeforeafter,width=7.1cm,height=2.1cm,valign=center,colback=crel!6,colframe=crel,arc=2.5mm,boxrule=0.9pt,left=3mm]{\Large\bfseries\color{crel}b)}\quad{\Large $\dfrac{1}{6}$}\end{tcolorbox}\par\vspace{2.5mm}
\begin{tcolorbox}[enhanced,nobeforeafter,width=7.1cm,height=2.1cm,valign=center,colback=crel!6,colframe=crel,arc=2.5mm,boxrule=0.9pt,left=3mm]{\Large\bfseries\color{crel}c)}\quad{\Large $\dfrac{2}{5}$}\end{tcolorbox}\hfill\begin{tcolorbox}[enhanced,nobeforeafter,width=7.1cm,height=2.1cm,valign=center,colback=crel!6,colframe=crel,arc=2.5mm,boxrule=0.9pt,left=3mm]{\Large\bfseries\color{crel}d)}\quad{\Large $\dfrac{1}{8}$}\end{tcolorbox}
\end{tcolorbox}
\newpage
\begin{tcolorbox}[enhanced,width=16cm,height=14cm,colback=white,colframe=crel,boxrule=2pt,arc=0mm,
left=6mm,right=6mm,top=0mm,bottom=4mm,sharp corners,
overlay={\fill[crel] (frame.north west) rectangle ([yshift=-1.25cm]frame.north east);
\node[anchor=west,text=white,font=\sffamily\bfseries\LARGE] at ([xshift=5mm,yshift=-0.62cm]frame.north west) {Question 8 B};
\node[anchor=east,text=white,font=\sffamily\itshape\large] at ([xshift=-5mm,yshift=-0.62cm]frame.north east) {Relatifs et fractions};}]
\vspace*{1.75cm}
\begin{minipage}[t][6.9cm][c]{\linewidth}\LARGE\centering
\begin{minipage}{0.95\linewidth}\raggedright\linespread{1.25}\selectfont Que vaut $\dfrac{1}{6}+\dfrac{2}{9}$~?\end{minipage}
\end{minipage}

\vfill
\begin{tcolorbox}[enhanced,nobeforeafter,width=7.1cm,height=2.1cm,valign=center,colback=crel!6,colframe=crel,arc=2.5mm,boxrule=0.9pt,left=3mm]{\Large\bfseries\color{crel}a)}\quad{\Large $\dfrac{1}{27}$}\end{tcolorbox}\hfill\begin{tcolorbox}[enhanced,nobeforeafter,width=7.1cm,height=2.1cm,valign=center,colback=crel!6,colframe=crel,arc=2.5mm,boxrule=0.9pt,left=3mm]{\Large\bfseries\color{crel}b)}\quad{\Large $\dfrac{7}{18}$}\end{tcolorbox}\par\vspace{2.5mm}
\begin{tcolorbox}[enhanced,nobeforeafter,width=7.1cm,height=2.1cm,valign=center,colback=crel!6,colframe=crel,arc=2.5mm,boxrule=0.9pt,left=3mm]{\Large\bfseries\color{crel}c)}\quad{\Large $\dfrac{1}{18}$}\end{tcolorbox}\hfill\begin{tcolorbox}[enhanced,nobeforeafter,width=7.1cm,height=2.1cm,valign=center,colback=crel!6,colframe=crel,arc=2.5mm,boxrule=0.9pt,left=3mm]{\Large\bfseries\color{crel}d)}\quad{\Large $\dfrac{1}{5}$}\end{tcolorbox}
\end{tcolorbox}
\newpage
\begin{tcolorbox}[enhanced,width=16cm,height=14cm,colback=white,colframe=crel,boxrule=2pt,arc=0mm,
left=6mm,right=6mm,top=0mm,bottom=4mm,sharp corners,
overlay={\fill[crel] (frame.north west) rectangle ([yshift=-1.25cm]frame.north east);
\node[anchor=west,text=white,font=\sffamily\bfseries\LARGE] at ([xshift=5mm,yshift=-0.62cm]frame.north west) {Question 9 B};
\node[anchor=east,text=white,font=\sffamily\itshape\large] at ([xshift=-5mm,yshift=-0.62cm]frame.north east) {Relatifs et fractions};}]
\vspace*{1.75cm}
\begin{minipage}[t][6.9cm][c]{\linewidth}\LARGE\centering
\begin{minipage}{0.95\linewidth}\raggedright\linespread{1.25}\selectfont Que vaut $\dfrac{4}{7}\times\left(-\dfrac{21}{8}\right)$~?\end{minipage}
\end{minipage}

\vfill
\begin{tcolorbox}[enhanced,nobeforeafter,width=7.1cm,height=2.1cm,valign=center,colback=crel!6,colframe=crel,arc=2.5mm,boxrule=0.9pt,left=3mm]{\Large\bfseries\color{crel}a)}\quad{\Large $-\dfrac{17}{15}$}\end{tcolorbox}\hfill\begin{tcolorbox}[enhanced,nobeforeafter,width=7.1cm,height=2.1cm,valign=center,colback=crel!6,colframe=crel,arc=2.5mm,boxrule=0.9pt,left=3mm]{\Large\bfseries\color{crel}b)}\quad{\Large $-\dfrac{32}{147}$}\end{tcolorbox}\par\vspace{2.5mm}
\begin{tcolorbox}[enhanced,nobeforeafter,width=7.1cm,height=2.1cm,valign=center,colback=crel!6,colframe=crel,arc=2.5mm,boxrule=0.9pt,left=3mm]{\Large\bfseries\color{crel}c)}\quad{\Large $\dfrac{3}{2}$}\end{tcolorbox}\hfill\begin{tcolorbox}[enhanced,nobeforeafter,width=7.1cm,height=2.1cm,valign=center,colback=crel!6,colframe=crel,arc=2.5mm,boxrule=0.9pt,left=3mm]{\Large\bfseries\color{crel}d)}\quad{\Large $-\dfrac{3}{2}$}\end{tcolorbox}
\end{tcolorbox}
\newpage
\begin{tcolorbox}[enhanced,width=16cm,height=14cm,colback=white,colframe=crel,boxrule=2pt,arc=0mm,
left=6mm,right=6mm,top=0mm,bottom=4mm,sharp corners,
overlay={\fill[crel] (frame.north west) rectangle ([yshift=-1.25cm]frame.north east);
\node[anchor=west,text=white,font=\sffamily\bfseries\LARGE] at ([xshift=5mm,yshift=-0.62cm]frame.north west) {Question 10 B};
\node[anchor=east,text=white,font=\sffamily\itshape\large] at ([xshift=-5mm,yshift=-0.62cm]frame.north east) {Relatifs et fractions};}]
\vspace*{1.75cm}
\begin{minipage}[t][6.9cm][c]{\linewidth}\LARGE\centering
\begin{minipage}{0.95\linewidth}\raggedright\linespread{1.25}\selectfont Que vaut $\left(-\dfrac{5}{6}\right)\times\left(-\dfrac{3}{10}\right)$~?\end{minipage}
\end{minipage}

\vfill
\begin{tcolorbox}[enhanced,nobeforeafter,width=7.1cm,height=2.1cm,valign=center,colback=crel!6,colframe=crel,arc=2.5mm,boxrule=0.9pt,left=3mm]{\Large\bfseries\color{crel}a)}\quad{\Large $-\dfrac{1}{2}$}\end{tcolorbox}\hfill\begin{tcolorbox}[enhanced,nobeforeafter,width=7.1cm,height=2.1cm,valign=center,colback=crel!6,colframe=crel,arc=2.5mm,boxrule=0.9pt,left=3mm]{\Large\bfseries\color{crel}b)}\quad{\Large $\dfrac{25}{9}$}\end{tcolorbox}\par\vspace{2.5mm}
\begin{tcolorbox}[enhanced,nobeforeafter,width=7.1cm,height=2.1cm,valign=center,colback=crel!6,colframe=crel,arc=2.5mm,boxrule=0.9pt,left=3mm]{\Large\bfseries\color{crel}c)}\quad{\Large $-\dfrac{1}{4}$}\end{tcolorbox}\hfill\begin{tcolorbox}[enhanced,nobeforeafter,width=7.1cm,height=2.1cm,valign=center,colback=crel!6,colframe=crel,arc=2.5mm,boxrule=0.9pt,left=3mm]{\Large\bfseries\color{crel}d)}\quad{\Large $\dfrac{1}{4}$}\end{tcolorbox}
\end{tcolorbox}
\newpage
\begin{tcolorbox}[enhanced,width=16cm,height=14cm,colback=white,colframe=crel,boxrule=2pt,arc=0mm,
left=6mm,right=6mm,top=0mm,bottom=4mm,sharp corners,
overlay={\fill[crel] (frame.north west) rectangle ([yshift=-1.25cm]frame.north east);
\node[anchor=west,text=white,font=\sffamily\bfseries\LARGE] at ([xshift=5mm,yshift=-0.62cm]frame.north west) {Question 11 B};
\node[anchor=east,text=white,font=\sffamily\itshape\large] at ([xshift=-5mm,yshift=-0.62cm]frame.north east) {Relatifs et fractions};}]
\vspace*{1.75cm}
\begin{minipage}[t][6.9cm][c]{\linewidth}\LARGE\centering
\begin{minipage}{0.95\linewidth}\raggedright\linespread{1.25}\selectfont Que vaut $\left(-\dfrac{5}{6}\right)\div\dfrac{10}{3}$~?\end{minipage}
\end{minipage}

\vfill
\begin{tcolorbox}[enhanced,nobeforeafter,width=7.1cm,height=2.1cm,valign=center,colback=crel!6,colframe=crel,arc=2.5mm,boxrule=0.9pt,left=3mm]{\Large\bfseries\color{crel}a)}\quad{\Large $\dfrac{1}{4}$}\end{tcolorbox}\hfill\begin{tcolorbox}[enhanced,nobeforeafter,width=7.1cm,height=2.1cm,valign=center,colback=crel!6,colframe=crel,arc=2.5mm,boxrule=0.9pt,left=3mm]{\Large\bfseries\color{crel}b)}\quad{\Large $-\dfrac{1}{4}$}\end{tcolorbox}\par\vspace{2.5mm}
\begin{tcolorbox}[enhanced,nobeforeafter,width=7.1cm,height=2.1cm,valign=center,colback=crel!6,colframe=crel,arc=2.5mm,boxrule=0.9pt,left=3mm]{\Large\bfseries\color{crel}c)}\quad{\Large $-\dfrac{25}{9}$}\end{tcolorbox}\hfill\begin{tcolorbox}[enhanced,nobeforeafter,width=7.1cm,height=2.1cm,valign=center,colback=crel!6,colframe=crel,arc=2.5mm,boxrule=0.9pt,left=3mm]{\Large\bfseries\color{crel}d)}\quad{\Large $-4$}\end{tcolorbox}
\end{tcolorbox}
\newpage
\begin{tcolorbox}[enhanced,width=16cm,height=14cm,colback=white,colframe=crel,boxrule=2pt,arc=0mm,
left=6mm,right=6mm,top=0mm,bottom=4mm,sharp corners,
overlay={\fill[crel] (frame.north west) rectangle ([yshift=-1.25cm]frame.north east);
\node[anchor=west,text=white,font=\sffamily\bfseries\LARGE] at ([xshift=5mm,yshift=-0.62cm]frame.north west) {Question 12 B};
\node[anchor=east,text=white,font=\sffamily\itshape\large] at ([xshift=-5mm,yshift=-0.62cm]frame.north east) {Relatifs et fractions};}]
\vspace*{1.75cm}
\begin{minipage}[t][6.9cm][c]{\linewidth}\LARGE\centering
\begin{minipage}{0.95\linewidth}\raggedright\linespread{1.25}\selectfont Que vaut $\dfrac{7}{10}\div\left(-\dfrac{21}{5}\right)$~?\end{minipage}
\end{minipage}

\vfill
\begin{tcolorbox}[enhanced,nobeforeafter,width=7.1cm,height=2.1cm,valign=center,colback=crel!6,colframe=crel,arc=2.5mm,boxrule=0.9pt,left=3mm]{\Large\bfseries\color{crel}a)}\quad{\Large $-6$}\end{tcolorbox}\hfill\begin{tcolorbox}[enhanced,nobeforeafter,width=7.1cm,height=2.1cm,valign=center,colback=crel!6,colframe=crel,arc=2.5mm,boxrule=0.9pt,left=3mm]{\Large\bfseries\color{crel}b)}\quad{\Large $-\dfrac{147}{50}$}\end{tcolorbox}\par\vspace{2.5mm}
\begin{tcolorbox}[enhanced,nobeforeafter,width=7.1cm,height=2.1cm,valign=center,colback=crel!6,colframe=crel,arc=2.5mm,boxrule=0.9pt,left=3mm]{\Large\bfseries\color{crel}c)}\quad{\Large $-\dfrac{1}{6}$}\end{tcolorbox}\hfill\begin{tcolorbox}[enhanced,nobeforeafter,width=7.1cm,height=2.1cm,valign=center,colback=crel!6,colframe=crel,arc=2.5mm,boxrule=0.9pt,left=3mm]{\Large\bfseries\color{crel}d)}\quad{\Large $\dfrac{1}{6}$}\end{tcolorbox}
\end{tcolorbox}
\newpage
\begin{tcolorbox}[enhanced,width=16cm,height=14cm,colback=white,colframe=crel,boxrule=2pt,arc=0mm,
left=6mm,right=6mm,top=0mm,bottom=4mm,sharp corners,
overlay={\fill[crel] (frame.north west) rectangle ([yshift=-1.25cm]frame.north east);
\node[anchor=west,text=white,font=\sffamily\bfseries\LARGE] at ([xshift=5mm,yshift=-0.62cm]frame.north west) {Question 13 B};
\node[anchor=east,text=white,font=\sffamily\itshape\large] at ([xshift=-5mm,yshift=-0.62cm]frame.north east) {Relatifs et fractions};}]
\vspace*{1.75cm}
\begin{minipage}[t][6.9cm][c]{\linewidth}\LARGE\centering
\begin{minipage}{0.95\linewidth}\raggedright\linespread{1.25}\selectfont Quel est l'inverse de $7$~?\end{minipage}
\end{minipage}

\vfill
\begin{tcolorbox}[enhanced,nobeforeafter,width=7.1cm,height=2.1cm,valign=center,colback=crel!6,colframe=crel,arc=2.5mm,boxrule=0.9pt,left=3mm]{\Large\bfseries\color{crel}a)}\quad{\Large $\dfrac{1}{7}$}\end{tcolorbox}\hfill\begin{tcolorbox}[enhanced,nobeforeafter,width=7.1cm,height=2.1cm,valign=center,colback=crel!6,colframe=crel,arc=2.5mm,boxrule=0.9pt,left=3mm]{\Large\bfseries\color{crel}b)}\quad{\Large $7$}\end{tcolorbox}\par\vspace{2.5mm}
\begin{tcolorbox}[enhanced,nobeforeafter,width=7.1cm,height=2.1cm,valign=center,colback=crel!6,colframe=crel,arc=2.5mm,boxrule=0.9pt,left=3mm]{\Large\bfseries\color{crel}c)}\quad{\Large $-\dfrac{1}{7}$}\end{tcolorbox}\hfill\begin{tcolorbox}[enhanced,nobeforeafter,width=7.1cm,height=2.1cm,valign=center,colback=crel!6,colframe=crel,arc=2.5mm,boxrule=0.9pt,left=3mm]{\Large\bfseries\color{crel}d)}\quad{\Large $-7$}\end{tcolorbox}
\end{tcolorbox}
\newpage
\begin{tcolorbox}[enhanced,width=16cm,height=14cm,colback=white,colframe=crel,boxrule=2pt,arc=0mm,
left=6mm,right=6mm,top=0mm,bottom=4mm,sharp corners,
overlay={\fill[crel] (frame.north west) rectangle ([yshift=-1.25cm]frame.north east);
\node[anchor=west,text=white,font=\sffamily\bfseries\LARGE] at ([xshift=5mm,yshift=-0.62cm]frame.north west) {Question 14 B};
\node[anchor=east,text=white,font=\sffamily\itshape\large] at ([xshift=-5mm,yshift=-0.62cm]frame.north east) {Relatifs et fractions};}]
\vspace*{1.75cm}
\begin{minipage}[t][6.9cm][c]{\linewidth}\LARGE\centering
\begin{minipage}{0.95\linewidth}\raggedright\linespread{1.25}\selectfont Quel est l'inverse de $\dfrac{1}{6}$~?\end{minipage}
\end{minipage}

\vfill
\begin{tcolorbox}[enhanced,nobeforeafter,width=7.1cm,height=2.1cm,valign=center,colback=crel!6,colframe=crel,arc=2.5mm,boxrule=0.9pt,left=3mm]{\Large\bfseries\color{crel}a)}\quad{\Large $\dfrac{1}{6}$}\end{tcolorbox}\hfill\begin{tcolorbox}[enhanced,nobeforeafter,width=7.1cm,height=2.1cm,valign=center,colback=crel!6,colframe=crel,arc=2.5mm,boxrule=0.9pt,left=3mm]{\Large\bfseries\color{crel}b)}\quad{\Large $-\dfrac{1}{6}$}\end{tcolorbox}\par\vspace{2.5mm}
\begin{tcolorbox}[enhanced,nobeforeafter,width=7.1cm,height=2.1cm,valign=center,colback=crel!6,colframe=crel,arc=2.5mm,boxrule=0.9pt,left=3mm]{\Large\bfseries\color{crel}c)}\quad{\Large $6$}\end{tcolorbox}\hfill\begin{tcolorbox}[enhanced,nobeforeafter,width=7.1cm,height=2.1cm,valign=center,colback=crel!6,colframe=crel,arc=2.5mm,boxrule=0.9pt,left=3mm]{\Large\bfseries\color{crel}d)}\quad{\Large $-6$}\end{tcolorbox}
\end{tcolorbox}
\newpage
\begin{tcolorbox}[enhanced,width=16cm,height=14cm,colback=white,colframe=crel,boxrule=2pt,arc=0mm,
left=6mm,right=6mm,top=0mm,bottom=4mm,sharp corners,
overlay={\fill[crel] (frame.north west) rectangle ([yshift=-1.25cm]frame.north east);
\node[anchor=west,text=white,font=\sffamily\bfseries\LARGE] at ([xshift=5mm,yshift=-0.62cm]frame.north west) {Question 15 B};
\node[anchor=east,text=white,font=\sffamily\itshape\large] at ([xshift=-5mm,yshift=-0.62cm]frame.north east) {Relatifs et fractions};}]
\vspace*{1.75cm}
\begin{minipage}[t][6.9cm][c]{\linewidth}\LARGE\centering
\begin{minipage}{0.95\linewidth}\raggedright\linespread{1.25}\selectfont Que vaut $\dfrac{5}{6}-\dfrac{1}{4}$~?\end{minipage}
\end{minipage}

\vfill
\begin{tcolorbox}[enhanced,nobeforeafter,width=7.1cm,height=2.1cm,valign=center,colback=crel!6,colframe=crel,arc=2.5mm,boxrule=0.9pt,left=3mm]{\Large\bfseries\color{crel}a)}\quad{\Large $-\dfrac{7}{12}$}\end{tcolorbox}\hfill\begin{tcolorbox}[enhanced,nobeforeafter,width=7.1cm,height=2.1cm,valign=center,colback=crel!6,colframe=crel,arc=2.5mm,boxrule=0.9pt,left=3mm]{\Large\bfseries\color{crel}b)}\quad{\Large $\dfrac{1}{6}$}\end{tcolorbox}\par\vspace{2.5mm}
\begin{tcolorbox}[enhanced,nobeforeafter,width=7.1cm,height=2.1cm,valign=center,colback=crel!6,colframe=crel,arc=2.5mm,boxrule=0.9pt,left=3mm]{\Large\bfseries\color{crel}c)}\quad{\Large $2$}\end{tcolorbox}\hfill\begin{tcolorbox}[enhanced,nobeforeafter,width=7.1cm,height=2.1cm,valign=center,colback=crel!6,colframe=crel,arc=2.5mm,boxrule=0.9pt,left=3mm]{\Large\bfseries\color{crel}d)}\quad{\Large $\dfrac{7}{12}$}\end{tcolorbox}
\end{tcolorbox}
\newpage
\begin{tcolorbox}[enhanced,width=16cm,height=14cm,colback=white,colframe=crel,boxrule=2pt,arc=0mm,
left=6mm,right=6mm,top=0mm,bottom=4mm,sharp corners,
overlay={\fill[crel] (frame.north west) rectangle ([yshift=-1.25cm]frame.north east);
\node[anchor=west,text=white,font=\sffamily\bfseries\LARGE] at ([xshift=5mm,yshift=-0.62cm]frame.north west) {Question 16 B};
\node[anchor=east,text=white,font=\sffamily\itshape\large] at ([xshift=-5mm,yshift=-0.62cm]frame.north east) {Relatifs et fractions};}]
\vspace*{1.75cm}
\begin{minipage}[t][6.9cm][c]{\linewidth}\LARGE\centering
\begin{minipage}{0.95\linewidth}\raggedright\linespread{1.25}\selectfont Que vaut $\dfrac{2}{7}-\dfrac{1}{2}$~?\end{minipage}
\end{minipage}

\vfill
\begin{tcolorbox}[enhanced,nobeforeafter,width=7.1cm,height=2.1cm,valign=center,colback=crel!6,colframe=crel,arc=2.5mm,boxrule=0.9pt,left=3mm]{\Large\bfseries\color{crel}a)}\quad{\Large $-\dfrac{3}{14}$}\end{tcolorbox}\hfill\begin{tcolorbox}[enhanced,nobeforeafter,width=7.1cm,height=2.1cm,valign=center,colback=crel!6,colframe=crel,arc=2.5mm,boxrule=0.9pt,left=3mm]{\Large\bfseries\color{crel}b)}\quad{\Large $\dfrac{1}{5}$}\end{tcolorbox}\par\vspace{2.5mm}
\begin{tcolorbox}[enhanced,nobeforeafter,width=7.1cm,height=2.1cm,valign=center,colback=crel!6,colframe=crel,arc=2.5mm,boxrule=0.9pt,left=3mm]{\Large\bfseries\color{crel}c)}\quad{\Large $\dfrac{1}{14}$}\end{tcolorbox}\hfill\begin{tcolorbox}[enhanced,nobeforeafter,width=7.1cm,height=2.1cm,valign=center,colback=crel!6,colframe=crel,arc=2.5mm,boxrule=0.9pt,left=3mm]{\Large\bfseries\color{crel}d)}\quad{\Large $\dfrac{3}{14}$}\end{tcolorbox}
\end{tcolorbox}
\newpage
\begin{tcolorbox}[enhanced,width=16cm,height=14cm,colback=white,colframe=cpui,boxrule=2pt,arc=0mm,
left=6mm,right=6mm,top=0mm,bottom=4mm,sharp corners,
overlay={\fill[cpui] (frame.north west) rectangle ([yshift=-1.25cm]frame.north east);
\node[anchor=west,text=white,font=\sffamily\bfseries\LARGE] at ([xshift=5mm,yshift=-0.62cm]frame.north west) {Question 17 B};
\node[anchor=east,text=white,font=\sffamily\itshape\large] at ([xshift=-5mm,yshift=-0.62cm]frame.north east) {Puissances et racines carrées};}]
\vspace*{1.75cm}
\begin{minipage}[t][6.9cm][c]{\linewidth}\LARGE\centering
\begin{minipage}{0.95\linewidth}\raggedright\linespread{1.25}\selectfont Que vaut $(-2)^{3}$~?\end{minipage}
\end{minipage}

\vfill
\begin{tcolorbox}[enhanced,nobeforeafter,width=7.1cm,height=2.1cm,valign=center,colback=cpui!6,colframe=cpui,arc=2.5mm,boxrule=0.9pt,left=3mm]{\Large\bfseries\color{cpui}a)}\quad{\Large $8$}\end{tcolorbox}\hfill\begin{tcolorbox}[enhanced,nobeforeafter,width=7.1cm,height=2.1cm,valign=center,colback=cpui!6,colframe=cpui,arc=2.5mm,boxrule=0.9pt,left=3mm]{\Large\bfseries\color{cpui}b)}\quad{\Large $4$}\end{tcolorbox}\par\vspace{2.5mm}
\begin{tcolorbox}[enhanced,nobeforeafter,width=7.1cm,height=2.1cm,valign=center,colback=cpui!6,colframe=cpui,arc=2.5mm,boxrule=0.9pt,left=3mm]{\Large\bfseries\color{cpui}c)}\quad{\Large $-8$}\end{tcolorbox}\hfill\begin{tcolorbox}[enhanced,nobeforeafter,width=7.1cm,height=2.1cm,valign=center,colback=cpui!6,colframe=cpui,arc=2.5mm,boxrule=0.9pt,left=3mm]{\Large\bfseries\color{cpui}d)}\quad{\Large $-6$}\end{tcolorbox}
\end{tcolorbox}
\newpage
\begin{tcolorbox}[enhanced,width=16cm,height=14cm,colback=white,colframe=cpui,boxrule=2pt,arc=0mm,
left=6mm,right=6mm,top=0mm,bottom=4mm,sharp corners,
overlay={\fill[cpui] (frame.north west) rectangle ([yshift=-1.25cm]frame.north east);
\node[anchor=west,text=white,font=\sffamily\bfseries\LARGE] at ([xshift=5mm,yshift=-0.62cm]frame.north west) {Question 18 B};
\node[anchor=east,text=white,font=\sffamily\itshape\large] at ([xshift=-5mm,yshift=-0.62cm]frame.north east) {Puissances et racines carrées};}]
\vspace*{1.75cm}
\begin{minipage}[t][6.9cm][c]{\linewidth}\LARGE\centering
\begin{minipage}{0.95\linewidth}\raggedright\linespread{1.25}\selectfont Que vaut $5^{3}$~?\end{minipage}
\end{minipage}

\vfill
\begin{tcolorbox}[enhanced,nobeforeafter,width=7.1cm,height=2.1cm,valign=center,colback=cpui!6,colframe=cpui,arc=2.5mm,boxrule=0.9pt,left=3mm]{\Large\bfseries\color{cpui}a)}\quad{\Large $-125$}\end{tcolorbox}\hfill\begin{tcolorbox}[enhanced,nobeforeafter,width=7.1cm,height=2.1cm,valign=center,colback=cpui!6,colframe=cpui,arc=2.5mm,boxrule=0.9pt,left=3mm]{\Large\bfseries\color{cpui}b)}\quad{\Large $15$}\end{tcolorbox}\par\vspace{2.5mm}
\begin{tcolorbox}[enhanced,nobeforeafter,width=7.1cm,height=2.1cm,valign=center,colback=cpui!6,colframe=cpui,arc=2.5mm,boxrule=0.9pt,left=3mm]{\Large\bfseries\color{cpui}c)}\quad{\Large $125$}\end{tcolorbox}\hfill\begin{tcolorbox}[enhanced,nobeforeafter,width=7.1cm,height=2.1cm,valign=center,colback=cpui!6,colframe=cpui,arc=2.5mm,boxrule=0.9pt,left=3mm]{\Large\bfseries\color{cpui}d)}\quad{\Large $25$}\end{tcolorbox}
\end{tcolorbox}
\newpage
\begin{tcolorbox}[enhanced,width=16cm,height=14cm,colback=white,colframe=cpui,boxrule=2pt,arc=0mm,
left=6mm,right=6mm,top=0mm,bottom=4mm,sharp corners,
overlay={\fill[cpui] (frame.north west) rectangle ([yshift=-1.25cm]frame.north east);
\node[anchor=west,text=white,font=\sffamily\bfseries\LARGE] at ([xshift=5mm,yshift=-0.62cm]frame.north west) {Question 19 B};
\node[anchor=east,text=white,font=\sffamily\itshape\large] at ([xshift=-5mm,yshift=-0.62cm]frame.north east) {Puissances et racines carrées};}]
\vspace*{1.75cm}
\begin{minipage}[t][6.9cm][c]{\linewidth}\LARGE\centering
\begin{minipage}{0.95\linewidth}\raggedright\linespread{1.25}\selectfont Écrire $10^{-3}\times10^{-4}$ sous la forme d'une seule puissance de $10$.\end{minipage}
\end{minipage}

\vfill
\begin{tcolorbox}[enhanced,nobeforeafter,width=7.1cm,height=2.1cm,valign=center,colback=cpui!6,colframe=cpui,arc=2.5mm,boxrule=0.9pt,left=3mm]{\Large\bfseries\color{cpui}a)}\quad{\Large $10^{12}$}\end{tcolorbox}\hfill\begin{tcolorbox}[enhanced,nobeforeafter,width=7.1cm,height=2.1cm,valign=center,colback=cpui!6,colframe=cpui,arc=2.5mm,boxrule=0.9pt,left=3mm]{\Large\bfseries\color{cpui}b)}\quad{\Large $100^{-7}$}\end{tcolorbox}\par\vspace{2.5mm}
\begin{tcolorbox}[enhanced,nobeforeafter,width=7.1cm,height=2.1cm,valign=center,colback=cpui!6,colframe=cpui,arc=2.5mm,boxrule=0.9pt,left=3mm]{\Large\bfseries\color{cpui}c)}\quad{\Large $10^{1}$}\end{tcolorbox}\hfill\begin{tcolorbox}[enhanced,nobeforeafter,width=7.1cm,height=2.1cm,valign=center,colback=cpui!6,colframe=cpui,arc=2.5mm,boxrule=0.9pt,left=3mm]{\Large\bfseries\color{cpui}d)}\quad{\Large $10^{-7}$}\end{tcolorbox}
\end{tcolorbox}
\newpage
\begin{tcolorbox}[enhanced,width=16cm,height=14cm,colback=white,colframe=cpui,boxrule=2pt,arc=0mm,
left=6mm,right=6mm,top=0mm,bottom=4mm,sharp corners,
overlay={\fill[cpui] (frame.north west) rectangle ([yshift=-1.25cm]frame.north east);
\node[anchor=west,text=white,font=\sffamily\bfseries\LARGE] at ([xshift=5mm,yshift=-0.62cm]frame.north west) {Question 20 B};
\node[anchor=east,text=white,font=\sffamily\itshape\large] at ([xshift=-5mm,yshift=-0.62cm]frame.north east) {Puissances et racines carrées};}]
\vspace*{1.75cm}
\begin{minipage}[t][6.9cm][c]{\linewidth}\LARGE\centering
\begin{minipage}{0.95\linewidth}\raggedright\linespread{1.25}\selectfont Écrire $10^{6}\times10^{2}$ sous la forme d'une seule puissance de $10$.\end{minipage}
\end{minipage}

\vfill
\begin{tcolorbox}[enhanced,nobeforeafter,width=7.1cm,height=2.1cm,valign=center,colback=cpui!6,colframe=cpui,arc=2.5mm,boxrule=0.9pt,left=3mm]{\Large\bfseries\color{cpui}a)}\quad{\Large $10^{12}$}\end{tcolorbox}\hfill\begin{tcolorbox}[enhanced,nobeforeafter,width=7.1cm,height=2.1cm,valign=center,colback=cpui!6,colframe=cpui,arc=2.5mm,boxrule=0.9pt,left=3mm]{\Large\bfseries\color{cpui}b)}\quad{\Large $100^{8}$}\end{tcolorbox}\par\vspace{2.5mm}
\begin{tcolorbox}[enhanced,nobeforeafter,width=7.1cm,height=2.1cm,valign=center,colback=cpui!6,colframe=cpui,arc=2.5mm,boxrule=0.9pt,left=3mm]{\Large\bfseries\color{cpui}c)}\quad{\Large $10^{4}$}\end{tcolorbox}\hfill\begin{tcolorbox}[enhanced,nobeforeafter,width=7.1cm,height=2.1cm,valign=center,colback=cpui!6,colframe=cpui,arc=2.5mm,boxrule=0.9pt,left=3mm]{\Large\bfseries\color{cpui}d)}\quad{\Large $10^{8}$}\end{tcolorbox}
\end{tcolorbox}
\newpage
\begin{tcolorbox}[enhanced,width=16cm,height=14cm,colback=white,colframe=cpui,boxrule=2pt,arc=0mm,
left=6mm,right=6mm,top=0mm,bottom=4mm,sharp corners,
overlay={\fill[cpui] (frame.north west) rectangle ([yshift=-1.25cm]frame.north east);
\node[anchor=west,text=white,font=\sffamily\bfseries\LARGE] at ([xshift=5mm,yshift=-0.62cm]frame.north west) {Question 21 B};
\node[anchor=east,text=white,font=\sffamily\itshape\large] at ([xshift=-5mm,yshift=-0.62cm]frame.north east) {Puissances et racines carrées};}]
\vspace*{1.75cm}
\begin{minipage}[t][6.9cm][c]{\linewidth}\LARGE\centering
\begin{minipage}{0.95\linewidth}\raggedright\linespread{1.25}\selectfont Quelle est l'écriture décimale de $10^{-1}$~?\end{minipage}
\end{minipage}

\vfill
\begin{tcolorbox}[enhanced,nobeforeafter,width=7.1cm,height=2.1cm,valign=center,colback=cpui!6,colframe=cpui,arc=2.5mm,boxrule=0.9pt,left=3mm]{\Large\bfseries\color{cpui}a)}\quad{\Large $0{,}1$}\end{tcolorbox}\hfill\begin{tcolorbox}[enhanced,nobeforeafter,width=7.1cm,height=2.1cm,valign=center,colback=cpui!6,colframe=cpui,arc=2.5mm,boxrule=0.9pt,left=3mm]{\Large\bfseries\color{cpui}b)}\quad{\Large $-10$}\end{tcolorbox}\par\vspace{2.5mm}
\begin{tcolorbox}[enhanced,nobeforeafter,width=7.1cm,height=2.1cm,valign=center,colback=cpui!6,colframe=cpui,arc=2.5mm,boxrule=0.9pt,left=3mm]{\Large\bfseries\color{cpui}c)}\quad{\Large $0{,}01$}\end{tcolorbox}\hfill\begin{tcolorbox}[enhanced,nobeforeafter,width=7.1cm,height=2.1cm,valign=center,colback=cpui!6,colframe=cpui,arc=2.5mm,boxrule=0.9pt,left=3mm]{\Large\bfseries\color{cpui}d)}\quad{\Large $10$}\end{tcolorbox}
\end{tcolorbox}
\newpage
\begin{tcolorbox}[enhanced,width=16cm,height=14cm,colback=white,colframe=cpui,boxrule=2pt,arc=0mm,
left=6mm,right=6mm,top=0mm,bottom=4mm,sharp corners,
overlay={\fill[cpui] (frame.north west) rectangle ([yshift=-1.25cm]frame.north east);
\node[anchor=west,text=white,font=\sffamily\bfseries\LARGE] at ([xshift=5mm,yshift=-0.62cm]frame.north west) {Question 22 B};
\node[anchor=east,text=white,font=\sffamily\itshape\large] at ([xshift=-5mm,yshift=-0.62cm]frame.north east) {Puissances et racines carrées};}]
\vspace*{1.75cm}
\begin{minipage}[t][6.9cm][c]{\linewidth}\LARGE\centering
\begin{minipage}{0.95\linewidth}\raggedright\linespread{1.25}\selectfont Quelle est l'écriture décimale de $10^{-4}$~?\end{minipage}
\end{minipage}

\vfill
\begin{tcolorbox}[enhanced,nobeforeafter,width=7.1cm,height=2.1cm,valign=center,colback=cpui!6,colframe=cpui,arc=2.5mm,boxrule=0.9pt,left=3mm]{\Large\bfseries\color{cpui}a)}\quad{\Large $0{,}0001$}\end{tcolorbox}\hfill\begin{tcolorbox}[enhanced,nobeforeafter,width=7.1cm,height=2.1cm,valign=center,colback=cpui!6,colframe=cpui,arc=2.5mm,boxrule=0.9pt,left=3mm]{\Large\bfseries\color{cpui}b)}\quad{\Large $0{,}00001$}\end{tcolorbox}\par\vspace{2.5mm}
\begin{tcolorbox}[enhanced,nobeforeafter,width=7.1cm,height=2.1cm,valign=center,colback=cpui!6,colframe=cpui,arc=2.5mm,boxrule=0.9pt,left=3mm]{\Large\bfseries\color{cpui}c)}\quad{\Large $0{,}001$}\end{tcolorbox}\hfill\begin{tcolorbox}[enhanced,nobeforeafter,width=7.1cm,height=2.1cm,valign=center,colback=cpui!6,colframe=cpui,arc=2.5mm,boxrule=0.9pt,left=3mm]{\Large\bfseries\color{cpui}d)}\quad{\Large $-10\,000$}\end{tcolorbox}
\end{tcolorbox}
\newpage
\begin{tcolorbox}[enhanced,width=16cm,height=14cm,colback=white,colframe=cpui,boxrule=2pt,arc=0mm,
left=6mm,right=6mm,top=0mm,bottom=4mm,sharp corners,
overlay={\fill[cpui] (frame.north west) rectangle ([yshift=-1.25cm]frame.north east);
\node[anchor=west,text=white,font=\sffamily\bfseries\LARGE] at ([xshift=5mm,yshift=-0.62cm]frame.north west) {Question 23 B};
\node[anchor=east,text=white,font=\sffamily\itshape\large] at ([xshift=-5mm,yshift=-0.62cm]frame.north east) {Puissances et racines carrées};}]
\vspace*{1.75cm}
\begin{minipage}[t][6.9cm][c]{\linewidth}\LARGE\centering
\begin{minipage}{0.95\linewidth}\raggedright\linespread{1.25}\selectfont Quelle est la notation scientifique de $72\,000$~?\end{minipage}
\end{minipage}

\vfill
\begin{tcolorbox}[enhanced,nobeforeafter,width=7.1cm,height=2.1cm,valign=center,colback=cpui!6,colframe=cpui,arc=2.5mm,boxrule=0.9pt,left=3mm]{\Large\bfseries\color{cpui}a)}\quad{\Large $72\times10^{4}$}\end{tcolorbox}\hfill\begin{tcolorbox}[enhanced,nobeforeafter,width=7.1cm,height=2.1cm,valign=center,colback=cpui!6,colframe=cpui,arc=2.5mm,boxrule=0.9pt,left=3mm]{\Large\bfseries\color{cpui}b)}\quad{\Large $7{,}2\times10^{-4}$}\end{tcolorbox}\par\vspace{2.5mm}
\begin{tcolorbox}[enhanced,nobeforeafter,width=7.1cm,height=2.1cm,valign=center,colback=cpui!6,colframe=cpui,arc=2.5mm,boxrule=0.9pt,left=3mm]{\Large\bfseries\color{cpui}c)}\quad{\Large $7{,}2\times10^{5}$}\end{tcolorbox}\hfill\begin{tcolorbox}[enhanced,nobeforeafter,width=7.1cm,height=2.1cm,valign=center,colback=cpui!6,colframe=cpui,arc=2.5mm,boxrule=0.9pt,left=3mm]{\Large\bfseries\color{cpui}d)}\quad{\Large $7{,}2\times10^{4}$}\end{tcolorbox}
\end{tcolorbox}
\newpage
\begin{tcolorbox}[enhanced,width=16cm,height=14cm,colback=white,colframe=cpui,boxrule=2pt,arc=0mm,
left=6mm,right=6mm,top=0mm,bottom=4mm,sharp corners,
overlay={\fill[cpui] (frame.north west) rectangle ([yshift=-1.25cm]frame.north east);
\node[anchor=west,text=white,font=\sffamily\bfseries\LARGE] at ([xshift=5mm,yshift=-0.62cm]frame.north west) {Question 24 B};
\node[anchor=east,text=white,font=\sffamily\itshape\large] at ([xshift=-5mm,yshift=-0.62cm]frame.north east) {Puissances et racines carrées};}]
\vspace*{1.75cm}
\begin{minipage}[t][6.9cm][c]{\linewidth}\LARGE\centering
\begin{minipage}{0.95\linewidth}\raggedright\linespread{1.25}\selectfont Quelle est la notation scientifique de $0{,}08$~?\end{minipage}
\end{minipage}

\vfill
\begin{tcolorbox}[enhanced,nobeforeafter,width=7.1cm,height=2.1cm,valign=center,colback=cpui!6,colframe=cpui,arc=2.5mm,boxrule=0.9pt,left=3mm]{\Large\bfseries\color{cpui}a)}\quad{\Large $8\times10^{2}$}\end{tcolorbox}\hfill\begin{tcolorbox}[enhanced,nobeforeafter,width=7.1cm,height=2.1cm,valign=center,colback=cpui!6,colframe=cpui,arc=2.5mm,boxrule=0.9pt,left=3mm]{\Large\bfseries\color{cpui}b)}\quad{\Large $8\times10^{-1}$}\end{tcolorbox}\par\vspace{2.5mm}
\begin{tcolorbox}[enhanced,nobeforeafter,width=7.1cm,height=2.1cm,valign=center,colback=cpui!6,colframe=cpui,arc=2.5mm,boxrule=0.9pt,left=3mm]{\Large\bfseries\color{cpui}c)}\quad{\Large $80\times10^{-2}$}\end{tcolorbox}\hfill\begin{tcolorbox}[enhanced,nobeforeafter,width=7.1cm,height=2.1cm,valign=center,colback=cpui!6,colframe=cpui,arc=2.5mm,boxrule=0.9pt,left=3mm]{\Large\bfseries\color{cpui}d)}\quad{\Large $8\times10^{-2}$}\end{tcolorbox}
\end{tcolorbox}
\newpage
\begin{tcolorbox}[enhanced,width=16cm,height=14cm,colback=white,colframe=cpui,boxrule=2pt,arc=0mm,
left=6mm,right=6mm,top=0mm,bottom=4mm,sharp corners,
overlay={\fill[cpui] (frame.north west) rectangle ([yshift=-1.25cm]frame.north east);
\node[anchor=west,text=white,font=\sffamily\bfseries\LARGE] at ([xshift=5mm,yshift=-0.62cm]frame.north west) {Question 25 B};
\node[anchor=east,text=white,font=\sffamily\itshape\large] at ([xshift=-5mm,yshift=-0.62cm]frame.north east) {Puissances et racines carrées};}]
\vspace*{1.75cm}
\begin{minipage}[t][6.9cm][c]{\linewidth}\LARGE\centering
\begin{minipage}{0.95\linewidth}\raggedright\linespread{1.25}\selectfont Que vaut $-5^{2}$~?\end{minipage}
\end{minipage}

\vfill
\begin{tcolorbox}[enhanced,nobeforeafter,width=7.1cm,height=2.1cm,valign=center,colback=cpui!6,colframe=cpui,arc=2.5mm,boxrule=0.9pt,left=3mm]{\Large\bfseries\color{cpui}a)}\quad{\Large $-25$}\end{tcolorbox}\hfill\begin{tcolorbox}[enhanced,nobeforeafter,width=7.1cm,height=2.1cm,valign=center,colback=cpui!6,colframe=cpui,arc=2.5mm,boxrule=0.9pt,left=3mm]{\Large\bfseries\color{cpui}b)}\quad{\Large $25$}\end{tcolorbox}\par\vspace{2.5mm}
\begin{tcolorbox}[enhanced,nobeforeafter,width=7.1cm,height=2.1cm,valign=center,colback=cpui!6,colframe=cpui,arc=2.5mm,boxrule=0.9pt,left=3mm]{\Large\bfseries\color{cpui}c)}\quad{\Large $10$}\end{tcolorbox}\hfill\begin{tcolorbox}[enhanced,nobeforeafter,width=7.1cm,height=2.1cm,valign=center,colback=cpui!6,colframe=cpui,arc=2.5mm,boxrule=0.9pt,left=3mm]{\Large\bfseries\color{cpui}d)}\quad{\Large $-10$}\end{tcolorbox}
\end{tcolorbox}
\newpage
\begin{tcolorbox}[enhanced,width=16cm,height=14cm,colback=white,colframe=cpui,boxrule=2pt,arc=0mm,
left=6mm,right=6mm,top=0mm,bottom=4mm,sharp corners,
overlay={\fill[cpui] (frame.north west) rectangle ([yshift=-1.25cm]frame.north east);
\node[anchor=west,text=white,font=\sffamily\bfseries\LARGE] at ([xshift=5mm,yshift=-0.62cm]frame.north west) {Question 26 B};
\node[anchor=east,text=white,font=\sffamily\itshape\large] at ([xshift=-5mm,yshift=-0.62cm]frame.north east) {Puissances et racines carrées};}]
\vspace*{1.75cm}
\begin{minipage}[t][6.9cm][c]{\linewidth}\LARGE\centering
\begin{minipage}{0.95\linewidth}\raggedright\linespread{1.25}\selectfont Que vaut $-6^{2}$~?\end{minipage}
\end{minipage}

\vfill
\begin{tcolorbox}[enhanced,nobeforeafter,width=7.1cm,height=2.1cm,valign=center,colback=cpui!6,colframe=cpui,arc=2.5mm,boxrule=0.9pt,left=3mm]{\Large\bfseries\color{cpui}a)}\quad{\Large $36$}\end{tcolorbox}\hfill\begin{tcolorbox}[enhanced,nobeforeafter,width=7.1cm,height=2.1cm,valign=center,colback=cpui!6,colframe=cpui,arc=2.5mm,boxrule=0.9pt,left=3mm]{\Large\bfseries\color{cpui}b)}\quad{\Large $-12$}\end{tcolorbox}\par\vspace{2.5mm}
\begin{tcolorbox}[enhanced,nobeforeafter,width=7.1cm,height=2.1cm,valign=center,colback=cpui!6,colframe=cpui,arc=2.5mm,boxrule=0.9pt,left=3mm]{\Large\bfseries\color{cpui}c)}\quad{\Large $12$}\end{tcolorbox}\hfill\begin{tcolorbox}[enhanced,nobeforeafter,width=7.1cm,height=2.1cm,valign=center,colback=cpui!6,colframe=cpui,arc=2.5mm,boxrule=0.9pt,left=3mm]{\Large\bfseries\color{cpui}d)}\quad{\Large $-36$}\end{tcolorbox}
\end{tcolorbox}
\newpage
\begin{tcolorbox}[enhanced,width=16cm,height=14cm,colback=white,colframe=cpui,boxrule=2pt,arc=0mm,
left=6mm,right=6mm,top=0mm,bottom=4mm,sharp corners,
overlay={\fill[cpui] (frame.north west) rectangle ([yshift=-1.25cm]frame.north east);
\node[anchor=west,text=white,font=\sffamily\bfseries\LARGE] at ([xshift=5mm,yshift=-0.62cm]frame.north west) {Question 27 B};
\node[anchor=east,text=white,font=\sffamily\itshape\large] at ([xshift=-5mm,yshift=-0.62cm]frame.north east) {Puissances et racines carrées};}]
\vspace*{1.75cm}
\begin{minipage}[t][6.9cm][c]{\linewidth}\LARGE\centering
\begin{minipage}{0.95\linewidth}\raggedright\linespread{1.25}\selectfont Écrire $\dfrac{10^{-2}}{10^{4}}$ sous la forme d'une seule puissance de $10$.\end{minipage}
\end{minipage}

\vfill
\begin{tcolorbox}[enhanced,nobeforeafter,width=7.1cm,height=2.1cm,valign=center,colback=cpui!6,colframe=cpui,arc=2.5mm,boxrule=0.9pt,left=3mm]{\Large\bfseries\color{cpui}a)}\quad{\Large $10^{-8}$}\end{tcolorbox}\hfill\begin{tcolorbox}[enhanced,nobeforeafter,width=7.1cm,height=2.1cm,valign=center,colback=cpui!6,colframe=cpui,arc=2.5mm,boxrule=0.9pt,left=3mm]{\Large\bfseries\color{cpui}b)}\quad{\Large $10^{-6}$}\end{tcolorbox}\par\vspace{2.5mm}
\begin{tcolorbox}[enhanced,nobeforeafter,width=7.1cm,height=2.1cm,valign=center,colback=cpui!6,colframe=cpui,arc=2.5mm,boxrule=0.9pt,left=3mm]{\Large\bfseries\color{cpui}c)}\quad{\Large $10^{2}$}\end{tcolorbox}\hfill\begin{tcolorbox}[enhanced,nobeforeafter,width=7.1cm,height=2.1cm,valign=center,colback=cpui!6,colframe=cpui,arc=2.5mm,boxrule=0.9pt,left=3mm]{\Large\bfseries\color{cpui}d)}\quad{\Large $10^{6}$}\end{tcolorbox}
\end{tcolorbox}
\newpage
\begin{tcolorbox}[enhanced,width=16cm,height=14cm,colback=white,colframe=cpui,boxrule=2pt,arc=0mm,
left=6mm,right=6mm,top=0mm,bottom=4mm,sharp corners,
overlay={\fill[cpui] (frame.north west) rectangle ([yshift=-1.25cm]frame.north east);
\node[anchor=west,text=white,font=\sffamily\bfseries\LARGE] at ([xshift=5mm,yshift=-0.62cm]frame.north west) {Question 28 B};
\node[anchor=east,text=white,font=\sffamily\itshape\large] at ([xshift=-5mm,yshift=-0.62cm]frame.north east) {Puissances et racines carrées};}]
\vspace*{1.75cm}
\begin{minipage}[t][6.9cm][c]{\linewidth}\LARGE\centering
\begin{minipage}{0.95\linewidth}\raggedright\linespread{1.25}\selectfont Écrire $\dfrac{10^{6}}{10^{-3}}$ sous la forme d'une seule puissance de $10$.\end{minipage}
\end{minipage}

\vfill
\begin{tcolorbox}[enhanced,nobeforeafter,width=7.1cm,height=2.1cm,valign=center,colback=cpui!6,colframe=cpui,arc=2.5mm,boxrule=0.9pt,left=3mm]{\Large\bfseries\color{cpui}a)}\quad{\Large $10^{-18}$}\end{tcolorbox}\hfill\begin{tcolorbox}[enhanced,nobeforeafter,width=7.1cm,height=2.1cm,valign=center,colback=cpui!6,colframe=cpui,arc=2.5mm,boxrule=0.9pt,left=3mm]{\Large\bfseries\color{cpui}b)}\quad{\Large $10^{-9}$}\end{tcolorbox}\par\vspace{2.5mm}
\begin{tcolorbox}[enhanced,nobeforeafter,width=7.1cm,height=2.1cm,valign=center,colback=cpui!6,colframe=cpui,arc=2.5mm,boxrule=0.9pt,left=3mm]{\Large\bfseries\color{cpui}c)}\quad{\Large $10^{3}$}\end{tcolorbox}\hfill\begin{tcolorbox}[enhanced,nobeforeafter,width=7.1cm,height=2.1cm,valign=center,colback=cpui!6,colframe=cpui,arc=2.5mm,boxrule=0.9pt,left=3mm]{\Large\bfseries\color{cpui}d)}\quad{\Large $10^{9}$}\end{tcolorbox}
\end{tcolorbox}
\newpage
\begin{tcolorbox}[enhanced,width=16cm,height=14cm,colback=white,colframe=cpui,boxrule=2pt,arc=0mm,
left=6mm,right=6mm,top=0mm,bottom=4mm,sharp corners,
overlay={\fill[cpui] (frame.north west) rectangle ([yshift=-1.25cm]frame.north east);
\node[anchor=west,text=white,font=\sffamily\bfseries\LARGE] at ([xshift=5mm,yshift=-0.62cm]frame.north west) {Question 29 B};
\node[anchor=east,text=white,font=\sffamily\itshape\large] at ([xshift=-5mm,yshift=-0.62cm]frame.north east) {Puissances et racines carrées};}]
\vspace*{1.75cm}
\begin{minipage}[t][6.9cm][c]{\linewidth}\LARGE\centering
\begin{minipage}{0.95\linewidth}\raggedright\linespread{1.25}\selectfont Que vaut $\sqrt{81}$~?\end{minipage}
\end{minipage}

\vfill
\begin{tcolorbox}[enhanced,nobeforeafter,width=7.1cm,height=2.1cm,valign=center,colback=cpui!6,colframe=cpui,arc=2.5mm,boxrule=0.9pt,left=3mm]{\Large\bfseries\color{cpui}a)}\quad{\Large $9$}\end{tcolorbox}\hfill\begin{tcolorbox}[enhanced,nobeforeafter,width=7.1cm,height=2.1cm,valign=center,colback=cpui!6,colframe=cpui,arc=2.5mm,boxrule=0.9pt,left=3mm]{\Large\bfseries\color{cpui}b)}\quad{\Large $18$}\end{tcolorbox}\par\vspace{2.5mm}
\begin{tcolorbox}[enhanced,nobeforeafter,width=7.1cm,height=2.1cm,valign=center,colback=cpui!6,colframe=cpui,arc=2.5mm,boxrule=0.9pt,left=3mm]{\Large\bfseries\color{cpui}c)}\quad{\Large $10$}\end{tcolorbox}\hfill\begin{tcolorbox}[enhanced,nobeforeafter,width=7.1cm,height=2.1cm,valign=center,colback=cpui!6,colframe=cpui,arc=2.5mm,boxrule=0.9pt,left=3mm]{\Large\bfseries\color{cpui}d)}\quad{\Large $40$}\end{tcolorbox}
\end{tcolorbox}
\newpage
\begin{tcolorbox}[enhanced,width=16cm,height=14cm,colback=white,colframe=cpui,boxrule=2pt,arc=0mm,
left=6mm,right=6mm,top=0mm,bottom=4mm,sharp corners,
overlay={\fill[cpui] (frame.north west) rectangle ([yshift=-1.25cm]frame.north east);
\node[anchor=west,text=white,font=\sffamily\bfseries\LARGE] at ([xshift=5mm,yshift=-0.62cm]frame.north west) {Question 30 B};
\node[anchor=east,text=white,font=\sffamily\itshape\large] at ([xshift=-5mm,yshift=-0.62cm]frame.north east) {Puissances et racines carrées};}]
\vspace*{1.75cm}
\begin{minipage}[t][6.9cm][c]{\linewidth}\LARGE\centering
\begin{minipage}{0.95\linewidth}\raggedright\linespread{1.25}\selectfont Que vaut $\sqrt{121}$~?\end{minipage}
\end{minipage}

\vfill
\begin{tcolorbox}[enhanced,nobeforeafter,width=7.1cm,height=2.1cm,valign=center,colback=cpui!6,colframe=cpui,arc=2.5mm,boxrule=0.9pt,left=3mm]{\Large\bfseries\color{cpui}a)}\quad{\Large $11$}\end{tcolorbox}\hfill\begin{tcolorbox}[enhanced,nobeforeafter,width=7.1cm,height=2.1cm,valign=center,colback=cpui!6,colframe=cpui,arc=2.5mm,boxrule=0.9pt,left=3mm]{\Large\bfseries\color{cpui}b)}\quad{\Large $22$}\end{tcolorbox}\par\vspace{2.5mm}
\begin{tcolorbox}[enhanced,nobeforeafter,width=7.1cm,height=2.1cm,valign=center,colback=cpui!6,colframe=cpui,arc=2.5mm,boxrule=0.9pt,left=3mm]{\Large\bfseries\color{cpui}c)}\quad{\Large $12$}\end{tcolorbox}\hfill\begin{tcolorbox}[enhanced,nobeforeafter,width=7.1cm,height=2.1cm,valign=center,colback=cpui!6,colframe=cpui,arc=2.5mm,boxrule=0.9pt,left=3mm]{\Large\bfseries\color{cpui}d)}\quad{\Large $60$}\end{tcolorbox}
\end{tcolorbox}
\newpage
\begin{tcolorbox}[enhanced,width=16cm,height=14cm,colback=white,colframe=cpui,boxrule=2pt,arc=0mm,
left=6mm,right=6mm,top=0mm,bottom=4mm,sharp corners,
overlay={\fill[cpui] (frame.north west) rectangle ([yshift=-1.25cm]frame.north east);
\node[anchor=west,text=white,font=\sffamily\bfseries\LARGE] at ([xshift=5mm,yshift=-0.62cm]frame.north west) {Question 31 B};
\node[anchor=east,text=white,font=\sffamily\itshape\large] at ([xshift=-5mm,yshift=-0.62cm]frame.north east) {Puissances et racines carrées};}]
\vspace*{1.75cm}
\begin{minipage}[t][6.9cm][c]{\linewidth}\LARGE\centering
\begin{minipage}{0.95\linewidth}\raggedright\linespread{1.25}\selectfont Entre quels nombres entiers consécutifs se trouve $\sqrt{70}$~?\end{minipage}
\end{minipage}

\vfill
\begin{tcolorbox}[enhanced,nobeforeafter,width=7.1cm,height=2.1cm,valign=center,colback=cpui!6,colframe=cpui,arc=2.5mm,boxrule=0.9pt,left=3mm]{\Large\bfseries\color{cpui}a)}\quad{\Large $8\text{ et }9$}\end{tcolorbox}\hfill\begin{tcolorbox}[enhanced,nobeforeafter,width=7.1cm,height=2.1cm,valign=center,colback=cpui!6,colframe=cpui,arc=2.5mm,boxrule=0.9pt,left=3mm]{\Large\bfseries\color{cpui}b)}\quad{\Large $9\text{ et }10$}\end{tcolorbox}\par\vspace{2.5mm}
\begin{tcolorbox}[enhanced,nobeforeafter,width=7.1cm,height=2.1cm,valign=center,colback=cpui!6,colframe=cpui,arc=2.5mm,boxrule=0.9pt,left=3mm]{\Large\bfseries\color{cpui}c)}\quad{\Large $7\text{ et }8$}\end{tcolorbox}\hfill\begin{tcolorbox}[enhanced,nobeforeafter,width=7.1cm,height=2.1cm,valign=center,colback=cpui!6,colframe=cpui,arc=2.5mm,boxrule=0.9pt,left=3mm]{\Large\bfseries\color{cpui}d)}\quad{\Large $35\text{ et }36$}\end{tcolorbox}
\end{tcolorbox}
\newpage
\begin{tcolorbox}[enhanced,width=16cm,height=14cm,colback=white,colframe=cpui,boxrule=2pt,arc=0mm,
left=6mm,right=6mm,top=0mm,bottom=4mm,sharp corners,
overlay={\fill[cpui] (frame.north west) rectangle ([yshift=-1.25cm]frame.north east);
\node[anchor=west,text=white,font=\sffamily\bfseries\LARGE] at ([xshift=5mm,yshift=-0.62cm]frame.north west) {Question 32 B};
\node[anchor=east,text=white,font=\sffamily\itshape\large] at ([xshift=-5mm,yshift=-0.62cm]frame.north east) {Puissances et racines carrées};}]
\vspace*{1.75cm}
\begin{minipage}[t][6.9cm][c]{\linewidth}\LARGE\centering
\begin{minipage}{0.95\linewidth}\raggedright\linespread{1.25}\selectfont Entre quels nombres entiers consécutifs se trouve $\sqrt{90}$~?\end{minipage}
\end{minipage}

\vfill
\begin{tcolorbox}[enhanced,nobeforeafter,width=7.1cm,height=2.1cm,valign=center,colback=cpui!6,colframe=cpui,arc=2.5mm,boxrule=0.9pt,left=3mm]{\Large\bfseries\color{cpui}a)}\quad{\Large $45\text{ et }46$}\end{tcolorbox}\hfill\begin{tcolorbox}[enhanced,nobeforeafter,width=7.1cm,height=2.1cm,valign=center,colback=cpui!6,colframe=cpui,arc=2.5mm,boxrule=0.9pt,left=3mm]{\Large\bfseries\color{cpui}b)}\quad{\Large $9\text{ et }10$}\end{tcolorbox}\par\vspace{2.5mm}
\begin{tcolorbox}[enhanced,nobeforeafter,width=7.1cm,height=2.1cm,valign=center,colback=cpui!6,colframe=cpui,arc=2.5mm,boxrule=0.9pt,left=3mm]{\Large\bfseries\color{cpui}c)}\quad{\Large $10\text{ et }11$}\end{tcolorbox}\hfill\begin{tcolorbox}[enhanced,nobeforeafter,width=7.1cm,height=2.1cm,valign=center,colback=cpui!6,colframe=cpui,arc=2.5mm,boxrule=0.9pt,left=3mm]{\Large\bfseries\color{cpui}d)}\quad{\Large $8\text{ et }9$}\end{tcolorbox}
\end{tcolorbox}
\newpage
\begin{tcolorbox}[enhanced,width=16cm,height=14cm,colback=white,colframe=clit,boxrule=2pt,arc=0mm,
left=6mm,right=6mm,top=0mm,bottom=4mm,sharp corners,
overlay={\fill[clit] (frame.north west) rectangle ([yshift=-1.25cm]frame.north east);
\node[anchor=west,text=white,font=\sffamily\bfseries\LARGE] at ([xshift=5mm,yshift=-0.62cm]frame.north west) {Question 33 B};
\node[anchor=east,text=white,font=\sffamily\itshape\large] at ([xshift=-5mm,yshift=-0.62cm]frame.north east) {Calcul littéral et équations};}]
\vspace*{1.75cm}
\begin{minipage}[t][6.9cm][c]{\linewidth}\LARGE\centering
\begin{minipage}{0.95\linewidth}\raggedright\linespread{1.25}\selectfont Développer $5(x-4)$.\end{minipage}
\end{minipage}

\vfill
\begin{tcolorbox}[enhanced,nobeforeafter,width=7.1cm,height=2.1cm,valign=center,colback=clit!6,colframe=clit,arc=2.5mm,boxrule=0.9pt,left=3mm]{\Large\bfseries\color{clit}a)}\quad{\Large $x-20$}\end{tcolorbox}\hfill\begin{tcolorbox}[enhanced,nobeforeafter,width=7.1cm,height=2.1cm,valign=center,colback=clit!6,colframe=clit,arc=2.5mm,boxrule=0.9pt,left=3mm]{\Large\bfseries\color{clit}b)}\quad{\Large $5x-4$}\end{tcolorbox}\par\vspace{2.5mm}
\begin{tcolorbox}[enhanced,nobeforeafter,width=7.1cm,height=2.1cm,valign=center,colback=clit!6,colframe=clit,arc=2.5mm,boxrule=0.9pt,left=3mm]{\Large\bfseries\color{clit}c)}\quad{\Large $5x+20$}\end{tcolorbox}\hfill\begin{tcolorbox}[enhanced,nobeforeafter,width=7.1cm,height=2.1cm,valign=center,colback=clit!6,colframe=clit,arc=2.5mm,boxrule=0.9pt,left=3mm]{\Large\bfseries\color{clit}d)}\quad{\Large $5x-20$}\end{tcolorbox}
\end{tcolorbox}
\newpage
\begin{tcolorbox}[enhanced,width=16cm,height=14cm,colback=white,colframe=clit,boxrule=2pt,arc=0mm,
left=6mm,right=6mm,top=0mm,bottom=4mm,sharp corners,
overlay={\fill[clit] (frame.north west) rectangle ([yshift=-1.25cm]frame.north east);
\node[anchor=west,text=white,font=\sffamily\bfseries\LARGE] at ([xshift=5mm,yshift=-0.62cm]frame.north west) {Question 34 B};
\node[anchor=east,text=white,font=\sffamily\itshape\large] at ([xshift=-5mm,yshift=-0.62cm]frame.north east) {Calcul littéral et équations};}]
\vspace*{1.75cm}
\begin{minipage}[t][6.9cm][c]{\linewidth}\LARGE\centering
\begin{minipage}{0.95\linewidth}\raggedright\linespread{1.25}\selectfont Développer $-3(-2x+7)$.\end{minipage}
\end{minipage}

\vfill
\begin{tcolorbox}[enhanced,nobeforeafter,width=7.1cm,height=2.1cm,valign=center,colback=clit!6,colframe=clit,arc=2.5mm,boxrule=0.9pt,left=3mm]{\Large\bfseries\color{clit}a)}\quad{\Large $6x-21$}\end{tcolorbox}\hfill\begin{tcolorbox}[enhanced,nobeforeafter,width=7.1cm,height=2.1cm,valign=center,colback=clit!6,colframe=clit,arc=2.5mm,boxrule=0.9pt,left=3mm]{\Large\bfseries\color{clit}b)}\quad{\Large $6x+7$}\end{tcolorbox}\par\vspace{2.5mm}
\begin{tcolorbox}[enhanced,nobeforeafter,width=7.1cm,height=2.1cm,valign=center,colback=clit!6,colframe=clit,arc=2.5mm,boxrule=0.9pt,left=3mm]{\Large\bfseries\color{clit}c)}\quad{\Large $6x+21$}\end{tcolorbox}\hfill\begin{tcolorbox}[enhanced,nobeforeafter,width=7.1cm,height=2.1cm,valign=center,colback=clit!6,colframe=clit,arc=2.5mm,boxrule=0.9pt,left=3mm]{\Large\bfseries\color{clit}d)}\quad{\Large $-2x-21$}\end{tcolorbox}
\end{tcolorbox}
\newpage
\begin{tcolorbox}[enhanced,width=16cm,height=14cm,colback=white,colframe=clit,boxrule=2pt,arc=0mm,
left=6mm,right=6mm,top=0mm,bottom=4mm,sharp corners,
overlay={\fill[clit] (frame.north west) rectangle ([yshift=-1.25cm]frame.north east);
\node[anchor=west,text=white,font=\sffamily\bfseries\LARGE] at ([xshift=5mm,yshift=-0.62cm]frame.north west) {Question 35 B};
\node[anchor=east,text=white,font=\sffamily\itshape\large] at ([xshift=-5mm,yshift=-0.62cm]frame.north east) {Calcul littéral et équations};}]
\vspace*{1.75cm}
\begin{minipage}[t][6.9cm][c]{\linewidth}\LARGE\centering
\begin{minipage}{0.95\linewidth}\raggedright\linespread{1.25}\selectfont Développer et réduire $(x-3)(x+5)$.\end{minipage}
\end{minipage}

\vfill
\begin{tcolorbox}[enhanced,nobeforeafter,width=7.1cm,height=2.1cm,valign=center,colback=clit!6,colframe=clit,arc=2.5mm,boxrule=0.9pt,left=3mm]{\Large\bfseries\color{clit}a)}\quad{\Large $x^2-15$}\end{tcolorbox}\hfill\begin{tcolorbox}[enhanced,nobeforeafter,width=7.1cm,height=2.1cm,valign=center,colback=clit!6,colframe=clit,arc=2.5mm,boxrule=0.9pt,left=3mm]{\Large\bfseries\color{clit}b)}\quad{\Large $x^2-15x+2$}\end{tcolorbox}\par\vspace{2.5mm}
\begin{tcolorbox}[enhanced,nobeforeafter,width=7.1cm,height=2.1cm,valign=center,colback=clit!6,colframe=clit,arc=2.5mm,boxrule=0.9pt,left=3mm]{\Large\bfseries\color{clit}c)}\quad{\Large $x^2+2x+2$}\end{tcolorbox}\hfill\begin{tcolorbox}[enhanced,nobeforeafter,width=7.1cm,height=2.1cm,valign=center,colback=clit!6,colframe=clit,arc=2.5mm,boxrule=0.9pt,left=3mm]{\Large\bfseries\color{clit}d)}\quad{\Large $x^2+2x-15$}\end{tcolorbox}
\end{tcolorbox}
\newpage
\begin{tcolorbox}[enhanced,width=16cm,height=14cm,colback=white,colframe=clit,boxrule=2pt,arc=0mm,
left=6mm,right=6mm,top=0mm,bottom=4mm,sharp corners,
overlay={\fill[clit] (frame.north west) rectangle ([yshift=-1.25cm]frame.north east);
\node[anchor=west,text=white,font=\sffamily\bfseries\LARGE] at ([xshift=5mm,yshift=-0.62cm]frame.north west) {Question 36 B};
\node[anchor=east,text=white,font=\sffamily\itshape\large] at ([xshift=-5mm,yshift=-0.62cm]frame.north east) {Calcul littéral et équations};}]
\vspace*{1.75cm}
\begin{minipage}[t][6.9cm][c]{\linewidth}\LARGE\centering
\begin{minipage}{0.95\linewidth}\raggedright\linespread{1.25}\selectfont Développer et réduire $(x-2)(x-6)$.\end{minipage}
\end{minipage}

\vfill
\begin{tcolorbox}[enhanced,nobeforeafter,width=7.1cm,height=2.1cm,valign=center,colback=clit!6,colframe=clit,arc=2.5mm,boxrule=0.9pt,left=3mm]{\Large\bfseries\color{clit}a)}\quad{\Large $x^2+12$}\end{tcolorbox}\hfill\begin{tcolorbox}[enhanced,nobeforeafter,width=7.1cm,height=2.1cm,valign=center,colback=clit!6,colframe=clit,arc=2.5mm,boxrule=0.9pt,left=3mm]{\Large\bfseries\color{clit}b)}\quad{\Large $x^2-8x-8$}\end{tcolorbox}\par\vspace{2.5mm}
\begin{tcolorbox}[enhanced,nobeforeafter,width=7.1cm,height=2.1cm,valign=center,colback=clit!6,colframe=clit,arc=2.5mm,boxrule=0.9pt,left=3mm]{\Large\bfseries\color{clit}c)}\quad{\Large $x^2-8x+12$}\end{tcolorbox}\hfill\begin{tcolorbox}[enhanced,nobeforeafter,width=7.1cm,height=2.1cm,valign=center,colback=clit!6,colframe=clit,arc=2.5mm,boxrule=0.9pt,left=3mm]{\Large\bfseries\color{clit}d)}\quad{\Large $x^2+12x-8$}\end{tcolorbox}
\end{tcolorbox}
\newpage
\begin{tcolorbox}[enhanced,width=16cm,height=14cm,colback=white,colframe=clit,boxrule=2pt,arc=0mm,
left=6mm,right=6mm,top=0mm,bottom=4mm,sharp corners,
overlay={\fill[clit] (frame.north west) rectangle ([yshift=-1.25cm]frame.north east);
\node[anchor=west,text=white,font=\sffamily\bfseries\LARGE] at ([xshift=5mm,yshift=-0.62cm]frame.north west) {Question 37 B};
\node[anchor=east,text=white,font=\sffamily\itshape\large] at ([xshift=-5mm,yshift=-0.62cm]frame.north east) {Calcul littéral et équations};}]
\vspace*{1.75cm}
\begin{minipage}[t][6.9cm][c]{\linewidth}\LARGE\centering
\begin{minipage}{0.95\linewidth}\raggedright\linespread{1.25}\selectfont Réduire l'expression $-2x+5+6x-9$.\end{minipage}
\end{minipage}

\vfill
\begin{tcolorbox}[enhanced,nobeforeafter,width=7.1cm,height=2.1cm,valign=center,colback=clit!6,colframe=clit,arc=2.5mm,boxrule=0.9pt,left=3mm]{\Large\bfseries\color{clit}a)}\quad{\Large $4x+14$}\end{tcolorbox}\hfill\begin{tcolorbox}[enhanced,nobeforeafter,width=7.1cm,height=2.1cm,valign=center,colback=clit!6,colframe=clit,arc=2.5mm,boxrule=0.9pt,left=3mm]{\Large\bfseries\color{clit}b)}\quad{\Large $-8x-4$}\end{tcolorbox}\par\vspace{2.5mm}
\begin{tcolorbox}[enhanced,nobeforeafter,width=7.1cm,height=2.1cm,valign=center,colback=clit!6,colframe=clit,arc=2.5mm,boxrule=0.9pt,left=3mm]{\Large\bfseries\color{clit}c)}\quad{\Large $0$}\end{tcolorbox}\hfill\begin{tcolorbox}[enhanced,nobeforeafter,width=7.1cm,height=2.1cm,valign=center,colback=clit!6,colframe=clit,arc=2.5mm,boxrule=0.9pt,left=3mm]{\Large\bfseries\color{clit}d)}\quad{\Large $4x-4$}\end{tcolorbox}
\end{tcolorbox}
\newpage
\begin{tcolorbox}[enhanced,width=16cm,height=14cm,colback=white,colframe=clit,boxrule=2pt,arc=0mm,
left=6mm,right=6mm,top=0mm,bottom=4mm,sharp corners,
overlay={\fill[clit] (frame.north west) rectangle ([yshift=-1.25cm]frame.north east);
\node[anchor=west,text=white,font=\sffamily\bfseries\LARGE] at ([xshift=5mm,yshift=-0.62cm]frame.north west) {Question 38 B};
\node[anchor=east,text=white,font=\sffamily\itshape\large] at ([xshift=-5mm,yshift=-0.62cm]frame.north east) {Calcul littéral et équations};}]
\vspace*{1.75cm}
\begin{minipage}[t][6.9cm][c]{\linewidth}\LARGE\centering
\begin{minipage}{0.95\linewidth}\raggedright\linespread{1.25}\selectfont Réduire l'expression $7x+1-7x+4$.\end{minipage}
\end{minipage}

\vfill
\begin{tcolorbox}[enhanced,nobeforeafter,width=7.1cm,height=2.1cm,valign=center,colback=clit!6,colframe=clit,arc=2.5mm,boxrule=0.9pt,left=3mm]{\Large\bfseries\color{clit}a)}\quad{\Large $14x+5$}\end{tcolorbox}\hfill\begin{tcolorbox}[enhanced,nobeforeafter,width=7.1cm,height=2.1cm,valign=center,colback=clit!6,colframe=clit,arc=2.5mm,boxrule=0.9pt,left=3mm]{\Large\bfseries\color{clit}b)}\quad{\Large $-3$}\end{tcolorbox}\par\vspace{2.5mm}
\begin{tcolorbox}[enhanced,nobeforeafter,width=7.1cm,height=2.1cm,valign=center,colback=clit!6,colframe=clit,arc=2.5mm,boxrule=0.9pt,left=3mm]{\Large\bfseries\color{clit}c)}\quad{\Large $5$}\end{tcolorbox}\hfill\begin{tcolorbox}[enhanced,nobeforeafter,width=7.1cm,height=2.1cm,valign=center,colback=clit!6,colframe=clit,arc=2.5mm,boxrule=0.9pt,left=3mm]{\Large\bfseries\color{clit}d)}\quad{\Large $5x$}\end{tcolorbox}
\end{tcolorbox}
\newpage
\begin{tcolorbox}[enhanced,width=16cm,height=14cm,colback=white,colframe=clit,boxrule=2pt,arc=0mm,
left=6mm,right=6mm,top=0mm,bottom=4mm,sharp corners,
overlay={\fill[clit] (frame.north west) rectangle ([yshift=-1.25cm]frame.north east);
\node[anchor=west,text=white,font=\sffamily\bfseries\LARGE] at ([xshift=5mm,yshift=-0.62cm]frame.north west) {Question 39 B};
\node[anchor=east,text=white,font=\sffamily\itshape\large] at ([xshift=-5mm,yshift=-0.62cm]frame.north east) {Calcul littéral et équations};}]
\vspace*{1.75cm}
\begin{minipage}[t][6.9cm][c]{\linewidth}\LARGE\centering
\begin{minipage}{0.95\linewidth}\raggedright\linespread{1.25}\selectfont Factoriser $5x+20$.\end{minipage}
\end{minipage}

\vfill
\begin{tcolorbox}[enhanced,nobeforeafter,width=7.1cm,height=2.1cm,valign=center,colback=clit!6,colframe=clit,arc=2.5mm,boxrule=0.9pt,left=3mm]{\Large\bfseries\color{clit}a)}\quad{\Large $5(5x+4)$}\end{tcolorbox}\hfill\begin{tcolorbox}[enhanced,nobeforeafter,width=7.1cm,height=2.1cm,valign=center,colback=clit!6,colframe=clit,arc=2.5mm,boxrule=0.9pt,left=3mm]{\Large\bfseries\color{clit}b)}\quad{\Large $5(x+4)$}\end{tcolorbox}\par\vspace{2.5mm}
\begin{tcolorbox}[enhanced,nobeforeafter,width=7.1cm,height=2.1cm,valign=center,colback=clit!6,colframe=clit,arc=2.5mm,boxrule=0.9pt,left=3mm]{\Large\bfseries\color{clit}c)}\quad{\Large $5x(x+4)$}\end{tcolorbox}\hfill\begin{tcolorbox}[enhanced,nobeforeafter,width=7.1cm,height=2.1cm,valign=center,colback=clit!6,colframe=clit,arc=2.5mm,boxrule=0.9pt,left=3mm]{\Large\bfseries\color{clit}d)}\quad{\Large $5(x+20)$}\end{tcolorbox}
\end{tcolorbox}
\newpage
\begin{tcolorbox}[enhanced,width=16cm,height=14cm,colback=white,colframe=clit,boxrule=2pt,arc=0mm,
left=6mm,right=6mm,top=0mm,bottom=4mm,sharp corners,
overlay={\fill[clit] (frame.north west) rectangle ([yshift=-1.25cm]frame.north east);
\node[anchor=west,text=white,font=\sffamily\bfseries\LARGE] at ([xshift=5mm,yshift=-0.62cm]frame.north west) {Question 40 B};
\node[anchor=east,text=white,font=\sffamily\itshape\large] at ([xshift=-5mm,yshift=-0.62cm]frame.north east) {Calcul littéral et équations};}]
\vspace*{1.75cm}
\begin{minipage}[t][6.9cm][c]{\linewidth}\LARGE\centering
\begin{minipage}{0.95\linewidth}\raggedright\linespread{1.25}\selectfont Factoriser $14x-21$.\end{minipage}
\end{minipage}

\vfill
\begin{tcolorbox}[enhanced,nobeforeafter,width=7.1cm,height=2.1cm,valign=center,colback=clit!6,colframe=clit,arc=2.5mm,boxrule=0.9pt,left=3mm]{\Large\bfseries\color{clit}a)}\quad{\Large $7(2x-3)$}\end{tcolorbox}\hfill\begin{tcolorbox}[enhanced,nobeforeafter,width=7.1cm,height=2.1cm,valign=center,colback=clit!6,colframe=clit,arc=2.5mm,boxrule=0.9pt,left=3mm]{\Large\bfseries\color{clit}b)}\quad{\Large $7(2x-21)$}\end{tcolorbox}\par\vspace{2.5mm}
\begin{tcolorbox}[enhanced,nobeforeafter,width=7.1cm,height=2.1cm,valign=center,colback=clit!6,colframe=clit,arc=2.5mm,boxrule=0.9pt,left=3mm]{\Large\bfseries\color{clit}c)}\quad{\Large $7(14x-3)$}\end{tcolorbox}\hfill\begin{tcolorbox}[enhanced,nobeforeafter,width=7.1cm,height=2.1cm,valign=center,colback=clit!6,colframe=clit,arc=2.5mm,boxrule=0.9pt,left=3mm]{\Large\bfseries\color{clit}d)}\quad{\Large $14(x-3)$}\end{tcolorbox}
\end{tcolorbox}
\newpage
\begin{tcolorbox}[enhanced,width=16cm,height=14cm,colback=white,colframe=clit,boxrule=2pt,arc=0mm,
left=6mm,right=6mm,top=0mm,bottom=4mm,sharp corners,
overlay={\fill[clit] (frame.north west) rectangle ([yshift=-1.25cm]frame.north east);
\node[anchor=west,text=white,font=\sffamily\bfseries\LARGE] at ([xshift=5mm,yshift=-0.62cm]frame.north west) {Question 41 B};
\node[anchor=east,text=white,font=\sffamily\itshape\large] at ([xshift=-5mm,yshift=-0.62cm]frame.north east) {Calcul littéral et équations};}]
\vspace*{1.75cm}
\begin{minipage}[t][6.9cm][c]{\linewidth}\LARGE\centering
\begin{minipage}{0.95\linewidth}\raggedright\linespread{1.25}\selectfont Que vaut $x^2-4x+6$ pour $x=3$~?\end{minipage}
\end{minipage}

\vfill
\begin{tcolorbox}[enhanced,nobeforeafter,width=7.1cm,height=2.1cm,valign=center,colback=clit!6,colframe=clit,arc=2.5mm,boxrule=0.9pt,left=3mm]{\Large\bfseries\color{clit}a)}\quad{\Large $0$}\end{tcolorbox}\hfill\begin{tcolorbox}[enhanced,nobeforeafter,width=7.1cm,height=2.1cm,valign=center,colback=clit!6,colframe=clit,arc=2.5mm,boxrule=0.9pt,left=3mm]{\Large\bfseries\color{clit}b)}\quad{\Large $-3$}\end{tcolorbox}\par\vspace{2.5mm}
\begin{tcolorbox}[enhanced,nobeforeafter,width=7.1cm,height=2.1cm,valign=center,colback=clit!6,colframe=clit,arc=2.5mm,boxrule=0.9pt,left=3mm]{\Large\bfseries\color{clit}c)}\quad{\Large $-15$}\end{tcolorbox}\hfill\begin{tcolorbox}[enhanced,nobeforeafter,width=7.1cm,height=2.1cm,valign=center,colback=clit!6,colframe=clit,arc=2.5mm,boxrule=0.9pt,left=3mm]{\Large\bfseries\color{clit}d)}\quad{\Large $3$}\end{tcolorbox}
\end{tcolorbox}
\newpage
\begin{tcolorbox}[enhanced,width=16cm,height=14cm,colback=white,colframe=clit,boxrule=2pt,arc=0mm,
left=6mm,right=6mm,top=0mm,bottom=4mm,sharp corners,
overlay={\fill[clit] (frame.north west) rectangle ([yshift=-1.25cm]frame.north east);
\node[anchor=west,text=white,font=\sffamily\bfseries\LARGE] at ([xshift=5mm,yshift=-0.62cm]frame.north west) {Question 42 B};
\node[anchor=east,text=white,font=\sffamily\itshape\large] at ([xshift=-5mm,yshift=-0.62cm]frame.north east) {Calcul littéral et équations};}]
\vspace*{1.75cm}
\begin{minipage}[t][6.9cm][c]{\linewidth}\LARGE\centering
\begin{minipage}{0.95\linewidth}\raggedright\linespread{1.25}\selectfont Que vaut $-x^2+3x+2$ pour $x=2$~?\end{minipage}
\end{minipage}

\vfill
\begin{tcolorbox}[enhanced,nobeforeafter,width=7.1cm,height=2.1cm,valign=center,colback=clit!6,colframe=clit,arc=2.5mm,boxrule=0.9pt,left=3mm]{\Large\bfseries\color{clit}a)}\quad{\Large $-4$}\end{tcolorbox}\hfill\begin{tcolorbox}[enhanced,nobeforeafter,width=7.1cm,height=2.1cm,valign=center,colback=clit!6,colframe=clit,arc=2.5mm,boxrule=0.9pt,left=3mm]{\Large\bfseries\color{clit}b)}\quad{\Large $4$}\end{tcolorbox}\par\vspace{2.5mm}
\begin{tcolorbox}[enhanced,nobeforeafter,width=7.1cm,height=2.1cm,valign=center,colback=clit!6,colframe=clit,arc=2.5mm,boxrule=0.9pt,left=3mm]{\Large\bfseries\color{clit}c)}\quad{\Large $5$}\end{tcolorbox}\hfill\begin{tcolorbox}[enhanced,nobeforeafter,width=7.1cm,height=2.1cm,valign=center,colback=clit!6,colframe=clit,arc=2.5mm,boxrule=0.9pt,left=3mm]{\Large\bfseries\color{clit}d)}\quad{\Large $12$}\end{tcolorbox}
\end{tcolorbox}
\newpage
\begin{tcolorbox}[enhanced,width=16cm,height=14cm,colback=white,colframe=clit,boxrule=2pt,arc=0mm,
left=6mm,right=6mm,top=0mm,bottom=4mm,sharp corners,
overlay={\fill[clit] (frame.north west) rectangle ([yshift=-1.25cm]frame.north east);
\node[anchor=west,text=white,font=\sffamily\bfseries\LARGE] at ([xshift=5mm,yshift=-0.62cm]frame.north west) {Question 43 B};
\node[anchor=east,text=white,font=\sffamily\itshape\large] at ([xshift=-5mm,yshift=-0.62cm]frame.north east) {Calcul littéral et équations};}]
\vspace*{1.75cm}
\begin{minipage}[t][6.9cm][c]{\linewidth}\LARGE\centering
\begin{minipage}{0.95\linewidth}\raggedright\linespread{1.25}\selectfont Quelle est la solution de l'équation $2x+9=-3$~?\end{minipage}
\end{minipage}

\vfill
\begin{tcolorbox}[enhanced,nobeforeafter,width=7.1cm,height=2.1cm,valign=center,colback=clit!6,colframe=clit,arc=2.5mm,boxrule=0.9pt,left=3mm]{\Large\bfseries\color{clit}a)}\quad{\Large $x=-6$}\end{tcolorbox}\hfill\begin{tcolorbox}[enhanced,nobeforeafter,width=7.1cm,height=2.1cm,valign=center,colback=clit!6,colframe=clit,arc=2.5mm,boxrule=0.9pt,left=3mm]{\Large\bfseries\color{clit}b)}\quad{\Large $x=3$}\end{tcolorbox}\par\vspace{2.5mm}
\begin{tcolorbox}[enhanced,nobeforeafter,width=7.1cm,height=2.1cm,valign=center,colback=clit!6,colframe=clit,arc=2.5mm,boxrule=0.9pt,left=3mm]{\Large\bfseries\color{clit}c)}\quad{\Large $x=-14$}\end{tcolorbox}\hfill\begin{tcolorbox}[enhanced,nobeforeafter,width=7.1cm,height=2.1cm,valign=center,colback=clit!6,colframe=clit,arc=2.5mm,boxrule=0.9pt,left=3mm]{\Large\bfseries\color{clit}d)}\quad{\Large $x=-\dfrac{1}{6}$}\end{tcolorbox}
\end{tcolorbox}
\newpage
\begin{tcolorbox}[enhanced,width=16cm,height=14cm,colback=white,colframe=clit,boxrule=2pt,arc=0mm,
left=6mm,right=6mm,top=0mm,bottom=4mm,sharp corners,
overlay={\fill[clit] (frame.north west) rectangle ([yshift=-1.25cm]frame.north east);
\node[anchor=west,text=white,font=\sffamily\bfseries\LARGE] at ([xshift=5mm,yshift=-0.62cm]frame.north west) {Question 44 B};
\node[anchor=east,text=white,font=\sffamily\itshape\large] at ([xshift=-5mm,yshift=-0.62cm]frame.north east) {Calcul littéral et équations};}]
\vspace*{1.75cm}
\begin{minipage}[t][6.9cm][c]{\linewidth}\LARGE\centering
\begin{minipage}{0.95\linewidth}\raggedright\linespread{1.25}\selectfont Quelle est la solution de l'équation $5x-3=27$~?\end{minipage}
\end{minipage}

\vfill
\begin{tcolorbox}[enhanced,nobeforeafter,width=7.1cm,height=2.1cm,valign=center,colback=clit!6,colframe=clit,arc=2.5mm,boxrule=0.9pt,left=3mm]{\Large\bfseries\color{clit}a)}\quad{\Large $x=\dfrac{24}{5}$}\end{tcolorbox}\hfill\begin{tcolorbox}[enhanced,nobeforeafter,width=7.1cm,height=2.1cm,valign=center,colback=clit!6,colframe=clit,arc=2.5mm,boxrule=0.9pt,left=3mm]{\Large\bfseries\color{clit}b)}\quad{\Large $x=25$}\end{tcolorbox}\par\vspace{2.5mm}
\begin{tcolorbox}[enhanced,nobeforeafter,width=7.1cm,height=2.1cm,valign=center,colback=clit!6,colframe=clit,arc=2.5mm,boxrule=0.9pt,left=3mm]{\Large\bfseries\color{clit}c)}\quad{\Large $x=6$}\end{tcolorbox}\hfill\begin{tcolorbox}[enhanced,nobeforeafter,width=7.1cm,height=2.1cm,valign=center,colback=clit!6,colframe=clit,arc=2.5mm,boxrule=0.9pt,left=3mm]{\Large\bfseries\color{clit}d)}\quad{\Large $x=\dfrac{1}{6}$}\end{tcolorbox}
\end{tcolorbox}
\newpage
\begin{tcolorbox}[enhanced,width=16cm,height=14cm,colback=white,colframe=clit,boxrule=2pt,arc=0mm,
left=6mm,right=6mm,top=0mm,bottom=4mm,sharp corners,
overlay={\fill[clit] (frame.north west) rectangle ([yshift=-1.25cm]frame.north east);
\node[anchor=west,text=white,font=\sffamily\bfseries\LARGE] at ([xshift=5mm,yshift=-0.62cm]frame.north west) {Question 45 B};
\node[anchor=east,text=white,font=\sffamily\itshape\large] at ([xshift=-5mm,yshift=-0.62cm]frame.north east) {Calcul littéral et équations};}]
\vspace*{1.75cm}
\begin{minipage}[t][6.9cm][c]{\linewidth}\LARGE\centering
\begin{minipage}{0.95\linewidth}\raggedright\linespread{1.25}\selectfont Quelle est la solution de l'équation $2x+9=5x-3$~?\end{minipage}
\end{minipage}

\vfill
\begin{tcolorbox}[enhanced,nobeforeafter,width=7.1cm,height=2.1cm,valign=center,colback=clit!6,colframe=clit,arc=2.5mm,boxrule=0.9pt,left=3mm]{\Large\bfseries\color{clit}a)}\quad{\Large $x=-2$}\end{tcolorbox}\hfill\begin{tcolorbox}[enhanced,nobeforeafter,width=7.1cm,height=2.1cm,valign=center,colback=clit!6,colframe=clit,arc=2.5mm,boxrule=0.9pt,left=3mm]{\Large\bfseries\color{clit}b)}\quad{\Large $x=-4$}\end{tcolorbox}\par\vspace{2.5mm}
\begin{tcolorbox}[enhanced,nobeforeafter,width=7.1cm,height=2.1cm,valign=center,colback=clit!6,colframe=clit,arc=2.5mm,boxrule=0.9pt,left=3mm]{\Large\bfseries\color{clit}c)}\quad{\Large $x=-\dfrac{12}{7}$}\end{tcolorbox}\hfill\begin{tcolorbox}[enhanced,nobeforeafter,width=7.1cm,height=2.1cm,valign=center,colback=clit!6,colframe=clit,arc=2.5mm,boxrule=0.9pt,left=3mm]{\Large\bfseries\color{clit}d)}\quad{\Large $x=4$}\end{tcolorbox}
\end{tcolorbox}
\newpage
\begin{tcolorbox}[enhanced,width=16cm,height=14cm,colback=white,colframe=clit,boxrule=2pt,arc=0mm,
left=6mm,right=6mm,top=0mm,bottom=4mm,sharp corners,
overlay={\fill[clit] (frame.north west) rectangle ([yshift=-1.25cm]frame.north east);
\node[anchor=west,text=white,font=\sffamily\bfseries\LARGE] at ([xshift=5mm,yshift=-0.62cm]frame.north west) {Question 46 B};
\node[anchor=east,text=white,font=\sffamily\itshape\large] at ([xshift=-5mm,yshift=-0.62cm]frame.north east) {Calcul littéral et équations};}]
\vspace*{1.75cm}
\begin{minipage}[t][6.9cm][c]{\linewidth}\LARGE\centering
\begin{minipage}{0.95\linewidth}\raggedright\linespread{1.25}\selectfont Quelle est la solution de l'équation $6x+1=-2x+17$~?\end{minipage}
\end{minipage}

\vfill
\begin{tcolorbox}[enhanced,nobeforeafter,width=7.1cm,height=2.1cm,valign=center,colback=clit!6,colframe=clit,arc=2.5mm,boxrule=0.9pt,left=3mm]{\Large\bfseries\color{clit}a)}\quad{\Large $x=2$}\end{tcolorbox}\hfill\begin{tcolorbox}[enhanced,nobeforeafter,width=7.1cm,height=2.1cm,valign=center,colback=clit!6,colframe=clit,arc=2.5mm,boxrule=0.9pt,left=3mm]{\Large\bfseries\color{clit}b)}\quad{\Large $x=4$}\end{tcolorbox}\par\vspace{2.5mm}
\begin{tcolorbox}[enhanced,nobeforeafter,width=7.1cm,height=2.1cm,valign=center,colback=clit!6,colframe=clit,arc=2.5mm,boxrule=0.9pt,left=3mm]{\Large\bfseries\color{clit}c)}\quad{\Large $x=-2$}\end{tcolorbox}\hfill\begin{tcolorbox}[enhanced,nobeforeafter,width=7.1cm,height=2.1cm,valign=center,colback=clit!6,colframe=clit,arc=2.5mm,boxrule=0.9pt,left=3mm]{\Large\bfseries\color{clit}d)}\quad{\Large $x=\dfrac{9}{4}$}\end{tcolorbox}
\end{tcolorbox}
\newpage
\begin{tcolorbox}[enhanced,width=16cm,height=14cm,colback=white,colframe=clit,boxrule=2pt,arc=0mm,
left=6mm,right=6mm,top=0mm,bottom=4mm,sharp corners,
overlay={\fill[clit] (frame.north west) rectangle ([yshift=-1.25cm]frame.north east);
\node[anchor=west,text=white,font=\sffamily\bfseries\LARGE] at ([xshift=5mm,yshift=-0.62cm]frame.north west) {Question 47 B};
\node[anchor=east,text=white,font=\sffamily\itshape\large] at ([xshift=-5mm,yshift=-0.62cm]frame.north east) {Calcul littéral et équations};}]
\vspace*{1.75cm}
\begin{minipage}[t][6.9cm][c]{\linewidth}\LARGE\centering
\begin{minipage}{0.95\linewidth}\raggedright\linespread{1.25}\selectfont Le quadruple d'un nombre, augmenté de 3, est égal à $27$. Quel est ce nombre~?\end{minipage}
\end{minipage}

\vfill
\begin{tcolorbox}[enhanced,nobeforeafter,width=7.1cm,height=2.1cm,valign=center,colback=clit!6,colframe=clit,arc=2.5mm,boxrule=0.9pt,left=3mm]{\Large\bfseries\color{clit}a)}\quad{\Large $6$}\end{tcolorbox}\hfill\begin{tcolorbox}[enhanced,nobeforeafter,width=7.1cm,height=2.1cm,valign=center,colback=clit!6,colframe=clit,arc=2.5mm,boxrule=0.9pt,left=3mm]{\Large\bfseries\color{clit}b)}\quad{\Large $\dfrac{27}{4}$}\end{tcolorbox}\par\vspace{2.5mm}
\begin{tcolorbox}[enhanced,nobeforeafter,width=7.1cm,height=2.1cm,valign=center,colback=clit!6,colframe=clit,arc=2.5mm,boxrule=0.9pt,left=3mm]{\Large\bfseries\color{clit}c)}\quad{\Large $\dfrac{15}{2}$}\end{tcolorbox}\hfill\begin{tcolorbox}[enhanced,nobeforeafter,width=7.1cm,height=2.1cm,valign=center,colback=clit!6,colframe=clit,arc=2.5mm,boxrule=0.9pt,left=3mm]{\Large\bfseries\color{clit}d)}\quad{\Large $24$}\end{tcolorbox}
\end{tcolorbox}
\newpage
\begin{tcolorbox}[enhanced,width=16cm,height=14cm,colback=white,colframe=clit,boxrule=2pt,arc=0mm,
left=6mm,right=6mm,top=0mm,bottom=4mm,sharp corners,
overlay={\fill[clit] (frame.north west) rectangle ([yshift=-1.25cm]frame.north east);
\node[anchor=west,text=white,font=\sffamily\bfseries\LARGE] at ([xshift=5mm,yshift=-0.62cm]frame.north west) {Question 48 B};
\node[anchor=east,text=white,font=\sffamily\itshape\large] at ([xshift=-5mm,yshift=-0.62cm]frame.north east) {Calcul littéral et équations};}]
\vspace*{1.75cm}
\begin{minipage}[t][6.9cm][c]{\linewidth}\LARGE\centering
\begin{minipage}{0.95\linewidth}\raggedright\linespread{1.25}\selectfont Le triple d'un nombre, diminué de 8, est égal à $16$. Quel est ce nombre~?\end{minipage}
\end{minipage}

\vfill
\begin{tcolorbox}[enhanced,nobeforeafter,width=7.1cm,height=2.1cm,valign=center,colback=clit!6,colframe=clit,arc=2.5mm,boxrule=0.9pt,left=3mm]{\Large\bfseries\color{clit}a)}\quad{\Large $\dfrac{8}{3}$}\end{tcolorbox}\hfill\begin{tcolorbox}[enhanced,nobeforeafter,width=7.1cm,height=2.1cm,valign=center,colback=clit!6,colframe=clit,arc=2.5mm,boxrule=0.9pt,left=3mm]{\Large\bfseries\color{clit}b)}\quad{\Large $\dfrac{16}{3}$}\end{tcolorbox}\par\vspace{2.5mm}
\begin{tcolorbox}[enhanced,nobeforeafter,width=7.1cm,height=2.1cm,valign=center,colback=clit!6,colframe=clit,arc=2.5mm,boxrule=0.9pt,left=3mm]{\Large\bfseries\color{clit}c)}\quad{\Large $8$}\end{tcolorbox}\hfill\begin{tcolorbox}[enhanced,nobeforeafter,width=7.1cm,height=2.1cm,valign=center,colback=clit!6,colframe=clit,arc=2.5mm,boxrule=0.9pt,left=3mm]{\Large\bfseries\color{clit}d)}\quad{\Large $24$}\end{tcolorbox}
\end{tcolorbox}
\newpage
\begin{tcolorbox}[enhanced,width=16cm,height=14cm,colback=white,colframe=cgeo,boxrule=2pt,arc=0mm,
left=6mm,right=6mm,top=0mm,bottom=4mm,sharp corners,
overlay={\fill[cgeo] (frame.north west) rectangle ([yshift=-1.25cm]frame.north east);
\node[anchor=west,text=white,font=\sffamily\bfseries\LARGE] at ([xshift=5mm,yshift=-0.62cm]frame.north west) {Question 49 B};
\node[anchor=east,text=white,font=\sffamily\itshape\large] at ([xshift=-5mm,yshift=-0.62cm]frame.north east) {Géométrie};}]
\vspace*{1.75cm}
\begin{minipage}[t][6.9cm][c]{\linewidth}\LARGE\centering
\begin{minipage}{0.95\linewidth}\raggedright\linespread{1.25}\selectfont Le triangle $ABC$ est rectangle en $A$, avec $AB=5$ cm et $AC=12$ cm. Combien mesure $BC$~?\end{minipage}
\end{minipage}

\vfill
\begin{tcolorbox}[enhanced,nobeforeafter,width=7.1cm,height=2.1cm,valign=center,colback=cgeo!6,colframe=cgeo,arc=2.5mm,boxrule=0.9pt,left=3mm]{\Large\bfseries\color{cgeo}a)}\quad{\Large $17~\text{cm}$}\end{tcolorbox}\hfill\begin{tcolorbox}[enhanced,nobeforeafter,width=7.1cm,height=2.1cm,valign=center,colback=cgeo!6,colframe=cgeo,arc=2.5mm,boxrule=0.9pt,left=3mm]{\Large\bfseries\color{cgeo}b)}\quad{\Large $169~\text{cm}$}\end{tcolorbox}\par\vspace{2.5mm}
\begin{tcolorbox}[enhanced,nobeforeafter,width=7.1cm,height=2.1cm,valign=center,colback=cgeo!6,colframe=cgeo,arc=2.5mm,boxrule=0.9pt,left=3mm]{\Large\bfseries\color{cgeo}c)}\quad{\Large $13~\text{cm}$}\end{tcolorbox}\hfill\begin{tcolorbox}[enhanced,nobeforeafter,width=7.1cm,height=2.1cm,valign=center,colback=cgeo!6,colframe=cgeo,arc=2.5mm,boxrule=0.9pt,left=3mm]{\Large\bfseries\color{cgeo}d)}\quad{\Large $\sqrt{119}~\text{cm}$}\end{tcolorbox}
\end{tcolorbox}
\newpage
\begin{tcolorbox}[enhanced,width=16cm,height=14cm,colback=white,colframe=cgeo,boxrule=2pt,arc=0mm,
left=6mm,right=6mm,top=0mm,bottom=4mm,sharp corners,
overlay={\fill[cgeo] (frame.north west) rectangle ([yshift=-1.25cm]frame.north east);
\node[anchor=west,text=white,font=\sffamily\bfseries\LARGE] at ([xshift=5mm,yshift=-0.62cm]frame.north west) {Question 50 B};
\node[anchor=east,text=white,font=\sffamily\itshape\large] at ([xshift=-5mm,yshift=-0.62cm]frame.north east) {Géométrie};}]
\vspace*{1.75cm}
\begin{minipage}[t][6.9cm][c]{\linewidth}\LARGE\centering
\begin{minipage}{0.95\linewidth}\raggedright\linespread{1.25}\selectfont Le triangle $ABC$ est rectangle en $A$, avec $AB=9$ cm et $AC=12$ cm. Combien mesure $BC$~?\end{minipage}
\end{minipage}

\vfill
\begin{tcolorbox}[enhanced,nobeforeafter,width=7.1cm,height=2.1cm,valign=center,colback=cgeo!6,colframe=cgeo,arc=2.5mm,boxrule=0.9pt,left=3mm]{\Large\bfseries\color{cgeo}a)}\quad{\Large $\sqrt{63}~\text{cm}$}\end{tcolorbox}\hfill\begin{tcolorbox}[enhanced,nobeforeafter,width=7.1cm,height=2.1cm,valign=center,colback=cgeo!6,colframe=cgeo,arc=2.5mm,boxrule=0.9pt,left=3mm]{\Large\bfseries\color{cgeo}b)}\quad{\Large $15~\text{cm}$}\end{tcolorbox}\par\vspace{2.5mm}
\begin{tcolorbox}[enhanced,nobeforeafter,width=7.1cm,height=2.1cm,valign=center,colback=cgeo!6,colframe=cgeo,arc=2.5mm,boxrule=0.9pt,left=3mm]{\Large\bfseries\color{cgeo}c)}\quad{\Large $225~\text{cm}$}\end{tcolorbox}\hfill\begin{tcolorbox}[enhanced,nobeforeafter,width=7.1cm,height=2.1cm,valign=center,colback=cgeo!6,colframe=cgeo,arc=2.5mm,boxrule=0.9pt,left=3mm]{\Large\bfseries\color{cgeo}d)}\quad{\Large $21~\text{cm}$}\end{tcolorbox}
\end{tcolorbox}
\newpage
\begin{tcolorbox}[enhanced,width=16cm,height=14cm,colback=white,colframe=cgeo,boxrule=2pt,arc=0mm,
left=6mm,right=6mm,top=0mm,bottom=4mm,sharp corners,
overlay={\fill[cgeo] (frame.north west) rectangle ([yshift=-1.25cm]frame.north east);
\node[anchor=west,text=white,font=\sffamily\bfseries\LARGE] at ([xshift=5mm,yshift=-0.62cm]frame.north west) {Question 51 B};
\node[anchor=east,text=white,font=\sffamily\itshape\large] at ([xshift=-5mm,yshift=-0.62cm]frame.north east) {Géométrie};}]
\vspace*{1.75cm}
\begin{minipage}[t][6.9cm][c]{\linewidth}\LARGE\centering
\begin{minipage}{0.95\linewidth}\raggedright\linespread{1.25}\selectfont Le triangle $ABC$ est rectangle en $A$, avec $BC=25$ cm et $AB=7$ cm. Combien mesure $AC$~?\end{minipage}
\end{minipage}

\vfill
\begin{tcolorbox}[enhanced,nobeforeafter,width=7.1cm,height=2.1cm,valign=center,colback=cgeo!6,colframe=cgeo,arc=2.5mm,boxrule=0.9pt,left=3mm]{\Large\bfseries\color{cgeo}a)}\quad{\Large $24~\text{cm}$}\end{tcolorbox}\hfill\begin{tcolorbox}[enhanced,nobeforeafter,width=7.1cm,height=2.1cm,valign=center,colback=cgeo!6,colframe=cgeo,arc=2.5mm,boxrule=0.9pt,left=3mm]{\Large\bfseries\color{cgeo}b)}\quad{\Large $18~\text{cm}$}\end{tcolorbox}\par\vspace{2.5mm}
\begin{tcolorbox}[enhanced,nobeforeafter,width=7.1cm,height=2.1cm,valign=center,colback=cgeo!6,colframe=cgeo,arc=2.5mm,boxrule=0.9pt,left=3mm]{\Large\bfseries\color{cgeo}c)}\quad{\Large $576~\text{cm}$}\end{tcolorbox}\hfill\begin{tcolorbox}[enhanced,nobeforeafter,width=7.1cm,height=2.1cm,valign=center,colback=cgeo!6,colframe=cgeo,arc=2.5mm,boxrule=0.9pt,left=3mm]{\Large\bfseries\color{cgeo}d)}\quad{\Large $\sqrt{674}~\text{cm}$}\end{tcolorbox}
\end{tcolorbox}
\newpage
\begin{tcolorbox}[enhanced,width=16cm,height=14cm,colback=white,colframe=cgeo,boxrule=2pt,arc=0mm,
left=6mm,right=6mm,top=0mm,bottom=4mm,sharp corners,
overlay={\fill[cgeo] (frame.north west) rectangle ([yshift=-1.25cm]frame.north east);
\node[anchor=west,text=white,font=\sffamily\bfseries\LARGE] at ([xshift=5mm,yshift=-0.62cm]frame.north west) {Question 52 B};
\node[anchor=east,text=white,font=\sffamily\itshape\large] at ([xshift=-5mm,yshift=-0.62cm]frame.north east) {Géométrie};}]
\vspace*{1.75cm}
\begin{minipage}[t][6.9cm][c]{\linewidth}\LARGE\centering
\begin{minipage}{0.95\linewidth}\raggedright\linespread{1.25}\selectfont Le triangle $ABC$ est rectangle en $A$, avec $BC=17$ cm et $AB=8$ cm. Combien mesure $AC$~?\end{minipage}
\end{minipage}

\vfill
\begin{tcolorbox}[enhanced,nobeforeafter,width=7.1cm,height=2.1cm,valign=center,colback=cgeo!6,colframe=cgeo,arc=2.5mm,boxrule=0.9pt,left=3mm]{\Large\bfseries\color{cgeo}a)}\quad{\Large $9~\text{cm}$}\end{tcolorbox}\hfill\begin{tcolorbox}[enhanced,nobeforeafter,width=7.1cm,height=2.1cm,valign=center,colback=cgeo!6,colframe=cgeo,arc=2.5mm,boxrule=0.9pt,left=3mm]{\Large\bfseries\color{cgeo}b)}\quad{\Large $\sqrt{353}~\text{cm}$}\end{tcolorbox}\par\vspace{2.5mm}
\begin{tcolorbox}[enhanced,nobeforeafter,width=7.1cm,height=2.1cm,valign=center,colback=cgeo!6,colframe=cgeo,arc=2.5mm,boxrule=0.9pt,left=3mm]{\Large\bfseries\color{cgeo}c)}\quad{\Large $15~\text{cm}$}\end{tcolorbox}\hfill\begin{tcolorbox}[enhanced,nobeforeafter,width=7.1cm,height=2.1cm,valign=center,colback=cgeo!6,colframe=cgeo,arc=2.5mm,boxrule=0.9pt,left=3mm]{\Large\bfseries\color{cgeo}d)}\quad{\Large $225~\text{cm}$}\end{tcolorbox}
\end{tcolorbox}
\newpage
\begin{tcolorbox}[enhanced,width=16cm,height=14cm,colback=white,colframe=cgeo,boxrule=2pt,arc=0mm,
left=6mm,right=6mm,top=0mm,bottom=4mm,sharp corners,
overlay={\fill[cgeo] (frame.north west) rectangle ([yshift=-1.25cm]frame.north east);
\node[anchor=west,text=white,font=\sffamily\bfseries\LARGE] at ([xshift=5mm,yshift=-0.62cm]frame.north west) {Question 53 B};
\node[anchor=east,text=white,font=\sffamily\itshape\large] at ([xshift=-5mm,yshift=-0.62cm]frame.north east) {Géométrie};}]
\vspace*{1.75cm}
\begin{minipage}[t][6.9cm][c]{\linewidth}\LARGE\centering
\begin{minipage}{0.95\linewidth}\raggedright\linespread{1.25}\selectfont Un triangle a des côtés de $9$ cm, $40$ cm et $41$ cm. Est-il rectangle~?\end{minipage}
\end{minipage}

\vfill
\begin{tcolorbox}[enhanced,nobeforeafter,width=7.1cm,height=2.1cm,valign=center,colback=cgeo!6,colframe=cgeo,arc=2.5mm,boxrule=0.9pt,left=3mm]{\Large\bfseries\color{cgeo}a)}\quad{\Large non}\end{tcolorbox}\hfill\begin{tcolorbox}[enhanced,nobeforeafter,width=7.1cm,height=2.1cm,valign=center,colback=cgeo!6,colframe=cgeo,arc=2.5mm,boxrule=0.9pt,left=3mm]{\Large\bfseries\color{cgeo}b)}\quad{\Large s'il est isocèle}\end{tcolorbox}\par\vspace{2.5mm}
\begin{tcolorbox}[enhanced,nobeforeafter,width=7.1cm,height=2.1cm,valign=center,colback=cgeo!6,colframe=cgeo,arc=2.5mm,boxrule=0.9pt,left=3mm]{\Large\bfseries\color{cgeo}c)}\quad{\Large on ne peut pas savoir}\end{tcolorbox}\hfill\begin{tcolorbox}[enhanced,nobeforeafter,width=7.1cm,height=2.1cm,valign=center,colback=cgeo!6,colframe=cgeo,arc=2.5mm,boxrule=0.9pt,left=3mm]{\Large\bfseries\color{cgeo}d)}\quad{\Large oui}\end{tcolorbox}
\end{tcolorbox}
\newpage
\begin{tcolorbox}[enhanced,width=16cm,height=14cm,colback=white,colframe=cgeo,boxrule=2pt,arc=0mm,
left=6mm,right=6mm,top=0mm,bottom=4mm,sharp corners,
overlay={\fill[cgeo] (frame.north west) rectangle ([yshift=-1.25cm]frame.north east);
\node[anchor=west,text=white,font=\sffamily\bfseries\LARGE] at ([xshift=5mm,yshift=-0.62cm]frame.north west) {Question 54 B};
\node[anchor=east,text=white,font=\sffamily\itshape\large] at ([xshift=-5mm,yshift=-0.62cm]frame.north east) {Géométrie};}]
\vspace*{1.75cm}
\begin{minipage}[t][6.9cm][c]{\linewidth}\LARGE\centering
\begin{minipage}{0.95\linewidth}\raggedright\linespread{1.25}\selectfont Un triangle a des côtés de $4$ cm, $6$ cm et $7$ cm. Est-il rectangle~?\end{minipage}
\end{minipage}

\vfill
\begin{tcolorbox}[enhanced,nobeforeafter,width=7.1cm,height=2.1cm,valign=center,colback=cgeo!6,colframe=cgeo,arc=2.5mm,boxrule=0.9pt,left=3mm]{\Large\bfseries\color{cgeo}a)}\quad{\Large s'il est isocèle}\end{tcolorbox}\hfill\begin{tcolorbox}[enhanced,nobeforeafter,width=7.1cm,height=2.1cm,valign=center,colback=cgeo!6,colframe=cgeo,arc=2.5mm,boxrule=0.9pt,left=3mm]{\Large\bfseries\color{cgeo}b)}\quad{\Large oui}\end{tcolorbox}\par\vspace{2.5mm}
\begin{tcolorbox}[enhanced,nobeforeafter,width=7.1cm,height=2.1cm,valign=center,colback=cgeo!6,colframe=cgeo,arc=2.5mm,boxrule=0.9pt,left=3mm]{\Large\bfseries\color{cgeo}c)}\quad{\Large non}\end{tcolorbox}\hfill\begin{tcolorbox}[enhanced,nobeforeafter,width=7.1cm,height=2.1cm,valign=center,colback=cgeo!6,colframe=cgeo,arc=2.5mm,boxrule=0.9pt,left=3mm]{\Large\bfseries\color{cgeo}d)}\quad{\Large on ne peut pas savoir}\end{tcolorbox}
\end{tcolorbox}
\newpage
\begin{tcolorbox}[enhanced,width=16cm,height=14cm,colback=white,colframe=cgeo,boxrule=2pt,arc=0mm,
left=6mm,right=6mm,top=0mm,bottom=4mm,sharp corners,
overlay={\fill[cgeo] (frame.north west) rectangle ([yshift=-1.25cm]frame.north east);
\node[anchor=west,text=white,font=\sffamily\bfseries\LARGE] at ([xshift=5mm,yshift=-0.62cm]frame.north west) {Question 55 B};
\node[anchor=east,text=white,font=\sffamily\itshape\large] at ([xshift=-5mm,yshift=-0.62cm]frame.north east) {Géométrie};}]
\vspace*{1.75cm}
\begin{minipage}[t][6.9cm][c]{\linewidth}\Large\centering
\begin{minipage}{0.95\linewidth}\raggedright\linespread{1.25}\selectfont Dans le triangle $ABC$, $M$ est sur $[AB]$, $N$ est sur $[AC]$ et les droites $(MN)$ et $(BC)$ sont parallèles. $AM=4$ cm, $AB=10$ cm et $AN=6$ cm. Combien mesure $AC$~?\end{minipage}
\end{minipage}

\vfill
\begin{tcolorbox}[enhanced,nobeforeafter,width=7.1cm,height=2.1cm,valign=center,colback=cgeo!6,colframe=cgeo,arc=2.5mm,boxrule=0.9pt,left=3mm]{\Large\bfseries\color{cgeo}a)}\quad{\Large $12~\text{cm}$}\end{tcolorbox}\hfill\begin{tcolorbox}[enhanced,nobeforeafter,width=7.1cm,height=2.1cm,valign=center,colback=cgeo!6,colframe=cgeo,arc=2.5mm,boxrule=0.9pt,left=3mm]{\Large\bfseries\color{cgeo}b)}\quad{\Large $15~\text{cm}$}\end{tcolorbox}\par\vspace{2.5mm}
\begin{tcolorbox}[enhanced,nobeforeafter,width=7.1cm,height=2.1cm,valign=center,colback=cgeo!6,colframe=cgeo,arc=2.5mm,boxrule=0.9pt,left=3mm]{\Large\bfseries\color{cgeo}c)}\quad{\Large $\dfrac{20}{3}~\text{cm}$}\end{tcolorbox}\hfill\begin{tcolorbox}[enhanced,nobeforeafter,width=7.1cm,height=2.1cm,valign=center,colback=cgeo!6,colframe=cgeo,arc=2.5mm,boxrule=0.9pt,left=3mm]{\Large\bfseries\color{cgeo}d)}\quad{\Large $2{,}4~\text{cm}$}\end{tcolorbox}
\end{tcolorbox}
\newpage
\begin{tcolorbox}[enhanced,width=16cm,height=14cm,colback=white,colframe=cgeo,boxrule=2pt,arc=0mm,
left=6mm,right=6mm,top=0mm,bottom=4mm,sharp corners,
overlay={\fill[cgeo] (frame.north west) rectangle ([yshift=-1.25cm]frame.north east);
\node[anchor=west,text=white,font=\sffamily\bfseries\LARGE] at ([xshift=5mm,yshift=-0.62cm]frame.north west) {Question 56 B};
\node[anchor=east,text=white,font=\sffamily\itshape\large] at ([xshift=-5mm,yshift=-0.62cm]frame.north east) {Géométrie};}]
\vspace*{1.75cm}
\begin{minipage}[t][6.9cm][c]{\linewidth}\Large\centering
\begin{minipage}{0.95\linewidth}\raggedright\linespread{1.25}\selectfont Dans le triangle $ABC$, $M$ est sur $[AB]$, $N$ est sur $[AC]$ et les droites $(MN)$ et $(BC)$ sont parallèles. $AM=5$ cm, $AB=8$ cm et $AN=10$ cm. Combien mesure $AC$~?\end{minipage}
\end{minipage}

\vfill
\begin{tcolorbox}[enhanced,nobeforeafter,width=7.1cm,height=2.1cm,valign=center,colback=cgeo!6,colframe=cgeo,arc=2.5mm,boxrule=0.9pt,left=3mm]{\Large\bfseries\color{cgeo}a)}\quad{\Large $4~\text{cm}$}\end{tcolorbox}\hfill\begin{tcolorbox}[enhanced,nobeforeafter,width=7.1cm,height=2.1cm,valign=center,colback=cgeo!6,colframe=cgeo,arc=2.5mm,boxrule=0.9pt,left=3mm]{\Large\bfseries\color{cgeo}b)}\quad{\Large $13~\text{cm}$}\end{tcolorbox}\par\vspace{2.5mm}
\begin{tcolorbox}[enhanced,nobeforeafter,width=7.1cm,height=2.1cm,valign=center,colback=cgeo!6,colframe=cgeo,arc=2.5mm,boxrule=0.9pt,left=3mm]{\Large\bfseries\color{cgeo}c)}\quad{\Large $16~\text{cm}$}\end{tcolorbox}\hfill\begin{tcolorbox}[enhanced,nobeforeafter,width=7.1cm,height=2.1cm,valign=center,colback=cgeo!6,colframe=cgeo,arc=2.5mm,boxrule=0.9pt,left=3mm]{\Large\bfseries\color{cgeo}d)}\quad{\Large $6{,}25~\text{cm}$}\end{tcolorbox}
\end{tcolorbox}
\newpage
\begin{tcolorbox}[enhanced,width=16cm,height=14cm,colback=white,colframe=cgeo,boxrule=2pt,arc=0mm,
left=6mm,right=6mm,top=0mm,bottom=4mm,sharp corners,
overlay={\fill[cgeo] (frame.north west) rectangle ([yshift=-1.25cm]frame.north east);
\node[anchor=west,text=white,font=\sffamily\bfseries\LARGE] at ([xshift=5mm,yshift=-0.62cm]frame.north west) {Question 57 B};
\node[anchor=east,text=white,font=\sffamily\itshape\large] at ([xshift=-5mm,yshift=-0.62cm]frame.north east) {Géométrie};}]
\vspace*{1.75cm}
\begin{minipage}[t][6.9cm][c]{\linewidth}\LARGE\centering
\begin{minipage}{0.95\linewidth}\raggedright\linespread{1.25}\selectfont Le triangle $ABC$ est rectangle en $A$, avec $BC=13$ cm et $AB=5$ cm. Que vaut $\cos\widehat{ABC}$~?\end{minipage}
\end{minipage}

\vfill
\begin{tcolorbox}[enhanced,nobeforeafter,width=7.1cm,height=2.1cm,valign=center,colback=cgeo!6,colframe=cgeo,arc=2.5mm,boxrule=0.9pt,left=3mm]{\Large\bfseries\color{cgeo}a)}\quad{\Large $\dfrac{5}{8}$}\end{tcolorbox}\hfill\begin{tcolorbox}[enhanced,nobeforeafter,width=7.1cm,height=2.1cm,valign=center,colback=cgeo!6,colframe=cgeo,arc=2.5mm,boxrule=0.9pt,left=3mm]{\Large\bfseries\color{cgeo}b)}\quad{\Large $\dfrac{8}{13}$}\end{tcolorbox}\par\vspace{2.5mm}
\begin{tcolorbox}[enhanced,nobeforeafter,width=7.1cm,height=2.1cm,valign=center,colback=cgeo!6,colframe=cgeo,arc=2.5mm,boxrule=0.9pt,left=3mm]{\Large\bfseries\color{cgeo}c)}\quad{\Large $\dfrac{13}{5}$}\end{tcolorbox}\hfill\begin{tcolorbox}[enhanced,nobeforeafter,width=7.1cm,height=2.1cm,valign=center,colback=cgeo!6,colframe=cgeo,arc=2.5mm,boxrule=0.9pt,left=3mm]{\Large\bfseries\color{cgeo}d)}\quad{\Large $\dfrac{5}{13}$}\end{tcolorbox}
\end{tcolorbox}
\newpage
\begin{tcolorbox}[enhanced,width=16cm,height=14cm,colback=white,colframe=cgeo,boxrule=2pt,arc=0mm,
left=6mm,right=6mm,top=0mm,bottom=4mm,sharp corners,
overlay={\fill[cgeo] (frame.north west) rectangle ([yshift=-1.25cm]frame.north east);
\node[anchor=west,text=white,font=\sffamily\bfseries\LARGE] at ([xshift=5mm,yshift=-0.62cm]frame.north west) {Question 58 B};
\node[anchor=east,text=white,font=\sffamily\itshape\large] at ([xshift=-5mm,yshift=-0.62cm]frame.north east) {Géométrie};}]
\vspace*{1.75cm}
\begin{minipage}[t][6.9cm][c]{\linewidth}\LARGE\centering
\begin{minipage}{0.95\linewidth}\raggedright\linespread{1.25}\selectfont Le triangle $ABC$ est rectangle en $A$, avec $BC=17$ cm et $AB=15$ cm. Que vaut $\cos\widehat{ABC}$~?\end{minipage}
\end{minipage}

\vfill
\begin{tcolorbox}[enhanced,nobeforeafter,width=7.1cm,height=2.1cm,valign=center,colback=cgeo!6,colframe=cgeo,arc=2.5mm,boxrule=0.9pt,left=3mm]{\Large\bfseries\color{cgeo}a)}\quad{\Large $\dfrac{15}{2}$}\end{tcolorbox}\hfill\begin{tcolorbox}[enhanced,nobeforeafter,width=7.1cm,height=2.1cm,valign=center,colback=cgeo!6,colframe=cgeo,arc=2.5mm,boxrule=0.9pt,left=3mm]{\Large\bfseries\color{cgeo}b)}\quad{\Large $\dfrac{15}{17}$}\end{tcolorbox}\par\vspace{2.5mm}
\begin{tcolorbox}[enhanced,nobeforeafter,width=7.1cm,height=2.1cm,valign=center,colback=cgeo!6,colframe=cgeo,arc=2.5mm,boxrule=0.9pt,left=3mm]{\Large\bfseries\color{cgeo}c)}\quad{\Large $\dfrac{17}{15}$}\end{tcolorbox}\hfill\begin{tcolorbox}[enhanced,nobeforeafter,width=7.1cm,height=2.1cm,valign=center,colback=cgeo!6,colframe=cgeo,arc=2.5mm,boxrule=0.9pt,left=3mm]{\Large\bfseries\color{cgeo}d)}\quad{\Large $\dfrac{2}{17}$}\end{tcolorbox}
\end{tcolorbox}
\newpage
\begin{tcolorbox}[enhanced,width=16cm,height=14cm,colback=white,colframe=cgeo,boxrule=2pt,arc=0mm,
left=6mm,right=6mm,top=0mm,bottom=4mm,sharp corners,
overlay={\fill[cgeo] (frame.north west) rectangle ([yshift=-1.25cm]frame.north east);
\node[anchor=west,text=white,font=\sffamily\bfseries\LARGE] at ([xshift=5mm,yshift=-0.62cm]frame.north west) {Question 59 B};
\node[anchor=east,text=white,font=\sffamily\itshape\large] at ([xshift=-5mm,yshift=-0.62cm]frame.north east) {Géométrie};}]
\vspace*{1.75cm}
\begin{minipage}[t][6.9cm][c]{\linewidth}\LARGE\centering
\begin{minipage}{0.95\linewidth}\raggedright\linespread{1.25}\selectfont Le triangle $ABC$ est rectangle en $A$, avec $BC=20$ cm et $\cos\widehat{ABC}=0{,}35$. Combien mesure $AB$~?\end{minipage}
\end{minipage}

\vfill
\begin{tcolorbox}[enhanced,nobeforeafter,width=7.1cm,height=2.1cm,valign=center,colback=cgeo!6,colframe=cgeo,arc=2.5mm,boxrule=0.9pt,left=3mm]{\Large\bfseries\color{cgeo}a)}\quad{\Large $8~\text{cm}$}\end{tcolorbox}\hfill\begin{tcolorbox}[enhanced,nobeforeafter,width=7.1cm,height=2.1cm,valign=center,colback=cgeo!6,colframe=cgeo,arc=2.5mm,boxrule=0.9pt,left=3mm]{\Large\bfseries\color{cgeo}b)}\quad{\Large $7~\text{cm}$}\end{tcolorbox}\par\vspace{2.5mm}
\begin{tcolorbox}[enhanced,nobeforeafter,width=7.1cm,height=2.1cm,valign=center,colback=cgeo!6,colframe=cgeo,arc=2.5mm,boxrule=0.9pt,left=3mm]{\Large\bfseries\color{cgeo}c)}\quad{\Large $13~\text{cm}$}\end{tcolorbox}\hfill\begin{tcolorbox}[enhanced,nobeforeafter,width=7.1cm,height=2.1cm,valign=center,colback=cgeo!6,colframe=cgeo,arc=2.5mm,boxrule=0.9pt,left=3mm]{\Large\bfseries\color{cgeo}d)}\quad{\Large $3{,}5~\text{cm}$}\end{tcolorbox}
\end{tcolorbox}
\newpage
\begin{tcolorbox}[enhanced,width=16cm,height=14cm,colback=white,colframe=cgeo,boxrule=2pt,arc=0mm,
left=6mm,right=6mm,top=0mm,bottom=4mm,sharp corners,
overlay={\fill[cgeo] (frame.north west) rectangle ([yshift=-1.25cm]frame.north east);
\node[anchor=west,text=white,font=\sffamily\bfseries\LARGE] at ([xshift=5mm,yshift=-0.62cm]frame.north west) {Question 60 B};
\node[anchor=east,text=white,font=\sffamily\itshape\large] at ([xshift=-5mm,yshift=-0.62cm]frame.north east) {Géométrie};}]
\vspace*{1.75cm}
\begin{minipage}[t][6.9cm][c]{\linewidth}\LARGE\centering
\begin{minipage}{0.95\linewidth}\raggedright\linespread{1.25}\selectfont Le triangle $ABC$ est rectangle en $A$, avec $BC=12$ cm et $\cos\widehat{ABC}=0{,}25$. Combien mesure $AB$~?\end{minipage}
\end{minipage}

\vfill
\begin{tcolorbox}[enhanced,nobeforeafter,width=7.1cm,height=2.1cm,valign=center,colback=cgeo!6,colframe=cgeo,arc=2.5mm,boxrule=0.9pt,left=3mm]{\Large\bfseries\color{cgeo}a)}\quad{\Large $1{,}5~\text{cm}$}\end{tcolorbox}\hfill\begin{tcolorbox}[enhanced,nobeforeafter,width=7.1cm,height=2.1cm,valign=center,colback=cgeo!6,colframe=cgeo,arc=2.5mm,boxrule=0.9pt,left=3mm]{\Large\bfseries\color{cgeo}b)}\quad{\Large $3~\text{cm}$}\end{tcolorbox}\par\vspace{2.5mm}
\begin{tcolorbox}[enhanced,nobeforeafter,width=7.1cm,height=2.1cm,valign=center,colback=cgeo!6,colframe=cgeo,arc=2.5mm,boxrule=0.9pt,left=3mm]{\Large\bfseries\color{cgeo}c)}\quad{\Large $4~\text{cm}$}\end{tcolorbox}\hfill\begin{tcolorbox}[enhanced,nobeforeafter,width=7.1cm,height=2.1cm,valign=center,colback=cgeo!6,colframe=cgeo,arc=2.5mm,boxrule=0.9pt,left=3mm]{\Large\bfseries\color{cgeo}d)}\quad{\Large $9~\text{cm}$}\end{tcolorbox}
\end{tcolorbox}
\newpage
\begin{tcolorbox}[enhanced,width=16cm,height=14cm,colback=white,colframe=cgeo,boxrule=2pt,arc=0mm,
left=6mm,right=6mm,top=0mm,bottom=4mm,sharp corners,
overlay={\fill[cgeo] (frame.north west) rectangle ([yshift=-1.25cm]frame.north east);
\node[anchor=west,text=white,font=\sffamily\bfseries\LARGE] at ([xshift=5mm,yshift=-0.62cm]frame.north west) {Question 61 B};
\node[anchor=east,text=white,font=\sffamily\itshape\large] at ([xshift=-5mm,yshift=-0.62cm]frame.north east) {Géométrie};}]
\vspace*{1.75cm}
\begin{minipage}[t][6.9cm][c]{\linewidth}\Large\centering
\begin{minipage}{0.95\linewidth}\raggedright\linespread{1.25}\selectfont Dans un repère, $A(0\,;\,0)$, $B(2\,;\,-3)$ et $C(-1\,;\,4)$. Quelles sont les coordonnées de l'image de $C$ par la translation qui transforme $A$ en $B$~?\end{minipage}
\end{minipage}

\vfill
\begin{tcolorbox}[enhanced,nobeforeafter,width=7.1cm,height=2.1cm,valign=center,colback=cgeo!6,colframe=cgeo,arc=2.5mm,boxrule=0.9pt,left=3mm]{\Large\bfseries\color{cgeo}a)}\quad{\Large $(1\,;\,1)$}\end{tcolorbox}\hfill\begin{tcolorbox}[enhanced,nobeforeafter,width=7.1cm,height=2.1cm,valign=center,colback=cgeo!6,colframe=cgeo,arc=2.5mm,boxrule=0.9pt,left=3mm]{\Large\bfseries\color{cgeo}b)}\quad{\Large $(1\,;\,7)$}\end{tcolorbox}\par\vspace{2.5mm}
\begin{tcolorbox}[enhanced,nobeforeafter,width=7.1cm,height=2.1cm,valign=center,colback=cgeo!6,colframe=cgeo,arc=2.5mm,boxrule=0.9pt,left=3mm]{\Large\bfseries\color{cgeo}c)}\quad{\Large $(-4\,;\,6)$}\end{tcolorbox}\hfill\begin{tcolorbox}[enhanced,nobeforeafter,width=7.1cm,height=2.1cm,valign=center,colback=cgeo!6,colframe=cgeo,arc=2.5mm,boxrule=0.9pt,left=3mm]{\Large\bfseries\color{cgeo}d)}\quad{\Large $(-3\,;\,7)$}\end{tcolorbox}
\end{tcolorbox}
\newpage
\begin{tcolorbox}[enhanced,width=16cm,height=14cm,colback=white,colframe=cgeo,boxrule=2pt,arc=0mm,
left=6mm,right=6mm,top=0mm,bottom=4mm,sharp corners,
overlay={\fill[cgeo] (frame.north west) rectangle ([yshift=-1.25cm]frame.north east);
\node[anchor=west,text=white,font=\sffamily\bfseries\LARGE] at ([xshift=5mm,yshift=-0.62cm]frame.north west) {Question 62 B};
\node[anchor=east,text=white,font=\sffamily\itshape\large] at ([xshift=-5mm,yshift=-0.62cm]frame.north east) {Géométrie};}]
\vspace*{1.75cm}
\begin{minipage}[t][6.9cm][c]{\linewidth}\Large\centering
\begin{minipage}{0.95\linewidth}\raggedright\linespread{1.25}\selectfont Dans un repère, $A(3\,;\,-1)$, $B(1\,;\,2)$ et $C(5\,;\,0)$. Quelles sont les coordonnées de l'image de $C$ par la translation qui transforme $A$ en $B$~?\end{minipage}
\end{minipage}

\vfill
\begin{tcolorbox}[enhanced,nobeforeafter,width=7.1cm,height=2.1cm,valign=center,colback=cgeo!6,colframe=cgeo,arc=2.5mm,boxrule=0.9pt,left=3mm]{\Large\bfseries\color{cgeo}a)}\quad{\Large $(7\,;\,-3)$}\end{tcolorbox}\hfill\begin{tcolorbox}[enhanced,nobeforeafter,width=7.1cm,height=2.1cm,valign=center,colback=cgeo!6,colframe=cgeo,arc=2.5mm,boxrule=0.9pt,left=3mm]{\Large\bfseries\color{cgeo}b)}\quad{\Large $(3\,;\,3)$}\end{tcolorbox}\par\vspace{2.5mm}
\begin{tcolorbox}[enhanced,nobeforeafter,width=7.1cm,height=2.1cm,valign=center,colback=cgeo!6,colframe=cgeo,arc=2.5mm,boxrule=0.9pt,left=3mm]{\Large\bfseries\color{cgeo}c)}\quad{\Large $(6\,;\,2)$}\end{tcolorbox}\hfill\begin{tcolorbox}[enhanced,nobeforeafter,width=7.1cm,height=2.1cm,valign=center,colback=cgeo!6,colframe=cgeo,arc=2.5mm,boxrule=0.9pt,left=3mm]{\Large\bfseries\color{cgeo}d)}\quad{\Large $(8\,;\,-2)$}\end{tcolorbox}
\end{tcolorbox}
\newpage
\begin{tcolorbox}[enhanced,width=16cm,height=14cm,colback=white,colframe=cgeo,boxrule=2pt,arc=0mm,
left=6mm,right=6mm,top=0mm,bottom=4mm,sharp corners,
overlay={\fill[cgeo] (frame.north west) rectangle ([yshift=-1.25cm]frame.north east);
\node[anchor=west,text=white,font=\sffamily\bfseries\LARGE] at ([xshift=5mm,yshift=-0.62cm]frame.north west) {Question 63 B};
\node[anchor=east,text=white,font=\sffamily\itshape\large] at ([xshift=-5mm,yshift=-0.62cm]frame.north east) {Géométrie};}]
\vspace*{1.75cm}
\begin{minipage}[t][6.9cm][c]{\linewidth}\LARGE\centering
\begin{minipage}{0.95\linewidth}\raggedright\linespread{1.25}\selectfont Une figure est agrandie à l'échelle $5$ (les longueurs sont multipliées par $5$). Par combien son aire est-elle multipliée~?\end{minipage}
\end{minipage}

\vfill
\begin{tcolorbox}[enhanced,nobeforeafter,width=7.1cm,height=2.1cm,valign=center,colback=cgeo!6,colframe=cgeo,arc=2.5mm,boxrule=0.9pt,left=3mm]{\Large\bfseries\color{cgeo}a)}\quad{\Large $5$}\end{tcolorbox}\hfill\begin{tcolorbox}[enhanced,nobeforeafter,width=7.1cm,height=2.1cm,valign=center,colback=cgeo!6,colframe=cgeo,arc=2.5mm,boxrule=0.9pt,left=3mm]{\Large\bfseries\color{cgeo}b)}\quad{\Large $125$}\end{tcolorbox}\par\vspace{2.5mm}
\begin{tcolorbox}[enhanced,nobeforeafter,width=7.1cm,height=2.1cm,valign=center,colback=cgeo!6,colframe=cgeo,arc=2.5mm,boxrule=0.9pt,left=3mm]{\Large\bfseries\color{cgeo}c)}\quad{\Large $25$}\end{tcolorbox}\hfill\begin{tcolorbox}[enhanced,nobeforeafter,width=7.1cm,height=2.1cm,valign=center,colback=cgeo!6,colframe=cgeo,arc=2.5mm,boxrule=0.9pt,left=3mm]{\Large\bfseries\color{cgeo}d)}\quad{\Large $10$}\end{tcolorbox}
\end{tcolorbox}
\newpage
\begin{tcolorbox}[enhanced,width=16cm,height=14cm,colback=white,colframe=cgeo,boxrule=2pt,arc=0mm,
left=6mm,right=6mm,top=0mm,bottom=4mm,sharp corners,
overlay={\fill[cgeo] (frame.north west) rectangle ([yshift=-1.25cm]frame.north east);
\node[anchor=west,text=white,font=\sffamily\bfseries\LARGE] at ([xshift=5mm,yshift=-0.62cm]frame.north west) {Question 64 B};
\node[anchor=east,text=white,font=\sffamily\itshape\large] at ([xshift=-5mm,yshift=-0.62cm]frame.north east) {Géométrie};}]
\vspace*{1.75cm}
\begin{minipage}[t][6.9cm][c]{\linewidth}\LARGE\centering
\begin{minipage}{0.95\linewidth}\raggedright\linespread{1.25}\selectfont Une figure est agrandie à l'échelle $4$ (les longueurs sont multipliées par $4$). Par combien son aire est-elle multipliée~?\end{minipage}
\end{minipage}

\vfill
\begin{tcolorbox}[enhanced,nobeforeafter,width=7.1cm,height=2.1cm,valign=center,colback=cgeo!6,colframe=cgeo,arc=2.5mm,boxrule=0.9pt,left=3mm]{\Large\bfseries\color{cgeo}a)}\quad{\Large $16$}\end{tcolorbox}\hfill\begin{tcolorbox}[enhanced,nobeforeafter,width=7.1cm,height=2.1cm,valign=center,colback=cgeo!6,colframe=cgeo,arc=2.5mm,boxrule=0.9pt,left=3mm]{\Large\bfseries\color{cgeo}b)}\quad{\Large $4$}\end{tcolorbox}\par\vspace{2.5mm}
\begin{tcolorbox}[enhanced,nobeforeafter,width=7.1cm,height=2.1cm,valign=center,colback=cgeo!6,colframe=cgeo,arc=2.5mm,boxrule=0.9pt,left=3mm]{\Large\bfseries\color{cgeo}c)}\quad{\Large $8$}\end{tcolorbox}\hfill\begin{tcolorbox}[enhanced,nobeforeafter,width=7.1cm,height=2.1cm,valign=center,colback=cgeo!6,colframe=cgeo,arc=2.5mm,boxrule=0.9pt,left=3mm]{\Large\bfseries\color{cgeo}d)}\quad{\Large $64$}\end{tcolorbox}
\end{tcolorbox}
\newpage
\begin{tcolorbox}[enhanced,width=16cm,height=14cm,colback=white,colframe=cdon,boxrule=2pt,arc=0mm,
left=6mm,right=6mm,top=0mm,bottom=4mm,sharp corners,
overlay={\fill[cdon] (frame.north west) rectangle ([yshift=-1.25cm]frame.north east);
\node[anchor=west,text=white,font=\sffamily\bfseries\LARGE] at ([xshift=5mm,yshift=-0.62cm]frame.north west) {Question 65 B};
\node[anchor=east,text=white,font=\sffamily\itshape\large] at ([xshift=-5mm,yshift=-0.62cm]frame.north east) {Proportionnalité, statistiques, probabilités};}]
\vspace*{1.75cm}
\begin{minipage}[t][6.9cm][c]{\linewidth}\LARGE\centering
\begin{minipage}{0.95\linewidth}\raggedright\linespread{1.25}\selectfont Ce tableau est un tableau de proportionnalité~: $\begin{array}{|c|c|}\hline 5 & 15\\\hline 8 &~?\\\hline\end{array}$. Quel nombre remplace «~?~»~?\end{minipage}
\end{minipage}

\vfill
\begin{tcolorbox}[enhanced,nobeforeafter,width=7.1cm,height=2.1cm,valign=center,colback=cdon!6,colframe=cdon,arc=2.5mm,boxrule=0.9pt,left=3mm]{\Large\bfseries\color{cdon}a)}\quad{\Large $18$}\end{tcolorbox}\hfill\begin{tcolorbox}[enhanced,nobeforeafter,width=7.1cm,height=2.1cm,valign=center,colback=cdon!6,colframe=cdon,arc=2.5mm,boxrule=0.9pt,left=3mm]{\Large\bfseries\color{cdon}b)}\quad{\Large $\dfrac{75}{8}$}\end{tcolorbox}\par\vspace{2.5mm}
\begin{tcolorbox}[enhanced,nobeforeafter,width=7.1cm,height=2.1cm,valign=center,colback=cdon!6,colframe=cdon,arc=2.5mm,boxrule=0.9pt,left=3mm]{\Large\bfseries\color{cdon}c)}\quad{\Large $24$}\end{tcolorbox}\hfill\begin{tcolorbox}[enhanced,nobeforeafter,width=7.1cm,height=2.1cm,valign=center,colback=cdon!6,colframe=cdon,arc=2.5mm,boxrule=0.9pt,left=3mm]{\Large\bfseries\color{cdon}d)}\quad{\Large $120$}\end{tcolorbox}
\end{tcolorbox}
\newpage
\begin{tcolorbox}[enhanced,width=16cm,height=14cm,colback=white,colframe=cdon,boxrule=2pt,arc=0mm,
left=6mm,right=6mm,top=0mm,bottom=4mm,sharp corners,
overlay={\fill[cdon] (frame.north west) rectangle ([yshift=-1.25cm]frame.north east);
\node[anchor=west,text=white,font=\sffamily\bfseries\LARGE] at ([xshift=5mm,yshift=-0.62cm]frame.north west) {Question 66 B};
\node[anchor=east,text=white,font=\sffamily\itshape\large] at ([xshift=-5mm,yshift=-0.62cm]frame.north east) {Proportionnalité, statistiques, probabilités};}]
\vspace*{1.75cm}
\begin{minipage}[t][6.9cm][c]{\linewidth}\LARGE\centering
\begin{minipage}{0.95\linewidth}\raggedright\linespread{1.25}\selectfont Ce tableau est un tableau de proportionnalité~: $\begin{array}{|c|c|}\hline 6 & 14\\\hline 9 &~?\\\hline\end{array}$. Quel nombre remplace «~?~»~?\end{minipage}
\end{minipage}

\vfill
\begin{tcolorbox}[enhanced,nobeforeafter,width=7.1cm,height=2.1cm,valign=center,colback=cdon!6,colframe=cdon,arc=2.5mm,boxrule=0.9pt,left=3mm]{\Large\bfseries\color{cdon}a)}\quad{\Large $21$}\end{tcolorbox}\hfill\begin{tcolorbox}[enhanced,nobeforeafter,width=7.1cm,height=2.1cm,valign=center,colback=cdon!6,colframe=cdon,arc=2.5mm,boxrule=0.9pt,left=3mm]{\Large\bfseries\color{cdon}b)}\quad{\Large $17$}\end{tcolorbox}\par\vspace{2.5mm}
\begin{tcolorbox}[enhanced,nobeforeafter,width=7.1cm,height=2.1cm,valign=center,colback=cdon!6,colframe=cdon,arc=2.5mm,boxrule=0.9pt,left=3mm]{\Large\bfseries\color{cdon}c)}\quad{\Large $126$}\end{tcolorbox}\hfill\begin{tcolorbox}[enhanced,nobeforeafter,width=7.1cm,height=2.1cm,valign=center,colback=cdon!6,colframe=cdon,arc=2.5mm,boxrule=0.9pt,left=3mm]{\Large\bfseries\color{cdon}d)}\quad{\Large $\dfrac{28}{3}$}\end{tcolorbox}
\end{tcolorbox}
\newpage
\begin{tcolorbox}[enhanced,width=16cm,height=14cm,colback=white,colframe=cdon,boxrule=2pt,arc=0mm,
left=6mm,right=6mm,top=0mm,bottom=4mm,sharp corners,
overlay={\fill[cdon] (frame.north west) rectangle ([yshift=-1.25cm]frame.north east);
\node[anchor=west,text=white,font=\sffamily\bfseries\LARGE] at ([xshift=5mm,yshift=-0.62cm]frame.north west) {Question 67 B};
\node[anchor=east,text=white,font=\sffamily\itshape\large] at ([xshift=-5mm,yshift=-0.62cm]frame.north east) {Proportionnalité, statistiques, probabilités};}]
\vspace*{1.75cm}
\begin{minipage}[t][6.9cm][c]{\linewidth}\LARGE\centering
\begin{minipage}{0.95\linewidth}\raggedright\linespread{1.25}\selectfont Que vaut $12\,\%$ de $50$~?\end{minipage}
\end{minipage}

\vfill
\begin{tcolorbox}[enhanced,nobeforeafter,width=7.1cm,height=2.1cm,valign=center,colback=cdon!6,colframe=cdon,arc=2.5mm,boxrule=0.9pt,left=3mm]{\Large\bfseries\color{cdon}a)}\quad{\Large $44$}\end{tcolorbox}\hfill\begin{tcolorbox}[enhanced,nobeforeafter,width=7.1cm,height=2.1cm,valign=center,colback=cdon!6,colframe=cdon,arc=2.5mm,boxrule=0.9pt,left=3mm]{\Large\bfseries\color{cdon}b)}\quad{\Large $6$}\end{tcolorbox}\par\vspace{2.5mm}
\begin{tcolorbox}[enhanced,nobeforeafter,width=7.1cm,height=2.1cm,valign=center,colback=cdon!6,colframe=cdon,arc=2.5mm,boxrule=0.9pt,left=3mm]{\Large\bfseries\color{cdon}c)}\quad{\Large $60$}\end{tcolorbox}\hfill\begin{tcolorbox}[enhanced,nobeforeafter,width=7.1cm,height=2.1cm,valign=center,colback=cdon!6,colframe=cdon,arc=2.5mm,boxrule=0.9pt,left=3mm]{\Large\bfseries\color{cdon}d)}\quad{\Large $18$}\end{tcolorbox}
\end{tcolorbox}
\newpage
\begin{tcolorbox}[enhanced,width=16cm,height=14cm,colback=white,colframe=cdon,boxrule=2pt,arc=0mm,
left=6mm,right=6mm,top=0mm,bottom=4mm,sharp corners,
overlay={\fill[cdon] (frame.north west) rectangle ([yshift=-1.25cm]frame.north east);
\node[anchor=west,text=white,font=\sffamily\bfseries\LARGE] at ([xshift=5mm,yshift=-0.62cm]frame.north west) {Question 68 B};
\node[anchor=east,text=white,font=\sffamily\itshape\large] at ([xshift=-5mm,yshift=-0.62cm]frame.north east) {Proportionnalité, statistiques, probabilités};}]
\vspace*{1.75cm}
\begin{minipage}[t][6.9cm][c]{\linewidth}\LARGE\centering
\begin{minipage}{0.95\linewidth}\raggedright\linespread{1.25}\selectfont Que vaut $45\,\%$ de $200$~?\end{minipage}
\end{minipage}

\vfill
\begin{tcolorbox}[enhanced,nobeforeafter,width=7.1cm,height=2.1cm,valign=center,colback=cdon!6,colframe=cdon,arc=2.5mm,boxrule=0.9pt,left=3mm]{\Large\bfseries\color{cdon}a)}\quad{\Large $90$}\end{tcolorbox}\hfill\begin{tcolorbox}[enhanced,nobeforeafter,width=7.1cm,height=2.1cm,valign=center,colback=cdon!6,colframe=cdon,arc=2.5mm,boxrule=0.9pt,left=3mm]{\Large\bfseries\color{cdon}b)}\quad{\Large $900$}\end{tcolorbox}\par\vspace{2.5mm}
\begin{tcolorbox}[enhanced,nobeforeafter,width=7.1cm,height=2.1cm,valign=center,colback=cdon!6,colframe=cdon,arc=2.5mm,boxrule=0.9pt,left=3mm]{\Large\bfseries\color{cdon}c)}\quad{\Large $135$}\end{tcolorbox}\hfill\begin{tcolorbox}[enhanced,nobeforeafter,width=7.1cm,height=2.1cm,valign=center,colback=cdon!6,colframe=cdon,arc=2.5mm,boxrule=0.9pt,left=3mm]{\Large\bfseries\color{cdon}d)}\quad{\Large $110$}\end{tcolorbox}
\end{tcolorbox}
\newpage
\begin{tcolorbox}[enhanced,width=16cm,height=14cm,colback=white,colframe=cdon,boxrule=2pt,arc=0mm,
left=6mm,right=6mm,top=0mm,bottom=4mm,sharp corners,
overlay={\fill[cdon] (frame.north west) rectangle ([yshift=-1.25cm]frame.north east);
\node[anchor=west,text=white,font=\sffamily\bfseries\LARGE] at ([xshift=5mm,yshift=-0.62cm]frame.north west) {Question 69 B};
\node[anchor=east,text=white,font=\sffamily\itshape\large] at ([xshift=-5mm,yshift=-0.62cm]frame.north east) {Proportionnalité, statistiques, probabilités};}]
\vspace*{1.75cm}
\begin{minipage}[t][6.9cm][c]{\linewidth}\LARGE\centering
\begin{minipage}{0.95\linewidth}\raggedright\linespread{1.25}\selectfont Dans une classe de $50$ élèves, $18$ sont demi-pensionnaires. Quel pourcentage des élèves cela représente-t-il~?\end{minipage}
\end{minipage}

\vfill
\begin{tcolorbox}[enhanced,nobeforeafter,width=7.1cm,height=2.1cm,valign=center,colback=cdon!6,colframe=cdon,arc=2.5mm,boxrule=0.9pt,left=3mm]{\Large\bfseries\color{cdon}a)}\quad{\Large $3{,}6\,\%$}\end{tcolorbox}\hfill\begin{tcolorbox}[enhanced,nobeforeafter,width=7.1cm,height=2.1cm,valign=center,colback=cdon!6,colframe=cdon,arc=2.5mm,boxrule=0.9pt,left=3mm]{\Large\bfseries\color{cdon}b)}\quad{\Large $18\,\%$}\end{tcolorbox}\par\vspace{2.5mm}
\begin{tcolorbox}[enhanced,nobeforeafter,width=7.1cm,height=2.1cm,valign=center,colback=cdon!6,colframe=cdon,arc=2.5mm,boxrule=0.9pt,left=3mm]{\Large\bfseries\color{cdon}c)}\quad{\Large $64\,\%$}\end{tcolorbox}\hfill\begin{tcolorbox}[enhanced,nobeforeafter,width=7.1cm,height=2.1cm,valign=center,colback=cdon!6,colframe=cdon,arc=2.5mm,boxrule=0.9pt,left=3mm]{\Large\bfseries\color{cdon}d)}\quad{\Large $36\,\%$}\end{tcolorbox}
\end{tcolorbox}
\newpage
\begin{tcolorbox}[enhanced,width=16cm,height=14cm,colback=white,colframe=cdon,boxrule=2pt,arc=0mm,
left=6mm,right=6mm,top=0mm,bottom=4mm,sharp corners,
overlay={\fill[cdon] (frame.north west) rectangle ([yshift=-1.25cm]frame.north east);
\node[anchor=west,text=white,font=\sffamily\bfseries\LARGE] at ([xshift=5mm,yshift=-0.62cm]frame.north west) {Question 70 B};
\node[anchor=east,text=white,font=\sffamily\itshape\large] at ([xshift=-5mm,yshift=-0.62cm]frame.north east) {Proportionnalité, statistiques, probabilités};}]
\vspace*{1.75cm}
\begin{minipage}[t][6.9cm][c]{\linewidth}\LARGE\centering
\begin{minipage}{0.95\linewidth}\raggedright\linespread{1.25}\selectfont Dans une classe de $40$ élèves, $30$ sont demi-pensionnaires. Quel pourcentage des élèves cela représente-t-il~?\end{minipage}
\end{minipage}

\vfill
\begin{tcolorbox}[enhanced,nobeforeafter,width=7.1cm,height=2.1cm,valign=center,colback=cdon!6,colframe=cdon,arc=2.5mm,boxrule=0.9pt,left=3mm]{\Large\bfseries\color{cdon}a)}\quad{\Large $75\,\%$}\end{tcolorbox}\hfill\begin{tcolorbox}[enhanced,nobeforeafter,width=7.1cm,height=2.1cm,valign=center,colback=cdon!6,colframe=cdon,arc=2.5mm,boxrule=0.9pt,left=3mm]{\Large\bfseries\color{cdon}b)}\quad{\Large $7{,}5\,\%$}\end{tcolorbox}\par\vspace{2.5mm}
\begin{tcolorbox}[enhanced,nobeforeafter,width=7.1cm,height=2.1cm,valign=center,colback=cdon!6,colframe=cdon,arc=2.5mm,boxrule=0.9pt,left=3mm]{\Large\bfseries\color{cdon}c)}\quad{\Large $25\,\%$}\end{tcolorbox}\hfill\begin{tcolorbox}[enhanced,nobeforeafter,width=7.1cm,height=2.1cm,valign=center,colback=cdon!6,colframe=cdon,arc=2.5mm,boxrule=0.9pt,left=3mm]{\Large\bfseries\color{cdon}d)}\quad{\Large $30\,\%$}\end{tcolorbox}
\end{tcolorbox}
\newpage
\begin{tcolorbox}[enhanced,width=16cm,height=14cm,colback=white,colframe=cdon,boxrule=2pt,arc=0mm,
left=6mm,right=6mm,top=0mm,bottom=4mm,sharp corners,
overlay={\fill[cdon] (frame.north west) rectangle ([yshift=-1.25cm]frame.north east);
\node[anchor=west,text=white,font=\sffamily\bfseries\LARGE] at ([xshift=5mm,yshift=-0.62cm]frame.north west) {Question 71 B};
\node[anchor=east,text=white,font=\sffamily\itshape\large] at ([xshift=-5mm,yshift=-0.62cm]frame.north east) {Proportionnalité, statistiques, probabilités};}]
\vspace*{1.75cm}
\begin{minipage}[t][6.9cm][c]{\linewidth}\LARGE\centering
\begin{minipage}{0.95\linewidth}\raggedright\linespread{1.25}\selectfont Quelle est la moyenne de la série $5\ ;\ 9\ ;\ 14\ ;\ 16$~?\end{minipage}
\end{minipage}

\vfill
\begin{tcolorbox}[enhanced,nobeforeafter,width=7.1cm,height=2.1cm,valign=center,colback=cdon!6,colframe=cdon,arc=2.5mm,boxrule=0.9pt,left=3mm]{\Large\bfseries\color{cdon}a)}\quad{\Large $44$}\end{tcolorbox}\hfill\begin{tcolorbox}[enhanced,nobeforeafter,width=7.1cm,height=2.1cm,valign=center,colback=cdon!6,colframe=cdon,arc=2.5mm,boxrule=0.9pt,left=3mm]{\Large\bfseries\color{cdon}b)}\quad{\Large $12$}\end{tcolorbox}\par\vspace{2.5mm}
\begin{tcolorbox}[enhanced,nobeforeafter,width=7.1cm,height=2.1cm,valign=center,colback=cdon!6,colframe=cdon,arc=2.5mm,boxrule=0.9pt,left=3mm]{\Large\bfseries\color{cdon}c)}\quad{\Large $10{,}5$}\end{tcolorbox}\hfill\begin{tcolorbox}[enhanced,nobeforeafter,width=7.1cm,height=2.1cm,valign=center,colback=cdon!6,colframe=cdon,arc=2.5mm,boxrule=0.9pt,left=3mm]{\Large\bfseries\color{cdon}d)}\quad{\Large $11$}\end{tcolorbox}
\end{tcolorbox}
\newpage
\begin{tcolorbox}[enhanced,width=16cm,height=14cm,colback=white,colframe=cdon,boxrule=2pt,arc=0mm,
left=6mm,right=6mm,top=0mm,bottom=4mm,sharp corners,
overlay={\fill[cdon] (frame.north west) rectangle ([yshift=-1.25cm]frame.north east);
\node[anchor=west,text=white,font=\sffamily\bfseries\LARGE] at ([xshift=5mm,yshift=-0.62cm]frame.north west) {Question 72 B};
\node[anchor=east,text=white,font=\sffamily\itshape\large] at ([xshift=-5mm,yshift=-0.62cm]frame.north east) {Proportionnalité, statistiques, probabilités};}]
\vspace*{1.75cm}
\begin{minipage}[t][6.9cm][c]{\linewidth}\LARGE\centering
\begin{minipage}{0.95\linewidth}\raggedright\linespread{1.25}\selectfont Quelle est la moyenne de la série $13\ ;\ 7\ ;\ 10\ ;\ 18\ ;\ 17$~?\end{minipage}
\end{minipage}

\vfill
\begin{tcolorbox}[enhanced,nobeforeafter,width=7.1cm,height=2.1cm,valign=center,colback=cdon!6,colframe=cdon,arc=2.5mm,boxrule=0.9pt,left=3mm]{\Large\bfseries\color{cdon}a)}\quad{\Large $65$}\end{tcolorbox}\hfill\begin{tcolorbox}[enhanced,nobeforeafter,width=7.1cm,height=2.1cm,valign=center,colback=cdon!6,colframe=cdon,arc=2.5mm,boxrule=0.9pt,left=3mm]{\Large\bfseries\color{cdon}b)}\quad{\Large $14$}\end{tcolorbox}\par\vspace{2.5mm}
\begin{tcolorbox}[enhanced,nobeforeafter,width=7.1cm,height=2.1cm,valign=center,colback=cdon!6,colframe=cdon,arc=2.5mm,boxrule=0.9pt,left=3mm]{\Large\bfseries\color{cdon}c)}\quad{\Large $12{,}5$}\end{tcolorbox}\hfill\begin{tcolorbox}[enhanced,nobeforeafter,width=7.1cm,height=2.1cm,valign=center,colback=cdon!6,colframe=cdon,arc=2.5mm,boxrule=0.9pt,left=3mm]{\Large\bfseries\color{cdon}d)}\quad{\Large $13$}\end{tcolorbox}
\end{tcolorbox}
\newpage
\begin{tcolorbox}[enhanced,width=16cm,height=14cm,colback=white,colframe=cdon,boxrule=2pt,arc=0mm,
left=6mm,right=6mm,top=0mm,bottom=4mm,sharp corners,
overlay={\fill[cdon] (frame.north west) rectangle ([yshift=-1.25cm]frame.north east);
\node[anchor=west,text=white,font=\sffamily\bfseries\LARGE] at ([xshift=5mm,yshift=-0.62cm]frame.north west) {Question 73 B};
\node[anchor=east,text=white,font=\sffamily\itshape\large] at ([xshift=-5mm,yshift=-0.62cm]frame.north east) {Proportionnalité, statistiques, probabilités};}]
\vspace*{1.75cm}
\begin{minipage}[t][6.9cm][c]{\linewidth}\LARGE\centering
\begin{minipage}{0.95\linewidth}\raggedright\linespread{1.25}\selectfont Quelle est la médiane de la série $7\ ;\ 1\ ;\ 9\ ;\ 3\ ;\ 11$~?\end{minipage}
\end{minipage}

\vfill
\begin{tcolorbox}[enhanced,nobeforeafter,width=7.1cm,height=2.1cm,valign=center,colback=cdon!6,colframe=cdon,arc=2.5mm,boxrule=0.9pt,left=3mm]{\Large\bfseries\color{cdon}a)}\quad{\Large $9$}\end{tcolorbox}\hfill\begin{tcolorbox}[enhanced,nobeforeafter,width=7.1cm,height=2.1cm,valign=center,colback=cdon!6,colframe=cdon,arc=2.5mm,boxrule=0.9pt,left=3mm]{\Large\bfseries\color{cdon}b)}\quad{\Large $7$}\end{tcolorbox}\par\vspace{2.5mm}
\begin{tcolorbox}[enhanced,nobeforeafter,width=7.1cm,height=2.1cm,valign=center,colback=cdon!6,colframe=cdon,arc=2.5mm,boxrule=0.9pt,left=3mm]{\Large\bfseries\color{cdon}c)}\quad{\Large $6{,}2$}\end{tcolorbox}\hfill\begin{tcolorbox}[enhanced,nobeforeafter,width=7.1cm,height=2.1cm,valign=center,colback=cdon!6,colframe=cdon,arc=2.5mm,boxrule=0.9pt,left=3mm]{\Large\bfseries\color{cdon}d)}\quad{\Large $3$}\end{tcolorbox}
\end{tcolorbox}
\newpage
\begin{tcolorbox}[enhanced,width=16cm,height=14cm,colback=white,colframe=cdon,boxrule=2pt,arc=0mm,
left=6mm,right=6mm,top=0mm,bottom=4mm,sharp corners,
overlay={\fill[cdon] (frame.north west) rectangle ([yshift=-1.25cm]frame.north east);
\node[anchor=west,text=white,font=\sffamily\bfseries\LARGE] at ([xshift=5mm,yshift=-0.62cm]frame.north west) {Question 74 B};
\node[anchor=east,text=white,font=\sffamily\itshape\large] at ([xshift=-5mm,yshift=-0.62cm]frame.north east) {Proportionnalité, statistiques, probabilités};}]
\vspace*{1.75cm}
\begin{minipage}[t][6.9cm][c]{\linewidth}\LARGE\centering
\begin{minipage}{0.95\linewidth}\raggedright\linespread{1.25}\selectfont Quelle est la médiane de la série $14\ ;\ 8\ ;\ 10\ ;\ 6$~?\end{minipage}
\end{minipage}

\vfill
\begin{tcolorbox}[enhanced,nobeforeafter,width=7.1cm,height=2.1cm,valign=center,colback=cdon!6,colframe=cdon,arc=2.5mm,boxrule=0.9pt,left=3mm]{\Large\bfseries\color{cdon}a)}\quad{\Large $9$}\end{tcolorbox}\hfill\begin{tcolorbox}[enhanced,nobeforeafter,width=7.1cm,height=2.1cm,valign=center,colback=cdon!6,colframe=cdon,arc=2.5mm,boxrule=0.9pt,left=3mm]{\Large\bfseries\color{cdon}b)}\quad{\Large $10$}\end{tcolorbox}\par\vspace{2.5mm}
\begin{tcolorbox}[enhanced,nobeforeafter,width=7.1cm,height=2.1cm,valign=center,colback=cdon!6,colframe=cdon,arc=2.5mm,boxrule=0.9pt,left=3mm]{\Large\bfseries\color{cdon}c)}\quad{\Large $2{,}5$}\end{tcolorbox}\hfill\begin{tcolorbox}[enhanced,nobeforeafter,width=7.1cm,height=2.1cm,valign=center,colback=cdon!6,colframe=cdon,arc=2.5mm,boxrule=0.9pt,left=3mm]{\Large\bfseries\color{cdon}d)}\quad{\Large $9{,}5$}\end{tcolorbox}
\end{tcolorbox}
\newpage
\begin{tcolorbox}[enhanced,width=16cm,height=14cm,colback=white,colframe=cdon,boxrule=2pt,arc=0mm,
left=6mm,right=6mm,top=0mm,bottom=4mm,sharp corners,
overlay={\fill[cdon] (frame.north west) rectangle ([yshift=-1.25cm]frame.north east);
\node[anchor=west,text=white,font=\sffamily\bfseries\LARGE] at ([xshift=5mm,yshift=-0.62cm]frame.north west) {Question 75 B};
\node[anchor=east,text=white,font=\sffamily\itshape\large] at ([xshift=-5mm,yshift=-0.62cm]frame.north east) {Proportionnalité, statistiques, probabilités};}]
\vspace*{1.75cm}
\begin{minipage}[t][6.9cm][c]{\linewidth}\LARGE\centering
\begin{minipage}{0.95\linewidth}\raggedright\linespread{1.25}\selectfont Quelle est l'étendue de la série $20\ ;\ 14\ ;\ 25\ ;\ 8$~?\end{minipage}
\end{minipage}

\vfill
\begin{tcolorbox}[enhanced,nobeforeafter,width=7.1cm,height=2.1cm,valign=center,colback=cdon!6,colframe=cdon,arc=2.5mm,boxrule=0.9pt,left=3mm]{\Large\bfseries\color{cdon}a)}\quad{\Large $17$}\end{tcolorbox}\hfill\begin{tcolorbox}[enhanced,nobeforeafter,width=7.1cm,height=2.1cm,valign=center,colback=cdon!6,colframe=cdon,arc=2.5mm,boxrule=0.9pt,left=3mm]{\Large\bfseries\color{cdon}b)}\quad{\Large $25$}\end{tcolorbox}\par\vspace{2.5mm}
\begin{tcolorbox}[enhanced,nobeforeafter,width=7.1cm,height=2.1cm,valign=center,colback=cdon!6,colframe=cdon,arc=2.5mm,boxrule=0.9pt,left=3mm]{\Large\bfseries\color{cdon}c)}\quad{\Large $33$}\end{tcolorbox}\hfill\begin{tcolorbox}[enhanced,nobeforeafter,width=7.1cm,height=2.1cm,valign=center,colback=cdon!6,colframe=cdon,arc=2.5mm,boxrule=0.9pt,left=3mm]{\Large\bfseries\color{cdon}d)}\quad{\Large $16{,}75$}\end{tcolorbox}
\end{tcolorbox}
\newpage
\begin{tcolorbox}[enhanced,width=16cm,height=14cm,colback=white,colframe=cdon,boxrule=2pt,arc=0mm,
left=6mm,right=6mm,top=0mm,bottom=4mm,sharp corners,
overlay={\fill[cdon] (frame.north west) rectangle ([yshift=-1.25cm]frame.north east);
\node[anchor=west,text=white,font=\sffamily\bfseries\LARGE] at ([xshift=5mm,yshift=-0.62cm]frame.north west) {Question 76 B};
\node[anchor=east,text=white,font=\sffamily\itshape\large] at ([xshift=-5mm,yshift=-0.62cm]frame.north east) {Proportionnalité, statistiques, probabilités};}]
\vspace*{1.75cm}
\begin{minipage}[t][6.9cm][c]{\linewidth}\LARGE\centering
\begin{minipage}{0.95\linewidth}\raggedright\linespread{1.25}\selectfont Quelle est l'étendue de la série $4\ ;\ 10\ ;\ 2\ ;\ 9\ ;\ 13$~?\end{minipage}
\end{minipage}

\vfill
\begin{tcolorbox}[enhanced,nobeforeafter,width=7.1cm,height=2.1cm,valign=center,colback=cdon!6,colframe=cdon,arc=2.5mm,boxrule=0.9pt,left=3mm]{\Large\bfseries\color{cdon}a)}\quad{\Large $13$}\end{tcolorbox}\hfill\begin{tcolorbox}[enhanced,nobeforeafter,width=7.1cm,height=2.1cm,valign=center,colback=cdon!6,colframe=cdon,arc=2.5mm,boxrule=0.9pt,left=3mm]{\Large\bfseries\color{cdon}b)}\quad{\Large $11$}\end{tcolorbox}\par\vspace{2.5mm}
\begin{tcolorbox}[enhanced,nobeforeafter,width=7.1cm,height=2.1cm,valign=center,colback=cdon!6,colframe=cdon,arc=2.5mm,boxrule=0.9pt,left=3mm]{\Large\bfseries\color{cdon}c)}\quad{\Large $7{,}6$}\end{tcolorbox}\hfill\begin{tcolorbox}[enhanced,nobeforeafter,width=7.1cm,height=2.1cm,valign=center,colback=cdon!6,colframe=cdon,arc=2.5mm,boxrule=0.9pt,left=3mm]{\Large\bfseries\color{cdon}d)}\quad{\Large $15$}\end{tcolorbox}
\end{tcolorbox}
\newpage
\begin{tcolorbox}[enhanced,width=16cm,height=14cm,colback=white,colframe=cdon,boxrule=2pt,arc=0mm,
left=6mm,right=6mm,top=0mm,bottom=4mm,sharp corners,
overlay={\fill[cdon] (frame.north west) rectangle ([yshift=-1.25cm]frame.north east);
\node[anchor=west,text=white,font=\sffamily\bfseries\LARGE] at ([xshift=5mm,yshift=-0.62cm]frame.north west) {Question 77 B};
\node[anchor=east,text=white,font=\sffamily\itshape\large] at ([xshift=-5mm,yshift=-0.62cm]frame.north east) {Proportionnalité, statistiques, probabilités};}]
\vspace*{1.75cm}
\begin{minipage}[t][6.9cm][c]{\linewidth}\Large\centering
\begin{minipage}{0.95\linewidth}\raggedright\linespread{1.25}\selectfont Un sac contient $6$ boules rouges, $3$ boules bleues et $3$ boules vertes. On tire une boule au hasard. Quelle est la probabilité de tirer une boule verte~?\end{minipage}
\end{minipage}

\vfill
\begin{tcolorbox}[enhanced,nobeforeafter,width=7.1cm,height=2.1cm,valign=center,colback=cdon!6,colframe=cdon,arc=2.5mm,boxrule=0.9pt,left=3mm]{\Large\bfseries\color{cdon}a)}\quad{\Large $\dfrac{1}{3}$}\end{tcolorbox}\hfill\begin{tcolorbox}[enhanced,nobeforeafter,width=7.1cm,height=2.1cm,valign=center,colback=cdon!6,colframe=cdon,arc=2.5mm,boxrule=0.9pt,left=3mm]{\Large\bfseries\color{cdon}b)}\quad{\Large $\dfrac{3}{13}$}\end{tcolorbox}\par\vspace{2.5mm}
\begin{tcolorbox}[enhanced,nobeforeafter,width=7.1cm,height=2.1cm,valign=center,colback=cdon!6,colframe=cdon,arc=2.5mm,boxrule=0.9pt,left=3mm]{\Large\bfseries\color{cdon}c)}\quad{\Large $\dfrac{1}{4}$}\end{tcolorbox}\hfill\begin{tcolorbox}[enhanced,nobeforeafter,width=7.1cm,height=2.1cm,valign=center,colback=cdon!6,colframe=cdon,arc=2.5mm,boxrule=0.9pt,left=3mm]{\Large\bfseries\color{cdon}d)}\quad{\Large $\dfrac{3}{4}$}\end{tcolorbox}
\end{tcolorbox}
\newpage
\begin{tcolorbox}[enhanced,width=16cm,height=14cm,colback=white,colframe=cdon,boxrule=2pt,arc=0mm,
left=6mm,right=6mm,top=0mm,bottom=4mm,sharp corners,
overlay={\fill[cdon] (frame.north west) rectangle ([yshift=-1.25cm]frame.north east);
\node[anchor=west,text=white,font=\sffamily\bfseries\LARGE] at ([xshift=5mm,yshift=-0.62cm]frame.north west) {Question 78 B};
\node[anchor=east,text=white,font=\sffamily\itshape\large] at ([xshift=-5mm,yshift=-0.62cm]frame.north east) {Proportionnalité, statistiques, probabilités};}]
\vspace*{1.75cm}
\begin{minipage}[t][6.9cm][c]{\linewidth}\Large\centering
\begin{minipage}{0.95\linewidth}\raggedright\linespread{1.25}\selectfont Un sac contient $2$ boules rouges, $7$ boules bleues et $1$ boule verte. On tire une boule au hasard. Quelle est la probabilité de tirer une boule bleue~?\end{minipage}
\end{minipage}

\vfill
\begin{tcolorbox}[enhanced,nobeforeafter,width=7.1cm,height=2.1cm,valign=center,colback=cdon!6,colframe=cdon,arc=2.5mm,boxrule=0.9pt,left=3mm]{\Large\bfseries\color{cdon}a)}\quad{\Large $\dfrac{1}{3}$}\end{tcolorbox}\hfill\begin{tcolorbox}[enhanced,nobeforeafter,width=7.1cm,height=2.1cm,valign=center,colback=cdon!6,colframe=cdon,arc=2.5mm,boxrule=0.9pt,left=3mm]{\Large\bfseries\color{cdon}b)}\quad{\Large $\dfrac{7}{10}$}\end{tcolorbox}\par\vspace{2.5mm}
\begin{tcolorbox}[enhanced,nobeforeafter,width=7.1cm,height=2.1cm,valign=center,colback=cdon!6,colframe=cdon,arc=2.5mm,boxrule=0.9pt,left=3mm]{\Large\bfseries\color{cdon}c)}\quad{\Large $\dfrac{3}{10}$}\end{tcolorbox}\hfill\begin{tcolorbox}[enhanced,nobeforeafter,width=7.1cm,height=2.1cm,valign=center,colback=cdon!6,colframe=cdon,arc=2.5mm,boxrule=0.9pt,left=3mm]{\Large\bfseries\color{cdon}d)}\quad{\Large $\dfrac{7}{3}$}\end{tcolorbox}
\end{tcolorbox}
\newpage
\begin{tcolorbox}[enhanced,width=16cm,height=14cm,colback=white,colframe=cdon,boxrule=2pt,arc=0mm,
left=6mm,right=6mm,top=0mm,bottom=4mm,sharp corners,
overlay={\fill[cdon] (frame.north west) rectangle ([yshift=-1.25cm]frame.north east);
\node[anchor=west,text=white,font=\sffamily\bfseries\LARGE] at ([xshift=5mm,yshift=-0.62cm]frame.north west) {Question 79 B};
\node[anchor=east,text=white,font=\sffamily\itshape\large] at ([xshift=-5mm,yshift=-0.62cm]frame.north east) {Proportionnalité, statistiques, probabilités};}]
\vspace*{1.75cm}
\begin{minipage}[t][6.9cm][c]{\linewidth}\LARGE\centering
\begin{minipage}{0.95\linewidth}\raggedright\linespread{1.25}\selectfont La probabilité qu'un événement $A$ se réalise est $\dfrac{2}{7}$. Quelle est la probabilité de l'événement contraire de $A$~?\end{minipage}
\end{minipage}

\vfill
\begin{tcolorbox}[enhanced,nobeforeafter,width=7.1cm,height=2.1cm,valign=center,colback=cdon!6,colframe=cdon,arc=2.5mm,boxrule=0.9pt,left=3mm]{\Large\bfseries\color{cdon}a)}\quad{\Large $\dfrac{5}{7}$}\end{tcolorbox}\hfill\begin{tcolorbox}[enhanced,nobeforeafter,width=7.1cm,height=2.1cm,valign=center,colback=cdon!6,colframe=cdon,arc=2.5mm,boxrule=0.9pt,left=3mm]{\Large\bfseries\color{cdon}b)}\quad{\Large $\dfrac{6}{7}$}\end{tcolorbox}\par\vspace{2.5mm}
\begin{tcolorbox}[enhanced,nobeforeafter,width=7.1cm,height=2.1cm,valign=center,colback=cdon!6,colframe=cdon,arc=2.5mm,boxrule=0.9pt,left=3mm]{\Large\bfseries\color{cdon}c)}\quad{\Large $\dfrac{3}{14}$}\end{tcolorbox}\hfill\begin{tcolorbox}[enhanced,nobeforeafter,width=7.1cm,height=2.1cm,valign=center,colback=cdon!6,colframe=cdon,arc=2.5mm,boxrule=0.9pt,left=3mm]{\Large\bfseries\color{cdon}d)}\quad{\Large $\dfrac{2}{7}$}\end{tcolorbox}
\end{tcolorbox}
\newpage
\begin{tcolorbox}[enhanced,width=16cm,height=14cm,colback=white,colframe=cdon,boxrule=2pt,arc=0mm,
left=6mm,right=6mm,top=0mm,bottom=4mm,sharp corners,
overlay={\fill[cdon] (frame.north west) rectangle ([yshift=-1.25cm]frame.north east);
\node[anchor=west,text=white,font=\sffamily\bfseries\LARGE] at ([xshift=5mm,yshift=-0.62cm]frame.north west) {Question 80 B};
\node[anchor=east,text=white,font=\sffamily\itshape\large] at ([xshift=-5mm,yshift=-0.62cm]frame.north east) {Proportionnalité, statistiques, probabilités};}]
\vspace*{1.75cm}
\begin{minipage}[t][6.9cm][c]{\linewidth}\LARGE\centering
\begin{minipage}{0.95\linewidth}\raggedright\linespread{1.25}\selectfont La probabilité qu'un événement $A$ se réalise est $0{,}08$. Quelle est la probabilité de l'événement contraire de $A$~?\end{minipage}
\end{minipage}

\vfill
\begin{tcolorbox}[enhanced,nobeforeafter,width=7.1cm,height=2.1cm,valign=center,colback=cdon!6,colframe=cdon,arc=2.5mm,boxrule=0.9pt,left=3mm]{\Large\bfseries\color{cdon}a)}\quad{\Large $0{,}42$}\end{tcolorbox}\hfill\begin{tcolorbox}[enhanced,nobeforeafter,width=7.1cm,height=2.1cm,valign=center,colback=cdon!6,colframe=cdon,arc=2.5mm,boxrule=0.9pt,left=3mm]{\Large\bfseries\color{cdon}b)}\quad{\Large $0{,}96$}\end{tcolorbox}\par\vspace{2.5mm}
\begin{tcolorbox}[enhanced,nobeforeafter,width=7.1cm,height=2.1cm,valign=center,colback=cdon!6,colframe=cdon,arc=2.5mm,boxrule=0.9pt,left=3mm]{\Large\bfseries\color{cdon}c)}\quad{\Large $0{,}92$}\end{tcolorbox}\hfill\begin{tcolorbox}[enhanced,nobeforeafter,width=7.1cm,height=2.1cm,valign=center,colback=cdon!6,colframe=cdon,arc=2.5mm,boxrule=0.9pt,left=3mm]{\Large\bfseries\color{cdon}d)}\quad{\Large $0{,}08$}\end{tcolorbox}
\end{tcolorbox}
\newpage
\begin{tcolorbox}[enhanced,width=16cm,height=14cm,colback=white,colframe=cgra,boxrule=2pt,arc=0mm,
left=6mm,right=6mm,top=0mm,bottom=4mm,sharp corners,
overlay={\fill[cgra] (frame.north west) rectangle ([yshift=-1.25cm]frame.north east);
\node[anchor=west,text=white,font=\sffamily\bfseries\LARGE] at ([xshift=5mm,yshift=-0.62cm]frame.north west) {Question 81 B};
\node[anchor=east,text=white,font=\sffamily\itshape\large] at ([xshift=-5mm,yshift=-0.62cm]frame.north east) {Grandeurs et mesures};}]
\vspace*{1.75cm}
\begin{minipage}[t][6.9cm][c]{\linewidth}\LARGE\centering
\begin{minipage}{0.95\linewidth}\raggedright\linespread{1.25}\selectfont Un véhicule parcourt $90$ km en 1 h 30 min. Quelle est sa vitesse moyenne~?\end{minipage}
\end{minipage}

\vfill
\begin{tcolorbox}[enhanced,nobeforeafter,width=7.1cm,height=2.1cm,valign=center,colback=cgra!6,colframe=cgra,arc=2.5mm,boxrule=0.9pt,left=3mm]{\Large\bfseries\color{cgra}a)}\quad{\Large $135~\text{km/h}$}\end{tcolorbox}\hfill\begin{tcolorbox}[enhanced,nobeforeafter,width=7.1cm,height=2.1cm,valign=center,colback=cgra!6,colframe=cgra,arc=2.5mm,boxrule=0.9pt,left=3mm]{\Large\bfseries\color{cgra}b)}\quad{\Large $60~\text{km/h}$}\end{tcolorbox}\par\vspace{2.5mm}
\begin{tcolorbox}[enhanced,nobeforeafter,width=7.1cm,height=2.1cm,valign=center,colback=cgra!6,colframe=cgra,arc=2.5mm,boxrule=0.9pt,left=3mm]{\Large\bfseries\color{cgra}c)}\quad{\Large $1~\text{km/h}$}\end{tcolorbox}\hfill\begin{tcolorbox}[enhanced,nobeforeafter,width=7.1cm,height=2.1cm,valign=center,colback=cgra!6,colframe=cgra,arc=2.5mm,boxrule=0.9pt,left=3mm]{\Large\bfseries\color{cgra}d)}\quad{\Large $70~\text{km/h}$}\end{tcolorbox}
\end{tcolorbox}
\newpage
\begin{tcolorbox}[enhanced,width=16cm,height=14cm,colback=white,colframe=cgra,boxrule=2pt,arc=0mm,
left=6mm,right=6mm,top=0mm,bottom=4mm,sharp corners,
overlay={\fill[cgra] (frame.north west) rectangle ([yshift=-1.25cm]frame.north east);
\node[anchor=west,text=white,font=\sffamily\bfseries\LARGE] at ([xshift=5mm,yshift=-0.62cm]frame.north west) {Question 82 B};
\node[anchor=east,text=white,font=\sffamily\itshape\large] at ([xshift=-5mm,yshift=-0.62cm]frame.north east) {Grandeurs et mesures};}]
\vspace*{1.75cm}
\begin{minipage}[t][6.9cm][c]{\linewidth}\LARGE\centering
\begin{minipage}{0.95\linewidth}\raggedright\linespread{1.25}\selectfont Un véhicule parcourt $45$ km en 30 min. Quelle est sa vitesse moyenne~?\end{minipage}
\end{minipage}

\vfill
\begin{tcolorbox}[enhanced,nobeforeafter,width=7.1cm,height=2.1cm,valign=center,colback=cgra!6,colframe=cgra,arc=2.5mm,boxrule=0.9pt,left=3mm]{\Large\bfseries\color{cgra}a)}\quad{\Large $1{,}5~\text{km/h}$}\end{tcolorbox}\hfill\begin{tcolorbox}[enhanced,nobeforeafter,width=7.1cm,height=2.1cm,valign=center,colback=cgra!6,colframe=cgra,arc=2.5mm,boxrule=0.9pt,left=3mm]{\Large\bfseries\color{cgra}b)}\quad{\Large $90~\text{km/h}$}\end{tcolorbox}\par\vspace{2.5mm}
\begin{tcolorbox}[enhanced,nobeforeafter,width=7.1cm,height=2.1cm,valign=center,colback=cgra!6,colframe=cgra,arc=2.5mm,boxrule=0.9pt,left=3mm]{\Large\bfseries\color{cgra}c)}\quad{\Large $100~\text{km/h}$}\end{tcolorbox}\hfill\begin{tcolorbox}[enhanced,nobeforeafter,width=7.1cm,height=2.1cm,valign=center,colback=cgra!6,colframe=cgra,arc=2.5mm,boxrule=0.9pt,left=3mm]{\Large\bfseries\color{cgra}d)}\quad{\Large $22{,}5~\text{km/h}$}\end{tcolorbox}
\end{tcolorbox}
\newpage
\begin{tcolorbox}[enhanced,width=16cm,height=14cm,colback=white,colframe=cgra,boxrule=2pt,arc=0mm,
left=6mm,right=6mm,top=0mm,bottom=4mm,sharp corners,
overlay={\fill[cgra] (frame.north west) rectangle ([yshift=-1.25cm]frame.north east);
\node[anchor=west,text=white,font=\sffamily\bfseries\LARGE] at ([xshift=5mm,yshift=-0.62cm]frame.north west) {Question 83 B};
\node[anchor=east,text=white,font=\sffamily\itshape\large] at ([xshift=-5mm,yshift=-0.62cm]frame.north east) {Grandeurs et mesures};}]
\vspace*{1.75cm}
\begin{minipage}[t][6.9cm][c]{\linewidth}\LARGE\centering
\begin{minipage}{0.95\linewidth}\raggedright\linespread{1.25}\selectfont Une voiture roule à la vitesse constante de $60$ km/h pendant $45$ min. Quelle distance parcourt-elle~?\end{minipage}
\end{minipage}

\vfill
\begin{tcolorbox}[enhanced,nobeforeafter,width=7.1cm,height=2.1cm,valign=center,colback=cgra!6,colframe=cgra,arc=2.5mm,boxrule=0.9pt,left=3mm]{\Large\bfseries\color{cgra}a)}\quad{\Large $45~\text{km}$}\end{tcolorbox}\hfill\begin{tcolorbox}[enhanced,nobeforeafter,width=7.1cm,height=2.1cm,valign=center,colback=cgra!6,colframe=cgra,arc=2.5mm,boxrule=0.9pt,left=3mm]{\Large\bfseries\color{cgra}b)}\quad{\Large $2\,700~\text{km}$}\end{tcolorbox}\par\vspace{2.5mm}
\begin{tcolorbox}[enhanced,nobeforeafter,width=7.1cm,height=2.1cm,valign=center,colback=cgra!6,colframe=cgra,arc=2.5mm,boxrule=0.9pt,left=3mm]{\Large\bfseries\color{cgra}c)}\quad{\Large $90~\text{km}$}\end{tcolorbox}\hfill\begin{tcolorbox}[enhanced,nobeforeafter,width=7.1cm,height=2.1cm,valign=center,colback=cgra!6,colframe=cgra,arc=2.5mm,boxrule=0.9pt,left=3mm]{\Large\bfseries\color{cgra}d)}\quad{\Large $27~\text{km}$}\end{tcolorbox}
\end{tcolorbox}
\newpage
\begin{tcolorbox}[enhanced,width=16cm,height=14cm,colback=white,colframe=cgra,boxrule=2pt,arc=0mm,
left=6mm,right=6mm,top=0mm,bottom=4mm,sharp corners,
overlay={\fill[cgra] (frame.north west) rectangle ([yshift=-1.25cm]frame.north east);
\node[anchor=west,text=white,font=\sffamily\bfseries\LARGE] at ([xshift=5mm,yshift=-0.62cm]frame.north west) {Question 84 B};
\node[anchor=east,text=white,font=\sffamily\itshape\large] at ([xshift=-5mm,yshift=-0.62cm]frame.north east) {Grandeurs et mesures};}]
\vspace*{1.75cm}
\begin{minipage}[t][6.9cm][c]{\linewidth}\LARGE\centering
\begin{minipage}{0.95\linewidth}\raggedright\linespread{1.25}\selectfont Une voiture roule à la vitesse constante de $90$ km/h pendant $20$ min. Quelle distance parcourt-elle~?\end{minipage}
\end{minipage}

\vfill
\begin{tcolorbox}[enhanced,nobeforeafter,width=7.1cm,height=2.1cm,valign=center,colback=cgra!6,colframe=cgra,arc=2.5mm,boxrule=0.9pt,left=3mm]{\Large\bfseries\color{cgra}a)}\quad{\Large $60~\text{km}$}\end{tcolorbox}\hfill\begin{tcolorbox}[enhanced,nobeforeafter,width=7.1cm,height=2.1cm,valign=center,colback=cgra!6,colframe=cgra,arc=2.5mm,boxrule=0.9pt,left=3mm]{\Large\bfseries\color{cgra}b)}\quad{\Large $18~\text{km}$}\end{tcolorbox}\par\vspace{2.5mm}
\begin{tcolorbox}[enhanced,nobeforeafter,width=7.1cm,height=2.1cm,valign=center,colback=cgra!6,colframe=cgra,arc=2.5mm,boxrule=0.9pt,left=3mm]{\Large\bfseries\color{cgra}c)}\quad{\Large $30~\text{km}$}\end{tcolorbox}\hfill\begin{tcolorbox}[enhanced,nobeforeafter,width=7.1cm,height=2.1cm,valign=center,colback=cgra!6,colframe=cgra,arc=2.5mm,boxrule=0.9pt,left=3mm]{\Large\bfseries\color{cgra}d)}\quad{\Large $1\,800~\text{km}$}\end{tcolorbox}
\end{tcolorbox}
\newpage
\begin{tcolorbox}[enhanced,width=16cm,height=14cm,colback=white,colframe=cgra,boxrule=2pt,arc=0mm,
left=6mm,right=6mm,top=0mm,bottom=4mm,sharp corners,
overlay={\fill[cgra] (frame.north west) rectangle ([yshift=-1.25cm]frame.north east);
\node[anchor=west,text=white,font=\sffamily\bfseries\LARGE] at ([xshift=5mm,yshift=-0.62cm]frame.north west) {Question 85 B};
\node[anchor=east,text=white,font=\sffamily\itshape\large] at ([xshift=-5mm,yshift=-0.62cm]frame.north east) {Grandeurs et mesures};}]
\vspace*{1.75cm}
\begin{minipage}[t][6.9cm][c]{\linewidth}\LARGE\centering
\begin{minipage}{0.95\linewidth}\raggedright\linespread{1.25}\selectfont Un train parcourt $45$ km à la vitesse moyenne de $90$ km/h. Combien de temps dure le trajet~?\end{minipage}
\end{minipage}

\vfill
\begin{tcolorbox}[enhanced,nobeforeafter,width=7.1cm,height=2.1cm,valign=center,colback=cgra!6,colframe=cgra,arc=2.5mm,boxrule=0.9pt,left=3mm]{\Large\bfseries\color{cgra}a)}\quad{\Large 15 min}\end{tcolorbox}\hfill\begin{tcolorbox}[enhanced,nobeforeafter,width=7.1cm,height=2.1cm,valign=center,colback=cgra!6,colframe=cgra,arc=2.5mm,boxrule=0.9pt,left=3mm]{\Large\bfseries\color{cgra}b)}\quad{\Large 50 min}\end{tcolorbox}\par\vspace{2.5mm}
\begin{tcolorbox}[enhanced,nobeforeafter,width=7.1cm,height=2.1cm,valign=center,colback=cgra!6,colframe=cgra,arc=2.5mm,boxrule=0.9pt,left=3mm]{\Large\bfseries\color{cgra}c)}\quad{\Large 1 h}\end{tcolorbox}\hfill\begin{tcolorbox}[enhanced,nobeforeafter,width=7.1cm,height=2.1cm,valign=center,colback=cgra!6,colframe=cgra,arc=2.5mm,boxrule=0.9pt,left=3mm]{\Large\bfseries\color{cgra}d)}\quad{\Large 30 min}\end{tcolorbox}
\end{tcolorbox}
\newpage
\begin{tcolorbox}[enhanced,width=16cm,height=14cm,colback=white,colframe=cgra,boxrule=2pt,arc=0mm,
left=6mm,right=6mm,top=0mm,bottom=4mm,sharp corners,
overlay={\fill[cgra] (frame.north west) rectangle ([yshift=-1.25cm]frame.north east);
\node[anchor=west,text=white,font=\sffamily\bfseries\LARGE] at ([xshift=5mm,yshift=-0.62cm]frame.north west) {Question 86 B};
\node[anchor=east,text=white,font=\sffamily\itshape\large] at ([xshift=-5mm,yshift=-0.62cm]frame.north east) {Grandeurs et mesures};}]
\vspace*{1.75cm}
\begin{minipage}[t][6.9cm][c]{\linewidth}\LARGE\centering
\begin{minipage}{0.95\linewidth}\raggedright\linespread{1.25}\selectfont Un train parcourt $350$ km à la vitesse moyenne de $100$ km/h. Combien de temps dure le trajet~?\end{minipage}
\end{minipage}

\vfill
\begin{tcolorbox}[enhanced,nobeforeafter,width=7.1cm,height=2.1cm,valign=center,colback=cgra!6,colframe=cgra,arc=2.5mm,boxrule=0.9pt,left=3mm]{\Large\bfseries\color{cgra}a)}\quad{\Large 3 h 50 min}\end{tcolorbox}\hfill\begin{tcolorbox}[enhanced,nobeforeafter,width=7.1cm,height=2.1cm,valign=center,colback=cgra!6,colframe=cgra,arc=2.5mm,boxrule=0.9pt,left=3mm]{\Large\bfseries\color{cgra}b)}\quad{\Large 4 h}\end{tcolorbox}\par\vspace{2.5mm}
\begin{tcolorbox}[enhanced,nobeforeafter,width=7.1cm,height=2.1cm,valign=center,colback=cgra!6,colframe=cgra,arc=2.5mm,boxrule=0.9pt,left=3mm]{\Large\bfseries\color{cgra}c)}\quad{\Large 3 h 30 min}\end{tcolorbox}\hfill\begin{tcolorbox}[enhanced,nobeforeafter,width=7.1cm,height=2.1cm,valign=center,colback=cgra!6,colframe=cgra,arc=2.5mm,boxrule=0.9pt,left=3mm]{\Large\bfseries\color{cgra}d)}\quad{\Large 3 h 15 min}\end{tcolorbox}
\end{tcolorbox}
\newpage
\begin{tcolorbox}[enhanced,width=16cm,height=14cm,colback=white,colframe=cgra,boxrule=2pt,arc=0mm,
left=6mm,right=6mm,top=0mm,bottom=4mm,sharp corners,
overlay={\fill[cgra] (frame.north west) rectangle ([yshift=-1.25cm]frame.north east);
\node[anchor=west,text=white,font=\sffamily\bfseries\LARGE] at ([xshift=5mm,yshift=-0.62cm]frame.north west) {Question 87 B};
\node[anchor=east,text=white,font=\sffamily\itshape\large] at ([xshift=-5mm,yshift=-0.62cm]frame.north east) {Grandeurs et mesures};}]
\vspace*{1.75cm}
\begin{minipage}[t][6.9cm][c]{\linewidth}\LARGE\centering
\begin{minipage}{0.95\linewidth}\raggedright\linespread{1.25}\selectfont Un véhicule roule à $90$ km/h. Quelle est sa vitesse en mètres par seconde~?\end{minipage}
\end{minipage}

\vfill
\begin{tcolorbox}[enhanced,nobeforeafter,width=7.1cm,height=2.1cm,valign=center,colback=cgra!6,colframe=cgra,arc=2.5mm,boxrule=0.9pt,left=3mm]{\Large\bfseries\color{cgra}a)}\quad{\Large $1{,}5~\text{m/s}$}\end{tcolorbox}\hfill\begin{tcolorbox}[enhanced,nobeforeafter,width=7.1cm,height=2.1cm,valign=center,colback=cgra!6,colframe=cgra,arc=2.5mm,boxrule=0.9pt,left=3mm]{\Large\bfseries\color{cgra}b)}\quad{\Large $324~\text{m/s}$}\end{tcolorbox}\par\vspace{2.5mm}
\begin{tcolorbox}[enhanced,nobeforeafter,width=7.1cm,height=2.1cm,valign=center,colback=cgra!6,colframe=cgra,arc=2.5mm,boxrule=0.9pt,left=3mm]{\Large\bfseries\color{cgra}c)}\quad{\Large $25~\text{m/s}$}\end{tcolorbox}\hfill\begin{tcolorbox}[enhanced,nobeforeafter,width=7.1cm,height=2.1cm,valign=center,colback=cgra!6,colframe=cgra,arc=2.5mm,boxrule=0.9pt,left=3mm]{\Large\bfseries\color{cgra}d)}\quad{\Large $90\,000~\text{m/s}$}\end{tcolorbox}
\end{tcolorbox}
\newpage
\begin{tcolorbox}[enhanced,width=16cm,height=14cm,colback=white,colframe=cgra,boxrule=2pt,arc=0mm,
left=6mm,right=6mm,top=0mm,bottom=4mm,sharp corners,
overlay={\fill[cgra] (frame.north west) rectangle ([yshift=-1.25cm]frame.north east);
\node[anchor=west,text=white,font=\sffamily\bfseries\LARGE] at ([xshift=5mm,yshift=-0.62cm]frame.north west) {Question 88 B};
\node[anchor=east,text=white,font=\sffamily\itshape\large] at ([xshift=-5mm,yshift=-0.62cm]frame.north east) {Grandeurs et mesures};}]
\vspace*{1.75cm}
\begin{minipage}[t][6.9cm][c]{\linewidth}\LARGE\centering
\begin{minipage}{0.95\linewidth}\raggedright\linespread{1.25}\selectfont Un véhicule roule à $54$ km/h. Quelle est sa vitesse en mètres par seconde~?\end{minipage}
\end{minipage}

\vfill
\begin{tcolorbox}[enhanced,nobeforeafter,width=7.1cm,height=2.1cm,valign=center,colback=cgra!6,colframe=cgra,arc=2.5mm,boxrule=0.9pt,left=3mm]{\Large\bfseries\color{cgra}a)}\quad{\Large $194{,}4~\text{m/s}$}\end{tcolorbox}\hfill\begin{tcolorbox}[enhanced,nobeforeafter,width=7.1cm,height=2.1cm,valign=center,colback=cgra!6,colframe=cgra,arc=2.5mm,boxrule=0.9pt,left=3mm]{\Large\bfseries\color{cgra}b)}\quad{\Large $0{,}9~\text{m/s}$}\end{tcolorbox}\par\vspace{2.5mm}
\begin{tcolorbox}[enhanced,nobeforeafter,width=7.1cm,height=2.1cm,valign=center,colback=cgra!6,colframe=cgra,arc=2.5mm,boxrule=0.9pt,left=3mm]{\Large\bfseries\color{cgra}c)}\quad{\Large $15~\text{m/s}$}\end{tcolorbox}\hfill\begin{tcolorbox}[enhanced,nobeforeafter,width=7.1cm,height=2.1cm,valign=center,colback=cgra!6,colframe=cgra,arc=2.5mm,boxrule=0.9pt,left=3mm]{\Large\bfseries\color{cgra}d)}\quad{\Large $54\,000~\text{m/s}$}\end{tcolorbox}
\end{tcolorbox}
\newpage
\begin{tcolorbox}[enhanced,width=16cm,height=14cm,colback=white,colframe=cgra,boxrule=2pt,arc=0mm,
left=6mm,right=6mm,top=0mm,bottom=4mm,sharp corners,
overlay={\fill[cgra] (frame.north west) rectangle ([yshift=-1.25cm]frame.north east);
\node[anchor=west,text=white,font=\sffamily\bfseries\LARGE] at ([xshift=5mm,yshift=-0.62cm]frame.north west) {Question 89 B};
\node[anchor=east,text=white,font=\sffamily\itshape\large] at ([xshift=-5mm,yshift=-0.62cm]frame.north east) {Grandeurs et mesures};}]
\vspace*{1.75cm}
\begin{minipage}[t][6.9cm][c]{\linewidth}\LARGE\centering
\begin{minipage}{0.95\linewidth}\raggedright\linespread{1.25}\selectfont Une pyramide a une base d'aire $18$ cm$^2$ et une hauteur de $10$ cm. Quel est son volume~?\end{minipage}
\end{minipage}

\vfill
\begin{tcolorbox}[enhanced,nobeforeafter,width=7.1cm,height=2.1cm,valign=center,colback=cgra!6,colframe=cgra,arc=2.5mm,boxrule=0.9pt,left=3mm]{\Large\bfseries\color{cgra}a)}\quad{\Large $60~\text{cm}^3$}\end{tcolorbox}\hfill\begin{tcolorbox}[enhanced,nobeforeafter,width=7.1cm,height=2.1cm,valign=center,colback=cgra!6,colframe=cgra,arc=2.5mm,boxrule=0.9pt,left=3mm]{\Large\bfseries\color{cgra}b)}\quad{\Large $180~\text{cm}^3$}\end{tcolorbox}\par\vspace{2.5mm}
\begin{tcolorbox}[enhanced,nobeforeafter,width=7.1cm,height=2.1cm,valign=center,colback=cgra!6,colframe=cgra,arc=2.5mm,boxrule=0.9pt,left=3mm]{\Large\bfseries\color{cgra}c)}\quad{\Large $28~\text{cm}^3$}\end{tcolorbox}\hfill\begin{tcolorbox}[enhanced,nobeforeafter,width=7.1cm,height=2.1cm,valign=center,colback=cgra!6,colframe=cgra,arc=2.5mm,boxrule=0.9pt,left=3mm]{\Large\bfseries\color{cgra}d)}\quad{\Large $90~\text{cm}^3$}\end{tcolorbox}
\end{tcolorbox}
\newpage
\begin{tcolorbox}[enhanced,width=16cm,height=14cm,colback=white,colframe=cgra,boxrule=2pt,arc=0mm,
left=6mm,right=6mm,top=0mm,bottom=4mm,sharp corners,
overlay={\fill[cgra] (frame.north west) rectangle ([yshift=-1.25cm]frame.north east);
\node[anchor=west,text=white,font=\sffamily\bfseries\LARGE] at ([xshift=5mm,yshift=-0.62cm]frame.north west) {Question 90 B};
\node[anchor=east,text=white,font=\sffamily\itshape\large] at ([xshift=-5mm,yshift=-0.62cm]frame.north east) {Grandeurs et mesures};}]
\vspace*{1.75cm}
\begin{minipage}[t][6.9cm][c]{\linewidth}\LARGE\centering
\begin{minipage}{0.95\linewidth}\raggedright\linespread{1.25}\selectfont Une pyramide a une base d'aire $9$ cm$^2$ et une hauteur de $7$ cm. Quel est son volume~?\end{minipage}
\end{minipage}

\vfill
\begin{tcolorbox}[enhanced,nobeforeafter,width=7.1cm,height=2.1cm,valign=center,colback=cgra!6,colframe=cgra,arc=2.5mm,boxrule=0.9pt,left=3mm]{\Large\bfseries\color{cgra}a)}\quad{\Large $31{,}5~\text{cm}^3$}\end{tcolorbox}\hfill\begin{tcolorbox}[enhanced,nobeforeafter,width=7.1cm,height=2.1cm,valign=center,colback=cgra!6,colframe=cgra,arc=2.5mm,boxrule=0.9pt,left=3mm]{\Large\bfseries\color{cgra}b)}\quad{\Large $63~\text{cm}^3$}\end{tcolorbox}\par\vspace{2.5mm}
\begin{tcolorbox}[enhanced,nobeforeafter,width=7.1cm,height=2.1cm,valign=center,colback=cgra!6,colframe=cgra,arc=2.5mm,boxrule=0.9pt,left=3mm]{\Large\bfseries\color{cgra}c)}\quad{\Large $16~\text{cm}^3$}\end{tcolorbox}\hfill\begin{tcolorbox}[enhanced,nobeforeafter,width=7.1cm,height=2.1cm,valign=center,colback=cgra!6,colframe=cgra,arc=2.5mm,boxrule=0.9pt,left=3mm]{\Large\bfseries\color{cgra}d)}\quad{\Large $21~\text{cm}^3$}\end{tcolorbox}
\end{tcolorbox}
\newpage
\begin{tcolorbox}[enhanced,width=16cm,height=14cm,colback=white,colframe=cgra,boxrule=2pt,arc=0mm,
left=6mm,right=6mm,top=0mm,bottom=4mm,sharp corners,
overlay={\fill[cgra] (frame.north west) rectangle ([yshift=-1.25cm]frame.north east);
\node[anchor=west,text=white,font=\sffamily\bfseries\LARGE] at ([xshift=5mm,yshift=-0.62cm]frame.north west) {Question 91 B};
\node[anchor=east,text=white,font=\sffamily\itshape\large] at ([xshift=-5mm,yshift=-0.62cm]frame.north east) {Grandeurs et mesures};}]
\vspace*{1.75cm}
\begin{minipage}[t][6.9cm][c]{\linewidth}\LARGE\centering
\begin{minipage}{0.95\linewidth}\raggedright\linespread{1.25}\selectfont Un cône de révolution a un rayon de base de $5$ cm et une hauteur de $3$ cm. Quel est son volume exact~?\end{minipage}
\end{minipage}

\vfill
\begin{tcolorbox}[enhanced,nobeforeafter,width=7.1cm,height=2.1cm,valign=center,colback=cgra!6,colframe=cgra,arc=2.5mm,boxrule=0.9pt,left=3mm]{\Large\bfseries\color{cgra}a)}\quad{\Large $100\pi~\text{cm}^3$}\end{tcolorbox}\hfill\begin{tcolorbox}[enhanced,nobeforeafter,width=7.1cm,height=2.1cm,valign=center,colback=cgra!6,colframe=cgra,arc=2.5mm,boxrule=0.9pt,left=3mm]{\Large\bfseries\color{cgra}b)}\quad{\Large $5\pi~\text{cm}^3$}\end{tcolorbox}\par\vspace{2.5mm}
\begin{tcolorbox}[enhanced,nobeforeafter,width=7.1cm,height=2.1cm,valign=center,colback=cgra!6,colframe=cgra,arc=2.5mm,boxrule=0.9pt,left=3mm]{\Large\bfseries\color{cgra}c)}\quad{\Large $25\pi~\text{cm}^3$}\end{tcolorbox}\hfill\begin{tcolorbox}[enhanced,nobeforeafter,width=7.1cm,height=2.1cm,valign=center,colback=cgra!6,colframe=cgra,arc=2.5mm,boxrule=0.9pt,left=3mm]{\Large\bfseries\color{cgra}d)}\quad{\Large $75\pi~\text{cm}^3$}\end{tcolorbox}
\end{tcolorbox}
\newpage
\begin{tcolorbox}[enhanced,width=16cm,height=14cm,colback=white,colframe=cgra,boxrule=2pt,arc=0mm,
left=6mm,right=6mm,top=0mm,bottom=4mm,sharp corners,
overlay={\fill[cgra] (frame.north west) rectangle ([yshift=-1.25cm]frame.north east);
\node[anchor=west,text=white,font=\sffamily\bfseries\LARGE] at ([xshift=5mm,yshift=-0.62cm]frame.north west) {Question 92 B};
\node[anchor=east,text=white,font=\sffamily\itshape\large] at ([xshift=-5mm,yshift=-0.62cm]frame.north east) {Grandeurs et mesures};}]
\vspace*{1.75cm}
\begin{minipage}[t][6.9cm][c]{\linewidth}\LARGE\centering
\begin{minipage}{0.95\linewidth}\raggedright\linespread{1.25}\selectfont Un cône de révolution a un rayon de base de $6$ cm et une hauteur de $2$ cm. Quel est son volume exact~?\end{minipage}
\end{minipage}

\vfill
\begin{tcolorbox}[enhanced,nobeforeafter,width=7.1cm,height=2.1cm,valign=center,colback=cgra!6,colframe=cgra,arc=2.5mm,boxrule=0.9pt,left=3mm]{\Large\bfseries\color{cgra}a)}\quad{\Large $24\pi~\text{cm}^3$}\end{tcolorbox}\hfill\begin{tcolorbox}[enhanced,nobeforeafter,width=7.1cm,height=2.1cm,valign=center,colback=cgra!6,colframe=cgra,arc=2.5mm,boxrule=0.9pt,left=3mm]{\Large\bfseries\color{cgra}b)}\quad{\Large $4\pi~\text{cm}^3$}\end{tcolorbox}\par\vspace{2.5mm}
\begin{tcolorbox}[enhanced,nobeforeafter,width=7.1cm,height=2.1cm,valign=center,colback=cgra!6,colframe=cgra,arc=2.5mm,boxrule=0.9pt,left=3mm]{\Large\bfseries\color{cgra}c)}\quad{\Large $96\pi~\text{cm}^3$}\end{tcolorbox}\hfill\begin{tcolorbox}[enhanced,nobeforeafter,width=7.1cm,height=2.1cm,valign=center,colback=cgra!6,colframe=cgra,arc=2.5mm,boxrule=0.9pt,left=3mm]{\Large\bfseries\color{cgra}d)}\quad{\Large $72\pi~\text{cm}^3$}\end{tcolorbox}
\end{tcolorbox}
\newpage
\begin{tcolorbox}[enhanced,width=16cm,height=14cm,colback=white,colframe=cgra,boxrule=2pt,arc=0mm,
left=6mm,right=6mm,top=0mm,bottom=4mm,sharp corners,
overlay={\fill[cgra] (frame.north west) rectangle ([yshift=-1.25cm]frame.north east);
\node[anchor=west,text=white,font=\sffamily\bfseries\LARGE] at ([xshift=5mm,yshift=-0.62cm]frame.north west) {Question 93 B};
\node[anchor=east,text=white,font=\sffamily\itshape\large] at ([xshift=-5mm,yshift=-0.62cm]frame.north east) {Grandeurs et mesures};}]
\vspace*{1.75cm}
\begin{minipage}[t][6.9cm][c]{\linewidth}\LARGE\centering
\begin{minipage}{0.95\linewidth}\raggedright\linespread{1.25}\selectfont Compléter~: $4~\text{L}=\;?\;\text{cm}^3$\end{minipage}
\end{minipage}

\vfill
\begin{tcolorbox}[enhanced,nobeforeafter,width=7.1cm,height=2.1cm,valign=center,colback=cgra!6,colframe=cgra,arc=2.5mm,boxrule=0.9pt,left=3mm]{\Large\bfseries\color{cgra}a)}\quad{\Large $400\,000~\text{cm}^3$}\end{tcolorbox}\hfill\begin{tcolorbox}[enhanced,nobeforeafter,width=7.1cm,height=2.1cm,valign=center,colback=cgra!6,colframe=cgra,arc=2.5mm,boxrule=0.9pt,left=3mm]{\Large\bfseries\color{cgra}b)}\quad{\Large $400~\text{cm}^3$}\end{tcolorbox}\par\vspace{2.5mm}
\begin{tcolorbox}[enhanced,nobeforeafter,width=7.1cm,height=2.1cm,valign=center,colback=cgra!6,colframe=cgra,arc=2.5mm,boxrule=0.9pt,left=3mm]{\Large\bfseries\color{cgra}c)}\quad{\Large $4\,000~\text{cm}^3$}\end{tcolorbox}\hfill\begin{tcolorbox}[enhanced,nobeforeafter,width=7.1cm,height=2.1cm,valign=center,colback=cgra!6,colframe=cgra,arc=2.5mm,boxrule=0.9pt,left=3mm]{\Large\bfseries\color{cgra}d)}\quad{\Large $40\,000~\text{cm}^3$}\end{tcolorbox}
\end{tcolorbox}
\newpage
\begin{tcolorbox}[enhanced,width=16cm,height=14cm,colback=white,colframe=cgra,boxrule=2pt,arc=0mm,
left=6mm,right=6mm,top=0mm,bottom=4mm,sharp corners,
overlay={\fill[cgra] (frame.north west) rectangle ([yshift=-1.25cm]frame.north east);
\node[anchor=west,text=white,font=\sffamily\bfseries\LARGE] at ([xshift=5mm,yshift=-0.62cm]frame.north west) {Question 94 B};
\node[anchor=east,text=white,font=\sffamily\itshape\large] at ([xshift=-5mm,yshift=-0.62cm]frame.north east) {Grandeurs et mesures};}]
\vspace*{1.75cm}
\begin{minipage}[t][6.9cm][c]{\linewidth}\LARGE\centering
\begin{minipage}{0.95\linewidth}\raggedright\linespread{1.25}\selectfont Compléter~: $250~\text{cm}^3=\;?\;\text{L}$\end{minipage}
\end{minipage}

\vfill
\begin{tcolorbox}[enhanced,nobeforeafter,width=7.1cm,height=2.1cm,valign=center,colback=cgra!6,colframe=cgra,arc=2.5mm,boxrule=0.9pt,left=3mm]{\Large\bfseries\color{cgra}a)}\quad{\Large $0{,}25~\text{L}$}\end{tcolorbox}\hfill\begin{tcolorbox}[enhanced,nobeforeafter,width=7.1cm,height=2.1cm,valign=center,colback=cgra!6,colframe=cgra,arc=2.5mm,boxrule=0.9pt,left=3mm]{\Large\bfseries\color{cgra}b)}\quad{\Large $0{,}025~\text{L}$}\end{tcolorbox}\par\vspace{2.5mm}
\begin{tcolorbox}[enhanced,nobeforeafter,width=7.1cm,height=2.1cm,valign=center,colback=cgra!6,colframe=cgra,arc=2.5mm,boxrule=0.9pt,left=3mm]{\Large\bfseries\color{cgra}c)}\quad{\Large $25~\text{L}$}\end{tcolorbox}\hfill\begin{tcolorbox}[enhanced,nobeforeafter,width=7.1cm,height=2.1cm,valign=center,colback=cgra!6,colframe=cgra,arc=2.5mm,boxrule=0.9pt,left=3mm]{\Large\bfseries\color{cgra}d)}\quad{\Large $2{,}5~\text{L}$}\end{tcolorbox}
\end{tcolorbox}
\newpage
\begin{tcolorbox}[enhanced,width=16cm,height=14cm,colback=white,colframe=cgra,boxrule=2pt,arc=0mm,
left=6mm,right=6mm,top=0mm,bottom=4mm,sharp corners,
overlay={\fill[cgra] (frame.north west) rectangle ([yshift=-1.25cm]frame.north east);
\node[anchor=west,text=white,font=\sffamily\bfseries\LARGE] at ([xshift=5mm,yshift=-0.62cm]frame.north west) {Question 95 B};
\node[anchor=east,text=white,font=\sffamily\itshape\large] at ([xshift=-5mm,yshift=-0.62cm]frame.north east) {Grandeurs et mesures};}]
\vspace*{1.75cm}
\begin{minipage}[t][6.9cm][c]{\linewidth}\LARGE\centering
\begin{minipage}{0.95\linewidth}\raggedright\linespread{1.25}\selectfont Compléter~: $45\,000~\text{cm}^2=\;?\;\text{m}^2$\end{minipage}
\end{minipage}

\vfill
\begin{tcolorbox}[enhanced,nobeforeafter,width=7.1cm,height=2.1cm,valign=center,colback=cgra!6,colframe=cgra,arc=2.5mm,boxrule=0.9pt,left=3mm]{\Large\bfseries\color{cgra}a)}\quad{\Large $4{,}5~\text{m}^2$}\end{tcolorbox}\hfill\begin{tcolorbox}[enhanced,nobeforeafter,width=7.1cm,height=2.1cm,valign=center,colback=cgra!6,colframe=cgra,arc=2.5mm,boxrule=0.9pt,left=3mm]{\Large\bfseries\color{cgra}b)}\quad{\Large $0{,}045~\text{m}^2$}\end{tcolorbox}\par\vspace{2.5mm}
\begin{tcolorbox}[enhanced,nobeforeafter,width=7.1cm,height=2.1cm,valign=center,colback=cgra!6,colframe=cgra,arc=2.5mm,boxrule=0.9pt,left=3mm]{\Large\bfseries\color{cgra}c)}\quad{\Large $450~\text{m}^2$}\end{tcolorbox}\hfill\begin{tcolorbox}[enhanced,nobeforeafter,width=7.1cm,height=2.1cm,valign=center,colback=cgra!6,colframe=cgra,arc=2.5mm,boxrule=0.9pt,left=3mm]{\Large\bfseries\color{cgra}d)}\quad{\Large $45~\text{m}^2$}\end{tcolorbox}
\end{tcolorbox}
\newpage
\begin{tcolorbox}[enhanced,width=16cm,height=14cm,colback=white,colframe=cgra,boxrule=2pt,arc=0mm,
left=6mm,right=6mm,top=0mm,bottom=4mm,sharp corners,
overlay={\fill[cgra] (frame.north west) rectangle ([yshift=-1.25cm]frame.north east);
\node[anchor=west,text=white,font=\sffamily\bfseries\LARGE] at ([xshift=5mm,yshift=-0.62cm]frame.north west) {Question 96 B};
\node[anchor=east,text=white,font=\sffamily\itshape\large] at ([xshift=-5mm,yshift=-0.62cm]frame.north east) {Grandeurs et mesures};}]
\vspace*{1.75cm}
\begin{minipage}[t][6.9cm][c]{\linewidth}\LARGE\centering
\begin{minipage}{0.95\linewidth}\raggedright\linespread{1.25}\selectfont Compléter~: $1{,}2~\text{km}^2=\;?\;\text{m}^2$\end{minipage}
\end{minipage}

\vfill
\begin{tcolorbox}[enhanced,nobeforeafter,width=7.1cm,height=2.1cm,valign=center,colback=cgra!6,colframe=cgra,arc=2.5mm,boxrule=0.9pt,left=3mm]{\Large\bfseries\color{cgra}a)}\quad{\Large $12\,000\,000~\text{m}^2$}\end{tcolorbox}\hfill\begin{tcolorbox}[enhanced,nobeforeafter,width=7.1cm,height=2.1cm,valign=center,colback=cgra!6,colframe=cgra,arc=2.5mm,boxrule=0.9pt,left=3mm]{\Large\bfseries\color{cgra}b)}\quad{\Large $1\,200~\text{m}^2$}\end{tcolorbox}\par\vspace{2.5mm}
\begin{tcolorbox}[enhanced,nobeforeafter,width=7.1cm,height=2.1cm,valign=center,colback=cgra!6,colframe=cgra,arc=2.5mm,boxrule=0.9pt,left=3mm]{\Large\bfseries\color{cgra}c)}\quad{\Large $1\,200\,000~\text{m}^2$}\end{tcolorbox}\hfill\begin{tcolorbox}[enhanced,nobeforeafter,width=7.1cm,height=2.1cm,valign=center,colback=cgra!6,colframe=cgra,arc=2.5mm,boxrule=0.9pt,left=3mm]{\Large\bfseries\color{cgra}d)}\quad{\Large $1\,200\,000\,000~\text{m}^2$}\end{tcolorbox}
\end{tcolorbox}
\newpage
\begin{tcolorbox}[enhanced,width=16cm,height=14cm,colback=white,colframe=crel,boxrule=2pt,arc=0mm,
left=6mm,right=6mm,top=0mm,bottom=4mm,sharp corners,
overlay={\fill[crel] (frame.north west) rectangle ([yshift=-1.25cm]frame.north east);
\node[anchor=west,text=white,font=\sffamily\bfseries\LARGE] at ([xshift=5mm,yshift=-0.62cm]frame.north west) {Question 1 C};
\node[anchor=east,text=white,font=\sffamily\itshape\large] at ([xshift=-5mm,yshift=-0.62cm]frame.north east) {Relatifs et fractions};}]
\vspace*{1.75cm}
\begin{minipage}[t][6.9cm][c]{\linewidth}\LARGE\centering
\begin{minipage}{0.95\linewidth}\raggedright\linespread{1.25}\selectfont Que vaut $(-5)\times(-2)\times(-3)$~?\end{minipage}
\end{minipage}

\vfill
\begin{tcolorbox}[enhanced,nobeforeafter,width=7.1cm,height=2.1cm,valign=center,colback=crel!6,colframe=crel,arc=2.5mm,boxrule=0.9pt,left=3mm]{\Large\bfseries\color{crel}a)}\quad{\Large $30$}\end{tcolorbox}\hfill\begin{tcolorbox}[enhanced,nobeforeafter,width=7.1cm,height=2.1cm,valign=center,colback=crel!6,colframe=crel,arc=2.5mm,boxrule=0.9pt,left=3mm]{\Large\bfseries\color{crel}b)}\quad{\Large $13$}\end{tcolorbox}\par\vspace{2.5mm}
\begin{tcolorbox}[enhanced,nobeforeafter,width=7.1cm,height=2.1cm,valign=center,colback=crel!6,colframe=crel,arc=2.5mm,boxrule=0.9pt,left=3mm]{\Large\bfseries\color{crel}c)}\quad{\Large $-10$}\end{tcolorbox}\hfill\begin{tcolorbox}[enhanced,nobeforeafter,width=7.1cm,height=2.1cm,valign=center,colback=crel!6,colframe=crel,arc=2.5mm,boxrule=0.9pt,left=3mm]{\Large\bfseries\color{crel}d)}\quad{\Large $-30$}\end{tcolorbox}
\end{tcolorbox}
\newpage
\begin{tcolorbox}[enhanced,width=16cm,height=14cm,colback=white,colframe=crel,boxrule=2pt,arc=0mm,
left=6mm,right=6mm,top=0mm,bottom=4mm,sharp corners,
overlay={\fill[crel] (frame.north west) rectangle ([yshift=-1.25cm]frame.north east);
\node[anchor=west,text=white,font=\sffamily\bfseries\LARGE] at ([xshift=5mm,yshift=-0.62cm]frame.north west) {Question 2 C};
\node[anchor=east,text=white,font=\sffamily\itshape\large] at ([xshift=-5mm,yshift=-0.62cm]frame.north east) {Relatifs et fractions};}]
\vspace*{1.75cm}
\begin{minipage}[t][6.9cm][c]{\linewidth}\LARGE\centering
\begin{minipage}{0.95\linewidth}\raggedright\linespread{1.25}\selectfont Que vaut $(+3)\times(-4)\times(+5)$~?\end{minipage}
\end{minipage}

\vfill
\begin{tcolorbox}[enhanced,nobeforeafter,width=7.1cm,height=2.1cm,valign=center,colback=crel!6,colframe=crel,arc=2.5mm,boxrule=0.9pt,left=3mm]{\Large\bfseries\color{crel}a)}\quad{\Large $60$}\end{tcolorbox}\hfill\begin{tcolorbox}[enhanced,nobeforeafter,width=7.1cm,height=2.1cm,valign=center,colback=crel!6,colframe=crel,arc=2.5mm,boxrule=0.9pt,left=3mm]{\Large\bfseries\color{crel}b)}\quad{\Large $-17$}\end{tcolorbox}\par\vspace{2.5mm}
\begin{tcolorbox}[enhanced,nobeforeafter,width=7.1cm,height=2.1cm,valign=center,colback=crel!6,colframe=crel,arc=2.5mm,boxrule=0.9pt,left=3mm]{\Large\bfseries\color{crel}c)}\quad{\Large $4$}\end{tcolorbox}\hfill\begin{tcolorbox}[enhanced,nobeforeafter,width=7.1cm,height=2.1cm,valign=center,colback=crel!6,colframe=crel,arc=2.5mm,boxrule=0.9pt,left=3mm]{\Large\bfseries\color{crel}d)}\quad{\Large $-60$}\end{tcolorbox}
\end{tcolorbox}
\newpage
\begin{tcolorbox}[enhanced,width=16cm,height=14cm,colback=white,colframe=crel,boxrule=2pt,arc=0mm,
left=6mm,right=6mm,top=0mm,bottom=4mm,sharp corners,
overlay={\fill[crel] (frame.north west) rectangle ([yshift=-1.25cm]frame.north east);
\node[anchor=west,text=white,font=\sffamily\bfseries\LARGE] at ([xshift=5mm,yshift=-0.62cm]frame.north west) {Question 3 C};
\node[anchor=east,text=white,font=\sffamily\itshape\large] at ([xshift=-5mm,yshift=-0.62cm]frame.north east) {Relatifs et fractions};}]
\vspace*{1.75cm}
\begin{minipage}[t][6.9cm][c]{\linewidth}\LARGE\centering
\begin{minipage}{0.95\linewidth}\raggedright\linespread{1.25}\selectfont Que vaut $(-63)\div(+7)$~?\end{minipage}
\end{minipage}

\vfill
\begin{tcolorbox}[enhanced,nobeforeafter,width=7.1cm,height=2.1cm,valign=center,colback=crel!6,colframe=crel,arc=2.5mm,boxrule=0.9pt,left=3mm]{\Large\bfseries\color{crel}a)}\quad{\Large $-9$}\end{tcolorbox}\hfill\begin{tcolorbox}[enhanced,nobeforeafter,width=7.1cm,height=2.1cm,valign=center,colback=crel!6,colframe=crel,arc=2.5mm,boxrule=0.9pt,left=3mm]{\Large\bfseries\color{crel}b)}\quad{\Large $-18$}\end{tcolorbox}\par\vspace{2.5mm}
\begin{tcolorbox}[enhanced,nobeforeafter,width=7.1cm,height=2.1cm,valign=center,colback=crel!6,colframe=crel,arc=2.5mm,boxrule=0.9pt,left=3mm]{\Large\bfseries\color{crel}c)}\quad{\Large $-8$}\end{tcolorbox}\hfill\begin{tcolorbox}[enhanced,nobeforeafter,width=7.1cm,height=2.1cm,valign=center,colback=crel!6,colframe=crel,arc=2.5mm,boxrule=0.9pt,left=3mm]{\Large\bfseries\color{crel}d)}\quad{\Large $9$}\end{tcolorbox}
\end{tcolorbox}
\newpage
\begin{tcolorbox}[enhanced,width=16cm,height=14cm,colback=white,colframe=crel,boxrule=2pt,arc=0mm,
left=6mm,right=6mm,top=0mm,bottom=4mm,sharp corners,
overlay={\fill[crel] (frame.north west) rectangle ([yshift=-1.25cm]frame.north east);
\node[anchor=west,text=white,font=\sffamily\bfseries\LARGE] at ([xshift=5mm,yshift=-0.62cm]frame.north west) {Question 4 C};
\node[anchor=east,text=white,font=\sffamily\itshape\large] at ([xshift=-5mm,yshift=-0.62cm]frame.north east) {Relatifs et fractions};}]
\vspace*{1.75cm}
\begin{minipage}[t][6.9cm][c]{\linewidth}\LARGE\centering
\begin{minipage}{0.95\linewidth}\raggedright\linespread{1.25}\selectfont Que vaut $(-48)\div(-6)$~?\end{minipage}
\end{minipage}

\vfill
\begin{tcolorbox}[enhanced,nobeforeafter,width=7.1cm,height=2.1cm,valign=center,colback=crel!6,colframe=crel,arc=2.5mm,boxrule=0.9pt,left=3mm]{\Large\bfseries\color{crel}a)}\quad{\Large $8$}\end{tcolorbox}\hfill\begin{tcolorbox}[enhanced,nobeforeafter,width=7.1cm,height=2.1cm,valign=center,colback=crel!6,colframe=crel,arc=2.5mm,boxrule=0.9pt,left=3mm]{\Large\bfseries\color{crel}b)}\quad{\Large $9$}\end{tcolorbox}\par\vspace{2.5mm}
\begin{tcolorbox}[enhanced,nobeforeafter,width=7.1cm,height=2.1cm,valign=center,colback=crel!6,colframe=crel,arc=2.5mm,boxrule=0.9pt,left=3mm]{\Large\bfseries\color{crel}c)}\quad{\Large $16$}\end{tcolorbox}\hfill\begin{tcolorbox}[enhanced,nobeforeafter,width=7.1cm,height=2.1cm,valign=center,colback=crel!6,colframe=crel,arc=2.5mm,boxrule=0.9pt,left=3mm]{\Large\bfseries\color{crel}d)}\quad{\Large $-8$}\end{tcolorbox}
\end{tcolorbox}
\newpage
\begin{tcolorbox}[enhanced,width=16cm,height=14cm,colback=white,colframe=crel,boxrule=2pt,arc=0mm,
left=6mm,right=6mm,top=0mm,bottom=4mm,sharp corners,
overlay={\fill[crel] (frame.north west) rectangle ([yshift=-1.25cm]frame.north east);
\node[anchor=west,text=white,font=\sffamily\bfseries\LARGE] at ([xshift=5mm,yshift=-0.62cm]frame.north west) {Question 5 C};
\node[anchor=east,text=white,font=\sffamily\itshape\large] at ([xshift=-5mm,yshift=-0.62cm]frame.north east) {Relatifs et fractions};}]
\vspace*{1.75cm}
\begin{minipage}[t][6.9cm][c]{\linewidth}\LARGE\centering
\begin{minipage}{0.95\linewidth}\raggedright\linespread{1.25}\selectfont Que vaut $-3-6\times(-2)$~?\end{minipage}
\end{minipage}

\vfill
\begin{tcolorbox}[enhanced,nobeforeafter,width=7.1cm,height=2.1cm,valign=center,colback=crel!6,colframe=crel,arc=2.5mm,boxrule=0.9pt,left=3mm]{\Large\bfseries\color{crel}a)}\quad{\Large $-9$}\end{tcolorbox}\hfill\begin{tcolorbox}[enhanced,nobeforeafter,width=7.1cm,height=2.1cm,valign=center,colback=crel!6,colframe=crel,arc=2.5mm,boxrule=0.9pt,left=3mm]{\Large\bfseries\color{crel}b)}\quad{\Large $9$}\end{tcolorbox}\par\vspace{2.5mm}
\begin{tcolorbox}[enhanced,nobeforeafter,width=7.1cm,height=2.1cm,valign=center,colback=crel!6,colframe=crel,arc=2.5mm,boxrule=0.9pt,left=3mm]{\Large\bfseries\color{crel}c)}\quad{\Large $-15$}\end{tcolorbox}\hfill\begin{tcolorbox}[enhanced,nobeforeafter,width=7.1cm,height=2.1cm,valign=center,colback=crel!6,colframe=crel,arc=2.5mm,boxrule=0.9pt,left=3mm]{\Large\bfseries\color{crel}d)}\quad{\Large $18$}\end{tcolorbox}
\end{tcolorbox}
\newpage
\begin{tcolorbox}[enhanced,width=16cm,height=14cm,colback=white,colframe=crel,boxrule=2pt,arc=0mm,
left=6mm,right=6mm,top=0mm,bottom=4mm,sharp corners,
overlay={\fill[crel] (frame.north west) rectangle ([yshift=-1.25cm]frame.north east);
\node[anchor=west,text=white,font=\sffamily\bfseries\LARGE] at ([xshift=5mm,yshift=-0.62cm]frame.north west) {Question 6 C};
\node[anchor=east,text=white,font=\sffamily\itshape\large] at ([xshift=-5mm,yshift=-0.62cm]frame.north east) {Relatifs et fractions};}]
\vspace*{1.75cm}
\begin{minipage}[t][6.9cm][c]{\linewidth}\LARGE\centering
\begin{minipage}{0.95\linewidth}\raggedright\linespread{1.25}\selectfont Que vaut $-6-4\times(-5)$~?\end{minipage}
\end{minipage}

\vfill
\begin{tcolorbox}[enhanced,nobeforeafter,width=7.1cm,height=2.1cm,valign=center,colback=crel!6,colframe=crel,arc=2.5mm,boxrule=0.9pt,left=3mm]{\Large\bfseries\color{crel}a)}\quad{\Large $50$}\end{tcolorbox}\hfill\begin{tcolorbox}[enhanced,nobeforeafter,width=7.1cm,height=2.1cm,valign=center,colback=crel!6,colframe=crel,arc=2.5mm,boxrule=0.9pt,left=3mm]{\Large\bfseries\color{crel}b)}\quad{\Large $-14$}\end{tcolorbox}\par\vspace{2.5mm}
\begin{tcolorbox}[enhanced,nobeforeafter,width=7.1cm,height=2.1cm,valign=center,colback=crel!6,colframe=crel,arc=2.5mm,boxrule=0.9pt,left=3mm]{\Large\bfseries\color{crel}c)}\quad{\Large $14$}\end{tcolorbox}\hfill\begin{tcolorbox}[enhanced,nobeforeafter,width=7.1cm,height=2.1cm,valign=center,colback=crel!6,colframe=crel,arc=2.5mm,boxrule=0.9pt,left=3mm]{\Large\bfseries\color{crel}d)}\quad{\Large $-26$}\end{tcolorbox}
\end{tcolorbox}
\newpage
\begin{tcolorbox}[enhanced,width=16cm,height=14cm,colback=white,colframe=crel,boxrule=2pt,arc=0mm,
left=6mm,right=6mm,top=0mm,bottom=4mm,sharp corners,
overlay={\fill[crel] (frame.north west) rectangle ([yshift=-1.25cm]frame.north east);
\node[anchor=west,text=white,font=\sffamily\bfseries\LARGE] at ([xshift=5mm,yshift=-0.62cm]frame.north west) {Question 7 C};
\node[anchor=east,text=white,font=\sffamily\itshape\large] at ([xshift=-5mm,yshift=-0.62cm]frame.north east) {Relatifs et fractions};}]
\vspace*{1.75cm}
\begin{minipage}[t][6.9cm][c]{\linewidth}\LARGE\centering
\begin{minipage}{0.95\linewidth}\raggedright\linespread{1.25}\selectfont Que vaut $\dfrac{5}{8}+\dfrac{1}{6}$~?\end{minipage}
\end{minipage}

\vfill
\begin{tcolorbox}[enhanced,nobeforeafter,width=7.1cm,height=2.1cm,valign=center,colback=crel!6,colframe=crel,arc=2.5mm,boxrule=0.9pt,left=3mm]{\Large\bfseries\color{crel}a)}\quad{\Large $\dfrac{19}{24}$}\end{tcolorbox}\hfill\begin{tcolorbox}[enhanced,nobeforeafter,width=7.1cm,height=2.1cm,valign=center,colback=crel!6,colframe=crel,arc=2.5mm,boxrule=0.9pt,left=3mm]{\Large\bfseries\color{crel}b)}\quad{\Large $\dfrac{1}{8}$}\end{tcolorbox}\par\vspace{2.5mm}
\begin{tcolorbox}[enhanced,nobeforeafter,width=7.1cm,height=2.1cm,valign=center,colback=crel!6,colframe=crel,arc=2.5mm,boxrule=0.9pt,left=3mm]{\Large\bfseries\color{crel}c)}\quad{\Large $\dfrac{5}{48}$}\end{tcolorbox}\hfill\begin{tcolorbox}[enhanced,nobeforeafter,width=7.1cm,height=2.1cm,valign=center,colback=crel!6,colframe=crel,arc=2.5mm,boxrule=0.9pt,left=3mm]{\Large\bfseries\color{crel}d)}\quad{\Large $\dfrac{3}{7}$}\end{tcolorbox}
\end{tcolorbox}
\newpage
\begin{tcolorbox}[enhanced,width=16cm,height=14cm,colback=white,colframe=crel,boxrule=2pt,arc=0mm,
left=6mm,right=6mm,top=0mm,bottom=4mm,sharp corners,
overlay={\fill[crel] (frame.north west) rectangle ([yshift=-1.25cm]frame.north east);
\node[anchor=west,text=white,font=\sffamily\bfseries\LARGE] at ([xshift=5mm,yshift=-0.62cm]frame.north west) {Question 8 C};
\node[anchor=east,text=white,font=\sffamily\itshape\large] at ([xshift=-5mm,yshift=-0.62cm]frame.north east) {Relatifs et fractions};}]
\vspace*{1.75cm}
\begin{minipage}[t][6.9cm][c]{\linewidth}\LARGE\centering
\begin{minipage}{0.95\linewidth}\raggedright\linespread{1.25}\selectfont Que vaut $\dfrac{2}{3}+\dfrac{3}{7}$~?\end{minipage}
\end{minipage}

\vfill
\begin{tcolorbox}[enhanced,nobeforeafter,width=7.1cm,height=2.1cm,valign=center,colback=crel!6,colframe=crel,arc=2.5mm,boxrule=0.9pt,left=3mm]{\Large\bfseries\color{crel}a)}\quad{\Large $\dfrac{2}{7}$}\end{tcolorbox}\hfill\begin{tcolorbox}[enhanced,nobeforeafter,width=7.1cm,height=2.1cm,valign=center,colback=crel!6,colframe=crel,arc=2.5mm,boxrule=0.9pt,left=3mm]{\Large\bfseries\color{crel}b)}\quad{\Large $\dfrac{1}{2}$}\end{tcolorbox}\par\vspace{2.5mm}
\begin{tcolorbox}[enhanced,nobeforeafter,width=7.1cm,height=2.1cm,valign=center,colback=crel!6,colframe=crel,arc=2.5mm,boxrule=0.9pt,left=3mm]{\Large\bfseries\color{crel}c)}\quad{\Large $\dfrac{23}{21}$}\end{tcolorbox}\hfill\begin{tcolorbox}[enhanced,nobeforeafter,width=7.1cm,height=2.1cm,valign=center,colback=crel!6,colframe=crel,arc=2.5mm,boxrule=0.9pt,left=3mm]{\Large\bfseries\color{crel}d)}\quad{\Large $\dfrac{5}{21}$}\end{tcolorbox}
\end{tcolorbox}
\newpage
\begin{tcolorbox}[enhanced,width=16cm,height=14cm,colback=white,colframe=crel,boxrule=2pt,arc=0mm,
left=6mm,right=6mm,top=0mm,bottom=4mm,sharp corners,
overlay={\fill[crel] (frame.north west) rectangle ([yshift=-1.25cm]frame.north east);
\node[anchor=west,text=white,font=\sffamily\bfseries\LARGE] at ([xshift=5mm,yshift=-0.62cm]frame.north west) {Question 9 C};
\node[anchor=east,text=white,font=\sffamily\itshape\large] at ([xshift=-5mm,yshift=-0.62cm]frame.north east) {Relatifs et fractions};}]
\vspace*{1.75cm}
\begin{minipage}[t][6.9cm][c]{\linewidth}\LARGE\centering
\begin{minipage}{0.95\linewidth}\raggedright\linespread{1.25}\selectfont Que vaut $\dfrac{7}{12}\times\dfrac{8}{21}$~?\end{minipage}
\end{minipage}

\vfill
\begin{tcolorbox}[enhanced,nobeforeafter,width=7.1cm,height=2.1cm,valign=center,colback=crel!6,colframe=crel,arc=2.5mm,boxrule=0.9pt,left=3mm]{\Large\bfseries\color{crel}a)}\quad{\Large $\dfrac{2}{9}$}\end{tcolorbox}\hfill\begin{tcolorbox}[enhanced,nobeforeafter,width=7.1cm,height=2.1cm,valign=center,colback=crel!6,colframe=crel,arc=2.5mm,boxrule=0.9pt,left=3mm]{\Large\bfseries\color{crel}b)}\quad{\Large $-\dfrac{2}{9}$}\end{tcolorbox}\par\vspace{2.5mm}
\begin{tcolorbox}[enhanced,nobeforeafter,width=7.1cm,height=2.1cm,valign=center,colback=crel!6,colframe=crel,arc=2.5mm,boxrule=0.9pt,left=3mm]{\Large\bfseries\color{crel}c)}\quad{\Large $\dfrac{49}{32}$}\end{tcolorbox}\hfill\begin{tcolorbox}[enhanced,nobeforeafter,width=7.1cm,height=2.1cm,valign=center,colback=crel!6,colframe=crel,arc=2.5mm,boxrule=0.9pt,left=3mm]{\Large\bfseries\color{crel}d)}\quad{\Large $\dfrac{5}{11}$}\end{tcolorbox}
\end{tcolorbox}
\newpage
\begin{tcolorbox}[enhanced,width=16cm,height=14cm,colback=white,colframe=crel,boxrule=2pt,arc=0mm,
left=6mm,right=6mm,top=0mm,bottom=4mm,sharp corners,
overlay={\fill[crel] (frame.north west) rectangle ([yshift=-1.25cm]frame.north east);
\node[anchor=west,text=white,font=\sffamily\bfseries\LARGE] at ([xshift=5mm,yshift=-0.62cm]frame.north west) {Question 10 C};
\node[anchor=east,text=white,font=\sffamily\itshape\large] at ([xshift=-5mm,yshift=-0.62cm]frame.north east) {Relatifs et fractions};}]
\vspace*{1.75cm}
\begin{minipage}[t][6.9cm][c]{\linewidth}\LARGE\centering
\begin{minipage}{0.95\linewidth}\raggedright\linespread{1.25}\selectfont Que vaut $\left(-\dfrac{2}{9}\right)\times\left(-\dfrac{3}{4}\right)$~?\end{minipage}
\end{minipage}

\vfill
\begin{tcolorbox}[enhanced,nobeforeafter,width=7.1cm,height=2.1cm,valign=center,colback=crel!6,colframe=crel,arc=2.5mm,boxrule=0.9pt,left=3mm]{\Large\bfseries\color{crel}a)}\quad{\Large $\dfrac{1}{5}$}\end{tcolorbox}\hfill\begin{tcolorbox}[enhanced,nobeforeafter,width=7.1cm,height=2.1cm,valign=center,colback=crel!6,colframe=crel,arc=2.5mm,boxrule=0.9pt,left=3mm]{\Large\bfseries\color{crel}b)}\quad{\Large $\dfrac{1}{6}$}\end{tcolorbox}\par\vspace{2.5mm}
\begin{tcolorbox}[enhanced,nobeforeafter,width=7.1cm,height=2.1cm,valign=center,colback=crel!6,colframe=crel,arc=2.5mm,boxrule=0.9pt,left=3mm]{\Large\bfseries\color{crel}c)}\quad{\Large $-\dfrac{1}{6}$}\end{tcolorbox}\hfill\begin{tcolorbox}[enhanced,nobeforeafter,width=7.1cm,height=2.1cm,valign=center,colback=crel!6,colframe=crel,arc=2.5mm,boxrule=0.9pt,left=3mm]{\Large\bfseries\color{crel}d)}\quad{\Large $\dfrac{8}{27}$}\end{tcolorbox}
\end{tcolorbox}
\newpage
\begin{tcolorbox}[enhanced,width=16cm,height=14cm,colback=white,colframe=crel,boxrule=2pt,arc=0mm,
left=6mm,right=6mm,top=0mm,bottom=4mm,sharp corners,
overlay={\fill[crel] (frame.north west) rectangle ([yshift=-1.25cm]frame.north east);
\node[anchor=west,text=white,font=\sffamily\bfseries\LARGE] at ([xshift=5mm,yshift=-0.62cm]frame.north west) {Question 11 C};
\node[anchor=east,text=white,font=\sffamily\itshape\large] at ([xshift=-5mm,yshift=-0.62cm]frame.north east) {Relatifs et fractions};}]
\vspace*{1.75cm}
\begin{minipage}[t][6.9cm][c]{\linewidth}\LARGE\centering
\begin{minipage}{0.95\linewidth}\raggedright\linespread{1.25}\selectfont Que vaut $\dfrac{4}{9}\div\dfrac{2}{15}$~?\end{minipage}
\end{minipage}

\vfill
\begin{tcolorbox}[enhanced,nobeforeafter,width=7.1cm,height=2.1cm,valign=center,colback=crel!6,colframe=crel,arc=2.5mm,boxrule=0.9pt,left=3mm]{\Large\bfseries\color{crel}a)}\quad{\Large $\dfrac{8}{135}$}\end{tcolorbox}\hfill\begin{tcolorbox}[enhanced,nobeforeafter,width=7.1cm,height=2.1cm,valign=center,colback=crel!6,colframe=crel,arc=2.5mm,boxrule=0.9pt,left=3mm]{\Large\bfseries\color{crel}b)}\quad{\Large $-\dfrac{10}{3}$}\end{tcolorbox}\par\vspace{2.5mm}
\begin{tcolorbox}[enhanced,nobeforeafter,width=7.1cm,height=2.1cm,valign=center,colback=crel!6,colframe=crel,arc=2.5mm,boxrule=0.9pt,left=3mm]{\Large\bfseries\color{crel}c)}\quad{\Large $\dfrac{3}{10}$}\end{tcolorbox}\hfill\begin{tcolorbox}[enhanced,nobeforeafter,width=7.1cm,height=2.1cm,valign=center,colback=crel!6,colframe=crel,arc=2.5mm,boxrule=0.9pt,left=3mm]{\Large\bfseries\color{crel}d)}\quad{\Large $\dfrac{10}{3}$}\end{tcolorbox}
\end{tcolorbox}
\newpage
\begin{tcolorbox}[enhanced,width=16cm,height=14cm,colback=white,colframe=crel,boxrule=2pt,arc=0mm,
left=6mm,right=6mm,top=0mm,bottom=4mm,sharp corners,
overlay={\fill[crel] (frame.north west) rectangle ([yshift=-1.25cm]frame.north east);
\node[anchor=west,text=white,font=\sffamily\bfseries\LARGE] at ([xshift=5mm,yshift=-0.62cm]frame.north west) {Question 12 C};
\node[anchor=east,text=white,font=\sffamily\itshape\large] at ([xshift=-5mm,yshift=-0.62cm]frame.north east) {Relatifs et fractions};}]
\vspace*{1.75cm}
\begin{minipage}[t][6.9cm][c]{\linewidth}\LARGE\centering
\begin{minipage}{0.95\linewidth}\raggedright\linespread{1.25}\selectfont Que vaut $\left(-\dfrac{3}{7}\right)\div\left(-\dfrac{3}{7}\right)$~?\end{minipage}
\end{minipage}

\vfill
\begin{tcolorbox}[enhanced,nobeforeafter,width=7.1cm,height=2.1cm,valign=center,colback=crel!6,colframe=crel,arc=2.5mm,boxrule=0.9pt,left=3mm]{\Large\bfseries\color{crel}a)}\quad{\Large $1$}\end{tcolorbox}\hfill\begin{tcolorbox}[enhanced,nobeforeafter,width=7.1cm,height=2.1cm,valign=center,colback=crel!6,colframe=crel,arc=2.5mm,boxrule=0.9pt,left=3mm]{\Large\bfseries\color{crel}b)}\quad{\Large $2$}\end{tcolorbox}\par\vspace{2.5mm}
\begin{tcolorbox}[enhanced,nobeforeafter,width=7.1cm,height=2.1cm,valign=center,colback=crel!6,colframe=crel,arc=2.5mm,boxrule=0.9pt,left=3mm]{\Large\bfseries\color{crel}c)}\quad{\Large $-1$}\end{tcolorbox}\hfill\begin{tcolorbox}[enhanced,nobeforeafter,width=7.1cm,height=2.1cm,valign=center,colback=crel!6,colframe=crel,arc=2.5mm,boxrule=0.9pt,left=3mm]{\Large\bfseries\color{crel}d)}\quad{\Large $\dfrac{9}{49}$}\end{tcolorbox}
\end{tcolorbox}
\newpage
\begin{tcolorbox}[enhanced,width=16cm,height=14cm,colback=white,colframe=crel,boxrule=2pt,arc=0mm,
left=6mm,right=6mm,top=0mm,bottom=4mm,sharp corners,
overlay={\fill[crel] (frame.north west) rectangle ([yshift=-1.25cm]frame.north east);
\node[anchor=west,text=white,font=\sffamily\bfseries\LARGE] at ([xshift=5mm,yshift=-0.62cm]frame.north west) {Question 13 C};
\node[anchor=east,text=white,font=\sffamily\itshape\large] at ([xshift=-5mm,yshift=-0.62cm]frame.north east) {Relatifs et fractions};}]
\vspace*{1.75cm}
\begin{minipage}[t][6.9cm][c]{\linewidth}\LARGE\centering
\begin{minipage}{0.95\linewidth}\raggedright\linespread{1.25}\selectfont Quel est l'inverse de $-4$~?\end{minipage}
\end{minipage}

\vfill
\begin{tcolorbox}[enhanced,nobeforeafter,width=7.1cm,height=2.1cm,valign=center,colback=crel!6,colframe=crel,arc=2.5mm,boxrule=0.9pt,left=3mm]{\Large\bfseries\color{crel}a)}\quad{\Large $-4$}\end{tcolorbox}\hfill\begin{tcolorbox}[enhanced,nobeforeafter,width=7.1cm,height=2.1cm,valign=center,colback=crel!6,colframe=crel,arc=2.5mm,boxrule=0.9pt,left=3mm]{\Large\bfseries\color{crel}b)}\quad{\Large $4$}\end{tcolorbox}\par\vspace{2.5mm}
\begin{tcolorbox}[enhanced,nobeforeafter,width=7.1cm,height=2.1cm,valign=center,colback=crel!6,colframe=crel,arc=2.5mm,boxrule=0.9pt,left=3mm]{\Large\bfseries\color{crel}c)}\quad{\Large $\dfrac{1}{4}$}\end{tcolorbox}\hfill\begin{tcolorbox}[enhanced,nobeforeafter,width=7.1cm,height=2.1cm,valign=center,colback=crel!6,colframe=crel,arc=2.5mm,boxrule=0.9pt,left=3mm]{\Large\bfseries\color{crel}d)}\quad{\Large $-\dfrac{1}{4}$}\end{tcolorbox}
\end{tcolorbox}
\newpage
\begin{tcolorbox}[enhanced,width=16cm,height=14cm,colback=white,colframe=crel,boxrule=2pt,arc=0mm,
left=6mm,right=6mm,top=0mm,bottom=4mm,sharp corners,
overlay={\fill[crel] (frame.north west) rectangle ([yshift=-1.25cm]frame.north east);
\node[anchor=west,text=white,font=\sffamily\bfseries\LARGE] at ([xshift=5mm,yshift=-0.62cm]frame.north west) {Question 14 C};
\node[anchor=east,text=white,font=\sffamily\itshape\large] at ([xshift=-5mm,yshift=-0.62cm]frame.north east) {Relatifs et fractions};}]
\vspace*{1.75cm}
\begin{minipage}[t][6.9cm][c]{\linewidth}\LARGE\centering
\begin{minipage}{0.95\linewidth}\raggedright\linespread{1.25}\selectfont Quel est l'inverse de $\dfrac{5}{3}$~?\end{minipage}
\end{minipage}

\vfill
\begin{tcolorbox}[enhanced,nobeforeafter,width=7.1cm,height=2.1cm,valign=center,colback=crel!6,colframe=crel,arc=2.5mm,boxrule=0.9pt,left=3mm]{\Large\bfseries\color{crel}a)}\quad{\Large $-\dfrac{5}{3}$}\end{tcolorbox}\hfill\begin{tcolorbox}[enhanced,nobeforeafter,width=7.1cm,height=2.1cm,valign=center,colback=crel!6,colframe=crel,arc=2.5mm,boxrule=0.9pt,left=3mm]{\Large\bfseries\color{crel}b)}\quad{\Large $-\dfrac{3}{5}$}\end{tcolorbox}\par\vspace{2.5mm}
\begin{tcolorbox}[enhanced,nobeforeafter,width=7.1cm,height=2.1cm,valign=center,colback=crel!6,colframe=crel,arc=2.5mm,boxrule=0.9pt,left=3mm]{\Large\bfseries\color{crel}c)}\quad{\Large $\dfrac{3}{5}$}\end{tcolorbox}\hfill\begin{tcolorbox}[enhanced,nobeforeafter,width=7.1cm,height=2.1cm,valign=center,colback=crel!6,colframe=crel,arc=2.5mm,boxrule=0.9pt,left=3mm]{\Large\bfseries\color{crel}d)}\quad{\Large $\dfrac{5}{3}$}\end{tcolorbox}
\end{tcolorbox}
\newpage
\begin{tcolorbox}[enhanced,width=16cm,height=14cm,colback=white,colframe=crel,boxrule=2pt,arc=0mm,
left=6mm,right=6mm,top=0mm,bottom=4mm,sharp corners,
overlay={\fill[crel] (frame.north west) rectangle ([yshift=-1.25cm]frame.north east);
\node[anchor=west,text=white,font=\sffamily\bfseries\LARGE] at ([xshift=5mm,yshift=-0.62cm]frame.north west) {Question 15 C};
\node[anchor=east,text=white,font=\sffamily\itshape\large] at ([xshift=-5mm,yshift=-0.62cm]frame.north east) {Relatifs et fractions};}]
\vspace*{1.75cm}
\begin{minipage}[t][6.9cm][c]{\linewidth}\LARGE\centering
\begin{minipage}{0.95\linewidth}\raggedright\linespread{1.25}\selectfont Que vaut $\dfrac{1}{3}-\dfrac{5}{6}$~?\end{minipage}
\end{minipage}

\vfill
\begin{tcolorbox}[enhanced,nobeforeafter,width=7.1cm,height=2.1cm,valign=center,colback=crel!6,colframe=crel,arc=2.5mm,boxrule=0.9pt,left=3mm]{\Large\bfseries\color{crel}a)}\quad{\Large $-\dfrac{4}{3}$}\end{tcolorbox}\hfill\begin{tcolorbox}[enhanced,nobeforeafter,width=7.1cm,height=2.1cm,valign=center,colback=crel!6,colframe=crel,arc=2.5mm,boxrule=0.9pt,left=3mm]{\Large\bfseries\color{crel}b)}\quad{\Large $-\dfrac{1}{2}$}\end{tcolorbox}\par\vspace{2.5mm}
\begin{tcolorbox}[enhanced,nobeforeafter,width=7.1cm,height=2.1cm,valign=center,colback=crel!6,colframe=crel,arc=2.5mm,boxrule=0.9pt,left=3mm]{\Large\bfseries\color{crel}c)}\quad{\Large $-\dfrac{2}{9}$}\end{tcolorbox}\hfill\begin{tcolorbox}[enhanced,nobeforeafter,width=7.1cm,height=2.1cm,valign=center,colback=crel!6,colframe=crel,arc=2.5mm,boxrule=0.9pt,left=3mm]{\Large\bfseries\color{crel}d)}\quad{\Large $\dfrac{1}{2}$}\end{tcolorbox}
\end{tcolorbox}
\newpage
\begin{tcolorbox}[enhanced,width=16cm,height=14cm,colback=white,colframe=crel,boxrule=2pt,arc=0mm,
left=6mm,right=6mm,top=0mm,bottom=4mm,sharp corners,
overlay={\fill[crel] (frame.north west) rectangle ([yshift=-1.25cm]frame.north east);
\node[anchor=west,text=white,font=\sffamily\bfseries\LARGE] at ([xshift=5mm,yshift=-0.62cm]frame.north west) {Question 16 C};
\node[anchor=east,text=white,font=\sffamily\itshape\large] at ([xshift=-5mm,yshift=-0.62cm]frame.north east) {Relatifs et fractions};}]
\vspace*{1.75cm}
\begin{minipage}[t][6.9cm][c]{\linewidth}\LARGE\centering
\begin{minipage}{0.95\linewidth}\raggedright\linespread{1.25}\selectfont Que vaut $\dfrac{7}{8}-\dfrac{5}{6}$~?\end{minipage}
\end{minipage}

\vfill
\begin{tcolorbox}[enhanced,nobeforeafter,width=7.1cm,height=2.1cm,valign=center,colback=crel!6,colframe=crel,arc=2.5mm,boxrule=0.9pt,left=3mm]{\Large\bfseries\color{crel}a)}\quad{\Large $\dfrac{1}{4}$}\end{tcolorbox}\hfill\begin{tcolorbox}[enhanced,nobeforeafter,width=7.1cm,height=2.1cm,valign=center,colback=crel!6,colframe=crel,arc=2.5mm,boxrule=0.9pt,left=3mm]{\Large\bfseries\color{crel}b)}\quad{\Large $\dfrac{1}{24}$}\end{tcolorbox}\par\vspace{2.5mm}
\begin{tcolorbox}[enhanced,nobeforeafter,width=7.1cm,height=2.1cm,valign=center,colback=crel!6,colframe=crel,arc=2.5mm,boxrule=0.9pt,left=3mm]{\Large\bfseries\color{crel}c)}\quad{\Large $-\dfrac{1}{24}$}\end{tcolorbox}\hfill\begin{tcolorbox}[enhanced,nobeforeafter,width=7.1cm,height=2.1cm,valign=center,colback=crel!6,colframe=crel,arc=2.5mm,boxrule=0.9pt,left=3mm]{\Large\bfseries\color{crel}d)}\quad{\Large $1$}\end{tcolorbox}
\end{tcolorbox}
\newpage
\begin{tcolorbox}[enhanced,width=16cm,height=14cm,colback=white,colframe=cpui,boxrule=2pt,arc=0mm,
left=6mm,right=6mm,top=0mm,bottom=4mm,sharp corners,
overlay={\fill[cpui] (frame.north west) rectangle ([yshift=-1.25cm]frame.north east);
\node[anchor=west,text=white,font=\sffamily\bfseries\LARGE] at ([xshift=5mm,yshift=-0.62cm]frame.north west) {Question 17 C};
\node[anchor=east,text=white,font=\sffamily\itshape\large] at ([xshift=-5mm,yshift=-0.62cm]frame.north east) {Puissances et racines carrées};}]
\vspace*{1.75cm}
\begin{minipage}[t][6.9cm][c]{\linewidth}\LARGE\centering
\begin{minipage}{0.95\linewidth}\raggedright\linespread{1.25}\selectfont Que vaut $(-3)^{2}$~?\end{minipage}
\end{minipage}

\vfill
\begin{tcolorbox}[enhanced,nobeforeafter,width=7.1cm,height=2.1cm,valign=center,colback=cpui!6,colframe=cpui,arc=2.5mm,boxrule=0.9pt,left=3mm]{\Large\bfseries\color{cpui}a)}\quad{\Large $-3$}\end{tcolorbox}\hfill\begin{tcolorbox}[enhanced,nobeforeafter,width=7.1cm,height=2.1cm,valign=center,colback=cpui!6,colframe=cpui,arc=2.5mm,boxrule=0.9pt,left=3mm]{\Large\bfseries\color{cpui}b)}\quad{\Large $9$}\end{tcolorbox}\par\vspace{2.5mm}
\begin{tcolorbox}[enhanced,nobeforeafter,width=7.1cm,height=2.1cm,valign=center,colback=cpui!6,colframe=cpui,arc=2.5mm,boxrule=0.9pt,left=3mm]{\Large\bfseries\color{cpui}c)}\quad{\Large $-9$}\end{tcolorbox}\hfill\begin{tcolorbox}[enhanced,nobeforeafter,width=7.1cm,height=2.1cm,valign=center,colback=cpui!6,colframe=cpui,arc=2.5mm,boxrule=0.9pt,left=3mm]{\Large\bfseries\color{cpui}d)}\quad{\Large $-6$}\end{tcolorbox}
\end{tcolorbox}
\newpage
\begin{tcolorbox}[enhanced,width=16cm,height=14cm,colback=white,colframe=cpui,boxrule=2pt,arc=0mm,
left=6mm,right=6mm,top=0mm,bottom=4mm,sharp corners,
overlay={\fill[cpui] (frame.north west) rectangle ([yshift=-1.25cm]frame.north east);
\node[anchor=west,text=white,font=\sffamily\bfseries\LARGE] at ([xshift=5mm,yshift=-0.62cm]frame.north west) {Question 18 C};
\node[anchor=east,text=white,font=\sffamily\itshape\large] at ([xshift=-5mm,yshift=-0.62cm]frame.north east) {Puissances et racines carrées};}]
\vspace*{1.75cm}
\begin{minipage}[t][6.9cm][c]{\linewidth}\LARGE\centering
\begin{minipage}{0.95\linewidth}\raggedright\linespread{1.25}\selectfont Que vaut $2^{7}$~?\end{minipage}
\end{minipage}

\vfill
\begin{tcolorbox}[enhanced,nobeforeafter,width=7.1cm,height=2.1cm,valign=center,colback=cpui!6,colframe=cpui,arc=2.5mm,boxrule=0.9pt,left=3mm]{\Large\bfseries\color{cpui}a)}\quad{\Large $128$}\end{tcolorbox}\hfill\begin{tcolorbox}[enhanced,nobeforeafter,width=7.1cm,height=2.1cm,valign=center,colback=cpui!6,colframe=cpui,arc=2.5mm,boxrule=0.9pt,left=3mm]{\Large\bfseries\color{cpui}b)}\quad{\Large $14$}\end{tcolorbox}\par\vspace{2.5mm}
\begin{tcolorbox}[enhanced,nobeforeafter,width=7.1cm,height=2.1cm,valign=center,colback=cpui!6,colframe=cpui,arc=2.5mm,boxrule=0.9pt,left=3mm]{\Large\bfseries\color{cpui}c)}\quad{\Large $64$}\end{tcolorbox}\hfill\begin{tcolorbox}[enhanced,nobeforeafter,width=7.1cm,height=2.1cm,valign=center,colback=cpui!6,colframe=cpui,arc=2.5mm,boxrule=0.9pt,left=3mm]{\Large\bfseries\color{cpui}d)}\quad{\Large $-128$}\end{tcolorbox}
\end{tcolorbox}
\newpage
\begin{tcolorbox}[enhanced,width=16cm,height=14cm,colback=white,colframe=cpui,boxrule=2pt,arc=0mm,
left=6mm,right=6mm,top=0mm,bottom=4mm,sharp corners,
overlay={\fill[cpui] (frame.north west) rectangle ([yshift=-1.25cm]frame.north east);
\node[anchor=west,text=white,font=\sffamily\bfseries\LARGE] at ([xshift=5mm,yshift=-0.62cm]frame.north west) {Question 19 C};
\node[anchor=east,text=white,font=\sffamily\itshape\large] at ([xshift=-5mm,yshift=-0.62cm]frame.north east) {Puissances et racines carrées};}]
\vspace*{1.75cm}
\begin{minipage}[t][6.9cm][c]{\linewidth}\LARGE\centering
\begin{minipage}{0.95\linewidth}\raggedright\linespread{1.25}\selectfont Écrire $10^{-5}\times10^{8}$ sous la forme d'une seule puissance de $10$.\end{minipage}
\end{minipage}

\vfill
\begin{tcolorbox}[enhanced,nobeforeafter,width=7.1cm,height=2.1cm,valign=center,colback=cpui!6,colframe=cpui,arc=2.5mm,boxrule=0.9pt,left=3mm]{\Large\bfseries\color{cpui}a)}\quad{\Large $10^{-40}$}\end{tcolorbox}\hfill\begin{tcolorbox}[enhanced,nobeforeafter,width=7.1cm,height=2.1cm,valign=center,colback=cpui!6,colframe=cpui,arc=2.5mm,boxrule=0.9pt,left=3mm]{\Large\bfseries\color{cpui}b)}\quad{\Large $10^{3}$}\end{tcolorbox}\par\vspace{2.5mm}
\begin{tcolorbox}[enhanced,nobeforeafter,width=7.1cm,height=2.1cm,valign=center,colback=cpui!6,colframe=cpui,arc=2.5mm,boxrule=0.9pt,left=3mm]{\Large\bfseries\color{cpui}c)}\quad{\Large $100^{3}$}\end{tcolorbox}\hfill\begin{tcolorbox}[enhanced,nobeforeafter,width=7.1cm,height=2.1cm,valign=center,colback=cpui!6,colframe=cpui,arc=2.5mm,boxrule=0.9pt,left=3mm]{\Large\bfseries\color{cpui}d)}\quad{\Large $10^{-13}$}\end{tcolorbox}
\end{tcolorbox}
\newpage
\begin{tcolorbox}[enhanced,width=16cm,height=14cm,colback=white,colframe=cpui,boxrule=2pt,arc=0mm,
left=6mm,right=6mm,top=0mm,bottom=4mm,sharp corners,
overlay={\fill[cpui] (frame.north west) rectangle ([yshift=-1.25cm]frame.north east);
\node[anchor=west,text=white,font=\sffamily\bfseries\LARGE] at ([xshift=5mm,yshift=-0.62cm]frame.north west) {Question 20 C};
\node[anchor=east,text=white,font=\sffamily\itshape\large] at ([xshift=-5mm,yshift=-0.62cm]frame.north east) {Puissances et racines carrées};}]
\vspace*{1.75cm}
\begin{minipage}[t][6.9cm][c]{\linewidth}\LARGE\centering
\begin{minipage}{0.95\linewidth}\raggedright\linespread{1.25}\selectfont Écrire $10^{2}\times10^{-7}$ sous la forme d'une seule puissance de $10$.\end{minipage}
\end{minipage}

\vfill
\begin{tcolorbox}[enhanced,nobeforeafter,width=7.1cm,height=2.1cm,valign=center,colback=cpui!6,colframe=cpui,arc=2.5mm,boxrule=0.9pt,left=3mm]{\Large\bfseries\color{cpui}a)}\quad{\Large $10^{9}$}\end{tcolorbox}\hfill\begin{tcolorbox}[enhanced,nobeforeafter,width=7.1cm,height=2.1cm,valign=center,colback=cpui!6,colframe=cpui,arc=2.5mm,boxrule=0.9pt,left=3mm]{\Large\bfseries\color{cpui}b)}\quad{\Large $10^{-14}$}\end{tcolorbox}\par\vspace{2.5mm}
\begin{tcolorbox}[enhanced,nobeforeafter,width=7.1cm,height=2.1cm,valign=center,colback=cpui!6,colframe=cpui,arc=2.5mm,boxrule=0.9pt,left=3mm]{\Large\bfseries\color{cpui}c)}\quad{\Large $10^{-5}$}\end{tcolorbox}\hfill\begin{tcolorbox}[enhanced,nobeforeafter,width=7.1cm,height=2.1cm,valign=center,colback=cpui!6,colframe=cpui,arc=2.5mm,boxrule=0.9pt,left=3mm]{\Large\bfseries\color{cpui}d)}\quad{\Large $100^{-5}$}\end{tcolorbox}
\end{tcolorbox}
\newpage
\begin{tcolorbox}[enhanced,width=16cm,height=14cm,colback=white,colframe=cpui,boxrule=2pt,arc=0mm,
left=6mm,right=6mm,top=0mm,bottom=4mm,sharp corners,
overlay={\fill[cpui] (frame.north west) rectangle ([yshift=-1.25cm]frame.north east);
\node[anchor=west,text=white,font=\sffamily\bfseries\LARGE] at ([xshift=5mm,yshift=-0.62cm]frame.north west) {Question 21 C};
\node[anchor=east,text=white,font=\sffamily\itshape\large] at ([xshift=-5mm,yshift=-0.62cm]frame.north east) {Puissances et racines carrées};}]
\vspace*{1.75cm}
\begin{minipage}[t][6.9cm][c]{\linewidth}\LARGE\centering
\begin{minipage}{0.95\linewidth}\raggedright\linespread{1.25}\selectfont Quelle est l'écriture décimale de $10^{-5}$~?\end{minipage}
\end{minipage}

\vfill
\begin{tcolorbox}[enhanced,nobeforeafter,width=7.1cm,height=2.1cm,valign=center,colback=cpui!6,colframe=cpui,arc=2.5mm,boxrule=0.9pt,left=3mm]{\Large\bfseries\color{cpui}a)}\quad{\Large $0{,}000001$}\end{tcolorbox}\hfill\begin{tcolorbox}[enhanced,nobeforeafter,width=7.1cm,height=2.1cm,valign=center,colback=cpui!6,colframe=cpui,arc=2.5mm,boxrule=0.9pt,left=3mm]{\Large\bfseries\color{cpui}b)}\quad{\Large $0{,}00001$}\end{tcolorbox}\par\vspace{2.5mm}
\begin{tcolorbox}[enhanced,nobeforeafter,width=7.1cm,height=2.1cm,valign=center,colback=cpui!6,colframe=cpui,arc=2.5mm,boxrule=0.9pt,left=3mm]{\Large\bfseries\color{cpui}c)}\quad{\Large $0{,}0001$}\end{tcolorbox}\hfill\begin{tcolorbox}[enhanced,nobeforeafter,width=7.1cm,height=2.1cm,valign=center,colback=cpui!6,colframe=cpui,arc=2.5mm,boxrule=0.9pt,left=3mm]{\Large\bfseries\color{cpui}d)}\quad{\Large $-100\,000$}\end{tcolorbox}
\end{tcolorbox}
\newpage
\begin{tcolorbox}[enhanced,width=16cm,height=14cm,colback=white,colframe=cpui,boxrule=2pt,arc=0mm,
left=6mm,right=6mm,top=0mm,bottom=4mm,sharp corners,
overlay={\fill[cpui] (frame.north west) rectangle ([yshift=-1.25cm]frame.north east);
\node[anchor=west,text=white,font=\sffamily\bfseries\LARGE] at ([xshift=5mm,yshift=-0.62cm]frame.north west) {Question 22 C};
\node[anchor=east,text=white,font=\sffamily\itshape\large] at ([xshift=-5mm,yshift=-0.62cm]frame.north east) {Puissances et racines carrées};}]
\vspace*{1.75cm}
\begin{minipage}[t][6.9cm][c]{\linewidth}\LARGE\centering
\begin{minipage}{0.95\linewidth}\raggedright\linespread{1.25}\selectfont Quelle est l'écriture décimale de $10^{-6}$~?\end{minipage}
\end{minipage}

\vfill
\begin{tcolorbox}[enhanced,nobeforeafter,width=7.1cm,height=2.1cm,valign=center,colback=cpui!6,colframe=cpui,arc=2.5mm,boxrule=0.9pt,left=3mm]{\Large\bfseries\color{cpui}a)}\quad{\Large $0{,}000001$}\end{tcolorbox}\hfill\begin{tcolorbox}[enhanced,nobeforeafter,width=7.1cm,height=2.1cm,valign=center,colback=cpui!6,colframe=cpui,arc=2.5mm,boxrule=0.9pt,left=3mm]{\Large\bfseries\color{cpui}b)}\quad{\Large $0{,}00001$}\end{tcolorbox}\par\vspace{2.5mm}
\begin{tcolorbox}[enhanced,nobeforeafter,width=7.1cm,height=2.1cm,valign=center,colback=cpui!6,colframe=cpui,arc=2.5mm,boxrule=0.9pt,left=3mm]{\Large\bfseries\color{cpui}c)}\quad{\Large $-1\,000\,000$}\end{tcolorbox}\hfill\begin{tcolorbox}[enhanced,nobeforeafter,width=7.1cm,height=2.1cm,valign=center,colback=cpui!6,colframe=cpui,arc=2.5mm,boxrule=0.9pt,left=3mm]{\Large\bfseries\color{cpui}d)}\quad{\Large $0{,}0000001$}\end{tcolorbox}
\end{tcolorbox}
\newpage
\begin{tcolorbox}[enhanced,width=16cm,height=14cm,colback=white,colframe=cpui,boxrule=2pt,arc=0mm,
left=6mm,right=6mm,top=0mm,bottom=4mm,sharp corners,
overlay={\fill[cpui] (frame.north west) rectangle ([yshift=-1.25cm]frame.north east);
\node[anchor=west,text=white,font=\sffamily\bfseries\LARGE] at ([xshift=5mm,yshift=-0.62cm]frame.north west) {Question 23 C};
\node[anchor=east,text=white,font=\sffamily\itshape\large] at ([xshift=-5mm,yshift=-0.62cm]frame.north east) {Puissances et racines carrées};}]
\vspace*{1.75cm}
\begin{minipage}[t][6.9cm][c]{\linewidth}\LARGE\centering
\begin{minipage}{0.95\linewidth}\raggedright\linespread{1.25}\selectfont Quelle est la notation scientifique de $605\,000\,000$~?\end{minipage}
\end{minipage}

\vfill
\begin{tcolorbox}[enhanced,nobeforeafter,width=7.1cm,height=2.1cm,valign=center,colback=cpui!6,colframe=cpui,arc=2.5mm,boxrule=0.9pt,left=3mm]{\Large\bfseries\color{cpui}a)}\quad{\Large $60{,}5\times10^{8}$}\end{tcolorbox}\hfill\begin{tcolorbox}[enhanced,nobeforeafter,width=7.1cm,height=2.1cm,valign=center,colback=cpui!6,colframe=cpui,arc=2.5mm,boxrule=0.9pt,left=3mm]{\Large\bfseries\color{cpui}b)}\quad{\Large $6{,}05\times10^{9}$}\end{tcolorbox}\par\vspace{2.5mm}
\begin{tcolorbox}[enhanced,nobeforeafter,width=7.1cm,height=2.1cm,valign=center,colback=cpui!6,colframe=cpui,arc=2.5mm,boxrule=0.9pt,left=3mm]{\Large\bfseries\color{cpui}c)}\quad{\Large $6{,}05\times10^{8}$}\end{tcolorbox}\hfill\begin{tcolorbox}[enhanced,nobeforeafter,width=7.1cm,height=2.1cm,valign=center,colback=cpui!6,colframe=cpui,arc=2.5mm,boxrule=0.9pt,left=3mm]{\Large\bfseries\color{cpui}d)}\quad{\Large $6{,}05\times10^{-8}$}\end{tcolorbox}
\end{tcolorbox}
\newpage
\begin{tcolorbox}[enhanced,width=16cm,height=14cm,colback=white,colframe=cpui,boxrule=2pt,arc=0mm,
left=6mm,right=6mm,top=0mm,bottom=4mm,sharp corners,
overlay={\fill[cpui] (frame.north west) rectangle ([yshift=-1.25cm]frame.north east);
\node[anchor=west,text=white,font=\sffamily\bfseries\LARGE] at ([xshift=5mm,yshift=-0.62cm]frame.north west) {Question 24 C};
\node[anchor=east,text=white,font=\sffamily\itshape\large] at ([xshift=-5mm,yshift=-0.62cm]frame.north east) {Puissances et racines carrées};}]
\vspace*{1.75cm}
\begin{minipage}[t][6.9cm][c]{\linewidth}\LARGE\centering
\begin{minipage}{0.95\linewidth}\raggedright\linespread{1.25}\selectfont Quelle est la notation scientifique de $0{,}0031$~?\end{minipage}
\end{minipage}

\vfill
\begin{tcolorbox}[enhanced,nobeforeafter,width=7.1cm,height=2.1cm,valign=center,colback=cpui!6,colframe=cpui,arc=2.5mm,boxrule=0.9pt,left=3mm]{\Large\bfseries\color{cpui}a)}\quad{\Large $31\times10^{-3}$}\end{tcolorbox}\hfill\begin{tcolorbox}[enhanced,nobeforeafter,width=7.1cm,height=2.1cm,valign=center,colback=cpui!6,colframe=cpui,arc=2.5mm,boxrule=0.9pt,left=3mm]{\Large\bfseries\color{cpui}b)}\quad{\Large $3{,}1\times10^{3}$}\end{tcolorbox}\par\vspace{2.5mm}
\begin{tcolorbox}[enhanced,nobeforeafter,width=7.1cm,height=2.1cm,valign=center,colback=cpui!6,colframe=cpui,arc=2.5mm,boxrule=0.9pt,left=3mm]{\Large\bfseries\color{cpui}c)}\quad{\Large $3{,}1\times10^{-3}$}\end{tcolorbox}\hfill\begin{tcolorbox}[enhanced,nobeforeafter,width=7.1cm,height=2.1cm,valign=center,colback=cpui!6,colframe=cpui,arc=2.5mm,boxrule=0.9pt,left=3mm]{\Large\bfseries\color{cpui}d)}\quad{\Large $3{,}1\times10^{-2}$}\end{tcolorbox}
\end{tcolorbox}
\newpage
\begin{tcolorbox}[enhanced,width=16cm,height=14cm,colback=white,colframe=cpui,boxrule=2pt,arc=0mm,
left=6mm,right=6mm,top=0mm,bottom=4mm,sharp corners,
overlay={\fill[cpui] (frame.north west) rectangle ([yshift=-1.25cm]frame.north east);
\node[anchor=west,text=white,font=\sffamily\bfseries\LARGE] at ([xshift=5mm,yshift=-0.62cm]frame.north west) {Question 25 C};
\node[anchor=east,text=white,font=\sffamily\itshape\large] at ([xshift=-5mm,yshift=-0.62cm]frame.north east) {Puissances et racines carrées};}]
\vspace*{1.75cm}
\begin{minipage}[t][6.9cm][c]{\linewidth}\LARGE\centering
\begin{minipage}{0.95\linewidth}\raggedright\linespread{1.25}\selectfont Que vaut $-3^{4}$~?\end{minipage}
\end{minipage}

\vfill
\begin{tcolorbox}[enhanced,nobeforeafter,width=7.1cm,height=2.1cm,valign=center,colback=cpui!6,colframe=cpui,arc=2.5mm,boxrule=0.9pt,left=3mm]{\Large\bfseries\color{cpui}a)}\quad{\Large $-81$}\end{tcolorbox}\hfill\begin{tcolorbox}[enhanced,nobeforeafter,width=7.1cm,height=2.1cm,valign=center,colback=cpui!6,colframe=cpui,arc=2.5mm,boxrule=0.9pt,left=3mm]{\Large\bfseries\color{cpui}b)}\quad{\Large $-12$}\end{tcolorbox}\par\vspace{2.5mm}
\begin{tcolorbox}[enhanced,nobeforeafter,width=7.1cm,height=2.1cm,valign=center,colback=cpui!6,colframe=cpui,arc=2.5mm,boxrule=0.9pt,left=3mm]{\Large\bfseries\color{cpui}c)}\quad{\Large $12$}\end{tcolorbox}\hfill\begin{tcolorbox}[enhanced,nobeforeafter,width=7.1cm,height=2.1cm,valign=center,colback=cpui!6,colframe=cpui,arc=2.5mm,boxrule=0.9pt,left=3mm]{\Large\bfseries\color{cpui}d)}\quad{\Large $81$}\end{tcolorbox}
\end{tcolorbox}
\newpage
\begin{tcolorbox}[enhanced,width=16cm,height=14cm,colback=white,colframe=cpui,boxrule=2pt,arc=0mm,
left=6mm,right=6mm,top=0mm,bottom=4mm,sharp corners,
overlay={\fill[cpui] (frame.north west) rectangle ([yshift=-1.25cm]frame.north east);
\node[anchor=west,text=white,font=\sffamily\bfseries\LARGE] at ([xshift=5mm,yshift=-0.62cm]frame.north west) {Question 26 C};
\node[anchor=east,text=white,font=\sffamily\itshape\large] at ([xshift=-5mm,yshift=-0.62cm]frame.north east) {Puissances et racines carrées};}]
\vspace*{1.75cm}
\begin{minipage}[t][6.9cm][c]{\linewidth}\LARGE\centering
\begin{minipage}{0.95\linewidth}\raggedright\linespread{1.25}\selectfont Que vaut $-10^{2}$~?\end{minipage}
\end{minipage}

\vfill
\begin{tcolorbox}[enhanced,nobeforeafter,width=7.1cm,height=2.1cm,valign=center,colback=cpui!6,colframe=cpui,arc=2.5mm,boxrule=0.9pt,left=3mm]{\Large\bfseries\color{cpui}a)}\quad{\Large $-20$}\end{tcolorbox}\hfill\begin{tcolorbox}[enhanced,nobeforeafter,width=7.1cm,height=2.1cm,valign=center,colback=cpui!6,colframe=cpui,arc=2.5mm,boxrule=0.9pt,left=3mm]{\Large\bfseries\color{cpui}b)}\quad{\Large $20$}\end{tcolorbox}\par\vspace{2.5mm}
\begin{tcolorbox}[enhanced,nobeforeafter,width=7.1cm,height=2.1cm,valign=center,colback=cpui!6,colframe=cpui,arc=2.5mm,boxrule=0.9pt,left=3mm]{\Large\bfseries\color{cpui}c)}\quad{\Large $100$}\end{tcolorbox}\hfill\begin{tcolorbox}[enhanced,nobeforeafter,width=7.1cm,height=2.1cm,valign=center,colback=cpui!6,colframe=cpui,arc=2.5mm,boxrule=0.9pt,left=3mm]{\Large\bfseries\color{cpui}d)}\quad{\Large $-100$}\end{tcolorbox}
\end{tcolorbox}
\newpage
\begin{tcolorbox}[enhanced,width=16cm,height=14cm,colback=white,colframe=cpui,boxrule=2pt,arc=0mm,
left=6mm,right=6mm,top=0mm,bottom=4mm,sharp corners,
overlay={\fill[cpui] (frame.north west) rectangle ([yshift=-1.25cm]frame.north east);
\node[anchor=west,text=white,font=\sffamily\bfseries\LARGE] at ([xshift=5mm,yshift=-0.62cm]frame.north west) {Question 27 C};
\node[anchor=east,text=white,font=\sffamily\itshape\large] at ([xshift=-5mm,yshift=-0.62cm]frame.north east) {Puissances et racines carrées};}]
\vspace*{1.75cm}
\begin{minipage}[t][6.9cm][c]{\linewidth}\LARGE\centering
\begin{minipage}{0.95\linewidth}\raggedright\linespread{1.25}\selectfont Écrire $\dfrac{10^{9}}{10^{2}}$ sous la forme d'une seule puissance de $10$.\end{minipage}
\end{minipage}

\vfill
\begin{tcolorbox}[enhanced,nobeforeafter,width=7.1cm,height=2.1cm,valign=center,colback=cpui!6,colframe=cpui,arc=2.5mm,boxrule=0.9pt,left=3mm]{\Large\bfseries\color{cpui}a)}\quad{\Large $10^{11}$}\end{tcolorbox}\hfill\begin{tcolorbox}[enhanced,nobeforeafter,width=7.1cm,height=2.1cm,valign=center,colback=cpui!6,colframe=cpui,arc=2.5mm,boxrule=0.9pt,left=3mm]{\Large\bfseries\color{cpui}b)}\quad{\Large $10^{18}$}\end{tcolorbox}\par\vspace{2.5mm}
\begin{tcolorbox}[enhanced,nobeforeafter,width=7.1cm,height=2.1cm,valign=center,colback=cpui!6,colframe=cpui,arc=2.5mm,boxrule=0.9pt,left=3mm]{\Large\bfseries\color{cpui}c)}\quad{\Large $10^{7}$}\end{tcolorbox}\hfill\begin{tcolorbox}[enhanced,nobeforeafter,width=7.1cm,height=2.1cm,valign=center,colback=cpui!6,colframe=cpui,arc=2.5mm,boxrule=0.9pt,left=3mm]{\Large\bfseries\color{cpui}d)}\quad{\Large $10^{-7}$}\end{tcolorbox}
\end{tcolorbox}
\newpage
\begin{tcolorbox}[enhanced,width=16cm,height=14cm,colback=white,colframe=cpui,boxrule=2pt,arc=0mm,
left=6mm,right=6mm,top=0mm,bottom=4mm,sharp corners,
overlay={\fill[cpui] (frame.north west) rectangle ([yshift=-1.25cm]frame.north east);
\node[anchor=west,text=white,font=\sffamily\bfseries\LARGE] at ([xshift=5mm,yshift=-0.62cm]frame.north west) {Question 28 C};
\node[anchor=east,text=white,font=\sffamily\itshape\large] at ([xshift=-5mm,yshift=-0.62cm]frame.north east) {Puissances et racines carrées};}]
\vspace*{1.75cm}
\begin{minipage}[t][6.9cm][c]{\linewidth}\LARGE\centering
\begin{minipage}{0.95\linewidth}\raggedright\linespread{1.25}\selectfont Écrire $\dfrac{10^{3}}{10^{-5}}$ sous la forme d'une seule puissance de $10$.\end{minipage}
\end{minipage}

\vfill
\begin{tcolorbox}[enhanced,nobeforeafter,width=7.1cm,height=2.1cm,valign=center,colback=cpui!6,colframe=cpui,arc=2.5mm,boxrule=0.9pt,left=3mm]{\Large\bfseries\color{cpui}a)}\quad{\Large $10^{-15}$}\end{tcolorbox}\hfill\begin{tcolorbox}[enhanced,nobeforeafter,width=7.1cm,height=2.1cm,valign=center,colback=cpui!6,colframe=cpui,arc=2.5mm,boxrule=0.9pt,left=3mm]{\Large\bfseries\color{cpui}b)}\quad{\Large $10^{-2}$}\end{tcolorbox}\par\vspace{2.5mm}
\begin{tcolorbox}[enhanced,nobeforeafter,width=7.1cm,height=2.1cm,valign=center,colback=cpui!6,colframe=cpui,arc=2.5mm,boxrule=0.9pt,left=3mm]{\Large\bfseries\color{cpui}c)}\quad{\Large $10^{8}$}\end{tcolorbox}\hfill\begin{tcolorbox}[enhanced,nobeforeafter,width=7.1cm,height=2.1cm,valign=center,colback=cpui!6,colframe=cpui,arc=2.5mm,boxrule=0.9pt,left=3mm]{\Large\bfseries\color{cpui}d)}\quad{\Large $10^{-8}$}\end{tcolorbox}
\end{tcolorbox}
\newpage
\begin{tcolorbox}[enhanced,width=16cm,height=14cm,colback=white,colframe=cpui,boxrule=2pt,arc=0mm,
left=6mm,right=6mm,top=0mm,bottom=4mm,sharp corners,
overlay={\fill[cpui] (frame.north west) rectangle ([yshift=-1.25cm]frame.north east);
\node[anchor=west,text=white,font=\sffamily\bfseries\LARGE] at ([xshift=5mm,yshift=-0.62cm]frame.north west) {Question 29 C};
\node[anchor=east,text=white,font=\sffamily\itshape\large] at ([xshift=-5mm,yshift=-0.62cm]frame.north east) {Puissances et racines carrées};}]
\vspace*{1.75cm}
\begin{minipage}[t][6.9cm][c]{\linewidth}\LARGE\centering
\begin{minipage}{0.95\linewidth}\raggedright\linespread{1.25}\selectfont Que vaut $\sqrt{169}$~?\end{minipage}
\end{minipage}

\vfill
\begin{tcolorbox}[enhanced,nobeforeafter,width=7.1cm,height=2.1cm,valign=center,colback=cpui!6,colframe=cpui,arc=2.5mm,boxrule=0.9pt,left=3mm]{\Large\bfseries\color{cpui}a)}\quad{\Large $26$}\end{tcolorbox}\hfill\begin{tcolorbox}[enhanced,nobeforeafter,width=7.1cm,height=2.1cm,valign=center,colback=cpui!6,colframe=cpui,arc=2.5mm,boxrule=0.9pt,left=3mm]{\Large\bfseries\color{cpui}b)}\quad{\Large $13$}\end{tcolorbox}\par\vspace{2.5mm}
\begin{tcolorbox}[enhanced,nobeforeafter,width=7.1cm,height=2.1cm,valign=center,colback=cpui!6,colframe=cpui,arc=2.5mm,boxrule=0.9pt,left=3mm]{\Large\bfseries\color{cpui}c)}\quad{\Large $84$}\end{tcolorbox}\hfill\begin{tcolorbox}[enhanced,nobeforeafter,width=7.1cm,height=2.1cm,valign=center,colback=cpui!6,colframe=cpui,arc=2.5mm,boxrule=0.9pt,left=3mm]{\Large\bfseries\color{cpui}d)}\quad{\Large $14$}\end{tcolorbox}
\end{tcolorbox}
\newpage
\begin{tcolorbox}[enhanced,width=16cm,height=14cm,colback=white,colframe=cpui,boxrule=2pt,arc=0mm,
left=6mm,right=6mm,top=0mm,bottom=4mm,sharp corners,
overlay={\fill[cpui] (frame.north west) rectangle ([yshift=-1.25cm]frame.north east);
\node[anchor=west,text=white,font=\sffamily\bfseries\LARGE] at ([xshift=5mm,yshift=-0.62cm]frame.north west) {Question 30 C};
\node[anchor=east,text=white,font=\sffamily\itshape\large] at ([xshift=-5mm,yshift=-0.62cm]frame.north east) {Puissances et racines carrées};}]
\vspace*{1.75cm}
\begin{minipage}[t][6.9cm][c]{\linewidth}\LARGE\centering
\begin{minipage}{0.95\linewidth}\raggedright\linespread{1.25}\selectfont Que vaut $\sqrt{225}$~?\end{minipage}
\end{minipage}

\vfill
\begin{tcolorbox}[enhanced,nobeforeafter,width=7.1cm,height=2.1cm,valign=center,colback=cpui!6,colframe=cpui,arc=2.5mm,boxrule=0.9pt,left=3mm]{\Large\bfseries\color{cpui}a)}\quad{\Large $112$}\end{tcolorbox}\hfill\begin{tcolorbox}[enhanced,nobeforeafter,width=7.1cm,height=2.1cm,valign=center,colback=cpui!6,colframe=cpui,arc=2.5mm,boxrule=0.9pt,left=3mm]{\Large\bfseries\color{cpui}b)}\quad{\Large $30$}\end{tcolorbox}\par\vspace{2.5mm}
\begin{tcolorbox}[enhanced,nobeforeafter,width=7.1cm,height=2.1cm,valign=center,colback=cpui!6,colframe=cpui,arc=2.5mm,boxrule=0.9pt,left=3mm]{\Large\bfseries\color{cpui}c)}\quad{\Large $15$}\end{tcolorbox}\hfill\begin{tcolorbox}[enhanced,nobeforeafter,width=7.1cm,height=2.1cm,valign=center,colback=cpui!6,colframe=cpui,arc=2.5mm,boxrule=0.9pt,left=3mm]{\Large\bfseries\color{cpui}d)}\quad{\Large $16$}\end{tcolorbox}
\end{tcolorbox}
\newpage
\begin{tcolorbox}[enhanced,width=16cm,height=14cm,colback=white,colframe=cpui,boxrule=2pt,arc=0mm,
left=6mm,right=6mm,top=0mm,bottom=4mm,sharp corners,
overlay={\fill[cpui] (frame.north west) rectangle ([yshift=-1.25cm]frame.north east);
\node[anchor=west,text=white,font=\sffamily\bfseries\LARGE] at ([xshift=5mm,yshift=-0.62cm]frame.north west) {Question 31 C};
\node[anchor=east,text=white,font=\sffamily\itshape\large] at ([xshift=-5mm,yshift=-0.62cm]frame.north east) {Puissances et racines carrées};}]
\vspace*{1.75cm}
\begin{minipage}[t][6.9cm][c]{\linewidth}\LARGE\centering
\begin{minipage}{0.95\linewidth}\raggedright\linespread{1.25}\selectfont Entre quels nombres entiers consécutifs se trouve $\sqrt{30}$~?\end{minipage}
\end{minipage}

\vfill
\begin{tcolorbox}[enhanced,nobeforeafter,width=7.1cm,height=2.1cm,valign=center,colback=cpui!6,colframe=cpui,arc=2.5mm,boxrule=0.9pt,left=3mm]{\Large\bfseries\color{cpui}a)}\quad{\Large $4\text{ et }5$}\end{tcolorbox}\hfill\begin{tcolorbox}[enhanced,nobeforeafter,width=7.1cm,height=2.1cm,valign=center,colback=cpui!6,colframe=cpui,arc=2.5mm,boxrule=0.9pt,left=3mm]{\Large\bfseries\color{cpui}b)}\quad{\Large $15\text{ et }16$}\end{tcolorbox}\par\vspace{2.5mm}
\begin{tcolorbox}[enhanced,nobeforeafter,width=7.1cm,height=2.1cm,valign=center,colback=cpui!6,colframe=cpui,arc=2.5mm,boxrule=0.9pt,left=3mm]{\Large\bfseries\color{cpui}c)}\quad{\Large $5\text{ et }6$}\end{tcolorbox}\hfill\begin{tcolorbox}[enhanced,nobeforeafter,width=7.1cm,height=2.1cm,valign=center,colback=cpui!6,colframe=cpui,arc=2.5mm,boxrule=0.9pt,left=3mm]{\Large\bfseries\color{cpui}d)}\quad{\Large $6\text{ et }7$}\end{tcolorbox}
\end{tcolorbox}
\newpage
\begin{tcolorbox}[enhanced,width=16cm,height=14cm,colback=white,colframe=cpui,boxrule=2pt,arc=0mm,
left=6mm,right=6mm,top=0mm,bottom=4mm,sharp corners,
overlay={\fill[cpui] (frame.north west) rectangle ([yshift=-1.25cm]frame.north east);
\node[anchor=west,text=white,font=\sffamily\bfseries\LARGE] at ([xshift=5mm,yshift=-0.62cm]frame.north west) {Question 32 C};
\node[anchor=east,text=white,font=\sffamily\itshape\large] at ([xshift=-5mm,yshift=-0.62cm]frame.north east) {Puissances et racines carrées};}]
\vspace*{1.75cm}
\begin{minipage}[t][6.9cm][c]{\linewidth}\LARGE\centering
\begin{minipage}{0.95\linewidth}\raggedright\linespread{1.25}\selectfont Entre quels nombres entiers consécutifs se trouve $\sqrt{110}$~?\end{minipage}
\end{minipage}

\vfill
\begin{tcolorbox}[enhanced,nobeforeafter,width=7.1cm,height=2.1cm,valign=center,colback=cpui!6,colframe=cpui,arc=2.5mm,boxrule=0.9pt,left=3mm]{\Large\bfseries\color{cpui}a)}\quad{\Large $11\text{ et }12$}\end{tcolorbox}\hfill\begin{tcolorbox}[enhanced,nobeforeafter,width=7.1cm,height=2.1cm,valign=center,colback=cpui!6,colframe=cpui,arc=2.5mm,boxrule=0.9pt,left=3mm]{\Large\bfseries\color{cpui}b)}\quad{\Large $9\text{ et }10$}\end{tcolorbox}\par\vspace{2.5mm}
\begin{tcolorbox}[enhanced,nobeforeafter,width=7.1cm,height=2.1cm,valign=center,colback=cpui!6,colframe=cpui,arc=2.5mm,boxrule=0.9pt,left=3mm]{\Large\bfseries\color{cpui}c)}\quad{\Large $55\text{ et }56$}\end{tcolorbox}\hfill\begin{tcolorbox}[enhanced,nobeforeafter,width=7.1cm,height=2.1cm,valign=center,colback=cpui!6,colframe=cpui,arc=2.5mm,boxrule=0.9pt,left=3mm]{\Large\bfseries\color{cpui}d)}\quad{\Large $10\text{ et }11$}\end{tcolorbox}
\end{tcolorbox}
\newpage
\begin{tcolorbox}[enhanced,width=16cm,height=14cm,colback=white,colframe=clit,boxrule=2pt,arc=0mm,
left=6mm,right=6mm,top=0mm,bottom=4mm,sharp corners,
overlay={\fill[clit] (frame.north west) rectangle ([yshift=-1.25cm]frame.north east);
\node[anchor=west,text=white,font=\sffamily\bfseries\LARGE] at ([xshift=5mm,yshift=-0.62cm]frame.north west) {Question 33 C};
\node[anchor=east,text=white,font=\sffamily\itshape\large] at ([xshift=-5mm,yshift=-0.62cm]frame.north east) {Calcul littéral et équations};}]
\vspace*{1.75cm}
\begin{minipage}[t][6.9cm][c]{\linewidth}\LARGE\centering
\begin{minipage}{0.95\linewidth}\raggedright\linespread{1.25}\selectfont Développer $4(3x-1)$.\end{minipage}
\end{minipage}

\vfill
\begin{tcolorbox}[enhanced,nobeforeafter,width=7.1cm,height=2.1cm,valign=center,colback=clit!6,colframe=clit,arc=2.5mm,boxrule=0.9pt,left=3mm]{\Large\bfseries\color{clit}a)}\quad{\Large $3x-4$}\end{tcolorbox}\hfill\begin{tcolorbox}[enhanced,nobeforeafter,width=7.1cm,height=2.1cm,valign=center,colback=clit!6,colframe=clit,arc=2.5mm,boxrule=0.9pt,left=3mm]{\Large\bfseries\color{clit}b)}\quad{\Large $12x+4$}\end{tcolorbox}\par\vspace{2.5mm}
\begin{tcolorbox}[enhanced,nobeforeafter,width=7.1cm,height=2.1cm,valign=center,colback=clit!6,colframe=clit,arc=2.5mm,boxrule=0.9pt,left=3mm]{\Large\bfseries\color{clit}c)}\quad{\Large $12x-4$}\end{tcolorbox}\hfill\begin{tcolorbox}[enhanced,nobeforeafter,width=7.1cm,height=2.1cm,valign=center,colback=clit!6,colframe=clit,arc=2.5mm,boxrule=0.9pt,left=3mm]{\Large\bfseries\color{clit}d)}\quad{\Large $12x-1$}\end{tcolorbox}
\end{tcolorbox}
\newpage
\begin{tcolorbox}[enhanced,width=16cm,height=14cm,colback=white,colframe=clit,boxrule=2pt,arc=0mm,
left=6mm,right=6mm,top=0mm,bottom=4mm,sharp corners,
overlay={\fill[clit] (frame.north west) rectangle ([yshift=-1.25cm]frame.north east);
\node[anchor=west,text=white,font=\sffamily\bfseries\LARGE] at ([xshift=5mm,yshift=-0.62cm]frame.north west) {Question 34 C};
\node[anchor=east,text=white,font=\sffamily\itshape\large] at ([xshift=-5mm,yshift=-0.62cm]frame.north east) {Calcul littéral et équations};}]
\vspace*{1.75cm}
\begin{minipage}[t][6.9cm][c]{\linewidth}\LARGE\centering
\begin{minipage}{0.95\linewidth}\raggedright\linespread{1.25}\selectfont Développer $-5(2x+3)$.\end{minipage}
\end{minipage}

\vfill
\begin{tcolorbox}[enhanced,nobeforeafter,width=7.1cm,height=2.1cm,valign=center,colback=clit!6,colframe=clit,arc=2.5mm,boxrule=0.9pt,left=3mm]{\Large\bfseries\color{clit}a)}\quad{\Large $-10x-15$}\end{tcolorbox}\hfill\begin{tcolorbox}[enhanced,nobeforeafter,width=7.1cm,height=2.1cm,valign=center,colback=clit!6,colframe=clit,arc=2.5mm,boxrule=0.9pt,left=3mm]{\Large\bfseries\color{clit}b)}\quad{\Large $-10x+3$}\end{tcolorbox}\par\vspace{2.5mm}
\begin{tcolorbox}[enhanced,nobeforeafter,width=7.1cm,height=2.1cm,valign=center,colback=clit!6,colframe=clit,arc=2.5mm,boxrule=0.9pt,left=3mm]{\Large\bfseries\color{clit}c)}\quad{\Large $-10x+15$}\end{tcolorbox}\hfill\begin{tcolorbox}[enhanced,nobeforeafter,width=7.1cm,height=2.1cm,valign=center,colback=clit!6,colframe=clit,arc=2.5mm,boxrule=0.9pt,left=3mm]{\Large\bfseries\color{clit}d)}\quad{\Large $2x-15$}\end{tcolorbox}
\end{tcolorbox}
\newpage
\begin{tcolorbox}[enhanced,width=16cm,height=14cm,colback=white,colframe=clit,boxrule=2pt,arc=0mm,
left=6mm,right=6mm,top=0mm,bottom=4mm,sharp corners,
overlay={\fill[clit] (frame.north west) rectangle ([yshift=-1.25cm]frame.north east);
\node[anchor=west,text=white,font=\sffamily\bfseries\LARGE] at ([xshift=5mm,yshift=-0.62cm]frame.north west) {Question 35 C};
\node[anchor=east,text=white,font=\sffamily\itshape\large] at ([xshift=-5mm,yshift=-0.62cm]frame.north east) {Calcul littéral et équations};}]
\vspace*{1.75cm}
\begin{minipage}[t][6.9cm][c]{\linewidth}\LARGE\centering
\begin{minipage}{0.95\linewidth}\raggedright\linespread{1.25}\selectfont Développer et réduire $(x+1)(x+7)$.\end{minipage}
\end{minipage}

\vfill
\begin{tcolorbox}[enhanced,nobeforeafter,width=7.1cm,height=2.1cm,valign=center,colback=clit!6,colframe=clit,arc=2.5mm,boxrule=0.9pt,left=3mm]{\Large\bfseries\color{clit}a)}\quad{\Large $x^2+8x+8$}\end{tcolorbox}\hfill\begin{tcolorbox}[enhanced,nobeforeafter,width=7.1cm,height=2.1cm,valign=center,colback=clit!6,colframe=clit,arc=2.5mm,boxrule=0.9pt,left=3mm]{\Large\bfseries\color{clit}b)}\quad{\Large $x^2+7$}\end{tcolorbox}\par\vspace{2.5mm}
\begin{tcolorbox}[enhanced,nobeforeafter,width=7.1cm,height=2.1cm,valign=center,colback=clit!6,colframe=clit,arc=2.5mm,boxrule=0.9pt,left=3mm]{\Large\bfseries\color{clit}c)}\quad{\Large $x^2+7x+8$}\end{tcolorbox}\hfill\begin{tcolorbox}[enhanced,nobeforeafter,width=7.1cm,height=2.1cm,valign=center,colback=clit!6,colframe=clit,arc=2.5mm,boxrule=0.9pt,left=3mm]{\Large\bfseries\color{clit}d)}\quad{\Large $x^2+8x+7$}\end{tcolorbox}
\end{tcolorbox}
\newpage
\begin{tcolorbox}[enhanced,width=16cm,height=14cm,colback=white,colframe=clit,boxrule=2pt,arc=0mm,
left=6mm,right=6mm,top=0mm,bottom=4mm,sharp corners,
overlay={\fill[clit] (frame.north west) rectangle ([yshift=-1.25cm]frame.north east);
\node[anchor=west,text=white,font=\sffamily\bfseries\LARGE] at ([xshift=5mm,yshift=-0.62cm]frame.north west) {Question 36 C};
\node[anchor=east,text=white,font=\sffamily\itshape\large] at ([xshift=-5mm,yshift=-0.62cm]frame.north east) {Calcul littéral et équations};}]
\vspace*{1.75cm}
\begin{minipage}[t][6.9cm][c]{\linewidth}\LARGE\centering
\begin{minipage}{0.95\linewidth}\raggedright\linespread{1.25}\selectfont Développer et réduire $(x+5)(x-4)$.\end{minipage}
\end{minipage}

\vfill
\begin{tcolorbox}[enhanced,nobeforeafter,width=7.1cm,height=2.1cm,valign=center,colback=clit!6,colframe=clit,arc=2.5mm,boxrule=0.9pt,left=3mm]{\Large\bfseries\color{clit}a)}\quad{\Large $x^2+x-20$}\end{tcolorbox}\hfill\begin{tcolorbox}[enhanced,nobeforeafter,width=7.1cm,height=2.1cm,valign=center,colback=clit!6,colframe=clit,arc=2.5mm,boxrule=0.9pt,left=3mm]{\Large\bfseries\color{clit}b)}\quad{\Large $x^2+x+1$}\end{tcolorbox}\par\vspace{2.5mm}
\begin{tcolorbox}[enhanced,nobeforeafter,width=7.1cm,height=2.1cm,valign=center,colback=clit!6,colframe=clit,arc=2.5mm,boxrule=0.9pt,left=3mm]{\Large\bfseries\color{clit}c)}\quad{\Large $x^2-20x+1$}\end{tcolorbox}\hfill\begin{tcolorbox}[enhanced,nobeforeafter,width=7.1cm,height=2.1cm,valign=center,colback=clit!6,colframe=clit,arc=2.5mm,boxrule=0.9pt,left=3mm]{\Large\bfseries\color{clit}d)}\quad{\Large $x^2-20$}\end{tcolorbox}
\end{tcolorbox}
\newpage
\begin{tcolorbox}[enhanced,width=16cm,height=14cm,colback=white,colframe=clit,boxrule=2pt,arc=0mm,
left=6mm,right=6mm,top=0mm,bottom=4mm,sharp corners,
overlay={\fill[clit] (frame.north west) rectangle ([yshift=-1.25cm]frame.north east);
\node[anchor=west,text=white,font=\sffamily\bfseries\LARGE] at ([xshift=5mm,yshift=-0.62cm]frame.north west) {Question 37 C};
\node[anchor=east,text=white,font=\sffamily\itshape\large] at ([xshift=-5mm,yshift=-0.62cm]frame.north east) {Calcul littéral et équations};}]
\vspace*{1.75cm}
\begin{minipage}[t][6.9cm][c]{\linewidth}\LARGE\centering
\begin{minipage}{0.95\linewidth}\raggedright\linespread{1.25}\selectfont Réduire l'expression $3x-8-5x+1$.\end{minipage}
\end{minipage}

\vfill
\begin{tcolorbox}[enhanced,nobeforeafter,width=7.1cm,height=2.1cm,valign=center,colback=clit!6,colframe=clit,arc=2.5mm,boxrule=0.9pt,left=3mm]{\Large\bfseries\color{clit}a)}\quad{\Large $-2x-7$}\end{tcolorbox}\hfill\begin{tcolorbox}[enhanced,nobeforeafter,width=7.1cm,height=2.1cm,valign=center,colback=clit!6,colframe=clit,arc=2.5mm,boxrule=0.9pt,left=3mm]{\Large\bfseries\color{clit}b)}\quad{\Large $8x-7$}\end{tcolorbox}\par\vspace{2.5mm}
\begin{tcolorbox}[enhanced,nobeforeafter,width=7.1cm,height=2.1cm,valign=center,colback=clit!6,colframe=clit,arc=2.5mm,boxrule=0.9pt,left=3mm]{\Large\bfseries\color{clit}c)}\quad{\Large $-9x$}\end{tcolorbox}\hfill\begin{tcolorbox}[enhanced,nobeforeafter,width=7.1cm,height=2.1cm,valign=center,colback=clit!6,colframe=clit,arc=2.5mm,boxrule=0.9pt,left=3mm]{\Large\bfseries\color{clit}d)}\quad{\Large $-2x-9$}\end{tcolorbox}
\end{tcolorbox}
\newpage
\begin{tcolorbox}[enhanced,width=16cm,height=14cm,colback=white,colframe=clit,boxrule=2pt,arc=0mm,
left=6mm,right=6mm,top=0mm,bottom=4mm,sharp corners,
overlay={\fill[clit] (frame.north west) rectangle ([yshift=-1.25cm]frame.north east);
\node[anchor=west,text=white,font=\sffamily\bfseries\LARGE] at ([xshift=5mm,yshift=-0.62cm]frame.north west) {Question 38 C};
\node[anchor=east,text=white,font=\sffamily\itshape\large] at ([xshift=-5mm,yshift=-0.62cm]frame.north east) {Calcul littéral et équations};}]
\vspace*{1.75cm}
\begin{minipage}[t][6.9cm][c]{\linewidth}\LARGE\centering
\begin{minipage}{0.95\linewidth}\raggedright\linespread{1.25}\selectfont Réduire l'expression $-4x+2+9x-6$.\end{minipage}
\end{minipage}

\vfill
\begin{tcolorbox}[enhanced,nobeforeafter,width=7.1cm,height=2.1cm,valign=center,colback=clit!6,colframe=clit,arc=2.5mm,boxrule=0.9pt,left=3mm]{\Large\bfseries\color{clit}a)}\quad{\Large $-13x-4$}\end{tcolorbox}\hfill\begin{tcolorbox}[enhanced,nobeforeafter,width=7.1cm,height=2.1cm,valign=center,colback=clit!6,colframe=clit,arc=2.5mm,boxrule=0.9pt,left=3mm]{\Large\bfseries\color{clit}b)}\quad{\Large $5x+8$}\end{tcolorbox}\par\vspace{2.5mm}
\begin{tcolorbox}[enhanced,nobeforeafter,width=7.1cm,height=2.1cm,valign=center,colback=clit!6,colframe=clit,arc=2.5mm,boxrule=0.9pt,left=3mm]{\Large\bfseries\color{clit}c)}\quad{\Large $x$}\end{tcolorbox}\hfill\begin{tcolorbox}[enhanced,nobeforeafter,width=7.1cm,height=2.1cm,valign=center,colback=clit!6,colframe=clit,arc=2.5mm,boxrule=0.9pt,left=3mm]{\Large\bfseries\color{clit}d)}\quad{\Large $5x-4$}\end{tcolorbox}
\end{tcolorbox}
\newpage
\begin{tcolorbox}[enhanced,width=16cm,height=14cm,colback=white,colframe=clit,boxrule=2pt,arc=0mm,
left=6mm,right=6mm,top=0mm,bottom=4mm,sharp corners,
overlay={\fill[clit] (frame.north west) rectangle ([yshift=-1.25cm]frame.north east);
\node[anchor=west,text=white,font=\sffamily\bfseries\LARGE] at ([xshift=5mm,yshift=-0.62cm]frame.north west) {Question 39 C};
\node[anchor=east,text=white,font=\sffamily\itshape\large] at ([xshift=-5mm,yshift=-0.62cm]frame.north east) {Calcul littéral et équations};}]
\vspace*{1.75cm}
\begin{minipage}[t][6.9cm][c]{\linewidth}\LARGE\centering
\begin{minipage}{0.95\linewidth}\raggedright\linespread{1.25}\selectfont Factoriser $6x-12$.\end{minipage}
\end{minipage}

\vfill
\begin{tcolorbox}[enhanced,nobeforeafter,width=7.1cm,height=2.1cm,valign=center,colback=clit!6,colframe=clit,arc=2.5mm,boxrule=0.9pt,left=3mm]{\Large\bfseries\color{clit}a)}\quad{\Large $6(x-12)$}\end{tcolorbox}\hfill\begin{tcolorbox}[enhanced,nobeforeafter,width=7.1cm,height=2.1cm,valign=center,colback=clit!6,colframe=clit,arc=2.5mm,boxrule=0.9pt,left=3mm]{\Large\bfseries\color{clit}b)}\quad{\Large $6x(x-2)$}\end{tcolorbox}\par\vspace{2.5mm}
\begin{tcolorbox}[enhanced,nobeforeafter,width=7.1cm,height=2.1cm,valign=center,colback=clit!6,colframe=clit,arc=2.5mm,boxrule=0.9pt,left=3mm]{\Large\bfseries\color{clit}c)}\quad{\Large $6(6x-2)$}\end{tcolorbox}\hfill\begin{tcolorbox}[enhanced,nobeforeafter,width=7.1cm,height=2.1cm,valign=center,colback=clit!6,colframe=clit,arc=2.5mm,boxrule=0.9pt,left=3mm]{\Large\bfseries\color{clit}d)}\quad{\Large $6(x-2)$}\end{tcolorbox}
\end{tcolorbox}
\newpage
\begin{tcolorbox}[enhanced,width=16cm,height=14cm,colback=white,colframe=clit,boxrule=2pt,arc=0mm,
left=6mm,right=6mm,top=0mm,bottom=4mm,sharp corners,
overlay={\fill[clit] (frame.north west) rectangle ([yshift=-1.25cm]frame.north east);
\node[anchor=west,text=white,font=\sffamily\bfseries\LARGE] at ([xshift=5mm,yshift=-0.62cm]frame.north west) {Question 40 C};
\node[anchor=east,text=white,font=\sffamily\itshape\large] at ([xshift=-5mm,yshift=-0.62cm]frame.north east) {Calcul littéral et équations};}]
\vspace*{1.75cm}
\begin{minipage}[t][6.9cm][c]{\linewidth}\LARGE\centering
\begin{minipage}{0.95\linewidth}\raggedright\linespread{1.25}\selectfont Factoriser $10x+6$.\end{minipage}
\end{minipage}

\vfill
\begin{tcolorbox}[enhanced,nobeforeafter,width=7.1cm,height=2.1cm,valign=center,colback=clit!6,colframe=clit,arc=2.5mm,boxrule=0.9pt,left=3mm]{\Large\bfseries\color{clit}a)}\quad{\Large $2(5x+3)$}\end{tcolorbox}\hfill\begin{tcolorbox}[enhanced,nobeforeafter,width=7.1cm,height=2.1cm,valign=center,colback=clit!6,colframe=clit,arc=2.5mm,boxrule=0.9pt,left=3mm]{\Large\bfseries\color{clit}b)}\quad{\Large $10(x+3)$}\end{tcolorbox}\par\vspace{2.5mm}
\begin{tcolorbox}[enhanced,nobeforeafter,width=7.1cm,height=2.1cm,valign=center,colback=clit!6,colframe=clit,arc=2.5mm,boxrule=0.9pt,left=3mm]{\Large\bfseries\color{clit}c)}\quad{\Large $2(10x+3)$}\end{tcolorbox}\hfill\begin{tcolorbox}[enhanced,nobeforeafter,width=7.1cm,height=2.1cm,valign=center,colback=clit!6,colframe=clit,arc=2.5mm,boxrule=0.9pt,left=3mm]{\Large\bfseries\color{clit}d)}\quad{\Large $2(5x+6)$}\end{tcolorbox}
\end{tcolorbox}
\newpage
\begin{tcolorbox}[enhanced,width=16cm,height=14cm,colback=white,colframe=clit,boxrule=2pt,arc=0mm,
left=6mm,right=6mm,top=0mm,bottom=4mm,sharp corners,
overlay={\fill[clit] (frame.north west) rectangle ([yshift=-1.25cm]frame.north east);
\node[anchor=west,text=white,font=\sffamily\bfseries\LARGE] at ([xshift=5mm,yshift=-0.62cm]frame.north west) {Question 41 C};
\node[anchor=east,text=white,font=\sffamily\itshape\large] at ([xshift=-5mm,yshift=-0.62cm]frame.north east) {Calcul littéral et équations};}]
\vspace*{1.75cm}
\begin{minipage}[t][6.9cm][c]{\linewidth}\LARGE\centering
\begin{minipage}{0.95\linewidth}\raggedright\linespread{1.25}\selectfont Que vaut $4x^2+x-5$ pour $x=-3$~?\end{minipage}
\end{minipage}

\vfill
\begin{tcolorbox}[enhanced,nobeforeafter,width=7.1cm,height=2.1cm,valign=center,colback=clit!6,colframe=clit,arc=2.5mm,boxrule=0.9pt,left=3mm]{\Large\bfseries\color{clit}a)}\quad{\Large $-28$}\end{tcolorbox}\hfill\begin{tcolorbox}[enhanced,nobeforeafter,width=7.1cm,height=2.1cm,valign=center,colback=clit!6,colframe=clit,arc=2.5mm,boxrule=0.9pt,left=3mm]{\Large\bfseries\color{clit}b)}\quad{\Large $-32$}\end{tcolorbox}\par\vspace{2.5mm}
\begin{tcolorbox}[enhanced,nobeforeafter,width=7.1cm,height=2.1cm,valign=center,colback=clit!6,colframe=clit,arc=2.5mm,boxrule=0.9pt,left=3mm]{\Large\bfseries\color{clit}c)}\quad{\Large $28$}\end{tcolorbox}\hfill\begin{tcolorbox}[enhanced,nobeforeafter,width=7.1cm,height=2.1cm,valign=center,colback=clit!6,colframe=clit,arc=2.5mm,boxrule=0.9pt,left=3mm]{\Large\bfseries\color{clit}d)}\quad{\Large $-44$}\end{tcolorbox}
\end{tcolorbox}
\newpage
\begin{tcolorbox}[enhanced,width=16cm,height=14cm,colback=white,colframe=clit,boxrule=2pt,arc=0mm,
left=6mm,right=6mm,top=0mm,bottom=4mm,sharp corners,
overlay={\fill[clit] (frame.north west) rectangle ([yshift=-1.25cm]frame.north east);
\node[anchor=west,text=white,font=\sffamily\bfseries\LARGE] at ([xshift=5mm,yshift=-0.62cm]frame.north west) {Question 42 C};
\node[anchor=east,text=white,font=\sffamily\itshape\large] at ([xshift=-5mm,yshift=-0.62cm]frame.north east) {Calcul littéral et équations};}]
\vspace*{1.75cm}
\begin{minipage}[t][6.9cm][c]{\linewidth}\LARGE\centering
\begin{minipage}{0.95\linewidth}\raggedright\linespread{1.25}\selectfont Que vaut $2x^2-3x$ pour $x=-4$~?\end{minipage}
\end{minipage}

\vfill
\begin{tcolorbox}[enhanced,nobeforeafter,width=7.1cm,height=2.1cm,valign=center,colback=clit!6,colframe=clit,arc=2.5mm,boxrule=0.9pt,left=3mm]{\Large\bfseries\color{clit}a)}\quad{\Large $-44$}\end{tcolorbox}\hfill\begin{tcolorbox}[enhanced,nobeforeafter,width=7.1cm,height=2.1cm,valign=center,colback=clit!6,colframe=clit,arc=2.5mm,boxrule=0.9pt,left=3mm]{\Large\bfseries\color{clit}b)}\quad{\Large $-4$}\end{tcolorbox}\par\vspace{2.5mm}
\begin{tcolorbox}[enhanced,nobeforeafter,width=7.1cm,height=2.1cm,valign=center,colback=clit!6,colframe=clit,arc=2.5mm,boxrule=0.9pt,left=3mm]{\Large\bfseries\color{clit}c)}\quad{\Large $-20$}\end{tcolorbox}\hfill\begin{tcolorbox}[enhanced,nobeforeafter,width=7.1cm,height=2.1cm,valign=center,colback=clit!6,colframe=clit,arc=2.5mm,boxrule=0.9pt,left=3mm]{\Large\bfseries\color{clit}d)}\quad{\Large $44$}\end{tcolorbox}
\end{tcolorbox}
\newpage
\begin{tcolorbox}[enhanced,width=16cm,height=14cm,colback=white,colframe=clit,boxrule=2pt,arc=0mm,
left=6mm,right=6mm,top=0mm,bottom=4mm,sharp corners,
overlay={\fill[clit] (frame.north west) rectangle ([yshift=-1.25cm]frame.north east);
\node[anchor=west,text=white,font=\sffamily\bfseries\LARGE] at ([xshift=5mm,yshift=-0.62cm]frame.north west) {Question 43 C};
\node[anchor=east,text=white,font=\sffamily\itshape\large] at ([xshift=-5mm,yshift=-0.62cm]frame.north east) {Calcul littéral et équations};}]
\vspace*{1.75cm}
\begin{minipage}[t][6.9cm][c]{\linewidth}\LARGE\centering
\begin{minipage}{0.95\linewidth}\raggedright\linespread{1.25}\selectfont Quelle est la solution de l'équation $6x+4=-20$~?\end{minipage}
\end{minipage}

\vfill
\begin{tcolorbox}[enhanced,nobeforeafter,width=7.1cm,height=2.1cm,valign=center,colback=clit!6,colframe=clit,arc=2.5mm,boxrule=0.9pt,left=3mm]{\Large\bfseries\color{clit}a)}\quad{\Large $x=-30$}\end{tcolorbox}\hfill\begin{tcolorbox}[enhanced,nobeforeafter,width=7.1cm,height=2.1cm,valign=center,colback=clit!6,colframe=clit,arc=2.5mm,boxrule=0.9pt,left=3mm]{\Large\bfseries\color{clit}b)}\quad{\Large $x=-\dfrac{8}{3}$}\end{tcolorbox}\par\vspace{2.5mm}
\begin{tcolorbox}[enhanced,nobeforeafter,width=7.1cm,height=2.1cm,valign=center,colback=clit!6,colframe=clit,arc=2.5mm,boxrule=0.9pt,left=3mm]{\Large\bfseries\color{clit}c)}\quad{\Large $x=-4$}\end{tcolorbox}\hfill\begin{tcolorbox}[enhanced,nobeforeafter,width=7.1cm,height=2.1cm,valign=center,colback=clit!6,colframe=clit,arc=2.5mm,boxrule=0.9pt,left=3mm]{\Large\bfseries\color{clit}d)}\quad{\Large $x=-\dfrac{1}{4}$}\end{tcolorbox}
\end{tcolorbox}
\newpage
\begin{tcolorbox}[enhanced,width=16cm,height=14cm,colback=white,colframe=clit,boxrule=2pt,arc=0mm,
left=6mm,right=6mm,top=0mm,bottom=4mm,sharp corners,
overlay={\fill[clit] (frame.north west) rectangle ([yshift=-1.25cm]frame.north east);
\node[anchor=west,text=white,font=\sffamily\bfseries\LARGE] at ([xshift=5mm,yshift=-0.62cm]frame.north west) {Question 44 C};
\node[anchor=east,text=white,font=\sffamily\itshape\large] at ([xshift=-5mm,yshift=-0.62cm]frame.north east) {Calcul littéral et équations};}]
\vspace*{1.75cm}
\begin{minipage}[t][6.9cm][c]{\linewidth}\LARGE\centering
\begin{minipage}{0.95\linewidth}\raggedright\linespread{1.25}\selectfont Quelle est la solution de l'équation $7x+2=30$~?\end{minipage}
\end{minipage}

\vfill
\begin{tcolorbox}[enhanced,nobeforeafter,width=7.1cm,height=2.1cm,valign=center,colback=clit!6,colframe=clit,arc=2.5mm,boxrule=0.9pt,left=3mm]{\Large\bfseries\color{clit}a)}\quad{\Large $x=21$}\end{tcolorbox}\hfill\begin{tcolorbox}[enhanced,nobeforeafter,width=7.1cm,height=2.1cm,valign=center,colback=clit!6,colframe=clit,arc=2.5mm,boxrule=0.9pt,left=3mm]{\Large\bfseries\color{clit}b)}\quad{\Large $x=4$}\end{tcolorbox}\par\vspace{2.5mm}
\begin{tcolorbox}[enhanced,nobeforeafter,width=7.1cm,height=2.1cm,valign=center,colback=clit!6,colframe=clit,arc=2.5mm,boxrule=0.9pt,left=3mm]{\Large\bfseries\color{clit}c)}\quad{\Large $x=\dfrac{1}{4}$}\end{tcolorbox}\hfill\begin{tcolorbox}[enhanced,nobeforeafter,width=7.1cm,height=2.1cm,valign=center,colback=clit!6,colframe=clit,arc=2.5mm,boxrule=0.9pt,left=3mm]{\Large\bfseries\color{clit}d)}\quad{\Large $x=\dfrac{32}{7}$}\end{tcolorbox}
\end{tcolorbox}
\newpage
\begin{tcolorbox}[enhanced,width=16cm,height=14cm,colback=white,colframe=clit,boxrule=2pt,arc=0mm,
left=6mm,right=6mm,top=0mm,bottom=4mm,sharp corners,
overlay={\fill[clit] (frame.north west) rectangle ([yshift=-1.25cm]frame.north east);
\node[anchor=west,text=white,font=\sffamily\bfseries\LARGE] at ([xshift=5mm,yshift=-0.62cm]frame.north west) {Question 45 C};
\node[anchor=east,text=white,font=\sffamily\itshape\large] at ([xshift=-5mm,yshift=-0.62cm]frame.north east) {Calcul littéral et équations};}]
\vspace*{1.75cm}
\begin{minipage}[t][6.9cm][c]{\linewidth}\LARGE\centering
\begin{minipage}{0.95\linewidth}\raggedright\linespread{1.25}\selectfont Quelle est la solution de l'équation $4x-5=9x+10$~?\end{minipage}
\end{minipage}

\vfill
\begin{tcolorbox}[enhanced,nobeforeafter,width=7.1cm,height=2.1cm,valign=center,colback=clit!6,colframe=clit,arc=2.5mm,boxrule=0.9pt,left=3mm]{\Large\bfseries\color{clit}a)}\quad{\Large $x=-1$}\end{tcolorbox}\hfill\begin{tcolorbox}[enhanced,nobeforeafter,width=7.1cm,height=2.1cm,valign=center,colback=clit!6,colframe=clit,arc=2.5mm,boxrule=0.9pt,left=3mm]{\Large\bfseries\color{clit}b)}\quad{\Large $x=-3$}\end{tcolorbox}\par\vspace{2.5mm}
\begin{tcolorbox}[enhanced,nobeforeafter,width=7.1cm,height=2.1cm,valign=center,colback=clit!6,colframe=clit,arc=2.5mm,boxrule=0.9pt,left=3mm]{\Large\bfseries\color{clit}c)}\quad{\Large $x=3$}\end{tcolorbox}\hfill\begin{tcolorbox}[enhanced,nobeforeafter,width=7.1cm,height=2.1cm,valign=center,colback=clit!6,colframe=clit,arc=2.5mm,boxrule=0.9pt,left=3mm]{\Large\bfseries\color{clit}d)}\quad{\Large $x=\dfrac{15}{13}$}\end{tcolorbox}
\end{tcolorbox}
\newpage
\begin{tcolorbox}[enhanced,width=16cm,height=14cm,colback=white,colframe=clit,boxrule=2pt,arc=0mm,
left=6mm,right=6mm,top=0mm,bottom=4mm,sharp corners,
overlay={\fill[clit] (frame.north west) rectangle ([yshift=-1.25cm]frame.north east);
\node[anchor=west,text=white,font=\sffamily\bfseries\LARGE] at ([xshift=5mm,yshift=-0.62cm]frame.north west) {Question 46 C};
\node[anchor=east,text=white,font=\sffamily\itshape\large] at ([xshift=-5mm,yshift=-0.62cm]frame.north east) {Calcul littéral et équations};}]
\vspace*{1.75cm}
\begin{minipage}[t][6.9cm][c]{\linewidth}\LARGE\centering
\begin{minipage}{0.95\linewidth}\raggedright\linespread{1.25}\selectfont Quelle est la solution de l'équation $3x+8=-x-4$~?\end{minipage}
\end{minipage}

\vfill
\begin{tcolorbox}[enhanced,nobeforeafter,width=7.1cm,height=2.1cm,valign=center,colback=clit!6,colframe=clit,arc=2.5mm,boxrule=0.9pt,left=3mm]{\Large\bfseries\color{clit}a)}\quad{\Large $x=-3$}\end{tcolorbox}\hfill\begin{tcolorbox}[enhanced,nobeforeafter,width=7.1cm,height=2.1cm,valign=center,colback=clit!6,colframe=clit,arc=2.5mm,boxrule=0.9pt,left=3mm]{\Large\bfseries\color{clit}b)}\quad{\Large $x=-6$}\end{tcolorbox}\par\vspace{2.5mm}
\begin{tcolorbox}[enhanced,nobeforeafter,width=7.1cm,height=2.1cm,valign=center,colback=clit!6,colframe=clit,arc=2.5mm,boxrule=0.9pt,left=3mm]{\Large\bfseries\color{clit}c)}\quad{\Large $x=1$}\end{tcolorbox}\hfill\begin{tcolorbox}[enhanced,nobeforeafter,width=7.1cm,height=2.1cm,valign=center,colback=clit!6,colframe=clit,arc=2.5mm,boxrule=0.9pt,left=3mm]{\Large\bfseries\color{clit}d)}\quad{\Large $x=3$}\end{tcolorbox}
\end{tcolorbox}
\newpage
\begin{tcolorbox}[enhanced,width=16cm,height=14cm,colback=white,colframe=clit,boxrule=2pt,arc=0mm,
left=6mm,right=6mm,top=0mm,bottom=4mm,sharp corners,
overlay={\fill[clit] (frame.north west) rectangle ([yshift=-1.25cm]frame.north east);
\node[anchor=west,text=white,font=\sffamily\bfseries\LARGE] at ([xshift=5mm,yshift=-0.62cm]frame.north west) {Question 47 C};
\node[anchor=east,text=white,font=\sffamily\itshape\large] at ([xshift=-5mm,yshift=-0.62cm]frame.north east) {Calcul littéral et équations};}]
\vspace*{1.75cm}
\begin{minipage}[t][6.9cm][c]{\linewidth}\LARGE\centering
\begin{minipage}{0.95\linewidth}\raggedright\linespread{1.25}\selectfont Le double d'un nombre, augmenté de 9, est égal à $35$. Quel est ce nombre~?\end{minipage}
\end{minipage}

\vfill
\begin{tcolorbox}[enhanced,nobeforeafter,width=7.1cm,height=2.1cm,valign=center,colback=clit!6,colframe=clit,arc=2.5mm,boxrule=0.9pt,left=3mm]{\Large\bfseries\color{clit}a)}\quad{\Large $22$}\end{tcolorbox}\hfill\begin{tcolorbox}[enhanced,nobeforeafter,width=7.1cm,height=2.1cm,valign=center,colback=clit!6,colframe=clit,arc=2.5mm,boxrule=0.9pt,left=3mm]{\Large\bfseries\color{clit}b)}\quad{\Large $26$}\end{tcolorbox}\par\vspace{2.5mm}
\begin{tcolorbox}[enhanced,nobeforeafter,width=7.1cm,height=2.1cm,valign=center,colback=clit!6,colframe=clit,arc=2.5mm,boxrule=0.9pt,left=3mm]{\Large\bfseries\color{clit}c)}\quad{\Large $\dfrac{35}{2}$}\end{tcolorbox}\hfill\begin{tcolorbox}[enhanced,nobeforeafter,width=7.1cm,height=2.1cm,valign=center,colback=clit!6,colframe=clit,arc=2.5mm,boxrule=0.9pt,left=3mm]{\Large\bfseries\color{clit}d)}\quad{\Large $13$}\end{tcolorbox}
\end{tcolorbox}
\newpage
\begin{tcolorbox}[enhanced,width=16cm,height=14cm,colback=white,colframe=clit,boxrule=2pt,arc=0mm,
left=6mm,right=6mm,top=0mm,bottom=4mm,sharp corners,
overlay={\fill[clit] (frame.north west) rectangle ([yshift=-1.25cm]frame.north east);
\node[anchor=west,text=white,font=\sffamily\bfseries\LARGE] at ([xshift=5mm,yshift=-0.62cm]frame.north west) {Question 48 C};
\node[anchor=east,text=white,font=\sffamily\itshape\large] at ([xshift=-5mm,yshift=-0.62cm]frame.north east) {Calcul littéral et équations};}]
\vspace*{1.75cm}
\begin{minipage}[t][6.9cm][c]{\linewidth}\LARGE\centering
\begin{minipage}{0.95\linewidth}\raggedright\linespread{1.25}\selectfont Le quadruple d'un nombre, diminué de 6, est égal à $30$. Quel est ce nombre~?\end{minipage}
\end{minipage}

\vfill
\begin{tcolorbox}[enhanced,nobeforeafter,width=7.1cm,height=2.1cm,valign=center,colback=clit!6,colframe=clit,arc=2.5mm,boxrule=0.9pt,left=3mm]{\Large\bfseries\color{clit}a)}\quad{\Large $\dfrac{15}{2}$}\end{tcolorbox}\hfill\begin{tcolorbox}[enhanced,nobeforeafter,width=7.1cm,height=2.1cm,valign=center,colback=clit!6,colframe=clit,arc=2.5mm,boxrule=0.9pt,left=3mm]{\Large\bfseries\color{clit}b)}\quad{\Large $9$}\end{tcolorbox}\par\vspace{2.5mm}
\begin{tcolorbox}[enhanced,nobeforeafter,width=7.1cm,height=2.1cm,valign=center,colback=clit!6,colframe=clit,arc=2.5mm,boxrule=0.9pt,left=3mm]{\Large\bfseries\color{clit}c)}\quad{\Large $6$}\end{tcolorbox}\hfill\begin{tcolorbox}[enhanced,nobeforeafter,width=7.1cm,height=2.1cm,valign=center,colback=clit!6,colframe=clit,arc=2.5mm,boxrule=0.9pt,left=3mm]{\Large\bfseries\color{clit}d)}\quad{\Large $36$}\end{tcolorbox}
\end{tcolorbox}
\newpage
\begin{tcolorbox}[enhanced,width=16cm,height=14cm,colback=white,colframe=cgeo,boxrule=2pt,arc=0mm,
left=6mm,right=6mm,top=0mm,bottom=4mm,sharp corners,
overlay={\fill[cgeo] (frame.north west) rectangle ([yshift=-1.25cm]frame.north east);
\node[anchor=west,text=white,font=\sffamily\bfseries\LARGE] at ([xshift=5mm,yshift=-0.62cm]frame.north west) {Question 49 C};
\node[anchor=east,text=white,font=\sffamily\itshape\large] at ([xshift=-5mm,yshift=-0.62cm]frame.north east) {Géométrie};}]
\vspace*{1.75cm}
\begin{minipage}[t][6.9cm][c]{\linewidth}\LARGE\centering
\begin{minipage}{0.95\linewidth}\raggedright\linespread{1.25}\selectfont Le triangle $ABC$ est rectangle en $A$, avec $AB=8$ cm et $AC=15$ cm. Combien mesure $BC$~?\end{minipage}
\end{minipage}

\vfill
\begin{tcolorbox}[enhanced,nobeforeafter,width=7.1cm,height=2.1cm,valign=center,colback=cgeo!6,colframe=cgeo,arc=2.5mm,boxrule=0.9pt,left=3mm]{\Large\bfseries\color{cgeo}a)}\quad{\Large $\sqrt{161}~\text{cm}$}\end{tcolorbox}\hfill\begin{tcolorbox}[enhanced,nobeforeafter,width=7.1cm,height=2.1cm,valign=center,colback=cgeo!6,colframe=cgeo,arc=2.5mm,boxrule=0.9pt,left=3mm]{\Large\bfseries\color{cgeo}b)}\quad{\Large $289~\text{cm}$}\end{tcolorbox}\par\vspace{2.5mm}
\begin{tcolorbox}[enhanced,nobeforeafter,width=7.1cm,height=2.1cm,valign=center,colback=cgeo!6,colframe=cgeo,arc=2.5mm,boxrule=0.9pt,left=3mm]{\Large\bfseries\color{cgeo}c)}\quad{\Large $17~\text{cm}$}\end{tcolorbox}\hfill\begin{tcolorbox}[enhanced,nobeforeafter,width=7.1cm,height=2.1cm,valign=center,colback=cgeo!6,colframe=cgeo,arc=2.5mm,boxrule=0.9pt,left=3mm]{\Large\bfseries\color{cgeo}d)}\quad{\Large $23~\text{cm}$}\end{tcolorbox}
\end{tcolorbox}
\newpage
\begin{tcolorbox}[enhanced,width=16cm,height=14cm,colback=white,colframe=cgeo,boxrule=2pt,arc=0mm,
left=6mm,right=6mm,top=0mm,bottom=4mm,sharp corners,
overlay={\fill[cgeo] (frame.north west) rectangle ([yshift=-1.25cm]frame.north east);
\node[anchor=west,text=white,font=\sffamily\bfseries\LARGE] at ([xshift=5mm,yshift=-0.62cm]frame.north west) {Question 50 C};
\node[anchor=east,text=white,font=\sffamily\itshape\large] at ([xshift=-5mm,yshift=-0.62cm]frame.north east) {Géométrie};}]
\vspace*{1.75cm}
\begin{minipage}[t][6.9cm][c]{\linewidth}\LARGE\centering
\begin{minipage}{0.95\linewidth}\raggedright\linespread{1.25}\selectfont Le triangle $ABC$ est rectangle en $A$, avec $AB=7$ cm et $AC=24$ cm. Combien mesure $BC$~?\end{minipage}
\end{minipage}

\vfill
\begin{tcolorbox}[enhanced,nobeforeafter,width=7.1cm,height=2.1cm,valign=center,colback=cgeo!6,colframe=cgeo,arc=2.5mm,boxrule=0.9pt,left=3mm]{\Large\bfseries\color{cgeo}a)}\quad{\Large $625~\text{cm}$}\end{tcolorbox}\hfill\begin{tcolorbox}[enhanced,nobeforeafter,width=7.1cm,height=2.1cm,valign=center,colback=cgeo!6,colframe=cgeo,arc=2.5mm,boxrule=0.9pt,left=3mm]{\Large\bfseries\color{cgeo}b)}\quad{\Large $25~\text{cm}$}\end{tcolorbox}\par\vspace{2.5mm}
\begin{tcolorbox}[enhanced,nobeforeafter,width=7.1cm,height=2.1cm,valign=center,colback=cgeo!6,colframe=cgeo,arc=2.5mm,boxrule=0.9pt,left=3mm]{\Large\bfseries\color{cgeo}c)}\quad{\Large $31~\text{cm}$}\end{tcolorbox}\hfill\begin{tcolorbox}[enhanced,nobeforeafter,width=7.1cm,height=2.1cm,valign=center,colback=cgeo!6,colframe=cgeo,arc=2.5mm,boxrule=0.9pt,left=3mm]{\Large\bfseries\color{cgeo}d)}\quad{\Large $\sqrt{527}~\text{cm}$}\end{tcolorbox}
\end{tcolorbox}
\newpage
\begin{tcolorbox}[enhanced,width=16cm,height=14cm,colback=white,colframe=cgeo,boxrule=2pt,arc=0mm,
left=6mm,right=6mm,top=0mm,bottom=4mm,sharp corners,
overlay={\fill[cgeo] (frame.north west) rectangle ([yshift=-1.25cm]frame.north east);
\node[anchor=west,text=white,font=\sffamily\bfseries\LARGE] at ([xshift=5mm,yshift=-0.62cm]frame.north west) {Question 51 C};
\node[anchor=east,text=white,font=\sffamily\itshape\large] at ([xshift=-5mm,yshift=-0.62cm]frame.north east) {Géométrie};}]
\vspace*{1.75cm}
\begin{minipage}[t][6.9cm][c]{\linewidth}\LARGE\centering
\begin{minipage}{0.95\linewidth}\raggedright\linespread{1.25}\selectfont Le triangle $ABC$ est rectangle en $A$, avec $BC=20$ cm et $AB=12$ cm. Combien mesure $AC$~?\end{minipage}
\end{minipage}

\vfill
\begin{tcolorbox}[enhanced,nobeforeafter,width=7.1cm,height=2.1cm,valign=center,colback=cgeo!6,colframe=cgeo,arc=2.5mm,boxrule=0.9pt,left=3mm]{\Large\bfseries\color{cgeo}a)}\quad{\Large $8~\text{cm}$}\end{tcolorbox}\hfill\begin{tcolorbox}[enhanced,nobeforeafter,width=7.1cm,height=2.1cm,valign=center,colback=cgeo!6,colframe=cgeo,arc=2.5mm,boxrule=0.9pt,left=3mm]{\Large\bfseries\color{cgeo}b)}\quad{\Large $16~\text{cm}$}\end{tcolorbox}\par\vspace{2.5mm}
\begin{tcolorbox}[enhanced,nobeforeafter,width=7.1cm,height=2.1cm,valign=center,colback=cgeo!6,colframe=cgeo,arc=2.5mm,boxrule=0.9pt,left=3mm]{\Large\bfseries\color{cgeo}c)}\quad{\Large $256~\text{cm}$}\end{tcolorbox}\hfill\begin{tcolorbox}[enhanced,nobeforeafter,width=7.1cm,height=2.1cm,valign=center,colback=cgeo!6,colframe=cgeo,arc=2.5mm,boxrule=0.9pt,left=3mm]{\Large\bfseries\color{cgeo}d)}\quad{\Large $\sqrt{544}~\text{cm}$}\end{tcolorbox}
\end{tcolorbox}
\newpage
\begin{tcolorbox}[enhanced,width=16cm,height=14cm,colback=white,colframe=cgeo,boxrule=2pt,arc=0mm,
left=6mm,right=6mm,top=0mm,bottom=4mm,sharp corners,
overlay={\fill[cgeo] (frame.north west) rectangle ([yshift=-1.25cm]frame.north east);
\node[anchor=west,text=white,font=\sffamily\bfseries\LARGE] at ([xshift=5mm,yshift=-0.62cm]frame.north west) {Question 52 C};
\node[anchor=east,text=white,font=\sffamily\itshape\large] at ([xshift=-5mm,yshift=-0.62cm]frame.north east) {Géométrie};}]
\vspace*{1.75cm}
\begin{minipage}[t][6.9cm][c]{\linewidth}\LARGE\centering
\begin{minipage}{0.95\linewidth}\raggedright\linespread{1.25}\selectfont Le triangle $ABC$ est rectangle en $A$, avec $BC=29$ cm et $AB=21$ cm. Combien mesure $AC$~?\end{minipage}
\end{minipage}

\vfill
\begin{tcolorbox}[enhanced,nobeforeafter,width=7.1cm,height=2.1cm,valign=center,colback=cgeo!6,colframe=cgeo,arc=2.5mm,boxrule=0.9pt,left=3mm]{\Large\bfseries\color{cgeo}a)}\quad{\Large $20~\text{cm}$}\end{tcolorbox}\hfill\begin{tcolorbox}[enhanced,nobeforeafter,width=7.1cm,height=2.1cm,valign=center,colback=cgeo!6,colframe=cgeo,arc=2.5mm,boxrule=0.9pt,left=3mm]{\Large\bfseries\color{cgeo}b)}\quad{\Large $400~\text{cm}$}\end{tcolorbox}\par\vspace{2.5mm}
\begin{tcolorbox}[enhanced,nobeforeafter,width=7.1cm,height=2.1cm,valign=center,colback=cgeo!6,colframe=cgeo,arc=2.5mm,boxrule=0.9pt,left=3mm]{\Large\bfseries\color{cgeo}c)}\quad{\Large $\sqrt{1282}~\text{cm}$}\end{tcolorbox}\hfill\begin{tcolorbox}[enhanced,nobeforeafter,width=7.1cm,height=2.1cm,valign=center,colback=cgeo!6,colframe=cgeo,arc=2.5mm,boxrule=0.9pt,left=3mm]{\Large\bfseries\color{cgeo}d)}\quad{\Large $8~\text{cm}$}\end{tcolorbox}
\end{tcolorbox}
\newpage
\begin{tcolorbox}[enhanced,width=16cm,height=14cm,colback=white,colframe=cgeo,boxrule=2pt,arc=0mm,
left=6mm,right=6mm,top=0mm,bottom=4mm,sharp corners,
overlay={\fill[cgeo] (frame.north west) rectangle ([yshift=-1.25cm]frame.north east);
\node[anchor=west,text=white,font=\sffamily\bfseries\LARGE] at ([xshift=5mm,yshift=-0.62cm]frame.north west) {Question 53 C};
\node[anchor=east,text=white,font=\sffamily\itshape\large] at ([xshift=-5mm,yshift=-0.62cm]frame.north east) {Géométrie};}]
\vspace*{1.75cm}
\begin{minipage}[t][6.9cm][c]{\linewidth}\LARGE\centering
\begin{minipage}{0.95\linewidth}\raggedright\linespread{1.25}\selectfont Un triangle a des côtés de $12$ cm, $35$ cm et $37$ cm. Est-il rectangle~?\end{minipage}
\end{minipage}

\vfill
\begin{tcolorbox}[enhanced,nobeforeafter,width=7.1cm,height=2.1cm,valign=center,colback=cgeo!6,colframe=cgeo,arc=2.5mm,boxrule=0.9pt,left=3mm]{\Large\bfseries\color{cgeo}a)}\quad{\Large s'il est isocèle}\end{tcolorbox}\hfill\begin{tcolorbox}[enhanced,nobeforeafter,width=7.1cm,height=2.1cm,valign=center,colback=cgeo!6,colframe=cgeo,arc=2.5mm,boxrule=0.9pt,left=3mm]{\Large\bfseries\color{cgeo}b)}\quad{\Large non}\end{tcolorbox}\par\vspace{2.5mm}
\begin{tcolorbox}[enhanced,nobeforeafter,width=7.1cm,height=2.1cm,valign=center,colback=cgeo!6,colframe=cgeo,arc=2.5mm,boxrule=0.9pt,left=3mm]{\Large\bfseries\color{cgeo}c)}\quad{\Large oui}\end{tcolorbox}\hfill\begin{tcolorbox}[enhanced,nobeforeafter,width=7.1cm,height=2.1cm,valign=center,colback=cgeo!6,colframe=cgeo,arc=2.5mm,boxrule=0.9pt,left=3mm]{\Large\bfseries\color{cgeo}d)}\quad{\Large on ne peut pas savoir}\end{tcolorbox}
\end{tcolorbox}
\newpage
\begin{tcolorbox}[enhanced,width=16cm,height=14cm,colback=white,colframe=cgeo,boxrule=2pt,arc=0mm,
left=6mm,right=6mm,top=0mm,bottom=4mm,sharp corners,
overlay={\fill[cgeo] (frame.north west) rectangle ([yshift=-1.25cm]frame.north east);
\node[anchor=west,text=white,font=\sffamily\bfseries\LARGE] at ([xshift=5mm,yshift=-0.62cm]frame.north west) {Question 54 C};
\node[anchor=east,text=white,font=\sffamily\itshape\large] at ([xshift=-5mm,yshift=-0.62cm]frame.north east) {Géométrie};}]
\vspace*{1.75cm}
\begin{minipage}[t][6.9cm][c]{\linewidth}\LARGE\centering
\begin{minipage}{0.95\linewidth}\raggedright\linespread{1.25}\selectfont Un triangle a des côtés de $8$ cm, $9$ cm et $12$ cm. Est-il rectangle~?\end{minipage}
\end{minipage}

\vfill
\begin{tcolorbox}[enhanced,nobeforeafter,width=7.1cm,height=2.1cm,valign=center,colback=cgeo!6,colframe=cgeo,arc=2.5mm,boxrule=0.9pt,left=3mm]{\Large\bfseries\color{cgeo}a)}\quad{\Large on ne peut pas savoir}\end{tcolorbox}\hfill\begin{tcolorbox}[enhanced,nobeforeafter,width=7.1cm,height=2.1cm,valign=center,colback=cgeo!6,colframe=cgeo,arc=2.5mm,boxrule=0.9pt,left=3mm]{\Large\bfseries\color{cgeo}b)}\quad{\Large s'il est isocèle}\end{tcolorbox}\par\vspace{2.5mm}
\begin{tcolorbox}[enhanced,nobeforeafter,width=7.1cm,height=2.1cm,valign=center,colback=cgeo!6,colframe=cgeo,arc=2.5mm,boxrule=0.9pt,left=3mm]{\Large\bfseries\color{cgeo}c)}\quad{\Large oui}\end{tcolorbox}\hfill\begin{tcolorbox}[enhanced,nobeforeafter,width=7.1cm,height=2.1cm,valign=center,colback=cgeo!6,colframe=cgeo,arc=2.5mm,boxrule=0.9pt,left=3mm]{\Large\bfseries\color{cgeo}d)}\quad{\Large non}\end{tcolorbox}
\end{tcolorbox}
\newpage
\begin{tcolorbox}[enhanced,width=16cm,height=14cm,colback=white,colframe=cgeo,boxrule=2pt,arc=0mm,
left=6mm,right=6mm,top=0mm,bottom=4mm,sharp corners,
overlay={\fill[cgeo] (frame.north west) rectangle ([yshift=-1.25cm]frame.north east);
\node[anchor=west,text=white,font=\sffamily\bfseries\LARGE] at ([xshift=5mm,yshift=-0.62cm]frame.north west) {Question 55 C};
\node[anchor=east,text=white,font=\sffamily\itshape\large] at ([xshift=-5mm,yshift=-0.62cm]frame.north east) {Géométrie};}]
\vspace*{1.75cm}
\begin{minipage}[t][6.9cm][c]{\linewidth}\Large\centering
\begin{minipage}{0.95\linewidth}\raggedright\linespread{1.25}\selectfont Dans le triangle $ABC$, $M$ est sur $[AB]$, $N$ est sur $[AC]$ et les droites $(MN)$ et $(BC)$ sont parallèles. $AM=6$ cm, $AB=9$ cm et $AN=4$ cm. Combien mesure $AC$~?\end{minipage}
\end{minipage}

\vfill
\begin{tcolorbox}[enhanced,nobeforeafter,width=7.1cm,height=2.1cm,valign=center,colback=cgeo!6,colframe=cgeo,arc=2.5mm,boxrule=0.9pt,left=3mm]{\Large\bfseries\color{cgeo}a)}\quad{\Large $7~\text{cm}$}\end{tcolorbox}\hfill\begin{tcolorbox}[enhanced,nobeforeafter,width=7.1cm,height=2.1cm,valign=center,colback=cgeo!6,colframe=cgeo,arc=2.5mm,boxrule=0.9pt,left=3mm]{\Large\bfseries\color{cgeo}b)}\quad{\Large $\dfrac{8}{3}~\text{cm}$}\end{tcolorbox}\par\vspace{2.5mm}
\begin{tcolorbox}[enhanced,nobeforeafter,width=7.1cm,height=2.1cm,valign=center,colback=cgeo!6,colframe=cgeo,arc=2.5mm,boxrule=0.9pt,left=3mm]{\Large\bfseries\color{cgeo}c)}\quad{\Large $6~\text{cm}$}\end{tcolorbox}\hfill\begin{tcolorbox}[enhanced,nobeforeafter,width=7.1cm,height=2.1cm,valign=center,colback=cgeo!6,colframe=cgeo,arc=2.5mm,boxrule=0.9pt,left=3mm]{\Large\bfseries\color{cgeo}d)}\quad{\Large $13{,}5~\text{cm}$}\end{tcolorbox}
\end{tcolorbox}
\newpage
\begin{tcolorbox}[enhanced,width=16cm,height=14cm,colback=white,colframe=cgeo,boxrule=2pt,arc=0mm,
left=6mm,right=6mm,top=0mm,bottom=4mm,sharp corners,
overlay={\fill[cgeo] (frame.north west) rectangle ([yshift=-1.25cm]frame.north east);
\node[anchor=west,text=white,font=\sffamily\bfseries\LARGE] at ([xshift=5mm,yshift=-0.62cm]frame.north west) {Question 56 C};
\node[anchor=east,text=white,font=\sffamily\itshape\large] at ([xshift=-5mm,yshift=-0.62cm]frame.north east) {Géométrie};}]
\vspace*{1.75cm}
\begin{minipage}[t][6.9cm][c]{\linewidth}\Large\centering
\begin{minipage}{0.95\linewidth}\raggedright\linespread{1.25}\selectfont Dans le triangle $ABC$, $M$ est sur $[AB]$, $N$ est sur $[AC]$ et les droites $(MN)$ et $(BC)$ sont parallèles. $AM=3$ cm, $AB=12$ cm et $AN=5$ cm. Combien mesure $AC$~?\end{minipage}
\end{minipage}

\vfill
\begin{tcolorbox}[enhanced,nobeforeafter,width=7.1cm,height=2.1cm,valign=center,colback=cgeo!6,colframe=cgeo,arc=2.5mm,boxrule=0.9pt,left=3mm]{\Large\bfseries\color{cgeo}a)}\quad{\Large $7{,}2~\text{cm}$}\end{tcolorbox}\hfill\begin{tcolorbox}[enhanced,nobeforeafter,width=7.1cm,height=2.1cm,valign=center,colback=cgeo!6,colframe=cgeo,arc=2.5mm,boxrule=0.9pt,left=3mm]{\Large\bfseries\color{cgeo}b)}\quad{\Large $14~\text{cm}$}\end{tcolorbox}\par\vspace{2.5mm}
\begin{tcolorbox}[enhanced,nobeforeafter,width=7.1cm,height=2.1cm,valign=center,colback=cgeo!6,colframe=cgeo,arc=2.5mm,boxrule=0.9pt,left=3mm]{\Large\bfseries\color{cgeo}c)}\quad{\Large $20~\text{cm}$}\end{tcolorbox}\hfill\begin{tcolorbox}[enhanced,nobeforeafter,width=7.1cm,height=2.1cm,valign=center,colback=cgeo!6,colframe=cgeo,arc=2.5mm,boxrule=0.9pt,left=3mm]{\Large\bfseries\color{cgeo}d)}\quad{\Large $1{,}25~\text{cm}$}\end{tcolorbox}
\end{tcolorbox}
\newpage
\begin{tcolorbox}[enhanced,width=16cm,height=14cm,colback=white,colframe=cgeo,boxrule=2pt,arc=0mm,
left=6mm,right=6mm,top=0mm,bottom=4mm,sharp corners,
overlay={\fill[cgeo] (frame.north west) rectangle ([yshift=-1.25cm]frame.north east);
\node[anchor=west,text=white,font=\sffamily\bfseries\LARGE] at ([xshift=5mm,yshift=-0.62cm]frame.north west) {Question 57 C};
\node[anchor=east,text=white,font=\sffamily\itshape\large] at ([xshift=-5mm,yshift=-0.62cm]frame.north east) {Géométrie};}]
\vspace*{1.75cm}
\begin{minipage}[t][6.9cm][c]{\linewidth}\LARGE\centering
\begin{minipage}{0.95\linewidth}\raggedright\linespread{1.25}\selectfont Le triangle $ABC$ est rectangle en $A$, avec $BC=25$ cm et $AB=7$ cm. Que vaut $\cos\widehat{ABC}$~?\end{minipage}
\end{minipage}

\vfill
\begin{tcolorbox}[enhanced,nobeforeafter,width=7.1cm,height=2.1cm,valign=center,colback=cgeo!6,colframe=cgeo,arc=2.5mm,boxrule=0.9pt,left=3mm]{\Large\bfseries\color{cgeo}a)}\quad{\Large $\dfrac{25}{7}$}\end{tcolorbox}\hfill\begin{tcolorbox}[enhanced,nobeforeafter,width=7.1cm,height=2.1cm,valign=center,colback=cgeo!6,colframe=cgeo,arc=2.5mm,boxrule=0.9pt,left=3mm]{\Large\bfseries\color{cgeo}b)}\quad{\Large $\dfrac{7}{18}$}\end{tcolorbox}\par\vspace{2.5mm}
\begin{tcolorbox}[enhanced,nobeforeafter,width=7.1cm,height=2.1cm,valign=center,colback=cgeo!6,colframe=cgeo,arc=2.5mm,boxrule=0.9pt,left=3mm]{\Large\bfseries\color{cgeo}c)}\quad{\Large $\dfrac{18}{25}$}\end{tcolorbox}\hfill\begin{tcolorbox}[enhanced,nobeforeafter,width=7.1cm,height=2.1cm,valign=center,colback=cgeo!6,colframe=cgeo,arc=2.5mm,boxrule=0.9pt,left=3mm]{\Large\bfseries\color{cgeo}d)}\quad{\Large $\dfrac{7}{25}$}\end{tcolorbox}
\end{tcolorbox}
\newpage
\begin{tcolorbox}[enhanced,width=16cm,height=14cm,colback=white,colframe=cgeo,boxrule=2pt,arc=0mm,
left=6mm,right=6mm,top=0mm,bottom=4mm,sharp corners,
overlay={\fill[cgeo] (frame.north west) rectangle ([yshift=-1.25cm]frame.north east);
\node[anchor=west,text=white,font=\sffamily\bfseries\LARGE] at ([xshift=5mm,yshift=-0.62cm]frame.north west) {Question 58 C};
\node[anchor=east,text=white,font=\sffamily\itshape\large] at ([xshift=-5mm,yshift=-0.62cm]frame.north east) {Géométrie};}]
\vspace*{1.75cm}
\begin{minipage}[t][6.9cm][c]{\linewidth}\LARGE\centering
\begin{minipage}{0.95\linewidth}\raggedright\linespread{1.25}\selectfont Le triangle $ABC$ est rectangle en $A$, avec $BC=4$ cm et $AB=2$ cm. Que vaut $\cos\widehat{ABC}$~?\end{minipage}
\end{minipage}

\vfill
\begin{tcolorbox}[enhanced,nobeforeafter,width=7.1cm,height=2.1cm,valign=center,colback=cgeo!6,colframe=cgeo,arc=2.5mm,boxrule=0.9pt,left=3mm]{\Large\bfseries\color{cgeo}a)}\quad{\Large $2$}\end{tcolorbox}\hfill\begin{tcolorbox}[enhanced,nobeforeafter,width=7.1cm,height=2.1cm,valign=center,colback=cgeo!6,colframe=cgeo,arc=2.5mm,boxrule=0.9pt,left=3mm]{\Large\bfseries\color{cgeo}b)}\quad{\Large $\dfrac{1}{2}$}\end{tcolorbox}\par\vspace{2.5mm}
\begin{tcolorbox}[enhanced,nobeforeafter,width=7.1cm,height=2.1cm,valign=center,colback=cgeo!6,colframe=cgeo,arc=2.5mm,boxrule=0.9pt,left=3mm]{\Large\bfseries\color{cgeo}c)}\quad{\Large $\dfrac{3}{5}$}\end{tcolorbox}\hfill\begin{tcolorbox}[enhanced,nobeforeafter,width=7.1cm,height=2.1cm,valign=center,colback=cgeo!6,colframe=cgeo,arc=2.5mm,boxrule=0.9pt,left=3mm]{\Large\bfseries\color{cgeo}d)}\quad{\Large $\dfrac{1}{3}$}\end{tcolorbox}
\end{tcolorbox}
\newpage
\begin{tcolorbox}[enhanced,width=16cm,height=14cm,colback=white,colframe=cgeo,boxrule=2pt,arc=0mm,
left=6mm,right=6mm,top=0mm,bottom=4mm,sharp corners,
overlay={\fill[cgeo] (frame.north west) rectangle ([yshift=-1.25cm]frame.north east);
\node[anchor=west,text=white,font=\sffamily\bfseries\LARGE] at ([xshift=5mm,yshift=-0.62cm]frame.north west) {Question 59 C};
\node[anchor=east,text=white,font=\sffamily\itshape\large] at ([xshift=-5mm,yshift=-0.62cm]frame.north east) {Géométrie};}]
\vspace*{1.75cm}
\begin{minipage}[t][6.9cm][c]{\linewidth}\LARGE\centering
\begin{minipage}{0.95\linewidth}\raggedright\linespread{1.25}\selectfont Le triangle $ABC$ est rectangle en $A$, avec $BC=15$ cm et $\cos\widehat{ABC}=0{,}8$. Combien mesure $AB$~?\end{minipage}
\end{minipage}

\vfill
\begin{tcolorbox}[enhanced,nobeforeafter,width=7.1cm,height=2.1cm,valign=center,colback=cgeo!6,colframe=cgeo,arc=2.5mm,boxrule=0.9pt,left=3mm]{\Large\bfseries\color{cgeo}a)}\quad{\Large $13~\text{cm}$}\end{tcolorbox}\hfill\begin{tcolorbox}[enhanced,nobeforeafter,width=7.1cm,height=2.1cm,valign=center,colback=cgeo!6,colframe=cgeo,arc=2.5mm,boxrule=0.9pt,left=3mm]{\Large\bfseries\color{cgeo}b)}\quad{\Large $12~\text{cm}$}\end{tcolorbox}\par\vspace{2.5mm}
\begin{tcolorbox}[enhanced,nobeforeafter,width=7.1cm,height=2.1cm,valign=center,colback=cgeo!6,colframe=cgeo,arc=2.5mm,boxrule=0.9pt,left=3mm]{\Large\bfseries\color{cgeo}c)}\quad{\Large $6~\text{cm}$}\end{tcolorbox}\hfill\begin{tcolorbox}[enhanced,nobeforeafter,width=7.1cm,height=2.1cm,valign=center,colback=cgeo!6,colframe=cgeo,arc=2.5mm,boxrule=0.9pt,left=3mm]{\Large\bfseries\color{cgeo}d)}\quad{\Large $3~\text{cm}$}\end{tcolorbox}
\end{tcolorbox}
\newpage
\begin{tcolorbox}[enhanced,width=16cm,height=14cm,colback=white,colframe=cgeo,boxrule=2pt,arc=0mm,
left=6mm,right=6mm,top=0mm,bottom=4mm,sharp corners,
overlay={\fill[cgeo] (frame.north west) rectangle ([yshift=-1.25cm]frame.north east);
\node[anchor=west,text=white,font=\sffamily\bfseries\LARGE] at ([xshift=5mm,yshift=-0.62cm]frame.north west) {Question 60 C};
\node[anchor=east,text=white,font=\sffamily\itshape\large] at ([xshift=-5mm,yshift=-0.62cm]frame.north east) {Géométrie};}]
\vspace*{1.75cm}
\begin{minipage}[t][6.9cm][c]{\linewidth}\LARGE\centering
\begin{minipage}{0.95\linewidth}\raggedright\linespread{1.25}\selectfont Le triangle $ABC$ est rectangle en $A$, avec $BC=5$ cm et $\cos\widehat{ABC}=0{,}9$. Combien mesure $AB$~?\end{minipage}
\end{minipage}

\vfill
\begin{tcolorbox}[enhanced,nobeforeafter,width=7.1cm,height=2.1cm,valign=center,colback=cgeo!6,colframe=cgeo,arc=2.5mm,boxrule=0.9pt,left=3mm]{\Large\bfseries\color{cgeo}a)}\quad{\Large $5{,}5~\text{cm}$}\end{tcolorbox}\hfill\begin{tcolorbox}[enhanced,nobeforeafter,width=7.1cm,height=2.1cm,valign=center,colback=cgeo!6,colframe=cgeo,arc=2.5mm,boxrule=0.9pt,left=3mm]{\Large\bfseries\color{cgeo}b)}\quad{\Large $2{,}25~\text{cm}$}\end{tcolorbox}\par\vspace{2.5mm}
\begin{tcolorbox}[enhanced,nobeforeafter,width=7.1cm,height=2.1cm,valign=center,colback=cgeo!6,colframe=cgeo,arc=2.5mm,boxrule=0.9pt,left=3mm]{\Large\bfseries\color{cgeo}c)}\quad{\Large $4{,}5~\text{cm}$}\end{tcolorbox}\hfill\begin{tcolorbox}[enhanced,nobeforeafter,width=7.1cm,height=2.1cm,valign=center,colback=cgeo!6,colframe=cgeo,arc=2.5mm,boxrule=0.9pt,left=3mm]{\Large\bfseries\color{cgeo}d)}\quad{\Large $0{,}5~\text{cm}$}\end{tcolorbox}
\end{tcolorbox}
\newpage
\begin{tcolorbox}[enhanced,width=16cm,height=14cm,colback=white,colframe=cgeo,boxrule=2pt,arc=0mm,
left=6mm,right=6mm,top=0mm,bottom=4mm,sharp corners,
overlay={\fill[cgeo] (frame.north west) rectangle ([yshift=-1.25cm]frame.north east);
\node[anchor=west,text=white,font=\sffamily\bfseries\LARGE] at ([xshift=5mm,yshift=-0.62cm]frame.north west) {Question 61 C};
\node[anchor=east,text=white,font=\sffamily\itshape\large] at ([xshift=-5mm,yshift=-0.62cm]frame.north east) {Géométrie};}]
\vspace*{1.75cm}
\begin{minipage}[t][6.9cm][c]{\linewidth}\Large\centering
\begin{minipage}{0.95\linewidth}\raggedright\linespread{1.25}\selectfont Dans un repère, $A(-1\,;\,-2)$, $B(2\,;\,-1)$ et $C(-3\,;\,2)$. Quelles sont les coordonnées de l'image de $C$ par la translation qui transforme $A$ en $B$~?\end{minipage}
\end{minipage}

\vfill
\begin{tcolorbox}[enhanced,nobeforeafter,width=7.1cm,height=2.1cm,valign=center,colback=cgeo!6,colframe=cgeo,arc=2.5mm,boxrule=0.9pt,left=3mm]{\Large\bfseries\color{cgeo}a)}\quad{\Large $(-2\,;\,5)$}\end{tcolorbox}\hfill\begin{tcolorbox}[enhanced,nobeforeafter,width=7.1cm,height=2.1cm,valign=center,colback=cgeo!6,colframe=cgeo,arc=2.5mm,boxrule=0.9pt,left=3mm]{\Large\bfseries\color{cgeo}b)}\quad{\Large $(-6\,;\,1)$}\end{tcolorbox}\par\vspace{2.5mm}
\begin{tcolorbox}[enhanced,nobeforeafter,width=7.1cm,height=2.1cm,valign=center,colback=cgeo!6,colframe=cgeo,arc=2.5mm,boxrule=0.9pt,left=3mm]{\Large\bfseries\color{cgeo}c)}\quad{\Large $(0\,;\,3)$}\end{tcolorbox}\hfill\begin{tcolorbox}[enhanced,nobeforeafter,width=7.1cm,height=2.1cm,valign=center,colback=cgeo!6,colframe=cgeo,arc=2.5mm,boxrule=0.9pt,left=3mm]{\Large\bfseries\color{cgeo}d)}\quad{\Large $(-1\,;\,1)$}\end{tcolorbox}
\end{tcolorbox}
\newpage
\begin{tcolorbox}[enhanced,width=16cm,height=14cm,colback=white,colframe=cgeo,boxrule=2pt,arc=0mm,
left=6mm,right=6mm,top=0mm,bottom=4mm,sharp corners,
overlay={\fill[cgeo] (frame.north west) rectangle ([yshift=-1.25cm]frame.north east);
\node[anchor=west,text=white,font=\sffamily\bfseries\LARGE] at ([xshift=5mm,yshift=-0.62cm]frame.north west) {Question 62 C};
\node[anchor=east,text=white,font=\sffamily\itshape\large] at ([xshift=-5mm,yshift=-0.62cm]frame.north east) {Géométrie};}]
\vspace*{1.75cm}
\begin{minipage}[t][6.9cm][c]{\linewidth}\Large\centering
\begin{minipage}{0.95\linewidth}\raggedright\linespread{1.25}\selectfont Dans un repère, $A(2\,;\,4)$, $B(-1\,;\,1)$ et $C(0\,;\,-2)$. Quelles sont les coordonnées de l'image de $C$ par la translation qui transforme $A$ en $B$~?\end{minipage}
\end{minipage}

\vfill
\begin{tcolorbox}[enhanced,nobeforeafter,width=7.1cm,height=2.1cm,valign=center,colback=cgeo!6,colframe=cgeo,arc=2.5mm,boxrule=0.9pt,left=3mm]{\Large\bfseries\color{cgeo}a)}\quad{\Large $(3\,;\,1)$}\end{tcolorbox}\hfill\begin{tcolorbox}[enhanced,nobeforeafter,width=7.1cm,height=2.1cm,valign=center,colback=cgeo!6,colframe=cgeo,arc=2.5mm,boxrule=0.9pt,left=3mm]{\Large\bfseries\color{cgeo}b)}\quad{\Large $(-3\,;\,1)$}\end{tcolorbox}\par\vspace{2.5mm}
\begin{tcolorbox}[enhanced,nobeforeafter,width=7.1cm,height=2.1cm,valign=center,colback=cgeo!6,colframe=cgeo,arc=2.5mm,boxrule=0.9pt,left=3mm]{\Large\bfseries\color{cgeo}c)}\quad{\Large $(-3\,;\,-5)$}\end{tcolorbox}\hfill\begin{tcolorbox}[enhanced,nobeforeafter,width=7.1cm,height=2.1cm,valign=center,colback=cgeo!6,colframe=cgeo,arc=2.5mm,boxrule=0.9pt,left=3mm]{\Large\bfseries\color{cgeo}d)}\quad{\Large $(-1\,;\,-1)$}\end{tcolorbox}
\end{tcolorbox}
\newpage
\begin{tcolorbox}[enhanced,width=16cm,height=14cm,colback=white,colframe=cgeo,boxrule=2pt,arc=0mm,
left=6mm,right=6mm,top=0mm,bottom=4mm,sharp corners,
overlay={\fill[cgeo] (frame.north west) rectangle ([yshift=-1.25cm]frame.north east);
\node[anchor=west,text=white,font=\sffamily\bfseries\LARGE] at ([xshift=5mm,yshift=-0.62cm]frame.north west) {Question 63 C};
\node[anchor=east,text=white,font=\sffamily\itshape\large] at ([xshift=-5mm,yshift=-0.62cm]frame.north east) {Géométrie};}]
\vspace*{1.75cm}
\begin{minipage}[t][6.9cm][c]{\linewidth}\LARGE\centering
\begin{minipage}{0.95\linewidth}\raggedright\linespread{1.25}\selectfont Une figure est agrandie à l'échelle $1{,}5$ (les longueurs sont multipliées par $1{,}5$). Par combien son aire est-elle multipliée~?\end{minipage}
\end{minipage}

\vfill
\begin{tcolorbox}[enhanced,nobeforeafter,width=7.1cm,height=2.1cm,valign=center,colback=cgeo!6,colframe=cgeo,arc=2.5mm,boxrule=0.9pt,left=3mm]{\Large\bfseries\color{cgeo}a)}\quad{\Large $2{,}25$}\end{tcolorbox}\hfill\begin{tcolorbox}[enhanced,nobeforeafter,width=7.1cm,height=2.1cm,valign=center,colback=cgeo!6,colframe=cgeo,arc=2.5mm,boxrule=0.9pt,left=3mm]{\Large\bfseries\color{cgeo}b)}\quad{\Large $1{,}5$}\end{tcolorbox}\par\vspace{2.5mm}
\begin{tcolorbox}[enhanced,nobeforeafter,width=7.1cm,height=2.1cm,valign=center,colback=cgeo!6,colframe=cgeo,arc=2.5mm,boxrule=0.9pt,left=3mm]{\Large\bfseries\color{cgeo}c)}\quad{\Large $3$}\end{tcolorbox}\hfill\begin{tcolorbox}[enhanced,nobeforeafter,width=7.1cm,height=2.1cm,valign=center,colback=cgeo!6,colframe=cgeo,arc=2.5mm,boxrule=0.9pt,left=3mm]{\Large\bfseries\color{cgeo}d)}\quad{\Large $3{,}375$}\end{tcolorbox}
\end{tcolorbox}
\newpage
\begin{tcolorbox}[enhanced,width=16cm,height=14cm,colback=white,colframe=cgeo,boxrule=2pt,arc=0mm,
left=6mm,right=6mm,top=0mm,bottom=4mm,sharp corners,
overlay={\fill[cgeo] (frame.north west) rectangle ([yshift=-1.25cm]frame.north east);
\node[anchor=west,text=white,font=\sffamily\bfseries\LARGE] at ([xshift=5mm,yshift=-0.62cm]frame.north west) {Question 64 C};
\node[anchor=east,text=white,font=\sffamily\itshape\large] at ([xshift=-5mm,yshift=-0.62cm]frame.north east) {Géométrie};}]
\vspace*{1.75cm}
\begin{minipage}[t][6.9cm][c]{\linewidth}\LARGE\centering
\begin{minipage}{0.95\linewidth}\raggedright\linespread{1.25}\selectfont Une figure est agrandie à l'échelle $10$ (les longueurs sont multipliées par $10$). Par combien son aire est-elle multipliée~?\end{minipage}
\end{minipage}

\vfill
\begin{tcolorbox}[enhanced,nobeforeafter,width=7.1cm,height=2.1cm,valign=center,colback=cgeo!6,colframe=cgeo,arc=2.5mm,boxrule=0.9pt,left=3mm]{\Large\bfseries\color{cgeo}a)}\quad{\Large $10$}\end{tcolorbox}\hfill\begin{tcolorbox}[enhanced,nobeforeafter,width=7.1cm,height=2.1cm,valign=center,colback=cgeo!6,colframe=cgeo,arc=2.5mm,boxrule=0.9pt,left=3mm]{\Large\bfseries\color{cgeo}b)}\quad{\Large $1\,000$}\end{tcolorbox}\par\vspace{2.5mm}
\begin{tcolorbox}[enhanced,nobeforeafter,width=7.1cm,height=2.1cm,valign=center,colback=cgeo!6,colframe=cgeo,arc=2.5mm,boxrule=0.9pt,left=3mm]{\Large\bfseries\color{cgeo}c)}\quad{\Large $20$}\end{tcolorbox}\hfill\begin{tcolorbox}[enhanced,nobeforeafter,width=7.1cm,height=2.1cm,valign=center,colback=cgeo!6,colframe=cgeo,arc=2.5mm,boxrule=0.9pt,left=3mm]{\Large\bfseries\color{cgeo}d)}\quad{\Large $100$}\end{tcolorbox}
\end{tcolorbox}
\newpage
\begin{tcolorbox}[enhanced,width=16cm,height=14cm,colback=white,colframe=cdon,boxrule=2pt,arc=0mm,
left=6mm,right=6mm,top=0mm,bottom=4mm,sharp corners,
overlay={\fill[cdon] (frame.north west) rectangle ([yshift=-1.25cm]frame.north east);
\node[anchor=west,text=white,font=\sffamily\bfseries\LARGE] at ([xshift=5mm,yshift=-0.62cm]frame.north west) {Question 65 C};
\node[anchor=east,text=white,font=\sffamily\itshape\large] at ([xshift=-5mm,yshift=-0.62cm]frame.north east) {Proportionnalité, statistiques, probabilités};}]
\vspace*{1.75cm}
\begin{minipage}[t][6.9cm][c]{\linewidth}\LARGE\centering
\begin{minipage}{0.95\linewidth}\raggedright\linespread{1.25}\selectfont Ce tableau est un tableau de proportionnalité~: $\begin{array}{|c|c|}\hline 8 & 12\\\hline 6 &~?\\\hline\end{array}$. Quel nombre remplace «~?~»~?\end{minipage}
\end{minipage}

\vfill
\begin{tcolorbox}[enhanced,nobeforeafter,width=7.1cm,height=2.1cm,valign=center,colback=cdon!6,colframe=cdon,arc=2.5mm,boxrule=0.9pt,left=3mm]{\Large\bfseries\color{cdon}a)}\quad{\Large $10$}\end{tcolorbox}\hfill\begin{tcolorbox}[enhanced,nobeforeafter,width=7.1cm,height=2.1cm,valign=center,colback=cdon!6,colframe=cdon,arc=2.5mm,boxrule=0.9pt,left=3mm]{\Large\bfseries\color{cdon}b)}\quad{\Large $9$}\end{tcolorbox}\par\vspace{2.5mm}
\begin{tcolorbox}[enhanced,nobeforeafter,width=7.1cm,height=2.1cm,valign=center,colback=cdon!6,colframe=cdon,arc=2.5mm,boxrule=0.9pt,left=3mm]{\Large\bfseries\color{cdon}c)}\quad{\Large $72$}\end{tcolorbox}\hfill\begin{tcolorbox}[enhanced,nobeforeafter,width=7.1cm,height=2.1cm,valign=center,colback=cdon!6,colframe=cdon,arc=2.5mm,boxrule=0.9pt,left=3mm]{\Large\bfseries\color{cdon}d)}\quad{\Large $16$}\end{tcolorbox}
\end{tcolorbox}
\newpage
\begin{tcolorbox}[enhanced,width=16cm,height=14cm,colback=white,colframe=cdon,boxrule=2pt,arc=0mm,
left=6mm,right=6mm,top=0mm,bottom=4mm,sharp corners,
overlay={\fill[cdon] (frame.north west) rectangle ([yshift=-1.25cm]frame.north east);
\node[anchor=west,text=white,font=\sffamily\bfseries\LARGE] at ([xshift=5mm,yshift=-0.62cm]frame.north west) {Question 66 C};
\node[anchor=east,text=white,font=\sffamily\itshape\large] at ([xshift=-5mm,yshift=-0.62cm]frame.north east) {Proportionnalité, statistiques, probabilités};}]
\vspace*{1.75cm}
\begin{minipage}[t][6.9cm][c]{\linewidth}\LARGE\centering
\begin{minipage}{0.95\linewidth}\raggedright\linespread{1.25}\selectfont Ce tableau est un tableau de proportionnalité~: $\begin{array}{|c|c|}\hline 2 & 5\\\hline 13 &~?\\\hline\end{array}$. Quel nombre remplace «~?~»~?\end{minipage}
\end{minipage}

\vfill
\begin{tcolorbox}[enhanced,nobeforeafter,width=7.1cm,height=2.1cm,valign=center,colback=cdon!6,colframe=cdon,arc=2.5mm,boxrule=0.9pt,left=3mm]{\Large\bfseries\color{cdon}a)}\quad{\Large $\dfrac{10}{13}$}\end{tcolorbox}\hfill\begin{tcolorbox}[enhanced,nobeforeafter,width=7.1cm,height=2.1cm,valign=center,colback=cdon!6,colframe=cdon,arc=2.5mm,boxrule=0.9pt,left=3mm]{\Large\bfseries\color{cdon}b)}\quad{\Large $65$}\end{tcolorbox}\par\vspace{2.5mm}
\begin{tcolorbox}[enhanced,nobeforeafter,width=7.1cm,height=2.1cm,valign=center,colback=cdon!6,colframe=cdon,arc=2.5mm,boxrule=0.9pt,left=3mm]{\Large\bfseries\color{cdon}c)}\quad{\Large $32{,}5$}\end{tcolorbox}\hfill\begin{tcolorbox}[enhanced,nobeforeafter,width=7.1cm,height=2.1cm,valign=center,colback=cdon!6,colframe=cdon,arc=2.5mm,boxrule=0.9pt,left=3mm]{\Large\bfseries\color{cdon}d)}\quad{\Large $16$}\end{tcolorbox}
\end{tcolorbox}
\newpage
\begin{tcolorbox}[enhanced,width=16cm,height=14cm,colback=white,colframe=cdon,boxrule=2pt,arc=0mm,
left=6mm,right=6mm,top=0mm,bottom=4mm,sharp corners,
overlay={\fill[cdon] (frame.north west) rectangle ([yshift=-1.25cm]frame.north east);
\node[anchor=west,text=white,font=\sffamily\bfseries\LARGE] at ([xshift=5mm,yshift=-0.62cm]frame.north west) {Question 67 C};
\node[anchor=east,text=white,font=\sffamily\itshape\large] at ([xshift=-5mm,yshift=-0.62cm]frame.north east) {Proportionnalité, statistiques, probabilités};}]
\vspace*{1.75cm}
\begin{minipage}[t][6.9cm][c]{\linewidth}\LARGE\centering
\begin{minipage}{0.95\linewidth}\raggedright\linespread{1.25}\selectfont Que vaut $8\,\%$ de $350$~?\end{minipage}
\end{minipage}

\vfill
\begin{tcolorbox}[enhanced,nobeforeafter,width=7.1cm,height=2.1cm,valign=center,colback=cdon!6,colframe=cdon,arc=2.5mm,boxrule=0.9pt,left=3mm]{\Large\bfseries\color{cdon}a)}\quad{\Large $28$}\end{tcolorbox}\hfill\begin{tcolorbox}[enhanced,nobeforeafter,width=7.1cm,height=2.1cm,valign=center,colback=cdon!6,colframe=cdon,arc=2.5mm,boxrule=0.9pt,left=3mm]{\Large\bfseries\color{cdon}b)}\quad{\Large $322$}\end{tcolorbox}\par\vspace{2.5mm}
\begin{tcolorbox}[enhanced,nobeforeafter,width=7.1cm,height=2.1cm,valign=center,colback=cdon!6,colframe=cdon,arc=2.5mm,boxrule=0.9pt,left=3mm]{\Large\bfseries\color{cdon}c)}\quad{\Large $36$}\end{tcolorbox}\hfill\begin{tcolorbox}[enhanced,nobeforeafter,width=7.1cm,height=2.1cm,valign=center,colback=cdon!6,colframe=cdon,arc=2.5mm,boxrule=0.9pt,left=3mm]{\Large\bfseries\color{cdon}d)}\quad{\Large $280$}\end{tcolorbox}
\end{tcolorbox}
\newpage
\begin{tcolorbox}[enhanced,width=16cm,height=14cm,colback=white,colframe=cdon,boxrule=2pt,arc=0mm,
left=6mm,right=6mm,top=0mm,bottom=4mm,sharp corners,
overlay={\fill[cdon] (frame.north west) rectangle ([yshift=-1.25cm]frame.north east);
\node[anchor=west,text=white,font=\sffamily\bfseries\LARGE] at ([xshift=5mm,yshift=-0.62cm]frame.north west) {Question 68 C};
\node[anchor=east,text=white,font=\sffamily\itshape\large] at ([xshift=-5mm,yshift=-0.62cm]frame.north east) {Proportionnalité, statistiques, probabilités};}]
\vspace*{1.75cm}
\begin{minipage}[t][6.9cm][c]{\linewidth}\LARGE\centering
\begin{minipage}{0.95\linewidth}\raggedright\linespread{1.25}\selectfont Que vaut $75\,\%$ de $64$~?\end{minipage}
\end{minipage}

\vfill
\begin{tcolorbox}[enhanced,nobeforeafter,width=7.1cm,height=2.1cm,valign=center,colback=cdon!6,colframe=cdon,arc=2.5mm,boxrule=0.9pt,left=3mm]{\Large\bfseries\color{cdon}a)}\quad{\Large $123$}\end{tcolorbox}\hfill\begin{tcolorbox}[enhanced,nobeforeafter,width=7.1cm,height=2.1cm,valign=center,colback=cdon!6,colframe=cdon,arc=2.5mm,boxrule=0.9pt,left=3mm]{\Large\bfseries\color{cdon}b)}\quad{\Large $480$}\end{tcolorbox}\par\vspace{2.5mm}
\begin{tcolorbox}[enhanced,nobeforeafter,width=7.1cm,height=2.1cm,valign=center,colback=cdon!6,colframe=cdon,arc=2.5mm,boxrule=0.9pt,left=3mm]{\Large\bfseries\color{cdon}c)}\quad{\Large $48$}\end{tcolorbox}\hfill\begin{tcolorbox}[enhanced,nobeforeafter,width=7.1cm,height=2.1cm,valign=center,colback=cdon!6,colframe=cdon,arc=2.5mm,boxrule=0.9pt,left=3mm]{\Large\bfseries\color{cdon}d)}\quad{\Large $16$}\end{tcolorbox}
\end{tcolorbox}
\newpage
\begin{tcolorbox}[enhanced,width=16cm,height=14cm,colback=white,colframe=cdon,boxrule=2pt,arc=0mm,
left=6mm,right=6mm,top=0mm,bottom=4mm,sharp corners,
overlay={\fill[cdon] (frame.north west) rectangle ([yshift=-1.25cm]frame.north east);
\node[anchor=west,text=white,font=\sffamily\bfseries\LARGE] at ([xshift=5mm,yshift=-0.62cm]frame.north west) {Question 69 C};
\node[anchor=east,text=white,font=\sffamily\itshape\large] at ([xshift=-5mm,yshift=-0.62cm]frame.north east) {Proportionnalité, statistiques, probabilités};}]
\vspace*{1.75cm}
\begin{minipage}[t][6.9cm][c]{\linewidth}\LARGE\centering
\begin{minipage}{0.95\linewidth}\raggedright\linespread{1.25}\selectfont Dans une classe de $30$ élèves, $21$ sont demi-pensionnaires. Quel pourcentage des élèves cela représente-t-il~?\end{minipage}
\end{minipage}

\vfill
\begin{tcolorbox}[enhanced,nobeforeafter,width=7.1cm,height=2.1cm,valign=center,colback=cdon!6,colframe=cdon,arc=2.5mm,boxrule=0.9pt,left=3mm]{\Large\bfseries\color{cdon}a)}\quad{\Large $70\,\%$}\end{tcolorbox}\hfill\begin{tcolorbox}[enhanced,nobeforeafter,width=7.1cm,height=2.1cm,valign=center,colback=cdon!6,colframe=cdon,arc=2.5mm,boxrule=0.9pt,left=3mm]{\Large\bfseries\color{cdon}b)}\quad{\Large $21\,\%$}\end{tcolorbox}\par\vspace{2.5mm}
\begin{tcolorbox}[enhanced,nobeforeafter,width=7.1cm,height=2.1cm,valign=center,colback=cdon!6,colframe=cdon,arc=2.5mm,boxrule=0.9pt,left=3mm]{\Large\bfseries\color{cdon}c)}\quad{\Large $30\,\%$}\end{tcolorbox}\hfill\begin{tcolorbox}[enhanced,nobeforeafter,width=7.1cm,height=2.1cm,valign=center,colback=cdon!6,colframe=cdon,arc=2.5mm,boxrule=0.9pt,left=3mm]{\Large\bfseries\color{cdon}d)}\quad{\Large $7\,\%$}\end{tcolorbox}
\end{tcolorbox}
\newpage
\begin{tcolorbox}[enhanced,width=16cm,height=14cm,colback=white,colframe=cdon,boxrule=2pt,arc=0mm,
left=6mm,right=6mm,top=0mm,bottom=4mm,sharp corners,
overlay={\fill[cdon] (frame.north west) rectangle ([yshift=-1.25cm]frame.north east);
\node[anchor=west,text=white,font=\sffamily\bfseries\LARGE] at ([xshift=5mm,yshift=-0.62cm]frame.north west) {Question 70 C};
\node[anchor=east,text=white,font=\sffamily\itshape\large] at ([xshift=-5mm,yshift=-0.62cm]frame.north east) {Proportionnalité, statistiques, probabilités};}]
\vspace*{1.75cm}
\begin{minipage}[t][6.9cm][c]{\linewidth}\LARGE\centering
\begin{minipage}{0.95\linewidth}\raggedright\linespread{1.25}\selectfont Sur $80$ élèves du collège, $20$ viennent en bus. Quel pourcentage des élèves cela représente-t-il~?\end{minipage}
\end{minipage}

\vfill
\begin{tcolorbox}[enhanced,nobeforeafter,width=7.1cm,height=2.1cm,valign=center,colback=cdon!6,colframe=cdon,arc=2.5mm,boxrule=0.9pt,left=3mm]{\Large\bfseries\color{cdon}a)}\quad{\Large $2{,}5\,\%$}\end{tcolorbox}\hfill\begin{tcolorbox}[enhanced,nobeforeafter,width=7.1cm,height=2.1cm,valign=center,colback=cdon!6,colframe=cdon,arc=2.5mm,boxrule=0.9pt,left=3mm]{\Large\bfseries\color{cdon}b)}\quad{\Large $75\,\%$}\end{tcolorbox}\par\vspace{2.5mm}
\begin{tcolorbox}[enhanced,nobeforeafter,width=7.1cm,height=2.1cm,valign=center,colback=cdon!6,colframe=cdon,arc=2.5mm,boxrule=0.9pt,left=3mm]{\Large\bfseries\color{cdon}c)}\quad{\Large $20\,\%$}\end{tcolorbox}\hfill\begin{tcolorbox}[enhanced,nobeforeafter,width=7.1cm,height=2.1cm,valign=center,colback=cdon!6,colframe=cdon,arc=2.5mm,boxrule=0.9pt,left=3mm]{\Large\bfseries\color{cdon}d)}\quad{\Large $25\,\%$}\end{tcolorbox}
\end{tcolorbox}
\newpage
\begin{tcolorbox}[enhanced,width=16cm,height=14cm,colback=white,colframe=cdon,boxrule=2pt,arc=0mm,
left=6mm,right=6mm,top=0mm,bottom=4mm,sharp corners,
overlay={\fill[cdon] (frame.north west) rectangle ([yshift=-1.25cm]frame.north east);
\node[anchor=west,text=white,font=\sffamily\bfseries\LARGE] at ([xshift=5mm,yshift=-0.62cm]frame.north west) {Question 71 C};
\node[anchor=east,text=white,font=\sffamily\itshape\large] at ([xshift=-5mm,yshift=-0.62cm]frame.north east) {Proportionnalité, statistiques, probabilités};}]
\vspace*{1.75cm}
\begin{minipage}[t][6.9cm][c]{\linewidth}\LARGE\centering
\begin{minipage}{0.95\linewidth}\raggedright\linespread{1.25}\selectfont Quelle est la moyenne de la série $20\ ;\ 14\ ;\ 11$~?\end{minipage}
\end{minipage}

\vfill
\begin{tcolorbox}[enhanced,nobeforeafter,width=7.1cm,height=2.1cm,valign=center,colback=cdon!6,colframe=cdon,arc=2.5mm,boxrule=0.9pt,left=3mm]{\Large\bfseries\color{cdon}a)}\quad{\Large $45$}\end{tcolorbox}\hfill\begin{tcolorbox}[enhanced,nobeforeafter,width=7.1cm,height=2.1cm,valign=center,colback=cdon!6,colframe=cdon,arc=2.5mm,boxrule=0.9pt,left=3mm]{\Large\bfseries\color{cdon}b)}\quad{\Large $15{,}5$}\end{tcolorbox}\par\vspace{2.5mm}
\begin{tcolorbox}[enhanced,nobeforeafter,width=7.1cm,height=2.1cm,valign=center,colback=cdon!6,colframe=cdon,arc=2.5mm,boxrule=0.9pt,left=3mm]{\Large\bfseries\color{cdon}c)}\quad{\Large $15$}\end{tcolorbox}\hfill\begin{tcolorbox}[enhanced,nobeforeafter,width=7.1cm,height=2.1cm,valign=center,colback=cdon!6,colframe=cdon,arc=2.5mm,boxrule=0.9pt,left=3mm]{\Large\bfseries\color{cdon}d)}\quad{\Large $16$}\end{tcolorbox}
\end{tcolorbox}
\newpage
\begin{tcolorbox}[enhanced,width=16cm,height=14cm,colback=white,colframe=cdon,boxrule=2pt,arc=0mm,
left=6mm,right=6mm,top=0mm,bottom=4mm,sharp corners,
overlay={\fill[cdon] (frame.north west) rectangle ([yshift=-1.25cm]frame.north east);
\node[anchor=west,text=white,font=\sffamily\bfseries\LARGE] at ([xshift=5mm,yshift=-0.62cm]frame.north west) {Question 72 C};
\node[anchor=east,text=white,font=\sffamily\itshape\large] at ([xshift=-5mm,yshift=-0.62cm]frame.north east) {Proportionnalité, statistiques, probabilités};}]
\vspace*{1.75cm}
\begin{minipage}[t][6.9cm][c]{\linewidth}\LARGE\centering
\begin{minipage}{0.95\linewidth}\raggedright\linespread{1.25}\selectfont Quelle est la moyenne de la série $6\ ;\ 15\ ;\ 9\ ;\ 12\ ;\ 8$~?\end{minipage}
\end{minipage}

\vfill
\begin{tcolorbox}[enhanced,nobeforeafter,width=7.1cm,height=2.1cm,valign=center,colback=cdon!6,colframe=cdon,arc=2.5mm,boxrule=0.9pt,left=3mm]{\Large\bfseries\color{cdon}a)}\quad{\Large $10{,}5$}\end{tcolorbox}\hfill\begin{tcolorbox}[enhanced,nobeforeafter,width=7.1cm,height=2.1cm,valign=center,colback=cdon!6,colframe=cdon,arc=2.5mm,boxrule=0.9pt,left=3mm]{\Large\bfseries\color{cdon}b)}\quad{\Large $10$}\end{tcolorbox}\par\vspace{2.5mm}
\begin{tcolorbox}[enhanced,nobeforeafter,width=7.1cm,height=2.1cm,valign=center,colback=cdon!6,colframe=cdon,arc=2.5mm,boxrule=0.9pt,left=3mm]{\Large\bfseries\color{cdon}c)}\quad{\Large $50$}\end{tcolorbox}\hfill\begin{tcolorbox}[enhanced,nobeforeafter,width=7.1cm,height=2.1cm,valign=center,colback=cdon!6,colframe=cdon,arc=2.5mm,boxrule=0.9pt,left=3mm]{\Large\bfseries\color{cdon}d)}\quad{\Large $11$}\end{tcolorbox}
\end{tcolorbox}
\newpage
\begin{tcolorbox}[enhanced,width=16cm,height=14cm,colback=white,colframe=cdon,boxrule=2pt,arc=0mm,
left=6mm,right=6mm,top=0mm,bottom=4mm,sharp corners,
overlay={\fill[cdon] (frame.north west) rectangle ([yshift=-1.25cm]frame.north east);
\node[anchor=west,text=white,font=\sffamily\bfseries\LARGE] at ([xshift=5mm,yshift=-0.62cm]frame.north west) {Question 73 C};
\node[anchor=east,text=white,font=\sffamily\itshape\large] at ([xshift=-5mm,yshift=-0.62cm]frame.north east) {Proportionnalité, statistiques, probabilités};}]
\vspace*{1.75cm}
\begin{minipage}[t][6.9cm][c]{\linewidth}\LARGE\centering
\begin{minipage}{0.95\linewidth}\raggedright\linespread{1.25}\selectfont Quelle est la médiane de la série $5\ ;\ 13\ ;\ 2\ ;\ 8\ ;\ 17\ ;\ 4\ ;\ 11$~?\end{minipage}
\end{minipage}

\vfill
\begin{tcolorbox}[enhanced,nobeforeafter,width=7.1cm,height=2.1cm,valign=center,colback=cdon!6,colframe=cdon,arc=2.5mm,boxrule=0.9pt,left=3mm]{\Large\bfseries\color{cdon}a)}\quad{\Large $4$}\end{tcolorbox}\hfill\begin{tcolorbox}[enhanced,nobeforeafter,width=7.1cm,height=2.1cm,valign=center,colback=cdon!6,colframe=cdon,arc=2.5mm,boxrule=0.9pt,left=3mm]{\Large\bfseries\color{cdon}b)}\quad{\Large $8$}\end{tcolorbox}\par\vspace{2.5mm}
\begin{tcolorbox}[enhanced,nobeforeafter,width=7.1cm,height=2.1cm,valign=center,colback=cdon!6,colframe=cdon,arc=2.5mm,boxrule=0.9pt,left=3mm]{\Large\bfseries\color{cdon}c)}\quad{\Large $9$}\end{tcolorbox}\hfill\begin{tcolorbox}[enhanced,nobeforeafter,width=7.1cm,height=2.1cm,valign=center,colback=cdon!6,colframe=cdon,arc=2.5mm,boxrule=0.9pt,left=3mm]{\Large\bfseries\color{cdon}d)}\quad{\Large $\dfrac{60}{7}$}\end{tcolorbox}
\end{tcolorbox}
\newpage
\begin{tcolorbox}[enhanced,width=16cm,height=14cm,colback=white,colframe=cdon,boxrule=2pt,arc=0mm,
left=6mm,right=6mm,top=0mm,bottom=4mm,sharp corners,
overlay={\fill[cdon] (frame.north west) rectangle ([yshift=-1.25cm]frame.north east);
\node[anchor=west,text=white,font=\sffamily\bfseries\LARGE] at ([xshift=5mm,yshift=-0.62cm]frame.north west) {Question 74 C};
\node[anchor=east,text=white,font=\sffamily\itshape\large] at ([xshift=-5mm,yshift=-0.62cm]frame.north east) {Proportionnalité, statistiques, probabilités};}]
\vspace*{1.75cm}
\begin{minipage}[t][6.9cm][c]{\linewidth}\LARGE\centering
\begin{minipage}{0.95\linewidth}\raggedright\linespread{1.25}\selectfont Quelle est la médiane de la série $20\ ;\ 15\ ;\ 18\ ;\ 12$~?\end{minipage}
\end{minipage}

\vfill
\begin{tcolorbox}[enhanced,nobeforeafter,width=7.1cm,height=2.1cm,valign=center,colback=cdon!6,colframe=cdon,arc=2.5mm,boxrule=0.9pt,left=3mm]{\Large\bfseries\color{cdon}a)}\quad{\Large $16{,}5$}\end{tcolorbox}\hfill\begin{tcolorbox}[enhanced,nobeforeafter,width=7.1cm,height=2.1cm,valign=center,colback=cdon!6,colframe=cdon,arc=2.5mm,boxrule=0.9pt,left=3mm]{\Large\bfseries\color{cdon}b)}\quad{\Large $17{,}5$}\end{tcolorbox}\par\vspace{2.5mm}
\begin{tcolorbox}[enhanced,nobeforeafter,width=7.1cm,height=2.1cm,valign=center,colback=cdon!6,colframe=cdon,arc=2.5mm,boxrule=0.9pt,left=3mm]{\Large\bfseries\color{cdon}c)}\quad{\Large $2{,}5$}\end{tcolorbox}\hfill\begin{tcolorbox}[enhanced,nobeforeafter,width=7.1cm,height=2.1cm,valign=center,colback=cdon!6,colframe=cdon,arc=2.5mm,boxrule=0.9pt,left=3mm]{\Large\bfseries\color{cdon}d)}\quad{\Large $16{,}25$}\end{tcolorbox}
\end{tcolorbox}
\newpage
\begin{tcolorbox}[enhanced,width=16cm,height=14cm,colback=white,colframe=cdon,boxrule=2pt,arc=0mm,
left=6mm,right=6mm,top=0mm,bottom=4mm,sharp corners,
overlay={\fill[cdon] (frame.north west) rectangle ([yshift=-1.25cm]frame.north east);
\node[anchor=west,text=white,font=\sffamily\bfseries\LARGE] at ([xshift=5mm,yshift=-0.62cm]frame.north west) {Question 75 C};
\node[anchor=east,text=white,font=\sffamily\itshape\large] at ([xshift=-5mm,yshift=-0.62cm]frame.north east) {Proportionnalité, statistiques, probabilités};}]
\vspace*{1.75cm}
\begin{minipage}[t][6.9cm][c]{\linewidth}\LARGE\centering
\begin{minipage}{0.95\linewidth}\raggedright\linespread{1.25}\selectfont Quelle est l'étendue de la série $16\ ;\ 26\ ;\ 11\ ;\ 19$~?\end{minipage}
\end{minipage}

\vfill
\begin{tcolorbox}[enhanced,nobeforeafter,width=7.1cm,height=2.1cm,valign=center,colback=cdon!6,colframe=cdon,arc=2.5mm,boxrule=0.9pt,left=3mm]{\Large\bfseries\color{cdon}a)}\quad{\Large $15$}\end{tcolorbox}\hfill\begin{tcolorbox}[enhanced,nobeforeafter,width=7.1cm,height=2.1cm,valign=center,colback=cdon!6,colframe=cdon,arc=2.5mm,boxrule=0.9pt,left=3mm]{\Large\bfseries\color{cdon}b)}\quad{\Large $18$}\end{tcolorbox}\par\vspace{2.5mm}
\begin{tcolorbox}[enhanced,nobeforeafter,width=7.1cm,height=2.1cm,valign=center,colback=cdon!6,colframe=cdon,arc=2.5mm,boxrule=0.9pt,left=3mm]{\Large\bfseries\color{cdon}c)}\quad{\Large $37$}\end{tcolorbox}\hfill\begin{tcolorbox}[enhanced,nobeforeafter,width=7.1cm,height=2.1cm,valign=center,colback=cdon!6,colframe=cdon,arc=2.5mm,boxrule=0.9pt,left=3mm]{\Large\bfseries\color{cdon}d)}\quad{\Large $26$}\end{tcolorbox}
\end{tcolorbox}
\newpage
\begin{tcolorbox}[enhanced,width=16cm,height=14cm,colback=white,colframe=cdon,boxrule=2pt,arc=0mm,
left=6mm,right=6mm,top=0mm,bottom=4mm,sharp corners,
overlay={\fill[cdon] (frame.north west) rectangle ([yshift=-1.25cm]frame.north east);
\node[anchor=west,text=white,font=\sffamily\bfseries\LARGE] at ([xshift=5mm,yshift=-0.62cm]frame.north west) {Question 76 C};
\node[anchor=east,text=white,font=\sffamily\itshape\large] at ([xshift=-5mm,yshift=-0.62cm]frame.north east) {Proportionnalité, statistiques, probabilités};}]
\vspace*{1.75cm}
\begin{minipage}[t][6.9cm][c]{\linewidth}\LARGE\centering
\begin{minipage}{0.95\linewidth}\raggedright\linespread{1.25}\selectfont Quelle est l'étendue de la série $3\ ;\ 8\ ;\ 1\ ;\ 15\ ;\ 6$~?\end{minipage}
\end{minipage}

\vfill
\begin{tcolorbox}[enhanced,nobeforeafter,width=7.1cm,height=2.1cm,valign=center,colback=cdon!6,colframe=cdon,arc=2.5mm,boxrule=0.9pt,left=3mm]{\Large\bfseries\color{cdon}a)}\quad{\Large $15$}\end{tcolorbox}\hfill\begin{tcolorbox}[enhanced,nobeforeafter,width=7.1cm,height=2.1cm,valign=center,colback=cdon!6,colframe=cdon,arc=2.5mm,boxrule=0.9pt,left=3mm]{\Large\bfseries\color{cdon}b)}\quad{\Large $6{,}6$}\end{tcolorbox}\par\vspace{2.5mm}
\begin{tcolorbox}[enhanced,nobeforeafter,width=7.1cm,height=2.1cm,valign=center,colback=cdon!6,colframe=cdon,arc=2.5mm,boxrule=0.9pt,left=3mm]{\Large\bfseries\color{cdon}c)}\quad{\Large $16$}\end{tcolorbox}\hfill\begin{tcolorbox}[enhanced,nobeforeafter,width=7.1cm,height=2.1cm,valign=center,colback=cdon!6,colframe=cdon,arc=2.5mm,boxrule=0.9pt,left=3mm]{\Large\bfseries\color{cdon}d)}\quad{\Large $14$}\end{tcolorbox}
\end{tcolorbox}
\newpage
\begin{tcolorbox}[enhanced,width=16cm,height=14cm,colback=white,colframe=cdon,boxrule=2pt,arc=0mm,
left=6mm,right=6mm,top=0mm,bottom=4mm,sharp corners,
overlay={\fill[cdon] (frame.north west) rectangle ([yshift=-1.25cm]frame.north east);
\node[anchor=west,text=white,font=\sffamily\bfseries\LARGE] at ([xshift=5mm,yshift=-0.62cm]frame.north west) {Question 77 C};
\node[anchor=east,text=white,font=\sffamily\itshape\large] at ([xshift=-5mm,yshift=-0.62cm]frame.north east) {Proportionnalité, statistiques, probabilités};}]
\vspace*{1.75cm}
\begin{minipage}[t][6.9cm][c]{\linewidth}\Large\centering
\begin{minipage}{0.95\linewidth}\raggedright\linespread{1.25}\selectfont Un sac contient $5$ boules rouges, $5$ boules bleues et $10$ boules vertes. On tire une boule au hasard. Quelle est la probabilité de tirer une boule verte~?\end{minipage}
\end{minipage}

\vfill
\begin{tcolorbox}[enhanced,nobeforeafter,width=7.1cm,height=2.1cm,valign=center,colback=cdon!6,colframe=cdon,arc=2.5mm,boxrule=0.9pt,left=3mm]{\Large\bfseries\color{cdon}a)}\quad{\Large $\dfrac{1}{2}$}\end{tcolorbox}\hfill\begin{tcolorbox}[enhanced,nobeforeafter,width=7.1cm,height=2.1cm,valign=center,colback=cdon!6,colframe=cdon,arc=2.5mm,boxrule=0.9pt,left=3mm]{\Large\bfseries\color{cdon}b)}\quad{\Large $\dfrac{1}{10}$}\end{tcolorbox}\par\vspace{2.5mm}
\begin{tcolorbox}[enhanced,nobeforeafter,width=7.1cm,height=2.1cm,valign=center,colback=cdon!6,colframe=cdon,arc=2.5mm,boxrule=0.9pt,left=3mm]{\Large\bfseries\color{cdon}c)}\quad{\Large $1$}\end{tcolorbox}\hfill\begin{tcolorbox}[enhanced,nobeforeafter,width=7.1cm,height=2.1cm,valign=center,colback=cdon!6,colframe=cdon,arc=2.5mm,boxrule=0.9pt,left=3mm]{\Large\bfseries\color{cdon}d)}\quad{\Large $\dfrac{1}{3}$}\end{tcolorbox}
\end{tcolorbox}
\newpage
\begin{tcolorbox}[enhanced,width=16cm,height=14cm,colback=white,colframe=cdon,boxrule=2pt,arc=0mm,
left=6mm,right=6mm,top=0mm,bottom=4mm,sharp corners,
overlay={\fill[cdon] (frame.north west) rectangle ([yshift=-1.25cm]frame.north east);
\node[anchor=west,text=white,font=\sffamily\bfseries\LARGE] at ([xshift=5mm,yshift=-0.62cm]frame.north west) {Question 78 C};
\node[anchor=east,text=white,font=\sffamily\itshape\large] at ([xshift=-5mm,yshift=-0.62cm]frame.north east) {Proportionnalité, statistiques, probabilités};}]
\vspace*{1.75cm}
\begin{minipage}[t][6.9cm][c]{\linewidth}\Large\centering
\begin{minipage}{0.95\linewidth}\raggedright\linespread{1.25}\selectfont Un sac contient $8$ boules rouges, $2$ boules bleues et $6$ boules vertes. On tire une boule au hasard. Quelle est la probabilité de tirer une boule rouge~?\end{minipage}
\end{minipage}

\vfill
\begin{tcolorbox}[enhanced,nobeforeafter,width=7.1cm,height=2.1cm,valign=center,colback=cdon!6,colframe=cdon,arc=2.5mm,boxrule=0.9pt,left=3mm]{\Large\bfseries\color{cdon}a)}\quad{\Large $1$}\end{tcolorbox}\hfill\begin{tcolorbox}[enhanced,nobeforeafter,width=7.1cm,height=2.1cm,valign=center,colback=cdon!6,colframe=cdon,arc=2.5mm,boxrule=0.9pt,left=3mm]{\Large\bfseries\color{cdon}b)}\quad{\Large $\dfrac{1}{3}$}\end{tcolorbox}\par\vspace{2.5mm}
\begin{tcolorbox}[enhanced,nobeforeafter,width=7.1cm,height=2.1cm,valign=center,colback=cdon!6,colframe=cdon,arc=2.5mm,boxrule=0.9pt,left=3mm]{\Large\bfseries\color{cdon}c)}\quad{\Large $\dfrac{1}{8}$}\end{tcolorbox}\hfill\begin{tcolorbox}[enhanced,nobeforeafter,width=7.1cm,height=2.1cm,valign=center,colback=cdon!6,colframe=cdon,arc=2.5mm,boxrule=0.9pt,left=3mm]{\Large\bfseries\color{cdon}d)}\quad{\Large $\dfrac{1}{2}$}\end{tcolorbox}
\end{tcolorbox}
\newpage
\begin{tcolorbox}[enhanced,width=16cm,height=14cm,colback=white,colframe=cdon,boxrule=2pt,arc=0mm,
left=6mm,right=6mm,top=0mm,bottom=4mm,sharp corners,
overlay={\fill[cdon] (frame.north west) rectangle ([yshift=-1.25cm]frame.north east);
\node[anchor=west,text=white,font=\sffamily\bfseries\LARGE] at ([xshift=5mm,yshift=-0.62cm]frame.north west) {Question 79 C};
\node[anchor=east,text=white,font=\sffamily\itshape\large] at ([xshift=-5mm,yshift=-0.62cm]frame.north east) {Proportionnalité, statistiques, probabilités};}]
\vspace*{1.75cm}
\begin{minipage}[t][6.9cm][c]{\linewidth}\LARGE\centering
\begin{minipage}{0.95\linewidth}\raggedright\linespread{1.25}\selectfont La probabilité qu'un événement $A$ se réalise est $\dfrac{5}{12}$. Quelle est la probabilité de l'événement contraire de $A$~?\end{minipage}
\end{minipage}

\vfill
\begin{tcolorbox}[enhanced,nobeforeafter,width=7.1cm,height=2.1cm,valign=center,colback=cdon!6,colframe=cdon,arc=2.5mm,boxrule=0.9pt,left=3mm]{\Large\bfseries\color{cdon}a)}\quad{\Large $\dfrac{19}{24}$}\end{tcolorbox}\hfill\begin{tcolorbox}[enhanced,nobeforeafter,width=7.1cm,height=2.1cm,valign=center,colback=cdon!6,colframe=cdon,arc=2.5mm,boxrule=0.9pt,left=3mm]{\Large\bfseries\color{cdon}b)}\quad{\Large $\dfrac{1}{12}$}\end{tcolorbox}\par\vspace{2.5mm}
\begin{tcolorbox}[enhanced,nobeforeafter,width=7.1cm,height=2.1cm,valign=center,colback=cdon!6,colframe=cdon,arc=2.5mm,boxrule=0.9pt,left=3mm]{\Large\bfseries\color{cdon}c)}\quad{\Large $\dfrac{7}{12}$}\end{tcolorbox}\hfill\begin{tcolorbox}[enhanced,nobeforeafter,width=7.1cm,height=2.1cm,valign=center,colback=cdon!6,colframe=cdon,arc=2.5mm,boxrule=0.9pt,left=3mm]{\Large\bfseries\color{cdon}d)}\quad{\Large $\dfrac{5}{12}$}\end{tcolorbox}
\end{tcolorbox}
\newpage
\begin{tcolorbox}[enhanced,width=16cm,height=14cm,colback=white,colframe=cdon,boxrule=2pt,arc=0mm,
left=6mm,right=6mm,top=0mm,bottom=4mm,sharp corners,
overlay={\fill[cdon] (frame.north west) rectangle ([yshift=-1.25cm]frame.north east);
\node[anchor=west,text=white,font=\sffamily\bfseries\LARGE] at ([xshift=5mm,yshift=-0.62cm]frame.north west) {Question 80 C};
\node[anchor=east,text=white,font=\sffamily\itshape\large] at ([xshift=-5mm,yshift=-0.62cm]frame.north east) {Proportionnalité, statistiques, probabilités};}]
\vspace*{1.75cm}
\begin{minipage}[t][6.9cm][c]{\linewidth}\LARGE\centering
\begin{minipage}{0.95\linewidth}\raggedright\linespread{1.25}\selectfont La probabilité qu'un événement $A$ se réalise est $0{,}42$. Quelle est la probabilité de l'événement contraire de $A$~?\end{minipage}
\end{minipage}

\vfill
\begin{tcolorbox}[enhanced,nobeforeafter,width=7.1cm,height=2.1cm,valign=center,colback=cdon!6,colframe=cdon,arc=2.5mm,boxrule=0.9pt,left=3mm]{\Large\bfseries\color{cdon}a)}\quad{\Large $0{,}08$}\end{tcolorbox}\hfill\begin{tcolorbox}[enhanced,nobeforeafter,width=7.1cm,height=2.1cm,valign=center,colback=cdon!6,colframe=cdon,arc=2.5mm,boxrule=0.9pt,left=3mm]{\Large\bfseries\color{cdon}b)}\quad{\Large $0{,}79$}\end{tcolorbox}\par\vspace{2.5mm}
\begin{tcolorbox}[enhanced,nobeforeafter,width=7.1cm,height=2.1cm,valign=center,colback=cdon!6,colframe=cdon,arc=2.5mm,boxrule=0.9pt,left=3mm]{\Large\bfseries\color{cdon}c)}\quad{\Large $0{,}42$}\end{tcolorbox}\hfill\begin{tcolorbox}[enhanced,nobeforeafter,width=7.1cm,height=2.1cm,valign=center,colback=cdon!6,colframe=cdon,arc=2.5mm,boxrule=0.9pt,left=3mm]{\Large\bfseries\color{cdon}d)}\quad{\Large $0{,}58$}\end{tcolorbox}
\end{tcolorbox}
\newpage
\begin{tcolorbox}[enhanced,width=16cm,height=14cm,colback=white,colframe=cgra,boxrule=2pt,arc=0mm,
left=6mm,right=6mm,top=0mm,bottom=4mm,sharp corners,
overlay={\fill[cgra] (frame.north west) rectangle ([yshift=-1.25cm]frame.north east);
\node[anchor=west,text=white,font=\sffamily\bfseries\LARGE] at ([xshift=5mm,yshift=-0.62cm]frame.north west) {Question 81 C};
\node[anchor=east,text=white,font=\sffamily\itshape\large] at ([xshift=-5mm,yshift=-0.62cm]frame.north east) {Grandeurs et mesures};}]
\vspace*{1.75cm}
\begin{minipage}[t][6.9cm][c]{\linewidth}\LARGE\centering
\begin{minipage}{0.95\linewidth}\raggedright\linespread{1.25}\selectfont Un véhicule parcourt $210$ km en 2 h 30 min. Quelle est sa vitesse moyenne~?\end{minipage}
\end{minipage}

\vfill
\begin{tcolorbox}[enhanced,nobeforeafter,width=7.1cm,height=2.1cm,valign=center,colback=cgra!6,colframe=cgra,arc=2.5mm,boxrule=0.9pt,left=3mm]{\Large\bfseries\color{cgra}a)}\quad{\Large $525~\text{km/h}$}\end{tcolorbox}\hfill\begin{tcolorbox}[enhanced,nobeforeafter,width=7.1cm,height=2.1cm,valign=center,colback=cgra!6,colframe=cgra,arc=2.5mm,boxrule=0.9pt,left=3mm]{\Large\bfseries\color{cgra}b)}\quad{\Large $1{,}4~\text{km/h}$}\end{tcolorbox}\par\vspace{2.5mm}
\begin{tcolorbox}[enhanced,nobeforeafter,width=7.1cm,height=2.1cm,valign=center,colback=cgra!6,colframe=cgra,arc=2.5mm,boxrule=0.9pt,left=3mm]{\Large\bfseries\color{cgra}c)}\quad{\Large $94~\text{km/h}$}\end{tcolorbox}\hfill\begin{tcolorbox}[enhanced,nobeforeafter,width=7.1cm,height=2.1cm,valign=center,colback=cgra!6,colframe=cgra,arc=2.5mm,boxrule=0.9pt,left=3mm]{\Large\bfseries\color{cgra}d)}\quad{\Large $84~\text{km/h}$}\end{tcolorbox}
\end{tcolorbox}
\newpage
\begin{tcolorbox}[enhanced,width=16cm,height=14cm,colback=white,colframe=cgra,boxrule=2pt,arc=0mm,
left=6mm,right=6mm,top=0mm,bottom=4mm,sharp corners,
overlay={\fill[cgra] (frame.north west) rectangle ([yshift=-1.25cm]frame.north east);
\node[anchor=west,text=white,font=\sffamily\bfseries\LARGE] at ([xshift=5mm,yshift=-0.62cm]frame.north west) {Question 82 C};
\node[anchor=east,text=white,font=\sffamily\itshape\large] at ([xshift=-5mm,yshift=-0.62cm]frame.north east) {Grandeurs et mesures};}]
\vspace*{1.75cm}
\begin{minipage}[t][6.9cm][c]{\linewidth}\LARGE\centering
\begin{minipage}{0.95\linewidth}\raggedright\linespread{1.25}\selectfont Un véhicule parcourt $12$ km en 15 min. Quelle est sa vitesse moyenne~?\end{minipage}
\end{minipage}

\vfill
\begin{tcolorbox}[enhanced,nobeforeafter,width=7.1cm,height=2.1cm,valign=center,colback=cgra!6,colframe=cgra,arc=2.5mm,boxrule=0.9pt,left=3mm]{\Large\bfseries\color{cgra}a)}\quad{\Large $3~\text{km/h}$}\end{tcolorbox}\hfill\begin{tcolorbox}[enhanced,nobeforeafter,width=7.1cm,height=2.1cm,valign=center,colback=cgra!6,colframe=cgra,arc=2.5mm,boxrule=0.9pt,left=3mm]{\Large\bfseries\color{cgra}b)}\quad{\Large $0{,}8~\text{km/h}$}\end{tcolorbox}\par\vspace{2.5mm}
\begin{tcolorbox}[enhanced,nobeforeafter,width=7.1cm,height=2.1cm,valign=center,colback=cgra!6,colframe=cgra,arc=2.5mm,boxrule=0.9pt,left=3mm]{\Large\bfseries\color{cgra}c)}\quad{\Large $58~\text{km/h}$}\end{tcolorbox}\hfill\begin{tcolorbox}[enhanced,nobeforeafter,width=7.1cm,height=2.1cm,valign=center,colback=cgra!6,colframe=cgra,arc=2.5mm,boxrule=0.9pt,left=3mm]{\Large\bfseries\color{cgra}d)}\quad{\Large $48~\text{km/h}$}\end{tcolorbox}
\end{tcolorbox}
\newpage
\begin{tcolorbox}[enhanced,width=16cm,height=14cm,colback=white,colframe=cgra,boxrule=2pt,arc=0mm,
left=6mm,right=6mm,top=0mm,bottom=4mm,sharp corners,
overlay={\fill[cgra] (frame.north west) rectangle ([yshift=-1.25cm]frame.north east);
\node[anchor=west,text=white,font=\sffamily\bfseries\LARGE] at ([xshift=5mm,yshift=-0.62cm]frame.north west) {Question 83 C};
\node[anchor=east,text=white,font=\sffamily\itshape\large] at ([xshift=-5mm,yshift=-0.62cm]frame.north east) {Grandeurs et mesures};}]
\vspace*{1.75cm}
\begin{minipage}[t][6.9cm][c]{\linewidth}\LARGE\centering
\begin{minipage}{0.95\linewidth}\raggedright\linespread{1.25}\selectfont Un cycliste roule à la vitesse constante de $24$ km/h pendant $15$ min. Quelle distance parcourt-il~?\end{minipage}
\end{minipage}

\vfill
\begin{tcolorbox}[enhanced,nobeforeafter,width=7.1cm,height=2.1cm,valign=center,colback=cgra!6,colframe=cgra,arc=2.5mm,boxrule=0.9pt,left=3mm]{\Large\bfseries\color{cgra}a)}\quad{\Large $12~\text{km}$}\end{tcolorbox}\hfill\begin{tcolorbox}[enhanced,nobeforeafter,width=7.1cm,height=2.1cm,valign=center,colback=cgra!6,colframe=cgra,arc=2.5mm,boxrule=0.9pt,left=3mm]{\Large\bfseries\color{cgra}b)}\quad{\Large $360~\text{km}$}\end{tcolorbox}\par\vspace{2.5mm}
\begin{tcolorbox}[enhanced,nobeforeafter,width=7.1cm,height=2.1cm,valign=center,colback=cgra!6,colframe=cgra,arc=2.5mm,boxrule=0.9pt,left=3mm]{\Large\bfseries\color{cgra}c)}\quad{\Large $6~\text{km}$}\end{tcolorbox}\hfill\begin{tcolorbox}[enhanced,nobeforeafter,width=7.1cm,height=2.1cm,valign=center,colback=cgra!6,colframe=cgra,arc=2.5mm,boxrule=0.9pt,left=3mm]{\Large\bfseries\color{cgra}d)}\quad{\Large $3{,}6~\text{km}$}\end{tcolorbox}
\end{tcolorbox}
\newpage
\begin{tcolorbox}[enhanced,width=16cm,height=14cm,colback=white,colframe=cgra,boxrule=2pt,arc=0mm,
left=6mm,right=6mm,top=0mm,bottom=4mm,sharp corners,
overlay={\fill[cgra] (frame.north west) rectangle ([yshift=-1.25cm]frame.north east);
\node[anchor=west,text=white,font=\sffamily\bfseries\LARGE] at ([xshift=5mm,yshift=-0.62cm]frame.north west) {Question 84 C};
\node[anchor=east,text=white,font=\sffamily\itshape\large] at ([xshift=-5mm,yshift=-0.62cm]frame.north east) {Grandeurs et mesures};}]
\vspace*{1.75cm}
\begin{minipage}[t][6.9cm][c]{\linewidth}\LARGE\centering
\begin{minipage}{0.95\linewidth}\raggedright\linespread{1.25}\selectfont Une voiture roule à la vitesse constante de $80$ km/h pendant $90$ min. Quelle distance parcourt-elle~?\end{minipage}
\end{minipage}

\vfill
\begin{tcolorbox}[enhanced,nobeforeafter,width=7.1cm,height=2.1cm,valign=center,colback=cgra!6,colframe=cgra,arc=2.5mm,boxrule=0.9pt,left=3mm]{\Large\bfseries\color{cgra}a)}\quad{\Large $7\,200~\text{km}$}\end{tcolorbox}\hfill\begin{tcolorbox}[enhanced,nobeforeafter,width=7.1cm,height=2.1cm,valign=center,colback=cgra!6,colframe=cgra,arc=2.5mm,boxrule=0.9pt,left=3mm]{\Large\bfseries\color{cgra}b)}\quad{\Large $120~\text{km}$}\end{tcolorbox}\par\vspace{2.5mm}
\begin{tcolorbox}[enhanced,nobeforeafter,width=7.1cm,height=2.1cm,valign=center,colback=cgra!6,colframe=cgra,arc=2.5mm,boxrule=0.9pt,left=3mm]{\Large\bfseries\color{cgra}c)}\quad{\Large $240~\text{km}$}\end{tcolorbox}\hfill\begin{tcolorbox}[enhanced,nobeforeafter,width=7.1cm,height=2.1cm,valign=center,colback=cgra!6,colframe=cgra,arc=2.5mm,boxrule=0.9pt,left=3mm]{\Large\bfseries\color{cgra}d)}\quad{\Large $72~\text{km}$}\end{tcolorbox}
\end{tcolorbox}
\newpage
\begin{tcolorbox}[enhanced,width=16cm,height=14cm,colback=white,colframe=cgra,boxrule=2pt,arc=0mm,
left=6mm,right=6mm,top=0mm,bottom=4mm,sharp corners,
overlay={\fill[cgra] (frame.north west) rectangle ([yshift=-1.25cm]frame.north east);
\node[anchor=west,text=white,font=\sffamily\bfseries\LARGE] at ([xshift=5mm,yshift=-0.62cm]frame.north west) {Question 85 C};
\node[anchor=east,text=white,font=\sffamily\itshape\large] at ([xshift=-5mm,yshift=-0.62cm]frame.north east) {Grandeurs et mesures};}]
\vspace*{1.75cm}
\begin{minipage}[t][6.9cm][c]{\linewidth}\LARGE\centering
\begin{minipage}{0.95\linewidth}\raggedright\linespread{1.25}\selectfont Un train parcourt $120$ km à la vitesse moyenne de $80$ km/h. Combien de temps dure le trajet~?\end{minipage}
\end{minipage}

\vfill
\begin{tcolorbox}[enhanced,nobeforeafter,width=7.1cm,height=2.1cm,valign=center,colback=cgra!6,colframe=cgra,arc=2.5mm,boxrule=0.9pt,left=3mm]{\Large\bfseries\color{cgra}a)}\quad{\Large 1 h 15 min}\end{tcolorbox}\hfill\begin{tcolorbox}[enhanced,nobeforeafter,width=7.1cm,height=2.1cm,valign=center,colback=cgra!6,colframe=cgra,arc=2.5mm,boxrule=0.9pt,left=3mm]{\Large\bfseries\color{cgra}b)}\quad{\Large 1 h 30 min}\end{tcolorbox}\par\vspace{2.5mm}
\begin{tcolorbox}[enhanced,nobeforeafter,width=7.1cm,height=2.1cm,valign=center,colback=cgra!6,colframe=cgra,arc=2.5mm,boxrule=0.9pt,left=3mm]{\Large\bfseries\color{cgra}c)}\quad{\Large 2 h}\end{tcolorbox}\hfill\begin{tcolorbox}[enhanced,nobeforeafter,width=7.1cm,height=2.1cm,valign=center,colback=cgra!6,colframe=cgra,arc=2.5mm,boxrule=0.9pt,left=3mm]{\Large\bfseries\color{cgra}d)}\quad{\Large 1 h 50 min}\end{tcolorbox}
\end{tcolorbox}
\newpage
\begin{tcolorbox}[enhanced,width=16cm,height=14cm,colback=white,colframe=cgra,boxrule=2pt,arc=0mm,
left=6mm,right=6mm,top=0mm,bottom=4mm,sharp corners,
overlay={\fill[cgra] (frame.north west) rectangle ([yshift=-1.25cm]frame.north east);
\node[anchor=west,text=white,font=\sffamily\bfseries\LARGE] at ([xshift=5mm,yshift=-0.62cm]frame.north west) {Question 86 C};
\node[anchor=east,text=white,font=\sffamily\itshape\large] at ([xshift=-5mm,yshift=-0.62cm]frame.north east) {Grandeurs et mesures};}]
\vspace*{1.75cm}
\begin{minipage}[t][6.9cm][c]{\linewidth}\LARGE\centering
\begin{minipage}{0.95\linewidth}\raggedright\linespread{1.25}\selectfont Un train parcourt $25$ km à la vitesse moyenne de $20$ km/h. Combien de temps dure le trajet~?\end{minipage}
\end{minipage}

\vfill
\begin{tcolorbox}[enhanced,nobeforeafter,width=7.1cm,height=2.1cm,valign=center,colback=cgra!6,colframe=cgra,arc=2.5mm,boxrule=0.9pt,left=3mm]{\Large\bfseries\color{cgra}a)}\quad{\Large 1 h}\end{tcolorbox}\hfill\begin{tcolorbox}[enhanced,nobeforeafter,width=7.1cm,height=2.1cm,valign=center,colback=cgra!6,colframe=cgra,arc=2.5mm,boxrule=0.9pt,left=3mm]{\Large\bfseries\color{cgra}b)}\quad{\Large 1 h 15 min}\end{tcolorbox}\par\vspace{2.5mm}
\begin{tcolorbox}[enhanced,nobeforeafter,width=7.1cm,height=2.1cm,valign=center,colback=cgra!6,colframe=cgra,arc=2.5mm,boxrule=0.9pt,left=3mm]{\Large\bfseries\color{cgra}c)}\quad{\Large 1 h 25 min}\end{tcolorbox}\hfill\begin{tcolorbox}[enhanced,nobeforeafter,width=7.1cm,height=2.1cm,valign=center,colback=cgra!6,colframe=cgra,arc=2.5mm,boxrule=0.9pt,left=3mm]{\Large\bfseries\color{cgra}d)}\quad{\Large 1 h 45 min}\end{tcolorbox}
\end{tcolorbox}
\newpage
\begin{tcolorbox}[enhanced,width=16cm,height=14cm,colback=white,colframe=cgra,boxrule=2pt,arc=0mm,
left=6mm,right=6mm,top=0mm,bottom=4mm,sharp corners,
overlay={\fill[cgra] (frame.north west) rectangle ([yshift=-1.25cm]frame.north east);
\node[anchor=west,text=white,font=\sffamily\bfseries\LARGE] at ([xshift=5mm,yshift=-0.62cm]frame.north west) {Question 87 C};
\node[anchor=east,text=white,font=\sffamily\itshape\large] at ([xshift=-5mm,yshift=-0.62cm]frame.north east) {Grandeurs et mesures};}]
\vspace*{1.75cm}
\begin{minipage}[t][6.9cm][c]{\linewidth}\LARGE\centering
\begin{minipage}{0.95\linewidth}\raggedright\linespread{1.25}\selectfont Un véhicule roule à $18$ km/h. Quelle est sa vitesse en mètres par seconde~?\end{minipage}
\end{minipage}

\vfill
\begin{tcolorbox}[enhanced,nobeforeafter,width=7.1cm,height=2.1cm,valign=center,colback=cgra!6,colframe=cgra,arc=2.5mm,boxrule=0.9pt,left=3mm]{\Large\bfseries\color{cgra}a)}\quad{\Large $5~\text{m/s}$}\end{tcolorbox}\hfill\begin{tcolorbox}[enhanced,nobeforeafter,width=7.1cm,height=2.1cm,valign=center,colback=cgra!6,colframe=cgra,arc=2.5mm,boxrule=0.9pt,left=3mm]{\Large\bfseries\color{cgra}b)}\quad{\Large $0{,}3~\text{m/s}$}\end{tcolorbox}\par\vspace{2.5mm}
\begin{tcolorbox}[enhanced,nobeforeafter,width=7.1cm,height=2.1cm,valign=center,colback=cgra!6,colframe=cgra,arc=2.5mm,boxrule=0.9pt,left=3mm]{\Large\bfseries\color{cgra}c)}\quad{\Large $18\,000~\text{m/s}$}\end{tcolorbox}\hfill\begin{tcolorbox}[enhanced,nobeforeafter,width=7.1cm,height=2.1cm,valign=center,colback=cgra!6,colframe=cgra,arc=2.5mm,boxrule=0.9pt,left=3mm]{\Large\bfseries\color{cgra}d)}\quad{\Large $64{,}8~\text{m/s}$}\end{tcolorbox}
\end{tcolorbox}
\newpage
\begin{tcolorbox}[enhanced,width=16cm,height=14cm,colback=white,colframe=cgra,boxrule=2pt,arc=0mm,
left=6mm,right=6mm,top=0mm,bottom=4mm,sharp corners,
overlay={\fill[cgra] (frame.north west) rectangle ([yshift=-1.25cm]frame.north east);
\node[anchor=west,text=white,font=\sffamily\bfseries\LARGE] at ([xshift=5mm,yshift=-0.62cm]frame.north west) {Question 88 C};
\node[anchor=east,text=white,font=\sffamily\itshape\large] at ([xshift=-5mm,yshift=-0.62cm]frame.north east) {Grandeurs et mesures};}]
\vspace*{1.75cm}
\begin{minipage}[t][6.9cm][c]{\linewidth}\LARGE\centering
\begin{minipage}{0.95\linewidth}\raggedright\linespread{1.25}\selectfont Un véhicule roule à $108$ km/h. Quelle est sa vitesse en mètres par seconde~?\end{minipage}
\end{minipage}

\vfill
\begin{tcolorbox}[enhanced,nobeforeafter,width=7.1cm,height=2.1cm,valign=center,colback=cgra!6,colframe=cgra,arc=2.5mm,boxrule=0.9pt,left=3mm]{\Large\bfseries\color{cgra}a)}\quad{\Large $108\,000~\text{m/s}$}\end{tcolorbox}\hfill\begin{tcolorbox}[enhanced,nobeforeafter,width=7.1cm,height=2.1cm,valign=center,colback=cgra!6,colframe=cgra,arc=2.5mm,boxrule=0.9pt,left=3mm]{\Large\bfseries\color{cgra}b)}\quad{\Large $1{,}8~\text{m/s}$}\end{tcolorbox}\par\vspace{2.5mm}
\begin{tcolorbox}[enhanced,nobeforeafter,width=7.1cm,height=2.1cm,valign=center,colback=cgra!6,colframe=cgra,arc=2.5mm,boxrule=0.9pt,left=3mm]{\Large\bfseries\color{cgra}c)}\quad{\Large $388{,}8~\text{m/s}$}\end{tcolorbox}\hfill\begin{tcolorbox}[enhanced,nobeforeafter,width=7.1cm,height=2.1cm,valign=center,colback=cgra!6,colframe=cgra,arc=2.5mm,boxrule=0.9pt,left=3mm]{\Large\bfseries\color{cgra}d)}\quad{\Large $30~\text{m/s}$}\end{tcolorbox}
\end{tcolorbox}
\newpage
\begin{tcolorbox}[enhanced,width=16cm,height=14cm,colback=white,colframe=cgra,boxrule=2pt,arc=0mm,
left=6mm,right=6mm,top=0mm,bottom=4mm,sharp corners,
overlay={\fill[cgra] (frame.north west) rectangle ([yshift=-1.25cm]frame.north east);
\node[anchor=west,text=white,font=\sffamily\bfseries\LARGE] at ([xshift=5mm,yshift=-0.62cm]frame.north west) {Question 89 C};
\node[anchor=east,text=white,font=\sffamily\itshape\large] at ([xshift=-5mm,yshift=-0.62cm]frame.north east) {Grandeurs et mesures};}]
\vspace*{1.75cm}
\begin{minipage}[t][6.9cm][c]{\linewidth}\LARGE\centering
\begin{minipage}{0.95\linewidth}\raggedright\linespread{1.25}\selectfont Une pyramide a une base d'aire $24$ cm$^2$ et une hauteur de $6$ cm. Quel est son volume~?\end{minipage}
\end{minipage}

\vfill
\begin{tcolorbox}[enhanced,nobeforeafter,width=7.1cm,height=2.1cm,valign=center,colback=cgra!6,colframe=cgra,arc=2.5mm,boxrule=0.9pt,left=3mm]{\Large\bfseries\color{cgra}a)}\quad{\Large $48~\text{cm}^3$}\end{tcolorbox}\hfill\begin{tcolorbox}[enhanced,nobeforeafter,width=7.1cm,height=2.1cm,valign=center,colback=cgra!6,colframe=cgra,arc=2.5mm,boxrule=0.9pt,left=3mm]{\Large\bfseries\color{cgra}b)}\quad{\Large $72~\text{cm}^3$}\end{tcolorbox}\par\vspace{2.5mm}
\begin{tcolorbox}[enhanced,nobeforeafter,width=7.1cm,height=2.1cm,valign=center,colback=cgra!6,colframe=cgra,arc=2.5mm,boxrule=0.9pt,left=3mm]{\Large\bfseries\color{cgra}c)}\quad{\Large $30~\text{cm}^3$}\end{tcolorbox}\hfill\begin{tcolorbox}[enhanced,nobeforeafter,width=7.1cm,height=2.1cm,valign=center,colback=cgra!6,colframe=cgra,arc=2.5mm,boxrule=0.9pt,left=3mm]{\Large\bfseries\color{cgra}d)}\quad{\Large $144~\text{cm}^3$}\end{tcolorbox}
\end{tcolorbox}
\newpage
\begin{tcolorbox}[enhanced,width=16cm,height=14cm,colback=white,colframe=cgra,boxrule=2pt,arc=0mm,
left=6mm,right=6mm,top=0mm,bottom=4mm,sharp corners,
overlay={\fill[cgra] (frame.north west) rectangle ([yshift=-1.25cm]frame.north east);
\node[anchor=west,text=white,font=\sffamily\bfseries\LARGE] at ([xshift=5mm,yshift=-0.62cm]frame.north west) {Question 90 C};
\node[anchor=east,text=white,font=\sffamily\itshape\large] at ([xshift=-5mm,yshift=-0.62cm]frame.north east) {Grandeurs et mesures};}]
\vspace*{1.75cm}
\begin{minipage}[t][6.9cm][c]{\linewidth}\LARGE\centering
\begin{minipage}{0.95\linewidth}\raggedright\linespread{1.25}\selectfont Une pyramide a une base d'aire $15$ cm$^2$ et une hauteur de $9$ cm. Quel est son volume~?\end{minipage}
\end{minipage}

\vfill
\begin{tcolorbox}[enhanced,nobeforeafter,width=7.1cm,height=2.1cm,valign=center,colback=cgra!6,colframe=cgra,arc=2.5mm,boxrule=0.9pt,left=3mm]{\Large\bfseries\color{cgra}a)}\quad{\Large $24~\text{cm}^3$}\end{tcolorbox}\hfill\begin{tcolorbox}[enhanced,nobeforeafter,width=7.1cm,height=2.1cm,valign=center,colback=cgra!6,colframe=cgra,arc=2.5mm,boxrule=0.9pt,left=3mm]{\Large\bfseries\color{cgra}b)}\quad{\Large $135~\text{cm}^3$}\end{tcolorbox}\par\vspace{2.5mm}
\begin{tcolorbox}[enhanced,nobeforeafter,width=7.1cm,height=2.1cm,valign=center,colback=cgra!6,colframe=cgra,arc=2.5mm,boxrule=0.9pt,left=3mm]{\Large\bfseries\color{cgra}c)}\quad{\Large $45~\text{cm}^3$}\end{tcolorbox}\hfill\begin{tcolorbox}[enhanced,nobeforeafter,width=7.1cm,height=2.1cm,valign=center,colback=cgra!6,colframe=cgra,arc=2.5mm,boxrule=0.9pt,left=3mm]{\Large\bfseries\color{cgra}d)}\quad{\Large $67{,}5~\text{cm}^3$}\end{tcolorbox}
\end{tcolorbox}
\newpage
\begin{tcolorbox}[enhanced,width=16cm,height=14cm,colback=white,colframe=cgra,boxrule=2pt,arc=0mm,
left=6mm,right=6mm,top=0mm,bottom=4mm,sharp corners,
overlay={\fill[cgra] (frame.north west) rectangle ([yshift=-1.25cm]frame.north east);
\node[anchor=west,text=white,font=\sffamily\bfseries\LARGE] at ([xshift=5mm,yshift=-0.62cm]frame.north west) {Question 91 C};
\node[anchor=east,text=white,font=\sffamily\itshape\large] at ([xshift=-5mm,yshift=-0.62cm]frame.north east) {Grandeurs et mesures};}]
\vspace*{1.75cm}
\begin{minipage}[t][6.9cm][c]{\linewidth}\LARGE\centering
\begin{minipage}{0.95\linewidth}\raggedright\linespread{1.25}\selectfont Un cône de révolution a un rayon de base de $4$ cm et une hauteur de $9$ cm. Quel est son volume exact~?\end{minipage}
\end{minipage}

\vfill
\begin{tcolorbox}[enhanced,nobeforeafter,width=7.1cm,height=2.1cm,valign=center,colback=cgra!6,colframe=cgra,arc=2.5mm,boxrule=0.9pt,left=3mm]{\Large\bfseries\color{cgra}a)}\quad{\Large $48\pi~\text{cm}^3$}\end{tcolorbox}\hfill\begin{tcolorbox}[enhanced,nobeforeafter,width=7.1cm,height=2.1cm,valign=center,colback=cgra!6,colframe=cgra,arc=2.5mm,boxrule=0.9pt,left=3mm]{\Large\bfseries\color{cgra}b)}\quad{\Large $144\pi~\text{cm}^3$}\end{tcolorbox}\par\vspace{2.5mm}
\begin{tcolorbox}[enhanced,nobeforeafter,width=7.1cm,height=2.1cm,valign=center,colback=cgra!6,colframe=cgra,arc=2.5mm,boxrule=0.9pt,left=3mm]{\Large\bfseries\color{cgra}c)}\quad{\Large $12\pi~\text{cm}^3$}\end{tcolorbox}\hfill\begin{tcolorbox}[enhanced,nobeforeafter,width=7.1cm,height=2.1cm,valign=center,colback=cgra!6,colframe=cgra,arc=2.5mm,boxrule=0.9pt,left=3mm]{\Large\bfseries\color{cgra}d)}\quad{\Large $192\pi~\text{cm}^3$}\end{tcolorbox}
\end{tcolorbox}
\newpage
\begin{tcolorbox}[enhanced,width=16cm,height=14cm,colback=white,colframe=cgra,boxrule=2pt,arc=0mm,
left=6mm,right=6mm,top=0mm,bottom=4mm,sharp corners,
overlay={\fill[cgra] (frame.north west) rectangle ([yshift=-1.25cm]frame.north east);
\node[anchor=west,text=white,font=\sffamily\bfseries\LARGE] at ([xshift=5mm,yshift=-0.62cm]frame.north west) {Question 92 C};
\node[anchor=east,text=white,font=\sffamily\itshape\large] at ([xshift=-5mm,yshift=-0.62cm]frame.north east) {Grandeurs et mesures};}]
\vspace*{1.75cm}
\begin{minipage}[t][6.9cm][c]{\linewidth}\LARGE\centering
\begin{minipage}{0.95\linewidth}\raggedright\linespread{1.25}\selectfont Un cône de révolution a un rayon de base de $3$ cm et une hauteur de $10$ cm. Quel est son volume exact~?\end{minipage}
\end{minipage}

\vfill
\begin{tcolorbox}[enhanced,nobeforeafter,width=7.1cm,height=2.1cm,valign=center,colback=cgra!6,colframe=cgra,arc=2.5mm,boxrule=0.9pt,left=3mm]{\Large\bfseries\color{cgra}a)}\quad{\Large $90\pi~\text{cm}^3$}\end{tcolorbox}\hfill\begin{tcolorbox}[enhanced,nobeforeafter,width=7.1cm,height=2.1cm,valign=center,colback=cgra!6,colframe=cgra,arc=2.5mm,boxrule=0.9pt,left=3mm]{\Large\bfseries\color{cgra}b)}\quad{\Large $120\pi~\text{cm}^3$}\end{tcolorbox}\par\vspace{2.5mm}
\begin{tcolorbox}[enhanced,nobeforeafter,width=7.1cm,height=2.1cm,valign=center,colback=cgra!6,colframe=cgra,arc=2.5mm,boxrule=0.9pt,left=3mm]{\Large\bfseries\color{cgra}c)}\quad{\Large $10\pi~\text{cm}^3$}\end{tcolorbox}\hfill\begin{tcolorbox}[enhanced,nobeforeafter,width=7.1cm,height=2.1cm,valign=center,colback=cgra!6,colframe=cgra,arc=2.5mm,boxrule=0.9pt,left=3mm]{\Large\bfseries\color{cgra}d)}\quad{\Large $30\pi~\text{cm}^3$}\end{tcolorbox}
\end{tcolorbox}
\newpage
\begin{tcolorbox}[enhanced,width=16cm,height=14cm,colback=white,colframe=cgra,boxrule=2pt,arc=0mm,
left=6mm,right=6mm,top=0mm,bottom=4mm,sharp corners,
overlay={\fill[cgra] (frame.north west) rectangle ([yshift=-1.25cm]frame.north east);
\node[anchor=west,text=white,font=\sffamily\bfseries\LARGE] at ([xshift=5mm,yshift=-0.62cm]frame.north west) {Question 93 C};
\node[anchor=east,text=white,font=\sffamily\itshape\large] at ([xshift=-5mm,yshift=-0.62cm]frame.north east) {Grandeurs et mesures};}]
\vspace*{1.75cm}
\begin{minipage}[t][6.9cm][c]{\linewidth}\LARGE\centering
\begin{minipage}{0.95\linewidth}\raggedright\linespread{1.25}\selectfont Compléter~: $0{,}2~\text{L}=\;?\;\text{cm}^3$\end{minipage}
\end{minipage}

\vfill
\begin{tcolorbox}[enhanced,nobeforeafter,width=7.1cm,height=2.1cm,valign=center,colback=cgra!6,colframe=cgra,arc=2.5mm,boxrule=0.9pt,left=3mm]{\Large\bfseries\color{cgra}a)}\quad{\Large $200~\text{cm}^3$}\end{tcolorbox}\hfill\begin{tcolorbox}[enhanced,nobeforeafter,width=7.1cm,height=2.1cm,valign=center,colback=cgra!6,colframe=cgra,arc=2.5mm,boxrule=0.9pt,left=3mm]{\Large\bfseries\color{cgra}b)}\quad{\Large $20~\text{cm}^3$}\end{tcolorbox}\par\vspace{2.5mm}
\begin{tcolorbox}[enhanced,nobeforeafter,width=7.1cm,height=2.1cm,valign=center,colback=cgra!6,colframe=cgra,arc=2.5mm,boxrule=0.9pt,left=3mm]{\Large\bfseries\color{cgra}c)}\quad{\Large $20\,000~\text{cm}^3$}\end{tcolorbox}\hfill\begin{tcolorbox}[enhanced,nobeforeafter,width=7.1cm,height=2.1cm,valign=center,colback=cgra!6,colframe=cgra,arc=2.5mm,boxrule=0.9pt,left=3mm]{\Large\bfseries\color{cgra}d)}\quad{\Large $2\,000~\text{cm}^3$}\end{tcolorbox}
\end{tcolorbox}
\newpage
\begin{tcolorbox}[enhanced,width=16cm,height=14cm,colback=white,colframe=cgra,boxrule=2pt,arc=0mm,
left=6mm,right=6mm,top=0mm,bottom=4mm,sharp corners,
overlay={\fill[cgra] (frame.north west) rectangle ([yshift=-1.25cm]frame.north east);
\node[anchor=west,text=white,font=\sffamily\bfseries\LARGE] at ([xshift=5mm,yshift=-0.62cm]frame.north west) {Question 94 C};
\node[anchor=east,text=white,font=\sffamily\itshape\large] at ([xshift=-5mm,yshift=-0.62cm]frame.north east) {Grandeurs et mesures};}]
\vspace*{1.75cm}
\begin{minipage}[t][6.9cm][c]{\linewidth}\LARGE\centering
\begin{minipage}{0.95\linewidth}\raggedright\linespread{1.25}\selectfont Compléter~: $3~\text{m}^3=\;?\;\text{L}$\end{minipage}
\end{minipage}

\vfill
\begin{tcolorbox}[enhanced,nobeforeafter,width=7.1cm,height=2.1cm,valign=center,colback=cgra!6,colframe=cgra,arc=2.5mm,boxrule=0.9pt,left=3mm]{\Large\bfseries\color{cgra}a)}\quad{\Large $3\,000~\text{L}$}\end{tcolorbox}\hfill\begin{tcolorbox}[enhanced,nobeforeafter,width=7.1cm,height=2.1cm,valign=center,colback=cgra!6,colframe=cgra,arc=2.5mm,boxrule=0.9pt,left=3mm]{\Large\bfseries\color{cgra}b)}\quad{\Large $300~\text{L}$}\end{tcolorbox}\par\vspace{2.5mm}
\begin{tcolorbox}[enhanced,nobeforeafter,width=7.1cm,height=2.1cm,valign=center,colback=cgra!6,colframe=cgra,arc=2.5mm,boxrule=0.9pt,left=3mm]{\Large\bfseries\color{cgra}c)}\quad{\Large $300\,000~\text{L}$}\end{tcolorbox}\hfill\begin{tcolorbox}[enhanced,nobeforeafter,width=7.1cm,height=2.1cm,valign=center,colback=cgra!6,colframe=cgra,arc=2.5mm,boxrule=0.9pt,left=3mm]{\Large\bfseries\color{cgra}d)}\quad{\Large $30\,000~\text{L}$}\end{tcolorbox}
\end{tcolorbox}
\newpage
\begin{tcolorbox}[enhanced,width=16cm,height=14cm,colback=white,colframe=cgra,boxrule=2pt,arc=0mm,
left=6mm,right=6mm,top=0mm,bottom=4mm,sharp corners,
overlay={\fill[cgra] (frame.north west) rectangle ([yshift=-1.25cm]frame.north east);
\node[anchor=west,text=white,font=\sffamily\bfseries\LARGE] at ([xshift=5mm,yshift=-0.62cm]frame.north west) {Question 95 C};
\node[anchor=east,text=white,font=\sffamily\itshape\large] at ([xshift=-5mm,yshift=-0.62cm]frame.north east) {Grandeurs et mesures};}]
\vspace*{1.75cm}
\begin{minipage}[t][6.9cm][c]{\linewidth}\LARGE\centering
\begin{minipage}{0.95\linewidth}\raggedright\linespread{1.25}\selectfont Compléter~: $0{,}8~\text{m}^2=\;?\;\text{dm}^2$\end{minipage}
\end{minipage}

\vfill
\begin{tcolorbox}[enhanced,nobeforeafter,width=7.1cm,height=2.1cm,valign=center,colback=cgra!6,colframe=cgra,arc=2.5mm,boxrule=0.9pt,left=3mm]{\Large\bfseries\color{cgra}a)}\quad{\Large $8~\text{dm}^2$}\end{tcolorbox}\hfill\begin{tcolorbox}[enhanced,nobeforeafter,width=7.1cm,height=2.1cm,valign=center,colback=cgra!6,colframe=cgra,arc=2.5mm,boxrule=0.9pt,left=3mm]{\Large\bfseries\color{cgra}b)}\quad{\Large $80~\text{dm}^2$}\end{tcolorbox}\par\vspace{2.5mm}
\begin{tcolorbox}[enhanced,nobeforeafter,width=7.1cm,height=2.1cm,valign=center,colback=cgra!6,colframe=cgra,arc=2.5mm,boxrule=0.9pt,left=3mm]{\Large\bfseries\color{cgra}c)}\quad{\Large $800~\text{dm}^2$}\end{tcolorbox}\hfill\begin{tcolorbox}[enhanced,nobeforeafter,width=7.1cm,height=2.1cm,valign=center,colback=cgra!6,colframe=cgra,arc=2.5mm,boxrule=0.9pt,left=3mm]{\Large\bfseries\color{cgra}d)}\quad{\Large $8\,000~\text{dm}^2$}\end{tcolorbox}
\end{tcolorbox}
\newpage
\begin{tcolorbox}[enhanced,width=16cm,height=14cm,colback=white,colframe=cgra,boxrule=2pt,arc=0mm,
left=6mm,right=6mm,top=0mm,bottom=4mm,sharp corners,
overlay={\fill[cgra] (frame.north west) rectangle ([yshift=-1.25cm]frame.north east);
\node[anchor=west,text=white,font=\sffamily\bfseries\LARGE] at ([xshift=5mm,yshift=-0.62cm]frame.north west) {Question 96 C};
\node[anchor=east,text=white,font=\sffamily\itshape\large] at ([xshift=-5mm,yshift=-0.62cm]frame.north east) {Grandeurs et mesures};}]
\vspace*{1.75cm}
\begin{minipage}[t][6.9cm][c]{\linewidth}\LARGE\centering
\begin{minipage}{0.95\linewidth}\raggedright\linespread{1.25}\selectfont Compléter~: $600~\text{mm}^2=\;?\;\text{cm}^2$\end{minipage}
\end{minipage}

\vfill
\begin{tcolorbox}[enhanced,nobeforeafter,width=7.1cm,height=2.1cm,valign=center,colback=cgra!6,colframe=cgra,arc=2.5mm,boxrule=0.9pt,left=3mm]{\Large\bfseries\color{cgra}a)}\quad{\Large $6~\text{cm}^2$}\end{tcolorbox}\hfill\begin{tcolorbox}[enhanced,nobeforeafter,width=7.1cm,height=2.1cm,valign=center,colback=cgra!6,colframe=cgra,arc=2.5mm,boxrule=0.9pt,left=3mm]{\Large\bfseries\color{cgra}b)}\quad{\Large $600~\text{cm}^2$}\end{tcolorbox}\par\vspace{2.5mm}
\begin{tcolorbox}[enhanced,nobeforeafter,width=7.1cm,height=2.1cm,valign=center,colback=cgra!6,colframe=cgra,arc=2.5mm,boxrule=0.9pt,left=3mm]{\Large\bfseries\color{cgra}c)}\quad{\Large $60~\text{cm}^2$}\end{tcolorbox}\hfill\begin{tcolorbox}[enhanced,nobeforeafter,width=7.1cm,height=2.1cm,valign=center,colback=cgra!6,colframe=cgra,arc=2.5mm,boxrule=0.9pt,left=3mm]{\Large\bfseries\color{cgra}d)}\quad{\Large $0{,}6~\text{cm}^2$}\end{tcolorbox}
\end{tcolorbox}
\newpage
\begin{tcolorbox}[enhanced,width=16cm,height=14cm,colback=white,colframe=crel,boxrule=2pt,arc=0mm,
left=6mm,right=6mm,top=0mm,bottom=4mm,sharp corners,
overlay={\fill[crel] (frame.north west) rectangle ([yshift=-1.25cm]frame.north east);
\node[anchor=west,text=white,font=\sffamily\bfseries\LARGE] at ([xshift=5mm,yshift=-0.62cm]frame.north west) {Question 1 D};
\node[anchor=east,text=white,font=\sffamily\itshape\large] at ([xshift=-5mm,yshift=-0.62cm]frame.north east) {Relatifs et fractions};}]
\vspace*{1.75cm}
\begin{minipage}[t][6.9cm][c]{\linewidth}\LARGE\centering
\begin{minipage}{0.95\linewidth}\raggedright\linespread{1.25}\selectfont Que vaut $(-7)\times(+2)\times(-1)$~?\end{minipage}
\end{minipage}

\vfill
\begin{tcolorbox}[enhanced,nobeforeafter,width=7.1cm,height=2.1cm,valign=center,colback=crel!6,colframe=crel,arc=2.5mm,boxrule=0.9pt,left=3mm]{\Large\bfseries\color{crel}a)}\quad{\Large $14$}\end{tcolorbox}\hfill\begin{tcolorbox}[enhanced,nobeforeafter,width=7.1cm,height=2.1cm,valign=center,colback=crel!6,colframe=crel,arc=2.5mm,boxrule=0.9pt,left=3mm]{\Large\bfseries\color{crel}b)}\quad{\Large $-6$}\end{tcolorbox}\par\vspace{2.5mm}
\begin{tcolorbox}[enhanced,nobeforeafter,width=7.1cm,height=2.1cm,valign=center,colback=crel!6,colframe=crel,arc=2.5mm,boxrule=0.9pt,left=3mm]{\Large\bfseries\color{crel}c)}\quad{\Large $-14$}\end{tcolorbox}\hfill\begin{tcolorbox}[enhanced,nobeforeafter,width=7.1cm,height=2.1cm,valign=center,colback=crel!6,colframe=crel,arc=2.5mm,boxrule=0.9pt,left=3mm]{\Large\bfseries\color{crel}d)}\quad{\Large $-13$}\end{tcolorbox}
\end{tcolorbox}
\newpage
\begin{tcolorbox}[enhanced,width=16cm,height=14cm,colback=white,colframe=crel,boxrule=2pt,arc=0mm,
left=6mm,right=6mm,top=0mm,bottom=4mm,sharp corners,
overlay={\fill[crel] (frame.north west) rectangle ([yshift=-1.25cm]frame.north east);
\node[anchor=west,text=white,font=\sffamily\bfseries\LARGE] at ([xshift=5mm,yshift=-0.62cm]frame.north west) {Question 2 D};
\node[anchor=east,text=white,font=\sffamily\itshape\large] at ([xshift=-5mm,yshift=-0.62cm]frame.north east) {Relatifs et fractions};}]
\vspace*{1.75cm}
\begin{minipage}[t][6.9cm][c]{\linewidth}\LARGE\centering
\begin{minipage}{0.95\linewidth}\raggedright\linespread{1.25}\selectfont Que vaut $(+2)\times(-5)\times(-3)$~?\end{minipage}
\end{minipage}

\vfill
\begin{tcolorbox}[enhanced,nobeforeafter,width=7.1cm,height=2.1cm,valign=center,colback=crel!6,colframe=crel,arc=2.5mm,boxrule=0.9pt,left=3mm]{\Large\bfseries\color{crel}a)}\quad{\Large $-30$}\end{tcolorbox}\hfill\begin{tcolorbox}[enhanced,nobeforeafter,width=7.1cm,height=2.1cm,valign=center,colback=crel!6,colframe=crel,arc=2.5mm,boxrule=0.9pt,left=3mm]{\Large\bfseries\color{crel}b)}\quad{\Large $-6$}\end{tcolorbox}\par\vspace{2.5mm}
\begin{tcolorbox}[enhanced,nobeforeafter,width=7.1cm,height=2.1cm,valign=center,colback=crel!6,colframe=crel,arc=2.5mm,boxrule=0.9pt,left=3mm]{\Large\bfseries\color{crel}c)}\quad{\Large $30$}\end{tcolorbox}\hfill\begin{tcolorbox}[enhanced,nobeforeafter,width=7.1cm,height=2.1cm,valign=center,colback=crel!6,colframe=crel,arc=2.5mm,boxrule=0.9pt,left=3mm]{\Large\bfseries\color{crel}d)}\quad{\Large $-7$}\end{tcolorbox}
\end{tcolorbox}
\newpage
\begin{tcolorbox}[enhanced,width=16cm,height=14cm,colback=white,colframe=crel,boxrule=2pt,arc=0mm,
left=6mm,right=6mm,top=0mm,bottom=4mm,sharp corners,
overlay={\fill[crel] (frame.north west) rectangle ([yshift=-1.25cm]frame.north east);
\node[anchor=west,text=white,font=\sffamily\bfseries\LARGE] at ([xshift=5mm,yshift=-0.62cm]frame.north west) {Question 3 D};
\node[anchor=east,text=white,font=\sffamily\itshape\large] at ([xshift=-5mm,yshift=-0.62cm]frame.north east) {Relatifs et fractions};}]
\vspace*{1.75cm}
\begin{minipage}[t][6.9cm][c]{\linewidth}\LARGE\centering
\begin{minipage}{0.95\linewidth}\raggedright\linespread{1.25}\selectfont Que vaut $(+72)\div(-8)$~?\end{minipage}
\end{minipage}

\vfill
\begin{tcolorbox}[enhanced,nobeforeafter,width=7.1cm,height=2.1cm,valign=center,colback=crel!6,colframe=crel,arc=2.5mm,boxrule=0.9pt,left=3mm]{\Large\bfseries\color{crel}a)}\quad{\Large $9$}\end{tcolorbox}\hfill\begin{tcolorbox}[enhanced,nobeforeafter,width=7.1cm,height=2.1cm,valign=center,colback=crel!6,colframe=crel,arc=2.5mm,boxrule=0.9pt,left=3mm]{\Large\bfseries\color{crel}b)}\quad{\Large $-9$}\end{tcolorbox}\par\vspace{2.5mm}
\begin{tcolorbox}[enhanced,nobeforeafter,width=7.1cm,height=2.1cm,valign=center,colback=crel!6,colframe=crel,arc=2.5mm,boxrule=0.9pt,left=3mm]{\Large\bfseries\color{crel}c)}\quad{\Large $-8$}\end{tcolorbox}\hfill\begin{tcolorbox}[enhanced,nobeforeafter,width=7.1cm,height=2.1cm,valign=center,colback=crel!6,colframe=crel,arc=2.5mm,boxrule=0.9pt,left=3mm]{\Large\bfseries\color{crel}d)}\quad{\Large $-18$}\end{tcolorbox}
\end{tcolorbox}
\newpage
\begin{tcolorbox}[enhanced,width=16cm,height=14cm,colback=white,colframe=crel,boxrule=2pt,arc=0mm,
left=6mm,right=6mm,top=0mm,bottom=4mm,sharp corners,
overlay={\fill[crel] (frame.north west) rectangle ([yshift=-1.25cm]frame.north east);
\node[anchor=west,text=white,font=\sffamily\bfseries\LARGE] at ([xshift=5mm,yshift=-0.62cm]frame.north west) {Question 4 D};
\node[anchor=east,text=white,font=\sffamily\itshape\large] at ([xshift=-5mm,yshift=-0.62cm]frame.north east) {Relatifs et fractions};}]
\vspace*{1.75cm}
\begin{minipage}[t][6.9cm][c]{\linewidth}\LARGE\centering
\begin{minipage}{0.95\linewidth}\raggedright\linespread{1.25}\selectfont Que vaut $(-35)\div(+5)$~?\end{minipage}
\end{minipage}

\vfill
\begin{tcolorbox}[enhanced,nobeforeafter,width=7.1cm,height=2.1cm,valign=center,colback=crel!6,colframe=crel,arc=2.5mm,boxrule=0.9pt,left=3mm]{\Large\bfseries\color{crel}a)}\quad{\Large $-6$}\end{tcolorbox}\hfill\begin{tcolorbox}[enhanced,nobeforeafter,width=7.1cm,height=2.1cm,valign=center,colback=crel!6,colframe=crel,arc=2.5mm,boxrule=0.9pt,left=3mm]{\Large\bfseries\color{crel}b)}\quad{\Large $7$}\end{tcolorbox}\par\vspace{2.5mm}
\begin{tcolorbox}[enhanced,nobeforeafter,width=7.1cm,height=2.1cm,valign=center,colback=crel!6,colframe=crel,arc=2.5mm,boxrule=0.9pt,left=3mm]{\Large\bfseries\color{crel}c)}\quad{\Large $-7$}\end{tcolorbox}\hfill\begin{tcolorbox}[enhanced,nobeforeafter,width=7.1cm,height=2.1cm,valign=center,colback=crel!6,colframe=crel,arc=2.5mm,boxrule=0.9pt,left=3mm]{\Large\bfseries\color{crel}d)}\quad{\Large $-14$}\end{tcolorbox}
\end{tcolorbox}
\newpage
\begin{tcolorbox}[enhanced,width=16cm,height=14cm,colback=white,colframe=crel,boxrule=2pt,arc=0mm,
left=6mm,right=6mm,top=0mm,bottom=4mm,sharp corners,
overlay={\fill[crel] (frame.north west) rectangle ([yshift=-1.25cm]frame.north east);
\node[anchor=west,text=white,font=\sffamily\bfseries\LARGE] at ([xshift=5mm,yshift=-0.62cm]frame.north west) {Question 5 D};
\node[anchor=east,text=white,font=\sffamily\itshape\large] at ([xshift=-5mm,yshift=-0.62cm]frame.north east) {Relatifs et fractions};}]
\vspace*{1.75cm}
\begin{minipage}[t][6.9cm][c]{\linewidth}\LARGE\centering
\begin{minipage}{0.95\linewidth}\raggedright\linespread{1.25}\selectfont Que vaut $9-2\times(-7)$~?\end{minipage}
\end{minipage}

\vfill
\begin{tcolorbox}[enhanced,nobeforeafter,width=7.1cm,height=2.1cm,valign=center,colback=crel!6,colframe=crel,arc=2.5mm,boxrule=0.9pt,left=3mm]{\Large\bfseries\color{crel}a)}\quad{\Large $-49$}\end{tcolorbox}\hfill\begin{tcolorbox}[enhanced,nobeforeafter,width=7.1cm,height=2.1cm,valign=center,colback=crel!6,colframe=crel,arc=2.5mm,boxrule=0.9pt,left=3mm]{\Large\bfseries\color{crel}b)}\quad{\Large $-5$}\end{tcolorbox}\par\vspace{2.5mm}
\begin{tcolorbox}[enhanced,nobeforeafter,width=7.1cm,height=2.1cm,valign=center,colback=crel!6,colframe=crel,arc=2.5mm,boxrule=0.9pt,left=3mm]{\Large\bfseries\color{crel}c)}\quad{\Large $23$}\end{tcolorbox}\hfill\begin{tcolorbox}[enhanced,nobeforeafter,width=7.1cm,height=2.1cm,valign=center,colback=crel!6,colframe=crel,arc=2.5mm,boxrule=0.9pt,left=3mm]{\Large\bfseries\color{crel}d)}\quad{\Large $-23$}\end{tcolorbox}
\end{tcolorbox}
\newpage
\begin{tcolorbox}[enhanced,width=16cm,height=14cm,colback=white,colframe=crel,boxrule=2pt,arc=0mm,
left=6mm,right=6mm,top=0mm,bottom=4mm,sharp corners,
overlay={\fill[crel] (frame.north west) rectangle ([yshift=-1.25cm]frame.north east);
\node[anchor=west,text=white,font=\sffamily\bfseries\LARGE] at ([xshift=5mm,yshift=-0.62cm]frame.north west) {Question 6 D};
\node[anchor=east,text=white,font=\sffamily\itshape\large] at ([xshift=-5mm,yshift=-0.62cm]frame.north east) {Relatifs et fractions};}]
\vspace*{1.75cm}
\begin{minipage}[t][6.9cm][c]{\linewidth}\LARGE\centering
\begin{minipage}{0.95\linewidth}\raggedright\linespread{1.25}\selectfont Que vaut $1-8\times(-3)$~?\end{minipage}
\end{minipage}

\vfill
\begin{tcolorbox}[enhanced,nobeforeafter,width=7.1cm,height=2.1cm,valign=center,colback=crel!6,colframe=crel,arc=2.5mm,boxrule=0.9pt,left=3mm]{\Large\bfseries\color{crel}a)}\quad{\Large $21$}\end{tcolorbox}\hfill\begin{tcolorbox}[enhanced,nobeforeafter,width=7.1cm,height=2.1cm,valign=center,colback=crel!6,colframe=crel,arc=2.5mm,boxrule=0.9pt,left=3mm]{\Large\bfseries\color{crel}b)}\quad{\Large $-25$}\end{tcolorbox}\par\vspace{2.5mm}
\begin{tcolorbox}[enhanced,nobeforeafter,width=7.1cm,height=2.1cm,valign=center,colback=crel!6,colframe=crel,arc=2.5mm,boxrule=0.9pt,left=3mm]{\Large\bfseries\color{crel}c)}\quad{\Large $25$}\end{tcolorbox}\hfill\begin{tcolorbox}[enhanced,nobeforeafter,width=7.1cm,height=2.1cm,valign=center,colback=crel!6,colframe=crel,arc=2.5mm,boxrule=0.9pt,left=3mm]{\Large\bfseries\color{crel}d)}\quad{\Large $-23$}\end{tcolorbox}
\end{tcolorbox}
\newpage
\begin{tcolorbox}[enhanced,width=16cm,height=14cm,colback=white,colframe=crel,boxrule=2pt,arc=0mm,
left=6mm,right=6mm,top=0mm,bottom=4mm,sharp corners,
overlay={\fill[crel] (frame.north west) rectangle ([yshift=-1.25cm]frame.north east);
\node[anchor=west,text=white,font=\sffamily\bfseries\LARGE] at ([xshift=5mm,yshift=-0.62cm]frame.north west) {Question 7 D};
\node[anchor=east,text=white,font=\sffamily\itshape\large] at ([xshift=-5mm,yshift=-0.62cm]frame.north east) {Relatifs et fractions};}]
\vspace*{1.75cm}
\begin{minipage}[t][6.9cm][c]{\linewidth}\LARGE\centering
\begin{minipage}{0.95\linewidth}\raggedright\linespread{1.25}\selectfont Que vaut $\dfrac{3}{10}+\dfrac{2}{15}$~?\end{minipage}
\end{minipage}

\vfill
\begin{tcolorbox}[enhanced,nobeforeafter,width=7.1cm,height=2.1cm,valign=center,colback=crel!6,colframe=crel,arc=2.5mm,boxrule=0.9pt,left=3mm]{\Large\bfseries\color{crel}a)}\quad{\Large $\dfrac{1}{30}$}\end{tcolorbox}\hfill\begin{tcolorbox}[enhanced,nobeforeafter,width=7.1cm,height=2.1cm,valign=center,colback=crel!6,colframe=crel,arc=2.5mm,boxrule=0.9pt,left=3mm]{\Large\bfseries\color{crel}b)}\quad{\Large $\dfrac{13}{30}$}\end{tcolorbox}\par\vspace{2.5mm}
\begin{tcolorbox}[enhanced,nobeforeafter,width=7.1cm,height=2.1cm,valign=center,colback=crel!6,colframe=crel,arc=2.5mm,boxrule=0.9pt,left=3mm]{\Large\bfseries\color{crel}c)}\quad{\Large $\dfrac{1}{25}$}\end{tcolorbox}\hfill\begin{tcolorbox}[enhanced,nobeforeafter,width=7.1cm,height=2.1cm,valign=center,colback=crel!6,colframe=crel,arc=2.5mm,boxrule=0.9pt,left=3mm]{\Large\bfseries\color{crel}d)}\quad{\Large $\dfrac{1}{5}$}\end{tcolorbox}
\end{tcolorbox}
\newpage
\begin{tcolorbox}[enhanced,width=16cm,height=14cm,colback=white,colframe=crel,boxrule=2pt,arc=0mm,
left=6mm,right=6mm,top=0mm,bottom=4mm,sharp corners,
overlay={\fill[crel] (frame.north west) rectangle ([yshift=-1.25cm]frame.north east);
\node[anchor=west,text=white,font=\sffamily\bfseries\LARGE] at ([xshift=5mm,yshift=-0.62cm]frame.north west) {Question 8 D};
\node[anchor=east,text=white,font=\sffamily\itshape\large] at ([xshift=-5mm,yshift=-0.62cm]frame.north east) {Relatifs et fractions};}]
\vspace*{1.75cm}
\begin{minipage}[t][6.9cm][c]{\linewidth}\LARGE\centering
\begin{minipage}{0.95\linewidth}\raggedright\linespread{1.25}\selectfont Que vaut $\dfrac{1}{4}+\dfrac{5}{6}$~?\end{minipage}
\end{minipage}

\vfill
\begin{tcolorbox}[enhanced,nobeforeafter,width=7.1cm,height=2.1cm,valign=center,colback=crel!6,colframe=crel,arc=2.5mm,boxrule=0.9pt,left=3mm]{\Large\bfseries\color{crel}a)}\quad{\Large $\dfrac{3}{5}$}\end{tcolorbox}\hfill\begin{tcolorbox}[enhanced,nobeforeafter,width=7.1cm,height=2.1cm,valign=center,colback=crel!6,colframe=crel,arc=2.5mm,boxrule=0.9pt,left=3mm]{\Large\bfseries\color{crel}b)}\quad{\Large $\dfrac{5}{24}$}\end{tcolorbox}\par\vspace{2.5mm}
\begin{tcolorbox}[enhanced,nobeforeafter,width=7.1cm,height=2.1cm,valign=center,colback=crel!6,colframe=crel,arc=2.5mm,boxrule=0.9pt,left=3mm]{\Large\bfseries\color{crel}c)}\quad{\Large $\dfrac{1}{4}$}\end{tcolorbox}\hfill\begin{tcolorbox}[enhanced,nobeforeafter,width=7.1cm,height=2.1cm,valign=center,colback=crel!6,colframe=crel,arc=2.5mm,boxrule=0.9pt,left=3mm]{\Large\bfseries\color{crel}d)}\quad{\Large $\dfrac{13}{12}$}\end{tcolorbox}
\end{tcolorbox}
\newpage
\begin{tcolorbox}[enhanced,width=16cm,height=14cm,colback=white,colframe=crel,boxrule=2pt,arc=0mm,
left=6mm,right=6mm,top=0mm,bottom=4mm,sharp corners,
overlay={\fill[crel] (frame.north west) rectangle ([yshift=-1.25cm]frame.north east);
\node[anchor=west,text=white,font=\sffamily\bfseries\LARGE] at ([xshift=5mm,yshift=-0.62cm]frame.north west) {Question 9 D};
\node[anchor=east,text=white,font=\sffamily\itshape\large] at ([xshift=-5mm,yshift=-0.62cm]frame.north east) {Relatifs et fractions};}]
\vspace*{1.75cm}
\begin{minipage}[t][6.9cm][c]{\linewidth}\LARGE\centering
\begin{minipage}{0.95\linewidth}\raggedright\linespread{1.25}\selectfont Que vaut $\dfrac{5}{4}\times\left(-\dfrac{8}{15}\right)$~?\end{minipage}
\end{minipage}

\vfill
\begin{tcolorbox}[enhanced,nobeforeafter,width=7.1cm,height=2.1cm,valign=center,colback=crel!6,colframe=crel,arc=2.5mm,boxrule=0.9pt,left=3mm]{\Large\bfseries\color{crel}a)}\quad{\Large $-\dfrac{2}{3}$}\end{tcolorbox}\hfill\begin{tcolorbox}[enhanced,nobeforeafter,width=7.1cm,height=2.1cm,valign=center,colback=crel!6,colframe=crel,arc=2.5mm,boxrule=0.9pt,left=3mm]{\Large\bfseries\color{crel}b)}\quad{\Large $-\dfrac{3}{19}$}\end{tcolorbox}\par\vspace{2.5mm}
\begin{tcolorbox}[enhanced,nobeforeafter,width=7.1cm,height=2.1cm,valign=center,colback=crel!6,colframe=crel,arc=2.5mm,boxrule=0.9pt,left=3mm]{\Large\bfseries\color{crel}c)}\quad{\Large $\dfrac{2}{3}$}\end{tcolorbox}\hfill\begin{tcolorbox}[enhanced,nobeforeafter,width=7.1cm,height=2.1cm,valign=center,colback=crel!6,colframe=crel,arc=2.5mm,boxrule=0.9pt,left=3mm]{\Large\bfseries\color{crel}d)}\quad{\Large $-\dfrac{75}{32}$}\end{tcolorbox}
\end{tcolorbox}
\newpage
\begin{tcolorbox}[enhanced,width=16cm,height=14cm,colback=white,colframe=crel,boxrule=2pt,arc=0mm,
left=6mm,right=6mm,top=0mm,bottom=4mm,sharp corners,
overlay={\fill[crel] (frame.north west) rectangle ([yshift=-1.25cm]frame.north east);
\node[anchor=west,text=white,font=\sffamily\bfseries\LARGE] at ([xshift=5mm,yshift=-0.62cm]frame.north west) {Question 10 D};
\node[anchor=east,text=white,font=\sffamily\itshape\large] at ([xshift=-5mm,yshift=-0.62cm]frame.north east) {Relatifs et fractions};}]
\vspace*{1.75cm}
\begin{minipage}[t][6.9cm][c]{\linewidth}\LARGE\centering
\begin{minipage}{0.95\linewidth}\raggedright\linespread{1.25}\selectfont Que vaut $\dfrac{3}{8}\times\left(-\dfrac{4}{9}\right)$~?\end{minipage}
\end{minipage}

\vfill
\begin{tcolorbox}[enhanced,nobeforeafter,width=7.1cm,height=2.1cm,valign=center,colback=crel!6,colframe=crel,arc=2.5mm,boxrule=0.9pt,left=3mm]{\Large\bfseries\color{crel}a)}\quad{\Large $-\dfrac{1}{17}$}\end{tcolorbox}\hfill\begin{tcolorbox}[enhanced,nobeforeafter,width=7.1cm,height=2.1cm,valign=center,colback=crel!6,colframe=crel,arc=2.5mm,boxrule=0.9pt,left=3mm]{\Large\bfseries\color{crel}b)}\quad{\Large $-\dfrac{1}{6}$}\end{tcolorbox}\par\vspace{2.5mm}
\begin{tcolorbox}[enhanced,nobeforeafter,width=7.1cm,height=2.1cm,valign=center,colback=crel!6,colframe=crel,arc=2.5mm,boxrule=0.9pt,left=3mm]{\Large\bfseries\color{crel}c)}\quad{\Large $-\dfrac{27}{32}$}\end{tcolorbox}\hfill\begin{tcolorbox}[enhanced,nobeforeafter,width=7.1cm,height=2.1cm,valign=center,colback=crel!6,colframe=crel,arc=2.5mm,boxrule=0.9pt,left=3mm]{\Large\bfseries\color{crel}d)}\quad{\Large $\dfrac{1}{6}$}\end{tcolorbox}
\end{tcolorbox}
\newpage
\begin{tcolorbox}[enhanced,width=16cm,height=14cm,colback=white,colframe=crel,boxrule=2pt,arc=0mm,
left=6mm,right=6mm,top=0mm,bottom=4mm,sharp corners,
overlay={\fill[crel] (frame.north west) rectangle ([yshift=-1.25cm]frame.north east);
\node[anchor=west,text=white,font=\sffamily\bfseries\LARGE] at ([xshift=5mm,yshift=-0.62cm]frame.north west) {Question 11 D};
\node[anchor=east,text=white,font=\sffamily\itshape\large] at ([xshift=-5mm,yshift=-0.62cm]frame.north east) {Relatifs et fractions};}]
\vspace*{1.75cm}
\begin{minipage}[t][6.9cm][c]{\linewidth}\LARGE\centering
\begin{minipage}{0.95\linewidth}\raggedright\linespread{1.25}\selectfont Que vaut $\dfrac{5}{8}\div\dfrac{15}{4}$~?\end{minipage}
\end{minipage}

\vfill
\begin{tcolorbox}[enhanced,nobeforeafter,width=7.1cm,height=2.1cm,valign=center,colback=crel!6,colframe=crel,arc=2.5mm,boxrule=0.9pt,left=3mm]{\Large\bfseries\color{crel}a)}\quad{\Large $\dfrac{1}{6}$}\end{tcolorbox}\hfill\begin{tcolorbox}[enhanced,nobeforeafter,width=7.1cm,height=2.1cm,valign=center,colback=crel!6,colframe=crel,arc=2.5mm,boxrule=0.9pt,left=3mm]{\Large\bfseries\color{crel}b)}\quad{\Large $\dfrac{75}{32}$}\end{tcolorbox}\par\vspace{2.5mm}
\begin{tcolorbox}[enhanced,nobeforeafter,width=7.1cm,height=2.1cm,valign=center,colback=crel!6,colframe=crel,arc=2.5mm,boxrule=0.9pt,left=3mm]{\Large\bfseries\color{crel}c)}\quad{\Large $6$}\end{tcolorbox}\hfill\begin{tcolorbox}[enhanced,nobeforeafter,width=7.1cm,height=2.1cm,valign=center,colback=crel!6,colframe=crel,arc=2.5mm,boxrule=0.9pt,left=3mm]{\Large\bfseries\color{crel}d)}\quad{\Large $-\dfrac{1}{6}$}\end{tcolorbox}
\end{tcolorbox}
\newpage
\begin{tcolorbox}[enhanced,width=16cm,height=14cm,colback=white,colframe=crel,boxrule=2pt,arc=0mm,
left=6mm,right=6mm,top=0mm,bottom=4mm,sharp corners,
overlay={\fill[crel] (frame.north west) rectangle ([yshift=-1.25cm]frame.north east);
\node[anchor=west,text=white,font=\sffamily\bfseries\LARGE] at ([xshift=5mm,yshift=-0.62cm]frame.north west) {Question 12 D};
\node[anchor=east,text=white,font=\sffamily\itshape\large] at ([xshift=-5mm,yshift=-0.62cm]frame.north east) {Relatifs et fractions};}]
\vspace*{1.75cm}
\begin{minipage}[t][6.9cm][c]{\linewidth}\LARGE\centering
\begin{minipage}{0.95\linewidth}\raggedright\linespread{1.25}\selectfont Que vaut $\dfrac{9}{10}\div\dfrac{3}{5}$~?\end{minipage}
\end{minipage}

\vfill
\begin{tcolorbox}[enhanced,nobeforeafter,width=7.1cm,height=2.1cm,valign=center,colback=crel!6,colframe=crel,arc=2.5mm,boxrule=0.9pt,left=3mm]{\Large\bfseries\color{crel}a)}\quad{\Large $\dfrac{2}{3}$}\end{tcolorbox}\hfill\begin{tcolorbox}[enhanced,nobeforeafter,width=7.1cm,height=2.1cm,valign=center,colback=crel!6,colframe=crel,arc=2.5mm,boxrule=0.9pt,left=3mm]{\Large\bfseries\color{crel}b)}\quad{\Large $\dfrac{3}{2}$}\end{tcolorbox}\par\vspace{2.5mm}
\begin{tcolorbox}[enhanced,nobeforeafter,width=7.1cm,height=2.1cm,valign=center,colback=crel!6,colframe=crel,arc=2.5mm,boxrule=0.9pt,left=3mm]{\Large\bfseries\color{crel}c)}\quad{\Large $\dfrac{27}{50}$}\end{tcolorbox}\hfill\begin{tcolorbox}[enhanced,nobeforeafter,width=7.1cm,height=2.1cm,valign=center,colback=crel!6,colframe=crel,arc=2.5mm,boxrule=0.9pt,left=3mm]{\Large\bfseries\color{crel}d)}\quad{\Large $-\dfrac{3}{2}$}\end{tcolorbox}
\end{tcolorbox}
\newpage
\begin{tcolorbox}[enhanced,width=16cm,height=14cm,colback=white,colframe=crel,boxrule=2pt,arc=0mm,
left=6mm,right=6mm,top=0mm,bottom=4mm,sharp corners,
overlay={\fill[crel] (frame.north west) rectangle ([yshift=-1.25cm]frame.north east);
\node[anchor=west,text=white,font=\sffamily\bfseries\LARGE] at ([xshift=5mm,yshift=-0.62cm]frame.north west) {Question 13 D};
\node[anchor=east,text=white,font=\sffamily\itshape\large] at ([xshift=-5mm,yshift=-0.62cm]frame.north east) {Relatifs et fractions};}]
\vspace*{1.75cm}
\begin{minipage}[t][6.9cm][c]{\linewidth}\LARGE\centering
\begin{minipage}{0.95\linewidth}\raggedright\linespread{1.25}\selectfont Quel est l'inverse de $-\dfrac{7}{2}$~?\end{minipage}
\end{minipage}

\vfill
\begin{tcolorbox}[enhanced,nobeforeafter,width=7.1cm,height=2.1cm,valign=center,colback=crel!6,colframe=crel,arc=2.5mm,boxrule=0.9pt,left=3mm]{\Large\bfseries\color{crel}a)}\quad{\Large $\dfrac{7}{2}$}\end{tcolorbox}\hfill\begin{tcolorbox}[enhanced,nobeforeafter,width=7.1cm,height=2.1cm,valign=center,colback=crel!6,colframe=crel,arc=2.5mm,boxrule=0.9pt,left=3mm]{\Large\bfseries\color{crel}b)}\quad{\Large $-\dfrac{7}{2}$}\end{tcolorbox}\par\vspace{2.5mm}
\begin{tcolorbox}[enhanced,nobeforeafter,width=7.1cm,height=2.1cm,valign=center,colback=crel!6,colframe=crel,arc=2.5mm,boxrule=0.9pt,left=3mm]{\Large\bfseries\color{crel}c)}\quad{\Large $\dfrac{2}{7}$}\end{tcolorbox}\hfill\begin{tcolorbox}[enhanced,nobeforeafter,width=7.1cm,height=2.1cm,valign=center,colback=crel!6,colframe=crel,arc=2.5mm,boxrule=0.9pt,left=3mm]{\Large\bfseries\color{crel}d)}\quad{\Large $-\dfrac{2}{7}$}\end{tcolorbox}
\end{tcolorbox}
\newpage
\begin{tcolorbox}[enhanced,width=16cm,height=14cm,colback=white,colframe=crel,boxrule=2pt,arc=0mm,
left=6mm,right=6mm,top=0mm,bottom=4mm,sharp corners,
overlay={\fill[crel] (frame.north west) rectangle ([yshift=-1.25cm]frame.north east);
\node[anchor=west,text=white,font=\sffamily\bfseries\LARGE] at ([xshift=5mm,yshift=-0.62cm]frame.north west) {Question 14 D};
\node[anchor=east,text=white,font=\sffamily\itshape\large] at ([xshift=-5mm,yshift=-0.62cm]frame.north east) {Relatifs et fractions};}]
\vspace*{1.75cm}
\begin{minipage}[t][6.9cm][c]{\linewidth}\LARGE\centering
\begin{minipage}{0.95\linewidth}\raggedright\linespread{1.25}\selectfont Quel est l'inverse de $\dfrac{2}{9}$~?\end{minipage}
\end{minipage}

\vfill
\begin{tcolorbox}[enhanced,nobeforeafter,width=7.1cm,height=2.1cm,valign=center,colback=crel!6,colframe=crel,arc=2.5mm,boxrule=0.9pt,left=3mm]{\Large\bfseries\color{crel}a)}\quad{\Large $\dfrac{2}{9}$}\end{tcolorbox}\hfill\begin{tcolorbox}[enhanced,nobeforeafter,width=7.1cm,height=2.1cm,valign=center,colback=crel!6,colframe=crel,arc=2.5mm,boxrule=0.9pt,left=3mm]{\Large\bfseries\color{crel}b)}\quad{\Large $-\dfrac{9}{2}$}\end{tcolorbox}\par\vspace{2.5mm}
\begin{tcolorbox}[enhanced,nobeforeafter,width=7.1cm,height=2.1cm,valign=center,colback=crel!6,colframe=crel,arc=2.5mm,boxrule=0.9pt,left=3mm]{\Large\bfseries\color{crel}c)}\quad{\Large $-\dfrac{2}{9}$}\end{tcolorbox}\hfill\begin{tcolorbox}[enhanced,nobeforeafter,width=7.1cm,height=2.1cm,valign=center,colback=crel!6,colframe=crel,arc=2.5mm,boxrule=0.9pt,left=3mm]{\Large\bfseries\color{crel}d)}\quad{\Large $\dfrac{9}{2}$}\end{tcolorbox}
\end{tcolorbox}
\newpage
\begin{tcolorbox}[enhanced,width=16cm,height=14cm,colback=white,colframe=crel,boxrule=2pt,arc=0mm,
left=6mm,right=6mm,top=0mm,bottom=4mm,sharp corners,
overlay={\fill[crel] (frame.north west) rectangle ([yshift=-1.25cm]frame.north east);
\node[anchor=west,text=white,font=\sffamily\bfseries\LARGE] at ([xshift=5mm,yshift=-0.62cm]frame.north west) {Question 15 D};
\node[anchor=east,text=white,font=\sffamily\itshape\large] at ([xshift=-5mm,yshift=-0.62cm]frame.north east) {Relatifs et fractions};}]
\vspace*{1.75cm}
\begin{minipage}[t][6.9cm][c]{\linewidth}\LARGE\centering
\begin{minipage}{0.95\linewidth}\raggedright\linespread{1.25}\selectfont Que vaut $\dfrac{2}{9}-\dfrac{1}{3}$~?\end{minipage}
\end{minipage}

\vfill
\begin{tcolorbox}[enhanced,nobeforeafter,width=7.1cm,height=2.1cm,valign=center,colback=crel!6,colframe=crel,arc=2.5mm,boxrule=0.9pt,left=3mm]{\Large\bfseries\color{crel}a)}\quad{\Large $\dfrac{1}{9}$}\end{tcolorbox}\hfill\begin{tcolorbox}[enhanced,nobeforeafter,width=7.1cm,height=2.1cm,valign=center,colback=crel!6,colframe=crel,arc=2.5mm,boxrule=0.9pt,left=3mm]{\Large\bfseries\color{crel}b)}\quad{\Large $-\dfrac{1}{9}$}\end{tcolorbox}\par\vspace{2.5mm}
\begin{tcolorbox}[enhanced,nobeforeafter,width=7.1cm,height=2.1cm,valign=center,colback=crel!6,colframe=crel,arc=2.5mm,boxrule=0.9pt,left=3mm]{\Large\bfseries\color{crel}c)}\quad{\Large $\dfrac{1}{27}$}\end{tcolorbox}\hfill\begin{tcolorbox}[enhanced,nobeforeafter,width=7.1cm,height=2.1cm,valign=center,colback=crel!6,colframe=crel,arc=2.5mm,boxrule=0.9pt,left=3mm]{\Large\bfseries\color{crel}d)}\quad{\Large $\dfrac{1}{6}$}\end{tcolorbox}
\end{tcolorbox}
\newpage
\begin{tcolorbox}[enhanced,width=16cm,height=14cm,colback=white,colframe=crel,boxrule=2pt,arc=0mm,
left=6mm,right=6mm,top=0mm,bottom=4mm,sharp corners,
overlay={\fill[crel] (frame.north west) rectangle ([yshift=-1.25cm]frame.north east);
\node[anchor=west,text=white,font=\sffamily\bfseries\LARGE] at ([xshift=5mm,yshift=-0.62cm]frame.north west) {Question 16 D};
\node[anchor=east,text=white,font=\sffamily\itshape\large] at ([xshift=-5mm,yshift=-0.62cm]frame.north east) {Relatifs et fractions};}]
\vspace*{1.75cm}
\begin{minipage}[t][6.9cm][c]{\linewidth}\LARGE\centering
\begin{minipage}{0.95\linewidth}\raggedright\linespread{1.25}\selectfont Que vaut $\dfrac{3}{4}-\dfrac{9}{10}$~?\end{minipage}
\end{minipage}

\vfill
\begin{tcolorbox}[enhanced,nobeforeafter,width=7.1cm,height=2.1cm,valign=center,colback=crel!6,colframe=crel,arc=2.5mm,boxrule=0.9pt,left=3mm]{\Large\bfseries\color{crel}a)}\quad{\Large $-1$}\end{tcolorbox}\hfill\begin{tcolorbox}[enhanced,nobeforeafter,width=7.1cm,height=2.1cm,valign=center,colback=crel!6,colframe=crel,arc=2.5mm,boxrule=0.9pt,left=3mm]{\Large\bfseries\color{crel}b)}\quad{\Large $\dfrac{3}{20}$}\end{tcolorbox}\par\vspace{2.5mm}
\begin{tcolorbox}[enhanced,nobeforeafter,width=7.1cm,height=2.1cm,valign=center,colback=crel!6,colframe=crel,arc=2.5mm,boxrule=0.9pt,left=3mm]{\Large\bfseries\color{crel}c)}\quad{\Large $-\dfrac{3}{20}$}\end{tcolorbox}\hfill\begin{tcolorbox}[enhanced,nobeforeafter,width=7.1cm,height=2.1cm,valign=center,colback=crel!6,colframe=crel,arc=2.5mm,boxrule=0.9pt,left=3mm]{\Large\bfseries\color{crel}d)}\quad{\Large $\dfrac{3}{10}$}\end{tcolorbox}
\end{tcolorbox}
\newpage
\begin{tcolorbox}[enhanced,width=16cm,height=14cm,colback=white,colframe=cpui,boxrule=2pt,arc=0mm,
left=6mm,right=6mm,top=0mm,bottom=4mm,sharp corners,
overlay={\fill[cpui] (frame.north west) rectangle ([yshift=-1.25cm]frame.north east);
\node[anchor=west,text=white,font=\sffamily\bfseries\LARGE] at ([xshift=5mm,yshift=-0.62cm]frame.north west) {Question 17 D};
\node[anchor=east,text=white,font=\sffamily\itshape\large] at ([xshift=-5mm,yshift=-0.62cm]frame.north east) {Puissances et racines carrées};}]
\vspace*{1.75cm}
\begin{minipage}[t][6.9cm][c]{\linewidth}\LARGE\centering
\begin{minipage}{0.95\linewidth}\raggedright\linespread{1.25}\selectfont Que vaut $10^{4}$~?\end{minipage}
\end{minipage}

\vfill
\begin{tcolorbox}[enhanced,nobeforeafter,width=7.1cm,height=2.1cm,valign=center,colback=cpui!6,colframe=cpui,arc=2.5mm,boxrule=0.9pt,left=3mm]{\Large\bfseries\color{cpui}a)}\quad{\Large $40$}\end{tcolorbox}\hfill\begin{tcolorbox}[enhanced,nobeforeafter,width=7.1cm,height=2.1cm,valign=center,colback=cpui!6,colframe=cpui,arc=2.5mm,boxrule=0.9pt,left=3mm]{\Large\bfseries\color{cpui}b)}\quad{\Large $10\,000$}\end{tcolorbox}\par\vspace{2.5mm}
\begin{tcolorbox}[enhanced,nobeforeafter,width=7.1cm,height=2.1cm,valign=center,colback=cpui!6,colframe=cpui,arc=2.5mm,boxrule=0.9pt,left=3mm]{\Large\bfseries\color{cpui}c)}\quad{\Large $-10\,000$}\end{tcolorbox}\hfill\begin{tcolorbox}[enhanced,nobeforeafter,width=7.1cm,height=2.1cm,valign=center,colback=cpui!6,colframe=cpui,arc=2.5mm,boxrule=0.9pt,left=3mm]{\Large\bfseries\color{cpui}d)}\quad{\Large $1\,000$}\end{tcolorbox}
\end{tcolorbox}
\newpage
\begin{tcolorbox}[enhanced,width=16cm,height=14cm,colback=white,colframe=cpui,boxrule=2pt,arc=0mm,
left=6mm,right=6mm,top=0mm,bottom=4mm,sharp corners,
overlay={\fill[cpui] (frame.north west) rectangle ([yshift=-1.25cm]frame.north east);
\node[anchor=west,text=white,font=\sffamily\bfseries\LARGE] at ([xshift=5mm,yshift=-0.62cm]frame.north west) {Question 18 D};
\node[anchor=east,text=white,font=\sffamily\itshape\large] at ([xshift=-5mm,yshift=-0.62cm]frame.north east) {Puissances et racines carrées};}]
\vspace*{1.75cm}
\begin{minipage}[t][6.9cm][c]{\linewidth}\LARGE\centering
\begin{minipage}{0.95\linewidth}\raggedright\linespread{1.25}\selectfont Que vaut $4^{3}$~?\end{minipage}
\end{minipage}

\vfill
\begin{tcolorbox}[enhanced,nobeforeafter,width=7.1cm,height=2.1cm,valign=center,colback=cpui!6,colframe=cpui,arc=2.5mm,boxrule=0.9pt,left=3mm]{\Large\bfseries\color{cpui}a)}\quad{\Large $-64$}\end{tcolorbox}\hfill\begin{tcolorbox}[enhanced,nobeforeafter,width=7.1cm,height=2.1cm,valign=center,colback=cpui!6,colframe=cpui,arc=2.5mm,boxrule=0.9pt,left=3mm]{\Large\bfseries\color{cpui}b)}\quad{\Large $16$}\end{tcolorbox}\par\vspace{2.5mm}
\begin{tcolorbox}[enhanced,nobeforeafter,width=7.1cm,height=2.1cm,valign=center,colback=cpui!6,colframe=cpui,arc=2.5mm,boxrule=0.9pt,left=3mm]{\Large\bfseries\color{cpui}c)}\quad{\Large $12$}\end{tcolorbox}\hfill\begin{tcolorbox}[enhanced,nobeforeafter,width=7.1cm,height=2.1cm,valign=center,colback=cpui!6,colframe=cpui,arc=2.5mm,boxrule=0.9pt,left=3mm]{\Large\bfseries\color{cpui}d)}\quad{\Large $64$}\end{tcolorbox}
\end{tcolorbox}
\newpage
\begin{tcolorbox}[enhanced,width=16cm,height=14cm,colback=white,colframe=cpui,boxrule=2pt,arc=0mm,
left=6mm,right=6mm,top=0mm,bottom=4mm,sharp corners,
overlay={\fill[cpui] (frame.north west) rectangle ([yshift=-1.25cm]frame.north east);
\node[anchor=west,text=white,font=\sffamily\bfseries\LARGE] at ([xshift=5mm,yshift=-0.62cm]frame.north west) {Question 19 D};
\node[anchor=east,text=white,font=\sffamily\itshape\large] at ([xshift=-5mm,yshift=-0.62cm]frame.north east) {Puissances et racines carrées};}]
\vspace*{1.75cm}
\begin{minipage}[t][6.9cm][c]{\linewidth}\LARGE\centering
\begin{minipage}{0.95\linewidth}\raggedright\linespread{1.25}\selectfont Écrire $10^{4}\times10^{4}$ sous la forme d'une seule puissance de $10$.\end{minipage}
\end{minipage}

\vfill
\begin{tcolorbox}[enhanced,nobeforeafter,width=7.1cm,height=2.1cm,valign=center,colback=cpui!6,colframe=cpui,arc=2.5mm,boxrule=0.9pt,left=3mm]{\Large\bfseries\color{cpui}a)}\quad{\Large $100^{8}$}\end{tcolorbox}\hfill\begin{tcolorbox}[enhanced,nobeforeafter,width=7.1cm,height=2.1cm,valign=center,colback=cpui!6,colframe=cpui,arc=2.5mm,boxrule=0.9pt,left=3mm]{\Large\bfseries\color{cpui}b)}\quad{\Large $10^{0}$}\end{tcolorbox}\par\vspace{2.5mm}
\begin{tcolorbox}[enhanced,nobeforeafter,width=7.1cm,height=2.1cm,valign=center,colback=cpui!6,colframe=cpui,arc=2.5mm,boxrule=0.9pt,left=3mm]{\Large\bfseries\color{cpui}c)}\quad{\Large $10^{8}$}\end{tcolorbox}\hfill\begin{tcolorbox}[enhanced,nobeforeafter,width=7.1cm,height=2.1cm,valign=center,colback=cpui!6,colframe=cpui,arc=2.5mm,boxrule=0.9pt,left=3mm]{\Large\bfseries\color{cpui}d)}\quad{\Large $10^{16}$}\end{tcolorbox}
\end{tcolorbox}
\newpage
\begin{tcolorbox}[enhanced,width=16cm,height=14cm,colback=white,colframe=cpui,boxrule=2pt,arc=0mm,
left=6mm,right=6mm,top=0mm,bottom=4mm,sharp corners,
overlay={\fill[cpui] (frame.north west) rectangle ([yshift=-1.25cm]frame.north east);
\node[anchor=west,text=white,font=\sffamily\bfseries\LARGE] at ([xshift=5mm,yshift=-0.62cm]frame.north west) {Question 20 D};
\node[anchor=east,text=white,font=\sffamily\itshape\large] at ([xshift=-5mm,yshift=-0.62cm]frame.north east) {Puissances et racines carrées};}]
\vspace*{1.75cm}
\begin{minipage}[t][6.9cm][c]{\linewidth}\LARGE\centering
\begin{minipage}{0.95\linewidth}\raggedright\linespread{1.25}\selectfont Écrire $10^{-1}\times10^{6}$ sous la forme d'une seule puissance de $10$.\end{minipage}
\end{minipage}

\vfill
\begin{tcolorbox}[enhanced,nobeforeafter,width=7.1cm,height=2.1cm,valign=center,colback=cpui!6,colframe=cpui,arc=2.5mm,boxrule=0.9pt,left=3mm]{\Large\bfseries\color{cpui}a)}\quad{\Large $100^{5}$}\end{tcolorbox}\hfill\begin{tcolorbox}[enhanced,nobeforeafter,width=7.1cm,height=2.1cm,valign=center,colback=cpui!6,colframe=cpui,arc=2.5mm,boxrule=0.9pt,left=3mm]{\Large\bfseries\color{cpui}b)}\quad{\Large $10^{-6}$}\end{tcolorbox}\par\vspace{2.5mm}
\begin{tcolorbox}[enhanced,nobeforeafter,width=7.1cm,height=2.1cm,valign=center,colback=cpui!6,colframe=cpui,arc=2.5mm,boxrule=0.9pt,left=3mm]{\Large\bfseries\color{cpui}c)}\quad{\Large $10^{-7}$}\end{tcolorbox}\hfill\begin{tcolorbox}[enhanced,nobeforeafter,width=7.1cm,height=2.1cm,valign=center,colback=cpui!6,colframe=cpui,arc=2.5mm,boxrule=0.9pt,left=3mm]{\Large\bfseries\color{cpui}d)}\quad{\Large $10^{5}$}\end{tcolorbox}
\end{tcolorbox}
\newpage
\begin{tcolorbox}[enhanced,width=16cm,height=14cm,colback=white,colframe=cpui,boxrule=2pt,arc=0mm,
left=6mm,right=6mm,top=0mm,bottom=4mm,sharp corners,
overlay={\fill[cpui] (frame.north west) rectangle ([yshift=-1.25cm]frame.north east);
\node[anchor=west,text=white,font=\sffamily\bfseries\LARGE] at ([xshift=5mm,yshift=-0.62cm]frame.north west) {Question 21 D};
\node[anchor=east,text=white,font=\sffamily\itshape\large] at ([xshift=-5mm,yshift=-0.62cm]frame.north east) {Puissances et racines carrées};}]
\vspace*{1.75cm}
\begin{minipage}[t][6.9cm][c]{\linewidth}\LARGE\centering
\begin{minipage}{0.95\linewidth}\raggedright\linespread{1.25}\selectfont Quelle est l'écriture décimale de $10^{-7}$~?\end{minipage}
\end{minipage}

\vfill
\begin{tcolorbox}[enhanced,nobeforeafter,width=7.1cm,height=2.1cm,valign=center,colback=cpui!6,colframe=cpui,arc=2.5mm,boxrule=0.9pt,left=3mm]{\Large\bfseries\color{cpui}a)}\quad{\Large $0{,}000001$}\end{tcolorbox}\hfill\begin{tcolorbox}[enhanced,nobeforeafter,width=7.1cm,height=2.1cm,valign=center,colback=cpui!6,colframe=cpui,arc=2.5mm,boxrule=0.9pt,left=3mm]{\Large\bfseries\color{cpui}b)}\quad{\Large $0{,}0000001$}\end{tcolorbox}\par\vspace{2.5mm}
\begin{tcolorbox}[enhanced,nobeforeafter,width=7.1cm,height=2.1cm,valign=center,colback=cpui!6,colframe=cpui,arc=2.5mm,boxrule=0.9pt,left=3mm]{\Large\bfseries\color{cpui}c)}\quad{\Large $-10\,000\,000$}\end{tcolorbox}\hfill\begin{tcolorbox}[enhanced,nobeforeafter,width=7.1cm,height=2.1cm,valign=center,colback=cpui!6,colframe=cpui,arc=2.5mm,boxrule=0.9pt,left=3mm]{\Large\bfseries\color{cpui}d)}\quad{\Large $0{,}00000001$}\end{tcolorbox}
\end{tcolorbox}
\newpage
\begin{tcolorbox}[enhanced,width=16cm,height=14cm,colback=white,colframe=cpui,boxrule=2pt,arc=0mm,
left=6mm,right=6mm,top=0mm,bottom=4mm,sharp corners,
overlay={\fill[cpui] (frame.north west) rectangle ([yshift=-1.25cm]frame.north east);
\node[anchor=west,text=white,font=\sffamily\bfseries\LARGE] at ([xshift=5mm,yshift=-0.62cm]frame.north west) {Question 22 D};
\node[anchor=east,text=white,font=\sffamily\itshape\large] at ([xshift=-5mm,yshift=-0.62cm]frame.north east) {Puissances et racines carrées};}]
\vspace*{1.75cm}
\begin{minipage}[t][6.9cm][c]{\linewidth}\LARGE\centering
\begin{minipage}{0.95\linewidth}\raggedright\linespread{1.25}\selectfont Quelle est l'écriture décimale de $10^{-8}$~?\end{minipage}
\end{minipage}

\vfill
\begin{tcolorbox}[enhanced,nobeforeafter,width=7.1cm,height=2.1cm,valign=center,colback=cpui!6,colframe=cpui,arc=2.5mm,boxrule=0.9pt,left=3mm]{\Large\bfseries\color{cpui}a)}\quad{\Large $0{,}0000001$}\end{tcolorbox}\hfill\begin{tcolorbox}[enhanced,nobeforeafter,width=7.1cm,height=2.1cm,valign=center,colback=cpui!6,colframe=cpui,arc=2.5mm,boxrule=0.9pt,left=3mm]{\Large\bfseries\color{cpui}b)}\quad{\Large $0{,}000000001$}\end{tcolorbox}\par\vspace{2.5mm}
\begin{tcolorbox}[enhanced,nobeforeafter,width=7.1cm,height=2.1cm,valign=center,colback=cpui!6,colframe=cpui,arc=2.5mm,boxrule=0.9pt,left=3mm]{\Large\bfseries\color{cpui}c)}\quad{\Large $-100\,000\,000$}\end{tcolorbox}\hfill\begin{tcolorbox}[enhanced,nobeforeafter,width=7.1cm,height=2.1cm,valign=center,colback=cpui!6,colframe=cpui,arc=2.5mm,boxrule=0.9pt,left=3mm]{\Large\bfseries\color{cpui}d)}\quad{\Large $0{,}00000001$}\end{tcolorbox}
\end{tcolorbox}
\newpage
\begin{tcolorbox}[enhanced,width=16cm,height=14cm,colback=white,colframe=cpui,boxrule=2pt,arc=0mm,
left=6mm,right=6mm,top=0mm,bottom=4mm,sharp corners,
overlay={\fill[cpui] (frame.north west) rectangle ([yshift=-1.25cm]frame.north east);
\node[anchor=west,text=white,font=\sffamily\bfseries\LARGE] at ([xshift=5mm,yshift=-0.62cm]frame.north west) {Question 23 D};
\node[anchor=east,text=white,font=\sffamily\itshape\large] at ([xshift=-5mm,yshift=-0.62cm]frame.north east) {Puissances et racines carrées};}]
\vspace*{1.75cm}
\begin{minipage}[t][6.9cm][c]{\linewidth}\LARGE\centering
\begin{minipage}{0.95\linewidth}\raggedright\linespread{1.25}\selectfont Quelle est la notation scientifique de $41\,000$~?\end{minipage}
\end{minipage}

\vfill
\begin{tcolorbox}[enhanced,nobeforeafter,width=7.1cm,height=2.1cm,valign=center,colback=cpui!6,colframe=cpui,arc=2.5mm,boxrule=0.9pt,left=3mm]{\Large\bfseries\color{cpui}a)}\quad{\Large $4{,}1\times10^{5}$}\end{tcolorbox}\hfill\begin{tcolorbox}[enhanced,nobeforeafter,width=7.1cm,height=2.1cm,valign=center,colback=cpui!6,colframe=cpui,arc=2.5mm,boxrule=0.9pt,left=3mm]{\Large\bfseries\color{cpui}b)}\quad{\Large $4{,}1\times10^{4}$}\end{tcolorbox}\par\vspace{2.5mm}
\begin{tcolorbox}[enhanced,nobeforeafter,width=7.1cm,height=2.1cm,valign=center,colback=cpui!6,colframe=cpui,arc=2.5mm,boxrule=0.9pt,left=3mm]{\Large\bfseries\color{cpui}c)}\quad{\Large $4{,}1\times10^{-4}$}\end{tcolorbox}\hfill\begin{tcolorbox}[enhanced,nobeforeafter,width=7.1cm,height=2.1cm,valign=center,colback=cpui!6,colframe=cpui,arc=2.5mm,boxrule=0.9pt,left=3mm]{\Large\bfseries\color{cpui}d)}\quad{\Large $41\times10^{4}$}\end{tcolorbox}
\end{tcolorbox}
\newpage
\begin{tcolorbox}[enhanced,width=16cm,height=14cm,colback=white,colframe=cpui,boxrule=2pt,arc=0mm,
left=6mm,right=6mm,top=0mm,bottom=4mm,sharp corners,
overlay={\fill[cpui] (frame.north west) rectangle ([yshift=-1.25cm]frame.north east);
\node[anchor=west,text=white,font=\sffamily\bfseries\LARGE] at ([xshift=5mm,yshift=-0.62cm]frame.north west) {Question 24 D};
\node[anchor=east,text=white,font=\sffamily\itshape\large] at ([xshift=-5mm,yshift=-0.62cm]frame.north east) {Puissances et racines carrées};}]
\vspace*{1.75cm}
\begin{minipage}[t][6.9cm][c]{\linewidth}\LARGE\centering
\begin{minipage}{0.95\linewidth}\raggedright\linespread{1.25}\selectfont Quelle est la notation scientifique de $0{,}000907$~?\end{minipage}
\end{minipage}

\vfill
\begin{tcolorbox}[enhanced,nobeforeafter,width=7.1cm,height=2.1cm,valign=center,colback=cpui!6,colframe=cpui,arc=2.5mm,boxrule=0.9pt,left=3mm]{\Large\bfseries\color{cpui}a)}\quad{\Large $9{,}07\times10^{4}$}\end{tcolorbox}\hfill\begin{tcolorbox}[enhanced,nobeforeafter,width=7.1cm,height=2.1cm,valign=center,colback=cpui!6,colframe=cpui,arc=2.5mm,boxrule=0.9pt,left=3mm]{\Large\bfseries\color{cpui}b)}\quad{\Large $9{,}07\times10^{-3}$}\end{tcolorbox}\par\vspace{2.5mm}
\begin{tcolorbox}[enhanced,nobeforeafter,width=7.1cm,height=2.1cm,valign=center,colback=cpui!6,colframe=cpui,arc=2.5mm,boxrule=0.9pt,left=3mm]{\Large\bfseries\color{cpui}c)}\quad{\Large $90{,}7\times10^{-4}$}\end{tcolorbox}\hfill\begin{tcolorbox}[enhanced,nobeforeafter,width=7.1cm,height=2.1cm,valign=center,colback=cpui!6,colframe=cpui,arc=2.5mm,boxrule=0.9pt,left=3mm]{\Large\bfseries\color{cpui}d)}\quad{\Large $9{,}07\times10^{-4}$}\end{tcolorbox}
\end{tcolorbox}
\newpage
\begin{tcolorbox}[enhanced,width=16cm,height=14cm,colback=white,colframe=cpui,boxrule=2pt,arc=0mm,
left=6mm,right=6mm,top=0mm,bottom=4mm,sharp corners,
overlay={\fill[cpui] (frame.north west) rectangle ([yshift=-1.25cm]frame.north east);
\node[anchor=west,text=white,font=\sffamily\bfseries\LARGE] at ([xshift=5mm,yshift=-0.62cm]frame.north west) {Question 25 D};
\node[anchor=east,text=white,font=\sffamily\itshape\large] at ([xshift=-5mm,yshift=-0.62cm]frame.north east) {Puissances et racines carrées};}]
\vspace*{1.75cm}
\begin{minipage}[t][6.9cm][c]{\linewidth}\LARGE\centering
\begin{minipage}{0.95\linewidth}\raggedright\linespread{1.25}\selectfont Que vaut $-7^{2}$~?\end{minipage}
\end{minipage}

\vfill
\begin{tcolorbox}[enhanced,nobeforeafter,width=7.1cm,height=2.1cm,valign=center,colback=cpui!6,colframe=cpui,arc=2.5mm,boxrule=0.9pt,left=3mm]{\Large\bfseries\color{cpui}a)}\quad{\Large $49$}\end{tcolorbox}\hfill\begin{tcolorbox}[enhanced,nobeforeafter,width=7.1cm,height=2.1cm,valign=center,colback=cpui!6,colframe=cpui,arc=2.5mm,boxrule=0.9pt,left=3mm]{\Large\bfseries\color{cpui}b)}\quad{\Large $-49$}\end{tcolorbox}\par\vspace{2.5mm}
\begin{tcolorbox}[enhanced,nobeforeafter,width=7.1cm,height=2.1cm,valign=center,colback=cpui!6,colframe=cpui,arc=2.5mm,boxrule=0.9pt,left=3mm]{\Large\bfseries\color{cpui}c)}\quad{\Large $14$}\end{tcolorbox}\hfill\begin{tcolorbox}[enhanced,nobeforeafter,width=7.1cm,height=2.1cm,valign=center,colback=cpui!6,colframe=cpui,arc=2.5mm,boxrule=0.9pt,left=3mm]{\Large\bfseries\color{cpui}d)}\quad{\Large $-14$}\end{tcolorbox}
\end{tcolorbox}
\newpage
\begin{tcolorbox}[enhanced,width=16cm,height=14cm,colback=white,colframe=cpui,boxrule=2pt,arc=0mm,
left=6mm,right=6mm,top=0mm,bottom=4mm,sharp corners,
overlay={\fill[cpui] (frame.north west) rectangle ([yshift=-1.25cm]frame.north east);
\node[anchor=west,text=white,font=\sffamily\bfseries\LARGE] at ([xshift=5mm,yshift=-0.62cm]frame.north west) {Question 26 D};
\node[anchor=east,text=white,font=\sffamily\itshape\large] at ([xshift=-5mm,yshift=-0.62cm]frame.north east) {Puissances et racines carrées};}]
\vspace*{1.75cm}
\begin{minipage}[t][6.9cm][c]{\linewidth}\LARGE\centering
\begin{minipage}{0.95\linewidth}\raggedright\linespread{1.25}\selectfont Que vaut $-2^{6}$~?\end{minipage}
\end{minipage}

\vfill
\begin{tcolorbox}[enhanced,nobeforeafter,width=7.1cm,height=2.1cm,valign=center,colback=cpui!6,colframe=cpui,arc=2.5mm,boxrule=0.9pt,left=3mm]{\Large\bfseries\color{cpui}a)}\quad{\Large $-64$}\end{tcolorbox}\hfill\begin{tcolorbox}[enhanced,nobeforeafter,width=7.1cm,height=2.1cm,valign=center,colback=cpui!6,colframe=cpui,arc=2.5mm,boxrule=0.9pt,left=3mm]{\Large\bfseries\color{cpui}b)}\quad{\Large $-12$}\end{tcolorbox}\par\vspace{2.5mm}
\begin{tcolorbox}[enhanced,nobeforeafter,width=7.1cm,height=2.1cm,valign=center,colback=cpui!6,colframe=cpui,arc=2.5mm,boxrule=0.9pt,left=3mm]{\Large\bfseries\color{cpui}c)}\quad{\Large $64$}\end{tcolorbox}\hfill\begin{tcolorbox}[enhanced,nobeforeafter,width=7.1cm,height=2.1cm,valign=center,colback=cpui!6,colframe=cpui,arc=2.5mm,boxrule=0.9pt,left=3mm]{\Large\bfseries\color{cpui}d)}\quad{\Large $12$}\end{tcolorbox}
\end{tcolorbox}
\newpage
\begin{tcolorbox}[enhanced,width=16cm,height=14cm,colback=white,colframe=cpui,boxrule=2pt,arc=0mm,
left=6mm,right=6mm,top=0mm,bottom=4mm,sharp corners,
overlay={\fill[cpui] (frame.north west) rectangle ([yshift=-1.25cm]frame.north east);
\node[anchor=west,text=white,font=\sffamily\bfseries\LARGE] at ([xshift=5mm,yshift=-0.62cm]frame.north west) {Question 27 D};
\node[anchor=east,text=white,font=\sffamily\itshape\large] at ([xshift=-5mm,yshift=-0.62cm]frame.north east) {Puissances et racines carrées};}]
\vspace*{1.75cm}
\begin{minipage}[t][6.9cm][c]{\linewidth}\LARGE\centering
\begin{minipage}{0.95\linewidth}\raggedright\linespread{1.25}\selectfont Écrire $\dfrac{10^{-4}}{10^{-6}}$ sous la forme d'une seule puissance de $10$.\end{minipage}
\end{minipage}

\vfill
\begin{tcolorbox}[enhanced,nobeforeafter,width=7.1cm,height=2.1cm,valign=center,colback=cpui!6,colframe=cpui,arc=2.5mm,boxrule=0.9pt,left=3mm]{\Large\bfseries\color{cpui}a)}\quad{\Large $10^{24}$}\end{tcolorbox}\hfill\begin{tcolorbox}[enhanced,nobeforeafter,width=7.1cm,height=2.1cm,valign=center,colback=cpui!6,colframe=cpui,arc=2.5mm,boxrule=0.9pt,left=3mm]{\Large\bfseries\color{cpui}b)}\quad{\Large $10^{2}$}\end{tcolorbox}\par\vspace{2.5mm}
\begin{tcolorbox}[enhanced,nobeforeafter,width=7.1cm,height=2.1cm,valign=center,colback=cpui!6,colframe=cpui,arc=2.5mm,boxrule=0.9pt,left=3mm]{\Large\bfseries\color{cpui}c)}\quad{\Large $10^{-10}$}\end{tcolorbox}\hfill\begin{tcolorbox}[enhanced,nobeforeafter,width=7.1cm,height=2.1cm,valign=center,colback=cpui!6,colframe=cpui,arc=2.5mm,boxrule=0.9pt,left=3mm]{\Large\bfseries\color{cpui}d)}\quad{\Large $10^{-2}$}\end{tcolorbox}
\end{tcolorbox}
\newpage
\begin{tcolorbox}[enhanced,width=16cm,height=14cm,colback=white,colframe=cpui,boxrule=2pt,arc=0mm,
left=6mm,right=6mm,top=0mm,bottom=4mm,sharp corners,
overlay={\fill[cpui] (frame.north west) rectangle ([yshift=-1.25cm]frame.north east);
\node[anchor=west,text=white,font=\sffamily\bfseries\LARGE] at ([xshift=5mm,yshift=-0.62cm]frame.north west) {Question 28 D};
\node[anchor=east,text=white,font=\sffamily\itshape\large] at ([xshift=-5mm,yshift=-0.62cm]frame.north east) {Puissances et racines carrées};}]
\vspace*{1.75cm}
\begin{minipage}[t][6.9cm][c]{\linewidth}\LARGE\centering
\begin{minipage}{0.95\linewidth}\raggedright\linespread{1.25}\selectfont Écrire $\dfrac{10^{8}}{10^{5}}$ sous la forme d'une seule puissance de $10$.\end{minipage}
\end{minipage}

\vfill
\begin{tcolorbox}[enhanced,nobeforeafter,width=7.1cm,height=2.1cm,valign=center,colback=cpui!6,colframe=cpui,arc=2.5mm,boxrule=0.9pt,left=3mm]{\Large\bfseries\color{cpui}a)}\quad{\Large $10^{-3}$}\end{tcolorbox}\hfill\begin{tcolorbox}[enhanced,nobeforeafter,width=7.1cm,height=2.1cm,valign=center,colback=cpui!6,colframe=cpui,arc=2.5mm,boxrule=0.9pt,left=3mm]{\Large\bfseries\color{cpui}b)}\quad{\Large $10^{13}$}\end{tcolorbox}\par\vspace{2.5mm}
\begin{tcolorbox}[enhanced,nobeforeafter,width=7.1cm,height=2.1cm,valign=center,colback=cpui!6,colframe=cpui,arc=2.5mm,boxrule=0.9pt,left=3mm]{\Large\bfseries\color{cpui}c)}\quad{\Large $10^{40}$}\end{tcolorbox}\hfill\begin{tcolorbox}[enhanced,nobeforeafter,width=7.1cm,height=2.1cm,valign=center,colback=cpui!6,colframe=cpui,arc=2.5mm,boxrule=0.9pt,left=3mm]{\Large\bfseries\color{cpui}d)}\quad{\Large $10^{3}$}\end{tcolorbox}
\end{tcolorbox}
\newpage
\begin{tcolorbox}[enhanced,width=16cm,height=14cm,colback=white,colframe=cpui,boxrule=2pt,arc=0mm,
left=6mm,right=6mm,top=0mm,bottom=4mm,sharp corners,
overlay={\fill[cpui] (frame.north west) rectangle ([yshift=-1.25cm]frame.north east);
\node[anchor=west,text=white,font=\sffamily\bfseries\LARGE] at ([xshift=5mm,yshift=-0.62cm]frame.north west) {Question 29 D};
\node[anchor=east,text=white,font=\sffamily\itshape\large] at ([xshift=-5mm,yshift=-0.62cm]frame.north east) {Puissances et racines carrées};}]
\vspace*{1.75cm}
\begin{minipage}[t][6.9cm][c]{\linewidth}\LARGE\centering
\begin{minipage}{0.95\linewidth}\raggedright\linespread{1.25}\selectfont Que vaut $\sqrt{400}$~?\end{minipage}
\end{minipage}

\vfill
\begin{tcolorbox}[enhanced,nobeforeafter,width=7.1cm,height=2.1cm,valign=center,colback=cpui!6,colframe=cpui,arc=2.5mm,boxrule=0.9pt,left=3mm]{\Large\bfseries\color{cpui}a)}\quad{\Large $200$}\end{tcolorbox}\hfill\begin{tcolorbox}[enhanced,nobeforeafter,width=7.1cm,height=2.1cm,valign=center,colback=cpui!6,colframe=cpui,arc=2.5mm,boxrule=0.9pt,left=3mm]{\Large\bfseries\color{cpui}b)}\quad{\Large $21$}\end{tcolorbox}\par\vspace{2.5mm}
\begin{tcolorbox}[enhanced,nobeforeafter,width=7.1cm,height=2.1cm,valign=center,colback=cpui!6,colframe=cpui,arc=2.5mm,boxrule=0.9pt,left=3mm]{\Large\bfseries\color{cpui}c)}\quad{\Large $20$}\end{tcolorbox}\hfill\begin{tcolorbox}[enhanced,nobeforeafter,width=7.1cm,height=2.1cm,valign=center,colback=cpui!6,colframe=cpui,arc=2.5mm,boxrule=0.9pt,left=3mm]{\Large\bfseries\color{cpui}d)}\quad{\Large $40$}\end{tcolorbox}
\end{tcolorbox}
\newpage
\begin{tcolorbox}[enhanced,width=16cm,height=14cm,colback=white,colframe=cpui,boxrule=2pt,arc=0mm,
left=6mm,right=6mm,top=0mm,bottom=4mm,sharp corners,
overlay={\fill[cpui] (frame.north west) rectangle ([yshift=-1.25cm]frame.north east);
\node[anchor=west,text=white,font=\sffamily\bfseries\LARGE] at ([xshift=5mm,yshift=-0.62cm]frame.north west) {Question 30 D};
\node[anchor=east,text=white,font=\sffamily\itshape\large] at ([xshift=-5mm,yshift=-0.62cm]frame.north east) {Puissances et racines carrées};}]
\vspace*{1.75cm}
\begin{minipage}[t][6.9cm][c]{\linewidth}\LARGE\centering
\begin{minipage}{0.95\linewidth}\raggedright\linespread{1.25}\selectfont Que vaut $\sqrt{64}$~?\end{minipage}
\end{minipage}

\vfill
\begin{tcolorbox}[enhanced,nobeforeafter,width=7.1cm,height=2.1cm,valign=center,colback=cpui!6,colframe=cpui,arc=2.5mm,boxrule=0.9pt,left=3mm]{\Large\bfseries\color{cpui}a)}\quad{\Large $16$}\end{tcolorbox}\hfill\begin{tcolorbox}[enhanced,nobeforeafter,width=7.1cm,height=2.1cm,valign=center,colback=cpui!6,colframe=cpui,arc=2.5mm,boxrule=0.9pt,left=3mm]{\Large\bfseries\color{cpui}b)}\quad{\Large $9$}\end{tcolorbox}\par\vspace{2.5mm}
\begin{tcolorbox}[enhanced,nobeforeafter,width=7.1cm,height=2.1cm,valign=center,colback=cpui!6,colframe=cpui,arc=2.5mm,boxrule=0.9pt,left=3mm]{\Large\bfseries\color{cpui}c)}\quad{\Large $8$}\end{tcolorbox}\hfill\begin{tcolorbox}[enhanced,nobeforeafter,width=7.1cm,height=2.1cm,valign=center,colback=cpui!6,colframe=cpui,arc=2.5mm,boxrule=0.9pt,left=3mm]{\Large\bfseries\color{cpui}d)}\quad{\Large $32$}\end{tcolorbox}
\end{tcolorbox}
\newpage
\begin{tcolorbox}[enhanced,width=16cm,height=14cm,colback=white,colframe=cpui,boxrule=2pt,arc=0mm,
left=6mm,right=6mm,top=0mm,bottom=4mm,sharp corners,
overlay={\fill[cpui] (frame.north west) rectangle ([yshift=-1.25cm]frame.north east);
\node[anchor=west,text=white,font=\sffamily\bfseries\LARGE] at ([xshift=5mm,yshift=-0.62cm]frame.north west) {Question 31 D};
\node[anchor=east,text=white,font=\sffamily\itshape\large] at ([xshift=-5mm,yshift=-0.62cm]frame.north east) {Puissances et racines carrées};}]
\vspace*{1.75cm}
\begin{minipage}[t][6.9cm][c]{\linewidth}\LARGE\centering
\begin{minipage}{0.95\linewidth}\raggedright\linespread{1.25}\selectfont Entre quels nombres entiers consécutifs se trouve $\sqrt{150}$~?\end{minipage}
\end{minipage}

\vfill
\begin{tcolorbox}[enhanced,nobeforeafter,width=7.1cm,height=2.1cm,valign=center,colback=cpui!6,colframe=cpui,arc=2.5mm,boxrule=0.9pt,left=3mm]{\Large\bfseries\color{cpui}a)}\quad{\Large $75\text{ et }76$}\end{tcolorbox}\hfill\begin{tcolorbox}[enhanced,nobeforeafter,width=7.1cm,height=2.1cm,valign=center,colback=cpui!6,colframe=cpui,arc=2.5mm,boxrule=0.9pt,left=3mm]{\Large\bfseries\color{cpui}b)}\quad{\Large $13\text{ et }14$}\end{tcolorbox}\par\vspace{2.5mm}
\begin{tcolorbox}[enhanced,nobeforeafter,width=7.1cm,height=2.1cm,valign=center,colback=cpui!6,colframe=cpui,arc=2.5mm,boxrule=0.9pt,left=3mm]{\Large\bfseries\color{cpui}c)}\quad{\Large $12\text{ et }13$}\end{tcolorbox}\hfill\begin{tcolorbox}[enhanced,nobeforeafter,width=7.1cm,height=2.1cm,valign=center,colback=cpui!6,colframe=cpui,arc=2.5mm,boxrule=0.9pt,left=3mm]{\Large\bfseries\color{cpui}d)}\quad{\Large $11\text{ et }12$}\end{tcolorbox}
\end{tcolorbox}
\newpage
\begin{tcolorbox}[enhanced,width=16cm,height=14cm,colback=white,colframe=cpui,boxrule=2pt,arc=0mm,
left=6mm,right=6mm,top=0mm,bottom=4mm,sharp corners,
overlay={\fill[cpui] (frame.north west) rectangle ([yshift=-1.25cm]frame.north east);
\node[anchor=west,text=white,font=\sffamily\bfseries\LARGE] at ([xshift=5mm,yshift=-0.62cm]frame.north west) {Question 32 D};
\node[anchor=east,text=white,font=\sffamily\itshape\large] at ([xshift=-5mm,yshift=-0.62cm]frame.north east) {Puissances et racines carrées};}]
\vspace*{1.75cm}
\begin{minipage}[t][6.9cm][c]{\linewidth}\LARGE\centering
\begin{minipage}{0.95\linewidth}\raggedright\linespread{1.25}\selectfont Entre quels nombres entiers consécutifs se trouve $\sqrt{200}$~?\end{minipage}
\end{minipage}

\vfill
\begin{tcolorbox}[enhanced,nobeforeafter,width=7.1cm,height=2.1cm,valign=center,colback=cpui!6,colframe=cpui,arc=2.5mm,boxrule=0.9pt,left=3mm]{\Large\bfseries\color{cpui}a)}\quad{\Large $13\text{ et }14$}\end{tcolorbox}\hfill\begin{tcolorbox}[enhanced,nobeforeafter,width=7.1cm,height=2.1cm,valign=center,colback=cpui!6,colframe=cpui,arc=2.5mm,boxrule=0.9pt,left=3mm]{\Large\bfseries\color{cpui}b)}\quad{\Large $100\text{ et }101$}\end{tcolorbox}\par\vspace{2.5mm}
\begin{tcolorbox}[enhanced,nobeforeafter,width=7.1cm,height=2.1cm,valign=center,colback=cpui!6,colframe=cpui,arc=2.5mm,boxrule=0.9pt,left=3mm]{\Large\bfseries\color{cpui}c)}\quad{\Large $15\text{ et }16$}\end{tcolorbox}\hfill\begin{tcolorbox}[enhanced,nobeforeafter,width=7.1cm,height=2.1cm,valign=center,colback=cpui!6,colframe=cpui,arc=2.5mm,boxrule=0.9pt,left=3mm]{\Large\bfseries\color{cpui}d)}\quad{\Large $14\text{ et }15$}\end{tcolorbox}
\end{tcolorbox}
\newpage
\begin{tcolorbox}[enhanced,width=16cm,height=14cm,colback=white,colframe=clit,boxrule=2pt,arc=0mm,
left=6mm,right=6mm,top=0mm,bottom=4mm,sharp corners,
overlay={\fill[clit] (frame.north west) rectangle ([yshift=-1.25cm]frame.north east);
\node[anchor=west,text=white,font=\sffamily\bfseries\LARGE] at ([xshift=5mm,yshift=-0.62cm]frame.north west) {Question 33 D};
\node[anchor=east,text=white,font=\sffamily\itshape\large] at ([xshift=-5mm,yshift=-0.62cm]frame.north east) {Calcul littéral et équations};}]
\vspace*{1.75cm}
\begin{minipage}[t][6.9cm][c]{\linewidth}\LARGE\centering
\begin{minipage}{0.95\linewidth}\raggedright\linespread{1.25}\selectfont Développer $6(-x+2)$.\end{minipage}
\end{minipage}

\vfill
\begin{tcolorbox}[enhanced,nobeforeafter,width=7.1cm,height=2.1cm,valign=center,colback=clit!6,colframe=clit,arc=2.5mm,boxrule=0.9pt,left=3mm]{\Large\bfseries\color{clit}a)}\quad{\Large $-6x-12$}\end{tcolorbox}\hfill\begin{tcolorbox}[enhanced,nobeforeafter,width=7.1cm,height=2.1cm,valign=center,colback=clit!6,colframe=clit,arc=2.5mm,boxrule=0.9pt,left=3mm]{\Large\bfseries\color{clit}b)}\quad{\Large $-6x+2$}\end{tcolorbox}\par\vspace{2.5mm}
\begin{tcolorbox}[enhanced,nobeforeafter,width=7.1cm,height=2.1cm,valign=center,colback=clit!6,colframe=clit,arc=2.5mm,boxrule=0.9pt,left=3mm]{\Large\bfseries\color{clit}c)}\quad{\Large $-x+12$}\end{tcolorbox}\hfill\begin{tcolorbox}[enhanced,nobeforeafter,width=7.1cm,height=2.1cm,valign=center,colback=clit!6,colframe=clit,arc=2.5mm,boxrule=0.9pt,left=3mm]{\Large\bfseries\color{clit}d)}\quad{\Large $-6x+12$}\end{tcolorbox}
\end{tcolorbox}
\newpage
\begin{tcolorbox}[enhanced,width=16cm,height=14cm,colback=white,colframe=clit,boxrule=2pt,arc=0mm,
left=6mm,right=6mm,top=0mm,bottom=4mm,sharp corners,
overlay={\fill[clit] (frame.north west) rectangle ([yshift=-1.25cm]frame.north east);
\node[anchor=west,text=white,font=\sffamily\bfseries\LARGE] at ([xshift=5mm,yshift=-0.62cm]frame.north west) {Question 34 D};
\node[anchor=east,text=white,font=\sffamily\itshape\large] at ([xshift=-5mm,yshift=-0.62cm]frame.north east) {Calcul littéral et équations};}]
\vspace*{1.75cm}
\begin{minipage}[t][6.9cm][c]{\linewidth}\LARGE\centering
\begin{minipage}{0.95\linewidth}\raggedright\linespread{1.25}\selectfont Développer $7(2x-3)$.\end{minipage}
\end{minipage}

\vfill
\begin{tcolorbox}[enhanced,nobeforeafter,width=7.1cm,height=2.1cm,valign=center,colback=clit!6,colframe=clit,arc=2.5mm,boxrule=0.9pt,left=3mm]{\Large\bfseries\color{clit}a)}\quad{\Large $14x+21$}\end{tcolorbox}\hfill\begin{tcolorbox}[enhanced,nobeforeafter,width=7.1cm,height=2.1cm,valign=center,colback=clit!6,colframe=clit,arc=2.5mm,boxrule=0.9pt,left=3mm]{\Large\bfseries\color{clit}b)}\quad{\Large $2x-21$}\end{tcolorbox}\par\vspace{2.5mm}
\begin{tcolorbox}[enhanced,nobeforeafter,width=7.1cm,height=2.1cm,valign=center,colback=clit!6,colframe=clit,arc=2.5mm,boxrule=0.9pt,left=3mm]{\Large\bfseries\color{clit}c)}\quad{\Large $14x-21$}\end{tcolorbox}\hfill\begin{tcolorbox}[enhanced,nobeforeafter,width=7.1cm,height=2.1cm,valign=center,colback=clit!6,colframe=clit,arc=2.5mm,boxrule=0.9pt,left=3mm]{\Large\bfseries\color{clit}d)}\quad{\Large $14x-3$}\end{tcolorbox}
\end{tcolorbox}
\newpage
\begin{tcolorbox}[enhanced,width=16cm,height=14cm,colback=white,colframe=clit,boxrule=2pt,arc=0mm,
left=6mm,right=6mm,top=0mm,bottom=4mm,sharp corners,
overlay={\fill[clit] (frame.north west) rectangle ([yshift=-1.25cm]frame.north east);
\node[anchor=west,text=white,font=\sffamily\bfseries\LARGE] at ([xshift=5mm,yshift=-0.62cm]frame.north west) {Question 35 D};
\node[anchor=east,text=white,font=\sffamily\itshape\large] at ([xshift=-5mm,yshift=-0.62cm]frame.north east) {Calcul littéral et équations};}]
\vspace*{1.75cm}
\begin{minipage}[t][6.9cm][c]{\linewidth}\LARGE\centering
\begin{minipage}{0.95\linewidth}\raggedright\linespread{1.25}\selectfont Développer et réduire $(x-1)(x-8)$.\end{minipage}
\end{minipage}

\vfill
\begin{tcolorbox}[enhanced,nobeforeafter,width=7.1cm,height=2.1cm,valign=center,colback=clit!6,colframe=clit,arc=2.5mm,boxrule=0.9pt,left=3mm]{\Large\bfseries\color{clit}a)}\quad{\Large $x^2+8$}\end{tcolorbox}\hfill\begin{tcolorbox}[enhanced,nobeforeafter,width=7.1cm,height=2.1cm,valign=center,colback=clit!6,colframe=clit,arc=2.5mm,boxrule=0.9pt,left=3mm]{\Large\bfseries\color{clit}b)}\quad{\Large $x^2-9x-9$}\end{tcolorbox}\par\vspace{2.5mm}
\begin{tcolorbox}[enhanced,nobeforeafter,width=7.1cm,height=2.1cm,valign=center,colback=clit!6,colframe=clit,arc=2.5mm,boxrule=0.9pt,left=3mm]{\Large\bfseries\color{clit}c)}\quad{\Large $x^2+8x-9$}\end{tcolorbox}\hfill\begin{tcolorbox}[enhanced,nobeforeafter,width=7.1cm,height=2.1cm,valign=center,colback=clit!6,colframe=clit,arc=2.5mm,boxrule=0.9pt,left=3mm]{\Large\bfseries\color{clit}d)}\quad{\Large $x^2-9x+8$}\end{tcolorbox}
\end{tcolorbox}
\newpage
\begin{tcolorbox}[enhanced,width=16cm,height=14cm,colback=white,colframe=clit,boxrule=2pt,arc=0mm,
left=6mm,right=6mm,top=0mm,bottom=4mm,sharp corners,
overlay={\fill[clit] (frame.north west) rectangle ([yshift=-1.25cm]frame.north east);
\node[anchor=west,text=white,font=\sffamily\bfseries\LARGE] at ([xshift=5mm,yshift=-0.62cm]frame.north west) {Question 36 D};
\node[anchor=east,text=white,font=\sffamily\itshape\large] at ([xshift=-5mm,yshift=-0.62cm]frame.north east) {Calcul littéral et équations};}]
\vspace*{1.75cm}
\begin{minipage}[t][6.9cm][c]{\linewidth}\LARGE\centering
\begin{minipage}{0.95\linewidth}\raggedright\linespread{1.25}\selectfont Développer et réduire $(x+3)(x+6)$.\end{minipage}
\end{minipage}

\vfill
\begin{tcolorbox}[enhanced,nobeforeafter,width=7.1cm,height=2.1cm,valign=center,colback=clit!6,colframe=clit,arc=2.5mm,boxrule=0.9pt,left=3mm]{\Large\bfseries\color{clit}a)}\quad{\Large $x^2+9x+18$}\end{tcolorbox}\hfill\begin{tcolorbox}[enhanced,nobeforeafter,width=7.1cm,height=2.1cm,valign=center,colback=clit!6,colframe=clit,arc=2.5mm,boxrule=0.9pt,left=3mm]{\Large\bfseries\color{clit}b)}\quad{\Large $x^2+18$}\end{tcolorbox}\par\vspace{2.5mm}
\begin{tcolorbox}[enhanced,nobeforeafter,width=7.1cm,height=2.1cm,valign=center,colback=clit!6,colframe=clit,arc=2.5mm,boxrule=0.9pt,left=3mm]{\Large\bfseries\color{clit}c)}\quad{\Large $x^2+18x+9$}\end{tcolorbox}\hfill\begin{tcolorbox}[enhanced,nobeforeafter,width=7.1cm,height=2.1cm,valign=center,colback=clit!6,colframe=clit,arc=2.5mm,boxrule=0.9pt,left=3mm]{\Large\bfseries\color{clit}d)}\quad{\Large $x^2+9x+9$}\end{tcolorbox}
\end{tcolorbox}
\newpage
\begin{tcolorbox}[enhanced,width=16cm,height=14cm,colback=white,colframe=clit,boxrule=2pt,arc=0mm,
left=6mm,right=6mm,top=0mm,bottom=4mm,sharp corners,
overlay={\fill[clit] (frame.north west) rectangle ([yshift=-1.25cm]frame.north east);
\node[anchor=west,text=white,font=\sffamily\bfseries\LARGE] at ([xshift=5mm,yshift=-0.62cm]frame.north west) {Question 37 D};
\node[anchor=east,text=white,font=\sffamily\itshape\large] at ([xshift=-5mm,yshift=-0.62cm]frame.north east) {Calcul littéral et équations};}]
\vspace*{1.75cm}
\begin{minipage}[t][6.9cm][c]{\linewidth}\LARGE\centering
\begin{minipage}{0.95\linewidth}\raggedright\linespread{1.25}\selectfont Réduire l'expression $6x+7-x-3$.\end{minipage}
\end{minipage}

\vfill
\begin{tcolorbox}[enhanced,nobeforeafter,width=7.1cm,height=2.1cm,valign=center,colback=clit!6,colframe=clit,arc=2.5mm,boxrule=0.9pt,left=3mm]{\Large\bfseries\color{clit}a)}\quad{\Large $5x+10$}\end{tcolorbox}\hfill\begin{tcolorbox}[enhanced,nobeforeafter,width=7.1cm,height=2.1cm,valign=center,colback=clit!6,colframe=clit,arc=2.5mm,boxrule=0.9pt,left=3mm]{\Large\bfseries\color{clit}b)}\quad{\Large $7x+4$}\end{tcolorbox}\par\vspace{2.5mm}
\begin{tcolorbox}[enhanced,nobeforeafter,width=7.1cm,height=2.1cm,valign=center,colback=clit!6,colframe=clit,arc=2.5mm,boxrule=0.9pt,left=3mm]{\Large\bfseries\color{clit}c)}\quad{\Large $5x+4$}\end{tcolorbox}\hfill\begin{tcolorbox}[enhanced,nobeforeafter,width=7.1cm,height=2.1cm,valign=center,colback=clit!6,colframe=clit,arc=2.5mm,boxrule=0.9pt,left=3mm]{\Large\bfseries\color{clit}d)}\quad{\Large $9x$}\end{tcolorbox}
\end{tcolorbox}
\newpage
\begin{tcolorbox}[enhanced,width=16cm,height=14cm,colback=white,colframe=clit,boxrule=2pt,arc=0mm,
left=6mm,right=6mm,top=0mm,bottom=4mm,sharp corners,
overlay={\fill[clit] (frame.north west) rectangle ([yshift=-1.25cm]frame.north east);
\node[anchor=west,text=white,font=\sffamily\bfseries\LARGE] at ([xshift=5mm,yshift=-0.62cm]frame.north west) {Question 38 D};
\node[anchor=east,text=white,font=\sffamily\itshape\large] at ([xshift=-5mm,yshift=-0.62cm]frame.north east) {Calcul littéral et équations};}]
\vspace*{1.75cm}
\begin{minipage}[t][6.9cm][c]{\linewidth}\LARGE\centering
\begin{minipage}{0.95\linewidth}\raggedright\linespread{1.25}\selectfont Réduire l'expression $2x-4+2x+9$.\end{minipage}
\end{minipage}

\vfill
\begin{tcolorbox}[enhanced,nobeforeafter,width=7.1cm,height=2.1cm,valign=center,colback=clit!6,colframe=clit,arc=2.5mm,boxrule=0.9pt,left=3mm]{\Large\bfseries\color{clit}a)}\quad{\Large $9x$}\end{tcolorbox}\hfill\begin{tcolorbox}[enhanced,nobeforeafter,width=7.1cm,height=2.1cm,valign=center,colback=clit!6,colframe=clit,arc=2.5mm,boxrule=0.9pt,left=3mm]{\Large\bfseries\color{clit}b)}\quad{\Large $4x+5$}\end{tcolorbox}\par\vspace{2.5mm}
\begin{tcolorbox}[enhanced,nobeforeafter,width=7.1cm,height=2.1cm,valign=center,colback=clit!6,colframe=clit,arc=2.5mm,boxrule=0.9pt,left=3mm]{\Large\bfseries\color{clit}c)}\quad{\Large $5$}\end{tcolorbox}\hfill\begin{tcolorbox}[enhanced,nobeforeafter,width=7.1cm,height=2.1cm,valign=center,colback=clit!6,colframe=clit,arc=2.5mm,boxrule=0.9pt,left=3mm]{\Large\bfseries\color{clit}d)}\quad{\Large $4x-13$}\end{tcolorbox}
\end{tcolorbox}
\newpage
\begin{tcolorbox}[enhanced,width=16cm,height=14cm,colback=white,colframe=clit,boxrule=2pt,arc=0mm,
left=6mm,right=6mm,top=0mm,bottom=4mm,sharp corners,
overlay={\fill[clit] (frame.north west) rectangle ([yshift=-1.25cm]frame.north east);
\node[anchor=west,text=white,font=\sffamily\bfseries\LARGE] at ([xshift=5mm,yshift=-0.62cm]frame.north west) {Question 39 D};
\node[anchor=east,text=white,font=\sffamily\itshape\large] at ([xshift=-5mm,yshift=-0.62cm]frame.north east) {Calcul littéral et équations};}]
\vspace*{1.75cm}
\begin{minipage}[t][6.9cm][c]{\linewidth}\LARGE\centering
\begin{minipage}{0.95\linewidth}\raggedright\linespread{1.25}\selectfont Factoriser $9x+18$.\end{minipage}
\end{minipage}

\vfill
\begin{tcolorbox}[enhanced,nobeforeafter,width=7.1cm,height=2.1cm,valign=center,colback=clit!6,colframe=clit,arc=2.5mm,boxrule=0.9pt,left=3mm]{\Large\bfseries\color{clit}a)}\quad{\Large $9(9x+2)$}\end{tcolorbox}\hfill\begin{tcolorbox}[enhanced,nobeforeafter,width=7.1cm,height=2.1cm,valign=center,colback=clit!6,colframe=clit,arc=2.5mm,boxrule=0.9pt,left=3mm]{\Large\bfseries\color{clit}b)}\quad{\Large $9x(x+2)$}\end{tcolorbox}\par\vspace{2.5mm}
\begin{tcolorbox}[enhanced,nobeforeafter,width=7.1cm,height=2.1cm,valign=center,colback=clit!6,colframe=clit,arc=2.5mm,boxrule=0.9pt,left=3mm]{\Large\bfseries\color{clit}c)}\quad{\Large $9(x+18)$}\end{tcolorbox}\hfill\begin{tcolorbox}[enhanced,nobeforeafter,width=7.1cm,height=2.1cm,valign=center,colback=clit!6,colframe=clit,arc=2.5mm,boxrule=0.9pt,left=3mm]{\Large\bfseries\color{clit}d)}\quad{\Large $9(x+2)$}\end{tcolorbox}
\end{tcolorbox}
\newpage
\begin{tcolorbox}[enhanced,width=16cm,height=14cm,colback=white,colframe=clit,boxrule=2pt,arc=0mm,
left=6mm,right=6mm,top=0mm,bottom=4mm,sharp corners,
overlay={\fill[clit] (frame.north west) rectangle ([yshift=-1.25cm]frame.north east);
\node[anchor=west,text=white,font=\sffamily\bfseries\LARGE] at ([xshift=5mm,yshift=-0.62cm]frame.north west) {Question 40 D};
\node[anchor=east,text=white,font=\sffamily\itshape\large] at ([xshift=-5mm,yshift=-0.62cm]frame.north east) {Calcul littéral et équations};}]
\vspace*{1.75cm}
\begin{minipage}[t][6.9cm][c]{\linewidth}\LARGE\centering
\begin{minipage}{0.95\linewidth}\raggedright\linespread{1.25}\selectfont Factoriser $24x+40$.\end{minipage}
\end{minipage}

\vfill
\begin{tcolorbox}[enhanced,nobeforeafter,width=7.1cm,height=2.1cm,valign=center,colback=clit!6,colframe=clit,arc=2.5mm,boxrule=0.9pt,left=3mm]{\Large\bfseries\color{clit}a)}\quad{\Large $8(3x+5)$}\end{tcolorbox}\hfill\begin{tcolorbox}[enhanced,nobeforeafter,width=7.1cm,height=2.1cm,valign=center,colback=clit!6,colframe=clit,arc=2.5mm,boxrule=0.9pt,left=3mm]{\Large\bfseries\color{clit}b)}\quad{\Large $24(x+5)$}\end{tcolorbox}\par\vspace{2.5mm}
\begin{tcolorbox}[enhanced,nobeforeafter,width=7.1cm,height=2.1cm,valign=center,colback=clit!6,colframe=clit,arc=2.5mm,boxrule=0.9pt,left=3mm]{\Large\bfseries\color{clit}c)}\quad{\Large $8(3x+40)$}\end{tcolorbox}\hfill\begin{tcolorbox}[enhanced,nobeforeafter,width=7.1cm,height=2.1cm,valign=center,colback=clit!6,colframe=clit,arc=2.5mm,boxrule=0.9pt,left=3mm]{\Large\bfseries\color{clit}d)}\quad{\Large $8(24x+5)$}\end{tcolorbox}
\end{tcolorbox}
\newpage
\begin{tcolorbox}[enhanced,width=16cm,height=14cm,colback=white,colframe=clit,boxrule=2pt,arc=0mm,
left=6mm,right=6mm,top=0mm,bottom=4mm,sharp corners,
overlay={\fill[clit] (frame.north west) rectangle ([yshift=-1.25cm]frame.north east);
\node[anchor=west,text=white,font=\sffamily\bfseries\LARGE] at ([xshift=5mm,yshift=-0.62cm]frame.north west) {Question 41 D};
\node[anchor=east,text=white,font=\sffamily\itshape\large] at ([xshift=-5mm,yshift=-0.62cm]frame.north east) {Calcul littéral et équations};}]
\vspace*{1.75cm}
\begin{minipage}[t][6.9cm][c]{\linewidth}\LARGE\centering
\begin{minipage}{0.95\linewidth}\raggedright\linespread{1.25}\selectfont Que vaut $-2x^2-x+7$ pour $x=3$~?\end{minipage}
\end{minipage}

\vfill
\begin{tcolorbox}[enhanced,nobeforeafter,width=7.1cm,height=2.1cm,valign=center,colback=clit!6,colframe=clit,arc=2.5mm,boxrule=0.9pt,left=3mm]{\Large\bfseries\color{clit}a)}\quad{\Large $-14$}\end{tcolorbox}\hfill\begin{tcolorbox}[enhanced,nobeforeafter,width=7.1cm,height=2.1cm,valign=center,colback=clit!6,colframe=clit,arc=2.5mm,boxrule=0.9pt,left=3mm]{\Large\bfseries\color{clit}b)}\quad{\Large $14$}\end{tcolorbox}\par\vspace{2.5mm}
\begin{tcolorbox}[enhanced,nobeforeafter,width=7.1cm,height=2.1cm,valign=center,colback=clit!6,colframe=clit,arc=2.5mm,boxrule=0.9pt,left=3mm]{\Large\bfseries\color{clit}c)}\quad{\Large $22$}\end{tcolorbox}\hfill\begin{tcolorbox}[enhanced,nobeforeafter,width=7.1cm,height=2.1cm,valign=center,colback=clit!6,colframe=clit,arc=2.5mm,boxrule=0.9pt,left=3mm]{\Large\bfseries\color{clit}d)}\quad{\Large $-8$}\end{tcolorbox}
\end{tcolorbox}
\newpage
\begin{tcolorbox}[enhanced,width=16cm,height=14cm,colback=white,colframe=clit,boxrule=2pt,arc=0mm,
left=6mm,right=6mm,top=0mm,bottom=4mm,sharp corners,
overlay={\fill[clit] (frame.north west) rectangle ([yshift=-1.25cm]frame.north east);
\node[anchor=west,text=white,font=\sffamily\bfseries\LARGE] at ([xshift=5mm,yshift=-0.62cm]frame.north west) {Question 42 D};
\node[anchor=east,text=white,font=\sffamily\itshape\large] at ([xshift=-5mm,yshift=-0.62cm]frame.north east) {Calcul littéral et équations};}]
\vspace*{1.75cm}
\begin{minipage}[t][6.9cm][c]{\linewidth}\LARGE\centering
\begin{minipage}{0.95\linewidth}\raggedright\linespread{1.25}\selectfont Que vaut $x^2+6x-2$ pour $x=-5$~?\end{minipage}
\end{minipage}

\vfill
\begin{tcolorbox}[enhanced,nobeforeafter,width=7.1cm,height=2.1cm,valign=center,colback=clit!6,colframe=clit,arc=2.5mm,boxrule=0.9pt,left=3mm]{\Large\bfseries\color{clit}a)}\quad{\Large $-42$}\end{tcolorbox}\hfill\begin{tcolorbox}[enhanced,nobeforeafter,width=7.1cm,height=2.1cm,valign=center,colback=clit!6,colframe=clit,arc=2.5mm,boxrule=0.9pt,left=3mm]{\Large\bfseries\color{clit}b)}\quad{\Large $-57$}\end{tcolorbox}\par\vspace{2.5mm}
\begin{tcolorbox}[enhanced,nobeforeafter,width=7.1cm,height=2.1cm,valign=center,colback=clit!6,colframe=clit,arc=2.5mm,boxrule=0.9pt,left=3mm]{\Large\bfseries\color{clit}c)}\quad{\Large $7$}\end{tcolorbox}\hfill\begin{tcolorbox}[enhanced,nobeforeafter,width=7.1cm,height=2.1cm,valign=center,colback=clit!6,colframe=clit,arc=2.5mm,boxrule=0.9pt,left=3mm]{\Large\bfseries\color{clit}d)}\quad{\Large $-7$}\end{tcolorbox}
\end{tcolorbox}
\newpage
\begin{tcolorbox}[enhanced,width=16cm,height=14cm,colback=white,colframe=clit,boxrule=2pt,arc=0mm,
left=6mm,right=6mm,top=0mm,bottom=4mm,sharp corners,
overlay={\fill[clit] (frame.north west) rectangle ([yshift=-1.25cm]frame.north east);
\node[anchor=west,text=white,font=\sffamily\bfseries\LARGE] at ([xshift=5mm,yshift=-0.62cm]frame.north west) {Question 43 D};
\node[anchor=east,text=white,font=\sffamily\itshape\large] at ([xshift=-5mm,yshift=-0.62cm]frame.north east) {Calcul littéral et équations};}]
\vspace*{1.75cm}
\begin{minipage}[t][6.9cm][c]{\linewidth}\LARGE\centering
\begin{minipage}{0.95\linewidth}\raggedright\linespread{1.25}\selectfont Quelle est la solution de l'équation $-3x+8=-7$~?\end{minipage}
\end{minipage}

\vfill
\begin{tcolorbox}[enhanced,nobeforeafter,width=7.1cm,height=2.1cm,valign=center,colback=clit!6,colframe=clit,arc=2.5mm,boxrule=0.9pt,left=3mm]{\Large\bfseries\color{clit}a)}\quad{\Large $x=5$}\end{tcolorbox}\hfill\begin{tcolorbox}[enhanced,nobeforeafter,width=7.1cm,height=2.1cm,valign=center,colback=clit!6,colframe=clit,arc=2.5mm,boxrule=0.9pt,left=3mm]{\Large\bfseries\color{clit}b)}\quad{\Large $x=-12$}\end{tcolorbox}\par\vspace{2.5mm}
\begin{tcolorbox}[enhanced,nobeforeafter,width=7.1cm,height=2.1cm,valign=center,colback=clit!6,colframe=clit,arc=2.5mm,boxrule=0.9pt,left=3mm]{\Large\bfseries\color{clit}c)}\quad{\Large $x=\dfrac{1}{5}$}\end{tcolorbox}\hfill\begin{tcolorbox}[enhanced,nobeforeafter,width=7.1cm,height=2.1cm,valign=center,colback=clit!6,colframe=clit,arc=2.5mm,boxrule=0.9pt,left=3mm]{\Large\bfseries\color{clit}d)}\quad{\Large $x=-\dfrac{1}{3}$}\end{tcolorbox}
\end{tcolorbox}
\newpage
\begin{tcolorbox}[enhanced,width=16cm,height=14cm,colback=white,colframe=clit,boxrule=2pt,arc=0mm,
left=6mm,right=6mm,top=0mm,bottom=4mm,sharp corners,
overlay={\fill[clit] (frame.north west) rectangle ([yshift=-1.25cm]frame.north east);
\node[anchor=west,text=white,font=\sffamily\bfseries\LARGE] at ([xshift=5mm,yshift=-0.62cm]frame.north west) {Question 44 D};
\node[anchor=east,text=white,font=\sffamily\itshape\large] at ([xshift=-5mm,yshift=-0.62cm]frame.north east) {Calcul littéral et équations};}]
\vspace*{1.75cm}
\begin{minipage}[t][6.9cm][c]{\linewidth}\LARGE\centering
\begin{minipage}{0.95\linewidth}\raggedright\linespread{1.25}\selectfont Quelle est la solution de l'équation $8x+3=9$~?\end{minipage}
\end{minipage}

\vfill
\begin{tcolorbox}[enhanced,nobeforeafter,width=7.1cm,height=2.1cm,valign=center,colback=clit!6,colframe=clit,arc=2.5mm,boxrule=0.9pt,left=3mm]{\Large\bfseries\color{clit}a)}\quad{\Large $x=\dfrac{4}{3}$}\end{tcolorbox}\hfill\begin{tcolorbox}[enhanced,nobeforeafter,width=7.1cm,height=2.1cm,valign=center,colback=clit!6,colframe=clit,arc=2.5mm,boxrule=0.9pt,left=3mm]{\Large\bfseries\color{clit}b)}\quad{\Large $x=-2$}\end{tcolorbox}\par\vspace{2.5mm}
\begin{tcolorbox}[enhanced,nobeforeafter,width=7.1cm,height=2.1cm,valign=center,colback=clit!6,colframe=clit,arc=2.5mm,boxrule=0.9pt,left=3mm]{\Large\bfseries\color{clit}c)}\quad{\Large $x=\dfrac{3}{4}$}\end{tcolorbox}\hfill\begin{tcolorbox}[enhanced,nobeforeafter,width=7.1cm,height=2.1cm,valign=center,colback=clit!6,colframe=clit,arc=2.5mm,boxrule=0.9pt,left=3mm]{\Large\bfseries\color{clit}d)}\quad{\Large $x=\dfrac{3}{2}$}\end{tcolorbox}
\end{tcolorbox}
\newpage
\begin{tcolorbox}[enhanced,width=16cm,height=14cm,colback=white,colframe=clit,boxrule=2pt,arc=0mm,
left=6mm,right=6mm,top=0mm,bottom=4mm,sharp corners,
overlay={\fill[clit] (frame.north west) rectangle ([yshift=-1.25cm]frame.north east);
\node[anchor=west,text=white,font=\sffamily\bfseries\LARGE] at ([xshift=5mm,yshift=-0.62cm]frame.north west) {Question 45 D};
\node[anchor=east,text=white,font=\sffamily\itshape\large] at ([xshift=-5mm,yshift=-0.62cm]frame.north east) {Calcul littéral et équations};}]
\vspace*{1.75cm}
\begin{minipage}[t][6.9cm][c]{\linewidth}\LARGE\centering
\begin{minipage}{0.95\linewidth}\raggedright\linespread{1.25}\selectfont Quelle est la solution de l'équation $8x+2=3x-13$~?\end{minipage}
\end{minipage}

\vfill
\begin{tcolorbox}[enhanced,nobeforeafter,width=7.1cm,height=2.1cm,valign=center,colback=clit!6,colframe=clit,arc=2.5mm,boxrule=0.9pt,left=3mm]{\Large\bfseries\color{clit}a)}\quad{\Large $x=-\dfrac{11}{5}$}\end{tcolorbox}\hfill\begin{tcolorbox}[enhanced,nobeforeafter,width=7.1cm,height=2.1cm,valign=center,colback=clit!6,colframe=clit,arc=2.5mm,boxrule=0.9pt,left=3mm]{\Large\bfseries\color{clit}b)}\quad{\Large $x=-\dfrac{15}{11}$}\end{tcolorbox}\par\vspace{2.5mm}
\begin{tcolorbox}[enhanced,nobeforeafter,width=7.1cm,height=2.1cm,valign=center,colback=clit!6,colframe=clit,arc=2.5mm,boxrule=0.9pt,left=3mm]{\Large\bfseries\color{clit}c)}\quad{\Large $x=3$}\end{tcolorbox}\hfill\begin{tcolorbox}[enhanced,nobeforeafter,width=7.1cm,height=2.1cm,valign=center,colback=clit!6,colframe=clit,arc=2.5mm,boxrule=0.9pt,left=3mm]{\Large\bfseries\color{clit}d)}\quad{\Large $x=-3$}\end{tcolorbox}
\end{tcolorbox}
\newpage
\begin{tcolorbox}[enhanced,width=16cm,height=14cm,colback=white,colframe=clit,boxrule=2pt,arc=0mm,
left=6mm,right=6mm,top=0mm,bottom=4mm,sharp corners,
overlay={\fill[clit] (frame.north west) rectangle ([yshift=-1.25cm]frame.north east);
\node[anchor=west,text=white,font=\sffamily\bfseries\LARGE] at ([xshift=5mm,yshift=-0.62cm]frame.north west) {Question 46 D};
\node[anchor=east,text=white,font=\sffamily\itshape\large] at ([xshift=-5mm,yshift=-0.62cm]frame.north east) {Calcul littéral et équations};}]
\vspace*{1.75cm}
\begin{minipage}[t][6.9cm][c]{\linewidth}\LARGE\centering
\begin{minipage}{0.95\linewidth}\raggedright\linespread{1.25}\selectfont Quelle est la solution de l'équation $-2x+7=4x-5$~?\end{minipage}
\end{minipage}

\vfill
\begin{tcolorbox}[enhanced,nobeforeafter,width=7.1cm,height=2.1cm,valign=center,colback=clit!6,colframe=clit,arc=2.5mm,boxrule=0.9pt,left=3mm]{\Large\bfseries\color{clit}a)}\quad{\Large $x=2$}\end{tcolorbox}\hfill\begin{tcolorbox}[enhanced,nobeforeafter,width=7.1cm,height=2.1cm,valign=center,colback=clit!6,colframe=clit,arc=2.5mm,boxrule=0.9pt,left=3mm]{\Large\bfseries\color{clit}b)}\quad{\Large $x=-6$}\end{tcolorbox}\par\vspace{2.5mm}
\begin{tcolorbox}[enhanced,nobeforeafter,width=7.1cm,height=2.1cm,valign=center,colback=clit!6,colframe=clit,arc=2.5mm,boxrule=0.9pt,left=3mm]{\Large\bfseries\color{clit}c)}\quad{\Large $x=-\dfrac{1}{3}$}\end{tcolorbox}\hfill\begin{tcolorbox}[enhanced,nobeforeafter,width=7.1cm,height=2.1cm,valign=center,colback=clit!6,colframe=clit,arc=2.5mm,boxrule=0.9pt,left=3mm]{\Large\bfseries\color{clit}d)}\quad{\Large $x=-2$}\end{tcolorbox}
\end{tcolorbox}
\newpage
\begin{tcolorbox}[enhanced,width=16cm,height=14cm,colback=white,colframe=clit,boxrule=2pt,arc=0mm,
left=6mm,right=6mm,top=0mm,bottom=4mm,sharp corners,
overlay={\fill[clit] (frame.north west) rectangle ([yshift=-1.25cm]frame.north east);
\node[anchor=west,text=white,font=\sffamily\bfseries\LARGE] at ([xshift=5mm,yshift=-0.62cm]frame.north west) {Question 47 D};
\node[anchor=east,text=white,font=\sffamily\itshape\large] at ([xshift=-5mm,yshift=-0.62cm]frame.north east) {Calcul littéral et équations};}]
\vspace*{1.75cm}
\begin{minipage}[t][6.9cm][c]{\linewidth}\LARGE\centering
\begin{minipage}{0.95\linewidth}\raggedright\linespread{1.25}\selectfont Le triple d'un nombre, augmenté de 11, est égal à $2$. Quel est ce nombre~?\end{minipage}
\end{minipage}

\vfill
\begin{tcolorbox}[enhanced,nobeforeafter,width=7.1cm,height=2.1cm,valign=center,colback=clit!6,colframe=clit,arc=2.5mm,boxrule=0.9pt,left=3mm]{\Large\bfseries\color{clit}a)}\quad{\Large $\dfrac{13}{3}$}\end{tcolorbox}\hfill\begin{tcolorbox}[enhanced,nobeforeafter,width=7.1cm,height=2.1cm,valign=center,colback=clit!6,colframe=clit,arc=2.5mm,boxrule=0.9pt,left=3mm]{\Large\bfseries\color{clit}b)}\quad{\Large $-3$}\end{tcolorbox}\par\vspace{2.5mm}
\begin{tcolorbox}[enhanced,nobeforeafter,width=7.1cm,height=2.1cm,valign=center,colback=clit!6,colframe=clit,arc=2.5mm,boxrule=0.9pt,left=3mm]{\Large\bfseries\color{clit}c)}\quad{\Large $-9$}\end{tcolorbox}\hfill\begin{tcolorbox}[enhanced,nobeforeafter,width=7.1cm,height=2.1cm,valign=center,colback=clit!6,colframe=clit,arc=2.5mm,boxrule=0.9pt,left=3mm]{\Large\bfseries\color{clit}d)}\quad{\Large $\dfrac{2}{3}$}\end{tcolorbox}
\end{tcolorbox}
\newpage
\begin{tcolorbox}[enhanced,width=16cm,height=14cm,colback=white,colframe=clit,boxrule=2pt,arc=0mm,
left=6mm,right=6mm,top=0mm,bottom=4mm,sharp corners,
overlay={\fill[clit] (frame.north west) rectangle ([yshift=-1.25cm]frame.north east);
\node[anchor=west,text=white,font=\sffamily\bfseries\LARGE] at ([xshift=5mm,yshift=-0.62cm]frame.north west) {Question 48 D};
\node[anchor=east,text=white,font=\sffamily\itshape\large] at ([xshift=-5mm,yshift=-0.62cm]frame.north east) {Calcul littéral et équations};}]
\vspace*{1.75cm}
\begin{minipage}[t][6.9cm][c]{\linewidth}\LARGE\centering
\begin{minipage}{0.95\linewidth}\raggedright\linespread{1.25}\selectfont Le double d'un nombre, diminué de 5, est égal à $-13$. Quel est ce nombre~?\end{minipage}
\end{minipage}

\vfill
\begin{tcolorbox}[enhanced,nobeforeafter,width=7.1cm,height=2.1cm,valign=center,colback=clit!6,colframe=clit,arc=2.5mm,boxrule=0.9pt,left=3mm]{\Large\bfseries\color{clit}a)}\quad{\Large $-\dfrac{13}{2}$}\end{tcolorbox}\hfill\begin{tcolorbox}[enhanced,nobeforeafter,width=7.1cm,height=2.1cm,valign=center,colback=clit!6,colframe=clit,arc=2.5mm,boxrule=0.9pt,left=3mm]{\Large\bfseries\color{clit}b)}\quad{\Large $-8$}\end{tcolorbox}\par\vspace{2.5mm}
\begin{tcolorbox}[enhanced,nobeforeafter,width=7.1cm,height=2.1cm,valign=center,colback=clit!6,colframe=clit,arc=2.5mm,boxrule=0.9pt,left=3mm]{\Large\bfseries\color{clit}c)}\quad{\Large $-4$}\end{tcolorbox}\hfill\begin{tcolorbox}[enhanced,nobeforeafter,width=7.1cm,height=2.1cm,valign=center,colback=clit!6,colframe=clit,arc=2.5mm,boxrule=0.9pt,left=3mm]{\Large\bfseries\color{clit}d)}\quad{\Large $-9$}\end{tcolorbox}
\end{tcolorbox}
\newpage
\begin{tcolorbox}[enhanced,width=16cm,height=14cm,colback=white,colframe=cgeo,boxrule=2pt,arc=0mm,
left=6mm,right=6mm,top=0mm,bottom=4mm,sharp corners,
overlay={\fill[cgeo] (frame.north west) rectangle ([yshift=-1.25cm]frame.north east);
\node[anchor=west,text=white,font=\sffamily\bfseries\LARGE] at ([xshift=5mm,yshift=-0.62cm]frame.north west) {Question 49 D};
\node[anchor=east,text=white,font=\sffamily\itshape\large] at ([xshift=-5mm,yshift=-0.62cm]frame.north east) {Géométrie};}]
\vspace*{1.75cm}
\begin{minipage}[t][6.9cm][c]{\linewidth}\LARGE\centering
\begin{minipage}{0.95\linewidth}\raggedright\linespread{1.25}\selectfont Le triangle $ABC$ est rectangle en $A$, avec $AB=12$ cm et $AC=16$ cm. Combien mesure $BC$~?\end{minipage}
\end{minipage}

\vfill
\begin{tcolorbox}[enhanced,nobeforeafter,width=7.1cm,height=2.1cm,valign=center,colback=cgeo!6,colframe=cgeo,arc=2.5mm,boxrule=0.9pt,left=3mm]{\Large\bfseries\color{cgeo}a)}\quad{\Large $400~\text{cm}$}\end{tcolorbox}\hfill\begin{tcolorbox}[enhanced,nobeforeafter,width=7.1cm,height=2.1cm,valign=center,colback=cgeo!6,colframe=cgeo,arc=2.5mm,boxrule=0.9pt,left=3mm]{\Large\bfseries\color{cgeo}b)}\quad{\Large $28~\text{cm}$}\end{tcolorbox}\par\vspace{2.5mm}
\begin{tcolorbox}[enhanced,nobeforeafter,width=7.1cm,height=2.1cm,valign=center,colback=cgeo!6,colframe=cgeo,arc=2.5mm,boxrule=0.9pt,left=3mm]{\Large\bfseries\color{cgeo}c)}\quad{\Large $20~\text{cm}$}\end{tcolorbox}\hfill\begin{tcolorbox}[enhanced,nobeforeafter,width=7.1cm,height=2.1cm,valign=center,colback=cgeo!6,colframe=cgeo,arc=2.5mm,boxrule=0.9pt,left=3mm]{\Large\bfseries\color{cgeo}d)}\quad{\Large $\sqrt{112}~\text{cm}$}\end{tcolorbox}
\end{tcolorbox}
\newpage
\begin{tcolorbox}[enhanced,width=16cm,height=14cm,colback=white,colframe=cgeo,boxrule=2pt,arc=0mm,
left=6mm,right=6mm,top=0mm,bottom=4mm,sharp corners,
overlay={\fill[cgeo] (frame.north west) rectangle ([yshift=-1.25cm]frame.north east);
\node[anchor=west,text=white,font=\sffamily\bfseries\LARGE] at ([xshift=5mm,yshift=-0.62cm]frame.north west) {Question 50 D};
\node[anchor=east,text=white,font=\sffamily\itshape\large] at ([xshift=-5mm,yshift=-0.62cm]frame.north east) {Géométrie};}]
\vspace*{1.75cm}
\begin{minipage}[t][6.9cm][c]{\linewidth}\LARGE\centering
\begin{minipage}{0.95\linewidth}\raggedright\linespread{1.25}\selectfont Le triangle $ABC$ est rectangle en $A$, avec $AB=2$ cm et $AC=3$ cm. Combien mesure $BC$~?\end{minipage}
\end{minipage}

\vfill
\begin{tcolorbox}[enhanced,nobeforeafter,width=7.1cm,height=2.1cm,valign=center,colback=cgeo!6,colframe=cgeo,arc=2.5mm,boxrule=0.9pt,left=3mm]{\Large\bfseries\color{cgeo}a)}\quad{\Large $\sqrt{5}~\text{cm}$}\end{tcolorbox}\hfill\begin{tcolorbox}[enhanced,nobeforeafter,width=7.1cm,height=2.1cm,valign=center,colback=cgeo!6,colframe=cgeo,arc=2.5mm,boxrule=0.9pt,left=3mm]{\Large\bfseries\color{cgeo}b)}\quad{\Large $13~\text{cm}$}\end{tcolorbox}\par\vspace{2.5mm}
\begin{tcolorbox}[enhanced,nobeforeafter,width=7.1cm,height=2.1cm,valign=center,colback=cgeo!6,colframe=cgeo,arc=2.5mm,boxrule=0.9pt,left=3mm]{\Large\bfseries\color{cgeo}c)}\quad{\Large $5~\text{cm}$}\end{tcolorbox}\hfill\begin{tcolorbox}[enhanced,nobeforeafter,width=7.1cm,height=2.1cm,valign=center,colback=cgeo!6,colframe=cgeo,arc=2.5mm,boxrule=0.9pt,left=3mm]{\Large\bfseries\color{cgeo}d)}\quad{\Large $\sqrt{13}~\text{cm}$}\end{tcolorbox}
\end{tcolorbox}
\newpage
\begin{tcolorbox}[enhanced,width=16cm,height=14cm,colback=white,colframe=cgeo,boxrule=2pt,arc=0mm,
left=6mm,right=6mm,top=0mm,bottom=4mm,sharp corners,
overlay={\fill[cgeo] (frame.north west) rectangle ([yshift=-1.25cm]frame.north east);
\node[anchor=west,text=white,font=\sffamily\bfseries\LARGE] at ([xshift=5mm,yshift=-0.62cm]frame.north west) {Question 51 D};
\node[anchor=east,text=white,font=\sffamily\itshape\large] at ([xshift=-5mm,yshift=-0.62cm]frame.north east) {Géométrie};}]
\vspace*{1.75cm}
\begin{minipage}[t][6.9cm][c]{\linewidth}\LARGE\centering
\begin{minipage}{0.95\linewidth}\raggedright\linespread{1.25}\selectfont Le triangle $ABC$ est rectangle en $A$, avec $BC=5$ cm et $AB=3$ cm. Combien mesure $AC$~?\end{minipage}
\end{minipage}

\vfill
\begin{tcolorbox}[enhanced,nobeforeafter,width=7.1cm,height=2.1cm,valign=center,colback=cgeo!6,colframe=cgeo,arc=2.5mm,boxrule=0.9pt,left=3mm]{\Large\bfseries\color{cgeo}a)}\quad{\Large $4~\text{cm}$}\end{tcolorbox}\hfill\begin{tcolorbox}[enhanced,nobeforeafter,width=7.1cm,height=2.1cm,valign=center,colback=cgeo!6,colframe=cgeo,arc=2.5mm,boxrule=0.9pt,left=3mm]{\Large\bfseries\color{cgeo}b)}\quad{\Large $2~\text{cm}$}\end{tcolorbox}\par\vspace{2.5mm}
\begin{tcolorbox}[enhanced,nobeforeafter,width=7.1cm,height=2.1cm,valign=center,colback=cgeo!6,colframe=cgeo,arc=2.5mm,boxrule=0.9pt,left=3mm]{\Large\bfseries\color{cgeo}c)}\quad{\Large $16~\text{cm}$}\end{tcolorbox}\hfill\begin{tcolorbox}[enhanced,nobeforeafter,width=7.1cm,height=2.1cm,valign=center,colback=cgeo!6,colframe=cgeo,arc=2.5mm,boxrule=0.9pt,left=3mm]{\Large\bfseries\color{cgeo}d)}\quad{\Large $\sqrt{34}~\text{cm}$}\end{tcolorbox}
\end{tcolorbox}
\newpage
\begin{tcolorbox}[enhanced,width=16cm,height=14cm,colback=white,colframe=cgeo,boxrule=2pt,arc=0mm,
left=6mm,right=6mm,top=0mm,bottom=4mm,sharp corners,
overlay={\fill[cgeo] (frame.north west) rectangle ([yshift=-1.25cm]frame.north east);
\node[anchor=west,text=white,font=\sffamily\bfseries\LARGE] at ([xshift=5mm,yshift=-0.62cm]frame.north west) {Question 52 D};
\node[anchor=east,text=white,font=\sffamily\itshape\large] at ([xshift=-5mm,yshift=-0.62cm]frame.north east) {Géométrie};}]
\vspace*{1.75cm}
\begin{minipage}[t][6.9cm][c]{\linewidth}\LARGE\centering
\begin{minipage}{0.95\linewidth}\raggedright\linespread{1.25}\selectfont Le triangle $ABC$ est rectangle en $A$, avec $BC=6$ cm et $AB=4$ cm. Combien mesure $AC$~?\end{minipage}
\end{minipage}

\vfill
\begin{tcolorbox}[enhanced,nobeforeafter,width=7.1cm,height=2.1cm,valign=center,colback=cgeo!6,colframe=cgeo,arc=2.5mm,boxrule=0.9pt,left=3mm]{\Large\bfseries\color{cgeo}a)}\quad{\Large $20~\text{cm}$}\end{tcolorbox}\hfill\begin{tcolorbox}[enhanced,nobeforeafter,width=7.1cm,height=2.1cm,valign=center,colback=cgeo!6,colframe=cgeo,arc=2.5mm,boxrule=0.9pt,left=3mm]{\Large\bfseries\color{cgeo}b)}\quad{\Large $\sqrt{52}~\text{cm}$}\end{tcolorbox}\par\vspace{2.5mm}
\begin{tcolorbox}[enhanced,nobeforeafter,width=7.1cm,height=2.1cm,valign=center,colback=cgeo!6,colframe=cgeo,arc=2.5mm,boxrule=0.9pt,left=3mm]{\Large\bfseries\color{cgeo}c)}\quad{\Large $2~\text{cm}$}\end{tcolorbox}\hfill\begin{tcolorbox}[enhanced,nobeforeafter,width=7.1cm,height=2.1cm,valign=center,colback=cgeo!6,colframe=cgeo,arc=2.5mm,boxrule=0.9pt,left=3mm]{\Large\bfseries\color{cgeo}d)}\quad{\Large $\sqrt{20}~\text{cm}$}\end{tcolorbox}
\end{tcolorbox}
\newpage
\begin{tcolorbox}[enhanced,width=16cm,height=14cm,colback=white,colframe=cgeo,boxrule=2pt,arc=0mm,
left=6mm,right=6mm,top=0mm,bottom=4mm,sharp corners,
overlay={\fill[cgeo] (frame.north west) rectangle ([yshift=-1.25cm]frame.north east);
\node[anchor=west,text=white,font=\sffamily\bfseries\LARGE] at ([xshift=5mm,yshift=-0.62cm]frame.north west) {Question 53 D};
\node[anchor=east,text=white,font=\sffamily\itshape\large] at ([xshift=-5mm,yshift=-0.62cm]frame.north east) {Géométrie};}]
\vspace*{1.75cm}
\begin{minipage}[t][6.9cm][c]{\linewidth}\LARGE\centering
\begin{minipage}{0.95\linewidth}\raggedright\linespread{1.25}\selectfont Un triangle a des côtés de $20$ cm, $21$ cm et $29$ cm. Est-il rectangle~?\end{minipage}
\end{minipage}

\vfill
\begin{tcolorbox}[enhanced,nobeforeafter,width=7.1cm,height=2.1cm,valign=center,colback=cgeo!6,colframe=cgeo,arc=2.5mm,boxrule=0.9pt,left=3mm]{\Large\bfseries\color{cgeo}a)}\quad{\Large oui}\end{tcolorbox}\hfill\begin{tcolorbox}[enhanced,nobeforeafter,width=7.1cm,height=2.1cm,valign=center,colback=cgeo!6,colframe=cgeo,arc=2.5mm,boxrule=0.9pt,left=3mm]{\Large\bfseries\color{cgeo}b)}\quad{\Large non}\end{tcolorbox}\par\vspace{2.5mm}
\begin{tcolorbox}[enhanced,nobeforeafter,width=7.1cm,height=2.1cm,valign=center,colback=cgeo!6,colframe=cgeo,arc=2.5mm,boxrule=0.9pt,left=3mm]{\Large\bfseries\color{cgeo}c)}\quad{\Large s'il est isocèle}\end{tcolorbox}\hfill\begin{tcolorbox}[enhanced,nobeforeafter,width=7.1cm,height=2.1cm,valign=center,colback=cgeo!6,colframe=cgeo,arc=2.5mm,boxrule=0.9pt,left=3mm]{\Large\bfseries\color{cgeo}d)}\quad{\Large on ne peut pas savoir}\end{tcolorbox}
\end{tcolorbox}
\newpage
\begin{tcolorbox}[enhanced,width=16cm,height=14cm,colback=white,colframe=cgeo,boxrule=2pt,arc=0mm,
left=6mm,right=6mm,top=0mm,bottom=4mm,sharp corners,
overlay={\fill[cgeo] (frame.north west) rectangle ([yshift=-1.25cm]frame.north east);
\node[anchor=west,text=white,font=\sffamily\bfseries\LARGE] at ([xshift=5mm,yshift=-0.62cm]frame.north west) {Question 54 D};
\node[anchor=east,text=white,font=\sffamily\itshape\large] at ([xshift=-5mm,yshift=-0.62cm]frame.north east) {Géométrie};}]
\vspace*{1.75cm}
\begin{minipage}[t][6.9cm][c]{\linewidth}\LARGE\centering
\begin{minipage}{0.95\linewidth}\raggedright\linespread{1.25}\selectfont Un triangle a des côtés de $7$ cm, $10$ cm et $12$ cm. Est-il rectangle~?\end{minipage}
\end{minipage}

\vfill
\begin{tcolorbox}[enhanced,nobeforeafter,width=7.1cm,height=2.1cm,valign=center,colback=cgeo!6,colframe=cgeo,arc=2.5mm,boxrule=0.9pt,left=3mm]{\Large\bfseries\color{cgeo}a)}\quad{\Large oui}\end{tcolorbox}\hfill\begin{tcolorbox}[enhanced,nobeforeafter,width=7.1cm,height=2.1cm,valign=center,colback=cgeo!6,colframe=cgeo,arc=2.5mm,boxrule=0.9pt,left=3mm]{\Large\bfseries\color{cgeo}b)}\quad{\Large non}\end{tcolorbox}\par\vspace{2.5mm}
\begin{tcolorbox}[enhanced,nobeforeafter,width=7.1cm,height=2.1cm,valign=center,colback=cgeo!6,colframe=cgeo,arc=2.5mm,boxrule=0.9pt,left=3mm]{\Large\bfseries\color{cgeo}c)}\quad{\Large s'il est isocèle}\end{tcolorbox}\hfill\begin{tcolorbox}[enhanced,nobeforeafter,width=7.1cm,height=2.1cm,valign=center,colback=cgeo!6,colframe=cgeo,arc=2.5mm,boxrule=0.9pt,left=3mm]{\Large\bfseries\color{cgeo}d)}\quad{\Large on ne peut pas savoir}\end{tcolorbox}
\end{tcolorbox}
\newpage
\begin{tcolorbox}[enhanced,width=16cm,height=14cm,colback=white,colframe=cgeo,boxrule=2pt,arc=0mm,
left=6mm,right=6mm,top=0mm,bottom=4mm,sharp corners,
overlay={\fill[cgeo] (frame.north west) rectangle ([yshift=-1.25cm]frame.north east);
\node[anchor=west,text=white,font=\sffamily\bfseries\LARGE] at ([xshift=5mm,yshift=-0.62cm]frame.north west) {Question 55 D};
\node[anchor=east,text=white,font=\sffamily\itshape\large] at ([xshift=-5mm,yshift=-0.62cm]frame.north east) {Géométrie};}]
\vspace*{1.75cm}
\begin{minipage}[t][6.9cm][c]{\linewidth}\Large\centering
\begin{minipage}{0.95\linewidth}\raggedright\linespread{1.25}\selectfont Dans le triangle $ABC$, $M$ est sur $[AB]$, $N$ est sur $[AC]$ et les droites $(MN)$ et $(BC)$ sont parallèles. $AM=7$ cm, $AB=14$ cm et $AN=9$ cm. Combien mesure $AC$~?\end{minipage}
\end{minipage}

\vfill
\begin{tcolorbox}[enhanced,nobeforeafter,width=7.1cm,height=2.1cm,valign=center,colback=cgeo!6,colframe=cgeo,arc=2.5mm,boxrule=0.9pt,left=3mm]{\Large\bfseries\color{cgeo}a)}\quad{\Large $18~\text{cm}$}\end{tcolorbox}\hfill\begin{tcolorbox}[enhanced,nobeforeafter,width=7.1cm,height=2.1cm,valign=center,colback=cgeo!6,colframe=cgeo,arc=2.5mm,boxrule=0.9pt,left=3mm]{\Large\bfseries\color{cgeo}b)}\quad{\Large $16~\text{cm}$}\end{tcolorbox}\par\vspace{2.5mm}
\begin{tcolorbox}[enhanced,nobeforeafter,width=7.1cm,height=2.1cm,valign=center,colback=cgeo!6,colframe=cgeo,arc=2.5mm,boxrule=0.9pt,left=3mm]{\Large\bfseries\color{cgeo}c)}\quad{\Large $4{,}5~\text{cm}$}\end{tcolorbox}\hfill\begin{tcolorbox}[enhanced,nobeforeafter,width=7.1cm,height=2.1cm,valign=center,colback=cgeo!6,colframe=cgeo,arc=2.5mm,boxrule=0.9pt,left=3mm]{\Large\bfseries\color{cgeo}d)}\quad{\Large $\dfrac{98}{9}~\text{cm}$}\end{tcolorbox}
\end{tcolorbox}
\newpage
\begin{tcolorbox}[enhanced,width=16cm,height=14cm,colback=white,colframe=cgeo,boxrule=2pt,arc=0mm,
left=6mm,right=6mm,top=0mm,bottom=4mm,sharp corners,
overlay={\fill[cgeo] (frame.north west) rectangle ([yshift=-1.25cm]frame.north east);
\node[anchor=west,text=white,font=\sffamily\bfseries\LARGE] at ([xshift=5mm,yshift=-0.62cm]frame.north west) {Question 56 D};
\node[anchor=east,text=white,font=\sffamily\itshape\large] at ([xshift=-5mm,yshift=-0.62cm]frame.north east) {Géométrie};}]
\vspace*{1.75cm}
\begin{minipage}[t][6.9cm][c]{\linewidth}\Large\centering
\begin{minipage}{0.95\linewidth}\raggedright\linespread{1.25}\selectfont Dans le triangle $ABC$, $M$ est sur $[AB]$, $N$ est sur $[AC]$ et les droites $(MN)$ et $(BC)$ sont parallèles. $AM=4$ cm, $AB=6$ cm et $AN=8$ cm. Combien mesure $AC$~?\end{minipage}
\end{minipage}

\vfill
\begin{tcolorbox}[enhanced,nobeforeafter,width=7.1cm,height=2.1cm,valign=center,colback=cgeo!6,colframe=cgeo,arc=2.5mm,boxrule=0.9pt,left=3mm]{\Large\bfseries\color{cgeo}a)}\quad{\Large $10~\text{cm}$}\end{tcolorbox}\hfill\begin{tcolorbox}[enhanced,nobeforeafter,width=7.1cm,height=2.1cm,valign=center,colback=cgeo!6,colframe=cgeo,arc=2.5mm,boxrule=0.9pt,left=3mm]{\Large\bfseries\color{cgeo}b)}\quad{\Large $12~\text{cm}$}\end{tcolorbox}\par\vspace{2.5mm}
\begin{tcolorbox}[enhanced,nobeforeafter,width=7.1cm,height=2.1cm,valign=center,colback=cgeo!6,colframe=cgeo,arc=2.5mm,boxrule=0.9pt,left=3mm]{\Large\bfseries\color{cgeo}c)}\quad{\Large $3~\text{cm}$}\end{tcolorbox}\hfill\begin{tcolorbox}[enhanced,nobeforeafter,width=7.1cm,height=2.1cm,valign=center,colback=cgeo!6,colframe=cgeo,arc=2.5mm,boxrule=0.9pt,left=3mm]{\Large\bfseries\color{cgeo}d)}\quad{\Large $\dfrac{16}{3}~\text{cm}$}\end{tcolorbox}
\end{tcolorbox}
\newpage
\begin{tcolorbox}[enhanced,width=16cm,height=14cm,colback=white,colframe=cgeo,boxrule=2pt,arc=0mm,
left=6mm,right=6mm,top=0mm,bottom=4mm,sharp corners,
overlay={\fill[cgeo] (frame.north west) rectangle ([yshift=-1.25cm]frame.north east);
\node[anchor=west,text=white,font=\sffamily\bfseries\LARGE] at ([xshift=5mm,yshift=-0.62cm]frame.north west) {Question 57 D};
\node[anchor=east,text=white,font=\sffamily\itshape\large] at ([xshift=-5mm,yshift=-0.62cm]frame.north east) {Géométrie};}]
\vspace*{1.75cm}
\begin{minipage}[t][6.9cm][c]{\linewidth}\LARGE\centering
\begin{minipage}{0.95\linewidth}\raggedright\linespread{1.25}\selectfont Le triangle $ABC$ est rectangle en $A$, avec $BC=8$ cm et $AB=6$ cm. Que vaut $\cos\widehat{ABC}$~?\end{minipage}
\end{minipage}

\vfill
\begin{tcolorbox}[enhanced,nobeforeafter,width=7.1cm,height=2.1cm,valign=center,colback=cgeo!6,colframe=cgeo,arc=2.5mm,boxrule=0.9pt,left=3mm]{\Large\bfseries\color{cgeo}a)}\quad{\Large $\dfrac{1}{4}$}\end{tcolorbox}\hfill\begin{tcolorbox}[enhanced,nobeforeafter,width=7.1cm,height=2.1cm,valign=center,colback=cgeo!6,colframe=cgeo,arc=2.5mm,boxrule=0.9pt,left=3mm]{\Large\bfseries\color{cgeo}b)}\quad{\Large $\dfrac{3}{4}$}\end{tcolorbox}\par\vspace{2.5mm}
\begin{tcolorbox}[enhanced,nobeforeafter,width=7.1cm,height=2.1cm,valign=center,colback=cgeo!6,colframe=cgeo,arc=2.5mm,boxrule=0.9pt,left=3mm]{\Large\bfseries\color{cgeo}c)}\quad{\Large $3$}\end{tcolorbox}\hfill\begin{tcolorbox}[enhanced,nobeforeafter,width=7.1cm,height=2.1cm,valign=center,colback=cgeo!6,colframe=cgeo,arc=2.5mm,boxrule=0.9pt,left=3mm]{\Large\bfseries\color{cgeo}d)}\quad{\Large $\dfrac{4}{3}$}\end{tcolorbox}
\end{tcolorbox}
\newpage
\begin{tcolorbox}[enhanced,width=16cm,height=14cm,colback=white,colframe=cgeo,boxrule=2pt,arc=0mm,
left=6mm,right=6mm,top=0mm,bottom=4mm,sharp corners,
overlay={\fill[cgeo] (frame.north west) rectangle ([yshift=-1.25cm]frame.north east);
\node[anchor=west,text=white,font=\sffamily\bfseries\LARGE] at ([xshift=5mm,yshift=-0.62cm]frame.north west) {Question 58 D};
\node[anchor=east,text=white,font=\sffamily\itshape\large] at ([xshift=-5mm,yshift=-0.62cm]frame.north east) {Géométrie};}]
\vspace*{1.75cm}
\begin{minipage}[t][6.9cm][c]{\linewidth}\LARGE\centering
\begin{minipage}{0.95\linewidth}\raggedright\linespread{1.25}\selectfont Le triangle $ABC$ est rectangle en $A$, avec $BC=9$ cm et $AB=4$ cm. Que vaut $\cos\widehat{ABC}$~?\end{minipage}
\end{minipage}

\vfill
\begin{tcolorbox}[enhanced,nobeforeafter,width=7.1cm,height=2.1cm,valign=center,colback=cgeo!6,colframe=cgeo,arc=2.5mm,boxrule=0.9pt,left=3mm]{\Large\bfseries\color{cgeo}a)}\quad{\Large $\dfrac{9}{4}$}\end{tcolorbox}\hfill\begin{tcolorbox}[enhanced,nobeforeafter,width=7.1cm,height=2.1cm,valign=center,colback=cgeo!6,colframe=cgeo,arc=2.5mm,boxrule=0.9pt,left=3mm]{\Large\bfseries\color{cgeo}b)}\quad{\Large $\dfrac{4}{9}$}\end{tcolorbox}\par\vspace{2.5mm}
\begin{tcolorbox}[enhanced,nobeforeafter,width=7.1cm,height=2.1cm,valign=center,colback=cgeo!6,colframe=cgeo,arc=2.5mm,boxrule=0.9pt,left=3mm]{\Large\bfseries\color{cgeo}c)}\quad{\Large $\dfrac{5}{9}$}\end{tcolorbox}\hfill\begin{tcolorbox}[enhanced,nobeforeafter,width=7.1cm,height=2.1cm,valign=center,colback=cgeo!6,colframe=cgeo,arc=2.5mm,boxrule=0.9pt,left=3mm]{\Large\bfseries\color{cgeo}d)}\quad{\Large $\dfrac{4}{5}$}\end{tcolorbox}
\end{tcolorbox}
\newpage
\begin{tcolorbox}[enhanced,width=16cm,height=14cm,colback=white,colframe=cgeo,boxrule=2pt,arc=0mm,
left=6mm,right=6mm,top=0mm,bottom=4mm,sharp corners,
overlay={\fill[cgeo] (frame.north west) rectangle ([yshift=-1.25cm]frame.north east);
\node[anchor=west,text=white,font=\sffamily\bfseries\LARGE] at ([xshift=5mm,yshift=-0.62cm]frame.north west) {Question 59 D};
\node[anchor=east,text=white,font=\sffamily\itshape\large] at ([xshift=-5mm,yshift=-0.62cm]frame.north east) {Géométrie};}]
\vspace*{1.75cm}
\begin{minipage}[t][6.9cm][c]{\linewidth}\LARGE\centering
\begin{minipage}{0.95\linewidth}\raggedright\linespread{1.25}\selectfont Le triangle $ABC$ est rectangle en $A$, avec $BC=40$ cm et $\cos\widehat{ABC}=0{,}45$. Combien mesure $AB$~?\end{minipage}
\end{minipage}

\vfill
\begin{tcolorbox}[enhanced,nobeforeafter,width=7.1cm,height=2.1cm,valign=center,colback=cgeo!6,colframe=cgeo,arc=2.5mm,boxrule=0.9pt,left=3mm]{\Large\bfseries\color{cgeo}a)}\quad{\Large $22~\text{cm}$}\end{tcolorbox}\hfill\begin{tcolorbox}[enhanced,nobeforeafter,width=7.1cm,height=2.1cm,valign=center,colback=cgeo!6,colframe=cgeo,arc=2.5mm,boxrule=0.9pt,left=3mm]{\Large\bfseries\color{cgeo}b)}\quad{\Large $9~\text{cm}$}\end{tcolorbox}\par\vspace{2.5mm}
\begin{tcolorbox}[enhanced,nobeforeafter,width=7.1cm,height=2.1cm,valign=center,colback=cgeo!6,colframe=cgeo,arc=2.5mm,boxrule=0.9pt,left=3mm]{\Large\bfseries\color{cgeo}c)}\quad{\Large $19~\text{cm}$}\end{tcolorbox}\hfill\begin{tcolorbox}[enhanced,nobeforeafter,width=7.1cm,height=2.1cm,valign=center,colback=cgeo!6,colframe=cgeo,arc=2.5mm,boxrule=0.9pt,left=3mm]{\Large\bfseries\color{cgeo}d)}\quad{\Large $18~\text{cm}$}\end{tcolorbox}
\end{tcolorbox}
\newpage
\begin{tcolorbox}[enhanced,width=16cm,height=14cm,colback=white,colframe=cgeo,boxrule=2pt,arc=0mm,
left=6mm,right=6mm,top=0mm,bottom=4mm,sharp corners,
overlay={\fill[cgeo] (frame.north west) rectangle ([yshift=-1.25cm]frame.north east);
\node[anchor=west,text=white,font=\sffamily\bfseries\LARGE] at ([xshift=5mm,yshift=-0.62cm]frame.north west) {Question 60 D};
\node[anchor=east,text=white,font=\sffamily\itshape\large] at ([xshift=-5mm,yshift=-0.62cm]frame.north east) {Géométrie};}]
\vspace*{1.75cm}
\begin{minipage}[t][6.9cm][c]{\linewidth}\LARGE\centering
\begin{minipage}{0.95\linewidth}\raggedright\linespread{1.25}\selectfont Le triangle $ABC$ est rectangle en $A$, avec $BC=16$ cm et $\cos\widehat{ABC}=0{,}625$. Combien mesure $AB$~?\end{minipage}
\end{minipage}

\vfill
\begin{tcolorbox}[enhanced,nobeforeafter,width=7.1cm,height=2.1cm,valign=center,colback=cgeo!6,colframe=cgeo,arc=2.5mm,boxrule=0.9pt,left=3mm]{\Large\bfseries\color{cgeo}a)}\quad{\Large $10~\text{cm}$}\end{tcolorbox}\hfill\begin{tcolorbox}[enhanced,nobeforeafter,width=7.1cm,height=2.1cm,valign=center,colback=cgeo!6,colframe=cgeo,arc=2.5mm,boxrule=0.9pt,left=3mm]{\Large\bfseries\color{cgeo}b)}\quad{\Large $11~\text{cm}$}\end{tcolorbox}\par\vspace{2.5mm}
\begin{tcolorbox}[enhanced,nobeforeafter,width=7.1cm,height=2.1cm,valign=center,colback=cgeo!6,colframe=cgeo,arc=2.5mm,boxrule=0.9pt,left=3mm]{\Large\bfseries\color{cgeo}c)}\quad{\Large $5~\text{cm}$}\end{tcolorbox}\hfill\begin{tcolorbox}[enhanced,nobeforeafter,width=7.1cm,height=2.1cm,valign=center,colback=cgeo!6,colframe=cgeo,arc=2.5mm,boxrule=0.9pt,left=3mm]{\Large\bfseries\color{cgeo}d)}\quad{\Large $6~\text{cm}$}\end{tcolorbox}
\end{tcolorbox}
\newpage
\begin{tcolorbox}[enhanced,width=16cm,height=14cm,colback=white,colframe=cgeo,boxrule=2pt,arc=0mm,
left=6mm,right=6mm,top=0mm,bottom=4mm,sharp corners,
overlay={\fill[cgeo] (frame.north west) rectangle ([yshift=-1.25cm]frame.north east);
\node[anchor=west,text=white,font=\sffamily\bfseries\LARGE] at ([xshift=5mm,yshift=-0.62cm]frame.north west) {Question 61 D};
\node[anchor=east,text=white,font=\sffamily\itshape\large] at ([xshift=-5mm,yshift=-0.62cm]frame.north east) {Géométrie};}]
\vspace*{1.75cm}
\begin{minipage}[t][6.9cm][c]{\linewidth}\Large\centering
\begin{minipage}{0.95\linewidth}\raggedright\linespread{1.25}\selectfont Dans un repère, $A(1\,;\,-3)$, $B(5\,;\,-1)$ et $C(-2\,;\,-4)$. Quelles sont les coordonnées de l'image de $C$ par la translation qui transforme $A$ en $B$~?\end{minipage}
\end{minipage}

\vfill
\begin{tcolorbox}[enhanced,nobeforeafter,width=7.1cm,height=2.1cm,valign=center,colback=cgeo!6,colframe=cgeo,arc=2.5mm,boxrule=0.9pt,left=3mm]{\Large\bfseries\color{cgeo}a)}\quad{\Large $(3\,;\,-5)$}\end{tcolorbox}\hfill\begin{tcolorbox}[enhanced,nobeforeafter,width=7.1cm,height=2.1cm,valign=center,colback=cgeo!6,colframe=cgeo,arc=2.5mm,boxrule=0.9pt,left=3mm]{\Large\bfseries\color{cgeo}b)}\quad{\Large $(0\,;\,0)$}\end{tcolorbox}\par\vspace{2.5mm}
\begin{tcolorbox}[enhanced,nobeforeafter,width=7.1cm,height=2.1cm,valign=center,colback=cgeo!6,colframe=cgeo,arc=2.5mm,boxrule=0.9pt,left=3mm]{\Large\bfseries\color{cgeo}c)}\quad{\Large $(-6\,;\,-6)$}\end{tcolorbox}\hfill\begin{tcolorbox}[enhanced,nobeforeafter,width=7.1cm,height=2.1cm,valign=center,colback=cgeo!6,colframe=cgeo,arc=2.5mm,boxrule=0.9pt,left=3mm]{\Large\bfseries\color{cgeo}d)}\quad{\Large $(2\,;\,-2)$}\end{tcolorbox}
\end{tcolorbox}
\newpage
\begin{tcolorbox}[enhanced,width=16cm,height=14cm,colback=white,colframe=cgeo,boxrule=2pt,arc=0mm,
left=6mm,right=6mm,top=0mm,bottom=4mm,sharp corners,
overlay={\fill[cgeo] (frame.north west) rectangle ([yshift=-1.25cm]frame.north east);
\node[anchor=west,text=white,font=\sffamily\bfseries\LARGE] at ([xshift=5mm,yshift=-0.62cm]frame.north west) {Question 62 D};
\node[anchor=east,text=white,font=\sffamily\itshape\large] at ([xshift=-5mm,yshift=-0.62cm]frame.north east) {Géométrie};}]
\vspace*{1.75cm}
\begin{minipage}[t][6.9cm][c]{\linewidth}\Large\centering
\begin{minipage}{0.95\linewidth}\raggedright\linespread{1.25}\selectfont Dans un repère, $A(-3\,;\,2)$, $B(-1\,;\,-2)$ et $C(4\,;\,1)$. Quelles sont les coordonnées de l'image de $C$ par la translation qui transforme $A$ en $B$~?\end{minipage}
\end{minipage}

\vfill
\begin{tcolorbox}[enhanced,nobeforeafter,width=7.1cm,height=2.1cm,valign=center,colback=cgeo!6,colframe=cgeo,arc=2.5mm,boxrule=0.9pt,left=3mm]{\Large\bfseries\color{cgeo}a)}\quad{\Large $(2\,;\,5)$}\end{tcolorbox}\hfill\begin{tcolorbox}[enhanced,nobeforeafter,width=7.1cm,height=2.1cm,valign=center,colback=cgeo!6,colframe=cgeo,arc=2.5mm,boxrule=0.9pt,left=3mm]{\Large\bfseries\color{cgeo}b)}\quad{\Large $(6\,;\,-3)$}\end{tcolorbox}\par\vspace{2.5mm}
\begin{tcolorbox}[enhanced,nobeforeafter,width=7.1cm,height=2.1cm,valign=center,colback=cgeo!6,colframe=cgeo,arc=2.5mm,boxrule=0.9pt,left=3mm]{\Large\bfseries\color{cgeo}c)}\quad{\Large $(3\,;\,-1)$}\end{tcolorbox}\hfill\begin{tcolorbox}[enhanced,nobeforeafter,width=7.1cm,height=2.1cm,valign=center,colback=cgeo!6,colframe=cgeo,arc=2.5mm,boxrule=0.9pt,left=3mm]{\Large\bfseries\color{cgeo}d)}\quad{\Large $(0\,;\,3)$}\end{tcolorbox}
\end{tcolorbox}
\newpage
\begin{tcolorbox}[enhanced,width=16cm,height=14cm,colback=white,colframe=cgeo,boxrule=2pt,arc=0mm,
left=6mm,right=6mm,top=0mm,bottom=4mm,sharp corners,
overlay={\fill[cgeo] (frame.north west) rectangle ([yshift=-1.25cm]frame.north east);
\node[anchor=west,text=white,font=\sffamily\bfseries\LARGE] at ([xshift=5mm,yshift=-0.62cm]frame.north west) {Question 63 D};
\node[anchor=east,text=white,font=\sffamily\itshape\large] at ([xshift=-5mm,yshift=-0.62cm]frame.north east) {Géométrie};}]
\vspace*{1.75cm}
\begin{minipage}[t][6.9cm][c]{\linewidth}\LARGE\centering
\begin{minipage}{0.95\linewidth}\raggedright\linespread{1.25}\selectfont Une figure est réduite à l'échelle $0{,}5$ (les longueurs sont multipliées par $0{,}5$). Par combien son aire est-elle multipliée~?\end{minipage}
\end{minipage}

\vfill
\begin{tcolorbox}[enhanced,nobeforeafter,width=7.1cm,height=2.1cm,valign=center,colback=cgeo!6,colframe=cgeo,arc=2.5mm,boxrule=0.9pt,left=3mm]{\Large\bfseries\color{cgeo}a)}\quad{\Large $0{,}125$}\end{tcolorbox}\hfill\begin{tcolorbox}[enhanced,nobeforeafter,width=7.1cm,height=2.1cm,valign=center,colback=cgeo!6,colframe=cgeo,arc=2.5mm,boxrule=0.9pt,left=3mm]{\Large\bfseries\color{cgeo}b)}\quad{\Large $0{,}5$}\end{tcolorbox}\par\vspace{2.5mm}
\begin{tcolorbox}[enhanced,nobeforeafter,width=7.1cm,height=2.1cm,valign=center,colback=cgeo!6,colframe=cgeo,arc=2.5mm,boxrule=0.9pt,left=3mm]{\Large\bfseries\color{cgeo}c)}\quad{\Large $0{,}25$}\end{tcolorbox}\hfill\begin{tcolorbox}[enhanced,nobeforeafter,width=7.1cm,height=2.1cm,valign=center,colback=cgeo!6,colframe=cgeo,arc=2.5mm,boxrule=0.9pt,left=3mm]{\Large\bfseries\color{cgeo}d)}\quad{\Large $1$}\end{tcolorbox}
\end{tcolorbox}
\newpage
\begin{tcolorbox}[enhanced,width=16cm,height=14cm,colback=white,colframe=cgeo,boxrule=2pt,arc=0mm,
left=6mm,right=6mm,top=0mm,bottom=4mm,sharp corners,
overlay={\fill[cgeo] (frame.north west) rectangle ([yshift=-1.25cm]frame.north east);
\node[anchor=west,text=white,font=\sffamily\bfseries\LARGE] at ([xshift=5mm,yshift=-0.62cm]frame.north west) {Question 64 D};
\node[anchor=east,text=white,font=\sffamily\itshape\large] at ([xshift=-5mm,yshift=-0.62cm]frame.north east) {Géométrie};}]
\vspace*{1.75cm}
\begin{minipage}[t][6.9cm][c]{\linewidth}\LARGE\centering
\begin{minipage}{0.95\linewidth}\raggedright\linespread{1.25}\selectfont Une figure est agrandie à l'échelle $2{,}5$ (les longueurs sont multipliées par $2{,}5$). Par combien son aire est-elle multipliée~?\end{minipage}
\end{minipage}

\vfill
\begin{tcolorbox}[enhanced,nobeforeafter,width=7.1cm,height=2.1cm,valign=center,colback=cgeo!6,colframe=cgeo,arc=2.5mm,boxrule=0.9pt,left=3mm]{\Large\bfseries\color{cgeo}a)}\quad{\Large $15{,}625$}\end{tcolorbox}\hfill\begin{tcolorbox}[enhanced,nobeforeafter,width=7.1cm,height=2.1cm,valign=center,colback=cgeo!6,colframe=cgeo,arc=2.5mm,boxrule=0.9pt,left=3mm]{\Large\bfseries\color{cgeo}b)}\quad{\Large $5$}\end{tcolorbox}\par\vspace{2.5mm}
\begin{tcolorbox}[enhanced,nobeforeafter,width=7.1cm,height=2.1cm,valign=center,colback=cgeo!6,colframe=cgeo,arc=2.5mm,boxrule=0.9pt,left=3mm]{\Large\bfseries\color{cgeo}c)}\quad{\Large $6{,}25$}\end{tcolorbox}\hfill\begin{tcolorbox}[enhanced,nobeforeafter,width=7.1cm,height=2.1cm,valign=center,colback=cgeo!6,colframe=cgeo,arc=2.5mm,boxrule=0.9pt,left=3mm]{\Large\bfseries\color{cgeo}d)}\quad{\Large $2{,}5$}\end{tcolorbox}
\end{tcolorbox}
\newpage
\begin{tcolorbox}[enhanced,width=16cm,height=14cm,colback=white,colframe=cdon,boxrule=2pt,arc=0mm,
left=6mm,right=6mm,top=0mm,bottom=4mm,sharp corners,
overlay={\fill[cdon] (frame.north west) rectangle ([yshift=-1.25cm]frame.north east);
\node[anchor=west,text=white,font=\sffamily\bfseries\LARGE] at ([xshift=5mm,yshift=-0.62cm]frame.north west) {Question 65 D};
\node[anchor=east,text=white,font=\sffamily\itshape\large] at ([xshift=-5mm,yshift=-0.62cm]frame.north east) {Proportionnalité, statistiques, probabilités};}]
\vspace*{1.75cm}
\begin{minipage}[t][6.9cm][c]{\linewidth}\LARGE\centering
\begin{minipage}{0.95\linewidth}\raggedright\linespread{1.25}\selectfont Ce tableau est un tableau de proportionnalité~: $\begin{array}{|c|c|}\hline 12 & 8\\\hline 30 &~?\\\hline\end{array}$. Quel nombre remplace «~?~»~?\end{minipage}
\end{minipage}

\vfill
\begin{tcolorbox}[enhanced,nobeforeafter,width=7.1cm,height=2.1cm,valign=center,colback=cdon!6,colframe=cdon,arc=2.5mm,boxrule=0.9pt,left=3mm]{\Large\bfseries\color{cdon}a)}\quad{\Large $20$}\end{tcolorbox}\hfill\begin{tcolorbox}[enhanced,nobeforeafter,width=7.1cm,height=2.1cm,valign=center,colback=cdon!6,colframe=cdon,arc=2.5mm,boxrule=0.9pt,left=3mm]{\Large\bfseries\color{cdon}b)}\quad{\Large $\dfrac{16}{5}$}\end{tcolorbox}\par\vspace{2.5mm}
\begin{tcolorbox}[enhanced,nobeforeafter,width=7.1cm,height=2.1cm,valign=center,colback=cdon!6,colframe=cdon,arc=2.5mm,boxrule=0.9pt,left=3mm]{\Large\bfseries\color{cdon}c)}\quad{\Large $240$}\end{tcolorbox}\hfill\begin{tcolorbox}[enhanced,nobeforeafter,width=7.1cm,height=2.1cm,valign=center,colback=cdon!6,colframe=cdon,arc=2.5mm,boxrule=0.9pt,left=3mm]{\Large\bfseries\color{cdon}d)}\quad{\Large $26$}\end{tcolorbox}
\end{tcolorbox}
\newpage
\begin{tcolorbox}[enhanced,width=16cm,height=14cm,colback=white,colframe=cdon,boxrule=2pt,arc=0mm,
left=6mm,right=6mm,top=0mm,bottom=4mm,sharp corners,
overlay={\fill[cdon] (frame.north west) rectangle ([yshift=-1.25cm]frame.north east);
\node[anchor=west,text=white,font=\sffamily\bfseries\LARGE] at ([xshift=5mm,yshift=-0.62cm]frame.north west) {Question 66 D};
\node[anchor=east,text=white,font=\sffamily\itshape\large] at ([xshift=-5mm,yshift=-0.62cm]frame.north east) {Proportionnalité, statistiques, probabilités};}]
\vspace*{1.75cm}
\begin{minipage}[t][6.9cm][c]{\linewidth}\LARGE\centering
\begin{minipage}{0.95\linewidth}\raggedright\linespread{1.25}\selectfont Ce tableau est un tableau de proportionnalité~: $\begin{array}{|c|c|}\hline 7 & 4\\\hline 21 &~?\\\hline\end{array}$. Quel nombre remplace «~?~»~?\end{minipage}
\end{minipage}

\vfill
\begin{tcolorbox}[enhanced,nobeforeafter,width=7.1cm,height=2.1cm,valign=center,colback=cdon!6,colframe=cdon,arc=2.5mm,boxrule=0.9pt,left=3mm]{\Large\bfseries\color{cdon}a)}\quad{\Large $\dfrac{4}{3}$}\end{tcolorbox}\hfill\begin{tcolorbox}[enhanced,nobeforeafter,width=7.1cm,height=2.1cm,valign=center,colback=cdon!6,colframe=cdon,arc=2.5mm,boxrule=0.9pt,left=3mm]{\Large\bfseries\color{cdon}b)}\quad{\Large $12$}\end{tcolorbox}\par\vspace{2.5mm}
\begin{tcolorbox}[enhanced,nobeforeafter,width=7.1cm,height=2.1cm,valign=center,colback=cdon!6,colframe=cdon,arc=2.5mm,boxrule=0.9pt,left=3mm]{\Large\bfseries\color{cdon}c)}\quad{\Large $18$}\end{tcolorbox}\hfill\begin{tcolorbox}[enhanced,nobeforeafter,width=7.1cm,height=2.1cm,valign=center,colback=cdon!6,colframe=cdon,arc=2.5mm,boxrule=0.9pt,left=3mm]{\Large\bfseries\color{cdon}d)}\quad{\Large $84$}\end{tcolorbox}
\end{tcolorbox}
\newpage
\begin{tcolorbox}[enhanced,width=16cm,height=14cm,colback=white,colframe=cdon,boxrule=2pt,arc=0mm,
left=6mm,right=6mm,top=0mm,bottom=4mm,sharp corners,
overlay={\fill[cdon] (frame.north west) rectangle ([yshift=-1.25cm]frame.north east);
\node[anchor=west,text=white,font=\sffamily\bfseries\LARGE] at ([xshift=5mm,yshift=-0.62cm]frame.north west) {Question 67 D};
\node[anchor=east,text=white,font=\sffamily\itshape\large] at ([xshift=-5mm,yshift=-0.62cm]frame.north east) {Proportionnalité, statistiques, probabilités};}]
\vspace*{1.75cm}
\begin{minipage}[t][6.9cm][c]{\linewidth}\LARGE\centering
\begin{minipage}{0.95\linewidth}\raggedright\linespread{1.25}\selectfont Que vaut $5\,\%$ de $460$~?\end{minipage}
\end{minipage}

\vfill
\begin{tcolorbox}[enhanced,nobeforeafter,width=7.1cm,height=2.1cm,valign=center,colback=cdon!6,colframe=cdon,arc=2.5mm,boxrule=0.9pt,left=3mm]{\Large\bfseries\color{cdon}a)}\quad{\Large $28$}\end{tcolorbox}\hfill\begin{tcolorbox}[enhanced,nobeforeafter,width=7.1cm,height=2.1cm,valign=center,colback=cdon!6,colframe=cdon,arc=2.5mm,boxrule=0.9pt,left=3mm]{\Large\bfseries\color{cdon}b)}\quad{\Large $437$}\end{tcolorbox}\par\vspace{2.5mm}
\begin{tcolorbox}[enhanced,nobeforeafter,width=7.1cm,height=2.1cm,valign=center,colback=cdon!6,colframe=cdon,arc=2.5mm,boxrule=0.9pt,left=3mm]{\Large\bfseries\color{cdon}c)}\quad{\Large $23$}\end{tcolorbox}\hfill\begin{tcolorbox}[enhanced,nobeforeafter,width=7.1cm,height=2.1cm,valign=center,colback=cdon!6,colframe=cdon,arc=2.5mm,boxrule=0.9pt,left=3mm]{\Large\bfseries\color{cdon}d)}\quad{\Large $230$}\end{tcolorbox}
\end{tcolorbox}
\newpage
\begin{tcolorbox}[enhanced,width=16cm,height=14cm,colback=white,colframe=cdon,boxrule=2pt,arc=0mm,
left=6mm,right=6mm,top=0mm,bottom=4mm,sharp corners,
overlay={\fill[cdon] (frame.north west) rectangle ([yshift=-1.25cm]frame.north east);
\node[anchor=west,text=white,font=\sffamily\bfseries\LARGE] at ([xshift=5mm,yshift=-0.62cm]frame.north west) {Question 68 D};
\node[anchor=east,text=white,font=\sffamily\itshape\large] at ([xshift=-5mm,yshift=-0.62cm]frame.north east) {Proportionnalité, statistiques, probabilités};}]
\vspace*{1.75cm}
\begin{minipage}[t][6.9cm][c]{\linewidth}\LARGE\centering
\begin{minipage}{0.95\linewidth}\raggedright\linespread{1.25}\selectfont Que vaut $60\,\%$ de $35$~?\end{minipage}
\end{minipage}

\vfill
\begin{tcolorbox}[enhanced,nobeforeafter,width=7.1cm,height=2.1cm,valign=center,colback=cdon!6,colframe=cdon,arc=2.5mm,boxrule=0.9pt,left=3mm]{\Large\bfseries\color{cdon}a)}\quad{\Large $210$}\end{tcolorbox}\hfill\begin{tcolorbox}[enhanced,nobeforeafter,width=7.1cm,height=2.1cm,valign=center,colback=cdon!6,colframe=cdon,arc=2.5mm,boxrule=0.9pt,left=3mm]{\Large\bfseries\color{cdon}b)}\quad{\Large $21$}\end{tcolorbox}\par\vspace{2.5mm}
\begin{tcolorbox}[enhanced,nobeforeafter,width=7.1cm,height=2.1cm,valign=center,colback=cdon!6,colframe=cdon,arc=2.5mm,boxrule=0.9pt,left=3mm]{\Large\bfseries\color{cdon}c)}\quad{\Large $14$}\end{tcolorbox}\hfill\begin{tcolorbox}[enhanced,nobeforeafter,width=7.1cm,height=2.1cm,valign=center,colback=cdon!6,colframe=cdon,arc=2.5mm,boxrule=0.9pt,left=3mm]{\Large\bfseries\color{cdon}d)}\quad{\Large $81$}\end{tcolorbox}
\end{tcolorbox}
\newpage
\begin{tcolorbox}[enhanced,width=16cm,height=14cm,colback=white,colframe=cdon,boxrule=2pt,arc=0mm,
left=6mm,right=6mm,top=0mm,bottom=4mm,sharp corners,
overlay={\fill[cdon] (frame.north west) rectangle ([yshift=-1.25cm]frame.north east);
\node[anchor=west,text=white,font=\sffamily\bfseries\LARGE] at ([xshift=5mm,yshift=-0.62cm]frame.north west) {Question 69 D};
\node[anchor=east,text=white,font=\sffamily\itshape\large] at ([xshift=-5mm,yshift=-0.62cm]frame.north east) {Proportionnalité, statistiques, probabilités};}]
\vspace*{1.75cm}
\begin{minipage}[t][6.9cm][c]{\linewidth}\LARGE\centering
\begin{minipage}{0.95\linewidth}\raggedright\linespread{1.25}\selectfont Sur $200$ élèves du collège, $46$ viennent en bus. Quel pourcentage des élèves cela représente-t-il~?\end{minipage}
\end{minipage}

\vfill
\begin{tcolorbox}[enhanced,nobeforeafter,width=7.1cm,height=2.1cm,valign=center,colback=cdon!6,colframe=cdon,arc=2.5mm,boxrule=0.9pt,left=3mm]{\Large\bfseries\color{cdon}a)}\quad{\Large $2{,}3\,\%$}\end{tcolorbox}\hfill\begin{tcolorbox}[enhanced,nobeforeafter,width=7.1cm,height=2.1cm,valign=center,colback=cdon!6,colframe=cdon,arc=2.5mm,boxrule=0.9pt,left=3mm]{\Large\bfseries\color{cdon}b)}\quad{\Large $23\,\%$}\end{tcolorbox}\par\vspace{2.5mm}
\begin{tcolorbox}[enhanced,nobeforeafter,width=7.1cm,height=2.1cm,valign=center,colback=cdon!6,colframe=cdon,arc=2.5mm,boxrule=0.9pt,left=3mm]{\Large\bfseries\color{cdon}c)}\quad{\Large $77\,\%$}\end{tcolorbox}\hfill\begin{tcolorbox}[enhanced,nobeforeafter,width=7.1cm,height=2.1cm,valign=center,colback=cdon!6,colframe=cdon,arc=2.5mm,boxrule=0.9pt,left=3mm]{\Large\bfseries\color{cdon}d)}\quad{\Large $46\,\%$}\end{tcolorbox}
\end{tcolorbox}
\newpage
\begin{tcolorbox}[enhanced,width=16cm,height=14cm,colback=white,colframe=cdon,boxrule=2pt,arc=0mm,
left=6mm,right=6mm,top=0mm,bottom=4mm,sharp corners,
overlay={\fill[cdon] (frame.north west) rectangle ([yshift=-1.25cm]frame.north east);
\node[anchor=west,text=white,font=\sffamily\bfseries\LARGE] at ([xshift=5mm,yshift=-0.62cm]frame.north west) {Question 70 D};
\node[anchor=east,text=white,font=\sffamily\itshape\large] at ([xshift=-5mm,yshift=-0.62cm]frame.north east) {Proportionnalité, statistiques, probabilités};}]
\vspace*{1.75cm}
\begin{minipage}[t][6.9cm][c]{\linewidth}\LARGE\centering
\begin{minipage}{0.95\linewidth}\raggedright\linespread{1.25}\selectfont Dans une classe de $16$ élèves, $14$ sont demi-pensionnaires. Quel pourcentage des élèves cela représente-t-il~?\end{minipage}
\end{minipage}

\vfill
\begin{tcolorbox}[enhanced,nobeforeafter,width=7.1cm,height=2.1cm,valign=center,colback=cdon!6,colframe=cdon,arc=2.5mm,boxrule=0.9pt,left=3mm]{\Large\bfseries\color{cdon}a)}\quad{\Large $8{,}75\,\%$}\end{tcolorbox}\hfill\begin{tcolorbox}[enhanced,nobeforeafter,width=7.1cm,height=2.1cm,valign=center,colback=cdon!6,colframe=cdon,arc=2.5mm,boxrule=0.9pt,left=3mm]{\Large\bfseries\color{cdon}b)}\quad{\Large $87{,}5\,\%$}\end{tcolorbox}\par\vspace{2.5mm}
\begin{tcolorbox}[enhanced,nobeforeafter,width=7.1cm,height=2.1cm,valign=center,colback=cdon!6,colframe=cdon,arc=2.5mm,boxrule=0.9pt,left=3mm]{\Large\bfseries\color{cdon}c)}\quad{\Large $14\,\%$}\end{tcolorbox}\hfill\begin{tcolorbox}[enhanced,nobeforeafter,width=7.1cm,height=2.1cm,valign=center,colback=cdon!6,colframe=cdon,arc=2.5mm,boxrule=0.9pt,left=3mm]{\Large\bfseries\color{cdon}d)}\quad{\Large $12{,}5\,\%$}\end{tcolorbox}
\end{tcolorbox}
\newpage
\begin{tcolorbox}[enhanced,width=16cm,height=14cm,colback=white,colframe=cdon,boxrule=2pt,arc=0mm,
left=6mm,right=6mm,top=0mm,bottom=4mm,sharp corners,
overlay={\fill[cdon] (frame.north west) rectangle ([yshift=-1.25cm]frame.north east);
\node[anchor=west,text=white,font=\sffamily\bfseries\LARGE] at ([xshift=5mm,yshift=-0.62cm]frame.north west) {Question 71 D};
\node[anchor=east,text=white,font=\sffamily\itshape\large] at ([xshift=-5mm,yshift=-0.62cm]frame.north east) {Proportionnalité, statistiques, probabilités};}]
\vspace*{1.75cm}
\begin{minipage}[t][6.9cm][c]{\linewidth}\LARGE\centering
\begin{minipage}{0.95\linewidth}\raggedright\linespread{1.25}\selectfont Quelle est la moyenne de la série $3\ ;\ 7\ ;\ 11\ ;\ 15$~?\end{minipage}
\end{minipage}

\vfill
\begin{tcolorbox}[enhanced,nobeforeafter,width=7.1cm,height=2.1cm,valign=center,colback=cdon!6,colframe=cdon,arc=2.5mm,boxrule=0.9pt,left=3mm]{\Large\bfseries\color{cdon}a)}\quad{\Large $36$}\end{tcolorbox}\hfill\begin{tcolorbox}[enhanced,nobeforeafter,width=7.1cm,height=2.1cm,valign=center,colback=cdon!6,colframe=cdon,arc=2.5mm,boxrule=0.9pt,left=3mm]{\Large\bfseries\color{cdon}b)}\quad{\Large $8$}\end{tcolorbox}\par\vspace{2.5mm}
\begin{tcolorbox}[enhanced,nobeforeafter,width=7.1cm,height=2.1cm,valign=center,colback=cdon!6,colframe=cdon,arc=2.5mm,boxrule=0.9pt,left=3mm]{\Large\bfseries\color{cdon}c)}\quad{\Large $9$}\end{tcolorbox}\hfill\begin{tcolorbox}[enhanced,nobeforeafter,width=7.1cm,height=2.1cm,valign=center,colback=cdon!6,colframe=cdon,arc=2.5mm,boxrule=0.9pt,left=3mm]{\Large\bfseries\color{cdon}d)}\quad{\Large $10$}\end{tcolorbox}
\end{tcolorbox}
\newpage
\begin{tcolorbox}[enhanced,width=16cm,height=14cm,colback=white,colframe=cdon,boxrule=2pt,arc=0mm,
left=6mm,right=6mm,top=0mm,bottom=4mm,sharp corners,
overlay={\fill[cdon] (frame.north west) rectangle ([yshift=-1.25cm]frame.north east);
\node[anchor=west,text=white,font=\sffamily\bfseries\LARGE] at ([xshift=5mm,yshift=-0.62cm]frame.north west) {Question 72 D};
\node[anchor=east,text=white,font=\sffamily\itshape\large] at ([xshift=-5mm,yshift=-0.62cm]frame.north east) {Proportionnalité, statistiques, probabilités};}]
\vspace*{1.75cm}
\begin{minipage}[t][6.9cm][c]{\linewidth}\LARGE\centering
\begin{minipage}{0.95\linewidth}\raggedright\linespread{1.25}\selectfont Quelle est la moyenne de la série $16\ ;\ 19\ ;\ 10\ ;\ 13$~?\end{minipage}
\end{minipage}

\vfill
\begin{tcolorbox}[enhanced,nobeforeafter,width=7.1cm,height=2.1cm,valign=center,colback=cdon!6,colframe=cdon,arc=2.5mm,boxrule=0.9pt,left=3mm]{\Large\bfseries\color{cdon}a)}\quad{\Large $58$}\end{tcolorbox}\hfill\begin{tcolorbox}[enhanced,nobeforeafter,width=7.1cm,height=2.1cm,valign=center,colback=cdon!6,colframe=cdon,arc=2.5mm,boxrule=0.9pt,left=3mm]{\Large\bfseries\color{cdon}b)}\quad{\Large $14{,}5$}\end{tcolorbox}\par\vspace{2.5mm}
\begin{tcolorbox}[enhanced,nobeforeafter,width=7.1cm,height=2.1cm,valign=center,colback=cdon!6,colframe=cdon,arc=2.5mm,boxrule=0.9pt,left=3mm]{\Large\bfseries\color{cdon}c)}\quad{\Large $15{,}5$}\end{tcolorbox}\hfill\begin{tcolorbox}[enhanced,nobeforeafter,width=7.1cm,height=2.1cm,valign=center,colback=cdon!6,colframe=cdon,arc=2.5mm,boxrule=0.9pt,left=3mm]{\Large\bfseries\color{cdon}d)}\quad{\Large $13{,}5$}\end{tcolorbox}
\end{tcolorbox}
\newpage
\begin{tcolorbox}[enhanced,width=16cm,height=14cm,colback=white,colframe=cdon,boxrule=2pt,arc=0mm,
left=6mm,right=6mm,top=0mm,bottom=4mm,sharp corners,
overlay={\fill[cdon] (frame.north west) rectangle ([yshift=-1.25cm]frame.north east);
\node[anchor=west,text=white,font=\sffamily\bfseries\LARGE] at ([xshift=5mm,yshift=-0.62cm]frame.north west) {Question 73 D};
\node[anchor=east,text=white,font=\sffamily\itshape\large] at ([xshift=-5mm,yshift=-0.62cm]frame.north east) {Proportionnalité, statistiques, probabilités};}]
\vspace*{1.75cm}
\begin{minipage}[t][6.9cm][c]{\linewidth}\LARGE\centering
\begin{minipage}{0.95\linewidth}\raggedright\linespread{1.25}\selectfont Quelle est la médiane de la série $3\ ;\ 3\ ;\ 6\ ;\ 10\ ;\ 1$~?\end{minipage}
\end{minipage}

\vfill
\begin{tcolorbox}[enhanced,nobeforeafter,width=7.1cm,height=2.1cm,valign=center,colback=cdon!6,colframe=cdon,arc=2.5mm,boxrule=0.9pt,left=3mm]{\Large\bfseries\color{cdon}a)}\quad{\Large $4$}\end{tcolorbox}\hfill\begin{tcolorbox}[enhanced,nobeforeafter,width=7.1cm,height=2.1cm,valign=center,colback=cdon!6,colframe=cdon,arc=2.5mm,boxrule=0.9pt,left=3mm]{\Large\bfseries\color{cdon}b)}\quad{\Large $4{,}6$}\end{tcolorbox}\par\vspace{2.5mm}
\begin{tcolorbox}[enhanced,nobeforeafter,width=7.1cm,height=2.1cm,valign=center,colback=cdon!6,colframe=cdon,arc=2.5mm,boxrule=0.9pt,left=3mm]{\Large\bfseries\color{cdon}c)}\quad{\Large $3$}\end{tcolorbox}\hfill\begin{tcolorbox}[enhanced,nobeforeafter,width=7.1cm,height=2.1cm,valign=center,colback=cdon!6,colframe=cdon,arc=2.5mm,boxrule=0.9pt,left=3mm]{\Large\bfseries\color{cdon}d)}\quad{\Large $6$}\end{tcolorbox}
\end{tcolorbox}
\newpage
\begin{tcolorbox}[enhanced,width=16cm,height=14cm,colback=white,colframe=cdon,boxrule=2pt,arc=0mm,
left=6mm,right=6mm,top=0mm,bottom=4mm,sharp corners,
overlay={\fill[cdon] (frame.north west) rectangle ([yshift=-1.25cm]frame.north east);
\node[anchor=west,text=white,font=\sffamily\bfseries\LARGE] at ([xshift=5mm,yshift=-0.62cm]frame.north west) {Question 74 D};
\node[anchor=east,text=white,font=\sffamily\itshape\large] at ([xshift=-5mm,yshift=-0.62cm]frame.north east) {Proportionnalité, statistiques, probabilités};}]
\vspace*{1.75cm}
\begin{minipage}[t][6.9cm][c]{\linewidth}\LARGE\centering
\begin{minipage}{0.95\linewidth}\raggedright\linespread{1.25}\selectfont Quelle est la médiane de la série $9\ ;\ 5\ ;\ 16\ ;\ 11\ ;\ 14\ ;\ 2$~?\end{minipage}
\end{minipage}

\vfill
\begin{tcolorbox}[enhanced,nobeforeafter,width=7.1cm,height=2.1cm,valign=center,colback=cdon!6,colframe=cdon,arc=2.5mm,boxrule=0.9pt,left=3mm]{\Large\bfseries\color{cdon}a)}\quad{\Large $10$}\end{tcolorbox}\hfill\begin{tcolorbox}[enhanced,nobeforeafter,width=7.1cm,height=2.1cm,valign=center,colback=cdon!6,colframe=cdon,arc=2.5mm,boxrule=0.9pt,left=3mm]{\Large\bfseries\color{cdon}b)}\quad{\Large $9{,}5$}\end{tcolorbox}\par\vspace{2.5mm}
\begin{tcolorbox}[enhanced,nobeforeafter,width=7.1cm,height=2.1cm,valign=center,colback=cdon!6,colframe=cdon,arc=2.5mm,boxrule=0.9pt,left=3mm]{\Large\bfseries\color{cdon}c)}\quad{\Large $3{,}5$}\end{tcolorbox}\hfill\begin{tcolorbox}[enhanced,nobeforeafter,width=7.1cm,height=2.1cm,valign=center,colback=cdon!6,colframe=cdon,arc=2.5mm,boxrule=0.9pt,left=3mm]{\Large\bfseries\color{cdon}d)}\quad{\Large $13{,}5$}\end{tcolorbox}
\end{tcolorbox}
\newpage
\begin{tcolorbox}[enhanced,width=16cm,height=14cm,colback=white,colframe=cdon,boxrule=2pt,arc=0mm,
left=6mm,right=6mm,top=0mm,bottom=4mm,sharp corners,
overlay={\fill[cdon] (frame.north west) rectangle ([yshift=-1.25cm]frame.north east);
\node[anchor=west,text=white,font=\sffamily\bfseries\LARGE] at ([xshift=5mm,yshift=-0.62cm]frame.north west) {Question 75 D};
\node[anchor=east,text=white,font=\sffamily\itshape\large] at ([xshift=-5mm,yshift=-0.62cm]frame.north east) {Proportionnalité, statistiques, probabilités};}]
\vspace*{1.75cm}
\begin{minipage}[t][6.9cm][c]{\linewidth}\LARGE\centering
\begin{minipage}{0.95\linewidth}\raggedright\linespread{1.25}\selectfont Quelle est l'étendue de la série $30\ ;\ 18\ ;\ 27\ ;\ 21$~?\end{minipage}
\end{minipage}

\vfill
\begin{tcolorbox}[enhanced,nobeforeafter,width=7.1cm,height=2.1cm,valign=center,colback=cdon!6,colframe=cdon,arc=2.5mm,boxrule=0.9pt,left=3mm]{\Large\bfseries\color{cdon}a)}\quad{\Large $12$}\end{tcolorbox}\hfill\begin{tcolorbox}[enhanced,nobeforeafter,width=7.1cm,height=2.1cm,valign=center,colback=cdon!6,colframe=cdon,arc=2.5mm,boxrule=0.9pt,left=3mm]{\Large\bfseries\color{cdon}b)}\quad{\Large $48$}\end{tcolorbox}\par\vspace{2.5mm}
\begin{tcolorbox}[enhanced,nobeforeafter,width=7.1cm,height=2.1cm,valign=center,colback=cdon!6,colframe=cdon,arc=2.5mm,boxrule=0.9pt,left=3mm]{\Large\bfseries\color{cdon}c)}\quad{\Large $24$}\end{tcolorbox}\hfill\begin{tcolorbox}[enhanced,nobeforeafter,width=7.1cm,height=2.1cm,valign=center,colback=cdon!6,colframe=cdon,arc=2.5mm,boxrule=0.9pt,left=3mm]{\Large\bfseries\color{cdon}d)}\quad{\Large $30$}\end{tcolorbox}
\end{tcolorbox}
\newpage
\begin{tcolorbox}[enhanced,width=16cm,height=14cm,colback=white,colframe=cdon,boxrule=2pt,arc=0mm,
left=6mm,right=6mm,top=0mm,bottom=4mm,sharp corners,
overlay={\fill[cdon] (frame.north west) rectangle ([yshift=-1.25cm]frame.north east);
\node[anchor=west,text=white,font=\sffamily\bfseries\LARGE] at ([xshift=5mm,yshift=-0.62cm]frame.north west) {Question 76 D};
\node[anchor=east,text=white,font=\sffamily\itshape\large] at ([xshift=-5mm,yshift=-0.62cm]frame.north east) {Proportionnalité, statistiques, probabilités};}]
\vspace*{1.75cm}
\begin{minipage}[t][6.9cm][c]{\linewidth}\LARGE\centering
\begin{minipage}{0.95\linewidth}\raggedright\linespread{1.25}\selectfont Quelle est l'étendue de la série $9\ ;\ 17\ ;\ 6\ ;\ 22\ ;\ 10$~?\end{minipage}
\end{minipage}

\vfill
\begin{tcolorbox}[enhanced,nobeforeafter,width=7.1cm,height=2.1cm,valign=center,colback=cdon!6,colframe=cdon,arc=2.5mm,boxrule=0.9pt,left=3mm]{\Large\bfseries\color{cdon}a)}\quad{\Large $28$}\end{tcolorbox}\hfill\begin{tcolorbox}[enhanced,nobeforeafter,width=7.1cm,height=2.1cm,valign=center,colback=cdon!6,colframe=cdon,arc=2.5mm,boxrule=0.9pt,left=3mm]{\Large\bfseries\color{cdon}b)}\quad{\Large $12{,}8$}\end{tcolorbox}\par\vspace{2.5mm}
\begin{tcolorbox}[enhanced,nobeforeafter,width=7.1cm,height=2.1cm,valign=center,colback=cdon!6,colframe=cdon,arc=2.5mm,boxrule=0.9pt,left=3mm]{\Large\bfseries\color{cdon}c)}\quad{\Large $22$}\end{tcolorbox}\hfill\begin{tcolorbox}[enhanced,nobeforeafter,width=7.1cm,height=2.1cm,valign=center,colback=cdon!6,colframe=cdon,arc=2.5mm,boxrule=0.9pt,left=3mm]{\Large\bfseries\color{cdon}d)}\quad{\Large $16$}\end{tcolorbox}
\end{tcolorbox}
\newpage
\begin{tcolorbox}[enhanced,width=16cm,height=14cm,colback=white,colframe=cdon,boxrule=2pt,arc=0mm,
left=6mm,right=6mm,top=0mm,bottom=4mm,sharp corners,
overlay={\fill[cdon] (frame.north west) rectangle ([yshift=-1.25cm]frame.north east);
\node[anchor=west,text=white,font=\sffamily\bfseries\LARGE] at ([xshift=5mm,yshift=-0.62cm]frame.north west) {Question 77 D};
\node[anchor=east,text=white,font=\sffamily\itshape\large] at ([xshift=-5mm,yshift=-0.62cm]frame.north east) {Proportionnalité, statistiques, probabilités};}]
\vspace*{1.75cm}
\begin{minipage}[t][6.9cm][c]{\linewidth}\Large\centering
\begin{minipage}{0.95\linewidth}\raggedright\linespread{1.25}\selectfont Un sac contient $1$ boule rouge, $4$ boules bleues et $7$ boules vertes. On tire une boule au hasard. Quelle est la probabilité de tirer une boule verte~?\end{minipage}
\end{minipage}

\vfill
\begin{tcolorbox}[enhanced,nobeforeafter,width=7.1cm,height=2.1cm,valign=center,colback=cdon!6,colframe=cdon,arc=2.5mm,boxrule=0.9pt,left=3mm]{\Large\bfseries\color{cdon}a)}\quad{\Large $\dfrac{7}{5}$}\end{tcolorbox}\hfill\begin{tcolorbox}[enhanced,nobeforeafter,width=7.1cm,height=2.1cm,valign=center,colback=cdon!6,colframe=cdon,arc=2.5mm,boxrule=0.9pt,left=3mm]{\Large\bfseries\color{cdon}b)}\quad{\Large $\dfrac{5}{12}$}\end{tcolorbox}\par\vspace{2.5mm}
\begin{tcolorbox}[enhanced,nobeforeafter,width=7.1cm,height=2.1cm,valign=center,colback=cdon!6,colframe=cdon,arc=2.5mm,boxrule=0.9pt,left=3mm]{\Large\bfseries\color{cdon}c)}\quad{\Large $\dfrac{1}{3}$}\end{tcolorbox}\hfill\begin{tcolorbox}[enhanced,nobeforeafter,width=7.1cm,height=2.1cm,valign=center,colback=cdon!6,colframe=cdon,arc=2.5mm,boxrule=0.9pt,left=3mm]{\Large\bfseries\color{cdon}d)}\quad{\Large $\dfrac{7}{12}$}\end{tcolorbox}
\end{tcolorbox}
\newpage
\begin{tcolorbox}[enhanced,width=16cm,height=14cm,colback=white,colframe=cdon,boxrule=2pt,arc=0mm,
left=6mm,right=6mm,top=0mm,bottom=4mm,sharp corners,
overlay={\fill[cdon] (frame.north west) rectangle ([yshift=-1.25cm]frame.north east);
\node[anchor=west,text=white,font=\sffamily\bfseries\LARGE] at ([xshift=5mm,yshift=-0.62cm]frame.north west) {Question 78 D};
\node[anchor=east,text=white,font=\sffamily\itshape\large] at ([xshift=-5mm,yshift=-0.62cm]frame.north east) {Proportionnalité, statistiques, probabilités};}]
\vspace*{1.75cm}
\begin{minipage}[t][6.9cm][c]{\linewidth}\Large\centering
\begin{minipage}{0.95\linewidth}\raggedright\linespread{1.25}\selectfont Un sac contient $9$ boules rouges, $3$ boules bleues et $4$ boules vertes. On tire une boule au hasard. Quelle est la probabilité de tirer une boule bleue~?\end{minipage}
\end{minipage}

\vfill
\begin{tcolorbox}[enhanced,nobeforeafter,width=7.1cm,height=2.1cm,valign=center,colback=cdon!6,colframe=cdon,arc=2.5mm,boxrule=0.9pt,left=3mm]{\Large\bfseries\color{cdon}a)}\quad{\Large $\dfrac{13}{16}$}\end{tcolorbox}\hfill\begin{tcolorbox}[enhanced,nobeforeafter,width=7.1cm,height=2.1cm,valign=center,colback=cdon!6,colframe=cdon,arc=2.5mm,boxrule=0.9pt,left=3mm]{\Large\bfseries\color{cdon}b)}\quad{\Large $\dfrac{3}{13}$}\end{tcolorbox}\par\vspace{2.5mm}
\begin{tcolorbox}[enhanced,nobeforeafter,width=7.1cm,height=2.1cm,valign=center,colback=cdon!6,colframe=cdon,arc=2.5mm,boxrule=0.9pt,left=3mm]{\Large\bfseries\color{cdon}c)}\quad{\Large $\dfrac{3}{16}$}\end{tcolorbox}\hfill\begin{tcolorbox}[enhanced,nobeforeafter,width=7.1cm,height=2.1cm,valign=center,colback=cdon!6,colframe=cdon,arc=2.5mm,boxrule=0.9pt,left=3mm]{\Large\bfseries\color{cdon}d)}\quad{\Large $\dfrac{1}{3}$}\end{tcolorbox}
\end{tcolorbox}
\newpage
\begin{tcolorbox}[enhanced,width=16cm,height=14cm,colback=white,colframe=cdon,boxrule=2pt,arc=0mm,
left=6mm,right=6mm,top=0mm,bottom=4mm,sharp corners,
overlay={\fill[cdon] (frame.north west) rectangle ([yshift=-1.25cm]frame.north east);
\node[anchor=west,text=white,font=\sffamily\bfseries\LARGE] at ([xshift=5mm,yshift=-0.62cm]frame.north west) {Question 79 D};
\node[anchor=east,text=white,font=\sffamily\itshape\large] at ([xshift=-5mm,yshift=-0.62cm]frame.north east) {Proportionnalité, statistiques, probabilités};}]
\vspace*{1.75cm}
\begin{minipage}[t][6.9cm][c]{\linewidth}\LARGE\centering
\begin{minipage}{0.95\linewidth}\raggedright\linespread{1.25}\selectfont La probabilité qu'un événement $A$ se réalise est $\dfrac{1}{9}$. Quelle est la probabilité de l'événement contraire de $A$~?\end{minipage}
\end{minipage}

\vfill
\begin{tcolorbox}[enhanced,nobeforeafter,width=7.1cm,height=2.1cm,valign=center,colback=cdon!6,colframe=cdon,arc=2.5mm,boxrule=0.9pt,left=3mm]{\Large\bfseries\color{cdon}a)}\quad{\Large $\dfrac{7}{18}$}\end{tcolorbox}\hfill\begin{tcolorbox}[enhanced,nobeforeafter,width=7.1cm,height=2.1cm,valign=center,colback=cdon!6,colframe=cdon,arc=2.5mm,boxrule=0.9pt,left=3mm]{\Large\bfseries\color{cdon}b)}\quad{\Large $\dfrac{1}{9}$}\end{tcolorbox}\par\vspace{2.5mm}
\begin{tcolorbox}[enhanced,nobeforeafter,width=7.1cm,height=2.1cm,valign=center,colback=cdon!6,colframe=cdon,arc=2.5mm,boxrule=0.9pt,left=3mm]{\Large\bfseries\color{cdon}c)}\quad{\Large $\dfrac{17}{18}$}\end{tcolorbox}\hfill\begin{tcolorbox}[enhanced,nobeforeafter,width=7.1cm,height=2.1cm,valign=center,colback=cdon!6,colframe=cdon,arc=2.5mm,boxrule=0.9pt,left=3mm]{\Large\bfseries\color{cdon}d)}\quad{\Large $\dfrac{8}{9}$}\end{tcolorbox}
\end{tcolorbox}
\newpage
\begin{tcolorbox}[enhanced,width=16cm,height=14cm,colback=white,colframe=cdon,boxrule=2pt,arc=0mm,
left=6mm,right=6mm,top=0mm,bottom=4mm,sharp corners,
overlay={\fill[cdon] (frame.north west) rectangle ([yshift=-1.25cm]frame.north east);
\node[anchor=west,text=white,font=\sffamily\bfseries\LARGE] at ([xshift=5mm,yshift=-0.62cm]frame.north west) {Question 80 D};
\node[anchor=east,text=white,font=\sffamily\itshape\large] at ([xshift=-5mm,yshift=-0.62cm]frame.north east) {Proportionnalité, statistiques, probabilités};}]
\vspace*{1.75cm}
\begin{minipage}[t][6.9cm][c]{\linewidth}\LARGE\centering
\begin{minipage}{0.95\linewidth}\raggedright\linespread{1.25}\selectfont La probabilité qu'un événement $A$ se réalise est $0{,}875$. Quelle est la probabilité de l'événement contraire de $A$~?\end{minipage}
\end{minipage}

\vfill
\begin{tcolorbox}[enhanced,nobeforeafter,width=7.1cm,height=2.1cm,valign=center,colback=cdon!6,colframe=cdon,arc=2.5mm,boxrule=0.9pt,left=3mm]{\Large\bfseries\color{cdon}a)}\quad{\Large $0{,}875$}\end{tcolorbox}\hfill\begin{tcolorbox}[enhanced,nobeforeafter,width=7.1cm,height=2.1cm,valign=center,colback=cdon!6,colframe=cdon,arc=2.5mm,boxrule=0.9pt,left=3mm]{\Large\bfseries\color{cdon}b)}\quad{\Large $0{,}5625$}\end{tcolorbox}\par\vspace{2.5mm}
\begin{tcolorbox}[enhanced,nobeforeafter,width=7.1cm,height=2.1cm,valign=center,colback=cdon!6,colframe=cdon,arc=2.5mm,boxrule=0.9pt,left=3mm]{\Large\bfseries\color{cdon}c)}\quad{\Large $0{,}125$}\end{tcolorbox}\hfill\begin{tcolorbox}[enhanced,nobeforeafter,width=7.1cm,height=2.1cm,valign=center,colback=cdon!6,colframe=cdon,arc=2.5mm,boxrule=0.9pt,left=3mm]{\Large\bfseries\color{cdon}d)}\quad{\Large $0{,}375$}\end{tcolorbox}
\end{tcolorbox}
\newpage
\begin{tcolorbox}[enhanced,width=16cm,height=14cm,colback=white,colframe=cgra,boxrule=2pt,arc=0mm,
left=6mm,right=6mm,top=0mm,bottom=4mm,sharp corners,
overlay={\fill[cgra] (frame.north west) rectangle ([yshift=-1.25cm]frame.north east);
\node[anchor=west,text=white,font=\sffamily\bfseries\LARGE] at ([xshift=5mm,yshift=-0.62cm]frame.north west) {Question 81 D};
\node[anchor=east,text=white,font=\sffamily\itshape\large] at ([xshift=-5mm,yshift=-0.62cm]frame.north east) {Grandeurs et mesures};}]
\vspace*{1.75cm}
\begin{minipage}[t][6.9cm][c]{\linewidth}\LARGE\centering
\begin{minipage}{0.95\linewidth}\raggedright\linespread{1.25}\selectfont Un véhicule parcourt $275$ km en 2 h 30 min. Quelle est sa vitesse moyenne~?\end{minipage}
\end{minipage}

\vfill
\begin{tcolorbox}[enhanced,nobeforeafter,width=7.1cm,height=2.1cm,valign=center,colback=cgra!6,colframe=cgra,arc=2.5mm,boxrule=0.9pt,left=3mm]{\Large\bfseries\color{cgra}a)}\quad{\Large $120~\text{km/h}$}\end{tcolorbox}\hfill\begin{tcolorbox}[enhanced,nobeforeafter,width=7.1cm,height=2.1cm,valign=center,colback=cgra!6,colframe=cgra,arc=2.5mm,boxrule=0.9pt,left=3mm]{\Large\bfseries\color{cgra}b)}\quad{\Large $220~\text{km/h}$}\end{tcolorbox}\par\vspace{2.5mm}
\begin{tcolorbox}[enhanced,nobeforeafter,width=7.1cm,height=2.1cm,valign=center,colback=cgra!6,colframe=cgra,arc=2.5mm,boxrule=0.9pt,left=3mm]{\Large\bfseries\color{cgra}c)}\quad{\Large $687{,}5~\text{km/h}$}\end{tcolorbox}\hfill\begin{tcolorbox}[enhanced,nobeforeafter,width=7.1cm,height=2.1cm,valign=center,colback=cgra!6,colframe=cgra,arc=2.5mm,boxrule=0.9pt,left=3mm]{\Large\bfseries\color{cgra}d)}\quad{\Large $110~\text{km/h}$}\end{tcolorbox}
\end{tcolorbox}
\newpage
\begin{tcolorbox}[enhanced,width=16cm,height=14cm,colback=white,colframe=cgra,boxrule=2pt,arc=0mm,
left=6mm,right=6mm,top=0mm,bottom=4mm,sharp corners,
overlay={\fill[cgra] (frame.north west) rectangle ([yshift=-1.25cm]frame.north east);
\node[anchor=west,text=white,font=\sffamily\bfseries\LARGE] at ([xshift=5mm,yshift=-0.62cm]frame.north west) {Question 82 D};
\node[anchor=east,text=white,font=\sffamily\itshape\large] at ([xshift=-5mm,yshift=-0.62cm]frame.north east) {Grandeurs et mesures};}]
\vspace*{1.75cm}
\begin{minipage}[t][6.9cm][c]{\linewidth}\LARGE\centering
\begin{minipage}{0.95\linewidth}\raggedright\linespread{1.25}\selectfont Un véhicule parcourt $52$ km en 48 min. Quelle est sa vitesse moyenne~?\end{minipage}
\end{minipage}

\vfill
\begin{tcolorbox}[enhanced,nobeforeafter,width=7.1cm,height=2.1cm,valign=center,colback=cgra!6,colframe=cgra,arc=2.5mm,boxrule=0.9pt,left=3mm]{\Large\bfseries\color{cgra}a)}\quad{\Large $75~\text{km/h}$}\end{tcolorbox}\hfill\begin{tcolorbox}[enhanced,nobeforeafter,width=7.1cm,height=2.1cm,valign=center,colback=cgra!6,colframe=cgra,arc=2.5mm,boxrule=0.9pt,left=3mm]{\Large\bfseries\color{cgra}b)}\quad{\Large $41{,}6~\text{km/h}$}\end{tcolorbox}\par\vspace{2.5mm}
\begin{tcolorbox}[enhanced,nobeforeafter,width=7.1cm,height=2.1cm,valign=center,colback=cgra!6,colframe=cgra,arc=2.5mm,boxrule=0.9pt,left=3mm]{\Large\bfseries\color{cgra}c)}\quad{\Large $130~\text{km/h}$}\end{tcolorbox}\hfill\begin{tcolorbox}[enhanced,nobeforeafter,width=7.1cm,height=2.1cm,valign=center,colback=cgra!6,colframe=cgra,arc=2.5mm,boxrule=0.9pt,left=3mm]{\Large\bfseries\color{cgra}d)}\quad{\Large $65~\text{km/h}$}\end{tcolorbox}
\end{tcolorbox}
\newpage
\begin{tcolorbox}[enhanced,width=16cm,height=14cm,colback=white,colframe=cgra,boxrule=2pt,arc=0mm,
left=6mm,right=6mm,top=0mm,bottom=4mm,sharp corners,
overlay={\fill[cgra] (frame.north west) rectangle ([yshift=-1.25cm]frame.north east);
\node[anchor=west,text=white,font=\sffamily\bfseries\LARGE] at ([xshift=5mm,yshift=-0.62cm]frame.north west) {Question 83 D};
\node[anchor=east,text=white,font=\sffamily\itshape\large] at ([xshift=-5mm,yshift=-0.62cm]frame.north east) {Grandeurs et mesures};}]
\vspace*{1.75cm}
\begin{minipage}[t][6.9cm][c]{\linewidth}\LARGE\centering
\begin{minipage}{0.95\linewidth}\raggedright\linespread{1.25}\selectfont Un cycliste roule à la vitesse constante de $36$ km/h pendant $25$ min. Quelle distance parcourt-il~?\end{minipage}
\end{minipage}

\vfill
\begin{tcolorbox}[enhanced,nobeforeafter,width=7.1cm,height=2.1cm,valign=center,colback=cgra!6,colframe=cgra,arc=2.5mm,boxrule=0.9pt,left=3mm]{\Large\bfseries\color{cgra}a)}\quad{\Large $900~\text{km}$}\end{tcolorbox}\hfill\begin{tcolorbox}[enhanced,nobeforeafter,width=7.1cm,height=2.1cm,valign=center,colback=cgra!6,colframe=cgra,arc=2.5mm,boxrule=0.9pt,left=3mm]{\Large\bfseries\color{cgra}b)}\quad{\Large $9~\text{km}$}\end{tcolorbox}\par\vspace{2.5mm}
\begin{tcolorbox}[enhanced,nobeforeafter,width=7.1cm,height=2.1cm,valign=center,colback=cgra!6,colframe=cgra,arc=2.5mm,boxrule=0.9pt,left=3mm]{\Large\bfseries\color{cgra}c)}\quad{\Large $15~\text{km}$}\end{tcolorbox}\hfill\begin{tcolorbox}[enhanced,nobeforeafter,width=7.1cm,height=2.1cm,valign=center,colback=cgra!6,colframe=cgra,arc=2.5mm,boxrule=0.9pt,left=3mm]{\Large\bfseries\color{cgra}d)}\quad{\Large $30~\text{km}$}\end{tcolorbox}
\end{tcolorbox}
\newpage
\begin{tcolorbox}[enhanced,width=16cm,height=14cm,colback=white,colframe=cgra,boxrule=2pt,arc=0mm,
left=6mm,right=6mm,top=0mm,bottom=4mm,sharp corners,
overlay={\fill[cgra] (frame.north west) rectangle ([yshift=-1.25cm]frame.north east);
\node[anchor=west,text=white,font=\sffamily\bfseries\LARGE] at ([xshift=5mm,yshift=-0.62cm]frame.north west) {Question 84 D};
\node[anchor=east,text=white,font=\sffamily\itshape\large] at ([xshift=-5mm,yshift=-0.62cm]frame.north east) {Grandeurs et mesures};}]
\vspace*{1.75cm}
\begin{minipage}[t][6.9cm][c]{\linewidth}\LARGE\centering
\begin{minipage}{0.95\linewidth}\raggedright\linespread{1.25}\selectfont Une voiture roule à la vitesse constante de $72$ km/h pendant $35$ min. Quelle distance parcourt-elle~?\end{minipage}
\end{minipage}

\vfill
\begin{tcolorbox}[enhanced,nobeforeafter,width=7.1cm,height=2.1cm,valign=center,colback=cgra!6,colframe=cgra,arc=2.5mm,boxrule=0.9pt,left=3mm]{\Large\bfseries\color{cgra}a)}\quad{\Large $42~\text{km}$}\end{tcolorbox}\hfill\begin{tcolorbox}[enhanced,nobeforeafter,width=7.1cm,height=2.1cm,valign=center,colback=cgra!6,colframe=cgra,arc=2.5mm,boxrule=0.9pt,left=3mm]{\Large\bfseries\color{cgra}b)}\quad{\Large $2\,520~\text{km}$}\end{tcolorbox}\par\vspace{2.5mm}
\begin{tcolorbox}[enhanced,nobeforeafter,width=7.1cm,height=2.1cm,valign=center,colback=cgra!6,colframe=cgra,arc=2.5mm,boxrule=0.9pt,left=3mm]{\Large\bfseries\color{cgra}c)}\quad{\Large $84~\text{km}$}\end{tcolorbox}\hfill\begin{tcolorbox}[enhanced,nobeforeafter,width=7.1cm,height=2.1cm,valign=center,colback=cgra!6,colframe=cgra,arc=2.5mm,boxrule=0.9pt,left=3mm]{\Large\bfseries\color{cgra}d)}\quad{\Large $25{,}2~\text{km}$}\end{tcolorbox}
\end{tcolorbox}
\newpage
\begin{tcolorbox}[enhanced,width=16cm,height=14cm,colback=white,colframe=cgra,boxrule=2pt,arc=0mm,
left=6mm,right=6mm,top=0mm,bottom=4mm,sharp corners,
overlay={\fill[cgra] (frame.north west) rectangle ([yshift=-1.25cm]frame.north east);
\node[anchor=west,text=white,font=\sffamily\bfseries\LARGE] at ([xshift=5mm,yshift=-0.62cm]frame.north west) {Question 85 D};
\node[anchor=east,text=white,font=\sffamily\itshape\large] at ([xshift=-5mm,yshift=-0.62cm]frame.north east) {Grandeurs et mesures};}]
\vspace*{1.75cm}
\begin{minipage}[t][6.9cm][c]{\linewidth}\LARGE\centering
\begin{minipage}{0.95\linewidth}\raggedright\linespread{1.25}\selectfont Un train parcourt $400$ km à la vitesse moyenne de $120$ km/h. Combien de temps dure le trajet~?\end{minipage}
\end{minipage}

\vfill
\begin{tcolorbox}[enhanced,nobeforeafter,width=7.1cm,height=2.1cm,valign=center,colback=cgra!6,colframe=cgra,arc=2.5mm,boxrule=0.9pt,left=3mm]{\Large\bfseries\color{cgra}a)}\quad{\Large 3 h 20 min}\end{tcolorbox}\hfill\begin{tcolorbox}[enhanced,nobeforeafter,width=7.1cm,height=2.1cm,valign=center,colback=cgra!6,colframe=cgra,arc=2.5mm,boxrule=0.9pt,left=3mm]{\Large\bfseries\color{cgra}b)}\quad{\Large 3 h 33 min}\end{tcolorbox}\par\vspace{2.5mm}
\begin{tcolorbox}[enhanced,nobeforeafter,width=7.1cm,height=2.1cm,valign=center,colback=cgra!6,colframe=cgra,arc=2.5mm,boxrule=0.9pt,left=3mm]{\Large\bfseries\color{cgra}c)}\quad{\Large 3 h 50 min}\end{tcolorbox}\hfill\begin{tcolorbox}[enhanced,nobeforeafter,width=7.1cm,height=2.1cm,valign=center,colback=cgra!6,colframe=cgra,arc=2.5mm,boxrule=0.9pt,left=3mm]{\Large\bfseries\color{cgra}d)}\quad{\Large 3 h 5 min}\end{tcolorbox}
\end{tcolorbox}
\newpage
\begin{tcolorbox}[enhanced,width=16cm,height=14cm,colback=white,colframe=cgra,boxrule=2pt,arc=0mm,
left=6mm,right=6mm,top=0mm,bottom=4mm,sharp corners,
overlay={\fill[cgra] (frame.north west) rectangle ([yshift=-1.25cm]frame.north east);
\node[anchor=west,text=white,font=\sffamily\bfseries\LARGE] at ([xshift=5mm,yshift=-0.62cm]frame.north west) {Question 86 D};
\node[anchor=east,text=white,font=\sffamily\itshape\large] at ([xshift=-5mm,yshift=-0.62cm]frame.north east) {Grandeurs et mesures};}]
\vspace*{1.75cm}
\begin{minipage}[t][6.9cm][c]{\linewidth}\LARGE\centering
\begin{minipage}{0.95\linewidth}\raggedright\linespread{1.25}\selectfont Un train parcourt $70$ km à la vitesse moyenne de $105$ km/h. Combien de temps dure le trajet~?\end{minipage}
\end{minipage}

\vfill
\begin{tcolorbox}[enhanced,nobeforeafter,width=7.1cm,height=2.1cm,valign=center,colback=cgra!6,colframe=cgra,arc=2.5mm,boxrule=0.9pt,left=3mm]{\Large\bfseries\color{cgra}a)}\quad{\Large 25 min}\end{tcolorbox}\hfill\begin{tcolorbox}[enhanced,nobeforeafter,width=7.1cm,height=2.1cm,valign=center,colback=cgra!6,colframe=cgra,arc=2.5mm,boxrule=0.9pt,left=3mm]{\Large\bfseries\color{cgra}b)}\quad{\Large 1 h 10 min}\end{tcolorbox}\par\vspace{2.5mm}
\begin{tcolorbox}[enhanced,nobeforeafter,width=7.1cm,height=2.1cm,valign=center,colback=cgra!6,colframe=cgra,arc=2.5mm,boxrule=0.9pt,left=3mm]{\Large\bfseries\color{cgra}c)}\quad{\Large 40 min}\end{tcolorbox}\hfill\begin{tcolorbox}[enhanced,nobeforeafter,width=7.1cm,height=2.1cm,valign=center,colback=cgra!6,colframe=cgra,arc=2.5mm,boxrule=0.9pt,left=3mm]{\Large\bfseries\color{cgra}d)}\quad{\Large 1 h 7 min}\end{tcolorbox}
\end{tcolorbox}
\newpage
\begin{tcolorbox}[enhanced,width=16cm,height=14cm,colback=white,colframe=cgra,boxrule=2pt,arc=0mm,
left=6mm,right=6mm,top=0mm,bottom=4mm,sharp corners,
overlay={\fill[cgra] (frame.north west) rectangle ([yshift=-1.25cm]frame.north east);
\node[anchor=west,text=white,font=\sffamily\bfseries\LARGE] at ([xshift=5mm,yshift=-0.62cm]frame.north west) {Question 87 D};
\node[anchor=east,text=white,font=\sffamily\itshape\large] at ([xshift=-5mm,yshift=-0.62cm]frame.north east) {Grandeurs et mesures};}]
\vspace*{1.75cm}
\begin{minipage}[t][6.9cm][c]{\linewidth}\LARGE\centering
\begin{minipage}{0.95\linewidth}\raggedright\linespread{1.25}\selectfont Un véhicule roule à $45$ km/h. Quelle est sa vitesse en mètres par seconde~?\end{minipage}
\end{minipage}

\vfill
\begin{tcolorbox}[enhanced,nobeforeafter,width=7.1cm,height=2.1cm,valign=center,colback=cgra!6,colframe=cgra,arc=2.5mm,boxrule=0.9pt,left=3mm]{\Large\bfseries\color{cgra}a)}\quad{\Large $0{,}75~\text{m/s}$}\end{tcolorbox}\hfill\begin{tcolorbox}[enhanced,nobeforeafter,width=7.1cm,height=2.1cm,valign=center,colback=cgra!6,colframe=cgra,arc=2.5mm,boxrule=0.9pt,left=3mm]{\Large\bfseries\color{cgra}b)}\quad{\Large $162~\text{m/s}$}\end{tcolorbox}\par\vspace{2.5mm}
\begin{tcolorbox}[enhanced,nobeforeafter,width=7.1cm,height=2.1cm,valign=center,colback=cgra!6,colframe=cgra,arc=2.5mm,boxrule=0.9pt,left=3mm]{\Large\bfseries\color{cgra}c)}\quad{\Large $45\,000~\text{m/s}$}\end{tcolorbox}\hfill\begin{tcolorbox}[enhanced,nobeforeafter,width=7.1cm,height=2.1cm,valign=center,colback=cgra!6,colframe=cgra,arc=2.5mm,boxrule=0.9pt,left=3mm]{\Large\bfseries\color{cgra}d)}\quad{\Large $12{,}5~\text{m/s}$}\end{tcolorbox}
\end{tcolorbox}
\newpage
\begin{tcolorbox}[enhanced,width=16cm,height=14cm,colback=white,colframe=cgra,boxrule=2pt,arc=0mm,
left=6mm,right=6mm,top=0mm,bottom=4mm,sharp corners,
overlay={\fill[cgra] (frame.north west) rectangle ([yshift=-1.25cm]frame.north east);
\node[anchor=west,text=white,font=\sffamily\bfseries\LARGE] at ([xshift=5mm,yshift=-0.62cm]frame.north west) {Question 88 D};
\node[anchor=east,text=white,font=\sffamily\itshape\large] at ([xshift=-5mm,yshift=-0.62cm]frame.north east) {Grandeurs et mesures};}]
\vspace*{1.75cm}
\begin{minipage}[t][6.9cm][c]{\linewidth}\LARGE\centering
\begin{minipage}{0.95\linewidth}\raggedright\linespread{1.25}\selectfont Un véhicule roule à $126$ km/h. Quelle est sa vitesse en mètres par seconde~?\end{minipage}
\end{minipage}

\vfill
\begin{tcolorbox}[enhanced,nobeforeafter,width=7.1cm,height=2.1cm,valign=center,colback=cgra!6,colframe=cgra,arc=2.5mm,boxrule=0.9pt,left=3mm]{\Large\bfseries\color{cgra}a)}\quad{\Large $126\,000~\text{m/s}$}\end{tcolorbox}\hfill\begin{tcolorbox}[enhanced,nobeforeafter,width=7.1cm,height=2.1cm,valign=center,colback=cgra!6,colframe=cgra,arc=2.5mm,boxrule=0.9pt,left=3mm]{\Large\bfseries\color{cgra}b)}\quad{\Large $2{,}1~\text{m/s}$}\end{tcolorbox}\par\vspace{2.5mm}
\begin{tcolorbox}[enhanced,nobeforeafter,width=7.1cm,height=2.1cm,valign=center,colback=cgra!6,colframe=cgra,arc=2.5mm,boxrule=0.9pt,left=3mm]{\Large\bfseries\color{cgra}c)}\quad{\Large $35~\text{m/s}$}\end{tcolorbox}\hfill\begin{tcolorbox}[enhanced,nobeforeafter,width=7.1cm,height=2.1cm,valign=center,colback=cgra!6,colframe=cgra,arc=2.5mm,boxrule=0.9pt,left=3mm]{\Large\bfseries\color{cgra}d)}\quad{\Large $453{,}6~\text{m/s}$}\end{tcolorbox}
\end{tcolorbox}
\newpage
\begin{tcolorbox}[enhanced,width=16cm,height=14cm,colback=white,colframe=cgra,boxrule=2pt,arc=0mm,
left=6mm,right=6mm,top=0mm,bottom=4mm,sharp corners,
overlay={\fill[cgra] (frame.north west) rectangle ([yshift=-1.25cm]frame.north east);
\node[anchor=west,text=white,font=\sffamily\bfseries\LARGE] at ([xshift=5mm,yshift=-0.62cm]frame.north west) {Question 89 D};
\node[anchor=east,text=white,font=\sffamily\itshape\large] at ([xshift=-5mm,yshift=-0.62cm]frame.north east) {Grandeurs et mesures};}]
\vspace*{1.75cm}
\begin{minipage}[t][6.9cm][c]{\linewidth}\LARGE\centering
\begin{minipage}{0.95\linewidth}\raggedright\linespread{1.25}\selectfont Une pyramide a une base d'aire $36$ cm$^2$ et une hauteur de $2$ cm. Quel est son volume~?\end{minipage}
\end{minipage}

\vfill
\begin{tcolorbox}[enhanced,nobeforeafter,width=7.1cm,height=2.1cm,valign=center,colback=cgra!6,colframe=cgra,arc=2.5mm,boxrule=0.9pt,left=3mm]{\Large\bfseries\color{cgra}a)}\quad{\Large $38~\text{cm}^3$}\end{tcolorbox}\hfill\begin{tcolorbox}[enhanced,nobeforeafter,width=7.1cm,height=2.1cm,valign=center,colback=cgra!6,colframe=cgra,arc=2.5mm,boxrule=0.9pt,left=3mm]{\Large\bfseries\color{cgra}b)}\quad{\Large $72~\text{cm}^3$}\end{tcolorbox}\par\vspace{2.5mm}
\begin{tcolorbox}[enhanced,nobeforeafter,width=7.1cm,height=2.1cm,valign=center,colback=cgra!6,colframe=cgra,arc=2.5mm,boxrule=0.9pt,left=3mm]{\Large\bfseries\color{cgra}c)}\quad{\Large $24~\text{cm}^3$}\end{tcolorbox}\hfill\begin{tcolorbox}[enhanced,nobeforeafter,width=7.1cm,height=2.1cm,valign=center,colback=cgra!6,colframe=cgra,arc=2.5mm,boxrule=0.9pt,left=3mm]{\Large\bfseries\color{cgra}d)}\quad{\Large $36~\text{cm}^3$}\end{tcolorbox}
\end{tcolorbox}
\newpage
\begin{tcolorbox}[enhanced,width=16cm,height=14cm,colback=white,colframe=cgra,boxrule=2pt,arc=0mm,
left=6mm,right=6mm,top=0mm,bottom=4mm,sharp corners,
overlay={\fill[cgra] (frame.north west) rectangle ([yshift=-1.25cm]frame.north east);
\node[anchor=west,text=white,font=\sffamily\bfseries\LARGE] at ([xshift=5mm,yshift=-0.62cm]frame.north west) {Question 90 D};
\node[anchor=east,text=white,font=\sffamily\itshape\large] at ([xshift=-5mm,yshift=-0.62cm]frame.north east) {Grandeurs et mesures};}]
\vspace*{1.75cm}
\begin{minipage}[t][6.9cm][c]{\linewidth}\LARGE\centering
\begin{minipage}{0.95\linewidth}\raggedright\linespread{1.25}\selectfont Une pyramide a une base d'aire $21$ cm$^2$ et une hauteur de $10$ cm. Quel est son volume~?\end{minipage}
\end{minipage}

\vfill
\begin{tcolorbox}[enhanced,nobeforeafter,width=7.1cm,height=2.1cm,valign=center,colback=cgra!6,colframe=cgra,arc=2.5mm,boxrule=0.9pt,left=3mm]{\Large\bfseries\color{cgra}a)}\quad{\Large $210~\text{cm}^3$}\end{tcolorbox}\hfill\begin{tcolorbox}[enhanced,nobeforeafter,width=7.1cm,height=2.1cm,valign=center,colback=cgra!6,colframe=cgra,arc=2.5mm,boxrule=0.9pt,left=3mm]{\Large\bfseries\color{cgra}b)}\quad{\Large $70~\text{cm}^3$}\end{tcolorbox}\par\vspace{2.5mm}
\begin{tcolorbox}[enhanced,nobeforeafter,width=7.1cm,height=2.1cm,valign=center,colback=cgra!6,colframe=cgra,arc=2.5mm,boxrule=0.9pt,left=3mm]{\Large\bfseries\color{cgra}c)}\quad{\Large $31~\text{cm}^3$}\end{tcolorbox}\hfill\begin{tcolorbox}[enhanced,nobeforeafter,width=7.1cm,height=2.1cm,valign=center,colback=cgra!6,colframe=cgra,arc=2.5mm,boxrule=0.9pt,left=3mm]{\Large\bfseries\color{cgra}d)}\quad{\Large $105~\text{cm}^3$}\end{tcolorbox}
\end{tcolorbox}
\newpage
\begin{tcolorbox}[enhanced,width=16cm,height=14cm,colback=white,colframe=cgra,boxrule=2pt,arc=0mm,
left=6mm,right=6mm,top=0mm,bottom=4mm,sharp corners,
overlay={\fill[cgra] (frame.north west) rectangle ([yshift=-1.25cm]frame.north east);
\node[anchor=west,text=white,font=\sffamily\bfseries\LARGE] at ([xshift=5mm,yshift=-0.62cm]frame.north west) {Question 91 D};
\node[anchor=east,text=white,font=\sffamily\itshape\large] at ([xshift=-5mm,yshift=-0.62cm]frame.north east) {Grandeurs et mesures};}]
\vspace*{1.75cm}
\begin{minipage}[t][6.9cm][c]{\linewidth}\LARGE\centering
\begin{minipage}{0.95\linewidth}\raggedright\linespread{1.25}\selectfont Un cône de révolution a un rayon de base de $1$ cm et une hauteur de $12$ cm. Quel est son volume exact~?\end{minipage}
\end{minipage}

\vfill
\begin{tcolorbox}[enhanced,nobeforeafter,width=7.1cm,height=2.1cm,valign=center,colback=cgra!6,colframe=cgra,arc=2.5mm,boxrule=0.9pt,left=3mm]{\Large\bfseries\color{cgra}a)}\quad{\Large $24\pi~\text{cm}^3$}\end{tcolorbox}\hfill\begin{tcolorbox}[enhanced,nobeforeafter,width=7.1cm,height=2.1cm,valign=center,colback=cgra!6,colframe=cgra,arc=2.5mm,boxrule=0.9pt,left=3mm]{\Large\bfseries\color{cgra}b)}\quad{\Large $4\pi~\text{cm}^3$}\end{tcolorbox}\par\vspace{2.5mm}
\begin{tcolorbox}[enhanced,nobeforeafter,width=7.1cm,height=2.1cm,valign=center,colback=cgra!6,colframe=cgra,arc=2.5mm,boxrule=0.9pt,left=3mm]{\Large\bfseries\color{cgra}c)}\quad{\Large $12\pi~\text{cm}^3$}\end{tcolorbox}\hfill\begin{tcolorbox}[enhanced,nobeforeafter,width=7.1cm,height=2.1cm,valign=center,colback=cgra!6,colframe=cgra,arc=2.5mm,boxrule=0.9pt,left=3mm]{\Large\bfseries\color{cgra}d)}\quad{\Large $16\pi~\text{cm}^3$}\end{tcolorbox}
\end{tcolorbox}
\newpage
\begin{tcolorbox}[enhanced,width=16cm,height=14cm,colback=white,colframe=cgra,boxrule=2pt,arc=0mm,
left=6mm,right=6mm,top=0mm,bottom=4mm,sharp corners,
overlay={\fill[cgra] (frame.north west) rectangle ([yshift=-1.25cm]frame.north east);
\node[anchor=west,text=white,font=\sffamily\bfseries\LARGE] at ([xshift=5mm,yshift=-0.62cm]frame.north west) {Question 92 D};
\node[anchor=east,text=white,font=\sffamily\itshape\large] at ([xshift=-5mm,yshift=-0.62cm]frame.north east) {Grandeurs et mesures};}]
\vspace*{1.75cm}
\begin{minipage}[t][6.9cm][c]{\linewidth}\LARGE\centering
\begin{minipage}{0.95\linewidth}\raggedright\linespread{1.25}\selectfont Un cône de révolution a un rayon de base de $6$ cm et une hauteur de $5$ cm. Quel est son volume exact~?\end{minipage}
\end{minipage}

\vfill
\begin{tcolorbox}[enhanced,nobeforeafter,width=7.1cm,height=2.1cm,valign=center,colback=cgra!6,colframe=cgra,arc=2.5mm,boxrule=0.9pt,left=3mm]{\Large\bfseries\color{cgra}a)}\quad{\Large $10\pi~\text{cm}^3$}\end{tcolorbox}\hfill\begin{tcolorbox}[enhanced,nobeforeafter,width=7.1cm,height=2.1cm,valign=center,colback=cgra!6,colframe=cgra,arc=2.5mm,boxrule=0.9pt,left=3mm]{\Large\bfseries\color{cgra}b)}\quad{\Large $60\pi~\text{cm}^3$}\end{tcolorbox}\par\vspace{2.5mm}
\begin{tcolorbox}[enhanced,nobeforeafter,width=7.1cm,height=2.1cm,valign=center,colback=cgra!6,colframe=cgra,arc=2.5mm,boxrule=0.9pt,left=3mm]{\Large\bfseries\color{cgra}c)}\quad{\Large $180\pi~\text{cm}^3$}\end{tcolorbox}\hfill\begin{tcolorbox}[enhanced,nobeforeafter,width=7.1cm,height=2.1cm,valign=center,colback=cgra!6,colframe=cgra,arc=2.5mm,boxrule=0.9pt,left=3mm]{\Large\bfseries\color{cgra}d)}\quad{\Large $240\pi~\text{cm}^3$}\end{tcolorbox}
\end{tcolorbox}
\newpage
\begin{tcolorbox}[enhanced,width=16cm,height=14cm,colback=white,colframe=cgra,boxrule=2pt,arc=0mm,
left=6mm,right=6mm,top=0mm,bottom=4mm,sharp corners,
overlay={\fill[cgra] (frame.north west) rectangle ([yshift=-1.25cm]frame.north east);
\node[anchor=west,text=white,font=\sffamily\bfseries\LARGE] at ([xshift=5mm,yshift=-0.62cm]frame.north west) {Question 93 D};
\node[anchor=east,text=white,font=\sffamily\itshape\large] at ([xshift=-5mm,yshift=-0.62cm]frame.north east) {Grandeurs et mesures};}]
\vspace*{1.75cm}
\begin{minipage}[t][6.9cm][c]{\linewidth}\LARGE\centering
\begin{minipage}{0.95\linewidth}\raggedright\linespread{1.25}\selectfont Compléter~: $12~\text{L}=\;?\;\text{dm}^3$\end{minipage}
\end{minipage}

\vfill
\begin{tcolorbox}[enhanced,nobeforeafter,width=7.1cm,height=2.1cm,valign=center,colback=cgra!6,colframe=cgra,arc=2.5mm,boxrule=0.9pt,left=3mm]{\Large\bfseries\color{cgra}a)}\quad{\Large $12~\text{dm}^3$}\end{tcolorbox}\hfill\begin{tcolorbox}[enhanced,nobeforeafter,width=7.1cm,height=2.1cm,valign=center,colback=cgra!6,colframe=cgra,arc=2.5mm,boxrule=0.9pt,left=3mm]{\Large\bfseries\color{cgra}b)}\quad{\Large $1\,200~\text{dm}^3$}\end{tcolorbox}\par\vspace{2.5mm}
\begin{tcolorbox}[enhanced,nobeforeafter,width=7.1cm,height=2.1cm,valign=center,colback=cgra!6,colframe=cgra,arc=2.5mm,boxrule=0.9pt,left=3mm]{\Large\bfseries\color{cgra}c)}\quad{\Large $0{,}012~\text{dm}^3$}\end{tcolorbox}\hfill\begin{tcolorbox}[enhanced,nobeforeafter,width=7.1cm,height=2.1cm,valign=center,colback=cgra!6,colframe=cgra,arc=2.5mm,boxrule=0.9pt,left=3mm]{\Large\bfseries\color{cgra}d)}\quad{\Large $12\,000~\text{dm}^3$}\end{tcolorbox}
\end{tcolorbox}
\newpage
\begin{tcolorbox}[enhanced,width=16cm,height=14cm,colback=white,colframe=cgra,boxrule=2pt,arc=0mm,
left=6mm,right=6mm,top=0mm,bottom=4mm,sharp corners,
overlay={\fill[cgra] (frame.north west) rectangle ([yshift=-1.25cm]frame.north east);
\node[anchor=west,text=white,font=\sffamily\bfseries\LARGE] at ([xshift=5mm,yshift=-0.62cm]frame.north west) {Question 94 D};
\node[anchor=east,text=white,font=\sffamily\itshape\large] at ([xshift=-5mm,yshift=-0.62cm]frame.north east) {Grandeurs et mesures};}]
\vspace*{1.75cm}
\begin{minipage}[t][6.9cm][c]{\linewidth}\LARGE\centering
\begin{minipage}{0.95\linewidth}\raggedright\linespread{1.25}\selectfont Compléter~: $5\,000~\text{cm}^3=\;?\;\text{L}$\end{minipage}
\end{minipage}

\vfill
\begin{tcolorbox}[enhanced,nobeforeafter,width=7.1cm,height=2.1cm,valign=center,colback=cgra!6,colframe=cgra,arc=2.5mm,boxrule=0.9pt,left=3mm]{\Large\bfseries\color{cgra}a)}\quad{\Large $5~\text{L}$}\end{tcolorbox}\hfill\begin{tcolorbox}[enhanced,nobeforeafter,width=7.1cm,height=2.1cm,valign=center,colback=cgra!6,colframe=cgra,arc=2.5mm,boxrule=0.9pt,left=3mm]{\Large\bfseries\color{cgra}b)}\quad{\Large $0{,}5~\text{L}$}\end{tcolorbox}\par\vspace{2.5mm}
\begin{tcolorbox}[enhanced,nobeforeafter,width=7.1cm,height=2.1cm,valign=center,colback=cgra!6,colframe=cgra,arc=2.5mm,boxrule=0.9pt,left=3mm]{\Large\bfseries\color{cgra}c)}\quad{\Large $50~\text{L}$}\end{tcolorbox}\hfill\begin{tcolorbox}[enhanced,nobeforeafter,width=7.1cm,height=2.1cm,valign=center,colback=cgra!6,colframe=cgra,arc=2.5mm,boxrule=0.9pt,left=3mm]{\Large\bfseries\color{cgra}d)}\quad{\Large $500~\text{L}$}\end{tcolorbox}
\end{tcolorbox}
\newpage
\begin{tcolorbox}[enhanced,width=16cm,height=14cm,colback=white,colframe=cgra,boxrule=2pt,arc=0mm,
left=6mm,right=6mm,top=0mm,bottom=4mm,sharp corners,
overlay={\fill[cgra] (frame.north west) rectangle ([yshift=-1.25cm]frame.north east);
\node[anchor=west,text=white,font=\sffamily\bfseries\LARGE] at ([xshift=5mm,yshift=-0.62cm]frame.north west) {Question 95 D};
\node[anchor=east,text=white,font=\sffamily\itshape\large] at ([xshift=-5mm,yshift=-0.62cm]frame.north east) {Grandeurs et mesures};}]
\vspace*{1.75cm}
\begin{minipage}[t][6.9cm][c]{\linewidth}\LARGE\centering
\begin{minipage}{0.95\linewidth}\raggedright\linespread{1.25}\selectfont Compléter~: $9~\text{m}^2=\;?\;\text{dm}^2$\end{minipage}
\end{minipage}

\vfill
\begin{tcolorbox}[enhanced,nobeforeafter,width=7.1cm,height=2.1cm,valign=center,colback=cgra!6,colframe=cgra,arc=2.5mm,boxrule=0.9pt,left=3mm]{\Large\bfseries\color{cgra}a)}\quad{\Large $900~\text{dm}^2$}\end{tcolorbox}\hfill\begin{tcolorbox}[enhanced,nobeforeafter,width=7.1cm,height=2.1cm,valign=center,colback=cgra!6,colframe=cgra,arc=2.5mm,boxrule=0.9pt,left=3mm]{\Large\bfseries\color{cgra}b)}\quad{\Large $9\,000~\text{dm}^2$}\end{tcolorbox}\par\vspace{2.5mm}
\begin{tcolorbox}[enhanced,nobeforeafter,width=7.1cm,height=2.1cm,valign=center,colback=cgra!6,colframe=cgra,arc=2.5mm,boxrule=0.9pt,left=3mm]{\Large\bfseries\color{cgra}c)}\quad{\Large $90\,000~\text{dm}^2$}\end{tcolorbox}\hfill\begin{tcolorbox}[enhanced,nobeforeafter,width=7.1cm,height=2.1cm,valign=center,colback=cgra!6,colframe=cgra,arc=2.5mm,boxrule=0.9pt,left=3mm]{\Large\bfseries\color{cgra}d)}\quad{\Large $90~\text{dm}^2$}\end{tcolorbox}
\end{tcolorbox}
\newpage
\begin{tcolorbox}[enhanced,width=16cm,height=14cm,colback=white,colframe=cgra,boxrule=2pt,arc=0mm,
left=6mm,right=6mm,top=0mm,bottom=4mm,sharp corners,
overlay={\fill[cgra] (frame.north west) rectangle ([yshift=-1.25cm]frame.north east);
\node[anchor=west,text=white,font=\sffamily\bfseries\LARGE] at ([xshift=5mm,yshift=-0.62cm]frame.north west) {Question 96 D};
\node[anchor=east,text=white,font=\sffamily\itshape\large] at ([xshift=-5mm,yshift=-0.62cm]frame.north east) {Grandeurs et mesures};}]
\vspace*{1.75cm}
\begin{minipage}[t][6.9cm][c]{\linewidth}\LARGE\centering
\begin{minipage}{0.95\linewidth}\raggedright\linespread{1.25}\selectfont Compléter~: $5~\text{cm}^2=\;?\;\text{mm}^2$\end{minipage}
\end{minipage}

\vfill
\begin{tcolorbox}[enhanced,nobeforeafter,width=7.1cm,height=2.1cm,valign=center,colback=cgra!6,colframe=cgra,arc=2.5mm,boxrule=0.9pt,left=3mm]{\Large\bfseries\color{cgra}a)}\quad{\Large $50~\text{mm}^2$}\end{tcolorbox}\hfill\begin{tcolorbox}[enhanced,nobeforeafter,width=7.1cm,height=2.1cm,valign=center,colback=cgra!6,colframe=cgra,arc=2.5mm,boxrule=0.9pt,left=3mm]{\Large\bfseries\color{cgra}b)}\quad{\Large $500~\text{mm}^2$}\end{tcolorbox}\par\vspace{2.5mm}
\begin{tcolorbox}[enhanced,nobeforeafter,width=7.1cm,height=2.1cm,valign=center,colback=cgra!6,colframe=cgra,arc=2.5mm,boxrule=0.9pt,left=3mm]{\Large\bfseries\color{cgra}c)}\quad{\Large $5\,000~\text{mm}^2$}\end{tcolorbox}\hfill\begin{tcolorbox}[enhanced,nobeforeafter,width=7.1cm,height=2.1cm,valign=center,colback=cgra!6,colframe=cgra,arc=2.5mm,boxrule=0.9pt,left=3mm]{\Large\bfseries\color{cgra}d)}\quad{\Large $50\,000~\text{mm}^2$}\end{tcolorbox}
\end{tcolorbox}
\newpage
\end{document}